\documentclass[11pt,a4paper]{article}
\usepackage[T1]{fontenc}
\usepackage[utf8]{inputenc}
\usepackage{lmodern}
\usepackage{amsmath,amssymb,amsthm,mathtools}
\usepackage[margin=27mm,headheight=14pt]{geometry}
\usepackage{microtype}
\usepackage{booktabs,longtable,array,multicol}
\usepackage{graphicx,xcolor,fancyhdr,fancyvrb,listings}
\usepackage{tikz}
\usetikzlibrary{arrows.meta,positioning}
\usepackage{xr-hyper}
\usepackage[unicode,colorlinks=true,linkcolor=blue!45!black,citecolor=blue!45!black,urlcolor=blue!45!black,filecolor=blue!45!black,bookmarksnumbered=true]{hyperref}
\hypersetup{pdfauthor={},pdfcreator={LaTeX},pdfsubject={Graceful labelings of spiders}}
\setlength{\parindent}{1.2em}
\setlength{\parskip}{0.18em}
\setlength{\emergencystretch}{2em}
\raggedbottom
\setcounter{tocdepth}{2}
\setcounter{secnumdepth}{3}
\numberwithin{equation}{section}
\theoremstyle{plain}
\newtheorem{theorem}{Theorem}[section]
\newtheorem{lemma}[theorem]{Lemma}
\newtheorem{proposition}[theorem]{Proposition}
\newtheorem{corollary}[theorem]{Corollary}
\theoremstyle{definition}
\newtheorem{definition}[theorem]{Definition}
\newtheorem{algorithm}[theorem]{Verification algorithm}
\newtheorem{example}[theorem]{Example}
\theoremstyle{remark}
\newtheorem{remark}[theorem]{Remark}
\newcommand{\Sp}[1]{S\!\left(#1\right)}
\newcommand{\Om}{\Omega}
\newcommand{\A}{\mathsf A}
\newcommand{\Hh}{\mathsf H}
\newcommand{\NN}{\mathbb N}
\newcommand{\ZZ}{\mathbb Z}
\newcommand{\QQ}{\mathbb Q}
\newcommand{\floor}[1]{\left\lfloor #1\right\rfloor}
\newcommand{\ceil}[1]{\left\lceil #1\right\rceil}
\newcommand{\set}[1]{\left\{#1\right\}}
\newcommand{\abs}[1]{\left|#1\right|}
\newcommand{\interval}[2]{[#1,#2]_{\ZZ}}
\newcommand{\file}[1]{{\small\path{#1}}}
\newcommand{\code}[1]{\texttt{\detokenize{#1}}}
\providecommand{\cref}[1]{\autoref{#1}}
\providecommand{\Cref}[1]{\autoref{#1}}
\newcommand{\sourcefile}[1]{\par\smallskip\noindent\textit{Computational data:} \file{#1}.\par\smallskip}
\newcommand{\publicationnote}[1]{\par\smallskip\noindent\textit{#1}\par\smallskip}
% GitHub repository address, used once (footnote in Section 1).
\newcommand{\githuburl}{https://github.com/mihai906/spider-gracefulness-supplement}

\newcommand{\providesupplementreference}[3]{%
  \expandafter\def\csname suppnumber@#1\endcsname{#2}%
  \expandafter\def\csname supppage@#1\endcsname{#3}}
\newcommand{\suppref}[2][Section]{#1~\csname suppnumber@#2\endcsname}
\newcommand{\As}{\mathsf A_7}
\newcommand{\Hs}{\mathsf H_7}
\newcommand{\Ae}{\mathsf A_8^{\epsilon}}
\newcommand{\He}{\mathsf H_8^{\epsilon}}
\newcommand{\Oe}{\mathsf O_8^{\epsilon}}
\pagestyle{fancy}
\fancyhf{}
\fancyhead[L]{\small Graceful labelings of spiders}
\fancyfoot[C]{\thepage}
\renewcommand{\headrulewidth}{0.3pt}
\lstset{basicstyle=\ttfamily\footnotesize,breaklines=true,columns=fullflexible,
 keepspaces=true,showstringspaces=false,frame=single,framerule=0.2pt,
 rulecolor=\color{black!20},aboveskip=0.7em,belowskip=0.7em}
\lstnewenvironment{Code}{\lstset{language={}}}{}

\newif\ifSevenVolume\SevenVolumetrue
\makeatletter
\renewcommand{\@pnumwidth}{2.3em}
\renewcommand{\@tocrmarg}{3em}
\renewcommand*\l@section[2]{\addvspace{0.4em}\@dottedtocline{1}{0em}{2.5em}{\bfseries #1}{\bfseries #2}}
\renewcommand*\l@subsection{\@dottedtocline{2}{2.5em}{4em}}
\makeatother
\renewcommand{\thesection}{S\arabic{section}}
\externaldocument[main-]{spider_gracefulness_v2}
\fancyhead[L]{\small Graceful labelings of spiders: technical supplement}
\title{Graceful Labelings of Seven-Arm Spiders\\and Two Eight-Arm Families\\[6pt]\large Technical supplement}
\author{Bogdan Dumitru \and Mihai Nacu}
\hypersetup{pdfauthor={Bogdan Dumitru; Mihai Nacu}}
% Author addresses, printed after the references.
\newcommand{\authoraddress}[2]{%
  \par\noindent #1\\
  Faculty of Mathematics and Computer Science, University of Bucharest, Romania\\
  \textit{Email address:} \href{mailto:#2}{\nolinkurl{#2}}\par\medskip}
\newcommand{\authoraddresses}{%
  \par\bigskip{\small
  \authoraddress{Bogdan Dumitru}{bogdan.dumitru@fmi.unibuc.ro}
  \authoraddress{Mihai Nacu}{mihainacu906@gmail.com}}}

\date{17 November 2026}
\hypersetup{pdftitle={Graceful Labelings of Seven-Arm Spiders and Two Eight-Arm Families: Technical Supplement}}
\begin{document}
\maketitle
\begin{abstract}
This supplement contains the detailed mathematical dependencies of
the accompanying article: rooted path constructions, lifts, offset
intervals, short-arm proofs, odd and mixed kernels, affine edge
trades, complete finite inventories, and exact arithmetic certificate
specifications. Results numbered S belong to this document. References
to the article are identified explicitly. Together the article and this supplement form the mathematical
manuscript; the supporting repository supplies the machine-readable
proof objects and their replay tools.
\end{abstract}

\subsection*{Notation and dependency guide}
Lengths count edges; $L^r$ means $r$ repeated arms. An arm word omits
its initial hub label. All label intervals are integer intervals,
and an interval with reversed endpoints is empty. The alpha cut of
a root-zero spider on arms indexed by $I$ is
$K(I)=\sum_{i\in I}\floor{L_i/2}$. The notation $\Om(T)$ denotes
the available rooted offsets. The predicates $\mathsf A_7,\mathsf H_7$
are defined in Article~\ref*{main-art:seven-predicates}. The eight-arm
predicates are defined in Section~\ref{eight:specification}; within
that section $\A,\mathsf O,\Hh$ abbreviate the corresponding
eight-arm predicates in its stated parity domain.

\begin{center}\small
\begin{tabular}{@{}p{.46\textwidth}p{.46\textwidth}@{}}
\toprule
Input used by the article & Location in this supplement\\
\midrule
Classical paths and compositions & Section~\ref{common:foundations}\\
Four arms with a short arm & Lemma~\ref{seven:four-short}\\
Small alpha triples & Section~\ref{app:small-alpha}\\
Ten alpha families & Sections~\ref{app:odd-kernels}--\ref{app:ten-alpha}\\
Five-to-four selection & Theorem~\ref{fs:selection}\\
All-parameter affine constructions & Section~\ref{app:template-certificates}\\
Complete finite words and splits & Sections~\ref{app:finite-bases}, \ref{app:auxiliary-finite}\\
Exact eight-arm rules & Section~\ref{eight:specification}\\
Arithmetic premise binding and replay & Section~\ref{app:seven-certificates}\\
\bottomrule
\end{tabular}
\end{center}
The construction proofs and complete inventories below provide the
inputs used by the article. The accompanying computational data and
replay tools are identified in Section~\ref{app:seven-certificates}.

\tableofcontents
\clearpage
\section{Rooted labeling interfaces and path inputs}
\label{common:foundations}

All graphs in this manuscript are finite. A path length is its number
of edges, while $P_n$ denotes a path on $n$ vertices. These conventions
are especially important when a path $P_L$ is joined to a hub by a new
edge: the new arm then has length $L$.

\begin{definition}
Let $T$ be a tree with $m$ edges. A \emph{graceful labeling} is a
bijection $f:V(T)\to\interval{0}{m}$ such that the edge differences
$\abs{f(u)-f(v)}$ are precisely $1,\ldots,m$, each once. It is an
\emph{$\alpha$-labeling of integer index $K$} if every edge has one endpoint
of label at most $K$ and the other of label greater than $K$,
where $0\le K\le m$. In particular, the singleton has index zero.
For a distinguished root $r$, \emph{root-zero} means $f(r)=0$.
\end{definition}

The rooted spider $\Sp{L_1,\ldots,L_d}$ is obtained by identifying
one endpoint of each of $d$ otherwise disjoint paths of positive edge
lengths $L_i$. Its distinguished root is the common hub. We also allow
$d\le2$, and allow $d=0$ for the one-vertex tree. Repeated lengths are
written $L^t$ when the multiplicity $t$ is specified.

\begin{definition}
For a rooted $m$-edge tree $T$, write $k\in\Om(T)$ if there is a
labeling with root zero, nonroot labels exactly
$\interval{k+1}{k+m}$, and differences exactly $\interval{1}{m}$.
Only integers $k\ge0$ are considered. A realization at $k=0$ is a
hub-zero graceful labeling when the root is a spider hub.
\end{definition}

\begin{lemma}[Forced index]\label{common:root-index}
A root-zero $\alpha$-labeling of a spider has index
\[
 K=\sum_{i=1}^d\floor{L_i/2}.
\]
\end{lemma}
\begin{proof}
The root belongs to the lower class. Its bipartition class consists of
the root and the vertices at positive even distance along the arms,
and hence has $1+\sum_i\floor{L_i/2}$ vertices. All labels from zero
through the index occur in this class, and every label occurs once.
Its size is therefore $K+1$.
\end{proof}

The classical graceful-tree problem asks for a graceful labeling of every tree. Spiders provide a natural family with one branching vertex. The published five-leg theorem is due to Panpa, Imnang and Wasuanankul~\cite{Panpa2025}; the six-leg rooted-offset approach appears in the 2026 preprint of Rhys~\cite{Rhys2026}. Related constructions appear in Bahls, Lake and Wertheim~\cite{BahlsLakeWertheim2010}, Shan and Zhong~\cite{ShanZhong2026}, and Dumitru and Nacu~\cite{DumitruNacu2026}. The present proofs use the classical path and amalgamation results below, and give the subsequent construction arguments explicitly. The six-leg conclusion itself is not a premise.\par\medskip

\subsection{Classical path results with precise normalizations}

We use the following established path theorem. The $\alpha$-statements
are due to Rosa and Cattell \cite{Rosa1977,Cattell2007}; the three
forms are restated as Theorems~2.1--2.3 of \cite{Panpa2025}.

\begin{theorem}[Prescribed labels on paths]\label{common:path-classical}
For $n\ge2$ the following hold.
\begin{enumerate}
\item Any specified vertex of $P_n$ can have label zero in an
$\alpha$-labeling, except the central vertex of $P_5$.
\item An endpoint of $P_n$ can have any specified graceful label in
$\interval{0}{n-1}$.
\item An endpoint can have any such label in an $\alpha$-labeling,
except when $n=4s+1$, $s\ge1$, and the label is $s$ or $3s$.
\end{enumerate}
The one-vertex path is treated separately with its unique label zero.
\end{theorem}

In construction formulas we often normalize a path to labels
$\interval{1}{n}$ instead. This adds one to every label and to the
cut, preserves differences, and changes the two exceptional labels
to $s+1$ and $3s+1$.

\begin{corollary}[Rooted pairs]\label{common:alpha-pair}
A rooted pair $\Sp{a,b}$ has a root-zero $\alpha$-labeling unless
$(a,b)=(2,2)$. Its index, when it exists, is
$\floor{a/2}+\floor{b/2}$. Every single arm is root-zero $\alpha$.
\end{corollary}
\begin{proof}
The pair is $P_{a+b+1}$ with the distinguished vertex at distances
$a,b$ from its endpoints. It is the exceptional rooted $P_5$ exactly
when $a=b=2$. A rooted single arm is covered by the same theorem, or
by the Walecki word $0,L,1,L-1,2,L-2,\ldots$.
\end{proof}

Every rooted pair is nevertheless ordinary hub-zero graceful. The only
case not supplied by the alpha statement is $\Sp{2,2}$, for which the
two outward words $(4,1)$ and $(2,3)$ give the colors $4,3,2,1$.

\begin{lemma}[Even-order path interface]\label{common:even-path}
For $r\ge1$ and $1\le v\le r$, there is an $\alpha$-labeled path
on $\interval{1}{2r}$, cut at $r$, starting at $r+v$ and ending at
$v$. There is no exceptional endpoint request.
\end{lemma}
\begin{proof}
Apply the classical endpoint theorem at $r+v$. The two classes of an
even-order path have equal size, so its normalized cut is $r$ and its
endpoints are in opposite classes. If their values are $E$ in the
upper class and $F$ in the lower class, summing positive differences
by vertex degrees gives
\[
 r(2r-1)=2\!\sum_{j=r+1}^{2r}j-2\!\sum_{j=1}^rj-E+F
             =2r^2-E+F.
\]
Thus $E-F=r$, which forces $F=v$.
\end{proof}

\begin{lemma}[Odd-order path interface]\label{common:odd-path}
For $u\ge2$ and $u\le X\le2u-1$, an $\alpha$-labeled path on
$\interval{1}{2u-1}$ with larger class upper, cut $u-1$, starts at
$X$ and ends at $Y=3u-1-X$, except when
\[
 2(X-1)=3(u-1).
\]
For $u=1$ the singleton word $(1)$ is allowed separately.
\end{lemma}
\begin{proof}
The two ends lie in the larger class. For $X>u$, an
$\alpha$-labeling supplied by the endpoint theorem must place that
class above the cut. If $X=u$, complementation by $2u-x$ fixes $X$
and interchanges the two classes if necessary. The forbidden upper
endpoint is exactly the displayed equality. The other endpoint is
forced by
\[
 (2u-1)(u-1)
 =2\!\sum_{j=u}^{2u-1}j-2\!\sum_{j=1}^{u-1}j-X-Y
 =2u^2-X-Y.
\]
This yields $X+Y=3u-1$.
\end{proof}

\begin{lemma}[Ordinary prescribed-first paths]\label{common:prescribed-first-path}
For every $L\ge1$ and $1\le j\le L$ there is a graceful
permutation of $\interval{1}{L}$ starting at $j$. Consequently a
single $L$-arm admits every offset $0\le k<L$.
\end{lemma}
\begin{proof}
The ordinary endpoint theorem supplies the permutation. For clarity,
even the two endpoint requests excluded only by the stronger
$\alpha$-theorem have explicit ordinary witnesses. When $L=4s+1$,
rotate the Walecki word
$(1,L,2,L-1,\ldots,2s+1)$ so that its entry at zero-based position
$2s+1$ is first. This entry is $3s+1$. The removed cut edge and the
added wrap edge both have difference $2s$, so every difference
$1,\ldots,L-1$ is retained. Complementation gives first label $s+1$.
The singleton case is immediate. To obtain an offset $k$, choose
first label $L-k$ and add $k$ to all $L$ labels. The new root edge
has color $L$, while the internal differences remain $1,\ldots,L-1$.
\end{proof}

\begin{theorem}[Three-arm hub-zero input]\label{common:hubzero-three}
\label{app:hubzero-three}
Every positive three-arm spider has an ordinary graceful labeling
with hub zero \cite[Theorem~3.2]{Panpa2025}. This input does not assert
a root-zero $\alpha$-labeling.
\end{theorem}

\subsection{Interval compositions}

\begin{lemma}[Alpha-offset amalgamation]\label{common:alpha-offset}
Let $G$ be root-zero $\alpha$ with $g$ edges and cut $K$, and let
$H$ have $h$ edges and an offset-$K$ realization. Their rooted
amalgamation is hub-zero graceful.
\end{lemma}
\begin{proof}
Keep lower labels of $G$, raise its upper labels by $h$, and keep
the offset labels of $H$. Identify roots at zero. The nonroot labels
occupy, without repetition,
\[
 \interval{1}{K},\quad\interval{K+1}{K+h},\quad
 \interval{K+h+1}{g+h}.
\]
The two palettes are $\interval{1}{h}$ and
$\interval{h+1}{h+g}$. This proves every required inventory.
\end{proof}

\begin{lemma}[Ordinary alpha amalgamation]\label{common:ordinary-amalgamation}
Let $G$ be root-zero $\alpha$ with $g$ edges and cut $K$, and let
$H$ be hub-zero graceful with $h$ edges. Their rooted amalgamation is
ordinary graceful, with final hub label $K$.
\end{lemma}
\begin{proof}
Reverse labels within the two classes of $G$ by
\[
 f^R(v)=\begin{cases}K-f(v),&f(v)\le K,\\
 g+K+1-f(v),&f(v)>K.\end{cases}
\]
Its root is $K$ and an old color $d$ becomes $g+1-d$. Raise its
uppers by $h$ and translate every label of $H$ by $K$. The common
root is $K$, and the remaining label blocks are
$\interval{0}{K-1}$, $\interval{K+1}{K+h}$ and
$\interval{K+h+1}{g+h}$. The colors are exactly
$\interval{1}{h}$ and $\interval{h+1}{g+h}$.
This is the interval form of the classical amalgamation theorem
\cite{HuangKotzigRosa1982}.
\end{proof}

\begin{lemma}[Unit stability]\label{common:unit-stability}
Arbitrarily many unit arms can be added to a hub-zero graceful
component, or to a root-zero $\alpha$ component while preserving its
index. An ordinary alpha/hub-zero decomposition therefore survives
any number of units added to the hub-zero piece before amalgamation.
\end{lemma}
\begin{proof}
For an $m$-edge component, label $v$ new leaves $m+1,\ldots,m+v$.
Their root edges have precisely these new colors. In the alpha case
each new leaf lies above the unchanged cut. The last assertion follows
by applying the preceding amalgamation to the enlarged hub-zero piece.
\end{proof}

\begin{corollary}[Ordinary hub-zero arm attachment]\label{common:ordinary-attachment}
A root-zero $\alpha$ piece of cut $K$ accepts an ordinary arm of
length $L>K$ to form a hub-zero graceful tree.
\end{corollary}
\begin{proof}
Use the single-arm offset $K$ and alpha-offset amalgamation.
\end{proof}

\begin{lemma}[Alpha-preserving arm attachment]\label{common:alpha-extension}
If a rooted tree has a root-zero $\alpha$-labeling of cut $K$, an
arm of length $L$ can be attached at the root, preserving root-zero
$\alpha$, whenever $L>2K$ and $L\ne4K+1$. The new cut is
$K+\floor{L/2}$.
\end{lemma}
\begin{proof}
Put $q=L-1$ and $\ell=\floor{(L-1)/2}$. On a separate $P_L$ choose
an $\alpha$-labeling on $\interval{0}{q}$ with lower endpoint
$j=\ell-K$. This is in the lower class because $L>2K$. The lower
exception from the endpoint theorem is exactly $L=4K+1$; the upper
exception is outside the requested lower half. Let $G$ have $g$ edges.
Keep its lower labels and raise its uppers by $L$. In the new path
map lower labels $y$ to $y+K+q-\ell+1$ and upper labels $y$ to
$y+K-\ell$. Its joining endpoint becomes $L$, and join it to the
old root zero. The four vertex blocks are
\[
 [0,K],\quad[K+1,K+q-\ell],\quad
 [K+q-\ell+1,K+q+1],\quad[K+L+1,g+L].
\]
Empty intervals are omitted. They partition $[0,g+L]$, and every
edge crosses $K+q-\ell=K+\floor{L/2}$. The new path colors are
$1,\ldots,L-1$, in reversed order; the joining color is $L$; and
the old tree supplies $L+1,\ldots,L+g$. This proves the claim.
\end{proof}

\subsection{Subdivision and the distinction between ordinary and alpha lifts}

\begin{theorem}[Uniform offset subdivision]\label{common:offset-subdivision}
If $s\in\Om(T)$ and every edge is replaced by a path of $q\ge1$
edges, then
\[
 \interval{qs}{q(s+1)-1}\ \subseteq\ \Om(T[q]).
\]
The complete column construction and signed-difference proof are given
in Section~\ref{app:subdivision}. The input need not be alpha.
\end{theorem}

In particular ordinary gracefulness is preserved under uniform
subdivision: root an ordinary labeling at its zero-labeled vertex and
take $s=0$. That vertex need not be the spider's hub. The underlying
graph operation multiplies all arm lengths by $q$, independently of
which vertex was used as the root during construction.

For alpha preservation a stronger hypothesis is used. Section~
\ref{app:lifts} proves balanced-index alpha subdivision and the ordinary
perfect-matching lift, with their exact root and index conditions.
An arbitrary alpha labeling is never silently assumed to survive
subdivision as an alpha labeling.

\section{Small-arm inputs and the short four-arm theorem}\label{shared:short-cores}
\subsection{Construction inputs for small arms}
\label{seven:small-inputs}

The following are the precise small-arm inputs needed in this section.
Their proofs and their complete finite tables are in
Sections~\ref{app:small-alpha}, \ref{app:pair-offsets}, and
\ref{app:finite-bases}.

\begin{proposition}\label{seven:small-alpha}
The spider $\Sp{s,a,b}$ has a root-zero alpha labeling in each of the
following cases:
\begin{enumerate}
\item $s=1$, provided $(a,b)\ne(2,2)$;
\item $s\in\{2,3,4,5\}$, $s\le a\le b$, and $b\ge7$;
\item $s\in\{2,3,4,5\}$ and $a=b=6$.
\end{enumerate}
\end{proposition}

The first assertion follows by adjoining a new maximum leaf to the
prescribed-zero alpha path on the other two arms.  For the second,
the appendix supplies explicit near-band constructions and their
resonance repairs.  To specify exactly what those proofs must establish,
put $e=\floor{s/2}$ and $k=e+\floor{a/2}$.  The alpha extension lemma
already gives every $b>2k$ other than $b=4k+1$.  Thus the additional
construction obligations are exactly
\begin{center}
\begin{tabular}{c|c|c|c}
$s$ & parity of $a$ & near band & resonance\\ \hline
$2,3$ & odd & $a\le b\le a+1$ & $b=2a+3$\\
$2,3$ & even & $a\le b\le a+2$ & $b=2a+5$\\
$4,5$ & odd & $a\le b\le a+3$ & $b=2a+7$\\
$4,5$ & even & $a\le b\le a+4$ & $b=2a+9$.
\end{tabular}
\end{center}
Only values with $b\ge7$ are asserted in the second part.  The
small-arm appendices prove these bands by explicit path formulas,
fixed affine edge trades with all-parameter topology checks, and their
exhaustive finite complements.  In particular, a finite experiment
inside a near band is not used as an existence argument outside it.
The third part consists of the four stated literal alpha labelings.

\begin{proposition}\label{seven:small-pairs}
For $x\in\{2,3,4,5\}$ and $n\ge x$,
\[
   \{0,\ldots,x-1\}\subseteq\Om(\Sp{x,n}).
\]
In addition, $2\in\Om(\Sp{2,n})$ for every $n\ge3$.
\end{proposition}

The pair appendix gives the cracked-permutation formulas, the ordinary
graceful tails, and every finite starting word needed here.  The prescribed-first ordinary path theorem supplies each graceful
tail in the stated range; the displayed prefixes partition all labels and colors.  Consequently these are uniform
statements for every $n$ in the displayed ranges.

We also use the hub-zero three-arm theorem
(Theorem~\ref{common:hubzero-three}).  The finite basis denoted below by
$F_4$ consists of all $\binom84=70$ sorted four-tuples with
entries in $\{2,\ldots,6\}$, each accompanied by a hub-zero graceful
labeling.  Verification of a finite labeling is explained in
Section~\ref{templates:finite-check}; the full words belong to
Section~\ref{app:finite-bases}.

\subsection{Four arms with a short arm}
\label{seven:four-short-section}

\begin{lemma}\label{seven:four-short}
Every four-arm spider with an arm of length at most five has a
hub-zero graceful labeling.
\end{lemma}
\begin{proof}
An arm of length one can be added by a new maximum leaf to the
hub-zero labeling of the other three arms.  It therefore remains to
treat minimum lengths two, three, four, and five.  In each case, if
the largest length is at most six, $F_4$ applies.  Hence throughout
the remaining argument the largest length $c$ is at least seven.

\medskip\noindent
\emph{Minimum two.}
Write the lengths as $2\le a\le b\le c$, with an additional arm of
length two.  Proposition~\ref{seven:small-alpha}, together with the
nine literal small positive cores in the finite appendix, supplies
$\Sp{2,a,b}$ except possibly when
\begin{equation}\label{seven:min2-exceptions}
 (a,b)\in\{(2,2),(2,3),(2,4),(2,6),(3,3),(4,4)\}.
\end{equation}
The nine positive cores are
\[
 (2,2,5),\ (2,3,4),\ (2,3,5),\ (2,3,6),\ (2,4,5),\
 (2,4,6),\ (2,5,5),\ (2,5,6),\ (2,6,6).
\]
No nonexistence claim about the exceptional cores is required.

For a supplied alpha core its index is
$J=1+\floor{a/2}+\floor{b/2}\le b+1$.  If $c\ge b+2$, then
$c>J$ and ordinary attachment applies.  If $c=b$, necessarily $b\ge7$.
When $a=2$, $J=2+\floor{b/2}<b$, so attachment still applies.
When $a\ge3$, use instead the alpha pair $\Sp{2,a}$, whose index
$1+\floor{a/2}\le b-1$ lies in the equal-pair offset interval for
$\Sp{b,b}$.  If $c=b+1$, attachment applies except when $a=b$ is
even.  In that remaining case $b\ge6$; the alpha pair $\Sp{2,c}$
has index $1+b/2\le b-1$ and can be combined with the equal pair
$\Sp{b,b}$.  This also includes $b=6,c=7$.

For the pairs $(2,2),(2,4),(2,6),(4,4)$ in
\eqref{seven:min2-exceptions}, at most one of the four arms is odd.
Halving every even length and flooring the possible odd length gives
a four-arm contree with a unit arm.  That contree is hub-zero graceful
by the first paragraph.  The ordinary full-subdivision or
perfect-matching lift of the foundations gives the required
root-zero labeling of the target.  For $(a,b)=(2,3)$, compose the
alpha pair $\Sp{2,3}$, of index two, with the offset-two pair
$\Sp{2,c}$.  For $(a,b)=(3,3)$, use the same alpha pair with the
offset-two pair $\Sp{3,c}$.  Proposition~\ref{seven:small-pairs}
provides both offsets.  This exhausts minimum two.

\medskip\noindent
\emph{Minimum three.}
Write the spider as $\Sp{3,a,b,c}$ with $3\le a\le b\le c$.
The supplied positive alpha cores $\Sp{3,a,b}$ include every case
except possibly $(a,b)=(3,3)$ and $(4,4)$.  For $b\le6$, the
eight positive literal cores are
\[
 (3,3,4),\ (3,3,5),\ (3,3,6),\ (3,4,5),\
 (3,4,6),\ (3,5,5),\ (3,5,6),\ (3,6,6).
\]
For a positive core, $J=1+\floor{a/2}+\floor{b/2}\le b+1$.
Thus $c\ge b+2$ is an ordinary attachment.  If $c=b$, the index
$1+\floor{a/2}$ of $\Sp{3,a}$ is at most $b-1$; combine it with
the equal offset pair $\Sp{b,b}$.  If $c=b+1$, attachment fails
only when $a=b$ is even.  Then $b\ge6$ and the alpha pair
$\Sp{3,c}$ has index $1+b/2\le b-1$, so the same equal-pair
composition applies.  For the exceptional core $(a,b)=(3,3)$, use
the alpha pair $\Sp{3,3}$ of index two and the offset-two pair
$\Sp{3,c}$.  For $(a,b)=(4,4)$, use the alpha pair $\Sp{3,4}$ of
index three and the offset-three pair $\Sp{4,c}$.  These offsets
are covered by Proposition~\ref{seven:small-pairs}.

\medskip\noindent
\emph{Minimum four or five.}
Write $\Sp{s,a,b,c}$, where $s\in\{4,5\}$ and
$s\le a\le b\le c$.  The small-alpha theorem and the six literal
cores
\[
 (4,5,5),\ (4,5,6),\ (4,6,6),\
 (5,5,5),\ (5,5,6),\ (5,6,6)
\]
give an alpha core $\Sp{s,a,b}$ except possibly for
\begin{equation}\label{seven:min4-exceptions}
 (s,a,b)=(4,4,4),\ (4,4,5),\ (4,4,6).
\end{equation}
For the first and third exceptional cores, there is at most one odd
arm in the target four-arm spider.  Its halved contree contains a
two-arm, so the already established minimum-two case and the ordinary
parity lift apply.  The same argument covers $\Sp{4,4,5,c}$ when
$c$ is even.  When $c$ is odd, compose the alpha pair $\Sp{4,4}$
of index four with the offset-four pair $\Sp{5,c}$; here $c\ge7$,
and Proposition~\ref{seven:small-pairs} applies.

Now suppose that the alpha core is supplied.  Put
\[
 k=2+\floor{a/2},\qquad J=k+\floor{b/2}.
\]
If $c>J$, attach $c$.  Hence assume $c\le J$.  We have $k\le b-1$:
the only possible failure with $s\le a\le b$ would be $a=b=4$,
already in \eqref{seven:min4-exceptions}.  If $c=b$, use the alpha
pair $\Sp{s,a}$ and the equal offset pair $\Sp{b,b}$.
Assume therefore $c>b$.  The inequalities
$b\le c\le k+\floor{b/2}$ give $\ceil{b/2}\le k$, and hence
$\floor{c/2}\le k$.  The HIGH pair lemma then gives the desired
offset $k$ on $\Sp{b,c}$ unless
\begin{equation}\label{seven:small-resonance}
 c=4q+1,\qquad k=b-1-q
\end{equation}
for an integer $q\ge1$.

We reduce this resonance exactly.  Equations
\eqref{seven:small-resonance} and $c\le J$ imply
\[
 \floor{a/2}=b-q-3,\qquad
 5q+2\le b+\floor{b/2}.
\]
Since $a\le b$, $b\le2q+6$.  Therefore
\[
 5q+2\le b+\floor{b/2}\le3q+9,
\]
which forces $q\le3$.  If $q=1$, then $c=5>b$ forces $b\le4$,
whereas $7\le b+\floor{b/2}$ forces $b\ge5$, a contradiction.
If $q=2$, then $c=9$, $b=8$, and $a\in\{6,7\}$.
If $q=3$, then $c=13$, $b=12$, and $a=12$.

The repeated case $a=b$ has another equal-pair composition.  In this
case $q=b-3-\floor{b/2}$, and the alpha pair $\Sp{s,c}$ has index
$2+2q$, equal to $b-4$ for even $b$ and $b-3$ for odd $b$.
This is an available offset on $\Sp{b,b}$.  Consequently only
\[
 \Sp{s,6,8,9},\qquad \Sp{s,7,8,9},\qquad s\in\{4,5\},
\]
remain.  For $\Sp{s,6,8,9}$ use the alpha pair $\Sp{s,8}$ of
index six and the following offset-six pair on $\Sp{6,9}$:
\begin{align*}
 &(15,13,14,20,7,21),\\
 &(12,16,11,8,19,9,18,10,17).
\end{align*}
Its labels are $[7,21]$ and its colors are $[1,15]$.  It can also
be obtained by factor-three offset subdivision of the offset-two
pair $(5,7)\mid(4,3,6)$ on $\Sp{2,3}$.

For $\Sp{4,7,8,9}$, the following words, read from hub zero, have
lengths $4,9,8,7$:
\begin{align*}
 &(28,1,27,2),\\
 &(24,3,26,4,23,5,25,8,22),\\
 &(16,6,21,9,20,11,19,12),\\
 &(13,7,10,15,17,18,14).
\end{align*}
They use every label and color from one through $28$ once.
Replace the first word by $(29,1,28,2,27)$ to obtain
$\Sp{5,7,8,9}$: the short color interval $[25,28]$ becomes
$[25,29]$, and all other edges are unchanged.  This finishes the
last two profiles, and the proof of the lemma.
\end{proof}


\section{Subdivision and matching lifts}
\label{app:lifts}

\subsection{Uniform offset subdivision}\label{app:subdivision}
We prove Theorem~\ref{common:offset-subdivision} in full. Let $T$ have
$m$ edges and an offset-$s$ labeling $g$. The zero-edge case is
immediate, so suppose $m\ge1$. Fix $q\ge1$ and $0\le e<q$. By
Lemma~\ref{common:prescribed-first-path}, choose a graceful permutation
$\rho=(\rho_1,\ldots,\rho_q)$ of $[1,q]$ with $\rho_1=q-e$.

Orient $T$ away from its distinguished root. For each old nonroot
vertex $v$, replace its parent edge by a block of $q$ vertices,
including the old child and excluding the old parent. Define the column labels in that block by
\begin{equation}\label{eq:subdivision-columns}
 B(v,h)=\begin{cases}
 qg(v)+\rho_h-q+e,&h\text{ odd},\\
 q(2s+m+1-g(v))+\rho_h-q+e,&h\text{ even}.
 \end{cases}
\end{equation}
Traverse the block in column order $1,\ldots,q$ when the old depth of
$v$ is odd and in the reverse order when that depth is even. Keep the
old root at zero. Column parity determines which formula is used;
depth parity determines only the traversal order.

\paragraph{Vertex labels.}
As $v$ ranges over the old nonroot vertices, both $g(v)$ and
$2s+m+1-g(v)$ range over $[s+1,s+m]$. Column $h$ therefore uses
\[
 qs+\rho_h+e,\ q(s+1)+\rho_h+e,\ldots,
 q(s+m-1)+\rho_h+e.
\]
The $q$ residues $\rho_h$ permute $[1,q]$, so all nonroot labels are
exactly $[qs+e+1,qs+e+qm]$, each once.

\paragraph{Connecting colors.}
The exit of a parent's block and entry of a child's block use the
same column, either $1$ or $q$. The same affine expression in $g$
therefore occurs at both ends, so their difference is $q$ times the
old edge difference. A child of the root starts in column one, whose
label is $qg(v)$ because $\rho_1=q-e$. Thus the $m$ connecting colors
are exactly $q,2q,\ldots,mq$.

\paragraph{Internal colors.}
Put $Z_m=\{1-m,3-m,\ldots,m-1\}$. For two consecutive columns,
the internal differences over all blocks are $\abs{qz+d}$,
$z\in Z_m$, where $|d|$ is the corresponding difference of $\rho$.
Since $Z_m=-Z_m$, the sign of $d$ is immaterial. The magnitudes of
the $q-1$ successive differences of $\rho$ are $1,\ldots,q-1$.
Hence the full internal multiset is
\begin{equation}\label{eq:subdivision-internal-colors}
 \{\abs{qz+d}:z\in Z_m,\ 1\le d<q\}.
\end{equation}
Every value is nonzero, smaller than $qm$, and not divisible by $q$.
If two signed expressions agree with the same sign, reduction modulo
$q$ gives $d=d'$ and then $z=z'$. Opposite-sign agreement would give
$d+d'=q$ and $z+z'=-1$; the latter is impossible because $z,z'$ have
the same parity. Thus all $m(q-1)$ values are distinct and are exactly
the required nonmultiples of $q$ in $[1,qm]$. Together with the
connecting colors, they give every color once. The block construction
is the actual uniform subdivision of $T$, proving offset $qs+e$.
All residues $e$ are allowed, completing the proof.

\begin{corollary}[Equal-arm intervals]\label{common:equal-offsets}
For all positive integers $L,d$, the equal-arm spider $\Sp{L^d}$
admits every offset $0,\ldots,L-1$.
\end{corollary}
\begin{proof}
Start with the root-zero unit star of degree $d$ and take $q=L$,
$s=0$. This includes $L=1$.
\end{proof}

\subsection{Perfect-matching lift}
\begin{theorem}\label{common:matching-lift}
Let $T$ have a perfect matching $M$, and let $C$ be the tree obtained
by contracting every matching edge. If $C$ has a graceful labeling,
then $T$ is graceful. If a distinguished matching pair of $C$ has
label zero, either prescribed vertex of that pair may be chosen to
receive zero in $T$.
\end{theorem}
\begin{proof}
Suppose $C$ has $m$ edges and graceful labeling $g$. Assign to the
two vertices of its pair $v$ the values
\[
 2g(v),\qquad 2m+1-2g(v).
\]
Choose the orientation of the root pair to put zero on the prescribed
vertex. Propagate orientations along the tree $C$ so that each
nonmatching edge joins values of the same parity. There is one parent
constraint per new pair and no cycle, so this is consistent. Vertex
labels partition $[0,2m+1]$. The nonmatching colors are the even
values $2|g(v)-g(w)|$, hence $2,4,\ldots,2m$. Matching colors are
$|4g(v)-(2m+1)|$. They are distinct: equality with opposite signs
would force $g(v)+g(w)=m+\tfrac12$. There are $m+1$ such odd
values in $[1,2m+1]$, giving precisely all odd colors.
\end{proof}

This is the standard Broersma--Hoede construction
\cite{BroersmaHoede1999}; the proof above specifies the root orientation
needed here. No alpha conclusion is imposed on an arbitrary contree.

\begin{corollary}\label{common:unique-odd-contree}
A spider with exactly one odd arm has a perfect matching whose contree
has arm lengths $\floor{L_i/2}$; zero-length arms are omitted.
A hub-zero graceful contree gives a hub-zero ordinary labeling of the
original spider.
\end{corollary}
\begin{proof}
Match the hub with the first vertex of the odd arm, then pair
successive vertices on every remaining part. Contracting gives the
stated lengths. Choose the matching-pair orientation from the theorem.
\end{proof}

\subsection{Balanced alpha subdivision}
\begin{theorem}\label{common:balanced-alpha-subdivision}
If a rooted $m$-edge tree has a root-zero alpha labeling $f$ with
index $k=\floor{m/2}$, its uniform $q$-subdivision is root-zero alpha
for every positive integer $q$, with index $\floor{qm/2}$.
\end{theorem}
\begin{proof}
Use \eqref{eq:subdivision-columns} with $s=e=0$ and the Walecki
permutation $\rho=(q,1,q-1,2,\ldots)$, which starts at $q$.
The preceding proof already gives the correct labels and colors.
Let $K=\floor{qm/2}$. When $m=2k$, the label in an odd column is
at most $K$ exactly when the old $f(v)\le k$; an even column reverses
that condition. For $m=2k+1$, the only possible boundary block is
resolved by
\[
 \rho_h\ge\floor{q/2}+1\quad(h\text{ odd}),\qquad
 \rho_h\le\floor{q/2}\quad(h\text{ even}).
\]
The same class rule follows. Consecutive columns therefore cross
$K$, and connecting edges do too because their old endpoints are in
opposite old classes. Root edges have value $qf(v)>K$. Thus every
edge crosses the new cut and the root remains zero.
\end{proof}

\begin{corollary}[Balanced alpha matching lift]\label{common:balanced-alpha-matching}
If the root-zero alpha contree has index $\floor{m/2}$, the
perfect-matching lift is root-zero alpha of index $m$.
\end{corollary}
\begin{proof}
The two labels of every matching pair sum to $2m+1$ and lie on
opposite sides of $m$. For an old edge the two $g$-values are on
opposite sides of $\floor{m/2}$; their even labels lie on opposite
sides of $m$, as do their odd labels. The same-parity orientation used
in the matching proof therefore makes every lifted edge cross $m$.
\end{proof}

\subsection{Leaf extensions with their exact deficits}
\begin{lemma}\label{common:lower-leaf-extension}
If a root-zero alpha tree has cut $K$ and a lower leaf of label
$K-\delta$, then an extension by $2r$ new vertices at that leaf
preserves root-zero alpha whenever $r\ge\delta+1$. The new cut is
$K+r$ and the new leaf is $K+r-\delta$.
\end{lemma}
\begin{proof}
Let $n$ be the old edge count. Leave old lower labels fixed and raise old upper labels by $2r$.
The unused interval is $[K+1,K+2r]$. Take an alpha path on
$[0,2r-1]$ whose first endpoint is the upper label
$i=2r-\delta-1$. This is an allowed upper label since
$r\ge\delta+1$, and an even-order path has no endpoint exception.
Its opposite endpoint is $i-r=r-\delta-1$ by
Lemma~\ref{common:even-path}. Add $K+1$ to all its labels and
join its first vertex to the old leaf $K-\delta$. The joining
color is $2r$, the internal new colors are $[1,2r-1]$, and the
shifted old colors are $[2r+1,2r+n]$. All labels are consecutive,
every edge crosses $K+r$, and the terminal label is
$K+1+(r-\delta-1)=K+r-\delta$. Root zero is unchanged.
\end{proof}

\begin{lemma}\label{common:upper-leaf-extension}
For an upper leaf of label $K+e$, $e\ge1$, extension by $2r$ new
vertices is possible whenever $r\ge e$. The new cut is $K+r$ and
the new leaf has label $K+r+e$.
\end{lemma}
\begin{proof}
Let $n$ be the old edge count. Leave old lower labels fixed, raise old upper labels by
$2r$, and fill the middle interval $[K+1,K+2r]$ with an alpha
path translated from $[0,2r-1]$. Request its first lower label
$j=e-1$, which is allowed because $r\ge e$. Its opposite upper
endpoint is $j+r$. The first new label is $K+e$, at distance
$2r$ from the raised old leaf $K+e+2r$. The internal new colors
are $[1,2r-1]$; the old colors are $[2r+1,2r+n]$. Every edge
crosses $K+r$, the root remains zero, and the final label is
$K+1+(e-1+r)=K+r+e$.
\end{proof}

These leaf bounds are essential. Neither lemma asserts that arbitrary
short independent extensions at several leaves are possible.

\section{Pair offsets and insertion constructions}\label{app:pair-offsets}
The generic high interval below is the rooted-offset construction of Rhys \cite[Lemma~5.1]{Rhys2026}. We include its full proof, the complementary low intervals, and the short-pair constructions used in the main theorems.

\subsection{Insertion and the high interval}
\begin{lemma}\label{po:insertion}
Let $T$ have $n$ edges and $s+\floor{t/2}\in\Om(T)$, where $t\ge2$ and $0\le s\le\ceil{t/2}-1$. Adjoining a $t$-edge arm at the root gives offset $s$, unless $t=4q+1$ and $s=q$.
\end{lemma}
\begin{proof}
Put $w=\floor{t/2}$. Choose an alpha path word $h$ on $t$ vertices with its first endpoint $i=t-s-1$ in the upper class and lower palette $[0,w-1]$. The endpoint theorem supplies it except at the displayed request. Indeed $i\ge w$, so when $t=4q+1$ only the forbidden upper endpoint $3q$ can occur. In the odd boundary case $i=w$, complementing the classes selects the required palette. Give the new arm labels
\[
F(v)=\begin{cases}s+1+h(v),&h(v)<w,\\s+n+1+h(v),&h(v)\ge w.\end{cases}
\]
The old and new nonroot label blocks are
$[s+1,s+w]$, $[s+w+1,s+w+n]$, $[s+w+n+1,s+n+t]$.
Old differences are $[1,n]$, new internal differences are $[n+1,n+t-1]$, and the joining root edge has difference $n+t$. This proves every inventory claim. The old root colors are unchanged.
\end{proof}

\begin{lemma}[High pair interval]\label{po:high}
For $2\le y\le z$ and $\floor{z/2}\le k\le y-1$, offset $k$ exists on $\Sp{y,z}$ unless
\[
z\equiv1\pmod4,\qquad 4k=4y-z-3.
\]
The construction has root-neighbor labels $y$ and $y+z$.
\end{lemma}
\begin{proof}
Set $w=\floor{z/2}$, $s=y-1-k$, $u=s+w$. Then $0\le u\le y-1$. A prescribed-first ordinary graceful word on $[1,y]$ beginning at $y-u$, translated by $u$, is an offset-$u$ $y$-arm whose root color is $y$. Insert the $z$-arm using Lemma~\ref{po:insertion}. Its capacity holds because $s\le y-1-w\le\ceil{z/2}-1$. The insertion exception is $z=4q+1$, $s=q$, precisely the displayed equality. The resulting offset-$s$ pair has root colors $y,N=y+z$. Reflect all nonroot labels by $N+y-v$. The interval becomes
\[
[N+y-(s+N),N+y-(s+1)]=[k+1,k+N].
\]
The two root colors exchange, and internal colors are fixed. This proves the conclusion.
\end{proof}

\begin{lemma}[Divisibility repair]\label{po:three-repair}
If $6\le y\le z$, $3\mid y,z$, and $\floor{z/2}\le k\le y-1$, the high interval has no exception.
\end{lemma}
\begin{proof}
Only $z\equiv1\pmod4$, $4k=4y-z-3$ needs repair. Modulo three it forces $3\mid k$. Write $y=3Y,z=3Z,k=3K$. Then $Y\ge2$, $Z\equiv3\pmod4$, $Z\ge Y$, and
$6K\ge3Z-1$, $3K\le3Y-1$. These imply $\floor{Z/2}\le K\le Y-1$. Lemma~\ref{po:high} has no exception at the quotient, since $Z\equiv3\pmod4$. Subdivision by three sends offset $K$ to $3K,3K+1,3K+2$, including $k$.
\end{proof}

\begin{lemma}[Low pair intervals]\label{po:low}
For $2\le x\le y$, every $0\le k\le\ceil{x/2}-1$ is available on $\Sp{x,y}$ except possibly $x=4q+1$, $k=q$. This exception is removed when $y\le6q+1$. Also every
\[
0\le k\le\ceil{y/2}-1,\qquad k+\floor{y/2}\le x-1
\]
is available unless $y=4q+1$, $k=q$.
\end{lemma}
\begin{proof}
Every $y$-arm has all offsets $0,\ldots,y-1$. Insert the $x$-arm with prior offset $k+\floor{x/2}\le x-1$. This gives the first claim. For its exceptional request with $x<y\le6q+1$, instead insert the $y$-arm into a single $x$-arm: $q+\floor{y/2}\le4q=x-1$, and the other capacity holds. Its exception would force $y=4q+1=x$, excluded. The case $x=y$ follows from equal-arm offsets. Finally exchange which arm is inserted to obtain the second displayed interval and its exact exception.
\end{proof}

\subsection{Complete intervals for minimum at most five}
\begin{theorem}\label{po:small}
For $1\le x\le5$ and $n\ge x$, every $0\le k<x$ is in $\Om(\Sp{x,n})$.
\end{theorem}
Throughout the proof let $P_M^{(r)}$ denote any ordinary graceful word on $[0,M-1]$ beginning at $M-r$, where $M\ge r$. The prescribed-first theorem proves its existence with no alpha condition. Consequently all the tails in the following constructions exist at every stated threshold.

\subsubsection{Minimum two and three}
For $n\ge9$ the cracked word
\[
Q_n=(2,n,3,n-1,0)\Vert(4+P_{n-5}^{(4)})
\]
uses $[0,n]\setminus\{1\}$, begins at $2$, and has internal colors $[1,n-1]$. Its five-vertex prefix supplies $n-2,n-3,n-4,n-1$; its tail begins at $n-5$, giving that joining color, and its internal tail supplies $[1,n-6]$. The small words are
\[
\begin{array}{c|l}
4&(2,0,3,4)\\5&(2,5,3,4,0)\\6&(2,6,4,3,0,5)\\
7&(2,7,3,0,6,4,5)\\8&(2,8,3,6,5,7,0,4).
\end{array}
\]
With $N=n+2$, the words $(N,2)$ and $N+1-Q_n$ give offset one: their labels are $[2,N+1]$ and colors are $N,N-2$, the long root color $N-1$, and $[1,N-3]$. The two smaller cases are $(4,3)\mid(2,5)$ for $n=2$, and $(5,4)\mid(2,6,3)$ for $n=3$.

Similarly, for $n\ge12$ set
\[
R_n=(3,n,2,n-2,4,n-1,0)\Vert(5+P_{n-7}^{(5)}).
\]
Its prefix colors are $n-3,n-2,n-4,n-6,n-5,n-1$, namely $[n-6,n-1]$; its joining color is $n-7$ and tail colors are $[1,n-8]$. Its labels are $[0,n]\setminus\{1\}$, first label $3$. The smaller words are
\[
\begin{array}{c|l}
3&(3,2,0)\\5&(3,4,0,2,5)\\6&(3,0,5,4,2,6)\\7&(3,0,6,2,7,5,4)\\
8&(3,0,7,2,8,4,6,5)\\9&(3,4,6,0,8,5,9,2,7)\\
10&(3,4,7,0,9,5,10,2,8,6)\\11&(3,5,8,0,10,4,9,2,11,7,6).
\end{array}
\]
With $N=n+3$, use $(N,3,N+2)$ and $N+1-R_n$. The short colors are $N,N-3,N-1$, the long root color is $N-2$, and the remaining colors are $[1,N-4]$. The labels are $[3,N+2]$. At $n=4$ use $(7,5,6)\mid(3,9,4,8)$ instead. This proves offset two for minimum three.
For offset zero at minimum two or three, take long arm $2+P_n^{(2)}$ and short arm $(n+2,1)$ or $(n+3,1,n+2)$, respectively. For offset one at minimum three take $3+P_n^{(3)}$ and $(n+3,2,n+4)$. Direct subtraction gives long colors $[1,n]$ and all remaining largest colors on the short arm. Thus the complete intervals follow.

\subsubsection{Minimum four}
Put $N=n+4$. Offsets zero and one use, respectively,
\[
(N,1,N-1,2)\mid(3+P_n^{(3)}),\qquad
(N,3,N+1,2)\mid(4+P_n^{(4)}).
\]
The long root color is $n$, its internal palette is $[1,n-1]$, and the short palette is $[n+1,n+4]$. Their labels are $[1,N]$ and $[2,N+1]$.
For offset two, $n\ge18$, use a cracked word with hole one and first five:
\[
Q=(5,n-2,3,n-3,6,n-4,4,n,2,n-1,0)
\Vert(7+P_{n-11}^{(7)}).
\]
The prefix uses low labels $\{0,2,3,4,5,6\}$ and high labels $[n-4,n]$. Its colors, subtracted from $n$, are $7,5,6,9,10,8,4,2,3,1$, exactly $[1,10]$. The tail fills $[7,n-5]$, begins at $n-11$, and supplies the remaining smaller colors. Use spider words $(N,3,N+2,4)\mid(n+5-Q)$. Their label interval is $[3,n+6]$; their colors are $N,N-3,N-1,N-2$ on the short arm, $n$ on the other root edge, and $[1,n-1]$ internally.
For offset three, $n\ge24$, use
\[
\begin{split}
Q={}&(5,n-2,3,n-1,1,n,6,n-6,8,n-5,4,n-4,7,n-3,0)\\
&\Vert(9+P_{n-15}^{(9)}).
\end{split}
\]
The low prefix labels are $\{0,1,3,4,5,6,7,8\}$, the high labels are $[n-6,n]$, and its colors subtracted from $n$ are
$7,5,4,2,1,6,12,14,13,9,8,11,10,3$, exactly $[1,14]$. The tail fills $[9,n-7]$, starts at $n-15$, and supplies the smaller colors. Take $(N,4,N+3,5)\mid(n+6-Q)$. Its labels are $[4,n+7]$; the short colors are $N,n,N-1,N-2$, the long root color is $n+1$, and the remaining colors are $[1,n-1]$. The finite offset-two cases $4\le n\le17$ and offset-three cases $4\le n\le23$ are all printed in Section~\ref{app:auxiliary-finite}.

\subsubsection{Minimum five}
Put $N=n+5$. Offsets zero and two are
\[
(N,1,N-1,2,N-2)\mid(3+P_n^{(3)}),\qquad
(N,4,N+1,3,N+2)\mid(5+P_n^{(5)}).
\]
Each has long colors $[1,n]$ and short colors $[n+1,N]$, with label intervals $[1,N]$ and $[3,N+2]$.
For offset one, $n\ge13$, put
\[
Q=(n-2,2,n,3,n-3,4,n-1,0)\Vert(5+P_{n-8}^{(5)}).
\]
This word uses $[0,n]\setminus\{1\}$, begins at $n-2$, and has colors $[1,n-1]$: the prefix, joining edge and tail give $[n-7,n-1]$, $n-8$, and $[1,n-9]$. Take $(N,2,N+1,4,N-1)\mid(3+Q)$. The short colors are $N,N-2,N-1,N-3,N-5$, the long root color is $N-4$, and all other colors are $[1,n-1]$. The labels are $[2,N+1]$.

For offsets three and four the following finite-prefix principle is convenient. Fix a hole $p$, a first label $c$, and $B=2p+c$. Let $L$ be a permutation of $[0,B-1]\setminus\{p\}$ starting at $c$ and ending at zero, and let $U$ permute $[0,B-3]$. If the adjacent sums $L_i+U_i,L_{i+1}+U_i$ are exactly $[1,2B-4]$, the alternating prefix
\[
L_0,n-U_0,L_1,n-U_1,\ldots,L_{B-2}
\]
has colors $[n-2B+4,n-1]$. For $n\ge3B-3$, append $B+P_{n-2B+3}^{(B)}$. This fills the unused interval $[B,n-B+2]$, begins at $n-2B+3$, and supplies every remaining color. Thus the result is a cracked word $Q$ with hole $p$ and first label $c$.
For offset three take
\[
p=1,\ c=6,\ B=8,\quad L=(6,7,3,5,4,2,0),\quad U=(5,3,2,4,0,1).
\]
The adjacent sums are $11,12,10,6,5,7,9,8,4,2,3,1$. For $n\ge21$ use $(N,4,N+3,5,N+2)\mid(N+1-Q)$. Its labels are $[4,N+3]$; the short colors are $N,N-4,N-1,N-2,N-3$, and the long arm supplies $[1,n]$.
For offset four take
\[
p=2,\ c=6,\ B=10,\quad L=(6,3,8,4,1,5,9,7,0),\quad U=(2,4,5,0,1,6,7,3).
\]
The adjacent sums are $8,5,7,12,13,9,4,1,2,6,11,15,16,14,10,3$. For $n\ge27$ use $(N,5,N+4,6,N+3)\mid(N+2-Q)$. Its labels are $[5,N+4]$; the short colors are $N,N-5,N-1,N-2,N-3$, the long root color is $N-4=n+1$, and the long internal colors are $[1,n-1]$. All $46$ finite complements, namely $5\le n\le12$ at offset one, $5\le n\le20$ at offset three, and $5\le n\le26$ at offset four, are printed in Section~\ref{app:auxiliary-finite}. Minimum one only asks for offset zero on a rooted path. This completes Theorem~\ref{po:small}.

\subsection{The additional offset two on a pair with a two-arm}\label{po:two-extra}
For every $n\ge3$, $2\in\Om(\Sp{2,n})$. For $n\ge15$ let
\[
Q=(3,n-1,1,n-2,4,n-3,5,n,0)\Vert(6+P_{n-9}^{(6)}).
\]
It uses $[0,n]\setminus\{2\}$, begins at $3$, and has colors $[1,n]\setminus\{n-1\}$. Indeed its prefix colors are $n-4,n-2,n-3,n-6,n-7,n-8,n-5,n$, the joining color is $n-9$, and the tail supplies $[1,n-10]$. Put $N=n+2$ and take $(N,3)\mid(N+2-Q)$. The nonroot labels are $[3,N+2]$, the short colors are $N,N-3$, and the long root color is $N-1$; its internal colors give all the other colors without duplication. The twelve finite words for $3\le n\le14$ are printed in Section~\ref{app:auxiliary-finite}. Thus this assertion is uniform for all $n\ge3$.

\subsection{Maximal offsets}\label{po:maximal}
\begin{theorem}
For positive $a\le b$, the largest attainable offset on $\Sp{a,b}$ is $a+b-2$, except that this value is unavailable for $(a,b)=(2,2)$ and $(3,3)$.
\end{theorem}
\begin{proof}
Write $m=a+b$. The two root neighbors are distinct labels at most $m$, so their interval beginning at $k+1$ forces $k\le m-2$. To attain equality, excluding temporarily $(1,1),(1,2),(2,2),(3,3)$, set $c=\ceil{a/2}+1$ and $\ell=b-c$. These exclusions give $1\le\ell\le b-1$. Take an ordinary graceful word $q$ on $[1,b-1]$ starting at $\ell$. Partition $[1,m]$ into the words
\[
A=(2,m,3,m-1,\ldots)\text{ truncated to length }a,
\qquad B=(1)\Vert(q+c).
\]
For $a=2w$, $A$ uses $[2,w+1]$ and $[m-w+1,m]$; for $a=2w+1$ its lower interval instead ends at $w+2$. In each case $B$ uses the complement. The internal colors of $A$ are $[b,m-2]$, the first internal color of $B$ is $\ell+c-1=b-1$, and its remaining colors are $[1,b-2]$. Add $m-2$ to every nonroot label; the two root colors become $m,m-1$. This supplies every color exactly once. The small pairs $(1,1)$ and $(1,2)$ use final words $(1)\mid(2)$ and $(2)\mid(3,4)$.

At maximal offset the two normalized words after subtracting $m-2$ must begin at $1,2$ and require internal differences $[1,m-2]$. For $(2,2)$ the two terminal assignments $(3,4)$ and $(4,3)$ yield colors $(2,2)$ and $(3,1)$, both invalid. For $(3,3)$, write the words as $(1,x,y)$ and $(2,z,w)$. The following complete table lists all $4!=24$ possibilities and their four internal colors. No row has color multiset $\{1,2,3,4\}$.
\begin{center}\small\begin{tabular}{cc|cc}
$(x,y,z,w)$ & Colors & $(x,y,z,w)$ & Colors\\\hline
$(3,4,5,6)$ & $2,1,3,1$ & $(5,3,4,6)$ & $4,2,2,2$\\
$(3,4,6,5)$ & $2,1,4,1$ & $(5,3,6,4)$ & $4,2,4,2$\\
$(3,5,4,6)$ & $2,2,2,2$ & $(5,4,3,6)$ & $4,1,1,3$\\
$(3,5,6,4)$ & $2,2,4,2$ & $(5,4,6,3)$ & $4,1,4,3$\\
$(3,6,4,5)$ & $2,3,2,1$ & $(5,6,3,4)$ & $4,1,1,1$\\
$(3,6,5,4)$ & $2,3,3,1$ & $(5,6,4,3)$ & $4,1,2,1$\\
$(4,3,5,6)$ & $3,1,3,1$ & $(6,3,4,5)$ & $5,3,2,1$\\
$(4,3,6,5)$ & $3,1,4,1$ & $(6,3,5,4)$ & $5,3,3,1$\\
$(4,5,3,6)$ & $3,1,1,3$ & $(6,4,3,5)$ & $5,2,1,2$\\
$(4,5,6,3)$ & $3,1,4,3$ & $(6,4,5,3)$ & $5,2,3,2$\\
$(4,6,3,5)$ & $3,2,1,2$ & $(6,5,3,4)$ & $5,1,1,1$\\
$(4,6,5,3)$ & $3,2,3,2$ & $(6,5,4,3)$ & $5,1,2,1$\\
\end{tabular}\end{center}

\end{proof}

\subsection{Three-arm offsets}\label{app:triple-offsets}
\begin{lemma}\label{po:triple-low}
Let $2\le x\le y\le z$, $w=\floor{x/2}$, and
$\floor{z/2}-w\le s\le\ceil{x/2}-1$, $s\ge0$. Then $s\in\Om(\Sp{x,y,z})$ unless either $x=4p+1,s=p$, or $z=4q+1,s+w=y-1-q$.
\end{lemma}
\begin{proof}
Put $u=s+w$. Its bounds give $\floor{z/2}\le u\le x-1\le y-1$, so apply the high pair construction to $(y,z)$, then insert the $x$-arm. The two exceptions are exactly those of the two operations.
\end{proof}
\begin{lemma}\label{po:repeated-triple}
For $2\le x\le y$ and $\floor{y/2}\le k\le x-1$, offset $k$ exists on $\Sp{x,y,y}$ except possibly when $y=4q+1$ and $k=3q$ or $k=x-1-q$. When $x=y$, equal-arm offsets remove both exceptions.
\end{lemma}
\begin{proof}
Set $s=x-1-k$, $u=s+\floor{y/2}$. The high pair lemma gives offset $u$ on $\Sp{x,y}$ with root colors $x,x+y$, apart from the first displayed exception. Inserting a $y$-arm gives offset $s$ and root colors $x,x+y,x+2y$, apart from the second. Put $N=x+2y$ and reflect all nonroot labels by $N+x-v$. These three root colors are permuted, internal colors are fixed, and the interval becomes $[k+1,k+N]$. Equal arms have every offset less than their common length.
\end{proof}

\subsection{Paired hub-zero four-arm spiders}\label{po:paired-four}
\begin{lemma}
For all positive $a,b$, $\Sp{a,a,b,b}$ is hub-zero graceful.
\end{lemma}
\begin{proof}
Assume $a\le b$. If $a\ne2$, the alpha pair $\Sp{a,a}$ has index $K=2\floor{a/2}$. For $a<b$, $K<b$; for $a=b$ odd, $K=a-1$. Combine this pair with equal-pair offset $K$ on $\Sp{b,b}$. Equal-even four arms are hub-zero by uniform subdivision of a star. If $a=2$ and $b\ge7$, Theorem~\ref{sa:main} gives alpha $\Sp{2,2,b}$ of index $2+\floor{b/2}<b$, so ordinary hub attachment adds the last arm. The remaining $b=2,3,4,5,6$ occur explicitly in the complete small four-arm table of Section~\ref{app:finite-bases}.
\end{proof}

\section{Alpha triples with a short arm}\label{app:small-alpha}
This section proves the input used when the shortest arm has length at most five. All finite words invoked below are printed in Section~\ref{app:auxiliary-finite}; their role is distinguished from the infinite affine constructions of Section~\ref{app:template-certificates}.

\begin{theorem}\label{sa:main}
If $1\le a\le b\le c$, $a\le5$, and $c\ge7$, then $\Sp{a,b,c}$ has a root-zero alpha labeling.
\end{theorem}

\subsection{Explicit fixed-endpoint path words}
All path words in this section are normalized to the labels $0,\ldots,n-1$. The word has $n$ vertices, and consequently $n-1$ internal differences.
\begin{lemma}\label{sa:endpoint-one}
An alpha word starting at lower label $1$ exists for $n\ge3$, $n\ne5$. It can be constructed recursively from
\[
(1,2,0),\quad(1,2,0,3),\quad(1,5,0,3,2,4),\quad
(1,8,0,6,2,7,4,5,3)
\]
of lengths $3,4,6,9$, by replacing a word $P$ of length $n$ with
\[
(1,n+2,0,n+3)\mathbin\Vert(P+2).
\]
For $n=2r+1$ its other endpoint is $r-1$; for $n=2r$ it is $r+1$.
\end{lemma}
\begin{proof}
The four new differences are $n,n+1,n+2,n+3$; the translated old differences are $1,\ldots,n-1$. The new labels $0,1,n+2,n+3$ and the translated old labels partition the full interval, and both classes remain consecutive. The four bases represent the four residues modulo $4$, beginning at the smallest admissible length in that residue. The endpoint identities follow from the signed endpoint sum for odd alpha words and the endpoint difference for even alpha words (Section~\ref{common:foundations}).
\end{proof}

The same construction has a useful general form. Fix $t\ge1$ and an alpha word $Q$ on $4t$ vertices starting at lower label $t$. Its opposite endpoint is $3t$. Given an alpha word $P$ of length $n$ starting at $t$, raise only the upper labels of $Q$ by $n$ and append $P+2t$. The joining difference is $n$, the first block supplies $n+1,\ldots,n+4t-1$, and the second supplies $1,\ldots,n-1$. The outside lower and upper blocks and the translated old blocks partition the labels. Thus this constructs a new alpha word of length $n+4t$ with the same prescribed endpoint.
For $t=2,3$, take one verified base in each residue, beginning at $2t+1$, except that the forbidden length $4t+1$ is replaced by $8t+1$. All these words, including $Q$, are printed in Section~\ref{app:auxiliary-finite}. For instance $Q=(2,5,3,4,0,7,1,6)$ works for $t=2$. Cattell's theorem shows that precisely the excluded length is unavailable at this endpoint; direct label and difference checks suffice for every printed base.

\subsection{Even balanced two-arm seeds}
\begin{lemma}\label{sa:even-balanced}
For every even $b\ge6$, $\Sp{2,b,b}$ is root-zero alpha.
\end{lemma}
\begin{proof}
First let $b=2k$ with $k\ge3$ odd. The root-zero alpha pair $\Sp{2,2k-2}$ has short word $(2k-1,1)$ and long word $(2k)\Vert(P+2)$, where $P$ is the endpoint-one word of length $2k-3$. Its index is $k$, its differences are $1,\ldots,2k$, and its first long vertex is $2k$. Raise its upper labels by $2k+2$ and attach the reversed Walecki arm
\[
D=(2k+2,2k+1,2k+3,2k,\ldots,3k+2,k+1).
\]
This yields an alpha spider with arms $2,2k-2,2k+2$, index $2k+1$, and all labels and differences $1,\ldots,4k+2$. The nonshort old arm $A$ is
\[
(4k+2)\ \Vert\!!\!\big\Vert_{j=0}^{(k-5)/2}
(2j+3,4k-1-2j,2j+2,4k-2j)\ \Vert\ (k,3k+3,k-1).
\]
When $k=3$ the concatenation of blocks is empty and $A=(14,3,12,2)$.

For $k\equiv3\pmod4$, replace the following edges, matching rows only as lists:
\[
\begin{array}{c|c}
\text{removed}&\text{added}\\\hline
(0,2k+2)&(0,3k+2)\\
((k+3)/2,(7k+7)/2)&((k+3)/2,(5k+5)/2)\\
(k+1,3k+2)&(k+1,2k+2)\\
((3k+3)/2,(5k+5)/2)&((3k+3)/2,(7k+7)/2)
\end{array}
\]
The removed colors are $2k+2,3k+2,2k+1,k+1$, and the added colors are the same multiset. Every degree and alpha class is preserved. The cuts on $A$ and $D$ have depths $k-1$ and $1,k+2,2k+2$, respectively. Write their resulting vertex segments as $A_0,A_1$ of lengths $k-2,k$, and $D_0,D_1,D_2$ of lengths $k+1,k,1$. The two new long arms are
\[
\operatorname{rev}(D_1)\Vert A_1,
\qquad A_0\Vert\operatorname{rev}(D_0)\Vert D_2,
\]
each of length $2k$.

For $k\equiv1\pmod4$, $k\ge5$, use instead
\[
\begin{array}{c|c}
\text{removed}&\text{added}\\\hline
(0,2k+2)&(0,3k+2)\\
((k-1)/2,(7k+3)/2)&((k-1)/2,(5k+5)/2)\\
(k,3k+3)&(k,3k+2)\\
(k+1,3k+2)&(k+1,2k+2)\\
(k+2,3k+2)&(k+2,3k+3)\\
((3k+3)/2,(5k+5)/2)&((3k+3)/2,(7k+3)/2)
\end{array}
\]
The removed colors are $2k+2,3k+2,2k+3,2k+1,2k,k+1$; the new colors are the same in another order. Degrees and classes agree. The $A$ cuts have depths $k-1,2k-3$ and give lengths $k-2,k-2,2$. The $D$ cuts have depths $1,k+2,2k+1,2k+2$ and give lengths $k+1,k-1,1,1$. The long arms become
\[
D_2\Vert\operatorname{rev}(A_1)\Vert D_1\Vert A_2,
\qquad A_0\Vert\operatorname{rev}(D_0)\Vert D_3.
\]
Their lengths are again $2k$. These explicit traversals also exclude an unwanted cycle.

If $b\equiv0\pmod4$, $b\ge8$, put $k=b/2$. The pair $\Sp{k,k}$ is alpha because $k\ge4$. Add a unit arm, retaining root zero and index $k$. This gives $\Sp{1,k,k}$ with $2k+1$ edges and balanced index $k$. Balanced subdivision (Section~\ref{app:lifts}) with factor $2$ supplies $\Sp{2,b,b}$. The two residues exhaust even $b\ge6$.
\end{proof}

\subsection{Minimum two}
For odd $b\ge5$ the required near bands are $c=b,b+1$; for even $b\ge6$ they are $c=b,b+1,b+2$. The sixteen explicit affine constructions in Section~\ref{app:template-certificates} supply all bands except the even diagonal, which is Lemma~\ref{sa:even-balanced}. Their source is the alpha spider $\Sp{2,b-2,b+\delta+2}$: its pair core is
\[
(b-1,1),\qquad (b)\Vert(\operatorname{endpoint\_one}(b-3)+2).
\]
Its index is $\floor{b/2}$. Raise upper labels by $n=b+\delta+2$ and fill the middle interval with an $n$-vertex alpha path starting at label $n$. A reversed endpoint-one word, complemented as well when $n$ is odd, gives this endpoint; the even diagonal uses reversed Walecki. The source has the full label and color inventory. The affine trades preserve colors, classes, degrees and the required topology for all parameters above the printed thresholds, by the exact verification theorem of that appendix.

For clarity the inclusive thresholds are
\[
\begin{array}{c|rrr}
b\bmod8&\delta=0&\delta=1&\delta=2\\\hline
0&\text{Lemma~\ref{sa:even-balanced}}&40&40\\
1&41&41&--\\
2&\text{Lemma~\ref{sa:even-balanced}}&42&42\\
3&35&35&--\\
4&\text{Lemma~\ref{sa:even-balanced}}&36&36\\
5&37&21&--\\
6&\text{Lemma~\ref{sa:even-balanced}}&38&30\\
7&23&23&--
\end{array}
\]
The exact finite complement has $53$ near-band profiles: $24$ for odd $b$ and $29$ for even $b$. Together with $(2,3,9)$ and $(2,4,13)$ these are the $55$ words printed in Section~\ref{app:auxiliary-finite}.

It remains to justify the two resonances omitted by ordinary alpha attachment. For an alpha spider with $N$ edges and index $K$, summing the positive signed edge differences gives
\begin{equation}\label{sa:signed-leaves}
\frac{N(N+1)}2
=2\left(\sum_{j=K+1}^{N}j-\sum_{j=0}^{K}j\right)
-\sum_{\text{upper leaves}} f(v)+\sum_{\text{lower leaves}}f(v).
\end{equation}
This follows because a nonleaf nonroot vertex has degree two, a leaf has degree one, and the root contributes zero.
For odd $b$, the seed $\Sp{2,b,b}$ has short leaf $1$ and upper long leaves $U,V$. Substitution gives $U+V=3b+4$. Since this is odd, one leaf has upper distance at most $(b+3)/2$ above index $b$. Extend that arm by the even number $b+3$, using the upper-leaf extension of Section~\ref{app:lifts}. Its length becomes $2b+3$.
For even $b$, the mixed seed $\Sp{2,b,b+1}$ has its unique upper leaf at distance $e$ from index $b+1$, with $2e\le b+4$. The infinite source and every degree-preserving trade retain upper leaf $3b/2+2$, giving $e=b/2+1$; finite words are checked for the stated sufficient bound. Extend that odd arm by $b+4$ to reach $2b+5$.

All other $c$ follow from the alpha pair $\Sp{2,b}$ and the strict attachment theorem: its index is $1+\floor{b/2}$. Thus the threshold is $c\ge b+2$ for odd $b$, or $c\ge b+3$ for even $b$, with precisely the resonances just proved. For $b=2$, use the published alpha construction for $\Sp{2,2,c}$, $c\ge7$ \cite[Lemma~3.6]{Panpa2025}. For $b=3,4$ attachment leaves exactly $c=9,13$, already printed. This proves Theorem~\ref{sa:main} when $a=2$.

\subsection{Minimum three}
The following fixed prefix replacement will be used without any search. If a size-$N$ alpha spider has short arm $(N-1,1)$ and another arm begins
\[
(N,3,N-3,2,N-2,5,N-5,4,N-4),
\]
replace the short arm with $(N,1)$ and that prefix with
\[
(N-2,2,N-1,4,N-3,3,N-5,5,N-4).
\]
The affected labels agree, and both complete affected color multisets are $[N-10,N]$. Classes and lengths agree. Now replace $(N,1)$ with $(N+1,1,N)$. This adds only label and difference $N+1$, retains the index, and increases the short length from two to three.

No affine two-arm seed trade cuts the short arm. Requiring every cut on the core long arm to have depth at least $10$ protects this prefix. The sixteen thresholds are unchanged except $(b\bmod8,\delta)=(5,1)$, whose threshold becomes $29$. Their exact coordinate inequalities are checked in the affine supplement. For $b\equiv0\pmod4$, $b\ge20$, the balanced-subdivision seed of Lemma~\ref{sa:even-balanced} already has pinned short arm $(N,1)$. For $b\equiv2\pmod4$, $b\ge22$, the explicit edge trade there has its first core cut at $b/2-1\ge10$, so it also admits the prefix replacement.

The remaining $29$ seeds for $b\ge18$ are printed as pinned $(2,b,b+\delta)$ words in Section~\ref{app:auxiliary-finite}. Every required even mixed seed additionally satisfies the same upper-leaf distance bound as above. Applying the short-two-to-three replacement gives every near band for $b\ge18$. For $3\le b\le17$, all $44$ missing near bands and resonances with $c\ge7$ are printed directly as $(3,b,c)$ words.

For odd $b\ge18$, the pinned diagonal two-arm seed has upper leaves summing to $3b+4$. The short upgrade retains those leaves and index $b$, so extension by $b+3$ gives $\Sp{3,b,2b+3}$. For even $b$, upgrade the mixed seed and append $b+4$ to its odd long arm, obtaining $\Sp{3,b,2b+5}$. All remaining longer arms follow from the alpha pair $\Sp{3,b}$, whose index is $1+\floor{b/2}$. This exhausts $a=3$.

\subsection{Minimum four and five}
For $b\ge9$, $b\ne13$, put $k=\floor{b/2}$, $y=b-4$, $n=b+\delta+4$, $N=2b+\delta+4$. A pair $\Sp{4,b-4}$ is labeled by
\[
(b,1,b-1,2),\qquad b-2-\operatorname{endpoint\_two}(b-4).
\]
The second arm uses $[3,b-2]$, starts at $b-4$, and has colors $[1,b-5]$. The first arm supplies $b,b-1,b-2,b-3$. Thus the pair is alpha with index $k$. The excluded $b=13$ is exactly the endpoint-two exception.
Raise its upper labels by $n$. For the new arm take a length-$n$ fixed-endpoint word with endpoint $t=\floor{n/2}-k-1$ when $n$ is even, or $t=\floor{n/2}-k$ when $n$ is odd; reverse it, and in the odd case also complement it, then add $k+1$. Its first label is $n$ by the endpoint identities. For the needed bands $t\in\{1,2,3\}$, and the recurrences above apply. This source is alpha and has short arm $(N,1,N-1,2)$.

The remaining bands are $\delta=0,1,2,3$ for odd $b$, and $\delta=0,1,2,3,4$ for even $b$. Parameterize $b=48h+r$. The $216$ affine recipes, all their inclusive thresholds and complete schema are in Section~\ref{app:template-certificates}. Each proves the target profile $\Sp{4,b,b+\delta}$ and preserves the four short-arm edges. The maximum threshold is $103$. The exact complement for $b\ge9$ comprises $139$ pinned finite words, all printed in Section~\ref{app:auxiliary-finite}; in particular it includes the exceptional source value $b=13$.

Upgrade any such size-$N$ word by replacing
\[
(N,1,N-1,2)\quad\hbox{with}\quad(N+1,1,N,2,N-1).
\]
The short color interval changes from $[N-3,N]$ to $[N-3,N+1]$. Its two old upper labels are shifted, one new upper label fills the gap, and other arms are fixed. Hence the new spider $\Sp{5,b,b+\delta}$ is alpha with unchanged index and unchanged long leaves.

For odd $b=2k+1$, the diagonal source has long upper leaves $3k+8$ and $3k+3$, with index $2k+2$. Their labels persist under degree-preserving trades. Extend the smaller leaf by $b+7=2k+8$, which is at least twice its upper distance. This gives resonance $c=2b+7$. For even $b=2k$, the $\delta=1$ source has upper long leaf $3k+3$, also at distance $k+1$ above index $2k+2$. Append $b+8=2k+8$ to obtain $c=2b+9$. The finite witnesses satisfy the sufficient bounds $2e\le b+7$ and $2e\le b+8$, respectively. Both arguments also apply after the four-to-five upgrade.

For $b=4$, the ordinary alpha witnesses $(4,4,7),(4,4,8),(4,4,17)$ supply the exceptional cases. For $b=5$, the words $(4,5,7),(5,5,7)$ and pinned helper $(4,5,8)$ give the near band; the pinned $(4,5,5)$ helper gives resonance $17$. For $b=6$, use pinned helpers with $c=7,8,9,10$ and the $c=7$ leaf interface for resonance $21$. For $b=7$, use $c=7,8,9,10$ and again resonance $21$. For $b=8$, use $c=8,9,10,11,12$ and the $c=9$ leaf interface for resonance $25$. These twenty small helper words and all needed interfaces are printed and verified in Section~\ref{app:auxiliary-finite}.

Every remaining $c$ follows from alpha attachment to $\Sp{a,b}$, whose index is $2+\floor{b/2}$ for either $a=4$ or $a=5$. Its strict threshold is $c>b+3$ for odd $b$, or $c>b+4$ for even $b$, and its sole exception is exactly the handled resonance. Finally, when $a=1$ and $c\ge7$, $\Sp{b,c}$ is an alpha pair; add the unit arm using Lemma~\ref{common:unit-stability}. This completes Theorem~\ref{sa:main}.

\section{Odd kernels and pair-offset constructions}
\label{ok:chapter}\label{app:odd-kernels}

This section supplies the path constructions behind the odd three-arm
$\alpha$ theorem and the additional odd-pair offset intervals. A word
lists the labels of a path in their order. When it describes an arm,
the hub label zero is understood to precede the word. All intervals
below are intervals of integers; an interval with its lower endpoint
larger than its upper endpoint is empty. Concatenation does not
identify the last entry of one word with the first entry of the next.

The only prescribed-path existence inputs are Lemmas~\ref{common:even-path}
and~\ref{common:odd-path}. Every use of the odd-order input checks its
exception explicitly. The small-minimum three-arm constructions are
proved in Section~\ref{app:small-alpha}. All other constructions used
here are given below, including their literal initial words.

\subsection{The kernel calculus}
\label{ok:calculus}

\begin{definition}[Odd kernel]\label{ok:kernel-definition}
For integral parameters $S,U,T,V$, with $T\ge1$, a word of type
$K(S,U,T,V)$ has $2T+1$ entries and the following inventories:
\begin{align}
 \text{lower labels}&=[1,T], \nonumber\\
 \text{upper labels}&=[T+1,2T+2]\setminus\{T+V\},\label{ok:kernel-labels}\\
 \text{internal differences}
   &=[1,2T+1]\setminus\{D\},
 \qquad D=S+T-U+V, \label{ok:kernel-colors}\\
 \text{ordered endpoints}&=(E,F),\nonumber\\
 E&=2T-S+2,\qquad F=2T+2S-U-V+3. \label{ok:kernel-ends}
\end{align}
Entries alternate between the two classes and both endpoints are
upper. In particular the displayed hole, missing color, and endpoints
must lie in their specified ranges. A kernel means an actual word
with these properties, rather than merely a feasible parameter tuple.
\end{definition}

The parameters of an auxiliary kernel need not be parameters of a
spider completion. In particular $U=0$ and $V>U$ are allowed in
source kernels. The restriction $1\le V\le U$ is imposed when an
even path of half-order $U$ is used to complete or extract a kernel.

\begin{lemma}[Completion of a kernel]\label{ok:completion}
Suppose $p$ is $K(s,u,t,v)$, where $s,t,u\ge1$ and $1\le v\le u$.
Then $\Sp{2s+1,2t+1,2u+1}$ has a root-zero $\alpha$-labeling of
index $s+t+u$. Its first arm can be chosen as the maximum Walecki word
\[
 W=(N,1,N-1,2,\ldots,N-s),\qquad N=2s+2t+2u+3.
\]
Every label on the other two arms is smaller than $N-s$.
\end{lemma}
\begin{proof}
Map a lower entry $x$ of $p$ to $x+s$, and an upper entry to
$x+s+2u$; call the resulting arm $A$. Take the even path on
$[1,2u]$ from upper endpoint $u+v$ to lower endpoint $v$, append
$2u+v$, and add $s+t$ to all of its entries; call this arm $B$.
Lemma~\ref{common:even-path} applies because $1\le v\le u$.

The lower labels of the three arms are respectively
$[1,s]$, $[s+1,s+t]$, and $[s+t+1,s+t+u]$. Put $K=s+t+u$.
The upper labels of $B$ before its appended entry are
$[K+1,s+t+2u]$. The upper labels of $A$ are
\[
 [s+t+2u+1,s+2t+2u+2]\setminus\{s+t+2u+v\};
 \]
the appended entry of $B$ fills that hole. Finally $W$ supplies
$[N-s,N]$. These are disjoint consecutive blocks, with the one
explicit hole filled once, so the nonroot labels are exactly $[1,N]$.
All edges cross $K$.

The internal colors of $A$ are
$[2u+1,2t+2u+1]\setminus\{D+2u\}$, and its root edge has color
$s+2u+E=2t+2u+2$. The internal colors of $B$ are $[1,2u]$:
its appended edge has color $2u$. Its root edge has color
\[
 s+t+u+v=D+2u.
 \]
The $2s+1$ edges of $W$, including its root edge, have consecutive
colors $[2t+2u+3,N]$. The three inventories partition $[1,N]$.
\end{proof}

\begin{lemma}[Diagonal growth]\label{ok:diagonal-growth}
Let $p$ be $K(S,U,T,V)$, put $d=T-S$, and assume $d\ge0$.
For every $P\ge d+2$ one obtains
\[
 K(S+P,U+P,T+P,V+P).
\]
\end{lemma}
\begin{proof}
Take an even path $q$ on $[1,2P]$ from upper $P+d+2$ to lower
$d+2$. Its endpoints are allowed because $1\le d+2\le P$.
Add $T$ to every entry of $q$. In the old word keep lower entries
fixed and raise upper entries by $2P$. Prepend the translated $q$.

The new lower labels are $[1,T]\cup[T+1,T+P]$.
The new upper blocks are $[T+P+1,T+2P]$ and
$[T+2P+1,2T+2P+2]\setminus\{T+V+2P\}$.
Thus the new hole is $(T+P)+(V+P)$. The prefix contributes
$[1,2P-1]$ and its last label is $T+d+2=E$.
The old first label became $E+2P$, so the join has color $2P$.
Every old color increases by $2P$. The new palette is therefore
$[1,2(T+P)+1]\setminus\{D+2P\}$, as required.
The new endpoints are
$T+P+d+2=2(T+P)-(S+P)+2$ and $F+2P$.
Substitution in \eqref{ok:kernel-ends} proves the asserted parameters.
\end{proof}

\begin{lemma}[Carrier prefix]\label{ok:carrier-prefix}
From $K(S,U,T,V)$ and any $w\ge S+1$ one obtains
\[
 K(S,U+w,T+w,V).
\]
\end{lemma}
\begin{proof}
Take the even path on $[1,2w]$ from upper $2w-S$ to lower $w-S$.
The inequality $w\ge S+1$ places both endpoints in their classes.
Raise its upper entries by $2T+2$. Prepend it to the old word with
every old entry raised by $w$.

The prefix lower block $[1,w]$ and old lower block $[w+1,w+T]$
are consecutive. The old upper block ends at $2T+w+2$ and the
prefix upper block is $[2T+w+3,2T+2w+2]$. The unique hole is
$T+V+w$. The prefix internal palette is
$[2T+3,2T+2w+1]$. Its final lower entry is $w-S$, while the old
first entry becomes $E+w$, so their join has color $2T+2$.
The old differences are unchanged. The first endpoint is
$2w-S+2T+2=2(T+w)-S+2$, and the last is $F+w$.
The missing color is unchanged since
$S+(T+w)-(U+w)+V=D$.
\end{proof}

\begin{lemma}[Asymmetric growth]\label{ok:asymmetric-growth}
From $K(S,U,T,V)$ and an integer
\[
 P\ge1,\qquad P\ge U+V-2S,
 \]
one obtains $K(S+P,U+2P,T+P,V)$.
\end{lemma}
\begin{proof}
Put
\[
 \ell=F+P-2T-2=P+2S-U-V+1.
 \]
The growth bound gives $\ell\ge1$. Since the old endpoint satisfies
$F\le2T+2$, we also have $\ell\le P$. Raise all old entries by $P$.
Append the even path on $[1,2P]$ from lower $\ell$ to upper
$\ell+P$, raising just its upper entries by $2T+2$.

The old labels lie in $[P+1,2T+P+2]$ except the translated hole.
The appended lower and upper labels fill respectively $[1,P]$ and
$[2T+P+3,2T+2P+2]$. They alternate across the new cut $T+P$.
The old colors are unchanged; the join has color
$(F+P)-\ell=2T+2$, and the new internal colors are
$[2T+3,2T+2P+1]$. The endpoints become $E+P,F+2P$.
These inventories have the stated parameters and the same missing
color $D$.
\end{proof}

\begin{lemma}[Reversal and class reflection]\label{ok:reversal}
\label{ok:reflection}
Reversing the word of a kernel gives the parameter transformation
\[
 J(S,U,T,V)=(U+V-2S-1,\ 2U+V-3S-1,\ T,\ V).
 \tag{J}
 \label{ok:J}
 \]
Complementing the labels separately within their classes, without
reversing the word, gives
\[
 R(S,U,T,V)=(T-S+1,\ 2T-U+2,\ T,\ T-V+3).
 \tag{R}
 \label{ok:R}
 \]
Both transformations are involutions on their realizable kernel
inventories.
\end{lemma}
\begin{proof}
For $J$, retain all labels. Substitution shows that the new first
endpoint is $F$, the new last endpoint is $E$, and the missing color
remains $D$. The classes and all other inventories are unchanged.
Applying the displayed affine transformation twice gives the identity.

For $R$ replace a lower entry $x$ by $T+1-x$ and an upper entry
$y$ by $3T+3-y$. This permutes each full class interval. The old
hole $T+V$ becomes $2T-V+3=T+(T-V+3)$. Every edge of color
$y-x$ becomes an edge of color $2T+2-(y-x)$. Its missing color is
therefore $2T+2-D$. The new endpoints are $3T+3-E$ and $3T+3-F$.
Substitution gives exactly \eqref{ok:R}; the label map is itself an
involution. Reversal is a separate operation: it is used in addition
to reflection only when explicitly stated.
\end{proof}

\subsubsection{Elementary seeds and a controlled-edge step}

\begin{lemma}[Three elementary sources]\label{ok:elementary-seeds}
The following kernels exist.
\begin{enumerate}
\item For $d\ge1$,
\[
 Z_d=(2d+2,1,2d+1,2,2d,\ldots,d,d+2)
       \quad\hbox{is }K(0,d,d,1).
 \]
\item For $r\ge2$, an even path on $[1,2r]$ from upper $r+2$
to lower $2$, followed by $2r+2$, is
$K(r,r,r,r+1)$.
\item For $h\ge2$, an even path on $[1,2h]$ from lower $2$ to
upper $h+2$, preceded by $2h+2$, is $K(0,0,h,h+1)$.
\end{enumerate}
\end{lemma}
\begin{proof}
In the first word the upper entries are $2d+2,2d+1,\ldots,d+2$;
its consecutive differences are $2d+1,2d,\ldots,2$.
Thus its label hole is $d+1$, color hole is $1$, and its endpoints
are $2d+2,d+2$, proving the first assertion.

In the second assertion the even path supplies $[1,2r-1]$ and
the appended edge has color $2r$. The omitted label is $2r+1$,
the missing color is $2r+1$, and the endpoints are $r+2,2r+2$.
In the third assertion the prepended edge has color $2h$.
The omitted label and missing color are both $2h+1$, and the
endpoints are $2h+2,h+2$. These agree with the kernel formulas.
The even-path requests have $2\le r$ and $2\le h$ respectively,
so all their prescribed endpoints are valid.
\end{proof}

\begin{lemma}[A controlled-edge enlargement]\label{ok:controlled-step}
Suppose a kernel $K(S,U,T,V)$ contains the edge $(2,T+1)$ and
$T\ge2$. There is a kernel
\[
 K(S+2,U+4,T+4,V+2)
 \]
which contains the new controlled edge $(2,T+5)$.
\end{lemma}
\begin{proof}
Raise old lower labels by $2$ and old upper labels by $6$.
The mapped cut is $T+4$. The old label blocks are $[3,T+2]$ and
$[T+7,2T+8]$ with the old hole raised by $6$. The fresh labels are
\[
 1,2,T+3,T+4;\qquad T+5,T+6,2T+9,2T+10.
 \]
The old controlled edge maps to $(4,T+7)$, with color $T+3$.
Replace this edge, using the reverse order if necessary, by
\begin{equation}
 (4,2T+10,1,2T+9,2,T+5,T+4,T+6,T+3,T+7).
 \label{ok:controlled-chain}
\end{equation}
Its nine colors in order are
\[
 (2T+6,2T+9,2T+8,2T+7,T+3,1,2,3,4).
 \]
The old controlled color is retained. Every other old color rose by
$4$, so the new colors added are exactly
$[1,4]\cup[2T+6,2T+9]$. This gives the new palette with missing
color $D+4$. All fresh entries occur once and cross the new cut.
The hole and both endpoints have risen by $6$.
These are precisely the changes specified by the new parameters.
The replacement contains $(2,T+5)$. Replacing one edge of a path
by a path through eight fresh vertices preserves a single path, so
there is no separate connectivity assumption.
\end{proof}

\subsection{All positive odd three-arm spiders}
\label{construction:all-odd-alpha3}

We first give the constructions for minimum arm at least five. The
two smaller minima and the single negative profile are treated at the
end of the section.

\subsubsection{Equal and repeated largest arms}

\begin{proposition}[Three equal odd arms]\label{ok:equal-odd-three}
For every $r\ge2$, $\Sp{(2r+1)^3}$ is root-zero $\alpha$, with
index $3r$, and its first arm is the maximum Walecki word.
\end{proposition}
\begin{proof}
At $r=2$ use
\[
 (15,1,14,2,13)\mid(10,3,8,4,12)\mid(9,6,7,5,11);
 \]
at $r=3$ use
\[
 (21,1,20,2,19,3,18)\mid
 (14,5,15,7,12,6,17)\mid(13,9,10,8,11,4,16).
 \]
Their lower labels are respectively $[1,6]$ and $[1,9]$,
and direct successive subtraction gives their full color intervals.

Now let $r\ge4$, $N=6r+3$, and $K=3r$.
For upper list $(U_0,\ldots,U_r)$ and lower list
$(L_0,\ldots,L_{r-1})$, interleaving means the word
$(U_0,L_0,U_1,\ldots,L_{r-1},U_r)$.
Take the first arm to be
\[
 W=(N,1,N-1,2,\ldots,r,N-r).
 \]
Its lower block is $[1,r]$, upper block $[5r+3,6r+3]$, and
palette $[4r+3,6r+3]$. For arm $A$ interleave
\begin{align*}
 U_A&=(4r+1,4r-1,4r+3,4r+4,\ldots,5r+1),\\
 L_A&=(2r-1,2r,2r-2,2r-3,\ldots,r+1).
\end{align*}
Its root color is $4r+1$, its first four internal colors are
$2r+2,2r,2r-1,2r+3$, and the remaining colors are consecutively
$2r+5,\ldots,4r$. Hence
\begin{equation}
 C_A=\{4r+1,2r-1,2r,2r+2,2r+3\}\cup[2r+5,4r].
 \label{ok:equal-three-A}
\end{equation}

If $r$ is even, take
\[
 L_C=(2r+1,\ldots,3r-3,3r-1,3r,3r-2).
 \]
The upper list begins with $4r+2$, then appends
$(4r-3-2j,4r-2j)$ for $0\le j\le r/2-3$, and finally appends
$(3r+2,3r+1,3r+4,5r+2)$. An empty loop is omitted.
The root color and first internal color are $4r+2,2r+1$.
Each loop packet contributes
\[
 2r-4-4j,\quad2r-5-4j,\quad2r-2-4j,\quad2r-3-4j.
 \]
Their union is $[7,2r-2]$. The seven final internal colors are
$5,3,2,1,4,6,2r+4$. Thus the complete palette is
\begin{equation}
 C_C=[1,2r-2]\cup\{2r+1,2r+4,4r+2\}.
 \label{ok:equal-three-C}
\end{equation}
If $r$ is odd, so $r\ge5$, use instead
\[
 L_C=(2r+1,\ldots,3r-3,3r,3r-1,3r-2).
 \]
The upper list begins with $4r+2$, then appends the same pairs
for $0\le j\le(r-7)/2$, then appends
$(3r+3,3r+5,3r+1,3r+2,5r+2)$.
The packets give $[9,2r-2]$ and the final nine colors are
$7,6,8,5,1,2,3,4,2r+4$, giving the same
\eqref{ok:equal-three-C}.

The lower lists partition $[1,3r]$. Arm $A$ has upper set
$\{4r-1,4r+1\}\cup[4r+3,5r+1]$; arm $C$ has upper set
\[
 ([3r+1,4r+2]\setminus\{4r-1,4r+1\})\cup\{5r+2\}.
 \]
Together with the upper block of $W$ these partition
$[3r+1,6r+3]$. Every edge crosses $3r$.
The three disjoint palettes
\eqref{ok:equal-three-A}, \eqref{ok:equal-three-C}, and that of
$W$ partition $[1,6r+3]$. This proves every inventory.
\end{proof}

\begin{lemma}[Repeated-arm kernel sources]\label{ok:repeated-kernels}
For $s\ge2$ and $t>s$, there is a kernel $K(s,t,t,v)$ with
$1\le v\le t$. More precisely:
\begin{enumerate}
\item If $s<t\le2s-2$, one can take $v=s+1$.
\item If $t\ge2s+1$, one can take $v=s+1$.
\item If $t=2s-1$, one can take $v=s$.
\item If $t=2s$ and $s\ge3$, one can take $v=s+3$.
At $(s,t)=(2,4)$ one can take $v=2$.
\end{enumerate}
\end{lemma}
\begin{proof}
In the first case put $d=t-s$, so $1\le d\le s-2$.
Apply diagonal growth with $P=s$ to $Z_d=K(0,d,d,1)$.
The required inequality is $s\ge d+2$, and the result is
$K(s,t,t,s+1)$.

In the second case start with $K(s,s,s,s+1)$ from
Lemma~\ref{ok:elementary-seeds} and apply carrier prefix with
$w=t-s\ge s+1$. This gives the same standard parameters.
The apparent failure $s+1>s$ in the initial source is harmless:
completion is used only after the prefix, when $s+1\le t$.

For $t=2s-1$ the desired label hole is $3s-1$, missing color
$2s$, and endpoints $3s,3s+2$. Two initial words are
\begin{align*}
 P_2&=(6,3,4,2,7,1,8),\\
 P_3&=(9,4,7,5,6,2,10,3,12,1,11).
\end{align*}
Both contain $(2,t+1)$ and have exactly the stated inventories.
Lemma~\ref{ok:controlled-step} changes $(s,t,t,s)$ to
$(s+2,t+4,t+4,s+2)$. Since $t+4=2(s+2)-1$, its two initial
words give every $s\ge2$.

For $t=2s$, $s\ge3$, take the standard source
$K(s-1,2s,2s,s)$ from the second case. First reflect within
classes and then reverse its word. The composite $J\circ R$ gives
$K(s,2s,2s,s+3)$: its hole is $3s+3$, missing color $2s+3$,
and ordered endpoints $3s+2,3s$. Here $s+3\le2s$.
For the isolated case $(2,4)$ the word
\[
 (8,3,5,4,7,1,10,2,9)
 \]
is $K(2,4,4,2)$; its hole is $6$, missing color $4$, and
endpoints $8,9$. These four cases exhaust $t>s$.
\end{proof}

\begin{corollary}\label{ok:repeated-alpha}
For odd $5\le a\le b$, $\Sp{a,b,b}$ is root-zero $\alpha$ and
admits a maximum Walecki $a$-arm. For $a<b$ this follows from
Lemma~\ref{ok:completion} and Lemma~\ref{ok:repeated-kernels};
for $a=b$ it follows from Proposition~\ref{ok:equal-odd-three}.
\end{corollary}

\subsubsection{The interior central cone and its arithmetic progressions}

\begin{lemma}[The non-AP central kernel]\label{ok:central-kernel}
Let
\[
 1\le s\le t\le u\le2s-2,\qquad u>s,\qquad u\ne2t-s.
 \]
There is a kernel $K(s,u,t,s+t-u+1)$.
\end{lemma}
\begin{proof}
Put
\[
 S=u-t,\quad T=u-s,\quad U=2u-s-t=S+T,\quad P=s+t-u.
 \]
Then $0\le S\le T$, $T\ge1$, $U\ge1$, and
$P\ge T-S+2$. Apply Lemma~\ref{common:odd-path} to the
$2T+1$-vertex path with upper first entry $2T-S+1$.
Its other endpoint is $T+S+1$. The first entry is between
$T+1$ and $2T+1$. The excluded request is
\[
 2(2T-S)=3T\quad\Longleftrightarrow\quad T=2S
 \quad\Longleftrightarrow\quad u=2t-s,
 \]
which is excluded in the hypothesis. The normalization in
Lemma~\ref{common:odd-path} also handles $S=T$, when the prescribed
entry is the first upper entry.

Raise every upper entry of this ordinary path by one. The new label
hole is $T+1$, the missing color is $1$, and its endpoints are
$2T-S+2,T+S+2$. Thus it is $K(S,U,T,1)$.
Diagonal growth by $P$ gives $K(s,u,t,P+1)$.
Finally $1\le P+1\le u$, since $U\ge1$ and $u=U+P$.
\end{proof}

\begin{lemma}[An explicit near-maximum-hole word]\label{ok:near-max-word}
For $q\ge2$ the following concatenation is a path word:
\[
 (3q+1),\quad
 (q+j,3q-j)_{j=0}^{q-1},\quad
 (2q,4q+2),\quad
 (j,4q+1-j)_{j=1}^{q-1}.
 \tag{*}\label{ok:near-max-formula}
 \]
Its lower labels are $[1,2q]$, its upper labels are
$[2q+1,4q+2]\setminus\{4q+1\}$, its colors are
$[1,4q+1]\setminus\{4q\}$, and its endpoints are $3q+1,3q+2$.
\end{lemma}
\begin{proof}
The lower packets are $[q,2q-1]$, $\{2q\}$, and $[1,q-1]$.
The upper packets are $\{3q+1\}$, $[2q+1,3q]$,
$\{4q+2\}$, and $[3q+2,4q]$. These are disjoint and give
the stated label sets.

Through the entry $2q$, the differences are exactly $[1,2q+1]$:
the first is $2q+1$, the first loop supplies alternating descending
even and odd differences down to $2$, and the following edge gives
$1$. The next edge gives $2q+2$. The remaining edges give all of
$[2q+3,4q+1]$ except $4q$: the within-pair differences are
$4q+1-2j$, and the joins give the intervening even differences.
All edges cross $2q$. The first and last entries are as asserted.
\end{proof}

\begin{lemma}[Pinned central AP kernels]\label{ok:ap-central}
If $d\ge2$ and $s\ge2d+2$, there is
\[
 K(s,s+2d,s+d,s-d+2).
 \]
For $d=1$ and every $s\ge2$ there is also a kernel completing
$\Sp{2s+1,2s+3,2s+5}$.
\end{lemma}
\begin{proof}
In \eqref{ok:near-max-formula} put $q=d$. Reflect its lower
labels by $2d+1-x$ and upper labels by $6d+3-x$. Its new hole
is $2d+2$, missing color $2$, and endpoints $3d+2,3d+1$.
These are exactly $K(d,3d,2d,2)$. Equivalently the reflected word
is the following explicit concatenation:
\[
 (3d+2),\quad
 (d+1-j,3d+3+j)_{j=0}^{d-1},\quad
 (1,2d+1),\quad
 (2d+1-j,2d+2+j)_{j=1}^{d-1}.
 \]
Apply diagonal growth with $P=s-d\ge d+2$. The result has
parameters $(s,s+2d,s+d,s-d+2)$, and its final hole parameter
is between $1$ and $s+2d$.

For $d=1$ the following three literal kernels supply $s=2,3,4$:
\begin{align*}
 K(2,4,3,2)&=(6,1,8,2,4,3,7),\\
 K(3,5,4,2)&=(7,4,5,3,8,2,9,1,10),\\
 K(4,6,5,4)&=(8,2,12,1,10,5,6,4,7,3,11).
\end{align*}
Their respective label holes are $5,6,9$ and missing colors
$3,4,7$, so direct subtraction checks all three.
For $s\ge5$, diagonally grow the first word by $P=s-2\ge3$.
It gives $K(s,s+2,s+1,s)$, with $1\le s\le s+2$.
The completion lemma applies in every case, with the Walecki
short arm retained.
\end{proof}

\begin{corollary}\label{ok:interior-cone}
If $5\le a\le b\le c\le2a-5$ are odd, then
$\Sp{a,b,c}$ is root-zero $\alpha$.
\end{corollary}
\begin{proof}
Write $a=2s+1$, $b=2t+1$, $c=2u+1$.
The equal case uses Proposition~\ref{ok:equal-odd-three}.
The non-AP case uses Lemmas~\ref{ok:central-kernel}
and~\ref{ok:completion}. In the remaining case put
$t=s+d$, $u=s+2d$, $d\ge1$.
The cone condition implies $s\ge2d+2$, so
Lemma~\ref{ok:ap-central} applies. This covers the complete
arithmetic-progression exception of the ordinary odd-path request.
\end{proof}

\subsubsection{The two central boundary layers}

\begin{lemma}[Odd-prefix boundary kernels]\label{ok:odd-prefix-boundaries}
For $\epsilon\in\{0,1\}$, suppose
\[
 u=2s-1+\epsilon,\qquad
 s\ge3+\epsilon,\qquad s\le t\le2s-3-\epsilon.
 \]
Then there is $K(s,u,t,t-s+3)$.
\end{lemma}
\begin{proof}
Put
\[
 T=s-1,\quad P=t-s+1,\quad L=2s-t-2-\epsilon,\quad v=P+2.
 \]
Here $T,P\ge1$ and $1\le L\le T$.
Take the even path on $[1,2T]$ from lower $L$ to upper $T+L$,
raising its upper entries by $2P+2$. Prepend the odd Walecki
word on $[1,2P+1]$, starting at $2P+1$ and ending at $P+1$,
translated by $T$.

The lower blocks are $[1,T]$ and $[T+1,t]$. The upper blocks
are $[t+1,T+2P+1]$ and $[T+2P+3,2t+2]$.
The label hole is $T+2P+2=t+v$.
The prefix colors are $[1,2P]$; the old even path contributes
$[2P+3,2t+1]$; and their join has color
\[
 (T+P+1)-L=2P+1+\epsilon.
 \]
Thus the sole missing color is $2P+2-\epsilon=s+t-u+v$.
The endpoints are
\[
 T+2P+1=2t-s+2,\qquad
 T+L+2P+2=t+s+1-\epsilon.
 \]
The second is $2t+2s-u-v+3$, proving the claimed kernel.
The inequalities also give $1\le v\le u$ when completion is wanted.
There is no odd-order prescribed endpoint in this argument: the odd
prefix is an explicit word.
\end{proof}

\begin{lemma}[The gap-two boundary kernels]\label{ok:gap-two-kernels}
For every $s\ge4$ there is
\[
 K(s,2s-1,2s-2,s+3).
 \]
Consequently $\Sp{a,2a-5,2a-3}$ is root-zero $\alpha$ for
every odd $a\ge5$.
\end{lemma}
\begin{proof}
For the displayed kernel the carrier is $T=2s-2$, the hole is
$3s+1$, missing color $2s+2$, and endpoints $3s-2,3s-3$.
At $s=4$ and $s=7$ use the controlled initial words
\begin{align*}
 P_4={}&(10,4,11,3,12,1,14,2,7,6,8,5,9),\\
 P_7={}&(19,7,20,6,21,4,25,1,26,3,23,5,24,2,\\
       &\hspace{9mm}13,12,14,11,15,10,16,9,17,8,18).
\end{align*}
They contain $(2,T+1)$ and have exactly the inventories just stated.
Lemma~\ref{ok:controlled-step} supplies $s\mapsto s+2$.
The $s=4$ word covers all even $s\ge4$, and the $s=7$ word
all odd $s\ge7$. The one remaining value, $s=5$, is supplied by
\[
 (13,6,9,8,10,4,18,1,17,2,15,5,14,3,11,7,12).
 \]
No controlled edge is required in this isolated word.
Since $s+3\le2s-1$ for $s\ge4$, completion is valid.

At $a=7$ the profile is $(7,9,11)$, supplied by the
$s=3,d=1$ word in Lemma~\ref{ok:ap-central}.
At $a=5$ it is $(5,5,7)$ and uses the complete
minimum-five theorem in Section~\ref{app:small-alpha}.
\end{proof}

\begin{lemma}[The reverse boundary]\label{ok:reverse-boundary}
For $s\ge2$ and $s\le t\le2s$, the spider
$\Sp{2s+1,2t+1,4s+1}$ is root-zero $\alpha$.
For $t<2s$ it has a completing kernel $K(s,2s,t,t-s+1)$.
\end{lemma}
\begin{proof}
First suppose $t<2s$ and put $r=t-s$, so $0\le r\le s-1$.
We construct the standard source $K(r,t,t,r+1)$.
For $r=0$ use $Z_t$. For $r=1$ use
\[
 B_3=(7,3,4,2,8,1,6),\qquad
 B_4=(9,3,5,4,8,1,10,2,7)
 \tag{B}\label{ok:B-small}
 \]
when $t=3,4$, respectively. They are $K(1,t,t,2)$.
For $t\ge5$, apply carrier prefix of size $w=t-3\ge2$ to
$B_3$; its even upper endpoint is $2w-1$.
For $r\ge2$, the inequality $t=s+r\ge2r+1$ permits the
standard short-wedge source in Lemma~\ref{ok:repeated-kernels}.
Reversal $J$ gives exactly $K(s,2s,t,r+1)$.
Its hole parameter satisfies $1\le r+1\le s\le2s$.
When $t=2s$, use the twice-carrier case of
Lemma~\ref{ok:repeated-kernels}, including its literal $(2,4)$
exception. Completion proves the assertion.
\end{proof}

\begin{proposition}[The complete central cone]\label{ok:central-cone}
For odd $5\le a\le b\le c\le2a-1$, the spider
$\Sp{a,b,c}$ is root-zero $\alpha$.
For $a\ge7$ these constructions retain the maximum Walecki $a$-arm.
\end{proposition}
\begin{proof}
The interior $c\le2a-5$ is Corollary~\ref{ok:interior-cone}.
At $c=2a-3$, either $b=c$, handled by the repeated-arm theorem,
or $b=c-2$, handled by Lemma~\ref{ok:gap-two-kernels}, or
$b\le c-4$, handled by $\epsilon=0$ in
Lemma~\ref{ok:odd-prefix-boundaries}.
At $c=2a-1$ use Lemma~\ref{ok:reverse-boundary}.
All three possibilities at each layer exhaust odd $b$.
For the small minimum $a=5$ the complete minimum-five theorem also
supplies any branch outside those displayed parameter domains.
Every construction for $a\ge7$ is a completion or the explicitly
pinned equal-arm construction. No Walecki pin is asserted for the
separate minimum-five input.
\end{proof}

\subsubsection{The far cone, including its resonance}

\begin{lemma}[Even extension at an upper leaf]\label{ok:upper-leaf-extension}
Suppose a root-zero $\alpha$-labeled tree has index $K$ and an
upper leaf labeled $K+e$. Any $2z$ vertices, with $z\ge e$, can be
appended to that leaf while preserving root-zero $\alpha$.
The new index is $K+z$ and the new leaf has label $K+z+e$.
\end{lemma}
\begin{proof}
Raise every old upper label by $2z$, leaving old lower labels fixed.
On the new middle interval use an even path on $[1,2z]$
from lower $e$ to upper $z+e$, adding $K$ to its entries.
Join its first entry to the former leaf. The join has color
$(K+e+2z)-(K+e)=2z$, the new internal colors are
$[1,2z-1]$, and all old colors have risen by $2z$.
The middle labels fill $[K+1,K+2z]$ exactly. Every edge crosses
$K+z$, and the asserted new leaf and distance follow.
\end{proof}

\begin{proposition}[The full odd far cone]\label{ok:far-cone}
If $7\le a\le b\le c$ are odd and $c\ge a+b-1$, then
$\Sp{a,b,c}$ is root-zero $\alpha$.
\end{proposition}
\begin{proof}
Write $a=2s+1$, $b=2t+1$, with $s\ge3$.
The rooted pair on arms $a,b$ is root-zero $\alpha$ of index
$s+t$, by Corollary~\ref{common:alpha-pair}.
Since $c>2(s+t)$, Lemma~\ref{common:alpha-extension}
attaches the $c$-arm except at
\[
 c_*=4(s+t)+1=2a+2b-3.
 \]
We now construct that value, rather than assume attachment at its
excluded endpoint.

Take the pinned repeated-arm construction $\Sp{a,b,b}$ from
Corollary~\ref{ok:repeated-alpha}, including its equal case.
Put $K=s+2t$ and $N=2K+3=2s+4t+3$.
All three leaves are upper. If their labels sum to $L$, summing the
positive edge differences by degrees gives
\[
 \frac{N(N+1)}2
   =2\sum_{j=K+1}^N j-2\sum_{j=1}^Kj-L,
 \]
because nonleaf nonroot vertices have degree two and the hub is zero.
Consequently
\begin{equation}
 L=\frac{N(N+1)}2-2K(K+1)=5K+6.
 \label{ok:signed-leaf-sum}
\end{equation}
The Walecki short leaf is $N-s=s+4t+3$.
The two long leaf labels therefore sum to $4s+6t+3$.
Their smaller label $U$ satisfies
\[
 U\le2s+3t+1,\qquad e=U-K\le s+t+1.
 \]
Choose that long arm and append $2(2s+t)$ vertices.
Since $2s+t\ge s+t+1\ge e$, Lemma~\ref{ok:upper-leaf-extension}
applies. Its length becomes
\[
 b+2(2s+t)=4s+4t+1=c_*.
 \]
This supplies the resonance.
The extension raises all old upper labels uniformly, so the short
arm remains the maximum Walecki word. The unextended lower labels
and all new middle labels stay below its final upper label.
\end{proof}

\subsubsection{The complete short wedge}

\begin{lemma}[Unequal short-wedge kernels]\label{ok:unequal-wedge-kernel}
Let $s\ge2$, $u\ge2s+2$, and $0\le\delta\le s-1$; put
$t=u+\delta$. There is a completing kernel $K(s,u,t,v)$, with
\[
 v=\begin{cases}
 s+\delta+3,&\delta\ne s-3,\\
 2s-1,&\delta=s-3.
 \end{cases}
 \]
\end{lemma}
\begin{proof}
In the nonexceptional case put $r=s+\delta+1$.
Apply the odd-path interface on $[1,2r+1]$ at upper endpoint
$2r-s+2$; its other endpoint is $2s+\delta+1$.
Both entries lie in the upper class. The sole excluded request is
\[
 2(2r-s+1)=3r
 \quad\Longleftrightarrow\quad r=2s-2
 \quad\Longleftrightarrow\quad\delta=s-3.
 \]
Thus the request is permitted here.
The word itself, with a new unused possible maximum $2r+2$,
is $K(s,s+1,r,s+\delta+3)$:
its hole is $r+v=2r+2$ and its missing color is $2r+1$.

If $\delta=s-3$, then $s\ge3$. Put $q=s-1$ in
Lemma~\ref{ok:near-max-word}. Its word is
$K(s,s+1,2s-2,2s-1)$, with missing color $4s-4$.
This is the needed replacement for the forbidden ordinary path.
In both cases apply carrier prefix with
$w=u-s-1\ge s+1$. The carrier becomes $t=u+\delta$,
the other half becomes $u$, and $v$ is unchanged.
Since $v\le2s+2\le u$ in the first case and $v=2s-1<u$
in the second, completion is allowed.
\end{proof}

\begin{proposition}[Both adjacent short-wedge boundaries]
\label{ok:short-wedges}
For odd $5\le a\le b\le c$ with $b\ge2a-1$,
$\Sp{a,b,c}$ is root-zero $\alpha$.
\end{proposition}
\begin{proof}
We can dispose of $a=5$ using Section~\ref{app:small-alpha};
write $a=2s+1$ with $s\ge3$ for the far-cone appeals below.

First let $b\ge2a+3$, so its half $u\ge2s+2$.
Write the largest half as $t=u+\delta$.
If $0\le\delta\le s-1$, use
Lemma~\ref{ok:unequal-wedge-kernel} and completion.
If $\delta\ge s$, then $c\ge a+b-1$ and
Proposition~\ref{ok:far-cone} applies.

Next let $b=2a+1$, with half $u=2s+1$.
For $2s+1\le t\le3s$ take the $\epsilon=0$ odd-prefix kernel
whose source parameters are
\[
 S_0=t-s+1,\quad U_0=2S_0-1,\quad T_0=t,\quad V_0=s+2.
 \]
Indeed $S_0\ge3$, $S_0\le t\le2S_0-3$,
as required by Lemma~\ref{ok:odd-prefix-boundaries}.
Reflecting it gives
\[
 K(s,2s+1,t,t-s+1).
 \]
The final hole parameter is at most $2s+1$ exactly because
$t\le3s$. For $t\ge3s+1$ the far cone applies.

Finally let $b=2a-1$, with half $u=2s$.
If $t=2s$, use the repeated-arm theorem.
If $t=2s+1$, reflect the gap-two source
$K(s+1,2s+1,2s,s+4)$ of
Lemma~\ref{ok:gap-two-kernels}. The result is
$K(s,2s+1,2s,s-1)$.
Completing it gives the desired unordered arm triple, with its two
long arms interchanged. The condition $s\ge3$ ensures $s-1\ge1$.
For $2s+2\le t\le3s-1$, use the $\epsilon=1$ odd-prefix source
with $S_0=t-s+1$, $U_0=2S_0$, $T_0=t$, $V_0=s+2$.
Its source inequalities hold, and reflection gives
$K(s,2s,t,t-s+1)$. Its final hole condition is exactly
$t\le3s-1$. For $t\ge3s$, the far cone applies.
All equality, gap-two, near, and far cases have now been included.
\end{proof}

\subsubsection{The last transition and its pinned AP ray}

\begin{lemma}[Transition kernels]\label{ok:transition-kernels}
The following kernels exist:
\begin{align}
 K(s,2s+g,s+g+h,h+1),
 &\quad g\ge1,\ h\ge2,\ s\ge g+h+1; \label{ok:transition-high}\\
 K(s,2s+g,s+g+1,3),
 &\quad g\ge1,\ s\ge g+3. \label{ok:transition-one}
\end{align}
\end{lemma}
\begin{proof}
For \eqref{ok:transition-high}, use the elementary pinned source
$K(0,0,h,h+1)$ from Lemma~\ref{ok:elementary-seeds}.
Carrier prefix of size $g$ gives $K(0,g,g+h,h+1)$.
The inequality $s\ge g+h+1=U+V$ permits asymmetric growth
by $P=s$, giving the asserted kernel. Its ordinary even appended
path has lower endpoint $s-g-h\ge1$ and upper endpoint
$2s-g-h\le2s$.

For \eqref{ok:transition-one}, the literal word
\[
 (6,1,4,2,3)=K(0,1,2,3)
 \]
has label hole $5$, missing color $4$, and endpoints $6,3$.
If $g>1$, use carrier prefix of size $g-1$; if $g=1$ do
nothing. The resulting kernel is $K(0,g,g+1,3)$.
Asymmetric growth by $s\ge g+3=U+V$ proves the second assertion.
The final hole parameters $h+1$ and $3$ lie in $[1,2s+g]$,
so either word can be completed.
\end{proof}

\begin{lemma}[The pinned limiting AP ray]\label{ok:pinned-ap-ray}
For every $s\ge3$ there is $K(s,3s-2,2s-1,2)$.
Its completion is a root-zero $\alpha$-labeling of
$\Sp{2s+1,4s-1,6s-3}$ with a maximum Walecki short arm.
\end{lemma}
\begin{proof}
The required word has label interval $[1,4s]\setminus\{2s+1\}$,
cut $2s-1$, palette $[1,4s-1]\setminus\{3\}$, and endpoints
$3s,3s+1$. At $s=3$ use
\[
 (9,5,6,4,11,3,8,2,12,1,10).
 \]
At $s=4,5$ use the two controlled initial words
\begin{align*}
 Q_4={}&(12,6,8,7,11,2,16,1,14,4,15,3,10,5,13),\\
 Q_5={}&(15,6,12,8,10,9,14,7,17,5,13,2,20,1,18,4,19,3,16).
\end{align*}
They contain the controlled edge $(2,2s+3)$.
These finite assertions follow by listing their absent label, endpoints,
and successive absolute differences: the latter give every integer
from $1$ to $4s-1$ except $3$.

For $s\ge4$ with a controlled word, raise lower entries by $2$
and upper entries by $6$. The mapped blocks are $[3,2s+1]$ and
$[2s+6,4s+6]\setminus\{2s+7\}$; the palette is
$[5,4s+3]\setminus\{7\}$. The controlled edge becomes
$(4,2s+9)$, of color $2s+5$. Replace it by
\[
 (4,4s+8,1,4s+7,2,2s+7,2s+3,2s+4,2s+2,2s+9),
 \]
or its reversal as dictated by the old edge orientation.
Its eight interior entries are the fresh lower labels
$1,2,2s+2,2s+3$ and upper labels
$2s+4,2s+7,4s+7,4s+8$.
Its colors in order are
\[
 (4s+4,4s+7,4s+6,4s+5,2s+5,4,1,2,7).
 \]
Thus the controlled color is retained; the added colors are
$1,2,4,7$ and $[4s+4,4s+7]$.
Together with the mapped old palette these give
$[1,4s+7]\setminus\{3\}$. The target hole is $2s+5$,
the new endpoints are $3s+6,3s+7$, and the new controlled edge
is $(2,2s+7)$. All edges cross $2s+3$.
The replacement uses fresh entries once and preserves the path topology.
This is the induction $s\mapsto s+2$; the two bases cover both
parities. Finally $1\le2\le3s-2$ permits completion.
\end{proof}

\begin{theorem}[Complete positive odd classification]\label{ok:all-odd-theorem}
A positive three-arm spider with all arm lengths odd is root-zero
$\alpha$ if and only if its profile is not $(3,3,3)$.
\end{theorem}
\begin{proof}
Order its lengths $a\le b\le c$.
First let $a\ge7$ and write $a=2s+1$, $b=2t+1$, $c=2u+1$.
The complete central cone $u\le2s$, complete short wedge
$t\ge2s$, and far cone $u\ge s+t$ have already been proved.
Outside their union,
\[
 u\ge2s+1,\qquad t\le2s-1,\qquad u\le s+t-1.
 \]
Put $g=u-2s$ and $h=s+t-u$. Then $g,h\ge1$ and
$t=s+g+h\le2s-1$, or $s\ge g+h+1$.
For $h\ge2$ use \eqref{ok:transition-high}.
For $h=1$ and $s\ge g+3$ use \eqref{ok:transition-one}.
The only remaining case is $h=1$, $s=g+2$, giving
$(a,b,c)=(2s+1,4s-1,6s-3)$; it is
Lemma~\ref{ok:pinned-ap-ray}. These alternatives exhaust every
positive size, without an upper bound or a total-size congruence.
For $a=5$, the equal profile $(5,5,5)$ is
Proposition~\ref{ok:equal-odd-three}; every other profile has
$c\ge7$ and is supplied by the minimum-five theorem in
Section~\ref{app:small-alpha}.

If $a=1$, the rooted odd pair $\Sp{b,c}$ is root-zero
$\alpha$ by Corollary~\ref{common:alpha-pair}; its only exceptional
rooted pair $(2,2)$ cannot occur. Adding a maximum-labeled unit
leaf preserves its index and proves the assertion.
If $a=3$ and $c\ge7$, use the complete minimum-three theorem
in Section~\ref{app:small-alpha}. If $c<7$, the two positive
profiles have the following explicit words:
\begin{align*}
 \Sp{3,3,5}:&
 \quad(11,1,10)\mid(6,2,9)\mid(8,3,5,4,7),\\
 \Sp{3,5,5}:&
 \quad(13,1,12)\mid(10,4,6,5,8)\mid(9,2,7,3,11).
\end{align*}
They have cuts $4,5$, respectively, consecutive labels, and
consecutive color inventories. The sole remaining profile is
$(3,3,3)$.

For this negative case a root-zero $\alpha$-labeling would have
cut $3$. Its depth-two labels would be $1,2,3$ and its first and
last labels $4,\ldots,9$. Equation~\eqref{ok:signed-leaf-sum}
at $K=3,N=9$ forces the leaf labels to sum to $21$.
Label $9$ must be a first label, since only an edge incident with
zero can have color $9$. The only three-element subset of
$[4,8]$ with sum $21$ is $\{6,7,8\}$; hence the first labels
are $4,5,9$.
Color $8$ forces $9$--$1$. Color $7$ then forces $1$--$8$:
the alternative $9$--$2$ would give degree three to the vertex
labeled $9$, while $7,8$ cannot be first labels.
Thus one entire arm is $(9,1,8)$.
The arm starting with $4$ cannot use middle $2$, because its final
edge to $6$ or $7$ would repeat root color $4$ or $5$.
It must therefore be $(4,3,6)$, since leaf $7$ would again repeat
color $4$. The last arm is then $(5,2,7)$ and repeats colors
$3,5$. This contradiction proves the exception.
The conclusion concerns root-zero $\alpha$ labeling; it does not
assert that $\Sp{3,3,3}$ is nongraceful.
\end{proof}

\subsection{Odd-pair offsets extracted from kernels}
\label{construction:odd-pair-kernels}

The basic pair offsets, including the low, high, reverse, equal, and
maximal constructions, are proved in Section~\ref{app:pair-offsets}.
The purpose of this section is to prove the additional odd-pair
intervals used in the eight-arm argument. Each statement is sufficient;
their union is not asserted to classify every possible offset.

\begin{lemma}[Kernel-to-pair extraction]\label{ok:pair-extraction}
Let $p$ be $K(r,u,t,V)$, with $t,u\ge1$ and $1\le V\le u$.
Then
\[
 k=t+u+V-1\ \in\ \Om(\Sp{2t+1,2u+1}).
 \]
No ordering assumption between $t$ and $u$ is needed.
\end{lemma}
\begin{proof}
Put $m=2t+2u+2$, $H=t+3-V$, and $\ell=u-V+1$.
Complement the kernel entries in $2t+3$. Its lower set is
$[1,t+2]\setminus\{H\}$, its upper set $[t+3,2t+2]$,
and its endpoints are $r+1,u-2r+V$.
Map its lower entries by $x\mapsto x+k$ and upper entries by
$x\mapsto x+k+2u$ to form arm $A$.
The even path on $[1,2u]$ from lower $\ell$ to upper
$u+\ell$ exists since $1\le\ell\le u$.
Translate it by $k+t+2$ and append $k+H$ to form arm $B$.

The two words have lengths $2t+1,2u+1$. Their nonroot labels
partition the following blocks, the indicated hole being filled last:
\begin{align*}
 &[k+1,k+t+2]\setminus\{k+H\},\\
 &[k+t+3,k+t+2u+2],\\
 &[k+t+2u+3,k+m],\qquad\{k+H\}.
\end{align*}
Thus their union is exactly $[k+1,k+m]$.
Complementation preserves the old differences, so the internal
colors of $A$ are
$[2u+1,m-1]\setminus\{D+2u\}$.
Its root color is $k+r+1=D+2u$.
The old even path in $B$ supplies $[1,2u-1]$,
its appended edge has color
\[
 (k+t+2+u+\ell)-(k+H)=2u,
 \]
and its root color is $k+t+2+\ell=m$.
Hence every color in $[1,m]$ appears once.
\end{proof}

\subsubsection{An extended interval for two equal odd arms}

\begin{lemma}[A variable-hole repeated kernel]\label{ok:variable-hole}
For $s\ge2$ and $1\le V\le s$, some integer $r\ge0$ admits
$K(r,s,s,V)$.
\end{lemma}
\begin{proof}
Normally put $r=V-1$, so the desired hole parameter is $r+1$.
For $r=0$, use $Z_s$. For $r=1$, the words in
\eqref{ok:B-small} apply at $s=3,4$, and carrier prefix of
size $s-3\ge2$ from $B_3$ applies at $s\ge5$.
For $r\ge2$, diagonal growth from $Z_{s-r}$ gives the standard
kernel if $s\le2r-2$, and carrier prefix from
$K(r,r,r,r+1)$ gives it if $s\ge2r+1$.
Here $r=V-1\le s-1$, so the diagonal source has positive gap
whenever it is used.

The two omitted carriers are $s=2(V-1)-1$ and $s=2(V-1)$.
In these cases change to $r=V-2$.
If $r=0$, necessarily $(s,V)=(2,2)$, and use
\[
 (6,1,5,2,3)=K(0,2,2,2).
 \]
If $r\ge1$, put $S=s-r+1$.
The two carrier alternatives give $S\ge3$ and
$S\le s\le2S-3$.
The $\epsilon=0$ odd-prefix source
of Lemma~\ref{ok:odd-prefix-boundaries} has parameters
$K(S,2S-1,s,r+2)$.
Reversing it by $J$ gives $K(r,s,s,r+2)$,
which has the desired $V$.
This repairs both excluded carriers without assuming a standard
kernel whose endpoints would coincide or hit its hole.
\end{proof}

\begin{proposition}[Equal odd upper interval]\label{ok:equal-odd-offsets}
For every odd $L=2s+1\ge5$,
\[
 [0,3s-1]\ \subseteq\ \Om(\Sp{L,L}).
 \]
\end{proposition}
\begin{proof}
The basic equal-arm construction in Section~\ref{app:pair-offsets}
supplies $0\le k<L$. It remains to cover
$2s\le k\le3s-1$, an interval overlapping the basic one.
Set $V=k-2s+1$, so $1\le V\le s$.
Lemma~\ref{ok:variable-hole} supplies $K(r,s,s,V)$,
and extraction gives $k=s+s+V-1$.
All endpoint and label/color conditions are exactly those of
Lemma~\ref{ok:pair-extraction}; no additional path request is hidden.
\end{proof}

\subsubsection{Consecutive odd lengths}

\begin{lemma}[Consecutive-carrier kernels]\label{ok:consecutive-kernels}
For every $t\ge3$ and $1\le V\le t+1$, there is some $r$
and a kernel $K(r,t+1,t,V)$.
\end{lemma}
\begin{proof}
The fifteen initial words in Table~\ref{ok:consecutive-table}
cover every required $V$ at $t=3,4,5$.
For a row with parameters $t,V,r$, the missing label is $t+V$,
the missing color is $r+V-1$, and the two ordered endpoints are
$2t-r+2,t+2r-V+2$.
Listing the entries and subtracting successive pairs gives precisely
those inventories. Thus the table consists of full literal kernels,
rather than declarations of a finite search outcome.

For $t\ge6$ first note the following ordinary source. For
$T\ge1$, $T\ne2$, take the odd path on $[1,2T+1]$ with upper
endpoints $2T,T+2$, and raise upper labels by one.
Its unique forbidden request would satisfy $2(2T-1)=3T$,
namely $T=2$. Thus this gives
\begin{equation}
 K(1,T+1,T,1)\qquad(T\ge1,\ T\ne2).
 \label{ok:consecutive-ordinary-seed}
\end{equation}
The literal source
\[
 M=(6,1,8,2,4,3,7)=K(2,4,3,2)
 \]
will repair the missing $T=2$ case where it arises.

At $V=1$, use \eqref{ok:consecutive-ordinary-seed} with $T=t$.
At $V=2$, use carrier prefix of size $w=t-3\ge3$ on $M$.
Reflection preserves $U=t+1$ and sends $V$ to $t-V+3$;
therefore this also gives $V=t+1$.

For $3\le V\le t$ and $2V\ge t+3$, put
$T=t-V+1$ and $P=V-1$.
Then $T\ge1$ and $P\ge T+1$.
If $V\ne t-1$, then $T\ne2$ and diagonal growth of
\eqref{ok:consecutive-ordinary-seed} gives the desired kernel:
its new carrier is $T+P=t$, its other half is $t+1$,
and its hole parameter is $1+P=V$.
If $V=t-1$, diagonally grow $M$ by $P=t-3\ge3$ instead.
Its gap $T-S$ is $1$, so this growth is permitted and gives
$U=t+1$, carrier $t$, and $V=t-1$.
Finally, every remaining $3\le V\le t$ has $2V<t+3$;
reflect the already constructed upper-half kernel at
$V'=t-V+3$. It has $3\le V'\le t$ and $2V'\ge t+3$.
All hole parameters have now been covered.
\end{proof}

\begin{longtable}{@{}rrr@{\quad}l@{}}
\caption{Literal consecutive-carrier kernels.}
\label{ok:consecutive-table}\\
\toprule
$t$&$V$&$r$&Word\\ \midrule
\endfirsthead
\toprule $t$&$V$&$r$&Word\\ \midrule
\endhead
\bottomrule
\endfoot
$3$&$1$&$1$&$(7,2,8,1,5,3,6)$\\
$3$&$2$&$2$&$(6,1,8,2,4,3,7)$\\
$3$&$3$&$1$&$(7,2,8,1,5,3,4)$\\
$3$&$4$&$2$&$(6,2,8,1,4,3,5)$\\
$4$&$1$&$1$&$(9,1,10,3,6,4,8,2,7)$\\
$4$&$2$&$3$&$(7,4,5,3,8,2,9,1,10)$\\
$4$&$3$&$1$&$(9,3,8,1,10,2,6,4,5)$\\
$4$&$4$&$4$&$(6,4,5,2,7,3,9,1,10)$\\
$4$&$5$&$2$&$(8,1,10,2,7,3,6,4,5)$\\
$5$&$1$&$1$&$(11,2,12,1,9,3,10,5,7,4,8)$\\
$5$&$2$&$2$&$(10,2,11,1,12,5,6,4,8,3,9)$\\
$5$&$3$&$1$&$(11,2,12,1,9,4,10,3,7,5,6)$\\
$5$&$4$&$2$&$(10,4,8,1,12,2,11,3,6,5,7)$\\
$5$&$5$&$3$&$(9,1,12,2,11,5,6,4,7,3,8)$\\
$5$&$6$&$3$&$(9,4,8,2,12,1,10,3,6,5,7)$\\
\end{longtable}

\begin{proposition}[Consecutive odd upper interval]\label{ok:consecutive-offsets}
For odd $a\ge7$,
\[
 [a,(3a-1)/2]\ \subseteq\ \Om(\Sp{a,a+2}).
 \]
\end{proposition}
\begin{proof}
Write $a=2t+1$, $u=t+1$. The preceding lemma supplies every
$K(r,u,t,V)$ with $1\le V\le t+1=u$.
Extraction gives $k=2t+V$, ranging from $a$ to $3t+1$.
\end{proof}

\subsubsection{Two bands from the central and short-wedge kernels}

\begin{proposition}[Short-wedge pair band]\label{ok:short-pair-band}
Let $1\le t\le u$ and put
\[
 S=k-2u-2,\qquad\delta=u-t,\qquad V=k-t-u+1.
 \]
If
\[
 S\ge2,\quad t\ge2S+2,\quad
 0\le\delta\le S-1,\quad\delta\ne S-3,\quad1\le V\le t,
 \]
then $k\in\Om(\Sp{2t+1,2u+1})$.
\end{proposition}
\begin{proof}
Lemma~\ref{ok:unequal-wedge-kernel}, in its nonexceptional branch,
gives $K(S,t,u,S+\delta+3)$.
Here $S+\delta+3=k-t-u+1=V$.
Its ordinary odd seed has carrier $r=S+\delta+1$ and upper
first entry $2r-S+2$; the explicitly excluded request is
exactly $\delta=S-3$.
Its even carrier prefix has half-order $w=t-S-1\ge S+1$.
Thus all the source hypotheses have been checked, not inferred
from a general three-arm $\alpha$ theorem.
Apply extraction with carrier $u$ and other half $t$.
The two lengths are produced in reversed order, and its hypothesis
is precisely $1\le V\le t$. The offset is $u+t+V-1=k$.
\end{proof}

\begin{proposition}[Central pair band]\label{ok:central-pair-band}
Let $1\le t\le u$ and put $S=k-2t$, $V=k-t-u+1$.
If
\[
 1\le S\le t\le u\le2S-2,\qquad
 u>S,\quad u\ne2t-S,\quad1\le V\le u,
 \]
then $k\in\Om(\Sp{2t+1,2u+1})$.
\end{proposition}
\begin{proof}
The source of Lemma~\ref{ok:central-kernel} has
\[
 S_0=u-t,\quad T_0=u-S,\quad
 U_0=2u-S-t,\quad P=S+t-u.
 \]
The stated inequalities give $0\le S_0\le T_0$,
$T_0,U_0\ge1$, and $P\ge T_0-S_0+2$.
Its odd-path exception is exactly $u=2t-S$, which has been
excluded, and diagonal growth gives
$K(S,u,t,P+1)=K(S,u,t,V)$.
Extraction now gives the desired pair in its original order and
offset $t+u+V-1=k$. The extra $u>S$ avoids the equal zero-carrier
source; no AP or equal exception is silently included.
\end{proof}

\subsubsection{Asymmetric, strict one-layer, and reflected bands}

\begin{proposition}[Asymmetric pair bands]\label{ok:asymmetric-pair-bands}
Let $1\le t\le u$ and $V=k-t-u+1$.
Each of the following conditions suffices for
$k\in\Om(\Sp{2t+1,2u+1})$.
\begin{enumerate}
\item Put $s=k-2t$, $g=u-2s$, and $h=k-t-u$.
Require
\[
 g\ge1,\quad h\ge2,\quad s\ge g+h+1,\quad1\le V\le u.
 \]
\item Put $s=u-t+1$ and $g=2t-u-2$.
Require
\[
 k=t+u+2,\quad g\ge1,\quad s\ge g+3,\quad u\ge3.
 \]
\item Put
\[
 s=3u-k+3,\quad h=t+2u-k+1,\quad
 g=2k-4u-t-4.
 \]
Require
\[
 g\ge1,\quad h\ge2,\quad s\ge g+h+1,\quad1\le V\le t.
 \]
\end{enumerate}
\end{proposition}
\begin{proof}
For the first case substitute the displayed parameters into
\eqref{ok:transition-high}. They satisfy
\[
 2s+g=u,\qquad s+g+h=t,\qquad h+1=V.
 \]
Thus the kernel is $K(s,u,t,V)$ and extraction gives the claim.

For the second case substitute into \eqref{ok:transition-one}.
Again $2s+g=u$ and $s+g+1=t$, while $V=3$.
The required completion/extraction inequality is $3\le u$,
which is explicitly assumed. This is the strict $h=1$ source:
the limiting AP ray is not included by weakening $s\ge g+3$.

For the third case the source \eqref{ok:transition-high} has
\[
 T=s+g+h=u,\qquad U=2s+g=2u-t+2,\qquad W=h+1.
 \]
Reflect it using $R$. The resulting carrier remains $u$, its
other half is $2u-U+2=t$, and its hole parameter is
\[
 u-W+3=u-h+2=k-t-u+1=V.
 \]
The new shortest parameter is $u-s+1$; its value is already
realizable because it comes from a genuine reflected word.
Extraction is used in reversed arm order, and therefore needs
$1\le V\le t$, not merely $V\le u$.
This essential guard is retained in the statement.
\end{proof}

\subsection{Proof dependencies and computational verification}
\label{ok:verification}

The infinite arguments of this section consist of interval
partitions, the two classical prescribed-path interfaces, and the
controlled two-step recurrences. No finite search is used to infer
the existence of a path for arbitrarily large parameters.
Every literal initial word required by these arguments is printed
above. The negative $(3,3,3)$ profile is excluded by the degree and
color argument in Theorem~\ref{ok:all-odd-theorem}.

The proofs
of the graph constructions are the formulas and inventory arguments
in this section; they are not claimed to be formalized in a proof
assistant.

\section{Equal alpha triples}\label{app:equal-alpha}
\begin{theorem}\label{ea:all}
For every integer $L\ge5$, $\Sp{L,L,L}$ has a root-zero alpha labeling.
\end{theorem}
\begin{proof}
The odd lengths have the explicit three-word construction of Proposition~\ref{ok:equal-odd-three}, including the literal lengths five and seven and the full interval inventories for both parities of $(L-1)/2$. It remains to prove the even lengths.

Let $L=2r$, with $r\ge5$ odd. Start from that equal-odd construction for length $2r+1$, in arm order Walecki, A, C. Delete the three leaves, whose labels are $5r+3,5r+1,5r+2$. Decrease remaining labels above $5r+3$ by three. The labels are now $[0,6r]$ with cut $3r$. The Walecki colors are $[4r+1,6r]$. The two other arms formerly supplied $[1,4r+2]$, and their deleted leaf colors are $4r,2r+4$. Thus the shortened source has duplicate colors $4r+1,4r+2$ and holes $4r,2r+4$.

Apply the class-preserving cycles
\[
(r-1,\ r+2,\ 3r-2,\ r+1,\ r),\qquad
(4r-1,\ 4r+2,\ 4r+1,\ 4r+3),
\]
where each entry maps to the next cyclically. The resulting source has signed excess
\[
+\{4r,2r+2,2r+1,2r+1\}
-\{4,2r+5,4r-3,4r-2\}.
\]
The complete fifteen affected incident-edge calculations, both residue-specific repair tables, and all their source coordinates are printed in Section~\ref{templates:equal-even}. For $r=4q+1$ and $r=4q+3$, $q\ge3$, the respective sixteen-edge and fifteen-edge trades cancel exactly this multiset, preserve all vertex degrees and alpha classes, and have a contracted topology consisting of three hub paths, each of affine length $2r$. Theorem~\ref{templates:soundness} proves these are labelings for every parameter in the stated half-lines. The four smaller values $r=5,7,9,11$ are supplied by the complete literal words in Section~\ref{app:auxiliary-finite}. Hence all $L\equiv2\pmod4$, $L\ge10$, are proved.

For the remaining even lengths the starting words are
\[
\begin{array}{c|l}
6&(18,1,17,5,16,2)\\
 &(15,6,12,7,14,4)\\
 &(13,9,10,8,11,3)\\[2pt]
8&(24,1,23,2,22,6,21,3)\\
 &(19,7,20,9,16,8,18,4)\\
 &(17,11,15,10,13,12,14,5).
\end{array}
\]
Each has the full label and color inventory and balanced index $3L/2$. Repeatedly halve an even $L\ge6$ while it is divisible by four, stopping at $6$ or $8$ if encountered. This process reaches $6$, $8$, or a value congruent to two modulo four and at least ten. Apply the balanced alpha subdivision theorem of Section~\ref{app:lifts} in reverse order to double those lengths back to $L$. The balanced hypothesis persists because all three arms are even. This proves every even length and the theorem.
\end{proof}

The corresponding root-zero alpha four-arm result is proved by its full seed formula, defect inventory, four explicit affine repair tables and twenty-one literal small bases in Section~\ref{templates:four-odd}. It is used only for four equal odd lengths at least seven. Ordinary equal-arm offsets, in contrast, are available for every common positive length and every number of arms by subdivision of a star; alpha separation is an additional conclusion requiring the separate constructions above.

\section{Mixed-parity alpha constructions}
\label{mc:chapter}
\label{app:ten-alpha}

This section supplies the mixed-parity constructions used in the
seven-arm arithmetic cover and the even-minimum eight-arm theorem.
All arm lengths count edges, and all displayed words omit their common
initial vertex, whose label is zero.  Intervals in formulas for labels
or colors are integer intervals.  The prescribed-endpoint path inputs
are exactly Lemmas~\ref{common:even-path} and~\ref{common:odd-path}.
Their existence is classical; all subsequent operations and inventories
are proved here.  In particular, an arbitrary alpha labeling is never
assumed to have a prescribed Walecki arm.

\subsection{A shortening operation with a prescribed arm}
\label{mc:shortening}

\begin{lemma}[Maximum-Walecki shortening]\label{mc:shorten}
Suppose a root-zero alpha labeling of a spider with $N$ edges contains
the arm
\[
 W=(N,1,N-1,2,\ldots,N-s)
\]
of length $2s+1$, and every label on the other arms is below $N-s$.
Deleting the last leaf of $W$ and subtracting one from every remaining
label above $N-s$ gives a root-zero alpha labeling with that arm
shortened to $2s$.  The alpha index is unchanged.
\end{lemma}
\begin{proof}
Write $h=N-s$.  Label compression is a bijection onto
$\interval{0}{N-1}$.  Only labels on $W$ are changed.  The deleted
edge has color $N-2s$; the remaining colors of $W$, formerly
$\interval{N-2s+1}{N}$, each decrease by one.  They now give
$\interval{N-2s}{N-1}$, adjoining the unchanged colors
$\interval{1}{N-2s-1}$ on the other arms.  Every surviving upper
label on $W$ is at least $h$, still above the old cut.  Lower labels
and the root are unchanged, so every edge still crosses that cut.
\end{proof}

\subsection{Odd--even--odd triples and their shortening}
\label{mc:oeo}

\begin{proposition}\label{mc:oeo-central}
Let $a<b<c$ have respective parities odd, even, odd.  If
$c\le2a-3$, then $\Sp{a,b,c}$ has a root-zero alpha labeling
with its shortest arm equal to the maximum Walecki word.
\end{proposition}
\begin{proof}
Write $(a,b,c)=(2s+1,2t,2u+1)$, so
$1\le s<t\le u\le2s-1$.  Put
\[
 S=u-t,\qquad T=u-s,\qquad P=s+t-u,\qquad v=P+1.
\]
Thus $0\le S<T$, $P\ge T-S+1$, and $1\le v\le u$.
Take an even-order alpha path on $\interval{1}{2T}$ from upper
endpoint $2T-S$ to lower endpoint $T-S$.  Raise its uppers by one.
It now omits label $T+1$ and has colors $\interval{2}{2T}$.
Raise its uppers by a further $2P$.  Before it place an even-order
alpha path on $\interval{1}{2P}$ from $P+T-S+1$ to $T-S+1$,
translated by $T$.  Both requests are legal since
$1\le T-S+1\le P$.  The join has color $2P$.

The resulting path $Q$ has lower labels $\interval{1}{t}$,
upper labels $\interval{t+1}{2t+1}\setminus\{t+v\}$, and colors
\[
 \interval{1}{2t}\setminus\{2P+1\}.
\]
Its first and last labels are $E=2t-s+1$ and $F=t-s$.
These assertions follow from the adjacent prefix intervals: the
prefix contributes colors $\interval{1}{2P-1}$ and the old path
contributes $\interval{2P+2}{2t}$, besides their joining color.

Map the lowers of $Q$ by $x\mapsto x+s$ and the uppers by
$x\mapsto x+s+2u$.  Call this arm $A$.  Its internal colors are
$\interval{2u+1}{2u+2t}\setminus\{2u+2P+1\}$ and its root
color is $2u+2t+1$.
For arm $B$, take an even-order alpha path on
$\interval{1}{2u}$ from $u+v$ to $v$, append $2u+v$, and
translate every label by $s+t$.  Its internal colors are
$\interval{1}{2u}$, and its root color is
$s+t+u+v=2u+2P+1$.

The lower labels of $A$ and $B$ are respectively
$\interval{s+1}{s+t}$ and $\interval{s+t+1}{s+t+u}$.
The upper block of $B$ is $\interval{s+t+u+1}{s+t+2u}$;
that of $A$ is $\interval{s+t+2u+1}{s+2t+2u+1}$ with the
single hole filled by the appended leaf of $B$.
Set $N=2s+2t+2u+2$ and adjoin
$W=(N,1,N-1,2,\ldots,N-s)$.  Its labels fill
$\interval{1}{s}\cup\interval{N-s}{N}$, and its colors are
$\interval{2t+2u+2}{N}$.  All labels and colors are now consecutive
without repetition.  All edges cross $s+t+u$, and every label off
$W$ is below its last leaf.  This proves the pinning as well.
\end{proof}

\begin{corollary}\label{mc:eeo-central}
If $a,b$ are even, $6\le a<b<c$, $c$ is odd, and
$c\le2a-1$, then $\Sp{a,b,c}$ is root-zero alpha.
\end{corollary}
\begin{proof}
Apply Proposition~\ref{mc:oeo-central} to $(a+1,b,c)$:
the parity inequalities imply $a+1<b<c$, and
$c\le2(a+1)-3$.  Its explicitly prescribed shortest arm can be
shortened by Lemma~\ref{mc:shorten}.  This proves exactly the bound
$c\le2a-1$; a larger bound is not inferred from this operation.
\end{proof}

\subsection{Even--odd--odd central and repeated triples}
\label{mc:eoo-central-section}

We use the upper-hole kernels $K(S,U,T,V)$ and their completion from
Definition~\ref{ok:kernel-definition} and Lemma~\ref{ok:completion}.
For clarity, their final completion has labels
\[
 \interval{1}{T}\ \cup\
 \bigl(\interval{T+1}{2T+2}\setminus\{T+V\}\bigr),
\]
missing color $S+T-U+V$, and ordered endpoints
$2T-S+2$, $2T+2S-U-V+3$.  Its three completed arm lengths are
$(2S+1,2T+1,2U+1)$.  That operation explicitly supplies the
maximum-Walecki shortest arm.  The auxiliary source kernels used in
growth need not themselves satisfy the bounds required for a final
completion.

\begin{lemma}[Pinned central AP kernels]\label{mc:eoo-ap}
For $d\ge1$, $s\ge2d+2$, the spider
\[
 \Sp{2s,2s+2d+1,2s+4d+1}
\]
is root-zero alpha.  For $d=1$ the stronger range $s\ge2$ holds.
\end{lemma}
\begin{proof}
For $d\ge2$, the following explicit word is a kernel
$K(d,3d,2d,2)$:
\[
\begin{split}
 &(3d+2),\\
 &(d+1-j,3d+3+j)\quad(0\le j\le d-1),\\
 &(1,2d+1),\\
 &(2d+1-j,2d+2+j)\quad(1\le j\le d-1).
\end{split}
\]
Its lowers partition $\interval{1}{2d}$; its uppers partition
$\interval{2d+1}{4d+2}\setminus\{2d+2\}$.  The endpoints
are $3d+2,3d+1$.  The word through its first label $1$ gives
colors $\interval{2d+1}{4d+1}$, the next edge gives $2d$, and
the final block, including its join, gives
$\interval{1}{2d-1}\setminus\{2\}$.  Thus its missing color
is precisely $2$.  The restriction $d\ge2$ excludes an endpoint
equal to the label hole.

Apply diagonal growth (Lemma~\ref{ok:diagonal-growth}) with
$P=s-d\ge d+2$.  In explicit terms its even-order path has upper
endpoint $P+d+2$ and lower endpoint $d+2$, is translated by $2d$,
and is put before the kernel after raising the latter's uppers by
$2P$.  The resulting parameters are
\[
 (S,U,T,V)=(s,s+2d,s+d,s-d+2).
\]
Here $1\le V\le U$.  Complete this kernel, then shorten its
prescribed Walecki arm by Lemma~\ref{mc:shorten}.  The resulting
index is $3s+3d$ and the arm lengths are those claimed.

For $d=1$, use these three literal kernels, in each case with
$T=s+1$ and $U=s+2$:
\begin{center}\begin{tabular}{ccp{0.72\textwidth}}
$s$&$V$&word\\\midrule
2&2&$(6,1,8,2,4,3,7)$\\
3&2&$(7,4,5,3,8,2,9,1,10)$\\
4&4&$(8,2,12,1,10,5,6,4,7,3,11)$.
\end{tabular}\end{center}
Direct subtraction checks the full kernel inventories.  For every
$s\ge5$, grow the first kernel by $P=s-2\ge3$, using even-path
endpoints $P+3,3$.  This gives parameters $(s,s+2,s+1,s)$.
Completion and shortening apply in all cases.
\end{proof}

\begin{proposition}\label{mc:eoo-central}
If $a\ge6$ is even, $b,c$ are odd, and
$a<b\le c\le2a+1$, then $\Sp{a,b,c}$ is root-zero alpha.
\end{proposition}
\begin{proof}
First assume $c\le2a-3$.  The odd source $(a+1,b,c)$ lies
in the pinned odd central construction of
Lemma~\ref{ok:central-kernel}, apart from its AP plane.
At $a+1=b=c$ use the explicit equal-odd construction in
Proposition~\ref{ok:equal-odd-three}; at $a+1<b=c$ one may use
the pinned repeated kernels of Lemma~\ref{ok:repeated-kernels}.
On the strict AP plane $a+1+c=2b$, write the lengths as
$(2s,2s+2d+1,2s+4d+1)$; its central inequality is
$s\ge2d+2$, so Lemma~\ref{mc:eoo-ap} supplies the missing pin.
Shortening the stated arm proves this first range.

For the two remaining odd layers write $a=2s$, $b=2t+1$,
$s\ge3$.  At $c=4s-1$ the source largest half is $2s-1$.
If $t\le2s-3$, use the explicit odd-prefix boundary construction
of Lemma~\ref{ok:odd-prefix-boundaries}.  If $t=2s-2$,
$s\ge4$, use Lemma~\ref{ok:gap-two-kernels}.  The exceptional
source $(7,9,11)$ uses the $s=3,d=1$ literal kernel of
Lemma~\ref{mc:eoo-ap}, rather than an unpinned AP existence
statement.  At $t=2s-1$, the explicit last-ray repeated kernel
has $V=s$; its even completion endpoint is $3s-1$ on
$\interval{1}{4s-2}$, an allowed request.

At $c=4s+1$ the reversed boundary construction of
Lemma~\ref{ok:reverse-boundary} covers $s\le t\le2s$.
Its repeated endpoint $t=2s$ is the explicit twice-ray kernel with
$V=s+3$ and even completion endpoint $3s+3$ on
$\interval{1}{4s}$; this is allowed since $s\ge3$.
Every cited construction completes by actually adjoining the standard
maximum-Walecki arm.  Apply Lemma~\ref{mc:shorten} in each case.
\end{proof}

\begin{proposition}\label{mc:eoo-repeated}
For every even $a\ge6$ and odd $b>a$, $\Sp{a,b,b}$ is
root-zero alpha.
\end{proposition}
\begin{proof}
Write $a=2s,b=2t+1$.  The pinned all-odd repeated source has
$V=s+1$ when $s<t\le2s-2$ or $t\ge2s+1$, $V=s$ at
$t=2s-1$, and $V=s+3$ at $t=2s$.  At $t=s$ use the explicit
equal-odd construction with its first Walecki arm.  These are the
complete repeated constructions of Proposition~\ref{ok:equal-odd-three}
and Lemma~\ref{ok:repeated-kernels}, and exhaust $t\ge s$.
All completion requests $t+V$ lie in $\interval{t+1}{2t}$ in
their stated domains.  Shortening the pinned arm gives the result.
In particular $b=a+1$ is included without a separate inference
about pinning of arbitrary all-odd alpha labelings.
\end{proof}

\subsection{Equal-even near triples}
\label{mc:eeo-near-section}

\begin{proposition}\label{mc:eeo-near}
For every even $a\ge6$, $\Sp{a,a,a+1}$ is root-zero alpha.
\end{proposition}
\begin{proof}
There is a useful endpoint operation.  In an $N$-edge root-zero
alpha labeling, suppose an even arm has first label $H$, last label
$l$, and $H+l=N+1$.  Replace its word $P$ by
$(N+1,\operatorname{rev}(P))$.  Reversal retains all old internal
colors.  The old root color $H$ is replaced by the joining color
$(N+1)-l=H$, and the new root supplies color $N+1$.
Alternation and the cut are retained, since the reversed even word
starts lower and ends upper.  This proves the operation.

Write $a=2r$.  If $r\ge7$ is odd, the appropriate equal-even construction in
Section~\ref{templates:equal-even} has complementary endpoints
$H=5r+1,l=r$ on one arm.  For $r\equiv3\pmod4$, this is the
E recipe of Section~\ref{templates:equal-r3}, with its finite
evaluations at $r=7,11$.  At $r=3$ the arm
$(15,6,12,7,14,4)$ has endpoint sum $19$.  For
$r\equiv1\pmod4$, $r\ge13$, use the alternative recipe of
Section~\ref{templates:alt}; its contracted component beginning with
$0,(5r+1)$ has length $2r$ and ends at $r$.
The same finite template is evaluated and checked at $r=9$.
The complete edge lists, original positions of their endpoints,
color-multiset cancellations, and finite weighted-component checks
are part of Section~\ref{templates:equal-even}; hence this endpoint
assertion is more than a degree or color calculation.  Apply the
operation just proved.

When $a$ is divisible by four and $a\ge12$, an equal-even source
with arms $a/2$ has balanced alpha index.  Corollary~\ref{common:balanced-alpha-matching} gives $\Sp{a,a,a+1}$:
its unique-odd-arm perfect matching has contree $\Sp{a/2,a/2,a/2}$
by Corollary~\ref{common:unique-odd-contree}.  The remaining targets $a=8,10$ have these
literal words, with cuts $12,15$, respectively:
\[
\begin{array}{l}
(25,1,24,2,23,3,21,4);\\
(19,6,20,8,17,10,18,7,22);\\
(16,11,14,12,13,9,15,5),
\end{array}
\qquad
\begin{array}{l}
(31,1,30,2,29,3,28,4,26,5);\\
(23,7,24,11,22,8,18,13,21,9,27);\\
(20,14,17,15,16,12,19,10,25,6).
\end{array}
\]
Their full label, color, and separation inventories are checked by
subtraction.  For completeness, the finite $r=9$ alternative evaluation
gives the following $a=18$ target with cut $27$:
\[
\begin{split}
 &(54,1,53,2,52,3,51,4,42,14,50,5,49,6,48,7,47,8);\\
 &(37,25,29,26,28,27,32,24,30,23,34,13,43,12,44,11,45,10);\\
 &(55,9,38,18,35,17,39,16,40,15,41,22,31,21,36,20,33,19,46).
\end{split}
\]
The cases just listed exhaust every even $a\ge6$.
\end{proof}

\subsection{Lower-hole kernels for two even long arms}
\label{mc:lower-kernels}

\begin{definition}\label{mc:lower-kernel}
A lower-hole kernel $L(s,u,t,H)$ is an alternating word of $2t$
vertices, starting upper and ending lower, with
\[
\begin{array}{ll}
\text{lower labels:}&\interval{1}{t+1}\setminus\{H\},\\
\text{upper labels:}&\interval{t+2}{2t+1},\\
\text{colors:}&\interval{1}{2t}\setminus\{D\},
 \quad D=s+2t-u+3-H,\\
\text{endpoints:}&E=2t-s+1,\quad F=t+u-2s-H.
\end{array}
\]
Only the indicated inventory conditions are imposed on auxiliary
source kernels.  The next lemma states the additional hypotheses
required to complete one to a spider.
\end{definition}

\begin{lemma}[Lower-hole completion]\label{mc:lower-completion}
Let $s\ge1$, $t\ge s$, $u\ge t$, and $1\le H\le t+1$.
Suppose $L(s,u,t,H)$ exists and
\[
 X=u+t-H+1\in\interval{u}{2u-1},\qquad
 2(X-1)\ne3(u-1).
\]
Then $\Sp{2s+1,2t,2u}$ is root-zero alpha with cut $s+t+u$,
and its shortest arm is the standard maximum Walecki word.
\end{lemma}
\begin{proof}
Map kernel lowers by $x\mapsto x+s$ and uppers by
$x\mapsto x+s+2u-1$, producing arm $A$.
Take the ordinary odd alpha path on $\interval{1}{2u-1}$
from upper endpoint $X$ to upper endpoint
$Y=3u-1-X=2u-t+H-2$.  Translate it by $s+t+1$ and append
the lower leaf $s+H$, obtaining arm $B$.

The internal colors of $A$ are
$\interval{2u}{2u+2t-1}\setminus\{R\}$, where
$R=D+2u-1=s+u+2t-H+2$.  Its root color is $2u+2t$.
The old internal colors of $B$ are $\interval{1}{2u-2}$,
its appended edge has color $Y+t+1-H=2u-1$, and its root color
is $X+s+t+1=R$.

For labels, $A$ supplies
\[
 (\interval{s+1}{s+t+1}\setminus\{s+H\})
 \ \cup\ \interval{s+t+2u+1}{s+2t+2u}.
\]
The lowers of $B$ fill the omitted $s+H$ and
$\interval{s+t+2}{s+t+u}$; its uppers fill
$\interval{s+t+u+1}{s+t+2u}$.
Thus $A,B$ partition $\interval{s+1}{s+2t+2u}$.
Set $N=2s+2t+2u+1$ and add
$W=(N,1,N-1,2,\ldots,N-s)$.  It supplies the two outer label
intervals and the color interval $\interval{2t+2u+1}{N}$.
All labels and colors occur once, and every edge crosses $s+t+u$.
The required strict bound on every label off $W$ also follows.
\end{proof}

\begin{lemma}[Diagonal growth of lower-hole kernels]
\label{mc:lower-growth}
Let $L(s_0,u_0,t_0,H)$ exist, put $g=t_0-s_0\ge1$, and
choose $P\ge g$.  Then $L(s_0+P,u_0+P,t_0+P,H)$ exists.
\end{lemma}
\begin{proof}
Take an even-order alpha path on $\interval{1}{2P}$ from $P+g$
to $g$, translate it by $t_0+1$, and prepend it to the old kernel
after raising the old uppers by $2P$.
Its last label is $t_0+1+g=2t_0-s_0+1$, and the mapped old
first label is $2P$ larger, so the join has color $2P$.
The colors are the new $\interval{1}{2P-1}$, that join, and
the old palette shifted by $2P$; the new missing color is $D_0+2P$.
The lower labels consist of the old lower interval with its hole
and $\interval{t_0+2}{t_0+P+1}$; the uppers form the two
adjacent intervals
\[
 \interval{t_0+P+2}{t_0+2P+1},\qquad
 \interval{t_0+2P+2}{2t_0+2P+1}.
\]
The new first endpoint is $E_0+P$ and its last is $F_0$.
These are exactly the required endpoints for the asserted parameters.
\end{proof}

\subsection{The complete distinct central OEE cone}
\label{mc:oee-central-section}

\begin{proposition}\label{mc:oee-central}
For positive odd $a$ and positive even $b,c$ satisfying
$a<b<c\le2a$, $\Sp{a,b,c}$ is root-zero alpha with its
shortest arm the standard maximum Walecki word.
\end{proposition}
\begin{proof}
Write $(a,b,c)=(2s+1,2t,2u)$, where
$1\le s<t<u\le2s+1$.
We first construct all points off the plane
\begin{equation}\label{mc:oee-plane}
 3u=2s+2t+3.
\end{equation}
Put $S=u-t-1$, $T=u-s-1$, $P=s+t-u+1$.
Then $0\le S<T$, $P\ge T-S\ge1$, $T+P=t$.
Take an even-order alpha path on $\interval{1}{2T}$
from $2T-S$ to $T-S$ and raise its uppers by $1+2P$.
Before it put an even-order alpha path on $\interval{1}{2P}$
from $P+T-S$ to $T-S$, translated by $T+1$.
The two endpoint requests are valid.  The join has color $2P$,
so the new word has missing color $D=2P+1$, lower hole
$H=T+1=u-s$, and endpoints $2t-s+1,t-s$.
Its full label classes are those of $L(s,u,t,H)$, since the
prefix supplies the intervening adjacent lower and upper blocks.

The completion request is $X=s+t+1$ on order $2u-1$.
Its endpoints $X$ and $3u-s-t-2$ lie in
$\interval{u}{2u-1}$.  The only forbidden upper request is
$2(X-1)=3(u-1)$, precisely \eqref{mc:oee-plane}.
Thus Lemma~\ref{mc:lower-completion} proves every off-plane point.

On the plane, first suppose $d=u-t\ge3$.  Set
\[
 R=s+t-u+2,\quad W=u-s-2,\quad A=t-s,\quad B=A+1.
\]
Here $R+W=t$, $1\le A\le R$ because $u\le2s+2$, and
$1\le B\le W$ because $d\ge3$.
Translate an even alpha path on $\interval{1}{2R}$ from
$R+A$ to $A$ by $W+1$.  Follow it with an even alpha path on
$\interval{1}{2W}$ from $W+B$ to $B$, keeping its lowers and
raising its uppers by $2R+1$.
The two lower intervals are $\interval{1}{W}$ and
$\interval{W+2}{t+1}$, so $H=W+1=u-s-1$.
The upper intervals adjoin to $\interval{t+2}{2t+1}$.
The internal color intervals are $\interval{1}{2R-1}$ and
$\interval{2R+2}{2t}$, while their join has color $2R+1$.
The missing color is $2R$, and the endpoints are
$2t-s+1$ and $t-s+1$, as required by $L(s,u,t,H)$.
The completion has $X=s+t+2$, $Y=3u-s-t-3$.
Both are upper; on \eqref{mc:oee-plane}, the forbidden equality
fails by two, since $2(X-1)=3u-1\ne3u-3$.
This fills every plane point with $d\ge3$.

For the remaining plane points put $j=s-(3d-2)$.  Then
\[
 s=3d-2+j,\quad t=3d-1+2j,\quad u=4d-1+2j,
 \qquad d\ge1,\ j\ge0.
\]
When $j\ge1$, $j\ne d+2,2d-2$, set
$T=d+j$, $P=2d-2+j$.
Take an odd alpha path on $\interval{1}{2T+1}$ with larger
class lower, from $d+2$ to $j$, and raise its uppers by $2P+2$.
These endpoints sum to $T+2$; its sole forbidden request is
$j=d+2$.  Before it put an odd alpha path on
$\interval{1}{2P+1}$ with larger class upper, from
$2d-1+2j$ to $4d-3+j$, translated by $T+2$.
Its endpoints sum to $3P+2$; its sole forbidden request is
$j=2d-2$.
The combined lower intervals are $\interval{1}{T+1}$ and
$\interval{T+3}{t+1}$, hence $H=T+2=u-s+1$.
The upper intervals adjoin to $\interval{t+2}{2t+1}$.
The two internal palettes are $\interval{1}{2P}$ and
$\interval{2P+3}{2t}$, and the join has color $2P+1$.
Thus the missing color is $2P+2$, with the required endpoints
$2t-s+1,j$.  The odd completion endpoints are
$6d-3+3j$ and $6d-1+3j$, both upper and distinct; their first
normalized endpoint differs by one from the forbidden midpoint.
There is no further exception.

For $d=1$ this leaves $j=0,3$; for $d=2$ it leaves $j=0,2,4$.
Four have the following literal lower-hole kernels:
\begin{center}\small
\begin{tabular}{ccl}
$(s,t,u)$&$H$&word\\\midrule
$(1,2,3)$&1&$(4,3,5,2)$\\
$(4,8,9)$&3&$(13,8,10,9,12,4,14,7,11,5,17,1,16,2,15,6)$\\
$(6,9,11)$&3&$(13,10,11,9,14,4,16,7,15,8,12,6,17,2,18,1,19,5)$\\
$(8,13,15)$&3&$(19,11,18,12,17,13,16,14,15,6,24,1,27,2,26,4,25,5,$\\
&&$22,8,23,7,20,10,21,9)$.
\end{tabular}\end{center}
Their completion orders and first labels are respectively
$(5,5),(17,15),(21,18),(29,26)$, all allowed.
The last point $(d,j)=(2,0)$ gives $(a,b,c)=(9,10,14)$;
its literal arm words, with cut $16$, are
\[
\begin{split}
 &(33,1,32,2,31,3,30,4,29);\\
 &(23,15,18,16,17,13,20,14,19,7);\\
 &(24,6,28,8,27,10,25,9,22,11,21,12,26,5).
\end{split}
\]
Full subtraction verifies every finite word above.  These branches
exhaust the plane, and all completions explicitly adjoin the required
Walecki arm.  Together with the off-plane construction they prove
the proposition, including its smallest instance $(3,4,6)$.
\end{proof}

\begin{corollary}\label{mc:eee-central}
Every all-even triple with $a<b<c\le2a+2$ is root-zero alpha.
\end{corollary}
\begin{proof}
The source $(a+1,b,c)$ satisfies the preceding proposition.
Shorten its prescribed Walecki arm by Lemma~\ref{mc:shorten}.
This includes $(2,4,6)$.
\end{proof}

\subsection{Repeated OEE triples}
\label{mc:oee-repeated-section}

\begin{lemma}[The adjacent repeated-even family]\label{mc:oee-near}
For even $b\ge6$, $\Sp{b-1,b,b}$ is root-zero alpha.
The shortest arm in this lemma is a modified Walecki word; the
standard-Walecki shortening lemma is not asserted for it.
\end{lemma}
\begin{proof}
Write $b=2u$, $u\ge3$, $N=6u-1$, $K=3u-1$.
If $u\ne5$, choose an ordinary alpha path on
$\interval{0}{2u-2}$ starting at $u+1$.  Its larger class is
upper and its other endpoint is $2u-4$.  Its only possible
forbidden endpoint request occurs at $u=5$: writing $u=2q+1$,
the upper forbidden label $3q$ equals $u+1=2q+2$ exactly at
$q=2$.  Add two to all labels and append the lower label $-1$.
The resulting word $Q$ has lowers
$\{-1\}\cup\interval{2}{u}$, uppers
$\interval{u+1}{2u}$, colors $\interval{1}{2u-1}$, first
label $u+3$, and last label $-1$.

Use three arms as follows.  First take the maximum Walecki word of
length $2u-1$, but raise its last lower label from $u-1$ to $u$:
\[
 W=(N,1,N-1,2,\ldots,5u+1,u,5u).
\]
Its lowers are $\interval{1}{u-2}\cup\{u\}$, its uppers
are $\interval{5u}{N}$, and its colors are
$\interval{4u+3}{N}\cup\{4u,4u+1\}$.
Map the lowers of $Q$ by $x\mapsto x+u$ and its uppers by
$x\mapsto x+3u-1$, giving an even arm $B$.
Its lowers are $\{u-1\}\cup\interval{u+2}{2u}$ and its
uppers are $\interval{4u}{5u-1}$.  It has internal colors
$\interval{2u}{4u-2}$ and root color $4u+2$.
For the other even arm $C$ use the pairs
\[
 (4u-1-j,2u+1+j)\quad(0\le j\le u-2),
 \qquad\text{followed by }(3u,u+1).
\]
Its lowers are $\interval{2u+1}{3u-1}\cup\{u+1\}$,
its uppers $\interval{3u}{4u-1}$, its internal colors
$\interval{1}{2u-1}$, and its root color $4u-1$.

These three lower sets partition $\interval{1}{K}$ and their
upper intervals partition $\interval{K+1}{N}$.  Their colors,
in order, partition
\[
 \interval{1}{2u-1},\quad\interval{2u}{4u-2},\quad
 \{4u-1,4u,4u+1,4u+2\},\quad\interval{4u+3}{N}.
\]
All edges cross $K$ and the root is zero.
At $u=5$ use the following literal words, with cut $14$:
\[
\begin{split}
 &(29,1,28,2,27,3,26,5,25);\\
 &(22,7,24,8,21,9,19,11,20,6);\\
 &(18,13,17,10,16,14,15,12,23,4).
\end{split}
\]
Their full inventories verify the omitted endpoint case directly.
\end{proof}

\begin{lemma}[Repeated central range]\label{mc:oee-repeated-central}
For $s\ge3$ and $2\le k\le s$, there is a root-zero alpha
labeling of $\Sp{2s+1,2s+2k,2s+2k}$ with standard Walecki
shortest arm.
\end{lemma}
\begin{proof}
For each $k\ge3$, the word
\[
 (2k+1-j,j+1)\quad(0\le j\le k-3),
 \qquad (k+3,k,k+2,k+1)
\]
is $L(0,k,k,k-1)$.  Its lower labels are
$\interval{1}{k-2}\cup\{k,k+1\}$; its uppers form
$\interval{k+2}{2k+1}$.  Its descending alternating initial
colors, the join to the final block, and that block itself give
\[
 \interval{6}{2k},\quad\{5\},\quad\{3,2,1\},
\]
respectively, with the obvious empty ranges at $k=3$.
Thus the missing color is four and its endpoints are $2k+1,k+1$.
Grow it by $P=s\ge k$ using Lemma~\ref{mc:lower-growth}.
The result has $t=u=s+k$, $H=k-1$, $D=2s+4$.
The completion request is
\[
 X=2s+k+2,\qquad Y=s+2k-3
\]
on order $2u-1$.  Both are in the upper interval because $k\ge3$;
the forbidden equality would require $k=s+5$, impossible when
$k\le s$.  Lemma~\ref{mc:lower-completion} applies.

For $k=2$ start instead with
$L(1,3,3,3)=(6,4,5,2,7,1)$ and grow by
$P=s-1\ge2$.  The resulting completion request on order
$2s+3$ has endpoints $2s+2,s+3$; its forbidden equality would
require $s=1$.  It is therefore allowed for every $s\ge3$.
\end{proof}

\begin{lemma}[The final repeated boundary]\label{mc:oee-repeated-boundary}
For $s\ge3$, $\Sp{2s+1,4s+2,4s+2}$ is root-zero alpha
with standard Walecki shortest arm.
\end{lemma}
\begin{proof}
Set $t=u=2s+1$.  Translate a standard maximum-Walecki odd path on
$\interval{1}{2s+1}$ by $s+2$.  It starts at $3s+3$, ends
at $2s+3$, has lower labels $\interval{s+3}{2s+2}$ and
upper labels $\interval{2s+3}{3s+3}$.
Follow it by an odd alpha path on $\interval{1}{2s+1}$ with
larger class lower, starting at two and ending at $s$, with its
uppers raised by $2s+2$.  Use the explicit endpoint-one word
of Section~\ref{templates:path-words}, shifted from zero-based
labels by one.  Its last three labels before the upper shift are
\[
 (s+1,s+2,s).
\]
For odd $s$ this follows from the final base $(1,2,0)$ shifted by
$s-1$ and then by one; for even $s$ it follows from the final
base tail $(4,5,3)$ shifted by $s-4$ and then by one.
These two cases apply for $s\ge3$ and $s\ge4$, respectively.
Its shifted internal colors are $\interval{2s+3}{4s+2}$.
The join has color $2s+1$; hence the combined kernel has hole
$H=s+2$, missing color $2s+2$, and endpoints $3s+3,s$.

Its last edge is $(3s+4,s)$, with color $2s+4$.
Change only its leaf $s$ to $s+2$.  This fills the old label hole
and opens the new hole $H'=s$.  The edge color becomes $2s+2$,
so the new missing color is $D'=2s+4$.
The new endpoint $s+2$ is exactly $t+u-2s-H'$.
Thus it is $L(s,2s+1,2s+1,s)$.

Its completion on order $4s+1$ has upper endpoints
$X=3s+3,Y=3s-1$.  The forbidden equality compares
$6s+4$ with $6s$, so cannot hold.  Apply
Lemma~\ref{mc:lower-completion} to obtain the claimed spider.
\end{proof}

\begin{lemma}[Two odd paths and reflection]\label{mc:two-odd-kernel}
Let $r,w\ge1$ and $\epsilon\in\{0,1\}$, and put
\[
 t=r+w+1,\quad u=s+w+1+\epsilon,\quad H=w+2.
\]
Define $X_A=2r+w-s+1$ and $X_B=s-r+2-\epsilon$.
If
\[
\begin{gathered}
 r+1\le X_A\le2r+1,\qquad 1\le X_B\le w+1,\\
 2(X_A-1)\ne3r,\qquad2(X_B-1)\ne w,
\end{gathered}
\]
then $L(s,u,t,H)$ exists.  Moreover, the class permutation
\[
 x\longmapsto t+2-x\ (x\le t+1),\qquad
 x\longmapsto3t+3-x\ (x\ge t+2)
\]
maps it to $L(t-s-1,2t-u+2,t,t+2-H)$.
\end{lemma}
\begin{proof}
Take an ordinary odd alpha path on $\interval{1}{2r+1}$
with larger class upper, starting at $X_A$ and ending at
$r-w+s+1$, and translate it by $w+2$.
Follow it by an odd alpha path on $\interval{1}{2w+1}$
with larger class lower, starting at $X_B$ and ending at
$w-s+r+\epsilon$; keep its lowers and raise its uppers by
$2r+2$.  The latter request is the complement of the ordinary
upper request, so its forbidden equality is exactly the one displayed.
The two paths give lower intervals $\interval{w+3}{t+1}$ and
$\interval{1}{w+1}$, and adjacent upper intervals
$\interval{t+2}{t+r+2}$ and
$\interval{t+r+3}{2t+1}$.
Their internal color intervals are $\interval{1}{2r}$ and
$\interval{2r+3}{2t}$, and the join has color $2r+1+\epsilon$.
Thus the missing color is $2r+2-\epsilon=D$, and direct substitution
gives the required endpoints.  The class permutation is a bijection
of the two label classes.  It changes each color $d$ to $2t+1-d$;
substitution in the endpoints, hole, and missing color gives exactly
the asserted new parameters.  It does not reverse the word.
\end{proof}

\begin{lemma}[The complete large repeated range]
\label{mc:oee-repeated-large}
For $S\ge3$, $T\ge2S+2$, $\Sp{2S+1,2T,2T}$ is root-zero
alpha with standard Walecki shortest arm.
\end{lemma}
\begin{proof}
Apply Lemma~\ref{mc:two-odd-kernel} with source parameters
\[
 s=T-S-1,\quad w=S+2-\epsilon,\quad r=T-w-1,
 \quad u=T+2.
\]
Its class reflection gives the target lower-hole parameters
$(S,T,T,H)$ with $H=r+1$.
The two source requests and the target completion are all valid if
\begin{equation}\label{mc:reflected-domain}
 T\ge2S+3-\epsilon,\quad T\ne3S+3-\epsilon,\quad
 S\ne4-3\epsilon,\quad T\ne2S+7-2\epsilon.
\end{equation}
For completeness, their first endpoints reduce to
$X_A=T-2+\epsilon$, $X_B=4-2\epsilon$, and
$X=T+S+3-\epsilon$.  The other completion endpoint is
$Y=2T-S-4+\epsilon$.
The three exclusions in \eqref{mc:reflected-domain} are respectively
the forbidden requests for the first source path, second source path,
and completion.  Thus none has been silently omitted.

For $\epsilon=1$, all $T\ge2S+2$ are covered except
$T=3S+2$ and $T=2S+5$.  For $T=3S+2$, the $\epsilon=0$
conditions hold unless $S=4$ or $S=5$.  For $T=2S+5$, they
hold unless $S=4$ (the other equality would force $S=2$).
Consequently exactly the following three finite kernels suffice:
\begin{center}\small
\begin{tabular}{ccl}
$(S,T)$&$H$&word\\\midrule
$(4,13)$&7&$(23,5,24,4,25,3,26,2,27,1,18,12,17,13,16,14,$\\
&&$15,8,22,6,21,9,20,10,19,11)$\\
$(4,14)$&9&$(25,7,24,5,26,6,22,8,23,1,29,2,28,3,27,4,$\\
&&$17,14,16,15,19,12,18,13,21,10,20,11)$\\
$(5,17)$&11&$(30,3,35,1,34,4,26,10,27,9,28,8,29,6,32,7,31,$\\
&&$2,33,5,20,18,19,16,21,17,23,14,22,15,25,12,24,13)$.
\end{tabular}\end{center}
These are $L(S,T,T,H)$ by direct complete inventory.
Their missing colors are $13,12,14$, respectively, and their
completion requests are $(2T-1,X)=(25,20),(27,20),(33,24)$.
All three requests are allowed.  Lemma~\ref{mc:lower-completion}
therefore fills these finite cases, proving the entire range.
\end{proof}

\begin{proposition}\label{mc:oee-repeated}
For every odd $a\ge7$ and even $b>a$, $\Sp{a,b,b}$ is
root-zero alpha.  Except at $b=a+1$, the constructions have the
standard maximum-Walecki shortest arm.
\end{proposition}
\begin{proof}
Write $a=2s+1,b=2u$, $s\ge3$.  If $u=s+1$ use
Lemma~\ref{mc:oee-near}.  If $s+2\le u\le2s$ use
Lemma~\ref{mc:oee-repeated-central}.  At $u=2s+1$ use
Lemma~\ref{mc:oee-repeated-boundary}.  For $u\ge2s+2$ use
Lemma~\ref{mc:oee-repeated-large}.  These disjoint ranges exhaust
all even $b>a$ and have the stated pinning properties.
\end{proof}

\begin{corollary}\label{mc:eee-repeated}
For even $a\ge6$ and even $b\ge a+4$, $\Sp{a,b,b}$ is
root-zero alpha.
\end{corollary}
\begin{proof}
Use Proposition~\ref{mc:oee-repeated} on $(a+1,b,b)$.
Here $b\ne(a+1)+1$, so the source has a standard Walecki
shortest arm and Lemma~\ref{mc:shorten} applies.  The modified
near-family word is not used in this deduction.
\end{proof}

\subsection{The complete EOO classification}
\label{construction:all-eoo-alpha3}
\label{mc:all-eoo-section}

\begin{theorem}\label{mc:all-eoo-large}
If $a\ge6$ is even and $a<b\le c$ are odd, then
$\Sp{a,b,c}$ is root-zero alpha.
\end{theorem}
\begin{proof}
Write $a=2s,b=2t+1,c=2u+1$, where $s\ge3$ and
$s\le t\le u$.  Every odd source shortened below is obtained
by a construction explicitly appending the maximum-Walecki arm of
length $2s+1$; an unpinned odd existence theorem is insufficient.
There are four exhaustive regions.

If $c\le2a+1$, Proposition~\ref{mc:eoo-central} applies.

If $c\ge a+b$, take the root-zero alpha pair $\Sp{a,b}$,
with cut $s+t$, and apply Lemma~\ref{common:alpha-extension}.
The only possible exception is
\begin{equation}\label{mc:eoo-far-resonance}
 c=4(s+t)+1.
\end{equation}
For it use the explicitly pinned repeated odd source
$\Sp{2s+1,2t+1,2t+1}$ from
Lemma~\ref{ok:repeated-kernels}.  It has a long upper leaf at
distance $e\le s+t+1$ above its cut $s+2t$.
Here is the full leaf-interface check.  At $t=s$, the equal-odd
formula has long leaves $5s+1,5s+2$ for $s\ge4$; at $s=3$
its literal equal-seven labeling has smaller long leaf $16$ and
cut nine.  At $t>s$, the generic repeated parameters give
\[
 e_A=t+2s-V+3,\qquad e_B=t+V.
\]
For $V=s+1$ choose $e_B=s+t+1$; for $V=s$ choose $e_B=s+t$;
for $V=s+3$ choose $e_A=s+t$.  These are exactly the three
parameter choices in the complete repeated construction.
Shorten its Walecki arm; all its other labels and these distances
are unchanged.  Append $2z$ vertices to the selected long leaf,
where $z=2s+t\ge e$, using
Lemma~\ref{common:upper-leaf-extension}.  The new length is
$2t+1+2(2s+t)=4(s+t)+1$, as required.

It remains that $c>2a+1$ and $c<a+b$.
Suppose first $b\ge2a+1$.  The explicit pinned families in
Lemma~\ref{ok:reverse-boundary} and Section~\ref{ok:short-wedges}
give this complete selection.  For $b=2a+1$, the middle half
is $2s$ and the largest half is $2s\le u\le3s-1$.
Its repeated endpoint uses the twice-ray kernel; the next point
$u=2s+1$ uses the class-reflected gap-two kernel, with its two
long carriers interchanged;
all $u\ge2s+2$ use the reversed wedge with its minus choice.
For $b=2a+3$, the middle half is $2s+1$ and the largest half
is at most $3s$; the plus-choice reversed wedge applies throughout.
For $b\ge2a+5$, put $U=(b-1)/2$, $\delta=(c-b)/2$.
Then $U\ge2s+2$ and $0\le\delta\le s-1$, exactly the
complete unequal-short-wedge domain, including its repaired
$\delta=s-3$ plane.  The indicated formulas, carrier orders,
and all their path bounds are proved in the cited sections.
Each completion pins the same Walecki short arm; shorten it.

Finally suppose $b<2a+1$ in that remaining region.  Then
\[
 u\ge2s+1,\quad t\le2s-1,\quad u\le s+t-1.
\]
Set $k=u-2s$, $h=s+t-u$.  Thus
\[
 k,h\ge1,\quad t=s+k+h,\quad u=2s+k,\quad s\ge k+h+1.
\]
The explicit transition construction in
Lemma~\ref{ok:transition-kernels} completes with a prescribed
Walecki shortest arm whenever $h\ge2$, or $h=1$ and
$s\ge k+3$.  For $h\ge2$ its source is the elementary kernel
$K(0,0,h,h+1)$, followed by carrier prefix and asymmetric growth.
The sole excluded boundary is $h=1,s=k+2$, whose source arms
are $(2s+1,4s-1,6s-3)$.
The controlled-edge induction in Lemma~\ref{ok:pinned-ap-ray}
constructs $K(s,3s-2,2s-1,2)$ for every $s\ge3$ and
completes it with the required pin.  This boundary also shortens.
The four regions exhaust the domain and prove the theorem.
\end{proof}

\begin{theorem}[Full EOO classification]\label{mc:all-eoo}
For positive integers $a<b\le c$, with $a$ even and $b,c$ odd,
$\Sp{a,b,c}$ is root-zero alpha if and only if
$(a,b,c)\ne(2,3,3)$.
\end{theorem}
\begin{proof}
The preceding theorem proves $a\ge6$.
For $a=2,4$ and $c\ge7$, use the complete minimum-two and
minimum-four constructions of Section~\ref{app:small-alpha}.
For $c<7$ the three positive cases have these literal arm words:
\[
\begin{array}{ll}
\Sp{2,3,5}:&(10,1);\ (6,2,9);\ (8,3,5,4,7),\\
\Sp{2,5,5}:&(12,1);\ (10,4,6,5,8);\ (9,2,7,3,11),\\
\Sp{4,5,5}:&(14,1,13,2);\ (10,3,8,4,12);\ (9,6,7,5,11).
\end{array}
\]
The full inventories follow by subtraction, with respective cuts
$4,5,6$ from the forced index formula.

For $\Sp{2,3,3}$ the cut would be three.
Every upper vertex is a neighbor of the root or a leaf, so none has
two nonroot lower neighbors.  Color one forces $(3,4)$.
Color two cannot be $(2,4)$ by this property, so it forces $(3,5)$.
Lower vertex three now has its full degree two occupied.
Color three can only be $(1,4)$, $(2,5)$, or $(3,6)$.
The first two violate the upper-vertex property, and the last gives
lower vertex three a third neighbor.  Each is impossible.
This excludes a root-zero alpha labeling of this profile and proves
the claimed classification.
\end{proof}

\subsection{The ten inputs and the five-to-four inputs}
\label{mc:input-summary}

The following table records exactly the constructions justifying
the seven-arm predicates.  The three lengths are sorted and at
least six.  Letters indicate their parities, and the forced cut is
always $\floor{a/2}+\floor{b/2}+\floor{c/2}$.
\begin{center}\small
\begin{tabular}{p{.17\linewidth}p{.42\linewidth}p{.32\linewidth}}
Input&Condition&Proof\\\midrule
Equal&$a=b=c$&Equal-triple formulas and equal-even repairs\\
Long&$c>2k$, $c\ne4k+1$, $k=\floor{a/2}+\floor{b/2}$
 &Lemma~\ref{common:alpha-extension}\\
OOO central&$c\le2a-1$&Proposition~\ref{ok:central-cone}\\
EOO central&$c\le2a+1$&Proposition~\ref{mc:eoo-central}\\
EEO central&$a<b$, $c\le2a-1$&Corollary~\ref{mc:eeo-central}\\
EEO near&$a=b$, $c=a+1$&Proposition~\ref{mc:eeo-near}\\
OEE central&$a<b<c\le2a$&Proposition~\ref{mc:oee-central}\\
OEE repeated&$a<b=c$&Proposition~\ref{mc:oee-repeated}\\
EEE central&$a<b<c\le2a+2$&Corollary~\ref{mc:eee-central}\\
EEE repeated&$a<b=c$, $b\ge a+4$&Corollary~\ref{mc:eee-repeated}.
\end{tabular}\end{center}
The equal-odd formula is in Proposition~\ref{ok:equal-odd-three};
equal-even repairs and their exact small bases are in
Section~\ref{templates:equal-even} and Section~\ref{app:finite-bases}.
None uses a seven-arm selection theorem.

The manual five-to-four proof additionally uses
Proposition~\ref{mc:oeo-central} and the EOO near family
$\Sp{a,a+1,a+1}$ from Proposition~\ref{mc:eoo-repeated}.
It can use the smaller pre-repair EOO central domain, which is
contained in Proposition~\ref{mc:eoo-central}.
Deriving the six-to-four bridge from that manual lemma therefore
introduces no circularity: the construction layer precedes every
seven-arm selection and unit-extension argument.

Every infinite assertion here reduces to prescribed-endpoint paths,
explicit interval operations, a complete induction, or the
finite-template verification theorem.  Printed finite words and
finite-template tables form part of those proofs.  Independent graph
and integer audits provide additional verification; bounded samples
are not used to infer an infinite family.

\section{Five-to-four hub-zero selection}\label{app:five-to-four}
\begin{theorem}\label{fs:selection}
Among any five positive arm lengths, four can be selected whose spider is hub-zero graceful.
\end{theorem}
The proof is an integer case argument. It uses the equal and long alpha triples; the OOO central cone $z\le2x-5$; the EOO central cone $z\le2x-3$, $x+1+z\ne2y$; the OEO cone $z\le2x-3$; the EEO cone $x<y$, $z\le2x-1$; and the near triples $(x,x+1,x+1)$ and $(x,x,x+1)$ for even $x$. All lengths in these uses are at least six. The equal-triple construction is in Section~\ref{app:equal-alpha}; the odd and mixed constructions are proved in Sections~\ref{app:odd-kernels} and \ref{app:ten-alpha}. Each triple's alpha index is the sum of its half-length floors.

\begin{proof}
If a length is at most five, select it and any three other arms and apply Lemma~\ref{seven:four-short}. Thus order the five lengths as $6\le a\le b\le c\le d\le e$ and put
\[
A=\floor{a/2},\ B=\floor{b/2},\ C=\floor{c/2},\
D=\floor{d/2},\ E=\floor{e/2},\quad K=A+B.
\]
In particular $A\ge3$.
If $a=b=c=x$, select these three and $e$. For $e>3\floor{x/2}$, use ordinary attachment to the equal alpha triple. Otherwise
\[
\floor{e/2}\le\floor{3\floor{x/2}/2}<x,
\]
so equal-triple offsets and insertion at residue zero apply. Hence assume $a,b,c$ are not all equal. Then $K\le c-1$: for odd $c$, $K\le2C=c-1$; for even $c$, failure would require $A=B=C$ and hence $a=b=c$ even.

\medskip\noindent\textit{Case 1: $d>2K$.}
If $d\ne4K+1$, the long alpha triple $(a,b,d)$ is available and accepts $e$. Indeed $d>2K$ implies $\ceil{d/2}>K$, so $e\ge d>K+D$, its alpha index.
Suppose $d=4K+1$. If $c=d$, combine alpha pair $(a,b)$ with equal-pair offset $K$ on $(c,d)$, since $K<c$. Assume $c<d$.
If $c>2K$, use the nonexceptional long alpha triple $(a,b,c)$ and attach $d$. If $c\le2K$ and $C>B$, use $(a,c,d)$: $C\le K$ gives $2(A+C)\le4K<d$, while $d=4(A+C)+1$ would force $C=B$. This triple accepts $e$ by the same ceiling inequality. If $c\le2K$, $C=B$, and $A<B$, use $(b,c,d)$: its long inequality is $d>4B$, and its exceptional equality would force $A=B$.
The only remaining possibility is $A=B=C=H$. Then $a,b,c$ lie in $\{2H,2H+1\}$, so their triple is equal, EEO near, or EOO near. In each case it is alpha with index $3H$, and $e\ge d=8H+1>3H$. Attach $e$. This finishes Case 1.

\medskip\noindent\textit{Case 2: $d\le2K$.}
Alpha pair $(a,b)$ and offset pair $(c,d)$ satisfy both high-interval bounds
\[
\floor{d/2}\le K\le c-1.
\]
They therefore compose unless
\begin{equation}\label{fs:R}
d=4c-4K-3.
\end{equation}
Assume this equality. It implies $d\equiv1\pmod4$, so $d$ is odd and $D$ is even. If $c=d$, equal-pair offsets work. If $b=c$, combine alpha pair $(a,d)$ with equal offset pair $(b,c)$: writing $b=2B+p$ with $p\in\{0,1\}$, equation~\eqref{fs:R} gives
\[
A+D=b+p-A-2<b.
\]
Thus assume $b<c<d$.

Since $A\le B$ and $d\le2K\le4B$, the congruence gives $d\le4B-3$. If $b$ is odd this is $d\le2b-5$; hence $(b,c,d)$ is alpha by the OOO cone when $c$ is odd and the OEO cone when $c$ is even. If $b,c$ are even, strict inequality gives $c\ge b+2$ and $d\le2b-3$, so EEO applies. If $b$ is even and $c$ odd, EOO applies except possibly when
\begin{equation}\label{fs:H}
b+1+d=2c.
\end{equation}
Together with~\eqref{fs:R}, this last equality gives $c=2A+B+1$, $d=4A+1$. Since $c$ is odd, $B$ is even; since $c<d$, $B<2A$, and therefore $B\le2A-2$. Thus $c\le4A-1$. If $a$ is odd, $(a,b,c)$ is OEO central because $4A-1=2a-3$. If $a$ is even, $d\le2K$ also gives $B\ge A+1$, so $b\ge a+2$ and $c\le2a-1$; $(a,b,c)$ is EEO central. We have found an alpha triple $(b,c,d)$, or in the sole exceptional subcase an alpha triple $(a,b,c)$.

If $e$ exceeds the selected triple's index, ordinary attachment finishes. Otherwise consider alpha pair $(a,d)$ with offset pair $(c,e)$ at index $A+D$. Equation~\eqref{fs:R} gives
\[
c=K+D/2+1.
\]
Since $D\le K\le2B$, this proves $A+D\le c-1$. For the lower bound, if $(a,b,c)$ was selected, failure of attachment gives $e\le A+B+C$. As $B,C\le D$, $\floor{e/2}\le A+D$ follows. If $(b,c,d)$ was selected, failure gives $e\le B+C+D$. Write $c=2C+p$. Equation~\eqref{fs:R} gives
\[
B+C-D=2A+3B-3C-2p+2.
\]
For $p=1$, $B\le C$ bounds this by $2A$. For $p=0$, $b<c$ forces $B\le C-1$, bounding it by $2A-1$. Hence again $e\le2(A+D)$ and $\floor{e/2}\le A+D$.
Finally the high-interval exception would require
\[
e=4c-4(A+D)-3=4B-2D+1.
\]
But $b<c<d$ implies $b\le d-2=2D-1$, so $B\le D-1$ and this value would be at most $2D-3=d-4$, contradicting $e\ge d$. The high pair construction is nonexceptional. Alpha-offset composition supplies the required hub-zero four-arm spider, finishing Case 2 and the proof.
\end{proof}

The theorem yields the same selection among any six or more arms by taking any five first. Its proof uses neither the seven-arm conclusion nor any solver assertion. This makes its use in the small-arm reduction logically independent of the unbounded seven-arm arithmetic certificate.

\section{Finite affine construction certificates}
\label{app:template-certificates}

Some construction families use a fixed edge trade whose vertices move
affinely with the arm lengths.  This section specifies the proof data,
the verification argument, and the exact parameter coverage of these
families.  In particular, fitting a formula to several examples is not
an existence argument here.  Every formula used in a construction is
checked against its source word on an entire integer ray.

The finite inputs are included in the accompanying machine-readable
supplement, with the names and hashes in Section~\ref{templates:ledger}.
All $16$ minimum-two templates and all $216$ minimum-four templates
are included in full.  The two equal-even repairs, the alternative
endpoint repair, and the four equal-odd four-arm repairs are also
printed in Section~\ref{templates:small-tables}, including every
changed vertex's source arm and depth.  The tables are construction
instructions; a search program is not needed to instantiate them.

\subsection{Finite word inventories}
\label{templates:finite-check}

A finite spider certificate specifies a hub label $z$, an ordered list
of arm lengths, and one complete outward word for each arm.  Put
$N$ equal to the sum of those lengths.  To check an ordinary graceful
certificate, check the word lengths and the two multiset identities
\[
 \{z\}\uplus\biguplus_i\{P_{i,j}\}=[0,N],\qquad
 \biguplus_{i,j}|P_{i,j}-P_{i,j-1}|=[1,N],\quad P_{i,0}=z.
\]
For a hub-zero certificate also require $z=0$.  For an offset-$k$
certificate require $z=0$ and replace the first identity on nonroot
labels by $\biguplus_i\{P_{i,j}\}=[k+1,k+N]$; the color identity
is unchanged.  For alpha, check $z=0$, the cut
$K=\sum_i\floor{L_i/2}$, and that depth $j$ is lower exactly when
$j$ is even.  A graceful permutation on $n$ vertices is checked by
labels $[0,n-1]$ and successive differences $[1,n-1]$, with the
stated first value and alpha condition when applicable.

All checks are finite list operations with integer arithmetic.  A
catalog must additionally have distinct profile keys and exactly
the required profile set.  Controlled interfaces, such as a pinned
short word or a specified upper-leaf distance, are checked directly
on the words.  Thus a finite record's search status or node count
has no role in its acceptance as a labeling proof.

\subsection{Affine data and source coordinates}
\label{templates:encoding}

An integer ray has the form $h\in\mathbb Z$, $h\ge h_0$.  A vertex
name is an affine function $Ah+B$, not a numerical vertex sampled at
a particular $h$.  The templates use the following three encodings.
\begin{enumerate}
\item In \texttt{S2o.jsonl}, \texttt{S2e.jsonl}, and
\texttt{S4.jsonl}, a pair \texttt{[A,B]} means $Ah+B$.
The first two files use $b=8h+r$; the third uses $b=48h+r$.
\item In \texttt{E.json} and \texttt{Alt.json}, a pair
\texttt{[A,B]} means $(Ar+B)/2$.  The parameter is
$r=4q+\epsilon$, $q\ge3$, with $\epsilon=1$ or $3$ as specified.
After substitution, every label and depth is integral affine in $q$.
\item In \texttt{Q0.json},\ldots,\texttt{Q3.json}, the nested pair
\texttt{[[a,b],[c,d]]} means $(a/b)q+c/d$.  Here $s=4q+r$,
$r\in\{0,1,2,3\}$, and $q\ge6$.  The rational representation is
exact; all actual label and depth coefficients become integers.
\end{enumerate}

Arm numbers are zero based.  Depth zero is the hub, and the edge of
depth $j$ joins depths $j-1$ and $j$.  In position tables the arm
number $-1$ denotes the hub, whose depth and label are zero.
For an explicitly defined source spider, write
$F_i(h,j)$ for the label at depth $j$ on source arm $i$, setting
$F_i(h,0)=0$.  The source formulas, including all exceptional small
path bases, are specified below and in the common construction
appendices.

\begin{center}
\begin{tabular}{p{.28\textwidth}p{.61\textwidth}}
Field & Mathematical meaning\\\hline
\texttt{residue}, \texttt{delta} & $r$ and the target long-arm gap $\delta$\\
\texttt{remove} & List of objects with fields \texttt{u}, \texttt{v},
\texttt{arm}, \texttt{depth}; these specify a source edge and its two affine labels\\
\texttt{add} & List of affine endpoint pairs specifying the replacement edges\\
\texttt{source\_b} & Parameter used during discovery; it is not the lower bound of the proved ray\\
\texttt{verification.h\_min} & The inclusive verified ray bound $h_0$\\
\texttt{verification.b\_min} & $8h_0+r$ or $48h_0+r$, respectively\\
\texttt{removed\_edges}, \texttt{added\_edges} & Remove/add lists in the equal-source and four-odd formats\\
\texttt{positions} & Affine label-to-source-coordinate table; a dictionary in E/Alt, an object list in the Q files\\
\texttt{q\_minimum} & Inclusive lower bound six in a Q recipe.
\end{tabular}
\end{center}
The cached \texttt{verification} output is descriptive.  A verifier
recomputes its obligations from the source words and the edge lists.
It does not accept the recorded word ``verified'' as a proof premise.
For the minimum-arm files its remaining output fields are
\texttt{removed\_edges}, \texttt{symbolic\_constraints}, and
\texttt{arm\_lengths}, together with copies of the residue and gap.
An E recipe is selected by the top-level key \texttt{"1"} or
\texttt{"3"}; an Alt recipe has the same three fields directly.
In their position dictionaries an entry
\texttt{"a,b":[i,[c,d]]} means label $(ar+b)/2$ at source arm $i$
and depth $(cr+d)/2$.  A Q position is instead an object with fields
\texttt{label}, \texttt{arm}, and \texttt{depth}, decoded using the
rational-coefficient convention above.

\subsection{Why the finite verification proves an infinite family}
\label{templates:soundness-section}

The source need not yet have the desired color inventory.  Let its
signed color defect be
\[
 D_0=\mathcal C(G_0)-[1,N],
\]
where $[1,N]$ denotes the multiset containing each integer once and
$\mathcal C(G_0)$ counts edge differences with multiplicity.  In the
applications below, an interval proof gives $D_0$ as a finite signed
sum of affine color functions.  For the minimum-two and minimum-four
sources, $D_0=0$.

\begin{theorem}[Soundness of an affine trade certificate]
\label{templates:soundness}
Suppose a source graph on labels $0,\ldots,N(h)$ is an explicitly
labeled spider for every integer $h\ge h_0$, all its edges cross a
specified cut $K(h)$, and its signed color defect is the proved
finite expression $D_0(h)$.  Let $R$ and $A$ be finite affine edge
lists.  The following checks imply that replacing $R$ by $A$ gives
a root-zero alpha labeling of the stated target spider for every
$h\ge h_0$.
\begin{enumerate}
\item All participating labels and depths are integral and in range.
Every recorded vertex is proved to equal its source coordinate
$F_i(h,j)$.  Distinct recorded vertices never collide on the ray.
\item Every edge of $R$ occurs at its recorded source arm and depth.
The lists contain no repeated edges, and every edge of $A$ is absent
from the retained source graph.  Each added edge crosses $K(h)$.
\item The endpoint-incidence multisets of $R$ and $A$ agree.  With
each edge oriented from its lower label to its upper label, its
color is the positive affine difference of its endpoint functions.
As a signed multiset of affine functions,
\begin{equation}\label{templates:color-identity}
 D_0-\mathcal C(R)+\mathcal C(A)=0.
\end{equation}
\item The weighted graph obtained by contracting all untouched path
intervals has exactly the stated number of hub-to-leaf paths, uses
every contracted vertex once, and has the stated target arm lengths
as affine identities.
\end{enumerate}
\end{theorem}

\begin{proof}
Fix an arbitrary integer $h\ge h_0$.  The coordinate checks identify
actual distinct vertices and actual removable edges of the source.
The retained and added edges therefore define the prescribed graph
on exactly the original vertex-label set.  In particular, every label
$0,\ldots,N(h)$ still occurs once, and the hub is still zero.
All retained and added edges cross the same cut.  Evaluating the
formal multiset identity \eqref{templates:color-identity} gives the
complete color multiset $[1,N(h)]$, even if two affine color functions
happen to coincide at that parameter.  The endpoint-incidence check
also preserves the degree of every labeled vertex.

To justify the topology check, on each source arm cut exactly the
edges in $R$.  Every remaining component is an interval of that path.
Keep its boundary vertices, and replace a nontrivial interval by a
weighted edge whose weight is its exact number of original edges.
A singleton component remains a vertex.  Different interval interiors
are disjoint.  Reintroduce the edges in $A$ with weight one.  Expanding
the weighted edges recovers precisely the traded graph.  Thus a
traversal proving that the weighted graph consists of the required
hub paths proves the same assertion for the full graph, with each
arm length equal to the sum of its weights.  Exhausting the weighted
vertices excludes any additional component or cycle.  This proves the
target topology, and hence the root-zero alpha labeling.  Since $h$
was arbitrary, the result holds on the whole ray.
\end{proof}

\subsubsection{Exact operations on an integer ray}

All affine assertions reduce to finite rational arithmetic.  Equality
means equality of both coefficients.  For $Ah+B\ge c$ on
$h\ge h_0$, it suffices and is necessary that $A\ge0$ and
$Ah_0+B\ge c$.  A strict positive inequality is treated with $c=1$
after integrality has been established.  Distinctness is checked by
solving $Ah+B=0$: a collision exists precisely when its root is an
integer at least $h_0$, apart from the constant cases.

When a source word has a piecewise coordinate formula, the verifier
first selects the appropriate finite base, block index, and residue
of the depth.  It then proves the inequalities defining that branch
on the full ray.  Any remaining periodic test is resolved by its
finitely many residue classes.  Ordering cut depths may be proposed
using a convenient parameter, but every consecutive difference in
the proposed order is subsequently proved positive on the whole
ray.  The proposed order alone supplies no mathematical inference.

The threshold computation used for the minimum-two and minimum-four
tables can be described explicitly.  Start with $h_0=0$.  Each
required inequality $Ah+B\ge c$ with $A>0$ raises the bound to
\[
 h_0\leftarrow\max\left(h_0,\ceil{(c-B)/A}\right).
\]
A constant inequality must already hold, and a negative slope fails
the selected eventual branch.  The inclusive bounds printed below
are sufficient certified bounds, not claims of optimal construction
thresholds.

\subsubsection{Verification pseudocode}

The finite procedure underlying Theorem~\ref{templates:soundness}
is the following.  Formal colors are stored as coefficient pairs
with signed multiplicities.
\begin{verbatim}
decode the parameter domain, source words, and affine records
for every participating vertex:
    prove integrality, source arm/depth, and label identity
prove all required ranges and vertex distinctness on the ray
for every removed edge:
    prove its two depths are consecutive on one source arm
for every added edge:
    prove distinctness, absence in the retained source, and cut crossing
compare the removed and added endpoint-incidence counters
cancel seed_defect - removed_color_counter + added_color_counter
on every source arm:
    sort and certify all cut depths
    contract each untouched interval with its exact affine length
insert each added edge with weight one
traverse every hub branch, rejecting repeated or omitted vertices
compare the affine path-length multiset with the target lengths
accept only if every preceding exact check succeeds
\end{verbatim}
For example, a removed edge with larger depth $d$ separates an
untouched interval ending at $d-1$ from one starting at $d$.
The weight between retained depths $u\le v$ is $v-u$; a zero
weight contributes a singleton, not a new edge.  This convention
avoids counting vertices as edges in the topology calculation.

\subsection{The fixed-endpoint path words used by the templates}
\label{templates:path-words}

Let $P_e(n)$ be an alpha word on $\{0,\ldots,n-1\}$ beginning at
the lower endpoint $e$, for $e\in\{1,2,3\}$.  The finite bases and
prefixes specify these words completely.  For $e=1$ the prefixes and
bases are
\[
 Q_1=(1,2,0,3),\qquad
 \begin{array}{c|l}
 n& P_1(n)\\\hline
 3&(1,2,0)\\
 4&(1,2,0,3)\\
 6&(1,5,0,3,2,4)\\
 9&(1,8,0,6,2,7,4,5,3).
 \end{array}
\]
For $e=2,3$, input \texttt{P.json} gives
\[
 Q_2=(2,5,3,4,0,7,1,6),\qquad
 Q_3=(3,11,0,10,1,8,2,7,4,6,5,9).
\]
Its base sizes are respectively
\[
 \begin{array}{c|l}
 e=2&5,6,7,8,10,11,12,17,\\
 e=3&7,8,9,10,11,12,14,15,16,17,18,25.
 \end{array}
\]
The file's top-level keys are \texttt{"2"} and \texttt{"3"}; each
has fields \texttt{prefix} and \texttt{bases}, the latter a dictionary
from the decimal base size to its entire vertex word.  Every one
of these finite words is checked by the ordinary permutation and
alpha inventories.

Here is a closed coordinate rule for all larger words.  Choose the
unique base size $\beta$ congruent to $n$ modulo $4e$, and write
$n=\beta+4ez$.  For $j=0,\ldots,z-1$, replace a value $x$ of
$Q_e$ by
\begin{equation}\label{templates:prefix-coordinate}
 \begin{cases}
 x+2ej,&x<2e,\\
 x+n-4e-2ej,&x\ge2e.
 \end{cases}
\end{equation}
Concatenate these $z$ blocks, followed by $P_e(\beta)+2ez$.
Equivalently, one prefix extends a word of order $n$ to order
$n+4e$: raise the upper prefix labels by $n$, then append the old
word translated by $2e$.  The prefix colors are
$n+1,\ldots,n+4e-1$, the joining color is $n$, and the old colors
are $1,\ldots,n-1$.  The lower and upper vertex blocks are
consecutive and disjoint.  This proves the recurrence, including
its first endpoint.  The complete permitted domains are
$n\ge2e+1$, $n\ne4e+1$; the exceptional residue begins instead at
$n=8e+1$.

\subsection{Sixteen minimum-two templates}
\label{templates:min2}

Write $b=8h+r$, $0\le r<8$, and let $\delta$ be the long-arm gap.
For odd $r$, $\delta\in\{0,1\}$; for even $r$,
$\delta\in\{1,2\}$.  Thus there are exactly $8+8=16$ templates.
The even repeated case $\delta=0$ is proved separately by the
balanced-even construction in Section~\ref{app:small-alpha}.

To define the source, put $n=b+\delta+2$ and $k=\floor{b/2}$.
Start with
\[
 (b-1,1),\qquad (b,\,P_1(b-3)+2),
\]
where the notation in the second word means concatenation of $b$
with the translated path.  Raise its labels above $k$ by $n$.
The added $n$-vertex arm is
\begin{equation}\label{templates:min2-newword}
 \begin{cases}
 (b+2+j,b+1-j)_{j=0}^{b/2},&b\text{ even},\ \delta=0,\\
 \operatorname{rev}(P_1(n))+k+1,&n\text{ even},\\
 n+k-\operatorname{rev}(P_1(n)),&n\text{ odd}.
 \end{cases}
\end{equation}
The first line specifies the separately used even repeated source;
the $16$ repair records use the appropriate other line.  The source
arm lengths are $(2,b-2,b+\delta+2)$.  Its labels, colors, and
alpha separation follow from the path recurrence and the middle
interval attachment; its defect is zero.  Its total size is
$N=2b+\delta+2$ and the target cut is
$K=1+\floor{b/2}+\floor{(b+\delta)/2}$.

Each record's \texttt{remove} entry identifies an actual edge in
one of these three source arms.  Applying its \texttt{add} list and
the checks of Theorem~\ref{templates:soundness} gives arm lengths
$(2,b,b+\delta)$ as exact affine identities.  The files
\texttt{S2o.jsonl} and \texttt{S2e.jsonl} have respectively eight
records; their exact inclusive thresholds and edge counts are in
Section~\ref{templates:thresholds}.  For illustration, the first
minimum-two odd record has $r=1$, $\delta=0$ and begins by removing
the depth-one edge of arm two with endpoints $0$ and $8h+3$.
This is the source root edge $0$--$n$, since $n=b+2=8h+3$.

The complete finite complement of the $16$ rays is specified without
an arbitrary cutoff.  For each row $(r,\delta,b_0)$ in the threshold
table take
\begin{equation}\label{templates:min2-finite}
 \{(2,b,b+\delta):b\equiv r\pmod8,\ b<b_0,\
 b\ge5\text{ if }r\text{ is odd},\ b\ge6\text{ otherwise}\},
\end{equation}
and add $(2,3,9)$ and $(2,4,13)$.  The resulting set has exactly
$55$ profiles, all recorded in \texttt{F2.jsonl}.  Balanced odd
seeds have short leaf one; mixed even seeds used for resonance have
an upper leaf $K+e$ satisfying $2e\le b+4$.  These interfaces are
checked in addition to the alpha inventory.  The appendix's leaf
extension then reaches each resonance, and alpha hub attachment
handles the remaining larger long arm.  This specifies every
infinite ray and its exact finite remainder.

\subsubsection{The additional pinned-prefix condition for minimum three}
\label{templates:min3-protection}

The minimum-three proof first uses the same minimum-two templates,
but preserves the canonical first nine vertices of source arm one.
Every remove record has arm number one or two.  For a removed edge
on arm one with depth $ah+d$, add the condition $ah+d\ge10$.
All such slopes $a$ are positive in the actual records, so the bound
becomes
\[
 h'_0=\max\left(h_0,\max_{\text{arm-one cuts}}
                      \ceil{(10-d)/a}\right).
\]
Recomputing this expression on every record changes only
$(r,\delta)=(5,1)$: its threshold $b_0=21$ becomes $29$.
The endpoint-incidence identity adds no new incidence to a vertex
whose incident edges were frozen, so the protected prefix remains
the stated source prefix after the trade.

For $b\ge18$, the exact required pinned finite set is
\[
 \{(2,18,18)\}\ \cup\!!!
 \bigcup_{(r,\delta)}
 \{(2,b,b+\delta):18\le b<8h'_0(r,\delta)+r,\ b\equiv r\pmod8\}.
\]
It has $29$ profiles and is supplied completely in
\texttt{F3p.jsonl}.  Each has the maximum-pinned short word $(N,1)$;
every even mixed profile also has the upper-leaf bound $2e\le b+4$.
The even diagonal outside this finite set is handled by the separate
balanced constructions: $4\mid b$, $b\ge20$, or
$b\equiv2\pmod4$, $b\ge22$.  The latter has its first cut at
$b/2-1\ge10$.

For $3\le b\le17$, the exact direct minimum-three finite set is
\begin{align*}
 &\{(3,b,b+\delta):3\le b\le17,\ b+\delta\ge7,\\
 &\hspace{22mm}0\le\delta\le1\ (b\text{ odd}),\quad
                    0\le\delta\le2\ (b\text{ even})\}\\
 &\qquad\cup\ \{(3,b,4(1+\floor{b/2})+1):3\le b\le17\}.
\end{align*}
It has $44$ profiles, all in \texttt{F3s.jsonl}.  These sets are
derived from the thresholds and the attachment exception; they do
not introduce an empirical bound on the theorem.

\subsection{Two hundred sixteen minimum-four templates}
\label{templates:min4}

Now write $b=48h+r$, $0\le r<48$.  For odd $r$ take
$\delta=0,1,2,3$; for even $r$ take $\delta=0,1,2,3,4$.
There are $24\cdot4+24\cdot5=216$ records, one for each required
residue and near-band gap.

Put $n=b+\delta+4$ and $k=\floor{b/2}$.  Before raising upper
labels, the two-arm core is
\[
 (b,1,b-1,2),\qquad b-2-P_2(b-4).
\]
Raise every core label above $k$ by $n$.  The added arm is
\[
 \begin{cases}
 \operatorname{rev}(P_e(n))+k+1,&n\text{ even},\\
 n+k-\operatorname{rev}(P_e(n)),&n\text{ odd},
 \end{cases}
\qquad
 e=\begin{cases}\floor{n/2}-k-1,&n\text{ even},\\
 \floor{n/2}-k,&n\text{ odd}.
 \end{cases}
\]
For the required bands $e\in\{1,2,3\}$.  This completely defines
the source $\Sp{4,b-4,b+\delta+4}$, of total size
$N=2b+\delta+4$.  Its short word is $(N,1,N-1,2)$.
The kernel exception $b-4=9$ is excluded by the source domain;
the threshold or its literal finite complement handles that value.

Every removed edge in \texttt{S4.jsonl} has arm number one or two.
The endpoint-incidence identity then prevents any new incidence on
the frozen short arm.  The $216$ trades satisfy the exact color and
weighted-topology conditions with target lengths $(4,b,b+\delta)$.
The threshold table prints every $b_0$; the largest is $103$.
The thresholds arise from the explicit coordinate branches, not from
specializing the discovered words at a large parameter.

The required finite remainder is exactly
\begin{equation}\label{templates:min4-finite}
 \begin{aligned}
 \{(4,b,b+\delta):{}&b\ge9,\quad b\equiv r\pmod{48},\quad b<b_0(r,\delta),\\
                  &0\le\delta<4\text{ for odd }r,\quad
                   0\le\delta<5\text{ for even }r\}.
 \end{aligned}
\end{equation}
It has $139$ profiles, given in \texttt{F4.jsonl}.  Every word has
the pinned short arm $(N,1,N-1,2)$.  For odd $b$ at $\delta=0$,
the smaller upper long leaf $K+e$ satisfies $2e\le b+7$; for
even $b$ at $\delta=1$, the odd long leaf satisfies $2e\le b+8$.
These are the exact sufficient inequalities for the resonance
extensions.  The $20$ small auxiliary records for $b=4,\ldots,8$
are given in \texttt{F45.jsonl}; their individual pinned or alpha-only
roles are specified in the small-arm proof.

Finally replace a pinned four-word by
\[
 (N,1,N-1,2)\longmapsto(N+1,1,N,2,N-1).
\]
The label set gains exactly $N+1$; the old short palette
$[N-3,N]$ becomes $[N-3,N+1]$.  The cut and the other leaf labels
are unchanged.  Thus every pinned four-arm seed supplies its
minimum-five seed with the same resonance interface.  The $216$
ray certificates and their stated finite complements are therefore
the full finite input for both minimum-four and minimum-five bands.

\subsection{Equal-even repairs and the near EEO family}
\label{templates:equal-even}

Let $r$ be odd, $r\ge5$.  Start from the explicitly proved equal-odd
source $\Sp{(2r+1)^3}$, with arm order Walecki, A, C as in
Proposition~\ref{ok:equal-odd-three}.  Delete the last vertex of each arm;
their labels are $5r+3,5r+1,5r+2$.  Decrease the remaining labels
above $5r+3$ by three.  The shortened source has labels $[0,6r]$,
cut $3r$, three lengths $2r$, and defect
\[
 +\{4r+1,4r+2\}-\{4r,2r+4\}.
\]
Apply the lower and upper label cycles
\[
 (r-1,\ r+2,\ 3r-2,\ r+1,\ r),\qquad
 (4r-1,\ 4r+2,\ 4r+1,\ 4r+3),
\]
where each displayed value maps to the next, cyclically.  These
permutations stay within their respective classes.  The resulting
source has the exact defect
\begin{equation}\label{templates:equal-defect}
 +\{4r,2r+2,2r+1,2r+1\}
 -\{4,2r+5,4r-3,4r-2\}.
\end{equation}
Only the incident edges of the moved labels need to be inventoried
to obtain this identity.  On W their depths are $2r-2,2r-1,2r$;
on A they are $1,\ldots,6,2r-2,2r-1,2r$; on C they are
$1,2,2r$.  The explicit word coordinates prove that this is the
complete affected set.

For $r=4q+1$ and $r=4q+3$, $q\ge3$, the two entries of
\texttt{E.json} replace respectively $16$ and $15$ edges.  Their
full edge and source-position tables are printed below.  The signed
identity \eqref{templates:color-identity} cancels
\eqref{templates:equal-defect}; the weighted graph has three paths
of length $2r$.  The six small equal sources with lengths
$6,8,10,14,18,22$ are given explicitly in \texttt{Esmall.json}.
Balanced subdivision from these sources and the two rays gives
every even equal length at least six, as proved in the equal-triple
appendix.

The alternative $15$-edge table \texttt{Alt.json} uses the same
fixed-label source, but only $r=4q+1$, $q\ge3$.  It has the same
color and arm-length verification and an additional endpoint check:
the hub branch beginning at $H=5r+1$ ends at label $r$ and has
length $2r$.  The standard E recipe for $r=4q+3$ has this endpoint
property as well.  It is checked by following that particular
weighted hub branch and summing its exact affine weights.

\begin{lemma}\label{templates:reverse-extend}
Let a root-zero alpha spider have $N$ edges and an even arm whose
first and last labels are $H$ and $\ell$, with $H+\ell=N+1$.
Replacing that word by $(N+1,\operatorname{rev}(P))$ extends the
arm by one edge and preserves root-zero alpha labeling.
\end{lemma}
\begin{proof}
The new root color is $N+1$.  The next edge has color
$N+1-\ell=H$, replacing the former root color; all old internal
colors are unchanged.  The new label is the new maximum.  Reversing
an even word begins in the lower class, so prepending the new upper
vertex restores the alternating classes without changing the cut.
\end{proof}

Apply the lemma with $N=6r$, $H=5r+1$, and $\ell=r$ to obtain
$\Sp{2r,2r,2r+1}$.  The remaining parameters are fully specified:
the standard E recipe also works at the finite half-lengths
$r=7,11$.  For $r=3$, the equal-six arm
$(15,6,12,7,14,4)$ has endpoint sum $19=6r+1$, so the same
extension lemma applies with $H=15$, $\ell=4$.
Alt works directly at $r=9$; $a=8,10$ are literal near profiles.
If $a\equiv0\pmod4$
and $a\ge12$, apply Corollary~\ref{common:balanced-alpha-matching}
to the balanced equal source with lengths $a/2$. The unique-odd-arm
contree description in Corollary~\ref{common:unique-odd-contree}
identifies its matching lift with $\Sp{a,a,a+1}$.  Thus
\[
 a=8,\quad a=10,\quad a\ge12\ (4\mid a),\quad
 a\equiv6\pmod8,\quad a\ge18\ (a\equiv2\pmod8)
\]
are a disjoint exhaustive partition of the even $a\ge6$.
The literal targets at $a=6,8,10,14,18,22,26$ are also recorded
in \texttt{Nsmall.jsonl}.  They supplement the stated universal
proof rather than replacing it.

\subsection{Four equal odd arms}
\label{templates:four-odd}\label{construction:equal-odd-alpha4}

Write $s=4q+r$, $r\in\{0,1,2,3\}$, $q\ge6$, and let
$L=2s+1$.  Let $R_{s,i}(d)$, $i=0,1,2$, be the fully specified
equal-odd three-arm words of length $2s+1$ in the equal-triple
appendix, with $i=0$ their Walecki arm.  The unpermuted four-arm
source has the following complete coordinate rule for $d\ge1$:
\begin{equation}\label{templates:four-source}
 \begin{aligned}
 F_0(s,d)&=\begin{cases}
 8s+4-(d-1)/2,&d\text{ odd},\\ d/2,&d\text{ even},
 \end{cases}\\
 F_{i+1}(s,d)&=s+R_{s,i}(2s+2-d),\qquad i=0,1,2.
 \end{aligned}
\end{equation}
Set $F_i(s,0)=0$.  Interchange the two pairs of upper labels
$5s+1\leftrightarrow7s+1$ and $5s+2\leftrightarrow7s+2$.
These swaps preserve the label set, classes, and four source path
lengths.  Before the swaps the root-color defect is
$+\{6s+1,6s+2\}-\{4s+1,4s+2\}$.  Inventorying every edge incident
to the four moved labels gives the swapped seed defect
\begin{equation}\label{templates:four-defect}
 +\{4s-2,4s-1,4s,4s+1,6s+2\}
 -\{2s+1,2s+2,6s-2,6s-1,6s\}.
\end{equation}
The two labels $5s+1,5s+2$ are leaves and have one incident edge
each; $7s+1,7s+2$ have two.  Their complete local source coordinates
prove that no incident color is omitted in this calculation.
Explicitly, the six affected edges and colors are
\begin{center}\small
\begin{tabular}{llrr}
Old edge & New edge & Old color & New color\\\hline
$(3s-1,5s+1)$ & $(3s-1,7s+1)$ & $2s+2$ & $4s+2$\\
$(3s+1,5s+2)$ & $(3s+1,7s+2)$ & $2s+1$ & $4s+1$\\
$(s+1,7s+2)$ & $(s+1,5s+2)$ & $6s+1$ & $4s+1$\\
$(s+2,7s+2)$ & $(s+2,5s+2)$ & $6s$ & $4s$\\
$(s+2,7s+1)$ & $(s+2,5s+1)$ & $6s-1$ & $4s-1$\\
$(s+3,7s+1)$ & $(s+3,5s+1)$ & $6s-2$ & $4s-2$.
\end{tabular}\end{center}
The first two edges are the leaf edges of source arms two and three.
The last four occur on source arm one, at the reversals of depths
two through six of the original Walecki arm.  Thus the table follows
directly from \eqref{templates:four-source} and is exhaustive.  Adding
its signed changes to the pre-swap defect proves
\eqref{templates:four-defect}.

For each $r$, file \texttt{Qr.json} records the actual removed and
added edges and their positions in this swapped seed.  The four
trade sizes are
\begin{center}
\begin{tabular}{r|r|r|r}
$r$ & $L$ & Removed/added edges & Contracted vertices\\\hline
0 & $8q+1$ & $40/40$ &63\\
1 & $8q+3$ & $51/51$ &76\\
2 & $8q+5$ & $48/48$ &73\\
3 & $8q+7$ & $38/38$ &62.
\end{tabular}
\end{center}
For every entry, source-coordinate identities hold for all $q\ge6$,
the exact added/removed degree counters agree, and the affine color
calculation cancels \eqref{templates:four-defect}.  The full position
lists and both edge columns are printed below.  Weighted traversal
uses every contracted vertex and returns four paths of length
$8q+2r+1$, each as an affine identity.  The cut is $4s$ and the
vertex labels are $[0,8s+4]$.  Theorem~\ref{templates:soundness}
therefore proves root-zero alpha labeling for every odd $L\ge49$.

The remaining odd $L=7,9,\ldots,47$ have $21$ full literal records
in \texttt{Qsmall.jsonl}, checked by vertex, color, and alpha
inventories.  Together they give the equal-four input for every odd
$L\ge7$.  These finite records and the four rays are distinct
obligations; proving a ray does not certify its small-parameter
complement.

\subsection{Exact ray thresholds}
\label{templates:thresholds}
\begin{longtable}{rrrrr}
$b\bmod8$ & $\delta$ & $h_0$ & $b_0$ & Edges removed/added\\\hline\endhead
0 & 1 & 5 & 40 & 25\\
0 & 2 & 5 & 40 & 23\\
1 & 0 & 5 & 41 & 23\\
1 & 1 & 5 & 41 & 31\\
2 & 1 & 5 & 42 & 29\\
2 & 2 & 5 & 42 & 28\\
3 & 0 & 4 & 35 & 26\\
3 & 1 & 4 & 35 & 25\\
4 & 1 & 4 & 36 & 17\\
4 & 2 & 4 & 36 & 22\\
5 & 0 & 4 & 37 & 26\\
5 & 1 & 2 & 21 & 17\\
6 & 1 & 4 & 38 & 23\\
6 & 2 & 3 & 30 & 15\\
7 & 0 & 2 & 23 & 14\\
7 & 1 & 2 & 23 & 16\\
\end{longtable}
For the minimum-four rays, each entry is the inclusive threshold $b_0$; $h_0=(b_0-r)/48$. A dash denotes a band not required for that residue.
\begin{longtable}{r|rrrrr}
$r=b\bmod48$ & $\delta=0$ & $1$ & $2$ & $3$ & $4$\\\hline\endhead
0 & 48 & 48 & 48 & 48 & 48\\
1 & 49 & 49 & 49 & 49 & ---\\
2 & 50 & 50 & 50 & 50 & 50\\
3 & 51 & 51 & 51 & 51 & ---\\
4 & 52 & 52 & 52 & 52 & 52\\
5 & 53 & 53 & 53 & 53 & ---\\
6 & 54 & 54 & 54 & 102 & 54\\
7 & 55 & 55 & 103 & 55 & ---\\
8 & 56 & 56 & 56 & 56 & 56\\
9 & 57 & 57 & 57 & 57 & ---\\
10 & 58 & 58 & 58 & 58 & 58\\
11 & 59 & 59 & 59 & 59 & ---\\
12 & 60 & 60 & 60 & 60 & 60\\
13 & 61 & 61 & 61 & 61 & ---\\
14 & 62 & 62 & 62 & 62 & 62\\
15 & 63 & 63 & 63 & 63 & ---\\
16 & 64 & 64 & 64 & 64 & 64\\
17 & 65 & 65 & 65 & 65 & ---\\
18 & 66 & 66 & 66 & 66 & 66\\
19 & 67 & 67 & 67 & 67 & ---\\
20 & 68 & 68 & 68 & 68 & 68\\
21 & 69 & 69 & 69 & 69 & ---\\
22 & 70 & 70 & 70 & 70 & 70\\
23 & 71 & 71 & 71 & 71 & ---\\
24 & 72 & 72 & 72 & 72 & 72\\
25 & 73 & 73 & 73 & 73 & ---\\
26 & 74 & 74 & 74 & 74 & 74\\
27 & 75 & 75 & 75 & 75 & ---\\
28 & 76 & 76 & 76 & 76 & 76\\
29 & 77 & 77 & 77 & 77 & ---\\
30 & 78 & 78 & 78 & 78 & 78\\
31 & 79 & 79 & 79 & 79 & ---\\
32 & 80 & 80 & 80 & 80 & 80\\
33 & 81 & 81 & 81 & 81 & ---\\
34 & 82 & 82 & 82 & 82 & 82\\
35 & 35 & 83 & 83 & 83 & ---\\
36 & 36 & 84 & 36 & 36 & 36\\
37 & 85 & 85 & 37 & 85 & ---\\
38 & 38 & 38 & 38 & 38 & 38\\
39 & 39 & 39 & 39 & 39 & ---\\
40 & 40 & 40 & 40 & 40 & 40\\
41 & 41 & 41 & 41 & 41 & ---\\
42 & 90 & 90 & 90 & 90 & 90\\
43 & 43 & 43 & 91 & 91 & ---\\
44 & 44 & 44 & 44 & 44 & 44\\
45 & 93 & 93 & 93 & 93 & ---\\
46 & 46 & 46 & 46 & 46 & 94\\
47 & 47 & 47 & 47 & 95 & ---\\
\end{longtable}

\subsection{The complete equal-source permutation inventory}
\label{templates:equal-local-colors}
The following are all edges incident to a moved label in the two class cycles. Coordinates and colors are exact affine expressions in $q$, $q\ge3$. Subtracting the old column and adding the new column to the shortened-source defect gives the stated fixed-label seed defect.
\subsubsection[r=4q+1]{$r=4q+1$}
\begin{longtable}{rrrr}
Source arm & Edge depth & Old color & New color\\\hline\endhead
0 & $8q$ & $16q+7$ & $16q+4$\\
0 & $8q+1$ & $16q+6$ & $16q+3$\\
0 & $8q+2$ & $16q+5$ & $16q+6$\\
1 & $1$ & $16q+5$ & $16q+7$\\
1 & $2$ & $8q+4$ & $8q+6$\\
1 & $3$ & $8q+2$ & $8q+5$\\
1 & $4$ & $8q+1$ & $8q+4$\\
1 & $5$ & $8q+5$ & $8q+1$\\
1 & $6$ & $8q+7$ & $8q+3$\\
1 & $8q$ & $16q+1$ & $8q+3$\\
1 & $8q+1$ & $16q+2$ & $8q+4$\\
1 & $8q+2$ & $16q+3$ & $16q+4$\\
2 & $1$ & $16q+6$ & $16q+5$\\
2 & $2$ & $8q+3$ & $8q+2$\\
2 & $8q+2$ & $4$ & $8q+3$\\
\end{longtable}
\subsubsection[r=4q+3]{$r=4q+3$}
\begin{longtable}{rrrr}
Source arm & Edge depth & Old color & New color\\\hline\endhead
0 & $8q+4$ & $16q+15$ & $16q+12$\\
0 & $8q+5$ & $16q+14$ & $16q+11$\\
0 & $8q+6$ & $16q+13$ & $16q+14$\\
1 & $1$ & $16q+13$ & $16q+15$\\
1 & $2$ & $8q+8$ & $8q+10$\\
1 & $3$ & $8q+6$ & $8q+9$\\
1 & $4$ & $8q+5$ & $8q+8$\\
1 & $5$ & $8q+9$ & $8q+5$\\
1 & $6$ & $8q+11$ & $8q+7$\\
1 & $8q+4$ & $16q+9$ & $8q+7$\\
1 & $8q+5$ & $16q+10$ & $8q+8$\\
1 & $8q+6$ & $16q+11$ & $16q+12$\\
2 & $1$ & $16q+14$ & $16q+13$\\
2 & $2$ & $8q+7$ & $8q+6$\\
2 & $8q+6$ & $4$ & $8q+7$\\
\end{longtable}

\subsection{Complete equal-source and four-odd repair tables}
\label{templates:small-tables}
The vertex and depth conventions are those of the preceding schema. All edges in each remove column are replaced collectively by its add column.
\subsubsection[Equal-even source, r=4q+1, q >= 3]{Equal-even source, $r=4q+1$, $q\ge3$}
\label{templates:equal-r1}
\begingroup\small
\begin{longtable}{ll}
Remove & Add\\\hline\endhead
$(0,4r+3)$ & $(0,5r+1)$\\
$(\frac{1}{2}r-\frac{1}{2},\frac{11}{2}r+\frac{1}{2})$ & $(\frac{1}{2}r-\frac{1}{2},\frac{9}{2}r+\frac{3}{2})$\\
$(r-1,5r+1)$ & $(r-1,5r+2)$\\
$(r,5r)$ & $(r,4r+2)$\\
$(r+1,3r+2)$ & $(r+1,5r)$\\
$(r+2,5r+1)$ & $(r+2,5r-1)$\\
$(r+2,5r+2)$ & $(r+2,5r)$\\
$(\frac{3}{2}r-\frac{1}{2},\frac{9}{2}r+\frac{3}{2})$ & $(\frac{3}{2}r-\frac{1}{2},\frac{7}{2}r+\frac{5}{2})$\\
$(\frac{3}{2}r+\frac{1}{2},\frac{9}{2}r+\frac{3}{2})$ & $(\frac{3}{2}r+\frac{1}{2},\frac{11}{2}r+\frac{1}{2})$\\
$(2r-2,4r-1)$ & $(2r-2,4r+3)$\\
$(2r-1,4r+2)$ & $(2r-1,4r-1)$\\
$(2r,4r+2)$ & $(2r,5r+1)$\\
$(2r+1,4r+1)$ & $(2r+1,4r+2)$\\
$(\frac{5}{2}r-\frac{1}{2},\frac{7}{2}r+\frac{5}{2})$ & $(\frac{5}{2}r-\frac{1}{2},\frac{9}{2}r+\frac{3}{2})$\\
$(3r-2,5r-1)$ & $(3r-2,3r+2)$\\
$(3r-2,5r)$ & $(3r-2,4r+1)$\\
\end{longtable}
Rows form one collective edge trade; a row does not assert individual color preservation.
\begin{longtable}{rrr}
Vertex label & Source arm & Source depth\\\hline\endhead
$0$ & $\mathrm{hub}$ & $0$\\
$4r+3$ & 1 & $1$\\
$\frac{1}{2}r-\frac{1}{2}$ & 0 & $r-1$\\
$\frac{11}{2}r+\frac{1}{2}$ & 0 & $r$\\
$r-1$ & 0 & $2r$\\
$5r+1$ & 0 & $2r-1$\\
$r$ & 1 & $2r$\\
$5r$ & 1 & $2r-1$\\
$r+1$ & 2 & $2r$\\
$3r+2$ & 2 & $2r-1$\\
$r+2$ & 0 & $2r-2$\\
$5r+2$ & 0 & $2r-3$\\
$\frac{3}{2}r-\frac{1}{2}$ & 1 & $r+3$\\
$\frac{9}{2}r+\frac{3}{2}$ & 1 & $r+2$\\
$\frac{3}{2}r+\frac{1}{2}$ & 1 & $r+1$\\
$2r-2$ & 1 & $6$\\
$4r-1$ & 1 & $5$\\
$2r-1$ & 1 & $2$\\
$4r+2$ & 1 & $3$\\
$2r$ & 1 & $4$\\
$2r+1$ & 2 & $2$\\
$4r+1$ & 2 & $1$\\
$\frac{5}{2}r-\frac{1}{2}$ & 2 & $r-1$\\
$\frac{7}{2}r+\frac{5}{2}$ & 2 & $r$\\
$3r-2$ & 1 & $2r-2$\\
$5r-1$ & 1 & $2r-3$\\
\end{longtable}
\endgroup

\subsubsection[Equal-even source, r=4q+3, q >= 3]{Equal-even source, $r=4q+3$, $q\ge3$}
\label{templates:equal-r3}
\begingroup\small
\begin{longtable}{ll}
Remove & Add\\\hline\endhead
$(0,4r+3)$ & $(0,5r+1)$\\
$(\frac{1}{2}r-\frac{1}{2},\frac{11}{2}r+\frac{1}{2})$ & $(\frac{1}{2}r-\frac{1}{2},\frac{9}{2}r+\frac{3}{2})$\\
$(r-1,5r+1)$ & $(r-1,5r+2)$\\
$(r,5r)$ & $(r,4r+2)$\\
$(r+1,3r+2)$ & $(r+1,5r)$\\
$(r+2,5r+1)$ & $(r+2,5r-1)$\\
$(r+2,5r+2)$ & $(r+2,5r)$\\
$(\frac{3}{2}r-\frac{1}{2},\frac{9}{2}r+\frac{3}{2})$ & $(\frac{3}{2}r-\frac{1}{2},\frac{7}{2}r+\frac{3}{2})$\\
$(\frac{3}{2}r+\frac{1}{2},\frac{9}{2}r+\frac{1}{2})$ & $(\frac{3}{2}r+\frac{1}{2},\frac{11}{2}r+\frac{1}{2})$\\
$(2r-2,4r-1)$ & $(2r-2,4r+3)$\\
$(2r,4r+2)$ & $(2r,4r+1)$\\
$(2r+1,4r+1)$ & $(2r+1,5r+1)$\\
$(\frac{5}{2}r+\frac{1}{2},\frac{7}{2}r+\frac{3}{2})$ & $(\frac{5}{2}r+\frac{1}{2},\frac{9}{2}r+\frac{1}{2})$\\
$(3r-2,5r-1)$ & $(3r-2,3r+2)$\\
$(3r-2,5r)$ & $(3r-2,4r-1)$\\
\end{longtable}
Rows form one collective edge trade; a row does not assert individual color preservation.
\begin{longtable}{rrr}
Vertex label & Source arm & Source depth\\\hline\endhead
$0$ & $\mathrm{hub}$ & $0$\\
$4r+3$ & 1 & $1$\\
$\frac{1}{2}r-\frac{1}{2}$ & 0 & $r-1$\\
$\frac{11}{2}r+\frac{1}{2}$ & 0 & $r$\\
$r-1$ & 0 & $2r$\\
$5r+1$ & 0 & $2r-1$\\
$r$ & 1 & $2r$\\
$5r$ & 1 & $2r-1$\\
$r+1$ & 2 & $2r$\\
$3r+2$ & 2 & $2r-1$\\
$r+2$ & 0 & $2r-2$\\
$5r+2$ & 0 & $2r-3$\\
$\frac{3}{2}r-\frac{1}{2}$ & 1 & $r+3$\\
$\frac{9}{2}r+\frac{3}{2}$ & 1 & $r+2$\\
$\frac{3}{2}r+\frac{1}{2}$ & 1 & $r+1$\\
$\frac{9}{2}r+\frac{1}{2}$ & 1 & $r$\\
$2r-2$ & 1 & $6$\\
$4r-1$ & 1 & $5$\\
$2r$ & 1 & $4$\\
$4r+2$ & 1 & $3$\\
$2r+1$ & 2 & $2$\\
$4r+1$ & 2 & $1$\\
$\frac{5}{2}r+\frac{1}{2}$ & 2 & $r+1$\\
$\frac{7}{2}r+\frac{3}{2}$ & 2 & $r+2$\\
$3r-2$ & 1 & $2r-2$\\
$5r-1$ & 1 & $2r-3$\\
\end{longtable}
\endgroup

\subsubsection[Alternative complementary endpoints, r=4q+1, q >= 3]{Alternative complementary endpoints, $r=4q+1$, $q\ge3$}
\label{templates:alt}
\begingroup\small
\begin{longtable}{ll}
Remove & Add\\\hline\endhead
$(0,4r+3)$ & $(0,5r+1)$\\
$(\frac{1}{2}r-\frac{1}{2},\frac{11}{2}r+\frac{1}{2})$ & $(\frac{1}{2}r-\frac{1}{2},\frac{9}{2}r+\frac{3}{2})$\\
$(r-1,5r+1)$ & $(r-1,5r+2)$\\
$(r,5r)$ & $(r,4r+2)$\\
$(r+1,3r+2)$ & $(r+1,5r)$\\
$(r+2,5r+1)$ & $(r+2,5r-1)$\\
$(r+2,5r+2)$ & $(r+2,5r)$\\
$(\frac{3}{2}r-\frac{1}{2},\frac{9}{2}r+\frac{3}{2})$ & $(\frac{3}{2}r-\frac{1}{2},\frac{7}{2}r+\frac{5}{2})$\\
$(\frac{3}{2}r+\frac{1}{2},\frac{9}{2}r+\frac{1}{2})$ & $(\frac{3}{2}r+\frac{1}{2},\frac{11}{2}r+\frac{1}{2})$\\
$(2r-2,4r-1)$ & $(2r-2,4r+3)$\\
$(2r-1,4r+2)$ & $(2r-1,4r-1)$\\
$(2r+1,4r+1)$ & $(2r+1,5r+1)$\\
$(\frac{5}{2}r-\frac{1}{2},\frac{7}{2}r+\frac{5}{2})$ & $(\frac{5}{2}r-\frac{1}{2},\frac{9}{2}r+\frac{1}{2})$\\
$(3r-2,5r-1)$ & $(3r-2,3r+2)$\\
$(3r-2,5r)$ & $(3r-2,4r+1)$\\
\end{longtable}
Rows form one collective edge trade; a row does not assert individual color preservation.
\begin{longtable}{rrr}
Vertex label & Source arm & Source depth\\\hline\endhead
$0$ & $\mathrm{hub}$ & $0$\\
$\frac{1}{2}r-\frac{1}{2}$ & 0 & $r-1$\\
$r-1$ & 0 & $2r$\\
$r$ & 1 & $2r$\\
$r+1$ & 2 & $2r$\\
$r+2$ & 0 & $2r-2$\\
$\frac{3}{2}r-\frac{1}{2}$ & 1 & $r+3$\\
$\frac{3}{2}r+\frac{1}{2}$ & 1 & $r+1$\\
$2r-2$ & 1 & $6$\\
$2r-1$ & 1 & $2$\\
$2r+1$ & 2 & $2$\\
$\frac{5}{2}r-\frac{1}{2}$ & 2 & $r-1$\\
$3r-2$ & 1 & $2r-2$\\
$3r+2$ & 2 & $2r-1$\\
$\frac{7}{2}r+\frac{5}{2}$ & 2 & $r$\\
$4r-1$ & 1 & $5$\\
$4r+1$ & 2 & $1$\\
$4r+2$ & 1 & $3$\\
$4r+3$ & 1 & $1$\\
$\frac{9}{2}r+\frac{1}{2}$ & 1 & $r$\\
$\frac{9}{2}r+\frac{3}{2}$ & 1 & $r+2$\\
$5r-1$ & 1 & $2r-3$\\
$5r$ & 1 & $2r-1$\\
$5r+1$ & 0 & $2r-1$\\
$5r+2$ & 0 & $2r-3$\\
$\frac{11}{2}r+\frac{1}{2}$ & 0 & $r$\\
\end{longtable}
\endgroup

\subsubsection[Four equal odd arms, s=4q+0, q >= 6]{Four equal odd arms, $s=4q+0$, $q\ge6$}
\label{templates:four-r0}
\begingroup\small
\begin{longtable}{ll}
Remove & Add\\\hline\endhead
$(0,24q+2)$ & $(0,24q-1)$\\
$(4q+1,20q+2)$ & $(4q+1,24q+4)$\\
$(4q+2,20q+1)$ & $(4q+2,24q-2)$\\
$(4q+2,20q+2)$ & $(4q+2,28q+2)$\\
$(4q+3,20q+1)$ & $(4q+3,28q+1)$\\
$(6q,26q+3)$ & $(6q,26q+2)$\\
$(6q+1,26q+3)$ & $(6q+1,26q+1)$\\
$(6q+2,26q+2)$ & $(6q+2,26q)$\\
$(6q+3,26q+1)$ & $(6q+3,22q)$\\
$(6q+4,26q)$ & $(6q+4,22q-1)$\\
$(8q-1,24q+4)$ & $(8q-1,24q+3)$\\
$(8q,24q+3)$ & $(8q,24q+1)$\\
$(8q,24q+4)$ & $(8q,24q+2)$\\
$(8q+1,24q+1)$ & $(8q+1,20q+2)$\\
$(8q+2,24q-1)$ & $(8q+2,24q+2)$\\
$(8q+3,24q-2)$ & $(8q+3,20q-3)$\\
$(10q-2,22q+3)$ & $(10q-2,26q+3)$\\
$(10q,22q+1)$ & $(10q,26q+3)$\\
$(10q+1,22q)$ & $(10q+1,22q-1)$\\
$(10q+2,22q-1)$ & $(10q+2,22q-2)$\\
$(10q+2,22q)$ & $(10q+2,22q+1)$\\
$(10q+3,22q-1)$ & $(10q+3,22q)$\\
$(10q+4,22q-2)$ & $(10q+4,18q-1)$\\
$(12q-1,20q-1)$ & $(12q-1,20q+3)$\\
$(12q-1,28q+1)$ & $(12q-1,24q+4)$\\
$(12q,20q-1)$ & $(12q,20q+1)$\\
$(12q,20q+3)$ & $(12q,20q+2)$\\
$(12q+1,20q-3)$ & $(12q+1,20q)$\\
$(12q+1,28q+2)$ & $(12q+1,20q+1)$\\
$(12q+2,20q-3)$ & $(12q+2,20q-1)$\\
$(12q+3,20q)$ & $(12q+3,20q-1)$\\
$(14q,18q-1)$ & $(14q,22q+3)$\\
$(16q-5,16q+3)$ & $(16q-5,16q+4)$\\
$(16q-4,16q+3)$ & $(16q-4,16q+4)$\\
$(16q-3,16q+6)$ & $(16q-3,16q+3)$\\
$(16q-2,16q+4)$ & $(16q-2,16q+1)$\\
$(16q-2,24q+2)$ & $(16q-2,20q-3)$\\
$(16q-1,16q+1)$ & $(16q-1,16q+3)$\\
$(16q-1,16q+2)$ & $(16q-1,16q+6)$\\
$(16q,16q+4)$ & $(16q,16q+2)$\\
\end{longtable}
Rows form one collective edge trade; a row does not assert individual color preservation.
\begin{longtable}{rrr}
Vertex label & Source arm & Source depth\\\hline\endhead
$0$ & $\mathrm{hub}$ & $0$\\
$4q+1$ & 1 & $8q$\\
$4q+2$ & 1 & $8q-2$\\
$4q+3$ & 1 & $8q-4$\\
$6q$ & 1 & $4q+2$\\
$6q+1$ & 1 & $4q$\\
$6q+2$ & 1 & $4q-2$\\
$6q+3$ & 1 & $4q-4$\\
$6q+4$ & 1 & $4q-6$\\
$8q-1$ & 1 & $4$\\
$8q$ & 1 & $2$\\
$8q+1$ & 2 & $2$\\
$8q+2$ & 2 & $4$\\
$8q+3$ & 2 & $6$\\
$10q-2$ & 2 & $4q-4$\\
$10q$ & 2 & $4q$\\
$10q+1$ & 2 & $4q+2$\\
$10q+2$ & 2 & $4q+4$\\
$10q+3$ & 2 & $4q+6$\\
$10q+4$ & 2 & $4q+8$\\
$12q-1$ & 2 & $8q$\\
$12q$ & 2 & $8q-2$\\
$12q+1$ & 3 & $8q$\\
$12q+2$ & 3 & $8q-2$\\
$12q+3$ & 3 & $8q-4$\\
$14q$ & 3 & $4q+2$\\
$16q-5$ & 3 & $12$\\
$16q-4$ & 3 & $10$\\
$16q-3$ & 3 & $8$\\
$16q-2$ & 3 & $2$\\
$16q-1$ & 3 & $6$\\
$16q$ & 3 & $4$\\
$16q+1$ & 3 & $5$\\
$16q+2$ & 3 & $7$\\
$16q+3$ & 3 & $11$\\
$16q+4$ & 3 & $3$\\
$16q+6$ & 3 & $9$\\
$18q-1$ & 3 & $4q+3$\\
$20q-3$ & 3 & $8q-1$\\
$20q-1$ & 2 & $8q-1$\\
$20q$ & 3 & $8q-3$\\
$20q+1$ & 1 & $8q-3$\\
$20q+2$ & 1 & $8q-1$\\
$20q+3$ & 2 & $8q-3$\\
$22q-2$ & 2 & $4q+7$\\
$22q-1$ & 2 & $4q+5$\\
$22q$ & 2 & $4q+3$\\
$22q+1$ & 2 & $4q+1$\\
$22q+3$ & 2 & $4q-3$\\
$24q-2$ & 2 & $7$\\
$24q-1$ & 2 & $5$\\
$24q+1$ & 2 & $1$\\
$24q+2$ & 3 & $1$\\
$24q+3$ & 1 & $1$\\
$24q+4$ & 1 & $3$\\
$26q$ & 1 & $4q-5$\\
$26q+1$ & 1 & $4q-3$\\
$26q+2$ & 1 & $4q-1$\\
$26q+3$ & 1 & $4q+1$\\
$28q+1$ & 2 & $8q+1$\\
$28q+2$ & 3 & $8q+1$\\
\end{longtable}
\endgroup

\subsubsection[Four equal odd arms, s=4q+1, q >= 6]{Four equal odd arms, $s=4q+1$, $q\ge6$}
\label{templates:four-r1}
\begingroup\small
\begin{longtable}{ll}
Remove & Add\\\hline\endhead
$(0,24q+8)$ & $(0,24q+5)$\\
$(4q+2,20q+7)$ & $(4q+2,24q+8)$\\
$(4q+3,20q+6)$ & $(4q+3,24q+4)$\\
$(4q+3,20q+7)$ & $(4q+3,28q+9)$\\
$(4q+4,20q+6)$ & $(4q+4,28q+8)$\\
$(6q+3,26q+8)$ & $(6q+3,26q+6)$\\
$(6q+3,26q+9)$ & $(6q+3,26q+7)$\\
$(6q+4,26q+7)$ & $(6q+4,22q+5)$\\
$(6q+4,26q+8)$ & $(6q+4,26q+9)$\\
$(6q+5,26q+6)$ & $(6q+5,22q+4)$\\
$(8q+2,24q+9)$ & $(8q+2,24q+7)$\\
$(8q+3,24q+7)$ & $(8q+3,20q+8)$\\
$(8q+4,24q+5)$ & $(8q+4,24q+8)$\\
$(8q+5,24q+4)$ & $(8q+5,20q+2)$\\
$(10q+1,22q+9)$ & $(10q+1,26q+8)$\\
$(10q+2,22q+7)$ & $(10q+2,22q+6)$\\
$(10q+2,22q+8)$ & $(10q+2,26q+8)$\\
$(10q+3,22q+7)$ & $(10q+3,22q+9)$\\
$(10q+4,22q+5)$ & $(10q+4,18q+4)$\\
$(10q+4,22q+6)$ & $(10q+4,22q+3)$\\
$(10q+5,22q+4)$ & $(10q+5,22q+7)$\\
$(10q+6,22q+3)$ & $(10q+6,22q+7)$\\
$(12q+1,20q+8)$ & $(12q+1,24q+9)$\\
$(12q+2,20q+4)$ & $(12q+2,20q+7)$\\
$(12q+2,28q+8)$ & $(12q+2,20q+8)$\\
$(12q+3,20q+4)$ & $(12q+3,20q+6)$\\
$(12q+3,20q+8)$ & $(12q+3,20q+7)$\\
$(12q+4,20q+2)$ & $(12q+4,20q+5)$\\
$(12q+4,28q+9)$ & $(12q+4,20q+6)$\\
$(12q+5,20q+5)$ & $(12q+5,20q+4)$\\
$(12q+6,20q+5)$ & $(12q+6,20q+4)$\\
$(14q,18q+6)$ & $(14q,18q+7)$\\
$(14q+1,18q+6)$ & $(14q+1,22q+8)$\\
$(14q+2,18q+4)$ & $(14q+2,18q+6)$\\
$(14q+2,18q+9)$ & $(14q+2,18q+7)$\\
$(14q+3,18q+4)$ & $(14q+3,18q+5)$\\
$(14q+3,18q+7)$ & $(14q+3,18q+9)$\\
$(14q+4,18q+2)$ & $(14q+4,18q+3)$\\
$(14q+4,18q+7)$ & $(14q+4,18q+4)$\\
$(14q+5,18q+2)$ & $(14q+5,18q)$\\
$(14q+5,18q+5)$ & $(14q+5,18q+6)$\\
$(14q+6,18q)$ & $(14q+6,18q+2)$\\
$(14q+6,18q+5)$ & $(14q+6,18q+3)$\\
$(14q+7,18q+3)$ & $(14q+7,18q+5)$\\
$(14q+8,18q+3)$ & $(14q+8,18q+2)$\\
$(16q-2,16q+8)$ & $(16q-2,16q+9)$\\
$(16q,16q+11)$ & $(16q,16q+6)$\\
$(16q+1,16q+7)$ & $(16q+1,16q+11)$\\
$(16q+2,16q+6)$ & $(16q+2,16q+7)$\\
$(16q+2,24q+8)$ & $(16q+2,20q+5)$\\
$(16q+4,16q+9)$ & $(16q+4,16q+8)$\\
\end{longtable}
Rows form one collective edge trade; a row does not assert individual color preservation.
\begin{longtable}{rrr}
Vertex label & Source arm & Source depth\\\hline\endhead
$0$ & $\mathrm{hub}$ & $0$\\
$4q+2$ & 1 & $8q+2$\\
$4q+3$ & 1 & $8q$\\
$4q+4$ & 1 & $8q-2$\\
$6q+3$ & 1 & $4q$\\
$6q+4$ & 1 & $4q-2$\\
$6q+5$ & 1 & $4q-4$\\
$8q+2$ & 1 & $2$\\
$8q+3$ & 2 & $2$\\
$8q+4$ & 2 & $4$\\
$8q+5$ & 2 & $6$\\
$10q+1$ & 2 & $4q-2$\\
$10q+2$ & 2 & $4q$\\
$10q+3$ & 2 & $4q+2$\\
$10q+4$ & 2 & $4q+4$\\
$10q+5$ & 2 & $4q+6$\\
$10q+6$ & 2 & $4q+8$\\
$12q+1$ & 2 & $8q-2$\\
$12q+2$ & 2 & $8q+2$\\
$12q+3$ & 2 & $8q$\\
$12q+4$ & 3 & $8q+2$\\
$12q+5$ & 3 & $8q$\\
$12q+6$ & 3 & $8q-2$\\
$14q$ & 3 & $4q+10$\\
$14q+1$ & 3 & $4q+8$\\
$14q+2$ & 3 & $4q+6$\\
$14q+3$ & 3 & $4q+4$\\
$14q+4$ & 3 & $4q+2$\\
$14q+5$ & 3 & $4q$\\
$14q+6$ & 3 & $4q-2$\\
$14q+7$ & 3 & $4q-4$\\
$14q+8$ & 3 & $4q-6$\\
$16q-2$ & 3 & $14$\\
$16q$ & 3 & $10$\\
$16q+1$ & 3 & $8$\\
$16q+2$ & 3 & $2$\\
$16q+4$ & 3 & $6$\\
$16q+6$ & 3 & $3$\\
$16q+7$ & 3 & $9$\\
$16q+8$ & 3 & $13$\\
$16q+9$ & 3 & $7$\\
$16q+11$ & 3 & $11$\\
$18q$ & 3 & $4q-3$\\
$18q+2$ & 3 & $4q+1$\\
$18q+3$ & 3 & $4q-5$\\
$18q+4$ & 3 & $4q+5$\\
$18q+5$ & 3 & $4q-1$\\
$18q+6$ & 3 & $4q+9$\\
$18q+7$ & 3 & $4q+3$\\
$18q+9$ & 3 & $4q+7$\\
$20q+2$ & 3 & $8q+1$\\
$20q+4$ & 2 & $8q+1$\\
$20q+5$ & 3 & $8q-1$\\
$20q+6$ & 1 & $8q-1$\\
$20q+7$ & 1 & $8q+1$\\
$20q+8$ & 2 & $8q-1$\\
$22q+3$ & 2 & $4q+9$\\
$22q+4$ & 2 & $4q+7$\\
$22q+5$ & 2 & $4q+5$\\
$22q+6$ & 2 & $4q+3$\\
$22q+7$ & 2 & $4q+1$\\
$22q+8$ & 2 & $4q-1$\\
$22q+9$ & 2 & $4q-3$\\
$24q+4$ & 2 & $7$\\
$24q+5$ & 2 & $5$\\
$24q+7$ & 2 & $1$\\
$24q+8$ & 3 & $1$\\
$24q+9$ & 1 & $1$\\
$26q+6$ & 1 & $4q-5$\\
$26q+7$ & 1 & $4q-3$\\
$26q+8$ & 1 & $4q-1$\\
$26q+9$ & 1 & $4q+1$\\
$28q+8$ & 2 & $8q+3$\\
$28q+9$ & 3 & $8q+3$\\
\end{longtable}
\endgroup

\subsubsection[Four equal odd arms, s=4q+2, q >= 6]{Four equal odd arms, $s=4q+2$, $q\ge6$}
\label{templates:four-r2}
\begingroup\small
\begin{longtable}{ll}
Remove & Add\\\hline\endhead
$(0,24q+13)$ & $(0,24q+10)$\\
$(4q+3,20q+12)$ & $(4q+3,24q+16)$\\
$(4q+3,28q+17)$ & $(4q+3,28q+16)$\\
$(4q+4,20q+11)$ & $(4q+4,24q+11)$\\
$(4q+4,20q+12)$ & $(4q+4,28q+15)$\\
$(4q+5,20q+11)$ & $(4q+5,28q+17)$\\
$(6q+3,26q+16)$ & $(6q+3,26q+14)$\\
$(6q+4,26q+15)$ & $(6q+4,26q+14)$\\
$(6q+5,26q+14)$ & $(6q+5,22q+10)$\\
$(6q+5,26q+15)$ & $(6q+5,26q+13)$\\
$(6q+6,26q+13)$ & $(6q+6,22q+9)$\\
$(6q+6,26q+14)$ & $(6q+6,26q+15)$\\
$(8q+3,24q+16)$ & $(8q+3,24q+13)$\\
$(8q+4,24q+16)$ & $(8q+4,24q+13)$\\
$(8q+5,24q+13)$ & $(8q+5,20q+12)$\\
$(8q+6,24q+11)$ & $(8q+6,24q+14)$\\
$(8q+7,24q+10)$ & $(8q+7,20q+7)$\\
$(10q+3,22q+14)$ & $(10q+3,26q+15)$\\
$(10q+3,22q+15)$ & $(10q+3,26q+16)$\\
$(10q+4,22q+13)$ & $(10q+4,22q+15)$\\
$(10q+5,22q+12)$ & $(10q+5,22q+14)$\\
$(10q+7,22q+10)$ & $(10q+7,18q+9)$\\
$(10q+9,22q+9)$ & $(10q+9,22q+12)$\\
$(12q+4,20q+13)$ & $(12q+4,24q+16)$\\
$(12q+5,20q+9)$ & $(12q+5,20q+12)$\\
$(12q+5,28q+15)$ & $(12q+5,20q+13)$\\
$(12q+6,20q+9)$ & $(12q+6,20q+10)$\\
$(12q+6,20q+13)$ & $(12q+6,20q+11)$\\
$(12q+7,20q+7)$ & $(12q+7,20q+10)$\\
$(12q+7,28q+16)$ & $(12q+7,20q+13)$\\
$(12q+8,20q+10)$ & $(12q+8,20q+9)$\\
$(12q+9,20q+10)$ & $(12q+9,20q+9)$\\
$(14q+4,18q+11)$ & $(14q+4,22q+13)$\\
$(14q+5,18q+9)$ & $(14q+5,18q+12)$\\
$(14q+6,18q+9)$ & $(14q+6,18q+7)$\\
$(14q+7,18q+7)$ & $(14q+7,18q+10)$\\
$(14q+7,18q+12)$ & $(14q+7,18q+11)$\\
$(14q+8,18q+7)$ & $(14q+8,18q+8)$\\
$(14q+9,18q+10)$ & $(14q+9,18q+8)$\\
$(14q+10,18q+8)$ & $(14q+10,18q+7)$\\
$(14q+11,18q+8)$ & $(14q+11,18q+9)$\\
$(16q+3,16q+11)$ & $(16q+3,16q+12)$\\
$(16q+4,16q+11)$ & $(16q+4,16q+12)$\\
$(16q+5,16q+14)$ & $(16q+5,16q+11)$\\
$(16q+6,16q+12)$ & $(16q+6,16q+10)$\\
$(16q+6,24q+14)$ & $(16q+6,20q+11)$\\
$(16q+7,16q+10)$ & $(16q+7,16q+14)$\\
$(16q+8,16q+12)$ & $(16q+8,16q+11)$\\
\end{longtable}
Rows form one collective edge trade; a row does not assert individual color preservation.
\begin{longtable}{rrr}
Vertex label & Source arm & Source depth\\\hline\endhead
$0$ & $\mathrm{hub}$ & $0$\\
$4q+3$ & 1 & $8q+4$\\
$4q+4$ & 1 & $8q+2$\\
$4q+5$ & 1 & $8q$\\
$6q+3$ & 1 & $4q+4$\\
$6q+4$ & 1 & $4q+2$\\
$6q+5$ & 1 & $4q$\\
$6q+6$ & 1 & $4q-2$\\
$8q+3$ & 1 & $4$\\
$8q+4$ & 1 & $2$\\
$8q+5$ & 2 & $2$\\
$8q+6$ & 2 & $4$\\
$8q+7$ & 2 & $6$\\
$10q+3$ & 2 & $4q-2$\\
$10q+4$ & 2 & $4q$\\
$10q+5$ & 2 & $4q+2$\\
$10q+7$ & 2 & $4q+6$\\
$10q+9$ & 2 & $4q+10$\\
$12q+4$ & 2 & $8q$\\
$12q+5$ & 2 & $8q+4$\\
$12q+6$ & 2 & $8q+2$\\
$12q+7$ & 3 & $8q+4$\\
$12q+8$ & 3 & $8q+2$\\
$12q+9$ & 3 & $8q$\\
$14q+4$ & 3 & $4q+10$\\
$14q+5$ & 3 & $4q+8$\\
$14q+6$ & 3 & $4q+6$\\
$14q+7$ & 3 & $4q+4$\\
$14q+8$ & 3 & $4q+2$\\
$14q+9$ & 3 & $4q$\\
$14q+10$ & 3 & $4q-2$\\
$14q+11$ & 3 & $4q-4$\\
$16q+3$ & 3 & $12$\\
$16q+4$ & 3 & $10$\\
$16q+5$ & 3 & $8$\\
$16q+6$ & 3 & $2$\\
$16q+7$ & 3 & $6$\\
$16q+8$ & 3 & $4$\\
$16q+10$ & 3 & $7$\\
$16q+11$ & 3 & $11$\\
$16q+12$ & 3 & $3$\\
$16q+14$ & 3 & $9$\\
$18q+7$ & 3 & $4q+3$\\
$18q+8$ & 3 & $4q-3$\\
$18q+9$ & 3 & $4q+7$\\
$18q+10$ & 3 & $4q+1$\\
$18q+11$ & 3 & $4q+11$\\
$18q+12$ & 3 & $4q+5$\\
$20q+7$ & 3 & $8q+3$\\
$20q+9$ & 2 & $8q+3$\\
$20q+10$ & 3 & $8q+1$\\
$20q+11$ & 1 & $8q+1$\\
$20q+12$ & 1 & $8q+3$\\
$20q+13$ & 2 & $8q+1$\\
$22q+9$ & 2 & $4q+9$\\
$22q+10$ & 2 & $4q+7$\\
$22q+12$ & 2 & $4q+3$\\
$22q+13$ & 2 & $4q+1$\\
$22q+14$ & 2 & $4q-1$\\
$22q+15$ & 2 & $4q-3$\\
$24q+10$ & 2 & $7$\\
$24q+11$ & 2 & $5$\\
$24q+13$ & 2 & $1$\\
$24q+14$ & 3 & $1$\\
$24q+16$ & 1 & $3$\\
$26q+13$ & 1 & $4q-3$\\
$26q+14$ & 1 & $4q-1$\\
$26q+15$ & 1 & $4q+1$\\
$26q+16$ & 1 & $4q+3$\\
$28q+15$ & 2 & $8q+5$\\
$28q+16$ & 3 & $8q+5$\\
$28q+17$ & 1 & $8q+5$\\
\end{longtable}
\endgroup

\subsubsection[Four equal odd arms, s=4q+3, q >= 6]{Four equal odd arms, $s=4q+3$, $q\ge6$}
\label{templates:four-r3}
\begingroup\small
\begin{longtable}{ll}
Remove & Add\\\hline\endhead
$(0,24q+19)$ & $(0,24q+16)$\\
$(4q+4,20q+17)$ & $(4q+4,24q+22)$\\
$(4q+4,28q+24)$ & $(4q+4,28q+22)$\\
$(4q+5,20q+16)$ & $(4q+5,24q+17)$\\
$(4q+5,20q+17)$ & $(4q+5,28q+24)$\\
$(4q+6,20q+16)$ & $(4q+6,28q+23)$\\
$(6q+5,26q+23)$ & $(6q+5,22q+19)$\\
$(6q+8,26q+20)$ & $(6q+8,22q+15)$\\
$(8q+5,24q+22)$ & $(8q+5,24q+20)$\\
$(8q+6,24q+21)$ & $(8q+6,24q+19)$\\
$(8q+8,24q+17)$ & $(8q+8,20q+16)$\\
$(8q+9,24q+16)$ & $(8q+9,20q+12)$\\
$(10q+6,22q+20)$ & $(10q+6,26q+23)$\\
$(10q+7,22q+19)$ & $(10q+7,22q+17)$\\
$(10q+8,22q+18)$ & $(10q+8,22q+20)$\\
$(10q+9,22q+17)$ & $(10q+9,22q+15)$\\
$(10q+10,22q+15)$ & $(10q+10,18q+14)$\\
$(10q+10,22q+16)$ & $(10q+10,22q+14)$\\
$(10q+11,22q+14)$ & $(10q+11,22q+16)$\\
$(10q+11,22q+15)$ & $(10q+11,26q+20)$\\
$(12q+7,20q+18)$ & $(12q+7,24q+21)$\\
$(12q+8,20q+14)$ & $(12q+8,20q+16)$\\
$(12q+8,28q+22)$ & $(12q+8,20q+18)$\\
$(12q+9,20q+14)$ & $(12q+9,20q+15)$\\
$(12q+10,20q+12)$ & $(12q+10,20q+15)$\\
$(12q+10,28q+23)$ & $(12q+10,20q+17)$\\
$(12q+11,20q+15)$ & $(12q+11,20q+14)$\\
$(12q+12,20q+15)$ & $(12q+12,20q+14)$\\
$(14q+7,18q+16)$ & $(14q+7,22q+18)$\\
$(14q+8,18q+14)$ & $(14q+8,18q+17)$\\
$(14q+9,18q+14)$ & $(14q+9,18q+12)$\\
$(14q+10,18q+12)$ & $(14q+10,18q+15)$\\
$(14q+10,18q+17)$ & $(14q+10,18q+16)$\\
$(14q+11,18q+12)$ & $(14q+11,18q+13)$\\
$(14q+12,18q+15)$ & $(14q+12,18q+13)$\\
$(14q+13,18q+13)$ & $(14q+13,18q+12)$\\
$(14q+14,18q+13)$ & $(14q+14,18q+14)$\\
$(16q+10,24q+20)$ & $(16q+10,20q+17)$\\
\end{longtable}
Rows form one collective edge trade; a row does not assert individual color preservation.
\begin{longtable}{rrr}
Vertex label & Source arm & Source depth\\\hline\endhead
$0$ & $\mathrm{hub}$ & $0$\\
$4q+4$ & 1 & $8q+6$\\
$4q+5$ & 1 & $8q+4$\\
$4q+6$ & 1 & $8q+2$\\
$6q+5$ & 1 & $4q+4$\\
$6q+8$ & 1 & $4q-2$\\
$8q+5$ & 1 & $4$\\
$8q+6$ & 1 & $2$\\
$8q+8$ & 2 & $4$\\
$8q+9$ & 2 & $6$\\
$10q+6$ & 2 & $4q$\\
$10q+7$ & 2 & $4q+2$\\
$10q+8$ & 2 & $4q+4$\\
$10q+9$ & 2 & $4q+6$\\
$10q+10$ & 2 & $4q+8$\\
$10q+11$ & 2 & $4q+10$\\
$12q+7$ & 2 & $8q+2$\\
$12q+8$ & 2 & $8q+6$\\
$12q+9$ & 2 & $8q+4$\\
$12q+10$ & 3 & $8q+6$\\
$12q+11$ & 3 & $8q+4$\\
$12q+12$ & 3 & $8q+2$\\
$14q+7$ & 3 & $4q+12$\\
$14q+8$ & 3 & $4q+10$\\
$14q+9$ & 3 & $4q+8$\\
$14q+10$ & 3 & $4q+6$\\
$14q+11$ & 3 & $4q+4$\\
$14q+12$ & 3 & $4q+2$\\
$14q+13$ & 3 & $4q$\\
$14q+14$ & 3 & $4q-2$\\
$16q+10$ & 3 & $2$\\
$18q+12$ & 3 & $4q+5$\\
$18q+13$ & 3 & $4q-1$\\
$18q+14$ & 3 & $4q+9$\\
$18q+15$ & 3 & $4q+3$\\
$18q+16$ & 3 & $4q+13$\\
$18q+17$ & 3 & $4q+7$\\
$20q+12$ & 3 & $8q+5$\\
$20q+14$ & 2 & $8q+5$\\
$20q+15$ & 3 & $8q+3$\\
$20q+16$ & 1 & $8q+3$\\
$20q+17$ & 1 & $8q+5$\\
$20q+18$ & 2 & $8q+3$\\
$22q+14$ & 2 & $4q+11$\\
$22q+15$ & 2 & $4q+9$\\
$22q+16$ & 2 & $4q+7$\\
$22q+17$ & 2 & $4q+5$\\
$22q+18$ & 2 & $4q+3$\\
$22q+19$ & 2 & $4q+1$\\
$22q+20$ & 2 & $4q-1$\\
$24q+16$ & 2 & $7$\\
$24q+17$ & 2 & $5$\\
$24q+19$ & 2 & $1$\\
$24q+20$ & 3 & $1$\\
$24q+21$ & 1 & $1$\\
$24q+22$ & 1 & $3$\\
$26q+20$ & 1 & $4q-1$\\
$26q+23$ & 1 & $4q+5$\\
$28q+22$ & 2 & $8q+7$\\
$28q+23$ & 3 & $8q+7$\\
$28q+24$ & 1 & $8q+7$\\
\end{longtable}
\endgroup


\subsection{Complete finite input ledger}
\label{templates:ledger}
The aliases below are the files in \texttt{data/affine/} of the supporting repository.
\begingroup\footnotesize\begin{longtable}{lr}
Input & SHA-256 (first and last halves)\\\hline\endhead
\texttt{S2o.jsonl} & \texttt{9f1f9481846b81323feae9828c6fcb71}\\
 & \texttt{4ba7f2be21cecbd9d09c385930005b2f}\\[3pt]
\texttt{S2e.jsonl} & \texttt{841cc9bfe43dd853a4c1185176a729f2}\\
 & \texttt{c30ca678ceb7abe6abd2bfda48f69ac3}\\[3pt]
\texttt{S4.jsonl} & \texttt{60b39f453abdc6fa130429c1df76457b}\\
 & \texttt{ca0a6294f9f93ea3bb43bfe00de3d229}\\[3pt]
\texttt{E.json} & \texttt{0eefe79d40a9df8aae39132d4b782681}\\
 & \texttt{bba19bf801c37e9bfd7bc8fd18dc37dc}\\[3pt]
\texttt{Alt.json} & \texttt{655183b5952106706804d3010b3be969}\\
 & \texttt{010cd5ff8a18d94871c55d56299804a7}\\[3pt]
\texttt{Q0.json} & \texttt{07b2679b47f3246af2d8776674679fed}\\
 & \texttt{d4535a262ee4ea2787c4c58afb7ebc9a}\\[3pt]
\texttt{Q1.json} & \texttt{d55627eca8bf2b27ff66657bb6e282a9}\\
 & \texttt{c9aae4ee841f344854883d002c675131}\\[3pt]
\texttt{Q2.json} & \texttt{0f165f3cea4ed0d71cf105cadffe8f2c}\\
 & \texttt{2ebb031e871ffecd6531db9490806317}\\[3pt]
\texttt{Q3.json} & \texttt{99cdc386444ce9f992a7266cd5e59642}\\
 & \texttt{34c5045cc6c7277ceb2f04154ae13582}\\[3pt]
\texttt{P.json} & \texttt{9b5d5d28c24bc60a52a63039c17c6786}\\
 & \texttt{bdcb0e13c135eac070e65a0aabccd859}\\[3pt]
\texttt{F2.jsonl} & \texttt{0e67af0313b85faac029eb0aca15c6a6}\\
 & \texttt{9b2c8b65a3df07eeac8927a6663aa909}\\[3pt]
\texttt{F3s.jsonl} & \texttt{28520669a674bf72dc9c2cb854ac4fe6}\\
 & \texttt{abed99be794d026b441453269db400ad}\\[3pt]
\texttt{F3p.jsonl} & \texttt{114e36ecda7dcba694628338bdde0aa3}\\
 & \texttt{cbf11eb4c002d708780303b39712eba7}\\[3pt]
\texttt{F4.jsonl} & \texttt{3df86054602a198e57ebe553fff61de7}\\
 & \texttt{26b77153e2bb8d08205170e430f59532}\\[3pt]
\texttt{F45.jsonl} & \texttt{0561aacc63e517c97f9c76e9f54f5394}\\
 & \texttt{90c6032551a66c89eeff40cadc4f65c8}\\[3pt]
\texttt{Nsmall.jsonl} & \texttt{600fd54ba59790f4315dd2c97f024e4e}\\
 & \texttt{c5d9c3a60520b7bd118df184fc4e9190}\\[3pt]
\texttt{Qsmall.jsonl} & \texttt{383ee17c3b04c362f435cc80aae39541}\\
 & \texttt{3d4483824d693bac23873d4c4fcb975b}\\[3pt]
\texttt{Esmall.json} & \texttt{8c9d47773f4c44b29b1a64aa8e95814c}\\
 & \texttt{122f6341838d040bd92bfa3c258767ef}\\[3pt]
\end{longtable}\endgroup


\section{Explicit finite construction certificates}\label{app:finite-bases}
All words below run outward from a common hub of label zero, unless a different hub is stated. Lengths count edges. A semicolon separates arms; it is not an additional vertex. Every printed labeling is checked by the elementary inventory algorithm of Section~\ref{app:template-certificates}.
\subsection{Small root-zero alpha bases}\label{app:small-finite-alpha}
These 24 bases suffice for every alpha component in the 311 seven-core decompositions. They also contain the 23 small positive alpha profiles required in the short-arm case analysis. The cut is the sum of the half-lengths. The printed words themselves, rather than an existence inference from bounded searching, are finite proof inputs.
\small
\begin{longtable}{@{}p{0.1\linewidth}@{\hspace{.02\linewidth}}p{0.24\linewidth}@{\hspace{.02\linewidth}}p{0.1\linewidth}@{\hspace{.02\linewidth}}p{0.5\linewidth}@{}}
\caption{Literal root-zero alpha labelings used by the small-core proof.}\label{tab:finite-alpha}\\
\toprule
ID & Lengths & Cut & Arm words\\\midrule
\endfirsthead
\toprule
ID & Lengths & Cut & Arm words\\\midrule
\endhead
\midrule\multicolumn{4}{r}{\footnotesize Continued on the next page}\\\endfoot
\bottomrule\endlastfoot
\(A_{1}\) & \((2,2,5)\) & \(4\) & \(\begin{array}{l}(9,\allowbreak 1)\\(6,\allowbreak 2)\\(7,\allowbreak 4,\allowbreak 5,\allowbreak 3,\allowbreak 8)\end{array}\)\\[.5em]
\(A_{2}\) & \((2,3,4)\) & \(4\) & \(\begin{array}{l}(9,\allowbreak 1)\\(7,\allowbreak 3,\allowbreak 8)\\(6,\allowbreak 4,\allowbreak 5,\allowbreak 2)\end{array}\)\\[.5em]
\(A_{3}\) & \((2,3,5)\) & \(4\) & \(\begin{array}{l}(10,\allowbreak 1)\\(6,\allowbreak 2,\allowbreak 9)\\(8,\allowbreak 3,\allowbreak 5,\allowbreak 4,\allowbreak 7)\end{array}\)\\[.5em]
\(A_{4}\) & \((2,3,6)\) & \(5\) & \(\begin{array}{l}(11,\allowbreak 1)\\(8,\allowbreak 4,\allowbreak 10)\\(9,\allowbreak 2,\allowbreak 7,\allowbreak 5,\allowbreak 6,\allowbreak 3)\end{array}\)\\[.5em]
\(A_{5}\) & \((2,4,5)\) & \(5\) & \(\begin{array}{l}(11,\allowbreak 1)\\(7,\allowbreak 5,\allowbreak 6,\allowbreak 3)\\(9,\allowbreak 4,\allowbreak 8,\allowbreak 2,\allowbreak 10)\end{array}\)\\[.5em]
\(A_{6}\) & \((2,4,6)\) & \(6\) & \(\begin{array}{l}(12,\allowbreak 1)\\(8,\allowbreak 4,\allowbreak 11,\allowbreak 2)\\(10,\allowbreak 5,\allowbreak 7,\allowbreak 6,\allowbreak 9,\allowbreak 3)\end{array}\)\\[.5em]
\(A_{7}\) & \((2,5,5)\) & \(5\) & \(\begin{array}{l}(12,\allowbreak 1)\\(10,\allowbreak 4,\allowbreak 6,\allowbreak 5,\allowbreak 8)\\(9,\allowbreak 2,\allowbreak 7,\allowbreak 3,\allowbreak 11)\end{array}\)\\[.5em]
\(A_{8}\) & \((2,5,6)\) & \(6\) & \(\begin{array}{l}(13,\allowbreak 1)\\(8,\allowbreak 5,\allowbreak 7,\allowbreak 6,\allowbreak 10)\\(11,\allowbreak 4,\allowbreak 9,\allowbreak 3,\allowbreak 12,\allowbreak 2)\end{array}\)\\[.5em]
\(A_{9}\) & \((2,6,6)\) & \(7\) & \(\begin{array}{l}(14,\allowbreak 1)\\(12,\allowbreak 3,\allowbreak 13,\allowbreak 5,\allowbreak 10,\allowbreak 4)\\(11,\allowbreak 7,\allowbreak 8,\allowbreak 6,\allowbreak 9,\allowbreak 2)\end{array}\)\\[.5em]
\(A_{10}\) & \((3,3,4)\) & \(4\) & \(\begin{array}{l}(10,\allowbreak 1,\allowbreak 9)\\(7,\allowbreak 3,\allowbreak 8)\\(6,\allowbreak 4,\allowbreak 5,\allowbreak 2)\end{array}\)\\[.5em]
\(A_{11}\) & \((3,3,5)\) & \(4\) & \(\begin{array}{l}(11,\allowbreak 1,\allowbreak 10)\\(6,\allowbreak 2,\allowbreak 9)\\(8,\allowbreak 3,\allowbreak 5,\allowbreak 4,\allowbreak 7)\end{array}\)\\[.5em]
\(A_{12}\) & \((3,3,6)\) & \(5\) & \(\begin{array}{l}(12,\allowbreak 1,\allowbreak 11)\\(8,\allowbreak 4,\allowbreak 10)\\(9,\allowbreak 2,\allowbreak 7,\allowbreak 5,\allowbreak 6,\allowbreak 3)\end{array}\)\\[.5em]
\(A_{13}\) & \((3,4,5)\) & \(5\) & \(\begin{array}{l}(12,\allowbreak 1,\allowbreak 11)\\(7,\allowbreak 5,\allowbreak 6,\allowbreak 3)\\(9,\allowbreak 4,\allowbreak 8,\allowbreak 2,\allowbreak 10)\end{array}\)\\[.5em]
\(A_{14}\) & \((3,4,6)\) & \(6\) & \(\begin{array}{l}(13,\allowbreak 1,\allowbreak 12)\\(8,\allowbreak 4,\allowbreak 11,\allowbreak 2)\\(10,\allowbreak 5,\allowbreak 7,\allowbreak 6,\allowbreak 9,\allowbreak 3)\end{array}\)\\[.5em]
\(A_{15}\) & \((3,5,5)\) & \(5\) & \(\begin{array}{l}(13,\allowbreak 1,\allowbreak 12)\\(10,\allowbreak 4,\allowbreak 6,\allowbreak 5,\allowbreak 8)\\(9,\allowbreak 2,\allowbreak 7,\allowbreak 3,\allowbreak 11)\end{array}\)\\[.5em]
\(A_{16}\) & \((3,5,6)\) & \(6\) & \(\begin{array}{l}(14,\allowbreak 1,\allowbreak 13)\\(8,\allowbreak 5,\allowbreak 7,\allowbreak 6,\allowbreak 10)\\(11,\allowbreak 4,\allowbreak 9,\allowbreak 3,\allowbreak 12,\allowbreak 2)\end{array}\)\\[.5em]
\(A_{17}\) & \((3,6,6)\) & \(7\) & \(\begin{array}{l}(15,\allowbreak 1,\allowbreak 14)\\(12,\allowbreak 3,\allowbreak 13,\allowbreak 5,\allowbreak 10,\allowbreak 4)\\(11,\allowbreak 7,\allowbreak 8,\allowbreak 6,\allowbreak 9,\allowbreak 2)\end{array}\)\\[.5em]
\(A_{18}\) & \((4,5,5)\) & \(6\) & \(\begin{array}{l}(14,\allowbreak 1,\allowbreak 13,\allowbreak 2)\\(10,\allowbreak 3,\allowbreak 8,\allowbreak 4,\allowbreak 12)\\(9,\allowbreak 6,\allowbreak 7,\allowbreak 5,\allowbreak 11)\end{array}\)\\[.5em]
\(A_{19}\) & \((4,5,6)\) & \(7\) & \(\begin{array}{l}(15,\allowbreak 1,\allowbreak 14,\allowbreak 2)\\(10,\allowbreak 7,\allowbreak 8,\allowbreak 6,\allowbreak 13)\\(11,\allowbreak 5,\allowbreak 9,\allowbreak 4,\allowbreak 12,\allowbreak 3)\end{array}\)\\[.5em]
\(A_{20}\) & \((4,6,6)\) & \(8\) & \(\begin{array}{l}(16,\allowbreak 2,\allowbreak 14,\allowbreak 1)\\(15,\allowbreak 5,\allowbreak 13,\allowbreak 6,\allowbreak 12,\allowbreak 3)\\(11,\allowbreak 7,\allowbreak 10,\allowbreak 8,\allowbreak 9,\allowbreak 4)\end{array}\)\\[.5em]
\(A_{21}\) & \((5,5,5)\) & \(6\) & \(\begin{array}{l}(15,\allowbreak 1,\allowbreak 14,\allowbreak 2,\allowbreak 13)\\(10,\allowbreak 3,\allowbreak 8,\allowbreak 4,\allowbreak 12)\\(9,\allowbreak 6,\allowbreak 7,\allowbreak 5,\allowbreak 11)\end{array}\)\\[.5em]
\(A_{22}\) & \((5,5,6)\) & \(7\) & \(\begin{array}{l}(16,\allowbreak 1,\allowbreak 15,\allowbreak 2,\allowbreak 14)\\(10,\allowbreak 7,\allowbreak 8,\allowbreak 6,\allowbreak 13)\\(11,\allowbreak 5,\allowbreak 9,\allowbreak 4,\allowbreak 12,\allowbreak 3)\end{array}\)\\[.5em]
\(A_{23}\) & \((5,6,6)\) & \(8\) & \(\begin{array}{l}(17,\allowbreak 1,\allowbreak 16,\allowbreak 3,\allowbreak 15)\\(14,\allowbreak 5,\allowbreak 13,\allowbreak 6,\allowbreak 12,\allowbreak 2)\\(11,\allowbreak 7,\allowbreak 10,\allowbreak 8,\allowbreak 9,\allowbreak 4)\end{array}\)\\[.5em]
\(A_{24}\) & \((6,6,6)\) & \(9\) & \(\begin{array}{l}(18,\allowbreak 1,\allowbreak 17,\allowbreak 5,\allowbreak 16,\allowbreak 2)\\(15,\allowbreak 6,\allowbreak 12,\allowbreak 7,\allowbreak 14,\allowbreak 4)\\(13,\allowbreak 9,\allowbreak 10,\allowbreak 8,\allowbreak 11,\allowbreak 3)\end{array}\)\\[.5em]
\end{longtable}
\normalsize
\subsection{Four-arm hub-zero bases}\label{app:four-finite-hubzero}
The following table covers exactly the 70 sorted four-tuples in $\{2,3,4,5,6\}$. A four-arm spider with a unit arm is obtained by appending a unit leaf to a hub-zero three-arm spider, so no additional finite table with a unit entry is needed.
\small
\begin{longtable}{@{}p{0.1\linewidth}@{\hspace{.02\linewidth}}p{0.24\linewidth}@{\hspace{.02\linewidth}}p{0.62\linewidth}@{}}
\caption{Complete hub-zero four-arm catalog for nonunit lengths at most six.}\label{tab:finite-four}\\
\toprule
ID & Lengths & Arm words\\\midrule
\endfirsthead
\toprule
ID & Lengths & Arm words\\\midrule
\endhead
\midrule\multicolumn{3}{r}{\footnotesize Continued on the next page}\\\endfoot
\bottomrule\endlastfoot
\(H_{1}\) & \((2,2,2,2)\) & \(\begin{array}{l}(8,\allowbreak 1)\\(6,\allowbreak 2)\\(5,\allowbreak 7)\\(3,\allowbreak 4)\end{array}\)\\[.5em]
\(H_{2}\) & \((2,2,2,3)\) & \(\begin{array}{l}(9,\allowbreak 1)\\(6,\allowbreak 8)\\(4,\allowbreak 3)\\(7,\allowbreak 2,\allowbreak 5)\end{array}\)\\[.5em]
\(H_{3}\) & \((2,2,2,4)\) & \(\begin{array}{l}(10,\allowbreak 1)\\(5,\allowbreak 7)\\(4,\allowbreak 3)\\(8,\allowbreak 2,\allowbreak 9,\allowbreak 6)\end{array}\)\\[.5em]
\(H_{4}\) & \((2,2,2,5)\) & \(\begin{array}{l}(11,\allowbreak 1)\\(6,\allowbreak 8)\\(3,\allowbreak 7)\\(9,\allowbreak 2,\allowbreak 10,\allowbreak 5,\allowbreak 4)\end{array}\)\\[.5em]
\(H_{5}\) & \((2,2,2,6)\) & \(\begin{array}{l}(12,\allowbreak 1)\\(6,\allowbreak 3)\\(5,\allowbreak 7)\\(10,\allowbreak 2,\allowbreak 11,\allowbreak 4,\allowbreak 8,\allowbreak 9)\end{array}\)\\[.5em]
\(H_{6}\) & \((2,2,3,3)\) & \(\begin{array}{l}(10,\allowbreak 1)\\(3,\allowbreak 7)\\(8,\allowbreak 2,\allowbreak 9)\\(5,\allowbreak 6,\allowbreak 4)\end{array}\)\\[.5em]
\(H_{7}\) & \((2,2,3,4)\) & \(\begin{array}{l}(11,\allowbreak 1)\\(3,\allowbreak 7)\\(5,\allowbreak 6,\allowbreak 8)\\(9,\allowbreak 2,\allowbreak 10,\allowbreak 4)\end{array}\)\\[.5em]
\(H_{8}\) & \((2,2,3,5)\) & \(\begin{array}{l}(12,\allowbreak 1)\\(3,\allowbreak 9)\\(5,\allowbreak 7,\allowbreak 6)\\(10,\allowbreak 2,\allowbreak 11,\allowbreak 4,\allowbreak 8)\end{array}\)\\[.5em]
\(H_{9}\) & \((2,2,3,6)\) & \(\begin{array}{l}(13,\allowbreak 1)\\(7,\allowbreak 5)\\(3,\allowbreak 8,\allowbreak 9)\\(11,\allowbreak 2,\allowbreak 12,\allowbreak 4,\allowbreak 10,\allowbreak 6)\end{array}\)\\[.5em]
\(H_{10}\) & \((2,2,4,4)\) & \(\begin{array}{l}(12,\allowbreak 1)\\(3,\allowbreak 8)\\(10,\allowbreak 2,\allowbreak 11,\allowbreak 4)\\(6,\allowbreak 7,\allowbreak 9,\allowbreak 5)\end{array}\)\\[.5em]
\(H_{11}\) & \((2,2,4,5)\) & \(\begin{array}{l}(13,\allowbreak 1)\\(3,\allowbreak 8)\\(7,\allowbreak 9,\allowbreak 5,\allowbreak 6)\\(11,\allowbreak 2,\allowbreak 12,\allowbreak 4,\allowbreak 10)\end{array}\)\\[.5em]
\(H_{12}\) & \((2,2,4,6)\) & \(\begin{array}{l}(14,\allowbreak 1)\\(5,\allowbreak 8)\\(6,\allowbreak 10,\allowbreak 9,\allowbreak 7)\\(12,\allowbreak 2,\allowbreak 13,\allowbreak 4,\allowbreak 11,\allowbreak 3)\end{array}\)\\[.5em]
\(H_{13}\) & \((2,2,5,5)\) & \(\begin{array}{l}(14,\allowbreak 1)\\(5,\allowbreak 6)\\(12,\allowbreak 2,\allowbreak 13,\allowbreak 4,\allowbreak 11)\\(8,\allowbreak 10,\allowbreak 7,\allowbreak 3,\allowbreak 9)\end{array}\)\\[.5em]
\(H_{14}\) & \((2,2,5,6)\) & \(\begin{array}{l}(15,\allowbreak 1)\\(5,\allowbreak 11)\\(7,\allowbreak 10,\allowbreak 6,\allowbreak 8,\allowbreak 9)\\(13,\allowbreak 2,\allowbreak 14,\allowbreak 4,\allowbreak 12,\allowbreak 3)\end{array}\)\\[.5em]
\(H_{15}\) & \((2,2,6,6)\) & \(\begin{array}{l}(16,\allowbreak 1)\\(6,\allowbreak 11)\\(14,\allowbreak 2,\allowbreak 15,\allowbreak 4,\allowbreak 13,\allowbreak 3)\\(8,\allowbreak 12,\allowbreak 5,\allowbreak 7,\allowbreak 10,\allowbreak 9)\end{array}\)\\[.5em]
\(H_{16}\) & \((2,3,3,3)\) & \(\begin{array}{l}(11,\allowbreak 1)\\(9,\allowbreak 2,\allowbreak 10)\\(6,\allowbreak 3,\allowbreak 8)\\(4,\allowbreak 5,\allowbreak 7)\end{array}\)\\[.5em]
\(H_{17}\) & \((2,3,3,4)\) & \(\begin{array}{l}(12,\allowbreak 1)\\(6,\allowbreak 9,\allowbreak 8)\\(5,\allowbreak 7,\allowbreak 3)\\(10,\allowbreak 2,\allowbreak 11,\allowbreak 4)\end{array}\)\\[.5em]
\(H_{18}\) & \((2,3,3,5)\) & \(\begin{array}{l}(13,\allowbreak 1)\\(6,\allowbreak 7,\allowbreak 9)\\(3,\allowbreak 10,\allowbreak 5)\\(11,\allowbreak 2,\allowbreak 12,\allowbreak 4,\allowbreak 8)\end{array}\)\\[.5em]
\(H_{19}\) & \((2,3,3,6)\) & \(\begin{array}{l}(14,\allowbreak 1)\\(8,\allowbreak 3,\allowbreak 5)\\(6,\allowbreak 9,\allowbreak 10)\\(12,\allowbreak 2,\allowbreak 13,\allowbreak 4,\allowbreak 11,\allowbreak 7)\end{array}\)\\[.5em]
\(H_{20}\) & \((2,3,4,4)\) & \(\begin{array}{l}(13,\allowbreak 1)\\(5,\allowbreak 8,\allowbreak 6)\\(11,\allowbreak 2,\allowbreak 12,\allowbreak 4)\\(7,\allowbreak 3,\allowbreak 9,\allowbreak 10)\end{array}\)\\[.5em]
\(H_{21}\) & \((2,3,4,5)\) & \(\begin{array}{l}(14,\allowbreak 1)\\(6,\allowbreak 9,\allowbreak 10)\\(8,\allowbreak 3,\allowbreak 7,\allowbreak 5)\\(12,\allowbreak 2,\allowbreak 13,\allowbreak 4,\allowbreak 11)\end{array}\)\\[.5em]
\(H_{22}\) & \((2,3,4,6)\) & \(\begin{array}{l}(15,\allowbreak 1)\\(6,\allowbreak 11,\allowbreak 10)\\(7,\allowbreak 9,\allowbreak 5,\allowbreak 8)\\(13,\allowbreak 2,\allowbreak 14,\allowbreak 4,\allowbreak 12,\allowbreak 3)\end{array}\)\\[.5em]
\(H_{23}\) & \((2,3,5,5)\) & \(\begin{array}{l}(15,\allowbreak 1)\\(5,\allowbreak 6,\allowbreak 8)\\(13,\allowbreak 2,\allowbreak 14,\allowbreak 4,\allowbreak 12)\\(9,\allowbreak 3,\allowbreak 10,\allowbreak 7,\allowbreak 11)\end{array}\)\\[.5em]
\(H_{24}\) & \((2,3,5,6)\) & \(\begin{array}{l}(16,\allowbreak 1)\\(6,\allowbreak 9,\allowbreak 7)\\(8,\allowbreak 12,\allowbreak 5,\allowbreak 10,\allowbreak 11)\\(14,\allowbreak 2,\allowbreak 15,\allowbreak 4,\allowbreak 13,\allowbreak 3)\end{array}\)\\[.5em]
\(H_{25}\) & \((2,3,6,6)\) & \(\begin{array}{l}(17,\allowbreak 1)\\(7,\allowbreak 8,\allowbreak 11)\\(15,\allowbreak 2,\allowbreak 16,\allowbreak 4,\allowbreak 14,\allowbreak 3)\\(9,\allowbreak 13,\allowbreak 5,\allowbreak 10,\allowbreak 12,\allowbreak 6)\end{array}\)\\[.5em]
\(H_{26}\) & \((2,4,4,4)\) & \(\begin{array}{l}(14,\allowbreak 1)\\(12,\allowbreak 2,\allowbreak 13,\allowbreak 4)\\(8,\allowbreak 3,\allowbreak 10,\allowbreak 7)\\(6,\allowbreak 5,\allowbreak 9,\allowbreak 11)\end{array}\)\\[.5em]
\(H_{27}\) & \((2,4,4,5)\) & \(\begin{array}{l}(15,\allowbreak 1)\\(9,\allowbreak 3,\allowbreak 10,\allowbreak 6)\\(5,\allowbreak 7,\allowbreak 8,\allowbreak 11)\\(13,\allowbreak 2,\allowbreak 14,\allowbreak 4,\allowbreak 12)\end{array}\)\\[.5em]
\(H_{28}\) & \((2,4,4,6)\) & \(\begin{array}{l}(16,\allowbreak 1)\\(8,\allowbreak 12,\allowbreak 6,\allowbreak 9)\\(7,\allowbreak 5,\allowbreak 10,\allowbreak 11)\\(14,\allowbreak 2,\allowbreak 15,\allowbreak 4,\allowbreak 13,\allowbreak 3)\end{array}\)\\[.5em]
\(H_{29}\) & \((2,4,5,5)\) & \(\begin{array}{l}(16,\allowbreak 1)\\(6,\allowbreak 8,\allowbreak 5,\allowbreak 9)\\(14,\allowbreak 2,\allowbreak 15,\allowbreak 4,\allowbreak 13)\\(10,\allowbreak 3,\allowbreak 11,\allowbreak 12,\allowbreak 7)\end{array}\)\\[.5em]
\(H_{30}\) & \((2,4,5,6)\) & \(\begin{array}{l}(17,\allowbreak 1)\\(6,\allowbreak 8,\allowbreak 11,\allowbreak 10)\\(9,\allowbreak 13,\allowbreak 5,\allowbreak 12,\allowbreak 7)\\(15,\allowbreak 2,\allowbreak 16,\allowbreak 4,\allowbreak 14,\allowbreak 3)\end{array}\)\\[.5em]
\(H_{31}\) & \((2,4,6,6)\) & \(\begin{array}{l}(18,\allowbreak 1)\\(7,\allowbreak 13,\allowbreak 12,\allowbreak 8)\\(16,\allowbreak 2,\allowbreak 17,\allowbreak 4,\allowbreak 15,\allowbreak 3)\\(10,\allowbreak 5,\allowbreak 14,\allowbreak 6,\allowbreak 9,\allowbreak 11)\end{array}\)\\[.5em]
\(H_{32}\) & \((2,5,5,5)\) & \(\begin{array}{l}(17,\allowbreak 1)\\(15,\allowbreak 2,\allowbreak 16,\allowbreak 4,\allowbreak 14)\\(11,\allowbreak 3,\allowbreak 12,\allowbreak 5,\allowbreak 9)\\(6,\allowbreak 7,\allowbreak 10,\allowbreak 8,\allowbreak 13)\end{array}\)\\[.5em]
\(H_{33}\) & \((2,5,5,6)\) & \(\begin{array}{l}(18,\allowbreak 1)\\(10,\allowbreak 5,\allowbreak 14,\allowbreak 6,\allowbreak 12)\\(7,\allowbreak 8,\allowbreak 11,\allowbreak 13,\allowbreak 9)\\(16,\allowbreak 2,\allowbreak 17,\allowbreak 4,\allowbreak 15,\allowbreak 3)\end{array}\)\\[.5em]
\(H_{34}\) & \((2,5,6,6)\) & \(\begin{array}{l}(19,\allowbreak 1)\\(7,\allowbreak 10,\allowbreak 8,\allowbreak 12,\allowbreak 13)\\(17,\allowbreak 2,\allowbreak 18,\allowbreak 4,\allowbreak 16,\allowbreak 3)\\(11,\allowbreak 5,\allowbreak 15,\allowbreak 6,\allowbreak 14,\allowbreak 9)\end{array}\)\\[.5em]
\(H_{35}\) & \((2,6,6,6)\) & \(\begin{array}{l}(20,\allowbreak 1)\\(18,\allowbreak 2,\allowbreak 19,\allowbreak 4,\allowbreak 17,\allowbreak 3)\\(12,\allowbreak 5,\allowbreak 16,\allowbreak 6,\allowbreak 15,\allowbreak 10)\\(8,\allowbreak 7,\allowbreak 13,\allowbreak 9,\allowbreak 11,\allowbreak 14)\end{array}\)\\[.5em]
\(H_{36}\) & \((3,3,3,3)\) & \(\begin{array}{l}(12,\allowbreak 1,\allowbreak 11)\\(9,\allowbreak 2,\allowbreak 10)\\(6,\allowbreak 3,\allowbreak 8)\\(4,\allowbreak 5,\allowbreak 7)\end{array}\)\\[.5em]
\(H_{37}\) & \((3,3,3,4)\) & \(\begin{array}{l}(13,\allowbreak 1,\allowbreak 12)\\(6,\allowbreak 9,\allowbreak 8)\\(5,\allowbreak 7,\allowbreak 3)\\(10,\allowbreak 2,\allowbreak 11,\allowbreak 4)\end{array}\)\\[.5em]
\(H_{38}\) & \((3,3,3,5)\) & \(\begin{array}{l}(14,\allowbreak 1,\allowbreak 13)\\(6,\allowbreak 7,\allowbreak 9)\\(3,\allowbreak 10,\allowbreak 5)\\(11,\allowbreak 2,\allowbreak 12,\allowbreak 4,\allowbreak 8)\end{array}\)\\[.5em]
\(H_{39}\) & \((3,3,3,6)\) & \(\begin{array}{l}(15,\allowbreak 1,\allowbreak 14)\\(8,\allowbreak 3,\allowbreak 5)\\(6,\allowbreak 9,\allowbreak 10)\\(12,\allowbreak 2,\allowbreak 13,\allowbreak 4,\allowbreak 11,\allowbreak 7)\end{array}\)\\[.5em]
\(H_{40}\) & \((3,3,4,4)\) & \(\begin{array}{l}(14,\allowbreak 1,\allowbreak 13)\\(5,\allowbreak 8,\allowbreak 6)\\(11,\allowbreak 2,\allowbreak 12,\allowbreak 4)\\(7,\allowbreak 3,\allowbreak 9,\allowbreak 10)\end{array}\)\\[.5em]
\(H_{41}\) & \((3,3,4,5)\) & \(\begin{array}{l}(15,\allowbreak 1,\allowbreak 14)\\(6,\allowbreak 9,\allowbreak 10)\\(8,\allowbreak 3,\allowbreak 7,\allowbreak 5)\\(12,\allowbreak 2,\allowbreak 13,\allowbreak 4,\allowbreak 11)\end{array}\)\\[.5em]
\(H_{42}\) & \((3,3,4,6)\) & \(\begin{array}{l}(16,\allowbreak 1,\allowbreak 15)\\(6,\allowbreak 11,\allowbreak 10)\\(7,\allowbreak 9,\allowbreak 5,\allowbreak 8)\\(13,\allowbreak 2,\allowbreak 14,\allowbreak 4,\allowbreak 12,\allowbreak 3)\end{array}\)\\[.5em]
\(H_{43}\) & \((3,3,5,5)\) & \(\begin{array}{l}(16,\allowbreak 1,\allowbreak 15)\\(5,\allowbreak 6,\allowbreak 8)\\(13,\allowbreak 2,\allowbreak 14,\allowbreak 4,\allowbreak 12)\\(9,\allowbreak 3,\allowbreak 10,\allowbreak 7,\allowbreak 11)\end{array}\)\\[.5em]
\(H_{44}\) & \((3,3,5,6)\) & \(\begin{array}{l}(17,\allowbreak 1,\allowbreak 16)\\(6,\allowbreak 9,\allowbreak 7)\\(8,\allowbreak 12,\allowbreak 5,\allowbreak 10,\allowbreak 11)\\(14,\allowbreak 2,\allowbreak 15,\allowbreak 4,\allowbreak 13,\allowbreak 3)\end{array}\)\\[.5em]
\(H_{45}\) & \((3,3,6,6)\) & \(\begin{array}{l}(18,\allowbreak 1,\allowbreak 17)\\(7,\allowbreak 8,\allowbreak 11)\\(15,\allowbreak 2,\allowbreak 16,\allowbreak 4,\allowbreak 14,\allowbreak 3)\\(9,\allowbreak 13,\allowbreak 5,\allowbreak 10,\allowbreak 12,\allowbreak 6)\end{array}\)\\[.5em]
\(H_{46}\) & \((3,4,4,4)\) & \(\begin{array}{l}(15,\allowbreak 1,\allowbreak 14)\\(12,\allowbreak 2,\allowbreak 13,\allowbreak 4)\\(8,\allowbreak 3,\allowbreak 10,\allowbreak 7)\\(6,\allowbreak 5,\allowbreak 9,\allowbreak 11)\end{array}\)\\[.5em]
\(H_{47}\) & \((3,4,4,5)\) & \(\begin{array}{l}(16,\allowbreak 1,\allowbreak 15)\\(9,\allowbreak 3,\allowbreak 10,\allowbreak 6)\\(5,\allowbreak 7,\allowbreak 8,\allowbreak 11)\\(13,\allowbreak 2,\allowbreak 14,\allowbreak 4,\allowbreak 12)\end{array}\)\\[.5em]
\(H_{48}\) & \((3,4,4,6)\) & \(\begin{array}{l}(17,\allowbreak 1,\allowbreak 16)\\(8,\allowbreak 12,\allowbreak 6,\allowbreak 9)\\(7,\allowbreak 5,\allowbreak 10,\allowbreak 11)\\(14,\allowbreak 2,\allowbreak 15,\allowbreak 4,\allowbreak 13,\allowbreak 3)\end{array}\)\\[.5em]
\(H_{49}\) & \((3,4,5,5)\) & \(\begin{array}{l}(17,\allowbreak 1,\allowbreak 16)\\(6,\allowbreak 8,\allowbreak 5,\allowbreak 9)\\(14,\allowbreak 2,\allowbreak 15,\allowbreak 4,\allowbreak 13)\\(10,\allowbreak 3,\allowbreak 11,\allowbreak 12,\allowbreak 7)\end{array}\)\\[.5em]
\(H_{50}\) & \((3,4,5,6)\) & \(\begin{array}{l}(18,\allowbreak 1,\allowbreak 17)\\(6,\allowbreak 8,\allowbreak 11,\allowbreak 10)\\(9,\allowbreak 13,\allowbreak 5,\allowbreak 12,\allowbreak 7)\\(15,\allowbreak 2,\allowbreak 16,\allowbreak 4,\allowbreak 14,\allowbreak 3)\end{array}\)\\[.5em]
\(H_{51}\) & \((3,4,6,6)\) & \(\begin{array}{l}(19,\allowbreak 1,\allowbreak 18)\\(7,\allowbreak 13,\allowbreak 12,\allowbreak 8)\\(16,\allowbreak 2,\allowbreak 17,\allowbreak 4,\allowbreak 15,\allowbreak 3)\\(10,\allowbreak 5,\allowbreak 14,\allowbreak 6,\allowbreak 9,\allowbreak 11)\end{array}\)\\[.5em]
\(H_{52}\) & \((3,5,5,5)\) & \(\begin{array}{l}(18,\allowbreak 1,\allowbreak 17)\\(15,\allowbreak 2,\allowbreak 16,\allowbreak 4,\allowbreak 14)\\(11,\allowbreak 3,\allowbreak 12,\allowbreak 5,\allowbreak 9)\\(6,\allowbreak 7,\allowbreak 10,\allowbreak 8,\allowbreak 13)\end{array}\)\\[.5em]
\(H_{53}\) & \((3,5,5,6)\) & \(\begin{array}{l}(19,\allowbreak 1,\allowbreak 18)\\(10,\allowbreak 5,\allowbreak 14,\allowbreak 6,\allowbreak 12)\\(7,\allowbreak 8,\allowbreak 11,\allowbreak 13,\allowbreak 9)\\(16,\allowbreak 2,\allowbreak 17,\allowbreak 4,\allowbreak 15,\allowbreak 3)\end{array}\)\\[.5em]
\(H_{54}\) & \((3,5,6,6)\) & \(\begin{array}{l}(20,\allowbreak 1,\allowbreak 19)\\(7,\allowbreak 10,\allowbreak 8,\allowbreak 12,\allowbreak 13)\\(17,\allowbreak 2,\allowbreak 18,\allowbreak 4,\allowbreak 16,\allowbreak 3)\\(11,\allowbreak 5,\allowbreak 15,\allowbreak 6,\allowbreak 14,\allowbreak 9)\end{array}\)\\[.5em]
\(H_{55}\) & \((3,6,6,6)\) & \(\begin{array}{l}(21,\allowbreak 1,\allowbreak 20)\\(18,\allowbreak 2,\allowbreak 19,\allowbreak 4,\allowbreak 17,\allowbreak 3)\\(12,\allowbreak 5,\allowbreak 16,\allowbreak 6,\allowbreak 15,\allowbreak 10)\\(8,\allowbreak 7,\allowbreak 13,\allowbreak 9,\allowbreak 11,\allowbreak 14)\end{array}\)\\[.5em]
\(H_{56}\) & \((4,4,4,4)\) & \(\begin{array}{l}(16,\allowbreak 1,\allowbreak 15,\allowbreak 2)\\(12,\allowbreak 3,\allowbreak 14,\allowbreak 4)\\(8,\allowbreak 13,\allowbreak 9,\allowbreak 6)\\(7,\allowbreak 5,\allowbreak 11,\allowbreak 10)\end{array}\)\\[.5em]
\(H_{57}\) & \((4,4,4,5)\) & \(\begin{array}{l}(17,\allowbreak 1,\allowbreak 16,\allowbreak 2)\\(9,\allowbreak 14,\allowbreak 7,\allowbreak 11)\\(6,\allowbreak 5,\allowbreak 8,\allowbreak 10)\\(13,\allowbreak 3,\allowbreak 15,\allowbreak 4,\allowbreak 12)\end{array}\)\\[.5em]
\(H_{58}\) & \((4,4,4,6)\) & \(\begin{array}{l}(18,\allowbreak 1,\allowbreak 17,\allowbreak 2)\\(10,\allowbreak 15,\allowbreak 12,\allowbreak 6)\\(7,\allowbreak 11,\allowbreak 9,\allowbreak 8)\\(14,\allowbreak 3,\allowbreak 16,\allowbreak 4,\allowbreak 13,\allowbreak 5)\end{array}\)\\[.5em]
\(H_{59}\) & \((4,4,5,5)\) & \(\begin{array}{l}(18,\allowbreak 1,\allowbreak 17,\allowbreak 2)\\(8,\allowbreak 6,\allowbreak 7,\allowbreak 11)\\(14,\allowbreak 3,\allowbreak 16,\allowbreak 4,\allowbreak 13)\\(10,\allowbreak 15,\allowbreak 9,\allowbreak 12,\allowbreak 5)\end{array}\)\\[.5em]
\(H_{60}\) & \((4,4,5,6)\) & \(\begin{array}{l}(19,\allowbreak 1,\allowbreak 18,\allowbreak 2)\\(7,\allowbreak 9,\allowbreak 10,\allowbreak 13)\\(11,\allowbreak 16,\allowbreak 8,\allowbreak 12,\allowbreak 6)\\(15,\allowbreak 3,\allowbreak 17,\allowbreak 4,\allowbreak 14,\allowbreak 5)\end{array}\)\\[.5em]
\(H_{61}\) & \((4,4,6,6)\) & \(\begin{array}{l}(20,\allowbreak 1,\allowbreak 19,\allowbreak 2)\\(7,\allowbreak 11,\allowbreak 13,\allowbreak 10)\\(16,\allowbreak 3,\allowbreak 18,\allowbreak 4,\allowbreak 15,\allowbreak 5)\\(12,\allowbreak 6,\allowbreak 14,\allowbreak 9,\allowbreak 8,\allowbreak 17)\end{array}\)\\[.5em]
\(H_{62}\) & \((4,5,5,5)\) & \(\begin{array}{l}(19,\allowbreak 1,\allowbreak 18,\allowbreak 2)\\(15,\allowbreak 3,\allowbreak 17,\allowbreak 4,\allowbreak 14)\\(11,\allowbreak 5,\allowbreak 13,\allowbreak 6,\allowbreak 8)\\(9,\allowbreak 10,\allowbreak 7,\allowbreak 12,\allowbreak 16)\end{array}\)\\[.5em]
\(H_{63}\) & \((4,5,5,6)\) & \(\begin{array}{l}(20,\allowbreak 1,\allowbreak 19,\allowbreak 2)\\(12,\allowbreak 6,\allowbreak 14,\allowbreak 9,\allowbreak 11)\\(7,\allowbreak 8,\allowbreak 17,\allowbreak 13,\allowbreak 10)\\(16,\allowbreak 3,\allowbreak 18,\allowbreak 4,\allowbreak 15,\allowbreak 5)\end{array}\)\\[.5em]
\(H_{64}\) & \((4,5,6,6)\) & \(\begin{array}{l}(21,\allowbreak 1,\allowbreak 20,\allowbreak 2)\\(10,\allowbreak 14,\allowbreak 11,\allowbreak 9,\allowbreak 8)\\(17,\allowbreak 3,\allowbreak 19,\allowbreak 4,\allowbreak 16,\allowbreak 5)\\(13,\allowbreak 6,\allowbreak 15,\allowbreak 7,\allowbreak 12,\allowbreak 18)\end{array}\)\\[.5em]
\(H_{65}\) & \((4,6,6,6)\) & \(\begin{array}{l}(22,\allowbreak 1,\allowbreak 21,\allowbreak 2)\\(18,\allowbreak 3,\allowbreak 20,\allowbreak 4,\allowbreak 17,\allowbreak 5)\\(14,\allowbreak 6,\allowbreak 16,\allowbreak 7,\allowbreak 12,\allowbreak 19)\\(11,\allowbreak 15,\allowbreak 9,\allowbreak 8,\allowbreak 10,\allowbreak 13)\end{array}\)\\[.5em]
\(H_{66}\) & \((5,5,5,5)\) & \(\begin{array}{l}(20,\allowbreak 1,\allowbreak 19,\allowbreak 2,\allowbreak 18)\\(15,\allowbreak 3,\allowbreak 17,\allowbreak 4,\allowbreak 14)\\(11,\allowbreak 5,\allowbreak 13,\allowbreak 6,\allowbreak 8)\\(9,\allowbreak 10,\allowbreak 7,\allowbreak 12,\allowbreak 16)\end{array}\)\\[.5em]
\(H_{67}\) & \((5,5,5,6)\) & \(\begin{array}{l}(21,\allowbreak 1,\allowbreak 20,\allowbreak 2,\allowbreak 19)\\(12,\allowbreak 6,\allowbreak 14,\allowbreak 9,\allowbreak 11)\\(7,\allowbreak 8,\allowbreak 17,\allowbreak 13,\allowbreak 10)\\(16,\allowbreak 3,\allowbreak 18,\allowbreak 4,\allowbreak 15,\allowbreak 5)\end{array}\)\\[.5em]
\(H_{68}\) & \((5,5,6,6)\) & \(\begin{array}{l}(22,\allowbreak 1,\allowbreak 21,\allowbreak 2,\allowbreak 20)\\(10,\allowbreak 14,\allowbreak 11,\allowbreak 9,\allowbreak 8)\\(17,\allowbreak 3,\allowbreak 19,\allowbreak 4,\allowbreak 16,\allowbreak 5)\\(13,\allowbreak 6,\allowbreak 15,\allowbreak 7,\allowbreak 12,\allowbreak 18)\end{array}\)\\[.5em]
\(H_{69}\) & \((5,6,6,6)\) & \(\begin{array}{l}(23,\allowbreak 1,\allowbreak 22,\allowbreak 2,\allowbreak 21)\\(18,\allowbreak 3,\allowbreak 20,\allowbreak 4,\allowbreak 17,\allowbreak 5)\\(14,\allowbreak 6,\allowbreak 16,\allowbreak 7,\allowbreak 12,\allowbreak 19)\\(11,\allowbreak 15,\allowbreak 9,\allowbreak 8,\allowbreak 10,\allowbreak 13)\end{array}\)\\[.5em]
\(H_{70}\) & \((6,6,6,6)\) & \(\begin{array}{l}(24,\allowbreak 1,\allowbreak 23,\allowbreak 2,\allowbreak 22,\allowbreak 3)\\(18,\allowbreak 4,\allowbreak 21,\allowbreak 5,\allowbreak 20,\allowbreak 7)\\(12,\allowbreak 19,\allowbreak 8,\allowbreak 17,\allowbreak 9,\allowbreak 14)\\(6,\allowbreak 16,\allowbreak 13,\allowbreak 15,\allowbreak 11,\allowbreak 10)\end{array}\)\\[.5em]
\end{longtable}
\normalsize
\ifSevenVolume
\subsection{The 311 alpha/hub-zero decompositions}\label{app:seven-finite-decompositions}
For each row, take the alpha basis $A_i$ and the hub-zero basis $H_j$ from the preceding tables. Their arm multisets combine to the displayed seven-tuple. Apply ordinary alpha amalgamation; add any unit arms to $H_j$ first. Thus each row proves an entire unbounded unit-arm family. Together with the next section these are exactly all $\binom{11}{7}=330$ sorted nonunit seven-tuples with maximum at most six.
\small
\begin{longtable}{@{}p{0.48\linewidth}@{\hspace{.02\linewidth}}p{0.15\linewidth}@{\hspace{.02\linewidth}}p{0.15\linewidth}@{\hspace{.02\linewidth}}p{0.16\linewidth}@{}}
\caption{Complete small seven-core decomposition table.}\label{tab:finite-seven}\\
\toprule
Lengths & Alpha base & Hub-zero base & Final hub\\\midrule
\endfirsthead
\toprule
Lengths & Alpha base & Hub-zero base & Final hub\\\midrule
\endhead
\midrule\multicolumn{4}{r}{\footnotesize Continued on the next page}\\\endfoot
\bottomrule\endlastfoot
\((2,2,2,2,2,2,5)\) & \(A_{1}\) & \(H_{1}\) & \(4\)\\
\((2,2,2,2,2,3,4)\) & \(A_{2}\) & \(H_{1}\) & \(4\)\\
\((2,2,2,2,2,3,5)\) & \(A_{1}\) & \(H_{2}\) & \(4\)\\
\((2,2,2,2,2,3,6)\) & \(A_{4}\) & \(H_{1}\) & \(5\)\\
\((2,2,2,2,2,4,5)\) & \(A_{1}\) & \(H_{3}\) & \(4\)\\
\((2,2,2,2,2,4,6)\) & \(A_{6}\) & \(H_{1}\) & \(6\)\\
\((2,2,2,2,2,5,5)\) & \(A_{1}\) & \(H_{4}\) & \(4\)\\
\((2,2,2,2,2,5,6)\) & \(A_{1}\) & \(H_{5}\) & \(4\)\\
\((2,2,2,2,2,6,6)\) & \(A_{9}\) & \(H_{1}\) & \(7\)\\
\((2,2,2,2,3,3,4)\) & \(A_{2}\) & \(H_{2}\) & \(4\)\\
\((2,2,2,2,3,3,5)\) & \(A_{1}\) & \(H_{6}\) & \(4\)\\
\((2,2,2,2,3,3,6)\) & \(A_{4}\) & \(H_{2}\) & \(5\)\\
\((2,2,2,2,3,4,4)\) & \(A_{2}\) & \(H_{3}\) & \(4\)\\
\((2,2,2,2,3,4,5)\) & \(A_{1}\) & \(H_{7}\) & \(4\)\\
\((2,2,2,2,3,4,6)\) & \(A_{2}\) & \(H_{5}\) & \(4\)\\
\((2,2,2,2,3,5,5)\) & \(A_{1}\) & \(H_{8}\) & \(4\)\\
\((2,2,2,2,3,5,6)\) & \(A_{1}\) & \(H_{9}\) & \(4\)\\
\((2,2,2,2,3,6,6)\) & \(A_{4}\) & \(H_{5}\) & \(5\)\\
\((2,2,2,2,4,4,5)\) & \(A_{1}\) & \(H_{10}\) & \(4\)\\
\((2,2,2,2,4,4,6)\) & \(A_{6}\) & \(H_{3}\) & \(6\)\\
\((2,2,2,2,4,5,5)\) & \(A_{1}\) & \(H_{11}\) & \(4\)\\
\((2,2,2,2,4,5,6)\) & \(A_{1}\) & \(H_{12}\) & \(4\)\\
\((2,2,2,2,4,6,6)\) & \(A_{6}\) & \(H_{5}\) & \(6\)\\
\((2,2,2,2,5,5,5)\) & \(A_{1}\) & \(H_{13}\) & \(4\)\\
\((2,2,2,2,5,5,6)\) & \(A_{1}\) & \(H_{14}\) & \(4\)\\
\((2,2,2,2,5,6,6)\) & \(A_{1}\) & \(H_{15}\) & \(4\)\\
\((2,2,2,2,6,6,6)\) & \(A_{9}\) & \(H_{5}\) & \(7\)\\
\((2,2,2,3,3,3,4)\) & \(A_{2}\) & \(H_{6}\) & \(4\)\\
\((2,2,2,3,3,3,5)\) & \(A_{1}\) & \(H_{16}\) & \(4\)\\
\((2,2,2,3,3,3,6)\) & \(A_{4}\) & \(H_{6}\) & \(5\)\\
\((2,2,2,3,3,4,4)\) & \(A_{2}\) & \(H_{7}\) & \(4\)\\
\((2,2,2,3,3,4,5)\) & \(A_{1}\) & \(H_{17}\) & \(4\)\\
\((2,2,2,3,3,4,6)\) & \(A_{2}\) & \(H_{9}\) & \(4\)\\
\((2,2,2,3,3,5,5)\) & \(A_{1}\) & \(H_{18}\) & \(4\)\\
\((2,2,2,3,3,5,6)\) & \(A_{1}\) & \(H_{19}\) & \(4\)\\
\((2,2,2,3,3,6,6)\) & \(A_{4}\) & \(H_{9}\) & \(5\)\\
\((2,2,2,3,4,4,4)\) & \(A_{2}\) & \(H_{10}\) & \(4\)\\
\((2,2,2,3,4,4,5)\) & \(A_{1}\) & \(H_{20}\) & \(4\)\\
\((2,2,2,3,4,4,6)\) & \(A_{2}\) & \(H_{12}\) & \(4\)\\
\((2,2,2,3,4,5,5)\) & \(A_{1}\) & \(H_{21}\) & \(4\)\\
\((2,2,2,3,4,5,6)\) & \(A_{1}\) & \(H_{22}\) & \(4\)\\
\((2,2,2,3,4,6,6)\) & \(A_{2}\) & \(H_{15}\) & \(4\)\\
\((2,2,2,3,5,5,5)\) & \(A_{1}\) & \(H_{23}\) & \(4\)\\
\((2,2,2,3,5,5,6)\) & \(A_{1}\) & \(H_{24}\) & \(4\)\\
\((2,2,2,3,5,6,6)\) & \(A_{1}\) & \(H_{25}\) & \(4\)\\
\((2,2,2,3,6,6,6)\) & \(A_{4}\) & \(H_{15}\) & \(5\)\\
\((2,2,2,4,4,4,5)\) & \(A_{1}\) & \(H_{26}\) & \(4\)\\
\((2,2,2,4,4,4,6)\) & \(A_{6}\) & \(H_{10}\) & \(6\)\\
\((2,2,2,4,4,5,5)\) & \(A_{1}\) & \(H_{27}\) & \(4\)\\
\((2,2,2,4,4,5,6)\) & \(A_{1}\) & \(H_{28}\) & \(4\)\\
\((2,2,2,4,4,6,6)\) & \(A_{6}\) & \(H_{12}\) & \(6\)\\
\((2,2,2,4,5,5,5)\) & \(A_{1}\) & \(H_{29}\) & \(4\)\\
\((2,2,2,4,5,5,6)\) & \(A_{1}\) & \(H_{30}\) & \(4\)\\
\((2,2,2,4,5,6,6)\) & \(A_{1}\) & \(H_{31}\) & \(4\)\\
\((2,2,2,4,6,6,6)\) & \(A_{6}\) & \(H_{15}\) & \(6\)\\
\((2,2,2,5,5,5,5)\) & \(A_{1}\) & \(H_{32}\) & \(4\)\\
\((2,2,2,5,5,5,6)\) & \(A_{1}\) & \(H_{33}\) & \(4\)\\
\((2,2,2,5,5,6,6)\) & \(A_{1}\) & \(H_{34}\) & \(4\)\\
\((2,2,2,5,6,6,6)\) & \(A_{1}\) & \(H_{35}\) & \(4\)\\
\((2,2,2,6,6,6,6)\) & \(A_{9}\) & \(H_{15}\) & \(7\)\\
\((2,2,3,3,3,3,4)\) & \(A_{2}\) & \(H_{16}\) & \(4\)\\
\((2,2,3,3,3,3,5)\) & \(A_{1}\) & \(H_{36}\) & \(4\)\\
\((2,2,3,3,3,3,6)\) & \(A_{4}\) & \(H_{16}\) & \(5\)\\
\((2,2,3,3,3,4,4)\) & \(A_{2}\) & \(H_{17}\) & \(4\)\\
\((2,2,3,3,3,4,5)\) & \(A_{1}\) & \(H_{37}\) & \(4\)\\
\((2,2,3,3,3,4,6)\) & \(A_{2}\) & \(H_{19}\) & \(4\)\\
\((2,2,3,3,3,5,5)\) & \(A_{1}\) & \(H_{38}\) & \(4\)\\
\((2,2,3,3,3,5,6)\) & \(A_{1}\) & \(H_{39}\) & \(4\)\\
\((2,2,3,3,3,6,6)\) & \(A_{4}\) & \(H_{19}\) & \(5\)\\
\((2,2,3,3,4,4,4)\) & \(A_{2}\) & \(H_{20}\) & \(4\)\\
\((2,2,3,3,4,4,5)\) & \(A_{1}\) & \(H_{40}\) & \(4\)\\
\((2,2,3,3,4,4,6)\) & \(A_{2}\) & \(H_{22}\) & \(4\)\\
\((2,2,3,3,4,5,5)\) & \(A_{1}\) & \(H_{41}\) & \(4\)\\
\((2,2,3,3,4,5,6)\) & \(A_{1}\) & \(H_{42}\) & \(4\)\\
\((2,2,3,3,4,6,6)\) & \(A_{2}\) & \(H_{25}\) & \(4\)\\
\((2,2,3,3,5,5,5)\) & \(A_{1}\) & \(H_{43}\) & \(4\)\\
\((2,2,3,3,5,5,6)\) & \(A_{1}\) & \(H_{44}\) & \(4\)\\
\((2,2,3,3,5,6,6)\) & \(A_{1}\) & \(H_{45}\) & \(4\)\\
\((2,2,3,3,6,6,6)\) & \(A_{4}\) & \(H_{25}\) & \(5\)\\
\((2,2,3,4,4,4,4)\) & \(A_{2}\) & \(H_{26}\) & \(4\)\\
\((2,2,3,4,4,4,5)\) & \(A_{1}\) & \(H_{46}\) & \(4\)\\
\((2,2,3,4,4,4,6)\) & \(A_{2}\) & \(H_{28}\) & \(4\)\\
\((2,2,3,4,4,5,5)\) & \(A_{1}\) & \(H_{47}\) & \(4\)\\
\((2,2,3,4,4,5,6)\) & \(A_{1}\) & \(H_{48}\) & \(4\)\\
\((2,2,3,4,4,6,6)\) & \(A_{2}\) & \(H_{31}\) & \(4\)\\
\((2,2,3,4,5,5,5)\) & \(A_{1}\) & \(H_{49}\) & \(4\)\\
\((2,2,3,4,5,5,6)\) & \(A_{1}\) & \(H_{50}\) & \(4\)\\
\((2,2,3,4,5,6,6)\) & \(A_{1}\) & \(H_{51}\) & \(4\)\\
\((2,2,3,4,6,6,6)\) & \(A_{2}\) & \(H_{35}\) & \(4\)\\
\((2,2,3,5,5,5,5)\) & \(A_{1}\) & \(H_{52}\) & \(4\)\\
\((2,2,3,5,5,5,6)\) & \(A_{1}\) & \(H_{53}\) & \(4\)\\
\((2,2,3,5,5,6,6)\) & \(A_{1}\) & \(H_{54}\) & \(4\)\\
\((2,2,3,5,6,6,6)\) & \(A_{1}\) & \(H_{55}\) & \(4\)\\
\((2,2,3,6,6,6,6)\) & \(A_{4}\) & \(H_{35}\) & \(5\)\\
\((2,2,4,4,4,4,5)\) & \(A_{1}\) & \(H_{56}\) & \(4\)\\
\((2,2,4,4,4,4,6)\) & \(A_{6}\) & \(H_{26}\) & \(6\)\\
\((2,2,4,4,4,5,5)\) & \(A_{1}\) & \(H_{57}\) & \(4\)\\
\((2,2,4,4,4,5,6)\) & \(A_{1}\) & \(H_{58}\) & \(4\)\\
\((2,2,4,4,4,6,6)\) & \(A_{6}\) & \(H_{28}\) & \(6\)\\
\((2,2,4,4,5,5,5)\) & \(A_{1}\) & \(H_{59}\) & \(4\)\\
\((2,2,4,4,5,5,6)\) & \(A_{1}\) & \(H_{60}\) & \(4\)\\
\((2,2,4,4,5,6,6)\) & \(A_{1}\) & \(H_{61}\) & \(4\)\\
\((2,2,4,4,6,6,6)\) & \(A_{6}\) & \(H_{31}\) & \(6\)\\
\((2,2,4,5,5,5,5)\) & \(A_{1}\) & \(H_{62}\) & \(4\)\\
\((2,2,4,5,5,5,6)\) & \(A_{1}\) & \(H_{63}\) & \(4\)\\
\((2,2,4,5,5,6,6)\) & \(A_{1}\) & \(H_{64}\) & \(4\)\\
\((2,2,4,5,6,6,6)\) & \(A_{1}\) & \(H_{65}\) & \(4\)\\
\((2,2,4,6,6,6,6)\) & \(A_{6}\) & \(H_{35}\) & \(6\)\\
\((2,2,5,5,5,5,5)\) & \(A_{1}\) & \(H_{66}\) & \(4\)\\
\((2,2,5,5,5,5,6)\) & \(A_{1}\) & \(H_{67}\) & \(4\)\\
\((2,2,5,5,5,6,6)\) & \(A_{1}\) & \(H_{68}\) & \(4\)\\
\((2,2,5,5,6,6,6)\) & \(A_{1}\) & \(H_{69}\) & \(4\)\\
\((2,2,5,6,6,6,6)\) & \(A_{1}\) & \(H_{70}\) & \(4\)\\
\((2,2,6,6,6,6,6)\) & \(A_{9}\) & \(H_{35}\) & \(7\)\\
\((2,3,3,3,3,3,4)\) & \(A_{2}\) & \(H_{36}\) & \(4\)\\
\((2,3,3,3,3,3,5)\) & \(A_{3}\) & \(H_{36}\) & \(4\)\\
\((2,3,3,3,3,3,6)\) & \(A_{4}\) & \(H_{36}\) & \(5\)\\
\((2,3,3,3,3,4,4)\) & \(A_{2}\) & \(H_{37}\) & \(4\)\\
\((2,3,3,3,3,4,5)\) & \(A_{2}\) & \(H_{38}\) & \(4\)\\
\((2,3,3,3,3,4,6)\) & \(A_{2}\) & \(H_{39}\) & \(4\)\\
\((2,3,3,3,3,5,5)\) & \(A_{3}\) & \(H_{38}\) & \(4\)\\
\((2,3,3,3,3,5,6)\) & \(A_{3}\) & \(H_{39}\) & \(4\)\\
\((2,3,3,3,3,6,6)\) & \(A_{4}\) & \(H_{39}\) & \(5\)\\
\((2,3,3,3,4,4,4)\) & \(A_{2}\) & \(H_{40}\) & \(4\)\\
\((2,3,3,3,4,4,5)\) & \(A_{2}\) & \(H_{41}\) & \(4\)\\
\((2,3,3,3,4,4,6)\) & \(A_{2}\) & \(H_{42}\) & \(4\)\\
\((2,3,3,3,4,5,5)\) & \(A_{2}\) & \(H_{43}\) & \(4\)\\
\((2,3,3,3,4,5,6)\) & \(A_{2}\) & \(H_{44}\) & \(4\)\\
\((2,3,3,3,4,6,6)\) & \(A_{2}\) & \(H_{45}\) & \(4\)\\
\((2,3,3,3,5,5,5)\) & \(A_{3}\) & \(H_{43}\) & \(4\)\\
\((2,3,3,3,5,5,6)\) & \(A_{3}\) & \(H_{44}\) & \(4\)\\
\((2,3,3,3,5,6,6)\) & \(A_{3}\) & \(H_{45}\) & \(4\)\\
\((2,3,3,3,6,6,6)\) & \(A_{4}\) & \(H_{45}\) & \(5\)\\
\((2,3,3,4,4,4,4)\) & \(A_{2}\) & \(H_{46}\) & \(4\)\\
\((2,3,3,4,4,4,5)\) & \(A_{2}\) & \(H_{47}\) & \(4\)\\
\((2,3,3,4,4,4,6)\) & \(A_{2}\) & \(H_{48}\) & \(4\)\\
\((2,3,3,4,4,5,5)\) & \(A_{2}\) & \(H_{49}\) & \(4\)\\
\((2,3,3,4,4,5,6)\) & \(A_{2}\) & \(H_{50}\) & \(4\)\\
\((2,3,3,4,4,6,6)\) & \(A_{2}\) & \(H_{51}\) & \(4\)\\
\((2,3,3,4,5,5,5)\) & \(A_{2}\) & \(H_{52}\) & \(4\)\\
\((2,3,3,4,5,5,6)\) & \(A_{2}\) & \(H_{53}\) & \(4\)\\
\((2,3,3,4,5,6,6)\) & \(A_{2}\) & \(H_{54}\) & \(4\)\\
\((2,3,3,4,6,6,6)\) & \(A_{2}\) & \(H_{55}\) & \(4\)\\
\((2,3,3,5,5,5,5)\) & \(A_{3}\) & \(H_{52}\) & \(4\)\\
\((2,3,3,5,5,5,6)\) & \(A_{3}\) & \(H_{53}\) & \(4\)\\
\((2,3,3,5,5,6,6)\) & \(A_{3}\) & \(H_{54}\) & \(4\)\\
\((2,3,3,5,6,6,6)\) & \(A_{3}\) & \(H_{55}\) & \(4\)\\
\((2,3,3,6,6,6,6)\) & \(A_{4}\) & \(H_{55}\) & \(5\)\\
\((2,3,4,4,4,4,4)\) & \(A_{2}\) & \(H_{56}\) & \(4\)\\
\((2,3,4,4,4,4,5)\) & \(A_{2}\) & \(H_{57}\) & \(4\)\\
\((2,3,4,4,4,4,6)\) & \(A_{2}\) & \(H_{58}\) & \(4\)\\
\((2,3,4,4,4,5,5)\) & \(A_{2}\) & \(H_{59}\) & \(4\)\\
\((2,3,4,4,4,5,6)\) & \(A_{2}\) & \(H_{60}\) & \(4\)\\
\((2,3,4,4,4,6,6)\) & \(A_{2}\) & \(H_{61}\) & \(4\)\\
\((2,3,4,4,5,5,5)\) & \(A_{2}\) & \(H_{62}\) & \(4\)\\
\((2,3,4,4,5,5,6)\) & \(A_{2}\) & \(H_{63}\) & \(4\)\\
\((2,3,4,4,5,6,6)\) & \(A_{2}\) & \(H_{64}\) & \(4\)\\
\((2,3,4,4,6,6,6)\) & \(A_{2}\) & \(H_{65}\) & \(4\)\\
\((2,3,4,5,5,5,5)\) & \(A_{2}\) & \(H_{66}\) & \(4\)\\
\((2,3,4,5,5,5,6)\) & \(A_{2}\) & \(H_{67}\) & \(4\)\\
\((2,3,4,5,5,6,6)\) & \(A_{2}\) & \(H_{68}\) & \(4\)\\
\((2,3,4,5,6,6,6)\) & \(A_{2}\) & \(H_{69}\) & \(4\)\\
\((2,3,4,6,6,6,6)\) & \(A_{2}\) & \(H_{70}\) & \(4\)\\
\((2,3,5,5,5,5,5)\) & \(A_{3}\) & \(H_{66}\) & \(4\)\\
\((2,3,5,5,5,5,6)\) & \(A_{3}\) & \(H_{67}\) & \(4\)\\
\((2,3,5,5,5,6,6)\) & \(A_{3}\) & \(H_{68}\) & \(4\)\\
\((2,3,5,5,6,6,6)\) & \(A_{3}\) & \(H_{69}\) & \(4\)\\
\((2,3,5,6,6,6,6)\) & \(A_{3}\) & \(H_{70}\) & \(4\)\\
\((2,3,6,6,6,6,6)\) & \(A_{4}\) & \(H_{70}\) & \(5\)\\
\((2,4,4,4,4,4,5)\) & \(A_{5}\) & \(H_{56}\) & \(5\)\\
\((2,4,4,4,4,4,6)\) & \(A_{6}\) & \(H_{56}\) & \(6\)\\
\((2,4,4,4,4,5,5)\) & \(A_{5}\) & \(H_{57}\) & \(5\)\\
\((2,4,4,4,4,5,6)\) & \(A_{5}\) & \(H_{58}\) & \(5\)\\
\((2,4,4,4,4,6,6)\) & \(A_{6}\) & \(H_{58}\) & \(6\)\\
\((2,4,4,4,5,5,5)\) & \(A_{5}\) & \(H_{59}\) & \(5\)\\
\((2,4,4,4,5,5,6)\) & \(A_{5}\) & \(H_{60}\) & \(5\)\\
\((2,4,4,4,5,6,6)\) & \(A_{5}\) & \(H_{61}\) & \(5\)\\
\((2,4,4,4,6,6,6)\) & \(A_{6}\) & \(H_{61}\) & \(6\)\\
\((2,4,4,5,5,5,5)\) & \(A_{5}\) & \(H_{62}\) & \(5\)\\
\((2,4,4,5,5,5,6)\) & \(A_{5}\) & \(H_{63}\) & \(5\)\\
\((2,4,4,5,5,6,6)\) & \(A_{5}\) & \(H_{64}\) & \(5\)\\
\((2,4,4,5,6,6,6)\) & \(A_{5}\) & \(H_{65}\) & \(5\)\\
\((2,4,4,6,6,6,6)\) & \(A_{6}\) & \(H_{65}\) & \(6\)\\
\((2,4,5,5,5,5,5)\) & \(A_{5}\) & \(H_{66}\) & \(5\)\\
\((2,4,5,5,5,5,6)\) & \(A_{5}\) & \(H_{67}\) & \(5\)\\
\((2,4,5,5,5,6,6)\) & \(A_{5}\) & \(H_{68}\) & \(5\)\\
\((2,4,5,5,6,6,6)\) & \(A_{5}\) & \(H_{69}\) & \(5\)\\
\((2,4,5,6,6,6,6)\) & \(A_{5}\) & \(H_{70}\) & \(5\)\\
\((2,4,6,6,6,6,6)\) & \(A_{6}\) & \(H_{70}\) & \(6\)\\
\((2,5,5,5,5,5,5)\) & \(A_{7}\) & \(H_{66}\) & \(5\)\\
\((2,5,5,5,5,5,6)\) & \(A_{7}\) & \(H_{67}\) & \(5\)\\
\((2,5,5,5,5,6,6)\) & \(A_{7}\) & \(H_{68}\) & \(5\)\\
\((2,5,5,5,6,6,6)\) & \(A_{7}\) & \(H_{69}\) & \(5\)\\
\((2,5,5,6,6,6,6)\) & \(A_{7}\) & \(H_{70}\) & \(5\)\\
\((2,5,6,6,6,6,6)\) & \(A_{8}\) & \(H_{70}\) & \(6\)\\
\((2,6,6,6,6,6,6)\) & \(A_{9}\) & \(H_{70}\) & \(7\)\\
\((3,3,3,3,3,3,4)\) & \(A_{10}\) & \(H_{36}\) & \(4\)\\
\((3,3,3,3,3,3,5)\) & \(A_{11}\) & \(H_{36}\) & \(4\)\\
\((3,3,3,3,3,3,6)\) & \(A_{12}\) & \(H_{36}\) & \(5\)\\
\((3,3,3,3,3,4,4)\) & \(A_{10}\) & \(H_{37}\) & \(4\)\\
\((3,3,3,3,3,4,5)\) & \(A_{10}\) & \(H_{38}\) & \(4\)\\
\((3,3,3,3,3,4,6)\) & \(A_{10}\) & \(H_{39}\) & \(4\)\\
\((3,3,3,3,3,5,5)\) & \(A_{11}\) & \(H_{38}\) & \(4\)\\
\((3,3,3,3,3,5,6)\) & \(A_{11}\) & \(H_{39}\) & \(4\)\\
\((3,3,3,3,3,6,6)\) & \(A_{12}\) & \(H_{39}\) & \(5\)\\
\((3,3,3,3,4,4,4)\) & \(A_{10}\) & \(H_{40}\) & \(4\)\\
\((3,3,3,3,4,4,5)\) & \(A_{10}\) & \(H_{41}\) & \(4\)\\
\((3,3,3,3,4,4,6)\) & \(A_{10}\) & \(H_{42}\) & \(4\)\\
\((3,3,3,3,4,5,5)\) & \(A_{10}\) & \(H_{43}\) & \(4\)\\
\((3,3,3,3,4,5,6)\) & \(A_{10}\) & \(H_{44}\) & \(4\)\\
\((3,3,3,3,4,6,6)\) & \(A_{10}\) & \(H_{45}\) & \(4\)\\
\((3,3,3,3,5,5,5)\) & \(A_{11}\) & \(H_{43}\) & \(4\)\\
\((3,3,3,3,5,5,6)\) & \(A_{11}\) & \(H_{44}\) & \(4\)\\
\((3,3,3,3,5,6,6)\) & \(A_{11}\) & \(H_{45}\) & \(4\)\\
\((3,3,3,3,6,6,6)\) & \(A_{12}\) & \(H_{45}\) & \(5\)\\
\((3,3,3,4,4,4,4)\) & \(A_{10}\) & \(H_{46}\) & \(4\)\\
\((3,3,3,4,4,4,5)\) & \(A_{10}\) & \(H_{47}\) & \(4\)\\
\((3,3,3,4,4,4,6)\) & \(A_{10}\) & \(H_{48}\) & \(4\)\\
\((3,3,3,4,4,5,5)\) & \(A_{10}\) & \(H_{49}\) & \(4\)\\
\((3,3,3,4,4,5,6)\) & \(A_{10}\) & \(H_{50}\) & \(4\)\\
\((3,3,3,4,4,6,6)\) & \(A_{10}\) & \(H_{51}\) & \(4\)\\
\((3,3,3,4,5,5,5)\) & \(A_{10}\) & \(H_{52}\) & \(4\)\\
\((3,3,3,4,5,5,6)\) & \(A_{10}\) & \(H_{53}\) & \(4\)\\
\((3,3,3,4,5,6,6)\) & \(A_{10}\) & \(H_{54}\) & \(4\)\\
\((3,3,3,4,6,6,6)\) & \(A_{10}\) & \(H_{55}\) & \(4\)\\
\((3,3,3,5,5,5,5)\) & \(A_{11}\) & \(H_{52}\) & \(4\)\\
\((3,3,3,5,5,5,6)\) & \(A_{11}\) & \(H_{53}\) & \(4\)\\
\((3,3,3,5,5,6,6)\) & \(A_{11}\) & \(H_{54}\) & \(4\)\\
\((3,3,3,5,6,6,6)\) & \(A_{11}\) & \(H_{55}\) & \(4\)\\
\((3,3,3,6,6,6,6)\) & \(A_{12}\) & \(H_{55}\) & \(5\)\\
\((3,3,4,4,4,4,4)\) & \(A_{10}\) & \(H_{56}\) & \(4\)\\
\((3,3,4,4,4,4,5)\) & \(A_{10}\) & \(H_{57}\) & \(4\)\\
\((3,3,4,4,4,4,6)\) & \(A_{10}\) & \(H_{58}\) & \(4\)\\
\((3,3,4,4,4,5,5)\) & \(A_{10}\) & \(H_{59}\) & \(4\)\\
\((3,3,4,4,4,5,6)\) & \(A_{10}\) & \(H_{60}\) & \(4\)\\
\((3,3,4,4,4,6,6)\) & \(A_{10}\) & \(H_{61}\) & \(4\)\\
\((3,3,4,4,5,5,5)\) & \(A_{10}\) & \(H_{62}\) & \(4\)\\
\((3,3,4,4,5,5,6)\) & \(A_{10}\) & \(H_{63}\) & \(4\)\\
\((3,3,4,4,5,6,6)\) & \(A_{10}\) & \(H_{64}\) & \(4\)\\
\((3,3,4,4,6,6,6)\) & \(A_{10}\) & \(H_{65}\) & \(4\)\\
\((3,3,4,5,5,5,5)\) & \(A_{10}\) & \(H_{66}\) & \(4\)\\
\((3,3,4,5,5,5,6)\) & \(A_{10}\) & \(H_{67}\) & \(4\)\\
\((3,3,4,5,5,6,6)\) & \(A_{10}\) & \(H_{68}\) & \(4\)\\
\((3,3,4,5,6,6,6)\) & \(A_{10}\) & \(H_{69}\) & \(4\)\\
\((3,3,4,6,6,6,6)\) & \(A_{10}\) & \(H_{70}\) & \(4\)\\
\((3,3,5,5,5,5,5)\) & \(A_{11}\) & \(H_{66}\) & \(4\)\\
\((3,3,5,5,5,5,6)\) & \(A_{11}\) & \(H_{67}\) & \(4\)\\
\((3,3,5,5,5,6,6)\) & \(A_{11}\) & \(H_{68}\) & \(4\)\\
\((3,3,5,5,6,6,6)\) & \(A_{11}\) & \(H_{69}\) & \(4\)\\
\((3,3,5,6,6,6,6)\) & \(A_{11}\) & \(H_{70}\) & \(4\)\\
\((3,3,6,6,6,6,6)\) & \(A_{12}\) & \(H_{70}\) & \(5\)\\
\((3,4,4,4,4,4,5)\) & \(A_{13}\) & \(H_{56}\) & \(5\)\\
\((3,4,4,4,4,4,6)\) & \(A_{14}\) & \(H_{56}\) & \(6\)\\
\((3,4,4,4,4,5,5)\) & \(A_{13}\) & \(H_{57}\) & \(5\)\\
\((3,4,4,4,4,5,6)\) & \(A_{13}\) & \(H_{58}\) & \(5\)\\
\((3,4,4,4,4,6,6)\) & \(A_{14}\) & \(H_{58}\) & \(6\)\\
\((3,4,4,4,5,5,5)\) & \(A_{13}\) & \(H_{59}\) & \(5\)\\
\((3,4,4,4,5,5,6)\) & \(A_{13}\) & \(H_{60}\) & \(5\)\\
\((3,4,4,4,5,6,6)\) & \(A_{13}\) & \(H_{61}\) & \(5\)\\
\((3,4,4,4,6,6,6)\) & \(A_{14}\) & \(H_{61}\) & \(6\)\\
\((3,4,4,5,5,5,5)\) & \(A_{13}\) & \(H_{62}\) & \(5\)\\
\((3,4,4,5,5,5,6)\) & \(A_{13}\) & \(H_{63}\) & \(5\)\\
\((3,4,4,5,5,6,6)\) & \(A_{13}\) & \(H_{64}\) & \(5\)\\
\((3,4,4,5,6,6,6)\) & \(A_{13}\) & \(H_{65}\) & \(5\)\\
\((3,4,4,6,6,6,6)\) & \(A_{14}\) & \(H_{65}\) & \(6\)\\
\((3,4,5,5,5,5,5)\) & \(A_{13}\) & \(H_{66}\) & \(5\)\\
\((3,4,5,5,5,5,6)\) & \(A_{13}\) & \(H_{67}\) & \(5\)\\
\((3,4,5,5,5,6,6)\) & \(A_{13}\) & \(H_{68}\) & \(5\)\\
\((3,4,5,5,6,6,6)\) & \(A_{13}\) & \(H_{69}\) & \(5\)\\
\((3,4,5,6,6,6,6)\) & \(A_{13}\) & \(H_{70}\) & \(5\)\\
\((3,4,6,6,6,6,6)\) & \(A_{14}\) & \(H_{70}\) & \(6\)\\
\((3,5,5,5,5,5,5)\) & \(A_{15}\) & \(H_{66}\) & \(5\)\\
\((3,5,5,5,5,5,6)\) & \(A_{15}\) & \(H_{67}\) & \(5\)\\
\((3,5,5,5,5,6,6)\) & \(A_{15}\) & \(H_{68}\) & \(5\)\\
\((3,5,5,5,6,6,6)\) & \(A_{15}\) & \(H_{69}\) & \(5\)\\
\((3,5,5,6,6,6,6)\) & \(A_{15}\) & \(H_{70}\) & \(5\)\\
\((3,5,6,6,6,6,6)\) & \(A_{16}\) & \(H_{70}\) & \(6\)\\
\((3,6,6,6,6,6,6)\) & \(A_{17}\) & \(H_{70}\) & \(7\)\\
\((4,4,4,4,4,5,5)\) & \(A_{18}\) & \(H_{56}\) & \(6\)\\
\((4,4,4,4,4,5,6)\) & \(A_{19}\) & \(H_{56}\) & \(7\)\\
\((4,4,4,4,4,6,6)\) & \(A_{20}\) & \(H_{56}\) & \(8\)\\
\((4,4,4,4,5,5,5)\) & \(A_{18}\) & \(H_{57}\) & \(6\)\\
\((4,4,4,4,5,5,6)\) & \(A_{18}\) & \(H_{58}\) & \(6\)\\
\((4,4,4,4,5,6,6)\) & \(A_{19}\) & \(H_{58}\) & \(7\)\\
\((4,4,4,4,6,6,6)\) & \(A_{20}\) & \(H_{58}\) & \(8\)\\
\((4,4,4,5,5,5,5)\) & \(A_{18}\) & \(H_{59}\) & \(6\)\\
\((4,4,4,5,5,5,6)\) & \(A_{18}\) & \(H_{60}\) & \(6\)\\
\((4,4,4,5,5,6,6)\) & \(A_{18}\) & \(H_{61}\) & \(6\)\\
\((4,4,4,5,6,6,6)\) & \(A_{19}\) & \(H_{61}\) & \(7\)\\
\((4,4,4,6,6,6,6)\) & \(A_{20}\) & \(H_{61}\) & \(8\)\\
\((4,4,5,5,5,5,5)\) & \(A_{18}\) & \(H_{62}\) & \(6\)\\
\((4,4,5,5,5,5,6)\) & \(A_{18}\) & \(H_{63}\) & \(6\)\\
\((4,4,5,5,5,6,6)\) & \(A_{18}\) & \(H_{64}\) & \(6\)\\
\((4,4,5,5,6,6,6)\) & \(A_{18}\) & \(H_{65}\) & \(6\)\\
\((4,4,5,6,6,6,6)\) & \(A_{19}\) & \(H_{65}\) & \(7\)\\
\((4,4,6,6,6,6,6)\) & \(A_{20}\) & \(H_{65}\) & \(8\)\\
\((4,5,5,5,5,5,5)\) & \(A_{18}\) & \(H_{66}\) & \(6\)\\
\((4,5,5,5,5,5,6)\) & \(A_{18}\) & \(H_{67}\) & \(6\)\\
\((4,5,5,5,5,6,6)\) & \(A_{18}\) & \(H_{68}\) & \(6\)\\
\((4,5,5,5,6,6,6)\) & \(A_{18}\) & \(H_{69}\) & \(6\)\\
\((4,5,5,6,6,6,6)\) & \(A_{18}\) & \(H_{70}\) & \(6\)\\
\((4,5,6,6,6,6,6)\) & \(A_{19}\) & \(H_{70}\) & \(7\)\\
\((4,6,6,6,6,6,6)\) & \(A_{20}\) & \(H_{70}\) & \(8\)\\
\((5,5,5,5,5,5,5)\) & \(A_{21}\) & \(H_{66}\) & \(6\)\\
\((5,5,5,5,5,5,6)\) & \(A_{21}\) & \(H_{67}\) & \(6\)\\
\((5,5,5,5,5,6,6)\) & \(A_{21}\) & \(H_{68}\) & \(6\)\\
\((5,5,5,5,6,6,6)\) & \(A_{21}\) & \(H_{69}\) & \(6\)\\
\((5,5,5,6,6,6,6)\) & \(A_{21}\) & \(H_{70}\) & \(6\)\\
\((5,5,6,6,6,6,6)\) & \(A_{22}\) & \(H_{70}\) & \(7\)\\
\((5,6,6,6,6,6,6)\) & \(A_{23}\) & \(H_{70}\) & \(8\)\\
\((6,6,6,6,6,6,6)\) & \(A_{24}\) & \(H_{70}\) & \(9\)\\
\end{longtable}
\normalsize
\subsection{The 19 direct hub-zero cores}\label{app:seven-direct-finite}
These remaining profiles are not negative alpha assertions. The direct hub-zero words displayed here suffice, and arbitrary unit leaves receive successive new maximum labels.
\small
\begin{longtable}{@{}p{0.3\linewidth}@{\hspace{.02\linewidth}}p{0.68\linewidth}@{}}
\caption{Direct hub-zero seven-core certificates.}\label{tab:finite-direct-seven}\\
\toprule
Lengths & Arm words\\\midrule
\endfirsthead
\toprule
Lengths & Arm words\\\midrule
\endhead
\midrule\multicolumn{2}{r}{\footnotesize Continued on the next page}\\\endfoot
\bottomrule\endlastfoot
\((2,2,2,2,2,2,2)\) & \(\begin{array}{l}(14,\allowbreak 1)\\(12,\allowbreak 2)\\(11,\allowbreak 3)\\(9,\allowbreak 13)\\(7,\allowbreak 10)\\(6,\allowbreak 8)\\(5,\allowbreak 4)\end{array}\)\\[.5em]
\((2,2,2,2,2,2,3)\) & \(\begin{array}{l}(15,\allowbreak 1)\\(10,\allowbreak 3)\\(9,\allowbreak 12)\\(8,\allowbreak 7)\\(5,\allowbreak 11)\\(4,\allowbreak 6)\\(13,\allowbreak 2,\allowbreak 14)\end{array}\)\\[.5em]
\((2,2,2,2,2,2,4)\) & \(\begin{array}{l}(16,\allowbreak 1)\\(10,\allowbreak 12)\\(9,\allowbreak 8)\\(7,\allowbreak 11)\\(6,\allowbreak 3)\\(5,\allowbreak 13)\\(14,\allowbreak 2,\allowbreak 15,\allowbreak 4)\end{array}\)\\[.5em]
\((2,2,2,2,2,2,6)\) & \(\begin{array}{l}(18,\allowbreak 1)\\(12,\allowbreak 14)\\(9,\allowbreak 13)\\(8,\allowbreak 7)\\(6,\allowbreak 11)\\(3,\allowbreak 10)\\(16,\allowbreak 2,\allowbreak 17,\allowbreak 4,\allowbreak 15,\allowbreak 5)\end{array}\)\\[.5em]
\((2,2,2,2,2,3,3)\) & \(\begin{array}{l}(16,\allowbreak 1)\\(9,\allowbreak 10)\\(6,\allowbreak 8)\\(5,\allowbreak 12)\\(4,\allowbreak 7)\\(14,\allowbreak 2,\allowbreak 15)\\(11,\allowbreak 3,\allowbreak 13)\end{array}\)\\[.5em]
\((2,2,2,2,2,4,4)\) & \(\begin{array}{l}(18,\allowbreak 1)\\(10,\allowbreak 8)\\(7,\allowbreak 11)\\(6,\allowbreak 9)\\(5,\allowbreak 13)\\(16,\allowbreak 2,\allowbreak 17,\allowbreak 4)\\(12,\allowbreak 3,\allowbreak 14,\allowbreak 15)\end{array}\)\\[.5em]
\((2,2,2,2,3,3,3)\) & \(\begin{array}{l}(17,\allowbreak 1)\\(8,\allowbreak 13)\\(6,\allowbreak 5)\\(4,\allowbreak 11)\\(15,\allowbreak 2,\allowbreak 16)\\(12,\allowbreak 3,\allowbreak 14)\\(10,\allowbreak 7,\allowbreak 9)\end{array}\)\\[.5em]
\((2,2,2,2,4,4,4)\) & \(\begin{array}{l}(20,\allowbreak 1)\\(9,\allowbreak 10)\\(8,\allowbreak 5)\\(7,\allowbreak 13)\\(18,\allowbreak 2,\allowbreak 19,\allowbreak 4)\\(14,\allowbreak 3,\allowbreak 16,\allowbreak 6)\\(12,\allowbreak 17,\allowbreak 15,\allowbreak 11)\end{array}\)\\[.5em]
\((2,2,2,3,3,3,3)\) & \(\begin{array}{l}(18,\allowbreak 1)\\(9,\allowbreak 14)\\(6,\allowbreak 5)\\(16,\allowbreak 2,\allowbreak 17)\\(13,\allowbreak 3,\allowbreak 15)\\(11,\allowbreak 4,\allowbreak 7)\\(8,\allowbreak 12,\allowbreak 10)\end{array}\)\\[.5em]
\((2,2,2,4,4,4,4)\) & \(\begin{array}{l}(22,\allowbreak 1)\\(11,\allowbreak 12)\\(7,\allowbreak 10)\\(20,\allowbreak 2,\allowbreak 21,\allowbreak 4)\\(16,\allowbreak 3,\allowbreak 18,\allowbreak 6)\\(14,\allowbreak 5,\allowbreak 15,\allowbreak 9)\\(8,\allowbreak 13,\allowbreak 17,\allowbreak 19)\end{array}\)\\[.5em]
\((2,2,3,3,3,3,3)\) & \(\begin{array}{l}(19,\allowbreak 1)\\(7,\allowbreak 5)\\(17,\allowbreak 2,\allowbreak 18)\\(14,\allowbreak 3,\allowbreak 16)\\(12,\allowbreak 4,\allowbreak 13)\\(10,\allowbreak 15,\allowbreak 11)\\(6,\allowbreak 9,\allowbreak 8)\end{array}\)\\[.5em]
\((2,2,4,4,4,4,4)\) & \(\begin{array}{l}(24,\allowbreak 1)\\(8,\allowbreak 10)\\(22,\allowbreak 2,\allowbreak 23,\allowbreak 4)\\(18,\allowbreak 3,\allowbreak 20,\allowbreak 6)\\(16,\allowbreak 5,\allowbreak 17,\allowbreak 7)\\(13,\allowbreak 19,\allowbreak 14,\allowbreak 21)\\(9,\allowbreak 12,\allowbreak 11,\allowbreak 15)\end{array}\)\\[.5em]
\((2,3,3,3,3,3,3)\) & \(\begin{array}{l}(20,\allowbreak 1)\\(18,\allowbreak 2,\allowbreak 19)\\(15,\allowbreak 3,\allowbreak 17)\\(13,\allowbreak 4,\allowbreak 14)\\(11,\allowbreak 16,\allowbreak 8)\\(7,\allowbreak 10,\allowbreak 12)\\(6,\allowbreak 5,\allowbreak 9)\end{array}\)\\[.5em]
\((2,4,4,4,4,4,4)\) & \(\begin{array}{l}(26,\allowbreak 1)\\(24,\allowbreak 2,\allowbreak 25,\allowbreak 4)\\(20,\allowbreak 3,\allowbreak 22,\allowbreak 6)\\(18,\allowbreak 5,\allowbreak 19,\allowbreak 7)\\(15,\allowbreak 23,\allowbreak 12,\allowbreak 14)\\(10,\allowbreak 11,\allowbreak 17,\allowbreak 21)\\(9,\allowbreak 16,\allowbreak 13,\allowbreak 8)\end{array}\)\\[.5em]
\((3,3,3,3,3,3,3)\) & \(\begin{array}{l}(21,\allowbreak 1,\allowbreak 20)\\(18,\allowbreak 2,\allowbreak 19)\\(15,\allowbreak 3,\allowbreak 17)\\(13,\allowbreak 4,\allowbreak 14)\\(11,\allowbreak 16,\allowbreak 8)\\(7,\allowbreak 10,\allowbreak 12)\\(6,\allowbreak 5,\allowbreak 9)\end{array}\)\\[.5em]
\((3,4,4,4,4,4,4)\) & \(\begin{array}{l}(27,\allowbreak 1,\allowbreak 26)\\(24,\allowbreak 2,\allowbreak 25,\allowbreak 4)\\(20,\allowbreak 3,\allowbreak 22,\allowbreak 6)\\(18,\allowbreak 5,\allowbreak 19,\allowbreak 7)\\(15,\allowbreak 23,\allowbreak 12,\allowbreak 14)\\(10,\allowbreak 11,\allowbreak 17,\allowbreak 21)\\(9,\allowbreak 16,\allowbreak 13,\allowbreak 8)\end{array}\)\\[.5em]
\((4,4,4,4,4,4,4)\) & \(\begin{array}{l}(28,\allowbreak 1,\allowbreak 27,\allowbreak 2)\\(24,\allowbreak 3,\allowbreak 26,\allowbreak 4)\\(20,\allowbreak 5,\allowbreak 23,\allowbreak 6)\\(19,\allowbreak 7,\allowbreak 21,\allowbreak 8)\\(16,\allowbreak 25,\allowbreak 17,\allowbreak 22)\\(11,\allowbreak 18,\allowbreak 12,\allowbreak 9)\\(10,\allowbreak 14,\allowbreak 15,\allowbreak 13)\end{array}\)\\[.5em]
\((4,4,4,4,4,4,5)\) & \(\begin{array}{l}(29,\allowbreak 1,\allowbreak 28,\allowbreak 2)\\(21,\allowbreak 5,\allowbreak 23,\allowbreak 6)\\(19,\allowbreak 7,\allowbreak 22,\allowbreak 8)\\(13,\allowbreak 20,\allowbreak 10,\allowbreak 16)\\(11,\allowbreak 12,\allowbreak 15,\allowbreak 17)\\(9,\allowbreak 14,\allowbreak 18,\allowbreak 26)\\(25,\allowbreak 3,\allowbreak 27,\allowbreak 4,\allowbreak 24)\end{array}\)\\[.5em]
\((4,4,4,4,4,4,6)\) & \(\begin{array}{l}(30,\allowbreak 1,\allowbreak 29,\allowbreak 2)\\(22,\allowbreak 6,\allowbreak 24,\allowbreak 7)\\(19,\allowbreak 8,\allowbreak 23,\allowbreak 9)\\(13,\allowbreak 21,\allowbreak 14,\allowbreak 17)\\(12,\allowbreak 11,\allowbreak 15,\allowbreak 20)\\(10,\allowbreak 16,\allowbreak 18,\allowbreak 27)\\(26,\allowbreak 3,\allowbreak 28,\allowbreak 4,\allowbreak 25,\allowbreak 5)\end{array}\)\\[.5em]
\end{longtable}
\normalsize
\fi
\subsection{Exact verification}
The tables are regenerated by \file{proof/finite/build_finite_tables.py}. It checks every printed base, checks the exact finite profile sets, inverts all 311 decompositions, and then independently rebuilds every seven-core labeling using the particular canonical bases printed above. The source data are \file{proof/finite/data/seven_closure_finite_base_792.jsonl} and \file{proof/finite/data/seven_closure_hubzero4_all_m6.jsonl}. Only the exact 330 nonunit profiles of the first file are used as proof inputs here.

\section{Auxiliary finite path and spider certificates}
\label{app:auxiliary-finite}
All spider words below are read outward from hub zero. The word lengths, complete vertex and color inventories, alpha separation, and the stated controlled interfaces have been checked by the integer-only procedure of Section~\ref{templates:finite-check}. A displayed alpha cut is the sum of the half-lengths. For an offset record the nonroot interval is $[k+1,k+N]$, not $[1,N]$. The builder \file{proof/finite/build_auxiliary_tables.py} rechecks every word and the equality of every catalog with its required finite domain.
\subsection{The 55 minimum-two finite cores}
\label{aux:min2}
\begingroup\footnotesize
\begin{longtable}{@{}p{25mm}p{115mm}@{}}
Profile / interface & Complete outward words\\\hline\endhead
{\raggedright\arraybackslash {\ttfamily\footnotesize (2,\allowbreak 3,\allowbreak 9)}\par $K=6$} & {\ttfamily\footnotesize (14,\allowbreak 1)\par\smallskip (12,\allowbreak 2,\allowbreak 10,\allowbreak 3,\allowbreak 8,\allowbreak 6,\allowbreak 7,\allowbreak 4,\allowbreak 13)\par\smallskip (11,\allowbreak 5,\allowbreak 9)}\tabularnewline[5pt]
{\raggedright\arraybackslash {\ttfamily\footnotesize (2,\allowbreak 4,\allowbreak 13)}\par $K=9$} & {\ttfamily\footnotesize (19,\allowbreak 1)\par\smallskip (17,\allowbreak 2,\allowbreak 18,\allowbreak 4,\allowbreak 16,\allowbreak 5,\allowbreak 15,\allowbreak 7,\allowbreak 11,\allowbreak 8,\allowbreak 10,\allowbreak 9,\allowbreak 14)\par\smallskip (13,\allowbreak 6,\allowbreak 12,\allowbreak 3)}\tabularnewline[5pt]
{\raggedright\arraybackslash {\ttfamily\footnotesize (2,\allowbreak 5,\allowbreak 5)}\par $K=5$\par short leaf $1$} & {\ttfamily\footnotesize (12,\allowbreak 1)\par\smallskip (10,\allowbreak 4,\allowbreak 6,\allowbreak 5,\allowbreak 8)\par\smallskip (9,\allowbreak 2,\allowbreak 7,\allowbreak 3,\allowbreak 11)}\tabularnewline[5pt]
{\raggedright\arraybackslash {\ttfamily\footnotesize (2,\allowbreak 5,\allowbreak 6)}\par $K=6$} & {\ttfamily\footnotesize (13,\allowbreak 1)\par\smallskip (11,\allowbreak 4,\allowbreak 9,\allowbreak 3,\allowbreak 12,\allowbreak 2)\par\smallskip (8,\allowbreak 5,\allowbreak 7,\allowbreak 6,\allowbreak 10)}\tabularnewline[5pt]
{\raggedright\arraybackslash {\ttfamily\footnotesize (2,\allowbreak 6,\allowbreak 7)}\par $K=7$\par $2e\le b+4$} & {\ttfamily\footnotesize (15,\allowbreak 1)\par\smallskip (13,\allowbreak 2,\allowbreak 14,\allowbreak 4,\allowbreak 10,\allowbreak 5,\allowbreak 12)\par\smallskip (9,\allowbreak 6,\allowbreak 8,\allowbreak 7,\allowbreak 11,\allowbreak 3)}\tabularnewline[5pt]
{\raggedright\arraybackslash {\ttfamily\footnotesize (2,\allowbreak 6,\allowbreak 8)}\par $K=8$} & {\ttfamily\footnotesize (16,\allowbreak 1)\par\smallskip (14,\allowbreak 2,\allowbreak 15,\allowbreak 5,\allowbreak 13,\allowbreak 6,\allowbreak 12,\allowbreak 3)\par\smallskip (11,\allowbreak 7,\allowbreak 10,\allowbreak 8,\allowbreak 9,\allowbreak 4)}\tabularnewline[5pt]
{\raggedright\arraybackslash {\ttfamily\footnotesize (2,\allowbreak 7,\allowbreak 7)}\par $K=7$\par short leaf $1$} & {\ttfamily\footnotesize (16,\allowbreak 1)\par\smallskip (14,\allowbreak 2,\allowbreak 15,\allowbreak 4,\allowbreak 11,\allowbreak 3,\allowbreak 13)\par\smallskip (9,\allowbreak 5,\allowbreak 10,\allowbreak 7,\allowbreak 8,\allowbreak 6,\allowbreak 12)}\tabularnewline[5pt]
{\raggedright\arraybackslash {\ttfamily\footnotesize (2,\allowbreak 7,\allowbreak 8)}\par $K=8$} & {\ttfamily\footnotesize (17,\allowbreak 1)\par\smallskip (15,\allowbreak 2,\allowbreak 16,\allowbreak 4,\allowbreak 12,\allowbreak 5,\allowbreak 14,\allowbreak 3)\par\smallskip (10,\allowbreak 6,\allowbreak 11,\allowbreak 8,\allowbreak 9,\allowbreak 7,\allowbreak 13)}\tabularnewline[5pt]
{\raggedright\arraybackslash {\ttfamily\footnotesize (2,\allowbreak 8,\allowbreak 9)}\par $K=9$\par $2e\le b+4$} & {\ttfamily\footnotesize (19,\allowbreak 1)\par\smallskip (17,\allowbreak 2,\allowbreak 18,\allowbreak 4,\allowbreak 16,\allowbreak 5,\allowbreak 15,\allowbreak 6,\allowbreak 14)\par\smallskip (13,\allowbreak 7,\allowbreak 12,\allowbreak 8,\allowbreak 11,\allowbreak 9,\allowbreak 10,\allowbreak 3)}\tabularnewline[5pt]
{\raggedright\arraybackslash {\ttfamily\footnotesize (2,\allowbreak 8,\allowbreak 10)}\par $K=10$} & {\ttfamily\footnotesize (20,\allowbreak 1)\par\smallskip (18,\allowbreak 2,\allowbreak 19,\allowbreak 4,\allowbreak 17,\allowbreak 5,\allowbreak 16,\allowbreak 8,\allowbreak 15,\allowbreak 6)\par\smallskip (14,\allowbreak 9,\allowbreak 12,\allowbreak 10,\allowbreak 11,\allowbreak 7,\allowbreak 13,\allowbreak 3)}\tabularnewline[5pt]
{\raggedright\arraybackslash {\ttfamily\footnotesize (2,\allowbreak 9,\allowbreak 9)}\par $K=9$\par short leaf $1$} & {\ttfamily\footnotesize (17,\allowbreak 6,\allowbreak 12,\allowbreak 8,\allowbreak 11,\allowbreak 9,\allowbreak 10,\allowbreak 5,\allowbreak 13)\par\smallskip (19,\allowbreak 1)\par\smallskip (20,\allowbreak 4,\allowbreak 14,\allowbreak 7,\allowbreak 16,\allowbreak 2,\allowbreak 15,\allowbreak 3,\allowbreak 18)}\tabularnewline[5pt]
{\raggedright\arraybackslash {\ttfamily\footnotesize (2,\allowbreak 9,\allowbreak 10)}\par $K=10$} & {\ttfamily\footnotesize (21,\allowbreak 1)\par\smallskip (19,\allowbreak 2,\allowbreak 20,\allowbreak 4,\allowbreak 18,\allowbreak 3,\allowbreak 16,\allowbreak 7,\allowbreak 15,\allowbreak 5)\par\smallskip (12,\allowbreak 8,\allowbreak 14,\allowbreak 9,\allowbreak 11,\allowbreak 10,\allowbreak 13,\allowbreak 6,\allowbreak 17)}\tabularnewline[5pt]
{\raggedright\arraybackslash {\ttfamily\footnotesize (2,\allowbreak 10,\allowbreak 11)}\par $K=11$\par $2e\le b+4$} & {\ttfamily\footnotesize (19,\allowbreak 3,\allowbreak 15,\allowbreak 10,\allowbreak 12,\allowbreak 11,\allowbreak 14,\allowbreak 8,\allowbreak 16,\allowbreak 7,\allowbreak 17)\par\smallskip (21,\allowbreak 6,\allowbreak 13,\allowbreak 9,\allowbreak 20,\allowbreak 2,\allowbreak 22,\allowbreak 5,\allowbreak 18,\allowbreak 4)\par\smallskip (23,\allowbreak 1)}\tabularnewline[5pt]
{\raggedright\arraybackslash {\ttfamily\footnotesize (2,\allowbreak 10,\allowbreak 12)}\par $K=12$} & {\ttfamily\footnotesize (24,\allowbreak 1)\par\smallskip (22,\allowbreak 2,\allowbreak 23,\allowbreak 4,\allowbreak 21,\allowbreak 3,\allowbreak 19,\allowbreak 7,\allowbreak 16,\allowbreak 8,\allowbreak 18,\allowbreak 5)\par\smallskip (15,\allowbreak 10,\allowbreak 17,\allowbreak 11,\allowbreak 14,\allowbreak 12,\allowbreak 13,\allowbreak 9,\allowbreak 20,\allowbreak 6)}\tabularnewline[5pt]
{\raggedright\arraybackslash {\ttfamily\footnotesize (2,\allowbreak 11,\allowbreak 11)}\par $K=11$\par short leaf $1$} & {\ttfamily\footnotesize (21,\allowbreak 1)\par\smallskip (23,\allowbreak 5,\allowbreak 16,\allowbreak 7,\allowbreak 12,\allowbreak 11,\allowbreak 13,\allowbreak 10,\allowbreak 14,\allowbreak 8,\allowbreak 15)\par\smallskip (24,\allowbreak 2,\allowbreak 19,\allowbreak 9,\allowbreak 17,\allowbreak 4,\allowbreak 20,\allowbreak 6,\allowbreak 18,\allowbreak 3,\allowbreak 22)}\tabularnewline[5pt]
{\raggedright\arraybackslash {\ttfamily\footnotesize (2,\allowbreak 11,\allowbreak 12)}\par $K=12$} & {\ttfamily\footnotesize (25,\allowbreak 1)\par\smallskip (23,\allowbreak 2,\allowbreak 24,\allowbreak 4,\allowbreak 22,\allowbreak 3,\allowbreak 20,\allowbreak 5,\allowbreak 19,\allowbreak 6,\allowbreak 18,\allowbreak 7)\par\smallskip (16,\allowbreak 8,\allowbreak 17,\allowbreak 10,\allowbreak 15,\allowbreak 9,\allowbreak 13,\allowbreak 12,\allowbreak 14,\allowbreak 11,\allowbreak 21)}\tabularnewline[5pt]
{\raggedright\arraybackslash {\ttfamily\footnotesize (2,\allowbreak 12,\allowbreak 13)}\par $K=13$\par $2e\le b+4$} & {\ttfamily\footnotesize (27,\allowbreak 1)\par\smallskip (25,\allowbreak 2,\allowbreak 26,\allowbreak 4,\allowbreak 24,\allowbreak 3,\allowbreak 22,\allowbreak 5,\allowbreak 23,\allowbreak 7,\allowbreak 19,\allowbreak 8,\allowbreak 21)\par\smallskip (15,\allowbreak 11,\allowbreak 18,\allowbreak 9,\allowbreak 17,\allowbreak 12,\allowbreak 14,\allowbreak 13,\allowbreak 16,\allowbreak 10,\allowbreak 20,\allowbreak 6)}\tabularnewline[5pt]
{\raggedright\arraybackslash {\ttfamily\footnotesize (2,\allowbreak 12,\allowbreak 14)}\par $K=14$} & {\ttfamily\footnotesize (28,\allowbreak 1)\par\smallskip (26,\allowbreak 2,\allowbreak 27,\allowbreak 4,\allowbreak 25,\allowbreak 3,\allowbreak 23,\allowbreak 5,\allowbreak 24,\allowbreak 8,\allowbreak 22,\allowbreak 9,\allowbreak 21,\allowbreak 6)\par\smallskip (17,\allowbreak 10,\allowbreak 20,\allowbreak 11,\allowbreak 19,\allowbreak 14,\allowbreak 15,\allowbreak 13,\allowbreak 16,\allowbreak 12,\allowbreak 18,\allowbreak 7)}\tabularnewline[5pt]
{\raggedright\arraybackslash {\ttfamily\footnotesize (2,\allowbreak 13,\allowbreak 13)}\par $K=13$\par short leaf $1$} & {\ttfamily\footnotesize (24,\allowbreak 2,\allowbreak 25,\allowbreak 4,\allowbreak 20,\allowbreak 9,\allowbreak 19,\allowbreak 10,\allowbreak 16,\allowbreak 13,\allowbreak 14,\allowbreak 12,\allowbreak 17)\par\smallskip (27,\allowbreak 1)\par\smallskip (28,\allowbreak 3,\allowbreak 22,\allowbreak 5,\allowbreak 23,\allowbreak 8,\allowbreak 21,\allowbreak 7,\allowbreak 15,\allowbreak 11,\allowbreak 18,\allowbreak 6,\allowbreak 26)}\tabularnewline[5pt]
{\raggedright\arraybackslash {\ttfamily\footnotesize (2,\allowbreak 13,\allowbreak 14)}\par $K=14$} & {\ttfamily\footnotesize (29,\allowbreak 1)\par\smallskip (27,\allowbreak 2,\allowbreak 28,\allowbreak 4,\allowbreak 26,\allowbreak 3,\allowbreak 24,\allowbreak 5,\allowbreak 25,\allowbreak 7,\allowbreak 23,\allowbreak 8,\allowbreak 20,\allowbreak 6)\par\smallskip (17,\allowbreak 10,\allowbreak 21,\allowbreak 11,\allowbreak 19,\allowbreak 14,\allowbreak 15,\allowbreak 13,\allowbreak 16,\allowbreak 12,\allowbreak 18,\allowbreak 9,\allowbreak 22)}\tabularnewline[5pt]
{\raggedright\arraybackslash {\ttfamily\footnotesize (2,\allowbreak 14,\allowbreak 15)}\par $K=15$\par $2e\le b+4$} & {\ttfamily\footnotesize (31,\allowbreak 1)\par\smallskip (29,\allowbreak 2,\allowbreak 30,\allowbreak 4,\allowbreak 28,\allowbreak 3,\allowbreak 26,\allowbreak 5,\allowbreak 27,\allowbreak 7,\allowbreak 25,\allowbreak 8,\allowbreak 24,\allowbreak 10,\allowbreak 23)\par\smallskip (19,\allowbreak 11,\allowbreak 22,\allowbreak 12,\allowbreak 17,\allowbreak 13,\allowbreak 20,\allowbreak 14,\allowbreak 16,\allowbreak 15,\allowbreak 18,\allowbreak 9,\allowbreak 21,\allowbreak 6)}\tabularnewline[5pt]
{\raggedright\arraybackslash {\ttfamily\footnotesize (2,\allowbreak 14,\allowbreak 16)}\par $K=16$} & {\ttfamily\footnotesize (32,\allowbreak 1)\par\smallskip (30,\allowbreak 2,\allowbreak 31,\allowbreak 4,\allowbreak 29,\allowbreak 3,\allowbreak 27,\allowbreak 5,\allowbreak 28,\allowbreak 7,\allowbreak 26,\allowbreak 8,\allowbreak 25,\allowbreak 10,\allowbreak 23,\allowbreak 9)\par\smallskip (20,\allowbreak 11,\allowbreak 21,\allowbreak 13,\allowbreak 24,\allowbreak 12,\allowbreak 19,\allowbreak 14,\allowbreak 18,\allowbreak 15,\allowbreak 17,\allowbreak 16,\allowbreak 22,\allowbreak 6)}\tabularnewline[5pt]
{\raggedright\arraybackslash {\ttfamily\footnotesize (2,\allowbreak 15,\allowbreak 15)}\par $K=15$\par short leaf $1$} & {\ttfamily\footnotesize (28,\allowbreak 4,\allowbreak 26,\allowbreak 10,\allowbreak 18,\allowbreak 14,\allowbreak 20,\allowbreak 9,\allowbreak 22,\allowbreak 13,\allowbreak 16,\allowbreak 15,\allowbreak 17,\allowbreak 12,\allowbreak 19)\par\smallskip (31,\allowbreak 1)\par\smallskip (32,\allowbreak 3,\allowbreak 29,\allowbreak 2,\allowbreak 27,\allowbreak 6,\allowbreak 24,\allowbreak 5,\allowbreak 25,\allowbreak 8,\allowbreak 23,\allowbreak 11,\allowbreak 21,\allowbreak 7,\allowbreak 30)}\tabularnewline[5pt]
{\raggedright\arraybackslash {\ttfamily\footnotesize (2,\allowbreak 15,\allowbreak 16)}\par $K=16$} & {\ttfamily\footnotesize (33,\allowbreak 1)\par\smallskip (31,\allowbreak 2,\allowbreak 32,\allowbreak 4,\allowbreak 30,\allowbreak 3,\allowbreak 28,\allowbreak 5,\allowbreak 29,\allowbreak 7,\allowbreak 27,\allowbreak 6,\allowbreak 25,\allowbreak 10,\allowbreak 24,\allowbreak 8)\par\smallskip (18,\allowbreak 14,\allowbreak 21,\allowbreak 11,\allowbreak 23,\allowbreak 12,\allowbreak 20,\allowbreak 15,\allowbreak 17,\allowbreak 16,\allowbreak 19,\allowbreak 13,\allowbreak 22,\allowbreak 9,\allowbreak 26)}\tabularnewline[5pt]
{\raggedright\arraybackslash {\ttfamily\footnotesize (2,\allowbreak 16,\allowbreak 17)}\par $K=17$\par $2e\le b+4$} & {\ttfamily\footnotesize (32,\allowbreak 6,\allowbreak 27,\allowbreak 10,\allowbreak 22,\allowbreak 16,\allowbreak 20,\allowbreak 12,\allowbreak 23,\allowbreak 13,\allowbreak 28,\allowbreak 5,\allowbreak 33,\allowbreak 3,\allowbreak 30,\allowbreak 8,\allowbreak 26)\par\smallskip (34,\allowbreak 1)\par\smallskip (35,\allowbreak 4,\allowbreak 29,\allowbreak 9,\allowbreak 25,\allowbreak 11,\allowbreak 24,\allowbreak 15,\allowbreak 18,\allowbreak 17,\allowbreak 19,\allowbreak 14,\allowbreak 21,\allowbreak 2,\allowbreak 31,\allowbreak 7)}\tabularnewline[5pt]
{\raggedright\arraybackslash {\ttfamily\footnotesize (2,\allowbreak 16,\allowbreak 18)}\par $K=18$} & {\ttfamily\footnotesize (30,\allowbreak 8,\allowbreak 31,\allowbreak 11,\allowbreak 25,\allowbreak 14,\allowbreak 21,\allowbreak 16,\allowbreak 22,\allowbreak 18,\allowbreak 19,\allowbreak 17,\allowbreak 20,\allowbreak 4,\allowbreak 32,\allowbreak 7)\par\smallskip (34,\allowbreak 5,\allowbreak 29,\allowbreak 3,\allowbreak 35,\allowbreak 2,\allowbreak 33,\allowbreak 6,\allowbreak 27,\allowbreak 12,\allowbreak 24,\allowbreak 15,\allowbreak 23,\allowbreak 13,\allowbreak 26,\allowbreak 9,\allowbreak 28,\allowbreak 10)\par\smallskip (36,\allowbreak 1)}\tabularnewline[5pt]
{\raggedright\arraybackslash {\ttfamily\footnotesize (2,\allowbreak 17,\allowbreak 17)}\par $K=17$\par short leaf $1$} & {\ttfamily\footnotesize (36,\allowbreak 1)\par\smallskip (34,\allowbreak 2,\allowbreak 35,\allowbreak 4,\allowbreak 33,\allowbreak 3,\allowbreak 31,\allowbreak 5,\allowbreak 32,\allowbreak 7,\allowbreak 30,\allowbreak 6,\allowbreak 28,\allowbreak 8,\allowbreak 27,\allowbreak 9,\allowbreak 26)\par\smallskip (21,\allowbreak 10,\allowbreak 25,\allowbreak 11,\allowbreak 24,\allowbreak 12,\allowbreak 22,\allowbreak 14,\allowbreak 23,\allowbreak 16,\allowbreak 20,\allowbreak 15,\allowbreak 18,\allowbreak 17,\allowbreak 19,\allowbreak 13,\allowbreak 29)}\tabularnewline[5pt]
{\raggedright\arraybackslash {\ttfamily\footnotesize (2,\allowbreak 17,\allowbreak 18)}\par $K=18$} & {\ttfamily\footnotesize (37,\allowbreak 1)\par\smallskip (35,\allowbreak 2,\allowbreak 36,\allowbreak 4,\allowbreak 34,\allowbreak 3,\allowbreak 32,\allowbreak 5,\allowbreak 33,\allowbreak 7,\allowbreak 31,\allowbreak 6,\allowbreak 29,\allowbreak 8,\allowbreak 28,\allowbreak 9,\allowbreak 27,\allowbreak 10)\par\smallskip (22,\allowbreak 11,\allowbreak 26,\allowbreak 12,\allowbreak 25,\allowbreak 13,\allowbreak 23,\allowbreak 15,\allowbreak 24,\allowbreak 17,\allowbreak 21,\allowbreak 16,\allowbreak 19,\allowbreak 18,\allowbreak 20,\allowbreak 14,\allowbreak 30)}\tabularnewline[5pt]
{\raggedright\arraybackslash {\ttfamily\footnotesize (2,\allowbreak 18,\allowbreak 19)}\par $K=19$\par $2e\le b+4$} & {\ttfamily\footnotesize (35,\allowbreak 9,\allowbreak 30,\allowbreak 10,\allowbreak 23,\allowbreak 15,\allowbreak 25,\allowbreak 2,\allowbreak 38,\allowbreak 4,\allowbreak 34,\allowbreak 6,\allowbreak 33,\allowbreak 11,\allowbreak 26,\allowbreak 14,\allowbreak 28,\allowbreak 12,\allowbreak 29)\par\smallskip (37,\allowbreak 5,\allowbreak 36,\allowbreak 3,\allowbreak 32,\allowbreak 7,\allowbreak 31,\allowbreak 13,\allowbreak 22,\allowbreak 18,\allowbreak 24,\allowbreak 17,\allowbreak 20,\allowbreak 19,\allowbreak 21,\allowbreak 16,\allowbreak 27,\allowbreak 8)\par\smallskip (39,\allowbreak 1)}\tabularnewline[5pt]
{\raggedright\arraybackslash {\ttfamily\footnotesize (2,\allowbreak 18,\allowbreak 20)}\par $K=20$} & {\ttfamily\footnotesize (37,\allowbreak 2,\allowbreak 36,\allowbreak 5,\allowbreak 38,\allowbreak 6,\allowbreak 35,\allowbreak 14,\allowbreak 31,\allowbreak 3,\allowbreak 30,\allowbreak 7,\allowbreak 33,\allowbreak 13,\allowbreak 27,\allowbreak 12,\allowbreak 25,\allowbreak 17,\allowbreak 29,\allowbreak 11)\par\smallskip (39,\allowbreak 1)\par\smallskip (40,\allowbreak 4,\allowbreak 34,\allowbreak 9,\allowbreak 28,\allowbreak 18,\allowbreak 23,\allowbreak 16,\allowbreak 22,\allowbreak 19,\allowbreak 21,\allowbreak 20,\allowbreak 24,\allowbreak 15,\allowbreak 26,\allowbreak 10,\allowbreak 32,\allowbreak 8)}\tabularnewline[5pt]
{\raggedright\arraybackslash {\ttfamily\footnotesize (2,\allowbreak 19,\allowbreak 19)}\par $K=19$\par short leaf $1$} & {\ttfamily\footnotesize (40,\allowbreak 1)\par\smallskip (38,\allowbreak 2,\allowbreak 39,\allowbreak 4,\allowbreak 37,\allowbreak 3,\allowbreak 35,\allowbreak 5,\allowbreak 36,\allowbreak 7,\allowbreak 34,\allowbreak 6,\allowbreak 32,\allowbreak 8,\allowbreak 33,\allowbreak 10,\allowbreak 29,\allowbreak 9,\allowbreak 31)\par\smallskip (21,\allowbreak 17,\allowbreak 24,\allowbreak 14,\allowbreak 27,\allowbreak 11,\allowbreak 28,\allowbreak 13,\allowbreak 25,\allowbreak 16,\allowbreak 22,\allowbreak 19,\allowbreak 20,\allowbreak 18,\allowbreak 23,\allowbreak 15,\allowbreak 26,\allowbreak 12,\allowbreak 30)}\tabularnewline[5pt]
{\raggedright\arraybackslash {\ttfamily\footnotesize (2,\allowbreak 19,\allowbreak 20)}\par $K=20$} & {\ttfamily\footnotesize (41,\allowbreak 1)\par\smallskip (39,\allowbreak 2,\allowbreak 40,\allowbreak 4,\allowbreak 38,\allowbreak 3,\allowbreak 36,\allowbreak 5,\allowbreak 37,\allowbreak 7,\allowbreak 35,\allowbreak 6,\allowbreak 33,\allowbreak 8,\allowbreak 34,\allowbreak 10,\allowbreak 32,\allowbreak 11,\allowbreak 29,\allowbreak 9)\par\smallskip (23,\allowbreak 15,\allowbreak 30,\allowbreak 13,\allowbreak 27,\allowbreak 14,\allowbreak 26,\allowbreak 16,\allowbreak 25,\allowbreak 18,\allowbreak 24,\allowbreak 19,\allowbreak 22,\allowbreak 20,\allowbreak 21,\allowbreak 17,\allowbreak 28,\allowbreak 12,\allowbreak 31)}\tabularnewline[5pt]
{\raggedright\arraybackslash {\ttfamily\footnotesize (2,\allowbreak 20,\allowbreak 21)}\par $K=21$\par $2e\le b+4$} & {\ttfamily\footnotesize (40,\allowbreak 8,\allowbreak 36,\allowbreak 7,\allowbreak 30,\allowbreak 17,\allowbreak 25,\allowbreak 15,\allowbreak 24,\allowbreak 20,\allowbreak 26,\allowbreak 19,\allowbreak 22,\allowbreak 21,\allowbreak 23,\allowbreak 18,\allowbreak 29,\allowbreak 14,\allowbreak 31,\allowbreak 6,\allowbreak 32)\par\smallskip (42,\allowbreak 1)\par\smallskip (43,\allowbreak 4,\allowbreak 41,\allowbreak 3,\allowbreak 34,\allowbreak 12,\allowbreak 33,\allowbreak 13,\allowbreak 27,\allowbreak 11,\allowbreak 35,\allowbreak 16,\allowbreak 28,\allowbreak 10,\allowbreak 37,\allowbreak 2,\allowbreak 38,\allowbreak 5,\allowbreak 39,\allowbreak 9)}\tabularnewline[5pt]
{\raggedright\arraybackslash {\ttfamily\footnotesize (2,\allowbreak 20,\allowbreak 22)}\par $K=22$} & {\ttfamily\footnotesize (42,\allowbreak 1)\par\smallskip (43,\allowbreak 3,\allowbreak 37,\allowbreak 8,\allowbreak 34,\allowbreak 11,\allowbreak 33,\allowbreak 18,\allowbreak 30,\allowbreak 14,\allowbreak 35,\allowbreak 4,\allowbreak 41,\allowbreak 16,\allowbreak 26,\allowbreak 19,\allowbreak 25,\allowbreak 17,\allowbreak 28,\allowbreak 15,\allowbreak 32,\allowbreak 12)\par\smallskip (44,\allowbreak 5,\allowbreak 40,\allowbreak 2,\allowbreak 38,\allowbreak 10,\allowbreak 29,\allowbreak 20,\allowbreak 24,\allowbreak 21,\allowbreak 23,\allowbreak 22,\allowbreak 27,\allowbreak 13,\allowbreak 31,\allowbreak 7,\allowbreak 39,\allowbreak 6,\allowbreak 36,\allowbreak 9)}\tabularnewline[5pt]
{\raggedright\arraybackslash {\ttfamily\footnotesize (2,\allowbreak 21,\allowbreak 21)}\par $K=21$\par short leaf $1$} & {\ttfamily\footnotesize (44,\allowbreak 1)\par\smallskip (42,\allowbreak 2,\allowbreak 43,\allowbreak 4,\allowbreak 41,\allowbreak 3,\allowbreak 39,\allowbreak 5,\allowbreak 40,\allowbreak 7,\allowbreak 38,\allowbreak 6,\allowbreak 36,\allowbreak 8,\allowbreak 37,\allowbreak 10,\allowbreak 35,\allowbreak 9,\allowbreak 32,\allowbreak 11,\allowbreak 33)\par\smallskip (24,\allowbreak 16,\allowbreak 31,\allowbreak 12,\allowbreak 30,\allowbreak 13,\allowbreak 29,\allowbreak 15,\allowbreak 28,\allowbreak 17,\allowbreak 27,\allowbreak 18,\allowbreak 25,\allowbreak 19,\allowbreak 23,\allowbreak 20,\allowbreak 22,\allowbreak 21,\allowbreak 26,\allowbreak 14,\allowbreak 34)}\tabularnewline[5pt]
{\raggedright\arraybackslash {\ttfamily\footnotesize (2,\allowbreak 22,\allowbreak 23)}\par $K=23$\par $2e\le b+4$} & {\ttfamily\footnotesize (43,\allowbreak 9,\allowbreak 42,\allowbreak 3,\allowbreak 41,\allowbreak 6,\allowbreak 46,\allowbreak 2,\allowbreak 44,\allowbreak 7,\allowbreak 33,\allowbreak 16,\allowbreak 27,\allowbreak 20,\allowbreak 25,\allowbreak 23,\allowbreak 24,\allowbreak 21,\allowbreak 31,\allowbreak 13,\allowbreak 32,\allowbreak 10)\par\smallskip (45,\allowbreak 4,\allowbreak 40,\allowbreak 8,\allowbreak 38,\allowbreak 18,\allowbreak 26,\allowbreak 22,\allowbreak 28,\allowbreak 19,\allowbreak 34,\allowbreak 5,\allowbreak 36,\allowbreak 15,\allowbreak 29,\allowbreak 17,\allowbreak 30,\allowbreak 14,\allowbreak 37,\allowbreak 12,\allowbreak 39,\allowbreak 11,\allowbreak 35)\par\smallskip (47,\allowbreak 1)}\tabularnewline[5pt]
{\raggedright\arraybackslash {\ttfamily\footnotesize (2,\allowbreak 22,\allowbreak 24)}\par $K=24$} & {\ttfamily\footnotesize (45,\allowbreak 5,\allowbreak 43,\allowbreak 12,\allowbreak 30,\allowbreak 16,\allowbreak 37,\allowbreak 15,\allowbreak 39,\allowbreak 14,\allowbreak 31,\allowbreak 22,\allowbreak 25,\allowbreak 24,\allowbreak 26,\allowbreak 20,\allowbreak 27,\allowbreak 23,\allowbreak 28,\allowbreak 18,\allowbreak 29,\allowbreak 21,\allowbreak 33,\allowbreak 13)\par\smallskip (47,\allowbreak 1)\par\smallskip (48,\allowbreak 6,\allowbreak 42,\allowbreak 9,\allowbreak 44,\allowbreak 3,\allowbreak 46,\allowbreak 2,\allowbreak 41,\allowbreak 4,\allowbreak 38,\allowbreak 8,\allowbreak 40,\allowbreak 11,\allowbreak 34,\allowbreak 7,\allowbreak 35,\allowbreak 19,\allowbreak 32,\allowbreak 17,\allowbreak 36,\allowbreak 10)}\tabularnewline[5pt]
{\raggedright\arraybackslash {\ttfamily\footnotesize (2,\allowbreak 24,\allowbreak 25)}\par $K=25$\par $2e\le b+4$} & {\ttfamily\footnotesize (48,\allowbreak 3,\allowbreak 49,\allowbreak 2,\allowbreak 44,\allowbreak 9,\allowbreak 42,\allowbreak 20,\allowbreak 33,\allowbreak 23,\allowbreak 28,\allowbreak 19,\allowbreak 31,\allowbreak 17,\allowbreak 37,\allowbreak 14,\allowbreak 45,\allowbreak 8,\allowbreak 46,\allowbreak 6,\allowbreak 47,\allowbreak 4,\allowbreak 43,\allowbreak 11)\par\smallskip (50,\allowbreak 1)\par\smallskip (51,\allowbreak 7,\allowbreak 41,\allowbreak 5,\allowbreak 34,\allowbreak 13,\allowbreak 39,\allowbreak 12,\allowbreak 36,\allowbreak 18,\allowbreak 35,\allowbreak 16,\allowbreak 32,\allowbreak 21,\allowbreak 29,\allowbreak 22,\allowbreak 26,\allowbreak 25,\allowbreak 27,\allowbreak 24,\allowbreak 30,\allowbreak 15,\allowbreak 40,\allowbreak 10,\allowbreak 38)}\tabularnewline[5pt]
{\raggedright\arraybackslash {\ttfamily\footnotesize (2,\allowbreak 24,\allowbreak 26)}\par $K=26$} & {\ttfamily\footnotesize (48,\allowbreak 10,\allowbreak 45,\allowbreak 3,\allowbreak 44,\allowbreak 18,\allowbreak 34,\allowbreak 19,\allowbreak 33,\allowbreak 20,\allowbreak 32,\allowbreak 21,\allowbreak 31,\allowbreak 22,\allowbreak 30,\allowbreak 23,\allowbreak 29,\allowbreak 24,\allowbreak 28,\allowbreak 25,\allowbreak 27,\allowbreak 26,\allowbreak 43,\allowbreak 11)\par\smallskip (50,\allowbreak 6,\allowbreak 46,\allowbreak 9,\allowbreak 42,\allowbreak 12,\allowbreak 40,\allowbreak 13,\allowbreak 38,\allowbreak 7,\allowbreak 41,\allowbreak 5,\allowbreak 51,\allowbreak 2,\allowbreak 49,\allowbreak 4,\allowbreak 47,\allowbreak 8,\allowbreak 37,\allowbreak 15,\allowbreak 39,\allowbreak 16,\allowbreak 36,\allowbreak 17,\allowbreak 35,\allowbreak 14)\par\smallskip (52,\allowbreak 1)}\tabularnewline[5pt]
{\raggedright\arraybackslash {\ttfamily\footnotesize (2,\allowbreak 25,\allowbreak 25)}\par $K=25$\par short leaf $1$} & {\ttfamily\footnotesize (49,\allowbreak 4,\allowbreak 51,\allowbreak 20,\allowbreak 34,\allowbreak 16,\allowbreak 36,\allowbreak 12,\allowbreak 33,\allowbreak 22,\allowbreak 30,\allowbreak 21,\allowbreak 28,\allowbreak 25,\allowbreak 26,\allowbreak 24,\allowbreak 29,\allowbreak 23,\allowbreak 27,\allowbreak 17,\allowbreak 32,\allowbreak 19,\allowbreak 31,\allowbreak 15,\allowbreak 38)\par\smallskip (50,\allowbreak 2,\allowbreak 48,\allowbreak 6,\allowbreak 43,\allowbreak 9,\allowbreak 37,\allowbreak 18,\allowbreak 35,\allowbreak 13,\allowbreak 40,\allowbreak 14,\allowbreak 39,\allowbreak 10,\allowbreak 42,\allowbreak 7,\allowbreak 45,\allowbreak 5,\allowbreak 46,\allowbreak 3,\allowbreak 47,\allowbreak 8,\allowbreak 44,\allowbreak 11,\allowbreak 41)\par\smallskip (52,\allowbreak 1)}\tabularnewline[5pt]
{\raggedright\arraybackslash {\ttfamily\footnotesize (2,\allowbreak 25,\allowbreak 26)}\par $K=26$} & {\ttfamily\footnotesize (50,\allowbreak 2,\allowbreak 45,\allowbreak 4,\allowbreak 48,\allowbreak 9,\allowbreak 39,\allowbreak 10,\allowbreak 41,\allowbreak 16,\allowbreak 38,\allowbreak 18,\allowbreak 30,\allowbreak 11,\allowbreak 49,\allowbreak 3,\allowbreak 52,\allowbreak 5,\allowbreak 47,\allowbreak 15,\allowbreak 36,\allowbreak 20,\allowbreak 35,\allowbreak 12,\allowbreak 40,\allowbreak 14)\par\smallskip (51,\allowbreak 6,\allowbreak 46,\allowbreak 13,\allowbreak 37,\allowbreak 19,\allowbreak 33,\allowbreak 23,\allowbreak 29,\allowbreak 22,\allowbreak 31,\allowbreak 26,\allowbreak 27,\allowbreak 25,\allowbreak 28,\allowbreak 24,\allowbreak 32,\allowbreak 21,\allowbreak 34,\allowbreak 17,\allowbreak 44,\allowbreak 7,\allowbreak 43,\allowbreak 8,\allowbreak 42)\par\smallskip (53,\allowbreak 1)}\tabularnewline[5pt]
{\raggedright\arraybackslash {\ttfamily\footnotesize (2,\allowbreak 26,\allowbreak 27)}\par $K=27$\par $2e\le b+4$} & {\ttfamily\footnotesize (49,\allowbreak 4,\allowbreak 46,\allowbreak 20,\allowbreak 39,\allowbreak 15,\allowbreak 45,\allowbreak 16,\allowbreak 36,\allowbreak 22,\allowbreak 32,\allowbreak 25,\allowbreak 29,\allowbreak 17,\allowbreak 40,\allowbreak 13,\allowbreak 38,\allowbreak 7,\allowbreak 53,\allowbreak 2,\allowbreak 52,\allowbreak 5,\allowbreak 43,\allowbreak 11,\allowbreak 48,\allowbreak 12)\par\smallskip (54,\allowbreak 1)\par\smallskip (55,\allowbreak 3,\allowbreak 51,\allowbreak 8,\allowbreak 47,\allowbreak 6,\allowbreak 50,\allowbreak 10,\allowbreak 44,\allowbreak 9,\allowbreak 42,\allowbreak 14,\allowbreak 35,\allowbreak 18,\allowbreak 33,\allowbreak 24,\allowbreak 30,\allowbreak 27,\allowbreak 28,\allowbreak 26,\allowbreak 31,\allowbreak 23,\allowbreak 34,\allowbreak 21,\allowbreak 37,\allowbreak 19,\allowbreak 41)}\tabularnewline[5pt]
{\raggedright\arraybackslash {\ttfamily\footnotesize (2,\allowbreak 26,\allowbreak 28)}\par $K=28$} & {\ttfamily\footnotesize (51,\allowbreak 4,\allowbreak 52,\allowbreak 6,\allowbreak 49,\allowbreak 9,\allowbreak 40,\allowbreak 16,\allowbreak 46,\allowbreak 2,\allowbreak 54,\allowbreak 5,\allowbreak 50,\allowbreak 8,\allowbreak 45,\allowbreak 18,\allowbreak 41,\allowbreak 19,\allowbreak 37,\allowbreak 22,\allowbreak 36,\allowbreak 20,\allowbreak 39,\allowbreak 13,\allowbreak 47,\allowbreak 12)\par\smallskip (55,\allowbreak 1)\par\smallskip (56,\allowbreak 3,\allowbreak 53,\allowbreak 14,\allowbreak 42,\allowbreak 17,\allowbreak 38,\allowbreak 21,\allowbreak 34,\allowbreak 24,\allowbreak 35,\allowbreak 23,\allowbreak 32,\allowbreak 27,\allowbreak 31,\allowbreak 25,\allowbreak 33,\allowbreak 26,\allowbreak 29,\allowbreak 28,\allowbreak 30,\allowbreak 10,\allowbreak 48,\allowbreak 7,\allowbreak 43,\allowbreak 11,\allowbreak 44,\allowbreak 15)}\tabularnewline[5pt]
{\raggedright\arraybackslash {\ttfamily\footnotesize (2,\allowbreak 27,\allowbreak 27)}\par $K=27$\par short leaf $1$} & {\ttfamily\footnotesize (53,\allowbreak 1)\par\smallskip (55,\allowbreak 6,\allowbreak 47,\allowbreak 9,\allowbreak 45,\allowbreak 15,\allowbreak 40,\allowbreak 13,\allowbreak 42,\allowbreak 16,\allowbreak 39,\allowbreak 17,\allowbreak 36,\allowbreak 19,\allowbreak 37,\allowbreak 23,\allowbreak 33,\allowbreak 12,\allowbreak 49,\allowbreak 14,\allowbreak 48,\allowbreak 4,\allowbreak 54,\allowbreak 3,\allowbreak 43,\allowbreak 11,\allowbreak 44)\par\smallskip (56,\allowbreak 2,\allowbreak 50,\allowbreak 8,\allowbreak 51,\allowbreak 5,\allowbreak 52,\allowbreak 7,\allowbreak 46,\allowbreak 18,\allowbreak 38,\allowbreak 22,\allowbreak 31,\allowbreak 26,\allowbreak 28,\allowbreak 27,\allowbreak 30,\allowbreak 24,\allowbreak 35,\allowbreak 20,\allowbreak 32,\allowbreak 25,\allowbreak 29,\allowbreak 21,\allowbreak 34,\allowbreak 10,\allowbreak 41)}\tabularnewline[5pt]
{\raggedright\arraybackslash {\ttfamily\footnotesize (2,\allowbreak 27,\allowbreak 28)}\par $K=28$} & {\ttfamily\footnotesize (55,\allowbreak 1)\par\smallskip (56,\allowbreak 6,\allowbreak 50,\allowbreak 3,\allowbreak 39,\allowbreak 22,\allowbreak 35,\allowbreak 20,\allowbreak 32,\allowbreak 23,\allowbreak 37,\allowbreak 21,\allowbreak 41,\allowbreak 12,\allowbreak 47,\allowbreak 10,\allowbreak 44,\allowbreak 16,\allowbreak 38,\allowbreak 19,\allowbreak 40,\allowbreak 17,\allowbreak 43,\allowbreak 13,\allowbreak 46,\allowbreak 14,\allowbreak 45)\par\smallskip (57,\allowbreak 4,\allowbreak 52,\allowbreak 9,\allowbreak 48,\allowbreak 7,\allowbreak 49,\allowbreak 11,\allowbreak 51,\allowbreak 5,\allowbreak 54,\allowbreak 2,\allowbreak 53,\allowbreak 8,\allowbreak 33,\allowbreak 27,\allowbreak 29,\allowbreak 28,\allowbreak 31,\allowbreak 24,\allowbreak 34,\allowbreak 26,\allowbreak 30,\allowbreak 25,\allowbreak 36,\allowbreak 18,\allowbreak 42,\allowbreak 15)}\tabularnewline[5pt]
{\raggedright\arraybackslash {\ttfamily\footnotesize (2,\allowbreak 28,\allowbreak 29)}\par $K=29$\par $2e\le b+4$} & {\ttfamily\footnotesize (55,\allowbreak 7,\allowbreak 49,\allowbreak 8,\allowbreak 52,\allowbreak 6,\allowbreak 56,\allowbreak 3,\allowbreak 46,\allowbreak 15,\allowbreak 40,\allowbreak 22,\allowbreak 35,\allowbreak 23,\allowbreak 32,\allowbreak 28,\allowbreak 34,\allowbreak 24,\allowbreak 38,\allowbreak 19,\allowbreak 39,\allowbreak 17,\allowbreak 50,\allowbreak 10,\allowbreak 45,\allowbreak 16,\allowbreak 43,\allowbreak 13)\par\smallskip (57,\allowbreak 5,\allowbreak 54,\allowbreak 9,\allowbreak 47,\allowbreak 11,\allowbreak 48,\allowbreak 14,\allowbreak 42,\allowbreak 18,\allowbreak 41,\allowbreak 20,\allowbreak 37,\allowbreak 21,\allowbreak 36,\allowbreak 25,\allowbreak 33,\allowbreak 26,\allowbreak 31,\allowbreak 29,\allowbreak 30,\allowbreak 27,\allowbreak 53,\allowbreak 2,\allowbreak 58,\allowbreak 4,\allowbreak 51,\allowbreak 12,\allowbreak 44)\par\smallskip (59,\allowbreak 1)}\tabularnewline[5pt]
{\raggedright\arraybackslash {\ttfamily\footnotesize (2,\allowbreak 28,\allowbreak 30)}\par $K=30$} & {\ttfamily\footnotesize (48,\allowbreak 12,\allowbreak 47,\allowbreak 14,\allowbreak 46,\allowbreak 17,\allowbreak 43,\allowbreak 18,\allowbreak 45,\allowbreak 15,\allowbreak 32,\allowbreak 29,\allowbreak 31,\allowbreak 30,\allowbreak 34,\allowbreak 27,\allowbreak 33,\allowbreak 28,\allowbreak 36,\allowbreak 25,\allowbreak 35,\allowbreak 26,\allowbreak 38,\allowbreak 22,\allowbreak 42,\allowbreak 19,\allowbreak 41,\allowbreak 20,\allowbreak 44,\allowbreak 16)\par\smallskip (60,\allowbreak 3,\allowbreak 57,\allowbreak 2,\allowbreak 58,\allowbreak 5,\allowbreak 55,\allowbreak 4,\allowbreak 56,\allowbreak 7,\allowbreak 54,\allowbreak 23,\allowbreak 37,\allowbreak 24,\allowbreak 39,\allowbreak 21,\allowbreak 40,\allowbreak 6,\allowbreak 52,\allowbreak 10,\allowbreak 51,\allowbreak 8,\allowbreak 53,\allowbreak 9,\allowbreak 49,\allowbreak 11,\allowbreak 50,\allowbreak 13)\par\smallskip (59,\allowbreak 1)}\tabularnewline[5pt]
{\raggedright\arraybackslash {\ttfamily\footnotesize (2,\allowbreak 29,\allowbreak 29)}\par $K=29$\par short leaf $1$} & {\ttfamily\footnotesize (55,\allowbreak 4,\allowbreak 57,\allowbreak 15,\allowbreak 51,\allowbreak 12,\allowbreak 42,\allowbreak 24,\allowbreak 36,\allowbreak 28,\allowbreak 30,\allowbreak 29,\allowbreak 32,\allowbreak 26,\allowbreak 31,\allowbreak 27,\allowbreak 34,\allowbreak 25,\allowbreak 38,\allowbreak 11,\allowbreak 48,\allowbreak 13,\allowbreak 46,\allowbreak 14,\allowbreak 35,\allowbreak 21,\allowbreak 40,\allowbreak 16,\allowbreak 47)\par\smallskip (58,\allowbreak 6,\allowbreak 56,\allowbreak 2,\allowbreak 59,\allowbreak 3,\allowbreak 49,\allowbreak 20,\allowbreak 37,\allowbreak 22,\allowbreak 33,\allowbreak 23,\allowbreak 39,\allowbreak 19,\allowbreak 41,\allowbreak 18,\allowbreak 43,\allowbreak 17,\allowbreak 45,\allowbreak 7,\allowbreak 54,\allowbreak 5,\allowbreak 53,\allowbreak 8,\allowbreak 52,\allowbreak 9,\allowbreak 50,\allowbreak 10,\allowbreak 44)\par\smallskip (60,\allowbreak 1)}\tabularnewline[5pt]
{\raggedright\arraybackslash {\ttfamily\footnotesize (2,\allowbreak 30,\allowbreak 31)}\par $K=31$\par $2e\le b+4$} & {\ttfamily\footnotesize (56,\allowbreak 8,\allowbreak 49,\allowbreak 16,\allowbreak 50,\allowbreak 13,\allowbreak 44,\allowbreak 6,\allowbreak 57,\allowbreak 5,\allowbreak 60,\allowbreak 7,\allowbreak 51,\allowbreak 15,\allowbreak 54,\allowbreak 4,\allowbreak 58,\allowbreak 9,\allowbreak 55,\allowbreak 10,\allowbreak 53,\allowbreak 11,\allowbreak 43,\allowbreak 24,\allowbreak 39,\allowbreak 21,\allowbreak 41,\allowbreak 20,\allowbreak 48,\allowbreak 18,\allowbreak 47)\par\smallskip (62,\allowbreak 1)\par\smallskip (63,\allowbreak 3,\allowbreak 61,\allowbreak 2,\allowbreak 59,\allowbreak 12,\allowbreak 52,\allowbreak 17,\allowbreak 42,\allowbreak 30,\allowbreak 33,\allowbreak 31,\allowbreak 32,\allowbreak 28,\allowbreak 36,\allowbreak 25,\allowbreak 38,\allowbreak 22,\allowbreak 46,\allowbreak 19,\allowbreak 45,\allowbreak 23,\allowbreak 40,\allowbreak 26,\allowbreak 35,\allowbreak 29,\allowbreak 34,\allowbreak 27,\allowbreak 37,\allowbreak 14)}\tabularnewline[5pt]
{\raggedright\arraybackslash {\ttfamily\footnotesize (2,\allowbreak 32,\allowbreak 33)}\par $K=33$\par $2e\le b+4$} & {\ttfamily\footnotesize (64,\allowbreak 1)\par\smallskip (66,\allowbreak 5,\allowbreak 65,\allowbreak 3,\allowbreak 58,\allowbreak 7,\allowbreak 54,\allowbreak 18,\allowbreak 49,\allowbreak 16,\allowbreak 55,\allowbreak 6,\allowbreak 62,\allowbreak 4,\allowbreak 63,\allowbreak 9,\allowbreak 57,\allowbreak 14,\allowbreak 51,\allowbreak 10,\allowbreak 60,\allowbreak 8,\allowbreak 61,\allowbreak 17,\allowbreak 46,\allowbreak 20,\allowbreak 44,\allowbreak 24,\allowbreak 45,\allowbreak 22,\allowbreak 52,\allowbreak 12,\allowbreak 50)\par\smallskip (67,\allowbreak 2,\allowbreak 59,\allowbreak 13,\allowbreak 48,\allowbreak 21,\allowbreak 43,\allowbreak 27,\allowbreak 38,\allowbreak 26,\allowbreak 40,\allowbreak 25,\allowbreak 42,\allowbreak 23,\allowbreak 41,\allowbreak 28,\allowbreak 37,\allowbreak 33,\allowbreak 34,\allowbreak 32,\allowbreak 35,\allowbreak 30,\allowbreak 36,\allowbreak 29,\allowbreak 39,\allowbreak 31,\allowbreak 56,\allowbreak 11,\allowbreak 53,\allowbreak 19,\allowbreak 47,\allowbreak 15)}\tabularnewline[5pt]
{\raggedright\arraybackslash {\ttfamily\footnotesize (2,\allowbreak 32,\allowbreak 34)}\par $K=34$} & {\ttfamily\footnotesize (54,\allowbreak 14,\allowbreak 53,\allowbreak 16,\allowbreak 36,\allowbreak 33,\allowbreak 35,\allowbreak 34,\allowbreak 38,\allowbreak 31,\allowbreak 37,\allowbreak 32,\allowbreak 40,\allowbreak 30,\allowbreak 39,\allowbreak 28,\allowbreak 41,\allowbreak 29,\allowbreak 43,\allowbreak 26,\allowbreak 42,\allowbreak 27,\allowbreak 45,\allowbreak 23,\allowbreak 46,\allowbreak 25,\allowbreak 44,\allowbreak 8,\allowbreak 59,\allowbreak 11,\allowbreak 56,\allowbreak 15)\par\smallskip (68,\allowbreak 3,\allowbreak 65,\allowbreak 2,\allowbreak 66,\allowbreak 5,\allowbreak 63,\allowbreak 4,\allowbreak 64,\allowbreak 7,\allowbreak 60,\allowbreak 10,\allowbreak 57,\allowbreak 13,\allowbreak 55,\allowbreak 12,\allowbreak 58,\allowbreak 9,\allowbreak 61,\allowbreak 6,\allowbreak 62,\allowbreak 24,\allowbreak 48,\allowbreak 21,\allowbreak 47,\allowbreak 22,\allowbreak 50,\allowbreak 19,\allowbreak 49,\allowbreak 20,\allowbreak 52,\allowbreak 17,\allowbreak 51,\allowbreak 18)\par\smallskip (67,\allowbreak 1)}\tabularnewline[5pt]
{\raggedright\arraybackslash {\ttfamily\footnotesize (2,\allowbreak 33,\allowbreak 33)}\par $K=33$\par short leaf $1$} & {\ttfamily\footnotesize (62,\allowbreak 6,\allowbreak 57,\allowbreak 8,\allowbreak 65,\allowbreak 2,\allowbreak 66,\allowbreak 14,\allowbreak 48,\allowbreak 21,\allowbreak 43,\allowbreak 24,\allowbreak 40,\allowbreak 23,\allowbreak 41,\allowbreak 28,\allowbreak 39,\allowbreak 29,\allowbreak 36,\allowbreak 27,\allowbreak 42,\allowbreak 22,\allowbreak 45,\allowbreak 16,\allowbreak 56,\allowbreak 19,\allowbreak 49,\allowbreak 17,\allowbreak 55,\allowbreak 12,\allowbreak 59,\allowbreak 9,\allowbreak 53)\par\smallskip (67,\allowbreak 1)\par\smallskip (68,\allowbreak 3,\allowbreak 64,\allowbreak 4,\allowbreak 63,\allowbreak 5,\allowbreak 60,\allowbreak 7,\allowbreak 61,\allowbreak 15,\allowbreak 54,\allowbreak 13,\allowbreak 58,\allowbreak 10,\allowbreak 52,\allowbreak 26,\allowbreak 47,\allowbreak 11,\allowbreak 46,\allowbreak 18,\allowbreak 51,\allowbreak 20,\allowbreak 44,\allowbreak 30,\allowbreak 38,\allowbreak 33,\allowbreak 34,\allowbreak 32,\allowbreak 35,\allowbreak 31,\allowbreak 37,\allowbreak 25,\allowbreak 50)}\tabularnewline[5pt]
{\raggedright\arraybackslash {\ttfamily\footnotesize (2,\allowbreak 33,\allowbreak 34)}\par $K=34$} & {\ttfamily\footnotesize (39,\allowbreak 26,\allowbreak 45,\allowbreak 28,\allowbreak 44,\allowbreak 24,\allowbreak 46,\allowbreak 23,\allowbreak 47,\allowbreak 22,\allowbreak 48,\allowbreak 21,\allowbreak 42,\allowbreak 27,\allowbreak 41,\allowbreak 30,\allowbreak 40,\allowbreak 32,\allowbreak 37,\allowbreak 33,\allowbreak 36,\allowbreak 34,\allowbreak 35,\allowbreak 29,\allowbreak 38,\allowbreak 31,\allowbreak 43,\allowbreak 25,\allowbreak 53,\allowbreak 20,\allowbreak 49,\allowbreak 19,\allowbreak 50,\allowbreak 18)\par\smallskip (68,\allowbreak 1)\par\smallskip (69,\allowbreak 3,\allowbreak 66,\allowbreak 2,\allowbreak 67,\allowbreak 5,\allowbreak 63,\allowbreak 6,\allowbreak 65,\allowbreak 4,\allowbreak 64,\allowbreak 8,\allowbreak 60,\allowbreak 9,\allowbreak 62,\allowbreak 7,\allowbreak 61,\allowbreak 11,\allowbreak 57,\allowbreak 13,\allowbreak 58,\allowbreak 10,\allowbreak 59,\allowbreak 12,\allowbreak 55,\allowbreak 17,\allowbreak 51,\allowbreak 16,\allowbreak 52,\allowbreak 15,\allowbreak 56,\allowbreak 14,\allowbreak 54)}\tabularnewline[5pt]
{\raggedright\arraybackslash {\ttfamily\footnotesize (2,\allowbreak 34,\allowbreak 35)}\par $K=35$\par $2e\le b+4$} & {\ttfamily\footnotesize (71,\allowbreak 3,\allowbreak 68,\allowbreak 2,\allowbreak 69,\allowbreak 5,\allowbreak 66,\allowbreak 4,\allowbreak 67,\allowbreak 7,\allowbreak 63,\allowbreak 27,\allowbreak 45,\allowbreak 25,\allowbreak 48,\allowbreak 24,\allowbreak 46,\allowbreak 8,\allowbreak 65,\allowbreak 6,\allowbreak 64,\allowbreak 10,\allowbreak 60,\allowbreak 13,\allowbreak 59,\allowbreak 11,\allowbreak 62,\allowbreak 9,\allowbreak 61,\allowbreak 12,\allowbreak 57,\allowbreak 14,\allowbreak 58,\allowbreak 16)\par\smallskip (55,\allowbreak 15,\allowbreak 56,\allowbreak 17,\allowbreak 54,\allowbreak 19,\allowbreak 52,\allowbreak 18,\allowbreak 37,\allowbreak 33,\allowbreak 38,\allowbreak 35,\allowbreak 36,\allowbreak 34,\allowbreak 40,\allowbreak 32,\allowbreak 39,\allowbreak 30,\allowbreak 42,\allowbreak 31,\allowbreak 41,\allowbreak 28,\allowbreak 44,\allowbreak 29,\allowbreak 43,\allowbreak 26,\allowbreak 47,\allowbreak 22,\allowbreak 50,\allowbreak 23,\allowbreak 49,\allowbreak 20,\allowbreak 51,\allowbreak 21,\allowbreak 53)\par\smallskip (70,\allowbreak 1)}\tabularnewline[5pt]
{\raggedright\arraybackslash {\ttfamily\footnotesize (2,\allowbreak 34,\allowbreak 36)}\par $K=36$} & {\ttfamily\footnotesize (72,\allowbreak 3,\allowbreak 69,\allowbreak 2,\allowbreak 70,\allowbreak 5,\allowbreak 67,\allowbreak 4,\allowbreak 68,\allowbreak 7,\allowbreak 64,\allowbreak 10,\allowbreak 61,\allowbreak 12,\allowbreak 62,\allowbreak 9,\allowbreak 65,\allowbreak 6,\allowbreak 66,\allowbreak 8,\allowbreak 47,\allowbreak 26,\allowbreak 46,\allowbreak 29,\allowbreak 43,\allowbreak 33,\allowbreak 40,\allowbreak 32,\allowbreak 41,\allowbreak 35,\allowbreak 39,\allowbreak 34,\allowbreak 37,\allowbreak 36,\allowbreak 38,\allowbreak 19)\par\smallskip (71,\allowbreak 1)\par\smallskip (55,\allowbreak 21,\allowbreak 52,\allowbreak 20,\allowbreak 53,\allowbreak 23,\allowbreak 50,\allowbreak 22,\allowbreak 51,\allowbreak 25,\allowbreak 48,\allowbreak 24,\allowbreak 49,\allowbreak 27,\allowbreak 45,\allowbreak 30,\allowbreak 42,\allowbreak 31,\allowbreak 44,\allowbreak 28,\allowbreak 63,\allowbreak 11,\allowbreak 59,\allowbreak 15,\allowbreak 58,\allowbreak 13,\allowbreak 60,\allowbreak 14,\allowbreak 56,\allowbreak 18,\allowbreak 54,\allowbreak 17,\allowbreak 57,\allowbreak 16)}\tabularnewline[5pt]
\end{longtable}
\endgroup

\subsection{The 44 direct minimum-three cores}
\label{aux:min3-small}
\begingroup\footnotesize
\begin{longtable}{@{}p{25mm}p{115mm}@{}}
Profile / interface & Complete outward words\\\hline\endhead
{\raggedright\arraybackslash {\ttfamily\footnotesize (3,\allowbreak 3,\allowbreak 9)}\par $K=6$} & {\ttfamily\footnotesize (15,\allowbreak 1,\allowbreak 14)\par\smallskip (12,\allowbreak 2,\allowbreak 10,\allowbreak 3,\allowbreak 8,\allowbreak 6,\allowbreak 7,\allowbreak 4,\allowbreak 13)\par\smallskip (11,\allowbreak 5,\allowbreak 9)}\tabularnewline[5pt]
{\raggedright\arraybackslash {\ttfamily\footnotesize (3,\allowbreak 4,\allowbreak 13)}\par $K=9$} & {\ttfamily\footnotesize (20,\allowbreak 1,\allowbreak 19)\par\smallskip (17,\allowbreak 2,\allowbreak 18,\allowbreak 4,\allowbreak 16,\allowbreak 5,\allowbreak 15,\allowbreak 7,\allowbreak 11,\allowbreak 8,\allowbreak 10,\allowbreak 9,\allowbreak 14)\par\smallskip (13,\allowbreak 6,\allowbreak 12,\allowbreak 3)}\tabularnewline[5pt]
{\raggedright\arraybackslash {\ttfamily\footnotesize (3,\allowbreak 5,\allowbreak 13)}\par $K=9$} & {\ttfamily\footnotesize (21,\allowbreak 1,\allowbreak 20)\par\smallskip (18,\allowbreak 2,\allowbreak 19,\allowbreak 4,\allowbreak 17,\allowbreak 3,\allowbreak 14,\allowbreak 5,\allowbreak 13,\allowbreak 7,\allowbreak 11,\allowbreak 6,\allowbreak 16)\par\smallskip (12,\allowbreak 9,\allowbreak 10,\allowbreak 8,\allowbreak 15)}\tabularnewline[5pt]
{\raggedright\arraybackslash {\ttfamily\footnotesize (3,\allowbreak 6,\allowbreak 7)}\par $K=7$} & {\ttfamily\footnotesize (16,\allowbreak 1,\allowbreak 15)\par\smallskip (13,\allowbreak 2,\allowbreak 14,\allowbreak 4,\allowbreak 10,\allowbreak 5,\allowbreak 12)\par\smallskip (9,\allowbreak 6,\allowbreak 8,\allowbreak 7,\allowbreak 11,\allowbreak 3)}\tabularnewline[5pt]
{\raggedright\arraybackslash {\ttfamily\footnotesize (3,\allowbreak 6,\allowbreak 8)}\par $K=8$} & {\ttfamily\footnotesize (17,\allowbreak 1,\allowbreak 16)\par\smallskip (14,\allowbreak 2,\allowbreak 15,\allowbreak 5,\allowbreak 13,\allowbreak 6,\allowbreak 12,\allowbreak 3)\par\smallskip (11,\allowbreak 7,\allowbreak 10,\allowbreak 8,\allowbreak 9,\allowbreak 4)}\tabularnewline[5pt]
{\raggedright\arraybackslash {\ttfamily\footnotesize (3,\allowbreak 6,\allowbreak 17)}\par $K=12$} & {\ttfamily\footnotesize (26,\allowbreak 1,\allowbreak 25)\par\smallskip (23,\allowbreak 2,\allowbreak 24,\allowbreak 4,\allowbreak 22,\allowbreak 3,\allowbreak 20,\allowbreak 5,\allowbreak 19,\allowbreak 6,\allowbreak 18,\allowbreak 8,\allowbreak 17,\allowbreak 9,\allowbreak 15,\allowbreak 10,\allowbreak 21)\par\smallskip (16,\allowbreak 12,\allowbreak 13,\allowbreak 11,\allowbreak 14,\allowbreak 7)}\tabularnewline[5pt]
{\raggedright\arraybackslash {\ttfamily\footnotesize (3,\allowbreak 7,\allowbreak 7)}\par $K=7$} & {\ttfamily\footnotesize (17,\allowbreak 1,\allowbreak 16)\par\smallskip (14,\allowbreak 2,\allowbreak 15,\allowbreak 4,\allowbreak 11,\allowbreak 3,\allowbreak 13)\par\smallskip (9,\allowbreak 5,\allowbreak 10,\allowbreak 7,\allowbreak 8,\allowbreak 6,\allowbreak 12)}\tabularnewline[5pt]
{\raggedright\arraybackslash {\ttfamily\footnotesize (3,\allowbreak 7,\allowbreak 8)}\par $K=8$} & {\ttfamily\footnotesize (18,\allowbreak 1,\allowbreak 17)\par\smallskip (15,\allowbreak 2,\allowbreak 16,\allowbreak 4,\allowbreak 12,\allowbreak 5,\allowbreak 14,\allowbreak 3)\par\smallskip (10,\allowbreak 6,\allowbreak 11,\allowbreak 8,\allowbreak 9,\allowbreak 7,\allowbreak 13)}\tabularnewline[5pt]
{\raggedright\arraybackslash {\ttfamily\footnotesize (3,\allowbreak 7,\allowbreak 17)}\par $K=12$} & {\ttfamily\footnotesize (27,\allowbreak 1,\allowbreak 26)\par\smallskip (24,\allowbreak 2,\allowbreak 25,\allowbreak 4,\allowbreak 23,\allowbreak 3,\allowbreak 21,\allowbreak 5,\allowbreak 20,\allowbreak 6,\allowbreak 19,\allowbreak 7,\allowbreak 16,\allowbreak 10,\allowbreak 15,\allowbreak 8,\allowbreak 18)\par\smallskip (17,\allowbreak 9,\allowbreak 13,\allowbreak 12,\allowbreak 14,\allowbreak 11,\allowbreak 22)}\tabularnewline[5pt]
{\raggedright\arraybackslash {\ttfamily\footnotesize (3,\allowbreak 8,\allowbreak 8)}\par $K=9$} & {\ttfamily\footnotesize (19,\allowbreak 1,\allowbreak 18)\par\smallskip (16,\allowbreak 2,\allowbreak 17,\allowbreak 4,\allowbreak 15,\allowbreak 7,\allowbreak 14,\allowbreak 5)\par\smallskip (12,\allowbreak 6,\allowbreak 11,\allowbreak 8,\allowbreak 10,\allowbreak 9,\allowbreak 13,\allowbreak 3)}\tabularnewline[5pt]
{\raggedright\arraybackslash {\ttfamily\footnotesize (3,\allowbreak 8,\allowbreak 9)}\par $K=9$} & {\ttfamily\footnotesize (20,\allowbreak 1,\allowbreak 19)\par\smallskip (17,\allowbreak 2,\allowbreak 18,\allowbreak 4,\allowbreak 16,\allowbreak 5,\allowbreak 15,\allowbreak 6,\allowbreak 14)\par\smallskip (13,\allowbreak 7,\allowbreak 12,\allowbreak 8,\allowbreak 11,\allowbreak 9,\allowbreak 10,\allowbreak 3)}\tabularnewline[5pt]
{\raggedright\arraybackslash {\ttfamily\footnotesize (3,\allowbreak 8,\allowbreak 10)}\par $K=10$} & {\ttfamily\footnotesize (21,\allowbreak 1,\allowbreak 20)\par\smallskip (18,\allowbreak 2,\allowbreak 19,\allowbreak 4,\allowbreak 17,\allowbreak 5,\allowbreak 16,\allowbreak 8,\allowbreak 15,\allowbreak 6)\par\smallskip (14,\allowbreak 9,\allowbreak 12,\allowbreak 10,\allowbreak 11,\allowbreak 7,\allowbreak 13,\allowbreak 3)}\tabularnewline[5pt]
{\raggedright\arraybackslash {\ttfamily\footnotesize (3,\allowbreak 8,\allowbreak 21)}\par $K=15$} & {\ttfamily\footnotesize (32,\allowbreak 1,\allowbreak 31)\par\smallskip (29,\allowbreak 2,\allowbreak 30,\allowbreak 4,\allowbreak 28,\allowbreak 3,\allowbreak 26,\allowbreak 5,\allowbreak 27,\allowbreak 7,\allowbreak 25,\allowbreak 8,\allowbreak 24,\allowbreak 9,\allowbreak 22,\allowbreak 10,\allowbreak 21,\allowbreak 11,\allowbreak 18,\allowbreak 14,\allowbreak 23)\par\smallskip (19,\allowbreak 13,\allowbreak 16,\allowbreak 15,\allowbreak 17,\allowbreak 12,\allowbreak 20,\allowbreak 6)}\tabularnewline[5pt]
{\raggedright\arraybackslash {\ttfamily\footnotesize (3,\allowbreak 9,\allowbreak 9)}\par $K=9$} & {\ttfamily\footnotesize (21,\allowbreak 1,\allowbreak 20)\par\smallskip (18,\allowbreak 2,\allowbreak 19,\allowbreak 4,\allowbreak 17,\allowbreak 3,\allowbreak 14,\allowbreak 5,\allowbreak 15)\par\smallskip (12,\allowbreak 7,\allowbreak 13,\allowbreak 6,\allowbreak 10,\allowbreak 9,\allowbreak 11,\allowbreak 8,\allowbreak 16)}\tabularnewline[5pt]
{\raggedright\arraybackslash {\ttfamily\footnotesize (3,\allowbreak 9,\allowbreak 10)}\par $K=10$} & {\ttfamily\footnotesize (22,\allowbreak 1,\allowbreak 21)\par\smallskip (19,\allowbreak 2,\allowbreak 20,\allowbreak 4,\allowbreak 18,\allowbreak 3,\allowbreak 16,\allowbreak 7,\allowbreak 15,\allowbreak 5)\par\smallskip (12,\allowbreak 8,\allowbreak 14,\allowbreak 9,\allowbreak 11,\allowbreak 10,\allowbreak 13,\allowbreak 6,\allowbreak 17)}\tabularnewline[5pt]
{\raggedright\arraybackslash {\ttfamily\footnotesize (3,\allowbreak 9,\allowbreak 21)}\par $K=15$} & {\ttfamily\footnotesize (33,\allowbreak 1,\allowbreak 32)\par\smallskip (30,\allowbreak 2,\allowbreak 31,\allowbreak 4,\allowbreak 29,\allowbreak 3,\allowbreak 27,\allowbreak 5,\allowbreak 28,\allowbreak 7,\allowbreak 26,\allowbreak 6,\allowbreak 23,\allowbreak 8,\allowbreak 22,\allowbreak 10,\allowbreak 21,\allowbreak 13,\allowbreak 20,\allowbreak 11,\allowbreak 24)\par\smallskip (18,\allowbreak 12,\allowbreak 17,\allowbreak 14,\allowbreak 16,\allowbreak 15,\allowbreak 19,\allowbreak 9,\allowbreak 25)}\tabularnewline[5pt]
{\raggedright\arraybackslash {\ttfamily\footnotesize (3,\allowbreak 10,\allowbreak 10)}\par $K=11$} & {\ttfamily\footnotesize (23,\allowbreak 1,\allowbreak 22)\par\smallskip (20,\allowbreak 2,\allowbreak 21,\allowbreak 4,\allowbreak 19,\allowbreak 5,\allowbreak 18,\allowbreak 6,\allowbreak 17,\allowbreak 7)\par\smallskip (16,\allowbreak 8,\allowbreak 15,\allowbreak 9,\allowbreak 14,\allowbreak 10,\allowbreak 13,\allowbreak 11,\allowbreak 12,\allowbreak 3)}\tabularnewline[5pt]
{\raggedright\arraybackslash {\ttfamily\footnotesize (3,\allowbreak 10,\allowbreak 11)}\par $K=11$} & {\ttfamily\footnotesize (24,\allowbreak 1,\allowbreak 23)\par\smallskip (21,\allowbreak 2,\allowbreak 22,\allowbreak 4,\allowbreak 20,\allowbreak 3,\allowbreak 18,\allowbreak 5,\allowbreak 17,\allowbreak 8,\allowbreak 19)\par\smallskip (14,\allowbreak 7,\allowbreak 15,\allowbreak 9,\allowbreak 13,\allowbreak 10,\allowbreak 12,\allowbreak 11,\allowbreak 16,\allowbreak 6)}\tabularnewline[5pt]
{\raggedright\arraybackslash {\ttfamily\footnotesize (3,\allowbreak 10,\allowbreak 12)}\par $K=12$} & {\ttfamily\footnotesize (25,\allowbreak 1,\allowbreak 24)\par\smallskip (22,\allowbreak 2,\allowbreak 23,\allowbreak 4,\allowbreak 21,\allowbreak 3,\allowbreak 19,\allowbreak 7,\allowbreak 16,\allowbreak 8,\allowbreak 18,\allowbreak 5)\par\smallskip (15,\allowbreak 10,\allowbreak 17,\allowbreak 11,\allowbreak 14,\allowbreak 12,\allowbreak 13,\allowbreak 9,\allowbreak 20,\allowbreak 6)}\tabularnewline[5pt]
{\raggedright\arraybackslash {\ttfamily\footnotesize (3,\allowbreak 10,\allowbreak 25)}\par $K=18$} & {\ttfamily\footnotesize (38,\allowbreak 1,\allowbreak 37)\par\smallskip (35,\allowbreak 2,\allowbreak 36,\allowbreak 4,\allowbreak 34,\allowbreak 3,\allowbreak 32,\allowbreak 5,\allowbreak 33,\allowbreak 7,\allowbreak 31,\allowbreak 6,\allowbreak 29,\allowbreak 8,\allowbreak 28,\allowbreak 9,\allowbreak 27,\allowbreak 11,\allowbreak 26,\allowbreak 12,\allowbreak 24,\allowbreak 16,\allowbreak 20,\allowbreak 13,\allowbreak 30)\par\smallskip (22,\allowbreak 17,\allowbreak 19,\allowbreak 18,\allowbreak 21,\allowbreak 15,\allowbreak 25,\allowbreak 14,\allowbreak 23,\allowbreak 10)}\tabularnewline[5pt]
{\raggedright\arraybackslash {\ttfamily\footnotesize (3,\allowbreak 11,\allowbreak 11)}\par $K=11$} & {\ttfamily\footnotesize (25,\allowbreak 1,\allowbreak 24)\par\smallskip (22,\allowbreak 2,\allowbreak 23,\allowbreak 4,\allowbreak 21,\allowbreak 3,\allowbreak 19,\allowbreak 5,\allowbreak 18,\allowbreak 6,\allowbreak 17)\par\smallskip (15,\allowbreak 7,\allowbreak 16,\allowbreak 9,\allowbreak 14,\allowbreak 8,\allowbreak 12,\allowbreak 11,\allowbreak 13,\allowbreak 10,\allowbreak 20)}\tabularnewline[5pt]
{\raggedright\arraybackslash {\ttfamily\footnotesize (3,\allowbreak 11,\allowbreak 12)}\par $K=12$} & {\ttfamily\footnotesize (26,\allowbreak 1,\allowbreak 25)\par\smallskip (23,\allowbreak 2,\allowbreak 24,\allowbreak 4,\allowbreak 22,\allowbreak 3,\allowbreak 20,\allowbreak 5,\allowbreak 19,\allowbreak 6,\allowbreak 18,\allowbreak 7)\par\smallskip (16,\allowbreak 8,\allowbreak 17,\allowbreak 10,\allowbreak 15,\allowbreak 9,\allowbreak 13,\allowbreak 12,\allowbreak 14,\allowbreak 11,\allowbreak 21)}\tabularnewline[5pt]
{\raggedright\arraybackslash {\ttfamily\footnotesize (3,\allowbreak 11,\allowbreak 25)}\par $K=18$} & {\ttfamily\footnotesize (39,\allowbreak 1,\allowbreak 38)\par\smallskip (36,\allowbreak 2,\allowbreak 37,\allowbreak 4,\allowbreak 35,\allowbreak 3,\allowbreak 33,\allowbreak 5,\allowbreak 34,\allowbreak 7,\allowbreak 32,\allowbreak 6,\allowbreak 30,\allowbreak 8,\allowbreak 29,\allowbreak 9,\allowbreak 28,\allowbreak 10,\allowbreak 26,\allowbreak 11,\allowbreak 25,\allowbreak 12,\allowbreak 24,\allowbreak 14,\allowbreak 31)\par\smallskip (23,\allowbreak 15,\allowbreak 22,\allowbreak 13,\allowbreak 19,\allowbreak 18,\allowbreak 20,\allowbreak 17,\allowbreak 21,\allowbreak 16,\allowbreak 27)}\tabularnewline[5pt]
{\raggedright\arraybackslash {\ttfamily\footnotesize (3,\allowbreak 12,\allowbreak 12)}\par $K=13$} & {\ttfamily\footnotesize (27,\allowbreak 1,\allowbreak 26)\par\smallskip (24,\allowbreak 2,\allowbreak 25,\allowbreak 4,\allowbreak 23,\allowbreak 3,\allowbreak 21,\allowbreak 6,\allowbreak 22,\allowbreak 8,\allowbreak 20,\allowbreak 7)\par\smallskip (17,\allowbreak 9,\allowbreak 19,\allowbreak 10,\allowbreak 15,\allowbreak 11,\allowbreak 18,\allowbreak 12,\allowbreak 14,\allowbreak 13,\allowbreak 16,\allowbreak 5)}\tabularnewline[5pt]
{\raggedright\arraybackslash {\ttfamily\footnotesize (3,\allowbreak 12,\allowbreak 13)}\par $K=13$} & {\ttfamily\footnotesize (28,\allowbreak 1,\allowbreak 27)\par\smallskip (25,\allowbreak 2,\allowbreak 26,\allowbreak 4,\allowbreak 24,\allowbreak 3,\allowbreak 22,\allowbreak 5,\allowbreak 23,\allowbreak 7,\allowbreak 19,\allowbreak 8,\allowbreak 21)\par\smallskip (15,\allowbreak 11,\allowbreak 18,\allowbreak 9,\allowbreak 17,\allowbreak 12,\allowbreak 14,\allowbreak 13,\allowbreak 16,\allowbreak 10,\allowbreak 20,\allowbreak 6)}\tabularnewline[5pt]
{\raggedright\arraybackslash {\ttfamily\footnotesize (3,\allowbreak 12,\allowbreak 14)}\par $K=14$} & {\ttfamily\footnotesize (29,\allowbreak 1,\allowbreak 28)\par\smallskip (26,\allowbreak 2,\allowbreak 27,\allowbreak 4,\allowbreak 25,\allowbreak 3,\allowbreak 23,\allowbreak 5,\allowbreak 24,\allowbreak 8,\allowbreak 22,\allowbreak 9,\allowbreak 21,\allowbreak 6)\par\smallskip (17,\allowbreak 10,\allowbreak 20,\allowbreak 11,\allowbreak 19,\allowbreak 14,\allowbreak 15,\allowbreak 13,\allowbreak 16,\allowbreak 12,\allowbreak 18,\allowbreak 7)}\tabularnewline[5pt]
{\raggedright\arraybackslash {\ttfamily\footnotesize (3,\allowbreak 12,\allowbreak 29)}\par $K=21$} & {\ttfamily\footnotesize (44,\allowbreak 1,\allowbreak 43)\par\smallskip (41,\allowbreak 2,\allowbreak 42,\allowbreak 4,\allowbreak 40,\allowbreak 3,\allowbreak 38,\allowbreak 5,\allowbreak 39,\allowbreak 7,\allowbreak 37,\allowbreak 6,\allowbreak 35,\allowbreak 8,\allowbreak 36,\allowbreak 10,\allowbreak 34,\allowbreak 11,\allowbreak 33,\allowbreak 12,\allowbreak 31,\allowbreak 13,\allowbreak 30,\allowbreak 14,\allowbreak 28,\allowbreak 15,\allowbreak 27,\allowbreak 17,\allowbreak 32)\par\smallskip (25,\allowbreak 16,\allowbreak 24,\allowbreak 19,\allowbreak 26,\allowbreak 20,\allowbreak 23,\allowbreak 21,\allowbreak 22,\allowbreak 18,\allowbreak 29,\allowbreak 9)}\tabularnewline[5pt]
{\raggedright\arraybackslash {\ttfamily\footnotesize (3,\allowbreak 13,\allowbreak 13)}\par $K=13$} & {\ttfamily\footnotesize (29,\allowbreak 1,\allowbreak 28)\par\smallskip (26,\allowbreak 2,\allowbreak 27,\allowbreak 4,\allowbreak 25,\allowbreak 3,\allowbreak 23,\allowbreak 5,\allowbreak 24,\allowbreak 7,\allowbreak 20,\allowbreak 6,\allowbreak 22)\par\smallskip (15,\allowbreak 11,\allowbreak 18,\allowbreak 8,\allowbreak 19,\allowbreak 10,\allowbreak 16,\allowbreak 13,\allowbreak 14,\allowbreak 12,\allowbreak 17,\allowbreak 9,\allowbreak 21)}\tabularnewline[5pt]
{\raggedright\arraybackslash {\ttfamily\footnotesize (3,\allowbreak 13,\allowbreak 14)}\par $K=14$} & {\ttfamily\footnotesize (30,\allowbreak 1,\allowbreak 29)\par\smallskip (27,\allowbreak 2,\allowbreak 28,\allowbreak 4,\allowbreak 26,\allowbreak 3,\allowbreak 24,\allowbreak 5,\allowbreak 25,\allowbreak 7,\allowbreak 23,\allowbreak 8,\allowbreak 20,\allowbreak 6)\par\smallskip (17,\allowbreak 10,\allowbreak 21,\allowbreak 11,\allowbreak 19,\allowbreak 14,\allowbreak 15,\allowbreak 13,\allowbreak 16,\allowbreak 12,\allowbreak 18,\allowbreak 9,\allowbreak 22)}\tabularnewline[5pt]
{\raggedright\arraybackslash {\ttfamily\footnotesize (3,\allowbreak 13,\allowbreak 29)}\par $K=21$} & {\ttfamily\footnotesize (45,\allowbreak 1,\allowbreak 44)\par\smallskip (42,\allowbreak 2,\allowbreak 43,\allowbreak 4,\allowbreak 41,\allowbreak 3,\allowbreak 39,\allowbreak 5,\allowbreak 40,\allowbreak 7,\allowbreak 38,\allowbreak 6,\allowbreak 36,\allowbreak 8,\allowbreak 37,\allowbreak 10,\allowbreak 35,\allowbreak 9,\allowbreak 32,\allowbreak 11,\allowbreak 31,\allowbreak 13,\allowbreak 30,\allowbreak 15,\allowbreak 29,\allowbreak 16,\allowbreak 28,\allowbreak 12,\allowbreak 34)\par\smallskip (24,\allowbreak 17,\allowbreak 27,\allowbreak 18,\allowbreak 26,\allowbreak 21,\allowbreak 22,\allowbreak 20,\allowbreak 23,\allowbreak 19,\allowbreak 25,\allowbreak 14,\allowbreak 33)}\tabularnewline[5pt]
{\raggedright\arraybackslash {\ttfamily\footnotesize (3,\allowbreak 14,\allowbreak 14)}\par $K=15$} & {\ttfamily\footnotesize (31,\allowbreak 1,\allowbreak 30)\par\smallskip (28,\allowbreak 2,\allowbreak 29,\allowbreak 4,\allowbreak 27,\allowbreak 3,\allowbreak 25,\allowbreak 5,\allowbreak 26,\allowbreak 7,\allowbreak 24,\allowbreak 10,\allowbreak 23,\allowbreak 8)\par\smallskip (18,\allowbreak 11,\allowbreak 21,\allowbreak 9,\allowbreak 20,\allowbreak 12,\allowbreak 17,\allowbreak 14,\allowbreak 16,\allowbreak 15,\allowbreak 19,\allowbreak 13,\allowbreak 22,\allowbreak 6)}\tabularnewline[5pt]
{\raggedright\arraybackslash {\ttfamily\footnotesize (3,\allowbreak 14,\allowbreak 15)}\par $K=15$} & {\ttfamily\footnotesize (32,\allowbreak 1,\allowbreak 31)\par\smallskip (29,\allowbreak 2,\allowbreak 30,\allowbreak 4,\allowbreak 28,\allowbreak 3,\allowbreak 26,\allowbreak 5,\allowbreak 27,\allowbreak 7,\allowbreak 25,\allowbreak 8,\allowbreak 24,\allowbreak 10,\allowbreak 23)\par\smallskip (19,\allowbreak 11,\allowbreak 22,\allowbreak 12,\allowbreak 17,\allowbreak 13,\allowbreak 20,\allowbreak 14,\allowbreak 16,\allowbreak 15,\allowbreak 18,\allowbreak 9,\allowbreak 21,\allowbreak 6)}\tabularnewline[5pt]
{\raggedright\arraybackslash {\ttfamily\footnotesize (3,\allowbreak 14,\allowbreak 16)}\par $K=16$} & {\ttfamily\footnotesize (33,\allowbreak 1,\allowbreak 32)\par\smallskip (30,\allowbreak 2,\allowbreak 31,\allowbreak 4,\allowbreak 29,\allowbreak 3,\allowbreak 27,\allowbreak 5,\allowbreak 28,\allowbreak 7,\allowbreak 26,\allowbreak 8,\allowbreak 25,\allowbreak 10,\allowbreak 23,\allowbreak 9)\par\smallskip (20,\allowbreak 11,\allowbreak 21,\allowbreak 13,\allowbreak 24,\allowbreak 12,\allowbreak 19,\allowbreak 14,\allowbreak 18,\allowbreak 15,\allowbreak 17,\allowbreak 16,\allowbreak 22,\allowbreak 6)}\tabularnewline[5pt]
{\raggedright\arraybackslash {\ttfamily\footnotesize (3,\allowbreak 14,\allowbreak 33)}\par $K=24$} & {\ttfamily\footnotesize (50,\allowbreak 1,\allowbreak 49)\par\smallskip (47,\allowbreak 2,\allowbreak 48,\allowbreak 4,\allowbreak 46,\allowbreak 3,\allowbreak 44,\allowbreak 5,\allowbreak 45,\allowbreak 7,\allowbreak 43,\allowbreak 6,\allowbreak 41,\allowbreak 8,\allowbreak 42,\allowbreak 10,\allowbreak 40,\allowbreak 9,\allowbreak 38,\allowbreak 11,\allowbreak 37,\allowbreak 12,\allowbreak 36,\allowbreak 14,\allowbreak 35,\allowbreak 15,\allowbreak 34,\allowbreak 17,\allowbreak 33,\allowbreak 18,\allowbreak 30,\allowbreak 16,\allowbreak 39)\par\smallskip (28,\allowbreak 19,\allowbreak 32,\allowbreak 22,\allowbreak 29,\allowbreak 21,\allowbreak 27,\allowbreak 23,\allowbreak 26,\allowbreak 24,\allowbreak 25,\allowbreak 20,\allowbreak 31,\allowbreak 13)}\tabularnewline[5pt]
{\raggedright\arraybackslash {\ttfamily\footnotesize (3,\allowbreak 15,\allowbreak 15)}\par $K=15$} & {\ttfamily\footnotesize (33,\allowbreak 1,\allowbreak 32)\par\smallskip (30,\allowbreak 2,\allowbreak 31,\allowbreak 4,\allowbreak 29,\allowbreak 3,\allowbreak 27,\allowbreak 5,\allowbreak 28,\allowbreak 7,\allowbreak 26,\allowbreak 6,\allowbreak 23,\allowbreak 8,\allowbreak 24)\par\smallskip (18,\allowbreak 10,\allowbreak 22,\allowbreak 9,\allowbreak 20,\allowbreak 13,\allowbreak 19,\allowbreak 14,\allowbreak 17,\allowbreak 15,\allowbreak 16,\allowbreak 12,\allowbreak 21,\allowbreak 11,\allowbreak 25)}\tabularnewline[5pt]
{\raggedright\arraybackslash {\ttfamily\footnotesize (3,\allowbreak 15,\allowbreak 16)}\par $K=16$} & {\ttfamily\footnotesize (34,\allowbreak 1,\allowbreak 33)\par\smallskip (31,\allowbreak 2,\allowbreak 32,\allowbreak 4,\allowbreak 30,\allowbreak 3,\allowbreak 28,\allowbreak 5,\allowbreak 29,\allowbreak 7,\allowbreak 27,\allowbreak 6,\allowbreak 25,\allowbreak 10,\allowbreak 24,\allowbreak 8)\par\smallskip (18,\allowbreak 14,\allowbreak 21,\allowbreak 11,\allowbreak 23,\allowbreak 12,\allowbreak 20,\allowbreak 15,\allowbreak 17,\allowbreak 16,\allowbreak 19,\allowbreak 13,\allowbreak 22,\allowbreak 9,\allowbreak 26)}\tabularnewline[5pt]
{\raggedright\arraybackslash {\ttfamily\footnotesize (3,\allowbreak 15,\allowbreak 33)}\par $K=24$} & {\ttfamily\footnotesize (51,\allowbreak 1,\allowbreak 50)\par\smallskip (48,\allowbreak 2,\allowbreak 49,\allowbreak 4,\allowbreak 47,\allowbreak 3,\allowbreak 45,\allowbreak 5,\allowbreak 46,\allowbreak 7,\allowbreak 44,\allowbreak 6,\allowbreak 42,\allowbreak 8,\allowbreak 43,\allowbreak 10,\allowbreak 41,\allowbreak 9,\allowbreak 39,\allowbreak 11,\allowbreak 38,\allowbreak 12,\allowbreak 37,\allowbreak 13,\allowbreak 35,\allowbreak 14,\allowbreak 34,\allowbreak 15,\allowbreak 33,\allowbreak 16,\allowbreak 32,\allowbreak 17,\allowbreak 40)\par\smallskip (29,\allowbreak 18,\allowbreak 31,\allowbreak 19,\allowbreak 28,\allowbreak 20,\allowbreak 30,\allowbreak 23,\allowbreak 27,\allowbreak 21,\allowbreak 26,\allowbreak 24,\allowbreak 25,\allowbreak 22,\allowbreak 36)}\tabularnewline[5pt]
{\raggedright\arraybackslash {\ttfamily\footnotesize (3,\allowbreak 16,\allowbreak 16)}\par $K=17$} & {\ttfamily\footnotesize (35,\allowbreak 1,\allowbreak 34)\par\smallskip (32,\allowbreak 2,\allowbreak 33,\allowbreak 4,\allowbreak 31,\allowbreak 3,\allowbreak 29,\allowbreak 5,\allowbreak 30,\allowbreak 7,\allowbreak 28,\allowbreak 8,\allowbreak 27,\allowbreak 9,\allowbreak 26,\allowbreak 10)\par\smallskip (22,\allowbreak 11,\allowbreak 25,\allowbreak 12,\allowbreak 24,\allowbreak 14,\allowbreak 23,\allowbreak 15,\allowbreak 20,\allowbreak 13,\allowbreak 19,\allowbreak 16,\allowbreak 18,\allowbreak 17,\allowbreak 21,\allowbreak 6)}\tabularnewline[5pt]
{\raggedright\arraybackslash {\ttfamily\footnotesize (3,\allowbreak 16,\allowbreak 17)}\par $K=17$} & {\ttfamily\footnotesize (36,\allowbreak 1,\allowbreak 35)\par\smallskip (33,\allowbreak 2,\allowbreak 34,\allowbreak 4,\allowbreak 32,\allowbreak 3,\allowbreak 30,\allowbreak 5,\allowbreak 31,\allowbreak 7,\allowbreak 29,\allowbreak 6,\allowbreak 27,\allowbreak 8,\allowbreak 26,\allowbreak 11,\allowbreak 28)\par\smallskip (20,\allowbreak 12,\allowbreak 24,\allowbreak 10,\allowbreak 23,\allowbreak 13,\allowbreak 22,\allowbreak 15,\allowbreak 21,\allowbreak 16,\allowbreak 19,\allowbreak 17,\allowbreak 18,\allowbreak 14,\allowbreak 25,\allowbreak 9)}\tabularnewline[5pt]
{\raggedright\arraybackslash {\ttfamily\footnotesize (3,\allowbreak 16,\allowbreak 18)}\par $K=18$} & {\ttfamily\footnotesize (37,\allowbreak 1,\allowbreak 36)\par\smallskip (34,\allowbreak 2,\allowbreak 35,\allowbreak 4,\allowbreak 33,\allowbreak 3,\allowbreak 31,\allowbreak 5,\allowbreak 32,\allowbreak 7,\allowbreak 30,\allowbreak 6,\allowbreak 28,\allowbreak 10,\allowbreak 25,\allowbreak 11,\allowbreak 27,\allowbreak 8)\par\smallskip (21,\allowbreak 13,\allowbreak 26,\allowbreak 14,\allowbreak 24,\allowbreak 15,\allowbreak 22,\allowbreak 16,\allowbreak 20,\allowbreak 17,\allowbreak 19,\allowbreak 18,\allowbreak 23,\allowbreak 12,\allowbreak 29,\allowbreak 9)}\tabularnewline[5pt]
{\raggedright\arraybackslash {\ttfamily\footnotesize (3,\allowbreak 16,\allowbreak 37)}\par $K=27$} & {\ttfamily\footnotesize (56,\allowbreak 1,\allowbreak 55)\par\smallskip (53,\allowbreak 2,\allowbreak 54,\allowbreak 4,\allowbreak 52,\allowbreak 3,\allowbreak 50,\allowbreak 5,\allowbreak 51,\allowbreak 7,\allowbreak 49,\allowbreak 6,\allowbreak 47,\allowbreak 8,\allowbreak 48,\allowbreak 10,\allowbreak 46,\allowbreak 9,\allowbreak 44,\allowbreak 11,\allowbreak 45,\allowbreak 13,\allowbreak 43,\allowbreak 14,\allowbreak 42,\allowbreak 15,\allowbreak 40,\allowbreak 16,\allowbreak 39,\allowbreak 17,\allowbreak 37,\allowbreak 18,\allowbreak 36,\allowbreak 19,\allowbreak 35,\allowbreak 20,\allowbreak 41)\par\smallskip (31,\allowbreak 21,\allowbreak 34,\allowbreak 22,\allowbreak 33,\allowbreak 25,\allowbreak 32,\allowbreak 23,\allowbreak 28,\allowbreak 27,\allowbreak 29,\allowbreak 26,\allowbreak 30,\allowbreak 24,\allowbreak 38,\allowbreak 12)}\tabularnewline[5pt]
{\raggedright\arraybackslash {\ttfamily\footnotesize (3,\allowbreak 17,\allowbreak 17)}\par $K=17$} & {\ttfamily\footnotesize (37,\allowbreak 1,\allowbreak 36)\par\smallskip (34,\allowbreak 2,\allowbreak 35,\allowbreak 4,\allowbreak 33,\allowbreak 3,\allowbreak 31,\allowbreak 5,\allowbreak 32,\allowbreak 7,\allowbreak 30,\allowbreak 6,\allowbreak 28,\allowbreak 8,\allowbreak 27,\allowbreak 9,\allowbreak 26)\par\smallskip (21,\allowbreak 10,\allowbreak 25,\allowbreak 11,\allowbreak 24,\allowbreak 12,\allowbreak 22,\allowbreak 14,\allowbreak 23,\allowbreak 16,\allowbreak 20,\allowbreak 15,\allowbreak 18,\allowbreak 17,\allowbreak 19,\allowbreak 13,\allowbreak 29)}\tabularnewline[5pt]
{\raggedright\arraybackslash {\ttfamily\footnotesize (3,\allowbreak 17,\allowbreak 18)}\par $K=18$} & {\ttfamily\footnotesize (38,\allowbreak 1,\allowbreak 37)\par\smallskip (35,\allowbreak 2,\allowbreak 36,\allowbreak 4,\allowbreak 34,\allowbreak 3,\allowbreak 32,\allowbreak 5,\allowbreak 33,\allowbreak 7,\allowbreak 31,\allowbreak 6,\allowbreak 29,\allowbreak 8,\allowbreak 28,\allowbreak 9,\allowbreak 27,\allowbreak 10)\par\smallskip (22,\allowbreak 11,\allowbreak 26,\allowbreak 12,\allowbreak 25,\allowbreak 13,\allowbreak 23,\allowbreak 15,\allowbreak 24,\allowbreak 17,\allowbreak 21,\allowbreak 16,\allowbreak 19,\allowbreak 18,\allowbreak 20,\allowbreak 14,\allowbreak 30)}\tabularnewline[5pt]
{\raggedright\arraybackslash {\ttfamily\footnotesize (3,\allowbreak 17,\allowbreak 37)}\par $K=27$} & {\ttfamily\footnotesize (57,\allowbreak 1,\allowbreak 56)\par\smallskip (54,\allowbreak 2,\allowbreak 55,\allowbreak 4,\allowbreak 53,\allowbreak 3,\allowbreak 51,\allowbreak 5,\allowbreak 52,\allowbreak 7,\allowbreak 50,\allowbreak 6,\allowbreak 48,\allowbreak 8,\allowbreak 47,\allowbreak 9,\allowbreak 46,\allowbreak 26,\allowbreak 29,\allowbreak 27,\allowbreak 28,\allowbreak 24,\allowbreak 31,\allowbreak 25,\allowbreak 30,\allowbreak 22,\allowbreak 33,\allowbreak 23,\allowbreak 32,\allowbreak 20,\allowbreak 35,\allowbreak 21,\allowbreak 34,\allowbreak 18,\allowbreak 37,\allowbreak 19,\allowbreak 36)\par\smallskip (41,\allowbreak 10,\allowbreak 45,\allowbreak 11,\allowbreak 44,\allowbreak 12,\allowbreak 42,\allowbreak 14,\allowbreak 43,\allowbreak 16,\allowbreak 40,\allowbreak 15,\allowbreak 38,\allowbreak 17,\allowbreak 39,\allowbreak 13,\allowbreak 49)}\tabularnewline[5pt]
\end{longtable}
\endgroup

\subsection{The 29 maximum-pinned short-two seeds for minimum three}
\label{aux:min3-pinned}
\begingroup\footnotesize
\begin{longtable}{@{}p{25mm}p{115mm}@{}}
Profile / interface & Complete outward words\\\hline\endhead
{\raggedright\arraybackslash {\ttfamily\footnotesize (2,\allowbreak 18,\allowbreak 18)}\par $K=19$\par pinned $(N,1)$} & {\ttfamily\footnotesize (35,\allowbreak 3,\allowbreak 37,\allowbreak 4,\allowbreak 34,\allowbreak 6,\allowbreak 33,\allowbreak 2,\allowbreak 31,\allowbreak 5,\allowbreak 29,\allowbreak 7,\allowbreak 30,\allowbreak 9,\allowbreak 26,\allowbreak 14,\allowbreak 23,\allowbreak 8)\par\smallskip (36,\allowbreak 11,\allowbreak 27,\allowbreak 13,\allowbreak 32,\allowbreak 12,\allowbreak 25,\allowbreak 15,\allowbreak 22,\allowbreak 19,\allowbreak 20,\allowbreak 18,\allowbreak 24,\allowbreak 16,\allowbreak 21,\allowbreak 17,\allowbreak 28,\allowbreak 10)\par\smallskip (38,\allowbreak 1)}\tabularnewline[5pt]
{\raggedright\arraybackslash {\ttfamily\footnotesize (2,\allowbreak 18,\allowbreak 19)}\par $K=19$\par pinned $(N,1)$\par $2e\le b+4$} & {\ttfamily\footnotesize (35,\allowbreak 9,\allowbreak 30,\allowbreak 10,\allowbreak 23,\allowbreak 15,\allowbreak 25,\allowbreak 2,\allowbreak 38,\allowbreak 4,\allowbreak 34,\allowbreak 6,\allowbreak 33,\allowbreak 11,\allowbreak 26,\allowbreak 14,\allowbreak 28,\allowbreak 12,\allowbreak 29)\par\smallskip (37,\allowbreak 5,\allowbreak 36,\allowbreak 3,\allowbreak 32,\allowbreak 7,\allowbreak 31,\allowbreak 13,\allowbreak 22,\allowbreak 18,\allowbreak 24,\allowbreak 17,\allowbreak 20,\allowbreak 19,\allowbreak 21,\allowbreak 16,\allowbreak 27,\allowbreak 8)\par\smallskip (39,\allowbreak 1)}\tabularnewline[5pt]
{\raggedright\arraybackslash {\ttfamily\footnotesize (2,\allowbreak 18,\allowbreak 20)}\par $K=20$\par pinned $(N,1)$} & {\ttfamily\footnotesize (37,\allowbreak 6,\allowbreak 32,\allowbreak 5,\allowbreak 24,\allowbreak 19,\allowbreak 23,\allowbreak 15,\allowbreak 25,\allowbreak 18,\allowbreak 21,\allowbreak 20,\allowbreak 22,\allowbreak 16,\allowbreak 28,\allowbreak 17,\allowbreak 26,\allowbreak 10,\allowbreak 31,\allowbreak 11)\par\smallskip (38,\allowbreak 4,\allowbreak 39,\allowbreak 3,\allowbreak 33,\allowbreak 9,\allowbreak 34,\allowbreak 2,\allowbreak 35,\allowbreak 13,\allowbreak 27,\allowbreak 14,\allowbreak 29,\allowbreak 12,\allowbreak 30,\allowbreak 7,\allowbreak 36,\allowbreak 8)\par\smallskip (40,\allowbreak 1)}\tabularnewline[5pt]
{\raggedright\arraybackslash {\ttfamily\footnotesize (2,\allowbreak 19,\allowbreak 19)}\par $K=19$\par pinned $(N,1)$} & {\ttfamily\footnotesize (40,\allowbreak 1)\par\smallskip (38,\allowbreak 2,\allowbreak 39,\allowbreak 4,\allowbreak 37,\allowbreak 3,\allowbreak 35,\allowbreak 5,\allowbreak 36,\allowbreak 7,\allowbreak 34,\allowbreak 6,\allowbreak 32,\allowbreak 8,\allowbreak 33,\allowbreak 10,\allowbreak 29,\allowbreak 9,\allowbreak 31)\par\smallskip (21,\allowbreak 17,\allowbreak 24,\allowbreak 14,\allowbreak 27,\allowbreak 11,\allowbreak 28,\allowbreak 13,\allowbreak 25,\allowbreak 16,\allowbreak 22,\allowbreak 19,\allowbreak 20,\allowbreak 18,\allowbreak 23,\allowbreak 15,\allowbreak 26,\allowbreak 12,\allowbreak 30)}\tabularnewline[5pt]
{\raggedright\arraybackslash {\ttfamily\footnotesize (2,\allowbreak 19,\allowbreak 20)}\par $K=20$\par pinned $(N,1)$} & {\ttfamily\footnotesize (41,\allowbreak 1)\par\smallskip (39,\allowbreak 2,\allowbreak 40,\allowbreak 4,\allowbreak 38,\allowbreak 3,\allowbreak 36,\allowbreak 5,\allowbreak 37,\allowbreak 7,\allowbreak 35,\allowbreak 6,\allowbreak 33,\allowbreak 8,\allowbreak 34,\allowbreak 10,\allowbreak 32,\allowbreak 11,\allowbreak 29,\allowbreak 9)\par\smallskip (23,\allowbreak 15,\allowbreak 30,\allowbreak 13,\allowbreak 27,\allowbreak 14,\allowbreak 26,\allowbreak 16,\allowbreak 25,\allowbreak 18,\allowbreak 24,\allowbreak 19,\allowbreak 22,\allowbreak 20,\allowbreak 21,\allowbreak 17,\allowbreak 28,\allowbreak 12,\allowbreak 31)}\tabularnewline[5pt]
{\raggedright\arraybackslash {\ttfamily\footnotesize (2,\allowbreak 20,\allowbreak 21)}\par $K=21$\par pinned $(N,1)$\par $2e\le b+4$} & {\ttfamily\footnotesize (39,\allowbreak 2,\allowbreak 42,\allowbreak 6,\allowbreak 40,\allowbreak 5,\allowbreak 35,\allowbreak 7,\allowbreak 38,\allowbreak 11,\allowbreak 33,\allowbreak 10,\allowbreak 31,\allowbreak 12,\allowbreak 36,\allowbreak 4,\allowbreak 37,\allowbreak 8,\allowbreak 34,\allowbreak 9)\par\smallskip (41,\allowbreak 3,\allowbreak 23,\allowbreak 21,\allowbreak 22,\allowbreak 19,\allowbreak 24,\allowbreak 20,\allowbreak 26,\allowbreak 18,\allowbreak 25,\allowbreak 16,\allowbreak 29,\allowbreak 13,\allowbreak 30,\allowbreak 15,\allowbreak 27,\allowbreak 17,\allowbreak 28,\allowbreak 14,\allowbreak 32)\par\smallskip (43,\allowbreak 1)}\tabularnewline[5pt]
{\raggedright\arraybackslash {\ttfamily\footnotesize (2,\allowbreak 20,\allowbreak 22)}\par $K=22$\par pinned $(N,1)$} & {\ttfamily\footnotesize (32,\allowbreak 13,\allowbreak 36,\allowbreak 14,\allowbreak 35,\allowbreak 10,\allowbreak 28,\allowbreak 15,\allowbreak 31,\allowbreak 11,\allowbreak 38,\allowbreak 5,\allowbreak 41,\allowbreak 4,\allowbreak 39,\allowbreak 8,\allowbreak 37,\allowbreak 7,\allowbreak 33,\allowbreak 9)\par\smallskip (42,\allowbreak 3,\allowbreak 43,\allowbreak 2,\allowbreak 40,\allowbreak 6,\allowbreak 34,\allowbreak 17,\allowbreak 24,\allowbreak 22,\allowbreak 23,\allowbreak 20,\allowbreak 26,\allowbreak 21,\allowbreak 25,\allowbreak 16,\allowbreak 30,\allowbreak 18,\allowbreak 29,\allowbreak 19,\allowbreak 27,\allowbreak 12)\par\smallskip (44,\allowbreak 1)}\tabularnewline[5pt]
{\raggedright\arraybackslash {\ttfamily\footnotesize (2,\allowbreak 21,\allowbreak 21)}\par $K=21$\par pinned $(N,1)$} & {\ttfamily\footnotesize (44,\allowbreak 1)\par\smallskip (42,\allowbreak 2,\allowbreak 43,\allowbreak 4,\allowbreak 41,\allowbreak 3,\allowbreak 39,\allowbreak 5,\allowbreak 40,\allowbreak 7,\allowbreak 38,\allowbreak 6,\allowbreak 36,\allowbreak 8,\allowbreak 37,\allowbreak 10,\allowbreak 35,\allowbreak 9,\allowbreak 32,\allowbreak 11,\allowbreak 33)\par\smallskip (24,\allowbreak 16,\allowbreak 31,\allowbreak 12,\allowbreak 30,\allowbreak 13,\allowbreak 29,\allowbreak 15,\allowbreak 28,\allowbreak 17,\allowbreak 27,\allowbreak 18,\allowbreak 25,\allowbreak 19,\allowbreak 23,\allowbreak 20,\allowbreak 22,\allowbreak 21,\allowbreak 26,\allowbreak 14,\allowbreak 34)}\tabularnewline[5pt]
{\raggedright\arraybackslash {\ttfamily\footnotesize (2,\allowbreak 21,\allowbreak 22)}\par $K=22$\par pinned $(N,1)$} & {\ttfamily\footnotesize (31,\allowbreak 20,\allowbreak 24,\allowbreak 19,\allowbreak 25,\allowbreak 22,\allowbreak 23,\allowbreak 21,\allowbreak 28,\allowbreak 16,\allowbreak 33,\allowbreak 11,\allowbreak 30,\allowbreak 14,\allowbreak 32,\allowbreak 17,\allowbreak 27,\allowbreak 18,\allowbreak 26,\allowbreak 13,\allowbreak 36)\par\smallskip (43,\allowbreak 3,\allowbreak 44,\allowbreak 2,\allowbreak 41,\allowbreak 6,\allowbreak 38,\allowbreak 10,\allowbreak 34,\allowbreak 7,\allowbreak 40,\allowbreak 4,\allowbreak 42,\allowbreak 5,\allowbreak 39,\allowbreak 9,\allowbreak 35,\allowbreak 15,\allowbreak 29,\allowbreak 8,\allowbreak 37,\allowbreak 12)\par\smallskip (45,\allowbreak 1)}\tabularnewline[5pt]
{\raggedright\arraybackslash {\ttfamily\footnotesize (2,\allowbreak 22,\allowbreak 23)}\par $K=23$\par pinned $(N,1)$\par $2e\le b+4$} & {\ttfamily\footnotesize (43,\allowbreak 9,\allowbreak 42,\allowbreak 3,\allowbreak 41,\allowbreak 6,\allowbreak 46,\allowbreak 2,\allowbreak 44,\allowbreak 7,\allowbreak 33,\allowbreak 16,\allowbreak 27,\allowbreak 20,\allowbreak 25,\allowbreak 23,\allowbreak 24,\allowbreak 21,\allowbreak 31,\allowbreak 13,\allowbreak 32,\allowbreak 10)\par\smallskip (45,\allowbreak 4,\allowbreak 40,\allowbreak 8,\allowbreak 38,\allowbreak 18,\allowbreak 26,\allowbreak 22,\allowbreak 28,\allowbreak 19,\allowbreak 34,\allowbreak 5,\allowbreak 36,\allowbreak 15,\allowbreak 29,\allowbreak 17,\allowbreak 30,\allowbreak 14,\allowbreak 37,\allowbreak 12,\allowbreak 39,\allowbreak 11,\allowbreak 35)\par\smallskip (47,\allowbreak 1)}\tabularnewline[5pt]
{\raggedright\arraybackslash {\ttfamily\footnotesize (2,\allowbreak 22,\allowbreak 24)}\par $K=24$\par pinned $(N,1)$} & {\ttfamily\footnotesize (40,\allowbreak 11,\allowbreak 30,\allowbreak 20,\allowbreak 28,\allowbreak 19,\allowbreak 34,\allowbreak 17,\allowbreak 31,\allowbreak 18,\allowbreak 29,\allowbreak 22,\allowbreak 27,\allowbreak 21,\allowbreak 25,\allowbreak 24,\allowbreak 26,\allowbreak 23,\allowbreak 35,\allowbreak 12,\allowbreak 37,\allowbreak 10)\par\smallskip (46,\allowbreak 3,\allowbreak 47,\allowbreak 2,\allowbreak 44,\allowbreak 8,\allowbreak 41,\allowbreak 6,\allowbreak 45,\allowbreak 4,\allowbreak 32,\allowbreak 16,\allowbreak 42,\allowbreak 5,\allowbreak 43,\allowbreak 9,\allowbreak 39,\allowbreak 7,\allowbreak 38,\allowbreak 14,\allowbreak 36,\allowbreak 15,\allowbreak 33,\allowbreak 13)\par\smallskip (48,\allowbreak 1)}\tabularnewline[5pt]
{\raggedright\arraybackslash {\ttfamily\footnotesize (2,\allowbreak 24,\allowbreak 25)}\par $K=25$\par pinned $(N,1)$\par $2e\le b+4$} & {\ttfamily\footnotesize (46,\allowbreak 3,\allowbreak 44,\allowbreak 6,\allowbreak 48,\allowbreak 4,\allowbreak 36,\allowbreak 20,\allowbreak 30,\allowbreak 24,\allowbreak 28,\allowbreak 21,\allowbreak 29,\allowbreak 10,\allowbreak 43,\allowbreak 13,\allowbreak 40,\allowbreak 9,\allowbreak 45,\allowbreak 8,\allowbreak 47,\allowbreak 7,\allowbreak 42,\allowbreak 14,\allowbreak 38)\par\smallskip (49,\allowbreak 2,\allowbreak 50,\allowbreak 5,\allowbreak 39,\allowbreak 17,\allowbreak 32,\allowbreak 23,\allowbreak 26,\allowbreak 25,\allowbreak 27,\allowbreak 22,\allowbreak 33,\allowbreak 19,\allowbreak 31,\allowbreak 18,\allowbreak 35,\allowbreak 15,\allowbreak 41,\allowbreak 12,\allowbreak 37,\allowbreak 16,\allowbreak 34,\allowbreak 11)\par\smallskip (51,\allowbreak 1)}\tabularnewline[5pt]
{\raggedright\arraybackslash {\ttfamily\footnotesize (2,\allowbreak 24,\allowbreak 26)}\par $K=26$\par pinned $(N,1)$} & {\ttfamily\footnotesize (48,\allowbreak 10,\allowbreak 45,\allowbreak 3,\allowbreak 44,\allowbreak 18,\allowbreak 34,\allowbreak 19,\allowbreak 33,\allowbreak 20,\allowbreak 32,\allowbreak 21,\allowbreak 31,\allowbreak 22,\allowbreak 30,\allowbreak 23,\allowbreak 29,\allowbreak 24,\allowbreak 28,\allowbreak 25,\allowbreak 27,\allowbreak 26,\allowbreak 43,\allowbreak 11)\par\smallskip (50,\allowbreak 6,\allowbreak 46,\allowbreak 9,\allowbreak 42,\allowbreak 12,\allowbreak 40,\allowbreak 13,\allowbreak 38,\allowbreak 7,\allowbreak 41,\allowbreak 5,\allowbreak 51,\allowbreak 2,\allowbreak 49,\allowbreak 4,\allowbreak 47,\allowbreak 8,\allowbreak 37,\allowbreak 15,\allowbreak 39,\allowbreak 16,\allowbreak 36,\allowbreak 17,\allowbreak 35,\allowbreak 14)\par\smallskip (52,\allowbreak 1)}\tabularnewline[5pt]
{\raggedright\arraybackslash {\ttfamily\footnotesize (2,\allowbreak 25,\allowbreak 25)}\par $K=25$\par pinned $(N,1)$} & {\ttfamily\footnotesize (49,\allowbreak 4,\allowbreak 51,\allowbreak 20,\allowbreak 34,\allowbreak 16,\allowbreak 36,\allowbreak 12,\allowbreak 33,\allowbreak 22,\allowbreak 30,\allowbreak 21,\allowbreak 28,\allowbreak 25,\allowbreak 26,\allowbreak 24,\allowbreak 29,\allowbreak 23,\allowbreak 27,\allowbreak 17,\allowbreak 32,\allowbreak 19,\allowbreak 31,\allowbreak 15,\allowbreak 38)\par\smallskip (50,\allowbreak 2,\allowbreak 48,\allowbreak 6,\allowbreak 43,\allowbreak 9,\allowbreak 37,\allowbreak 18,\allowbreak 35,\allowbreak 13,\allowbreak 40,\allowbreak 14,\allowbreak 39,\allowbreak 10,\allowbreak 42,\allowbreak 7,\allowbreak 45,\allowbreak 5,\allowbreak 46,\allowbreak 3,\allowbreak 47,\allowbreak 8,\allowbreak 44,\allowbreak 11,\allowbreak 41)\par\smallskip (52,\allowbreak 1)}\tabularnewline[5pt]
{\raggedright\arraybackslash {\ttfamily\footnotesize (2,\allowbreak 25,\allowbreak 26)}\par $K=26$\par pinned $(N,1)$} & {\ttfamily\footnotesize (50,\allowbreak 2,\allowbreak 45,\allowbreak 4,\allowbreak 48,\allowbreak 9,\allowbreak 39,\allowbreak 10,\allowbreak 41,\allowbreak 16,\allowbreak 38,\allowbreak 18,\allowbreak 30,\allowbreak 11,\allowbreak 49,\allowbreak 3,\allowbreak 52,\allowbreak 5,\allowbreak 47,\allowbreak 15,\allowbreak 36,\allowbreak 20,\allowbreak 35,\allowbreak 12,\allowbreak 40,\allowbreak 14)\par\smallskip (51,\allowbreak 6,\allowbreak 46,\allowbreak 13,\allowbreak 37,\allowbreak 19,\allowbreak 33,\allowbreak 23,\allowbreak 29,\allowbreak 22,\allowbreak 31,\allowbreak 26,\allowbreak 27,\allowbreak 25,\allowbreak 28,\allowbreak 24,\allowbreak 32,\allowbreak 21,\allowbreak 34,\allowbreak 17,\allowbreak 44,\allowbreak 7,\allowbreak 43,\allowbreak 8,\allowbreak 42)\par\smallskip (53,\allowbreak 1)}\tabularnewline[5pt]
{\raggedright\arraybackslash {\ttfamily\footnotesize (2,\allowbreak 26,\allowbreak 27)}\par $K=27$\par pinned $(N,1)$\par $2e\le b+4$} & {\ttfamily\footnotesize (49,\allowbreak 5,\allowbreak 46,\allowbreak 14,\allowbreak 42,\allowbreak 18,\allowbreak 36,\allowbreak 20,\allowbreak 37,\allowbreak 8,\allowbreak 51,\allowbreak 6,\allowbreak 48,\allowbreak 11,\allowbreak 44,\allowbreak 10,\allowbreak 50,\allowbreak 3,\allowbreak 54,\allowbreak 2,\allowbreak 52,\allowbreak 4,\allowbreak 43,\allowbreak 13,\allowbreak 40,\allowbreak 15,\allowbreak 41)\par\smallskip (53,\allowbreak 7,\allowbreak 45,\allowbreak 9,\allowbreak 29,\allowbreak 27,\allowbreak 28,\allowbreak 25,\allowbreak 32,\allowbreak 22,\allowbreak 35,\allowbreak 21,\allowbreak 33,\allowbreak 24,\allowbreak 30,\allowbreak 26,\allowbreak 31,\allowbreak 23,\allowbreak 34,\allowbreak 19,\allowbreak 38,\allowbreak 17,\allowbreak 39,\allowbreak 16,\allowbreak 47,\allowbreak 12)\par\smallskip (55,\allowbreak 1)}\tabularnewline[5pt]
{\raggedright\arraybackslash {\ttfamily\footnotesize (2,\allowbreak 26,\allowbreak 28)}\par $K=28$\par pinned $(N,1)$} & {\ttfamily\footnotesize (52,\allowbreak 5,\allowbreak 43,\allowbreak 25,\allowbreak 33,\allowbreak 28,\allowbreak 29,\allowbreak 27,\allowbreak 30,\allowbreak 26,\allowbreak 32,\allowbreak 23,\allowbreak 34,\allowbreak 22,\allowbreak 37,\allowbreak 24,\allowbreak 31,\allowbreak 21,\allowbreak 35,\allowbreak 16,\allowbreak 41,\allowbreak 11,\allowbreak 45,\allowbreak 14,\allowbreak 47,\allowbreak 12)\par\smallskip (54,\allowbreak 6,\allowbreak 46,\allowbreak 9,\allowbreak 48,\allowbreak 7,\allowbreak 51,\allowbreak 2,\allowbreak 55,\allowbreak 4,\allowbreak 49,\allowbreak 3,\allowbreak 53,\allowbreak 10,\allowbreak 38,\allowbreak 17,\allowbreak 44,\allowbreak 8,\allowbreak 50,\allowbreak 18,\allowbreak 40,\allowbreak 20,\allowbreak 36,\allowbreak 19,\allowbreak 42,\allowbreak 13,\allowbreak 39,\allowbreak 15)\par\smallskip (56,\allowbreak 1)}\tabularnewline[5pt]
{\raggedright\arraybackslash {\ttfamily\footnotesize (2,\allowbreak 27,\allowbreak 27)}\par $K=27$\par pinned $(N,1)$} & {\ttfamily\footnotesize (53,\allowbreak 2,\allowbreak 51,\allowbreak 6,\allowbreak 52,\allowbreak 5,\allowbreak 47,\allowbreak 8,\allowbreak 44,\allowbreak 11,\allowbreak 34,\allowbreak 21,\allowbreak 29,\allowbreak 25,\allowbreak 32,\allowbreak 20,\allowbreak 35,\allowbreak 24,\allowbreak 30,\allowbreak 27,\allowbreak 28,\allowbreak 26,\allowbreak 31,\allowbreak 22,\allowbreak 38,\allowbreak 18,\allowbreak 42)\par\smallskip (54,\allowbreak 4,\allowbreak 48,\allowbreak 14,\allowbreak 49,\allowbreak 12,\allowbreak 33,\allowbreak 23,\allowbreak 37,\allowbreak 19,\allowbreak 36,\allowbreak 17,\allowbreak 39,\allowbreak 13,\allowbreak 41,\allowbreak 9,\allowbreak 50,\allowbreak 10,\allowbreak 40,\allowbreak 15,\allowbreak 46,\allowbreak 3,\allowbreak 55,\allowbreak 7,\allowbreak 45,\allowbreak 16,\allowbreak 43)\par\smallskip (56,\allowbreak 1)}\tabularnewline[5pt]
{\raggedright\arraybackslash {\ttfamily\footnotesize (2,\allowbreak 27,\allowbreak 28)}\par $K=28$\par pinned $(N,1)$} & {\ttfamily\footnotesize (51,\allowbreak 9,\allowbreak 54,\allowbreak 2,\allowbreak 56,\allowbreak 3,\allowbreak 50,\allowbreak 6,\allowbreak 52,\allowbreak 4,\allowbreak 53,\allowbreak 16,\allowbreak 39,\allowbreak 19,\allowbreak 33,\allowbreak 26,\allowbreak 34,\allowbreak 24,\allowbreak 49,\allowbreak 8,\allowbreak 47,\allowbreak 7,\allowbreak 43,\allowbreak 11,\allowbreak 46,\allowbreak 12,\allowbreak 45)\par\smallskip (55,\allowbreak 5,\allowbreak 48,\allowbreak 10,\allowbreak 40,\allowbreak 18,\allowbreak 36,\allowbreak 17,\allowbreak 44,\allowbreak 13,\allowbreak 42,\allowbreak 14,\allowbreak 38,\allowbreak 23,\allowbreak 35,\allowbreak 22,\allowbreak 31,\allowbreak 25,\allowbreak 29,\allowbreak 28,\allowbreak 30,\allowbreak 27,\allowbreak 32,\allowbreak 21,\allowbreak 37,\allowbreak 20,\allowbreak 41,\allowbreak 15)\par\smallskip (57,\allowbreak 1)}\tabularnewline[5pt]
{\raggedright\arraybackslash {\ttfamily\footnotesize (2,\allowbreak 28,\allowbreak 29)}\par $K=29$\par pinned $(N,1)$\par $2e\le b+4$} & {\ttfamily\footnotesize (55,\allowbreak 7,\allowbreak 49,\allowbreak 8,\allowbreak 52,\allowbreak 6,\allowbreak 56,\allowbreak 3,\allowbreak 46,\allowbreak 15,\allowbreak 40,\allowbreak 22,\allowbreak 35,\allowbreak 23,\allowbreak 32,\allowbreak 28,\allowbreak 34,\allowbreak 24,\allowbreak 38,\allowbreak 19,\allowbreak 39,\allowbreak 17,\allowbreak 50,\allowbreak 10,\allowbreak 45,\allowbreak 16,\allowbreak 43,\allowbreak 13)\par\smallskip (57,\allowbreak 5,\allowbreak 54,\allowbreak 9,\allowbreak 47,\allowbreak 11,\allowbreak 48,\allowbreak 14,\allowbreak 42,\allowbreak 18,\allowbreak 41,\allowbreak 20,\allowbreak 37,\allowbreak 21,\allowbreak 36,\allowbreak 25,\allowbreak 33,\allowbreak 26,\allowbreak 31,\allowbreak 29,\allowbreak 30,\allowbreak 27,\allowbreak 53,\allowbreak 2,\allowbreak 58,\allowbreak 4,\allowbreak 51,\allowbreak 12,\allowbreak 44)\par\smallskip (59,\allowbreak 1)}\tabularnewline[5pt]
{\raggedright\arraybackslash {\ttfamily\footnotesize (2,\allowbreak 28,\allowbreak 30)}\par $K=30$\par pinned $(N,1)$} & {\ttfamily\footnotesize (48,\allowbreak 12,\allowbreak 47,\allowbreak 14,\allowbreak 46,\allowbreak 17,\allowbreak 43,\allowbreak 18,\allowbreak 45,\allowbreak 15,\allowbreak 32,\allowbreak 29,\allowbreak 31,\allowbreak 30,\allowbreak 34,\allowbreak 27,\allowbreak 33,\allowbreak 28,\allowbreak 36,\allowbreak 25,\allowbreak 35,\allowbreak 26,\allowbreak 38,\allowbreak 22,\allowbreak 42,\allowbreak 19,\allowbreak 41,\allowbreak 20,\allowbreak 44,\allowbreak 16)\par\smallskip (58,\allowbreak 2,\allowbreak 59,\allowbreak 4,\allowbreak 57,\allowbreak 3,\allowbreak 55,\allowbreak 5,\allowbreak 56,\allowbreak 7,\allowbreak 54,\allowbreak 23,\allowbreak 37,\allowbreak 24,\allowbreak 39,\allowbreak 21,\allowbreak 40,\allowbreak 6,\allowbreak 52,\allowbreak 10,\allowbreak 51,\allowbreak 8,\allowbreak 53,\allowbreak 9,\allowbreak 49,\allowbreak 11,\allowbreak 50,\allowbreak 13)\par\smallskip (60,\allowbreak 1)}\tabularnewline[5pt]
{\raggedright\arraybackslash {\ttfamily\footnotesize (2,\allowbreak 29,\allowbreak 29)}\par $K=29$\par pinned $(N,1)$} & {\ttfamily\footnotesize (55,\allowbreak 4,\allowbreak 57,\allowbreak 15,\allowbreak 51,\allowbreak 12,\allowbreak 42,\allowbreak 24,\allowbreak 36,\allowbreak 28,\allowbreak 30,\allowbreak 29,\allowbreak 32,\allowbreak 26,\allowbreak 31,\allowbreak 27,\allowbreak 34,\allowbreak 25,\allowbreak 38,\allowbreak 11,\allowbreak 48,\allowbreak 13,\allowbreak 46,\allowbreak 14,\allowbreak 35,\allowbreak 21,\allowbreak 40,\allowbreak 16,\allowbreak 47)\par\smallskip (58,\allowbreak 6,\allowbreak 56,\allowbreak 2,\allowbreak 59,\allowbreak 3,\allowbreak 49,\allowbreak 20,\allowbreak 37,\allowbreak 22,\allowbreak 33,\allowbreak 23,\allowbreak 39,\allowbreak 19,\allowbreak 41,\allowbreak 18,\allowbreak 43,\allowbreak 17,\allowbreak 45,\allowbreak 7,\allowbreak 54,\allowbreak 5,\allowbreak 53,\allowbreak 8,\allowbreak 52,\allowbreak 9,\allowbreak 50,\allowbreak 10,\allowbreak 44)\par\smallskip (60,\allowbreak 1)}\tabularnewline[5pt]
{\raggedright\arraybackslash {\ttfamily\footnotesize (2,\allowbreak 30,\allowbreak 31)}\par $K=31$\par pinned $(N,1)$\par $2e\le b+4$} & {\ttfamily\footnotesize (60,\allowbreak 2,\allowbreak 59,\allowbreak 4,\allowbreak 56,\allowbreak 9,\allowbreak 53,\allowbreak 17,\allowbreak 44,\allowbreak 18,\allowbreak 43,\allowbreak 22,\allowbreak 40,\allowbreak 20,\allowbreak 49,\allowbreak 15,\allowbreak 55,\allowbreak 12,\allowbreak 51,\allowbreak 13,\allowbreak 50,\allowbreak 28,\allowbreak 34,\allowbreak 27,\allowbreak 35,\allowbreak 25,\allowbreak 39,\allowbreak 23,\allowbreak 42,\allowbreak 19,\allowbreak 47)\par\smallskip (61,\allowbreak 7,\allowbreak 57,\allowbreak 6,\allowbreak 62,\allowbreak 3,\allowbreak 48,\allowbreak 16,\allowbreak 46,\allowbreak 11,\allowbreak 52,\allowbreak 10,\allowbreak 58,\allowbreak 5,\allowbreak 54,\allowbreak 8,\allowbreak 41,\allowbreak 24,\allowbreak 37,\allowbreak 26,\allowbreak 38,\allowbreak 29,\allowbreak 33,\allowbreak 30,\allowbreak 32,\allowbreak 31,\allowbreak 36,\allowbreak 21,\allowbreak 45,\allowbreak 14)\par\smallskip (63,\allowbreak 1)}\tabularnewline[5pt]
{\raggedright\arraybackslash {\ttfamily\footnotesize (2,\allowbreak 32,\allowbreak 33)}\par $K=33$\par pinned $(N,1)$\par $2e\le b+4$} & {\ttfamily\footnotesize (64,\allowbreak 2,\allowbreak 59,\allowbreak 13,\allowbreak 48,\allowbreak 21,\allowbreak 43,\allowbreak 27,\allowbreak 38,\allowbreak 26,\allowbreak 40,\allowbreak 25,\allowbreak 42,\allowbreak 23,\allowbreak 41,\allowbreak 28,\allowbreak 37,\allowbreak 33,\allowbreak 34,\allowbreak 32,\allowbreak 35,\allowbreak 30,\allowbreak 36,\allowbreak 29,\allowbreak 39,\allowbreak 31,\allowbreak 56,\allowbreak 11,\allowbreak 53,\allowbreak 19,\allowbreak 47,\allowbreak 15)\par\smallskip (65,\allowbreak 5,\allowbreak 66,\allowbreak 3,\allowbreak 58,\allowbreak 7,\allowbreak 54,\allowbreak 18,\allowbreak 49,\allowbreak 16,\allowbreak 55,\allowbreak 6,\allowbreak 62,\allowbreak 4,\allowbreak 63,\allowbreak 9,\allowbreak 57,\allowbreak 14,\allowbreak 51,\allowbreak 10,\allowbreak 60,\allowbreak 8,\allowbreak 61,\allowbreak 17,\allowbreak 46,\allowbreak 20,\allowbreak 44,\allowbreak 24,\allowbreak 45,\allowbreak 22,\allowbreak 52,\allowbreak 12,\allowbreak 50)\par\smallskip (67,\allowbreak 1)}\tabularnewline[5pt]
{\raggedright\arraybackslash {\ttfamily\footnotesize (2,\allowbreak 32,\allowbreak 34)}\par $K=34$\par pinned $(N,1)$} & {\ttfamily\footnotesize (54,\allowbreak 14,\allowbreak 53,\allowbreak 16,\allowbreak 36,\allowbreak 33,\allowbreak 35,\allowbreak 34,\allowbreak 38,\allowbreak 31,\allowbreak 37,\allowbreak 32,\allowbreak 40,\allowbreak 30,\allowbreak 39,\allowbreak 28,\allowbreak 41,\allowbreak 29,\allowbreak 43,\allowbreak 26,\allowbreak 42,\allowbreak 27,\allowbreak 45,\allowbreak 23,\allowbreak 46,\allowbreak 25,\allowbreak 44,\allowbreak 8,\allowbreak 59,\allowbreak 11,\allowbreak 56,\allowbreak 15)\par\smallskip (66,\allowbreak 2,\allowbreak 67,\allowbreak 4,\allowbreak 65,\allowbreak 3,\allowbreak 63,\allowbreak 5,\allowbreak 64,\allowbreak 7,\allowbreak 60,\allowbreak 10,\allowbreak 57,\allowbreak 13,\allowbreak 55,\allowbreak 12,\allowbreak 58,\allowbreak 9,\allowbreak 61,\allowbreak 6,\allowbreak 62,\allowbreak 24,\allowbreak 48,\allowbreak 21,\allowbreak 47,\allowbreak 22,\allowbreak 50,\allowbreak 19,\allowbreak 49,\allowbreak 20,\allowbreak 52,\allowbreak 17,\allowbreak 51,\allowbreak 18)\par\smallskip (68,\allowbreak 1)}\tabularnewline[5pt]
{\raggedright\arraybackslash {\ttfamily\footnotesize (2,\allowbreak 33,\allowbreak 33)}\par $K=33$\par pinned $(N,1)$} & {\ttfamily\footnotesize (64,\allowbreak 4,\allowbreak 63,\allowbreak 8,\allowbreak 59,\allowbreak 9,\allowbreak 58,\allowbreak 10,\allowbreak 57,\allowbreak 11,\allowbreak 56,\allowbreak 12,\allowbreak 55,\allowbreak 14,\allowbreak 54,\allowbreak 16,\allowbreak 49,\allowbreak 23,\allowbreak 47,\allowbreak 20,\allowbreak 43,\allowbreak 24,\allowbreak 45,\allowbreak 25,\allowbreak 38,\allowbreak 21,\allowbreak 46,\allowbreak 18,\allowbreak 48,\allowbreak 17,\allowbreak 52,\allowbreak 13,\allowbreak 50)\par\smallskip (66,\allowbreak 5,\allowbreak 62,\allowbreak 6,\allowbreak 60,\allowbreak 7,\allowbreak 39,\allowbreak 28,\allowbreak 36,\allowbreak 31,\allowbreak 40,\allowbreak 30,\allowbreak 37,\allowbreak 33,\allowbreak 34,\allowbreak 32,\allowbreak 35,\allowbreak 29,\allowbreak 41,\allowbreak 27,\allowbreak 42,\allowbreak 26,\allowbreak 44,\allowbreak 22,\allowbreak 51,\allowbreak 15,\allowbreak 67,\allowbreak 2,\allowbreak 65,\allowbreak 3,\allowbreak 61,\allowbreak 19,\allowbreak 53)\par\smallskip (68,\allowbreak 1)}\tabularnewline[5pt]
{\raggedright\arraybackslash {\ttfamily\footnotesize (2,\allowbreak 33,\allowbreak 34)}\par $K=34$\par pinned $(N,1)$} & {\ttfamily\footnotesize (44,\allowbreak 25,\allowbreak 47,\allowbreak 24,\allowbreak 49,\allowbreak 20,\allowbreak 52,\allowbreak 21,\allowbreak 48,\allowbreak 15,\allowbreak 50,\allowbreak 29,\allowbreak 41,\allowbreak 28,\allowbreak 43,\allowbreak 32,\allowbreak 36,\allowbreak 30,\allowbreak 39,\allowbreak 31,\allowbreak 38,\allowbreak 33,\allowbreak 35,\allowbreak 34,\allowbreak 37,\allowbreak 27,\allowbreak 45,\allowbreak 19,\allowbreak 55,\allowbreak 14,\allowbreak 53,\allowbreak 13,\allowbreak 56,\allowbreak 18)\par\smallskip (67,\allowbreak 4,\allowbreak 64,\allowbreak 6,\allowbreak 65,\allowbreak 3,\allowbreak 68,\allowbreak 2,\allowbreak 66,\allowbreak 5,\allowbreak 62,\allowbreak 8,\allowbreak 63,\allowbreak 7,\allowbreak 60,\allowbreak 11,\allowbreak 59,\allowbreak 9,\allowbreak 61,\allowbreak 10,\allowbreak 57,\allowbreak 12,\allowbreak 58,\allowbreak 16,\allowbreak 46,\allowbreak 22,\allowbreak 42,\allowbreak 26,\allowbreak 40,\allowbreak 23,\allowbreak 51,\allowbreak 17,\allowbreak 54)\par\smallskip (69,\allowbreak 1)}\tabularnewline[5pt]
{\raggedright\arraybackslash {\ttfamily\footnotesize (2,\allowbreak 34,\allowbreak 35)}\par $K=35$\par pinned $(N,1)$\par $2e\le b+4$} & {\ttfamily\footnotesize (69,\allowbreak 2,\allowbreak 70,\allowbreak 4,\allowbreak 68,\allowbreak 3,\allowbreak 66,\allowbreak 5,\allowbreak 67,\allowbreak 7,\allowbreak 63,\allowbreak 27,\allowbreak 45,\allowbreak 25,\allowbreak 48,\allowbreak 24,\allowbreak 46,\allowbreak 8,\allowbreak 65,\allowbreak 6,\allowbreak 64,\allowbreak 10,\allowbreak 60,\allowbreak 13,\allowbreak 59,\allowbreak 11,\allowbreak 62,\allowbreak 9,\allowbreak 61,\allowbreak 12,\allowbreak 57,\allowbreak 14,\allowbreak 58,\allowbreak 16)\par\smallskip (55,\allowbreak 15,\allowbreak 56,\allowbreak 17,\allowbreak 54,\allowbreak 19,\allowbreak 52,\allowbreak 18,\allowbreak 37,\allowbreak 33,\allowbreak 38,\allowbreak 35,\allowbreak 36,\allowbreak 34,\allowbreak 40,\allowbreak 32,\allowbreak 39,\allowbreak 30,\allowbreak 42,\allowbreak 31,\allowbreak 41,\allowbreak 28,\allowbreak 44,\allowbreak 29,\allowbreak 43,\allowbreak 26,\allowbreak 47,\allowbreak 22,\allowbreak 50,\allowbreak 23,\allowbreak 49,\allowbreak 20,\allowbreak 51,\allowbreak 21,\allowbreak 53)\par\smallskip (71,\allowbreak 1)}\tabularnewline[5pt]
{\raggedright\arraybackslash {\ttfamily\footnotesize (2,\allowbreak 34,\allowbreak 36)}\par $K=36$\par pinned $(N,1)$} & {\ttfamily\footnotesize (70,\allowbreak 2,\allowbreak 71,\allowbreak 4,\allowbreak 69,\allowbreak 3,\allowbreak 67,\allowbreak 5,\allowbreak 68,\allowbreak 7,\allowbreak 64,\allowbreak 10,\allowbreak 61,\allowbreak 12,\allowbreak 62,\allowbreak 9,\allowbreak 65,\allowbreak 6,\allowbreak 66,\allowbreak 8,\allowbreak 47,\allowbreak 26,\allowbreak 46,\allowbreak 29,\allowbreak 43,\allowbreak 33,\allowbreak 40,\allowbreak 32,\allowbreak 41,\allowbreak 35,\allowbreak 39,\allowbreak 34,\allowbreak 37,\allowbreak 36,\allowbreak 38,\allowbreak 19)\par\smallskip (72,\allowbreak 1)\par\smallskip (55,\allowbreak 21,\allowbreak 52,\allowbreak 20,\allowbreak 53,\allowbreak 23,\allowbreak 50,\allowbreak 22,\allowbreak 51,\allowbreak 25,\allowbreak 48,\allowbreak 24,\allowbreak 49,\allowbreak 27,\allowbreak 45,\allowbreak 30,\allowbreak 42,\allowbreak 31,\allowbreak 44,\allowbreak 28,\allowbreak 63,\allowbreak 11,\allowbreak 59,\allowbreak 15,\allowbreak 58,\allowbreak 13,\allowbreak 60,\allowbreak 14,\allowbreak 56,\allowbreak 18,\allowbreak 54,\allowbreak 17,\allowbreak 57,\allowbreak 16)}\tabularnewline[5pt]
\end{longtable}
\endgroup

\subsection{The 139 pinned minimum-four remainder cores}
\label{aux:min4}
\begingroup\footnotesize
\begin{longtable}{@{}p{25mm}p{115mm}@{}}
Profile / interface & Complete outward words\\\hline\endhead
{\raggedright\arraybackslash {\ttfamily\footnotesize (4,\allowbreak 9,\allowbreak 9)}\par $K=10$\par pinned W4\par $2e\le b+7$} & {\ttfamily\footnotesize (22,\allowbreak 1,\allowbreak 21,\allowbreak 2)\par\smallskip (18,\allowbreak 3,\allowbreak 20,\allowbreak 4,\allowbreak 17,\allowbreak 5,\allowbreak 15,\allowbreak 7,\allowbreak 16)\par\smallskip (14,\allowbreak 9,\allowbreak 13,\allowbreak 6,\allowbreak 12,\allowbreak 10,\allowbreak 11,\allowbreak 8,\allowbreak 19)}\tabularnewline[5pt]
{\raggedright\arraybackslash {\ttfamily\footnotesize (4,\allowbreak 9,\allowbreak 10)}\par $K=11$\par pinned W4} & {\ttfamily\footnotesize (23,\allowbreak 1,\allowbreak 22,\allowbreak 2)\par\smallskip (19,\allowbreak 3,\allowbreak 21,\allowbreak 4,\allowbreak 18,\allowbreak 5,\allowbreak 17,\allowbreak 8,\allowbreak 16,\allowbreak 6)\par\smallskip (15,\allowbreak 10,\allowbreak 14,\allowbreak 7,\allowbreak 13,\allowbreak 11,\allowbreak 12,\allowbreak 9,\allowbreak 20)}\tabularnewline[5pt]
{\raggedright\arraybackslash {\ttfamily\footnotesize (4,\allowbreak 9,\allowbreak 11)}\par $K=11$\par pinned W4} & {\ttfamily\footnotesize (24,\allowbreak 1,\allowbreak 23,\allowbreak 2)\par\smallskip (20,\allowbreak 3,\allowbreak 22,\allowbreak 4,\allowbreak 19,\allowbreak 5,\allowbreak 18,\allowbreak 6,\allowbreak 15,\allowbreak 7,\allowbreak 17)\par\smallskip (16,\allowbreak 9,\allowbreak 14,\allowbreak 8,\allowbreak 12,\allowbreak 11,\allowbreak 13,\allowbreak 10,\allowbreak 21)}\tabularnewline[5pt]
{\raggedright\arraybackslash {\ttfamily\footnotesize (4,\allowbreak 9,\allowbreak 12)}\par $K=12$\par pinned W4} & {\ttfamily\footnotesize (25,\allowbreak 1,\allowbreak 24,\allowbreak 2)\par\smallskip (21,\allowbreak 3,\allowbreak 23,\allowbreak 4,\allowbreak 20,\allowbreak 5,\allowbreak 19,\allowbreak 6,\allowbreak 18,\allowbreak 8,\allowbreak 16,\allowbreak 7)\par\smallskip (17,\allowbreak 10,\allowbreak 15,\allowbreak 9,\allowbreak 13,\allowbreak 12,\allowbreak 14,\allowbreak 11,\allowbreak 22)}\tabularnewline[5pt]
{\raggedright\arraybackslash {\ttfamily\footnotesize (4,\allowbreak 10,\allowbreak 10)}\par $K=12$\par pinned W4} & {\ttfamily\footnotesize (24,\allowbreak 1,\allowbreak 23,\allowbreak 2)\par\smallskip (20,\allowbreak 3,\allowbreak 22,\allowbreak 7,\allowbreak 21,\allowbreak 5,\allowbreak 17,\allowbreak 8,\allowbreak 19,\allowbreak 6)\par\smallskip (18,\allowbreak 10,\allowbreak 15,\allowbreak 9,\allowbreak 16,\allowbreak 12,\allowbreak 13,\allowbreak 11,\allowbreak 14,\allowbreak 4)}\tabularnewline[5pt]
{\raggedright\arraybackslash {\ttfamily\footnotesize (4,\allowbreak 10,\allowbreak 11)}\par $K=12$\par pinned W4\par $2e\le b+8$} & {\ttfamily\footnotesize (25,\allowbreak 1,\allowbreak 24,\allowbreak 2)\par\smallskip (20,\allowbreak 7,\allowbreak 23,\allowbreak 5,\allowbreak 17,\allowbreak 6,\allowbreak 16,\allowbreak 10,\allowbreak 14,\allowbreak 9,\allowbreak 18)\par\smallskip (21,\allowbreak 4,\allowbreak 19,\allowbreak 11,\allowbreak 13,\allowbreak 12,\allowbreak 15,\allowbreak 8,\allowbreak 22,\allowbreak 3)}\tabularnewline[5pt]
{\raggedright\arraybackslash {\ttfamily\footnotesize (4,\allowbreak 10,\allowbreak 12)}\par $K=13$\par pinned W4} & {\ttfamily\footnotesize (26,\allowbreak 1,\allowbreak 25,\allowbreak 2)\par\smallskip (22,\allowbreak 3,\allowbreak 24,\allowbreak 4,\allowbreak 21,\allowbreak 7,\allowbreak 23,\allowbreak 8,\allowbreak 20,\allowbreak 9,\allowbreak 19,\allowbreak 6)\par\smallskip (18,\allowbreak 10,\allowbreak 17,\allowbreak 11,\allowbreak 16,\allowbreak 12,\allowbreak 15,\allowbreak 13,\allowbreak 14,\allowbreak 5)}\tabularnewline[5pt]
{\raggedright\arraybackslash {\ttfamily\footnotesize (4,\allowbreak 10,\allowbreak 13)}\par $K=13$\par pinned W4} & {\ttfamily\footnotesize (27,\allowbreak 1,\allowbreak 26,\allowbreak 2)\par\smallskip (23,\allowbreak 3,\allowbreak 25,\allowbreak 4,\allowbreak 22,\allowbreak 5,\allowbreak 21,\allowbreak 6,\allowbreak 20,\allowbreak 9,\allowbreak 18,\allowbreak 11,\allowbreak 24)\par\smallskip (19,\allowbreak 7,\allowbreak 17,\allowbreak 12,\allowbreak 15,\allowbreak 13,\allowbreak 14,\allowbreak 10,\allowbreak 16,\allowbreak 8)}\tabularnewline[5pt]
{\raggedright\arraybackslash {\ttfamily\footnotesize (4,\allowbreak 10,\allowbreak 14)}\par $K=14$\par pinned W4} & {\ttfamily\footnotesize (28,\allowbreak 1,\allowbreak 27,\allowbreak 2)\par\smallskip (24,\allowbreak 3,\allowbreak 26,\allowbreak 4,\allowbreak 23,\allowbreak 6,\allowbreak 22,\allowbreak 8,\allowbreak 21,\allowbreak 9,\allowbreak 20,\allowbreak 10,\allowbreak 25,\allowbreak 5)\par\smallskip (18,\allowbreak 11,\allowbreak 19,\allowbreak 13,\allowbreak 17,\allowbreak 12,\allowbreak 15,\allowbreak 14,\allowbreak 16,\allowbreak 7)}\tabularnewline[5pt]
{\raggedright\arraybackslash {\ttfamily\footnotesize (4,\allowbreak 11,\allowbreak 11)}\par $K=12$\par pinned W4\par $2e\le b+7$} & {\ttfamily\footnotesize (26,\allowbreak 1,\allowbreak 25,\allowbreak 2)\par\smallskip (22,\allowbreak 3,\allowbreak 24,\allowbreak 4,\allowbreak 21,\allowbreak 6,\allowbreak 20,\allowbreak 7,\allowbreak 19,\allowbreak 8,\allowbreak 18)\par\smallskip (16,\allowbreak 9,\allowbreak 17,\allowbreak 11,\allowbreak 15,\allowbreak 10,\allowbreak 13,\allowbreak 12,\allowbreak 14,\allowbreak 5,\allowbreak 23)}\tabularnewline[5pt]
{\raggedright\arraybackslash {\ttfamily\footnotesize (4,\allowbreak 11,\allowbreak 12)}\par $K=13$\par pinned W4} & {\ttfamily\footnotesize (27,\allowbreak 1,\allowbreak 26,\allowbreak 2)\par\smallskip (23,\allowbreak 3,\allowbreak 25,\allowbreak 4,\allowbreak 22,\allowbreak 5,\allowbreak 21,\allowbreak 6,\allowbreak 20,\allowbreak 7,\allowbreak 18,\allowbreak 8)\par\smallskip (19,\allowbreak 10,\allowbreak 17,\allowbreak 9,\allowbreak 15,\allowbreak 11,\allowbreak 16,\allowbreak 13,\allowbreak 14,\allowbreak 12,\allowbreak 24)}\tabularnewline[5pt]
{\raggedright\arraybackslash {\ttfamily\footnotesize (4,\allowbreak 11,\allowbreak 13)}\par $K=13$\par pinned W4} & {\ttfamily\footnotesize (28,\allowbreak 1,\allowbreak 27,\allowbreak 2)\par\smallskip (24,\allowbreak 3,\allowbreak 26,\allowbreak 4,\allowbreak 23,\allowbreak 5,\allowbreak 22,\allowbreak 6,\allowbreak 21,\allowbreak 7,\allowbreak 18,\allowbreak 9,\allowbreak 19)\par\smallskip (20,\allowbreak 8,\allowbreak 16,\allowbreak 11,\allowbreak 17,\allowbreak 10,\allowbreak 14,\allowbreak 13,\allowbreak 15,\allowbreak 12,\allowbreak 25)}\tabularnewline[5pt]
{\raggedright\arraybackslash {\ttfamily\footnotesize (4,\allowbreak 11,\allowbreak 14)}\par $K=14$\par pinned W4} & {\ttfamily\footnotesize (29,\allowbreak 1,\allowbreak 28,\allowbreak 2)\par\smallskip (25,\allowbreak 3,\allowbreak 27,\allowbreak 4,\allowbreak 24,\allowbreak 5,\allowbreak 26,\allowbreak 8,\allowbreak 22,\allowbreak 9,\allowbreak 20,\allowbreak 11,\allowbreak 21,\allowbreak 6)\par\smallskip (17,\allowbreak 10,\allowbreak 18,\allowbreak 12,\allowbreak 16,\allowbreak 13,\allowbreak 15,\allowbreak 14,\allowbreak 19,\allowbreak 7,\allowbreak 23)}\tabularnewline[5pt]
{\raggedright\arraybackslash {\ttfamily\footnotesize (4,\allowbreak 12,\allowbreak 12)}\par $K=14$\par pinned W4} & {\ttfamily\footnotesize (28,\allowbreak 1,\allowbreak 27,\allowbreak 2)\par\smallskip (24,\allowbreak 3,\allowbreak 26,\allowbreak 4,\allowbreak 23,\allowbreak 6,\allowbreak 22,\allowbreak 8,\allowbreak 21,\allowbreak 10,\allowbreak 25,\allowbreak 5)\par\smallskip (18,\allowbreak 11,\allowbreak 20,\allowbreak 12,\allowbreak 17,\allowbreak 13,\allowbreak 16,\allowbreak 14,\allowbreak 15,\allowbreak 9,\allowbreak 19,\allowbreak 7)}\tabularnewline[5pt]
{\raggedright\arraybackslash {\ttfamily\footnotesize (4,\allowbreak 12,\allowbreak 13)}\par $K=14$\par pinned W4\par $2e\le b+8$} & {\ttfamily\footnotesize (29,\allowbreak 1,\allowbreak 28,\allowbreak 2)\par\smallskip (25,\allowbreak 3,\allowbreak 27,\allowbreak 4,\allowbreak 24,\allowbreak 5,\allowbreak 26,\allowbreak 8,\allowbreak 22,\allowbreak 9,\allowbreak 19,\allowbreak 7,\allowbreak 23)\par\smallskip (17,\allowbreak 12,\allowbreak 20,\allowbreak 11,\allowbreak 18,\allowbreak 14,\allowbreak 15,\allowbreak 13,\allowbreak 16,\allowbreak 10,\allowbreak 21,\allowbreak 6)}\tabularnewline[5pt]
{\raggedright\arraybackslash {\ttfamily\footnotesize (4,\allowbreak 12,\allowbreak 14)}\par $K=15$\par pinned W4} & {\ttfamily\footnotesize (30,\allowbreak 1,\allowbreak 29,\allowbreak 2)\par\smallskip (26,\allowbreak 3,\allowbreak 28,\allowbreak 4,\allowbreak 25,\allowbreak 5,\allowbreak 27,\allowbreak 8,\allowbreak 24,\allowbreak 9,\allowbreak 22,\allowbreak 10,\allowbreak 21,\allowbreak 7)\par\smallskip (18,\allowbreak 11,\allowbreak 20,\allowbreak 12,\allowbreak 17,\allowbreak 14,\allowbreak 16,\allowbreak 15,\allowbreak 19,\allowbreak 13,\allowbreak 23,\allowbreak 6)}\tabularnewline[5pt]
{\raggedright\arraybackslash {\ttfamily\footnotesize (4,\allowbreak 12,\allowbreak 15)}\par $K=15$\par pinned W4} & {\ttfamily\footnotesize (31,\allowbreak 1,\allowbreak 30,\allowbreak 2)\par\smallskip (27,\allowbreak 3,\allowbreak 29,\allowbreak 4,\allowbreak 26,\allowbreak 5,\allowbreak 28,\allowbreak 8,\allowbreak 25,\allowbreak 7,\allowbreak 23,\allowbreak 9,\allowbreak 22,\allowbreak 12,\allowbreak 24)\par\smallskip (19,\allowbreak 11,\allowbreak 20,\allowbreak 13,\allowbreak 18,\allowbreak 14,\allowbreak 17,\allowbreak 15,\allowbreak 16,\allowbreak 10,\allowbreak 21,\allowbreak 6)}\tabularnewline[5pt]
{\raggedright\arraybackslash {\ttfamily\footnotesize (4,\allowbreak 12,\allowbreak 16)}\par $K=16$\par pinned W4} & {\ttfamily\footnotesize (32,\allowbreak 1,\allowbreak 31,\allowbreak 2)\par\smallskip (28,\allowbreak 3,\allowbreak 30,\allowbreak 4,\allowbreak 27,\allowbreak 5,\allowbreak 29,\allowbreak 9,\allowbreak 26,\allowbreak 7,\allowbreak 25,\allowbreak 10,\allowbreak 24,\allowbreak 11,\allowbreak 22,\allowbreak 6)\par\smallskip (21,\allowbreak 12,\allowbreak 19,\allowbreak 15,\allowbreak 23,\allowbreak 13,\allowbreak 18,\allowbreak 16,\allowbreak 17,\allowbreak 14,\allowbreak 20,\allowbreak 8)}\tabularnewline[5pt]
{\raggedright\arraybackslash {\ttfamily\footnotesize (4,\allowbreak 13,\allowbreak 13)}\par $K=14$\par pinned W4\par $2e\le b+7$} & {\ttfamily\footnotesize (30,\allowbreak 1,\allowbreak 29,\allowbreak 2)\par\smallskip (26,\allowbreak 3,\allowbreak 28,\allowbreak 4,\allowbreak 25,\allowbreak 5,\allowbreak 27,\allowbreak 8,\allowbreak 22,\allowbreak 6,\allowbreak 21,\allowbreak 10,\allowbreak 23)\par\smallskip (18,\allowbreak 11,\allowbreak 20,\allowbreak 12,\allowbreak 17,\allowbreak 13,\allowbreak 16,\allowbreak 14,\allowbreak 15,\allowbreak 9,\allowbreak 19,\allowbreak 7,\allowbreak 24)}\tabularnewline[5pt]
{\raggedright\arraybackslash {\ttfamily\footnotesize (4,\allowbreak 13,\allowbreak 14)}\par $K=15$\par pinned W4} & {\ttfamily\footnotesize (31,\allowbreak 1,\allowbreak 30,\allowbreak 2)\par\smallskip (27,\allowbreak 3,\allowbreak 29,\allowbreak 4,\allowbreak 26,\allowbreak 5,\allowbreak 28,\allowbreak 8,\allowbreak 25,\allowbreak 7,\allowbreak 23,\allowbreak 9,\allowbreak 21,\allowbreak 6)\par\smallskip (19,\allowbreak 10,\allowbreak 20,\allowbreak 12,\allowbreak 18,\allowbreak 13,\allowbreak 17,\allowbreak 14,\allowbreak 16,\allowbreak 15,\allowbreak 22,\allowbreak 11,\allowbreak 24)}\tabularnewline[5pt]
{\raggedright\arraybackslash {\ttfamily\footnotesize (4,\allowbreak 13,\allowbreak 15)}\par $K=15$\par pinned W4} & {\ttfamily\footnotesize (32,\allowbreak 1,\allowbreak 31,\allowbreak 2)\par\smallskip (28,\allowbreak 3,\allowbreak 30,\allowbreak 4,\allowbreak 27,\allowbreak 5,\allowbreak 29,\allowbreak 8,\allowbreak 25,\allowbreak 7,\allowbreak 23,\allowbreak 9,\allowbreak 22,\allowbreak 12,\allowbreak 24)\par\smallskip (19,\allowbreak 11,\allowbreak 20,\allowbreak 13,\allowbreak 18,\allowbreak 14,\allowbreak 17,\allowbreak 15,\allowbreak 16,\allowbreak 10,\allowbreak 21,\allowbreak 6,\allowbreak 26)}\tabularnewline[5pt]
{\raggedright\arraybackslash {\ttfamily\footnotesize (4,\allowbreak 13,\allowbreak 16)}\par $K=16$\par pinned W4} & {\ttfamily\footnotesize (33,\allowbreak 1,\allowbreak 32,\allowbreak 2)\par\smallskip (29,\allowbreak 3,\allowbreak 31,\allowbreak 4,\allowbreak 28,\allowbreak 5,\allowbreak 30,\allowbreak 8,\allowbreak 27,\allowbreak 6,\allowbreak 24,\allowbreak 10,\allowbreak 25,\allowbreak 12,\allowbreak 23,\allowbreak 7)\par\smallskip (20,\allowbreak 13,\allowbreak 22,\allowbreak 14,\allowbreak 19,\allowbreak 15,\allowbreak 18,\allowbreak 16,\allowbreak 17,\allowbreak 11,\allowbreak 21,\allowbreak 9,\allowbreak 26)}\tabularnewline[5pt]
{\raggedright\arraybackslash {\ttfamily\footnotesize (4,\allowbreak 14,\allowbreak 14)}\par $K=16$\par pinned W4} & {\ttfamily\footnotesize (32,\allowbreak 1,\allowbreak 31,\allowbreak 2)\par\smallskip (28,\allowbreak 3,\allowbreak 30,\allowbreak 4,\allowbreak 27,\allowbreak 5,\allowbreak 29,\allowbreak 9,\allowbreak 26,\allowbreak 7,\allowbreak 25,\allowbreak 10,\allowbreak 24,\allowbreak 8)\par\smallskip (21,\allowbreak 11,\allowbreak 23,\allowbreak 12,\allowbreak 20,\allowbreak 13,\allowbreak 22,\allowbreak 16,\allowbreak 17,\allowbreak 15,\allowbreak 18,\allowbreak 14,\allowbreak 19,\allowbreak 6)}\tabularnewline[5pt]
{\raggedright\arraybackslash {\ttfamily\footnotesize (4,\allowbreak 14,\allowbreak 15)}\par $K=16$\par pinned W4\par $2e\le b+8$} & {\ttfamily\footnotesize (33,\allowbreak 1,\allowbreak 32,\allowbreak 2)\par\smallskip (29,\allowbreak 3,\allowbreak 31,\allowbreak 4,\allowbreak 28,\allowbreak 5,\allowbreak 30,\allowbreak 8,\allowbreak 27,\allowbreak 6,\allowbreak 24,\allowbreak 11,\allowbreak 25,\allowbreak 9,\allowbreak 26)\par\smallskip (20,\allowbreak 12,\allowbreak 23,\allowbreak 13,\allowbreak 18,\allowbreak 14,\allowbreak 21,\allowbreak 15,\allowbreak 17,\allowbreak 16,\allowbreak 19,\allowbreak 10,\allowbreak 22,\allowbreak 7)}\tabularnewline[5pt]
{\raggedright\arraybackslash {\ttfamily\footnotesize (4,\allowbreak 14,\allowbreak 16)}\par $K=17$\par pinned W4} & {\ttfamily\footnotesize (34,\allowbreak 1,\allowbreak 33,\allowbreak 2)\par\smallskip (30,\allowbreak 3,\allowbreak 32,\allowbreak 4,\allowbreak 29,\allowbreak 5,\allowbreak 31,\allowbreak 8,\allowbreak 28,\allowbreak 7,\allowbreak 26,\allowbreak 10,\allowbreak 27,\allowbreak 12,\allowbreak 23,\allowbreak 9)\par\smallskip (22,\allowbreak 13,\allowbreak 25,\allowbreak 15,\allowbreak 20,\allowbreak 14,\allowbreak 21,\allowbreak 17,\allowbreak 18,\allowbreak 16,\allowbreak 19,\allowbreak 11,\allowbreak 24,\allowbreak 6)}\tabularnewline[5pt]
{\raggedright\arraybackslash {\ttfamily\footnotesize (4,\allowbreak 14,\allowbreak 17)}\par $K=17$\par pinned W4} & {\ttfamily\footnotesize (35,\allowbreak 1,\allowbreak 34,\allowbreak 2)\par\smallskip (23,\allowbreak 12,\allowbreak 22,\allowbreak 17,\allowbreak 18,\allowbreak 16,\allowbreak 19,\allowbreak 15,\allowbreak 24,\allowbreak 10,\allowbreak 28,\allowbreak 7,\allowbreak 20,\allowbreak 14,\allowbreak 21,\allowbreak 13,\allowbreak 25)\par\smallskip (31,\allowbreak 3,\allowbreak 33,\allowbreak 4,\allowbreak 29,\allowbreak 9,\allowbreak 26,\allowbreak 11,\allowbreak 27,\allowbreak 8,\allowbreak 30,\allowbreak 6,\allowbreak 32,\allowbreak 5)}\tabularnewline[5pt]
{\raggedright\arraybackslash {\ttfamily\footnotesize (4,\allowbreak 14,\allowbreak 18)}\par $K=18$\par pinned W4} & {\ttfamily\footnotesize (36,\allowbreak 1,\allowbreak 35,\allowbreak 2)\par\smallskip (28,\allowbreak 9,\allowbreak 23,\allowbreak 13,\allowbreak 21,\allowbreak 17,\allowbreak 20,\allowbreak 18,\allowbreak 19,\allowbreak 14,\allowbreak 26,\allowbreak 6,\allowbreak 29,\allowbreak 8,\allowbreak 33,\allowbreak 3,\allowbreak 34,\allowbreak 5)\par\smallskip (32,\allowbreak 10,\allowbreak 25,\allowbreak 12,\allowbreak 30,\allowbreak 4,\allowbreak 31,\allowbreak 7,\allowbreak 24,\allowbreak 15,\allowbreak 22,\allowbreak 16,\allowbreak 27,\allowbreak 11)}\tabularnewline[5pt]
{\raggedright\arraybackslash {\ttfamily\footnotesize (4,\allowbreak 15,\allowbreak 15)}\par $K=16$\par pinned W4\par $2e\le b+7$} & {\ttfamily\footnotesize (34,\allowbreak 1,\allowbreak 33,\allowbreak 2)\par\smallskip (30,\allowbreak 3,\allowbreak 32,\allowbreak 4,\allowbreak 29,\allowbreak 5,\allowbreak 31,\allowbreak 8,\allowbreak 28,\allowbreak 6,\allowbreak 25,\allowbreak 7,\allowbreak 24,\allowbreak 10,\allowbreak 26)\par\smallskip (21,\allowbreak 9,\allowbreak 22,\allowbreak 11,\allowbreak 20,\allowbreak 13,\allowbreak 23,\allowbreak 15,\allowbreak 19,\allowbreak 14,\allowbreak 17,\allowbreak 16,\allowbreak 18,\allowbreak 12,\allowbreak 27)}\tabularnewline[5pt]
{\raggedright\arraybackslash {\ttfamily\footnotesize (4,\allowbreak 15,\allowbreak 16)}\par $K=17$\par pinned W4} & {\ttfamily\footnotesize (35,\allowbreak 1,\allowbreak 34,\allowbreak 2)\par\smallskip (31,\allowbreak 3,\allowbreak 33,\allowbreak 4,\allowbreak 30,\allowbreak 5,\allowbreak 32,\allowbreak 8,\allowbreak 29,\allowbreak 6,\allowbreak 28,\allowbreak 9,\allowbreak 26,\allowbreak 10,\allowbreak 25,\allowbreak 7)\par\smallskip (20,\allowbreak 12,\allowbreak 24,\allowbreak 11,\allowbreak 22,\allowbreak 15,\allowbreak 21,\allowbreak 16,\allowbreak 19,\allowbreak 17,\allowbreak 18,\allowbreak 14,\allowbreak 23,\allowbreak 13,\allowbreak 27)}\tabularnewline[5pt]
{\raggedright\arraybackslash {\ttfamily\footnotesize (4,\allowbreak 15,\allowbreak 17)}\par $K=17$\par pinned W4} & {\ttfamily\footnotesize (36,\allowbreak 1,\allowbreak 35,\allowbreak 2)\par\smallskip (30,\allowbreak 3,\allowbreak 34,\allowbreak 5,\allowbreak 26,\allowbreak 13,\allowbreak 21,\allowbreak 12,\allowbreak 24,\allowbreak 7,\allowbreak 33,\allowbreak 9,\allowbreak 29,\allowbreak 10,\allowbreak 28,\allowbreak 6,\allowbreak 31)\par\smallskip (32,\allowbreak 4,\allowbreak 27,\allowbreak 11,\allowbreak 22,\allowbreak 8,\allowbreak 23,\allowbreak 16,\allowbreak 18,\allowbreak 17,\allowbreak 20,\allowbreak 14,\allowbreak 19,\allowbreak 15,\allowbreak 25)}\tabularnewline[5pt]
{\raggedright\arraybackslash {\ttfamily\footnotesize (4,\allowbreak 15,\allowbreak 18)}\par $K=18$\par pinned W4} & {\ttfamily\footnotesize (37,\allowbreak 1,\allowbreak 36,\allowbreak 2)\par\smallskip (30,\allowbreak 8,\allowbreak 27,\allowbreak 9,\allowbreak 23,\allowbreak 16,\allowbreak 28,\allowbreak 4,\allowbreak 31,\allowbreak 5,\allowbreak 34,\allowbreak 3,\allowbreak 35,\allowbreak 7,\allowbreak 32)\par\smallskip (33,\allowbreak 12,\allowbreak 29,\allowbreak 6,\allowbreak 26,\allowbreak 10,\allowbreak 25,\allowbreak 15,\allowbreak 21,\allowbreak 17,\allowbreak 20,\allowbreak 18,\allowbreak 19,\allowbreak 14,\allowbreak 22,\allowbreak 13,\allowbreak 24,\allowbreak 11)}\tabularnewline[5pt]
{\raggedright\arraybackslash {\ttfamily\footnotesize (4,\allowbreak 16,\allowbreak 16)}\par $K=18$\par pinned W4} & {\ttfamily\footnotesize (36,\allowbreak 1,\allowbreak 35,\allowbreak 2)\par\smallskip (32,\allowbreak 3,\allowbreak 34,\allowbreak 4,\allowbreak 31,\allowbreak 5,\allowbreak 33,\allowbreak 8,\allowbreak 30,\allowbreak 6,\allowbreak 29,\allowbreak 10,\allowbreak 28,\allowbreak 11,\allowbreak 27,\allowbreak 7)\par\smallskip (21,\allowbreak 13,\allowbreak 25,\allowbreak 12,\allowbreak 26,\allowbreak 15,\allowbreak 22,\allowbreak 16,\allowbreak 20,\allowbreak 17,\allowbreak 19,\allowbreak 18,\allowbreak 23,\allowbreak 14,\allowbreak 24,\allowbreak 9)}\tabularnewline[5pt]
{\raggedright\arraybackslash {\ttfamily\footnotesize (4,\allowbreak 16,\allowbreak 17)}\par $K=18$\par pinned W4\par $2e\le b+8$} & {\ttfamily\footnotesize (37,\allowbreak 1,\allowbreak 36,\allowbreak 2)\par\smallskip (28,\allowbreak 11,\allowbreak 23,\allowbreak 18,\allowbreak 19,\allowbreak 17,\allowbreak 20,\allowbreak 16,\allowbreak 22,\allowbreak 9,\allowbreak 31,\allowbreak 8,\allowbreak 29,\allowbreak 5,\allowbreak 32,\allowbreak 6)\par\smallskip (33,\allowbreak 4,\allowbreak 34,\allowbreak 3,\allowbreak 35,\allowbreak 10,\allowbreak 25,\allowbreak 15,\allowbreak 24,\allowbreak 13,\allowbreak 21,\allowbreak 14,\allowbreak 30,\allowbreak 12,\allowbreak 26,\allowbreak 7,\allowbreak 27)}\tabularnewline[5pt]
{\raggedright\arraybackslash {\ttfamily\footnotesize (4,\allowbreak 16,\allowbreak 18)}\par $K=19$\par pinned W4} & {\ttfamily\footnotesize (38,\allowbreak 1,\allowbreak 37,\allowbreak 2)\par\smallskip (27,\allowbreak 10,\allowbreak 29,\allowbreak 14,\allowbreak 22,\allowbreak 16,\allowbreak 21,\allowbreak 18,\allowbreak 20,\allowbreak 19,\allowbreak 23,\allowbreak 12,\allowbreak 28,\allowbreak 4,\allowbreak 35,\allowbreak 3,\allowbreak 36,\allowbreak 6)\par\smallskip (34,\allowbreak 5,\allowbreak 33,\allowbreak 7,\allowbreak 32,\allowbreak 9,\allowbreak 30,\allowbreak 8,\allowbreak 26,\allowbreak 13,\allowbreak 25,\allowbreak 15,\allowbreak 24,\allowbreak 17,\allowbreak 31,\allowbreak 11)}\tabularnewline[5pt]
{\raggedright\arraybackslash {\ttfamily\footnotesize (4,\allowbreak 16,\allowbreak 19)}\par $K=19$\par pinned W4} & {\ttfamily\footnotesize (39,\allowbreak 1,\allowbreak 38,\allowbreak 2)\par\smallskip (34,\allowbreak 7,\allowbreak 25,\allowbreak 16,\allowbreak 29,\allowbreak 14,\allowbreak 22,\allowbreak 18,\allowbreak 21,\allowbreak 19,\allowbreak 20,\allowbreak 15,\allowbreak 26,\allowbreak 10,\allowbreak 36,\allowbreak 5,\allowbreak 30,\allowbreak 11,\allowbreak 28)\par\smallskip (35,\allowbreak 3,\allowbreak 33,\allowbreak 4,\allowbreak 37,\allowbreak 9,\allowbreak 31,\allowbreak 8,\allowbreak 32,\allowbreak 12,\allowbreak 24,\allowbreak 17,\allowbreak 23,\allowbreak 13,\allowbreak 27,\allowbreak 6)}\tabularnewline[5pt]
{\raggedright\arraybackslash {\ttfamily\footnotesize (4,\allowbreak 16,\allowbreak 20)}\par $K=20$\par pinned W4} & {\ttfamily\footnotesize (40,\allowbreak 1,\allowbreak 39,\allowbreak 2)\par\smallskip (34,\allowbreak 7,\allowbreak 32,\allowbreak 4,\allowbreak 37,\allowbreak 8,\allowbreak 38,\allowbreak 3,\allowbreak 35,\allowbreak 15,\allowbreak 24,\allowbreak 13,\allowbreak 30,\allowbreak 11,\allowbreak 27,\allowbreak 6)\par\smallskip (36,\allowbreak 5,\allowbreak 31,\allowbreak 9,\allowbreak 33,\allowbreak 10,\allowbreak 28,\allowbreak 14,\allowbreak 29,\allowbreak 17,\allowbreak 23,\allowbreak 19,\allowbreak 22,\allowbreak 20,\allowbreak 21,\allowbreak 16,\allowbreak 26,\allowbreak 18,\allowbreak 25,\allowbreak 12)}\tabularnewline[5pt]
{\raggedright\arraybackslash {\ttfamily\footnotesize (4,\allowbreak 17,\allowbreak 17)}\par $K=18$\par pinned W4\par $2e\le b+7$} & {\ttfamily\footnotesize (38,\allowbreak 1,\allowbreak 37,\allowbreak 2)\par\smallskip (33,\allowbreak 5,\allowbreak 30,\allowbreak 9,\allowbreak 31,\allowbreak 12,\allowbreak 23,\allowbreak 16,\allowbreak 21,\allowbreak 15,\allowbreak 19,\allowbreak 18,\allowbreak 20,\allowbreak 17,\allowbreak 26,\allowbreak 10,\allowbreak 27)\par\smallskip (34,\allowbreak 7,\allowbreak 36,\allowbreak 6,\allowbreak 29,\allowbreak 3,\allowbreak 35,\allowbreak 4,\allowbreak 22,\allowbreak 14,\allowbreak 24,\allowbreak 11,\allowbreak 25,\allowbreak 13,\allowbreak 28,\allowbreak 8,\allowbreak 32)}\tabularnewline[5pt]
{\raggedright\arraybackslash {\ttfamily\footnotesize (4,\allowbreak 17,\allowbreak 18)}\par $K=19$\par pinned W4} & {\ttfamily\footnotesize (39,\allowbreak 1,\allowbreak 38,\allowbreak 2)\par\smallskip (27,\allowbreak 10,\allowbreak 29,\allowbreak 9,\allowbreak 30,\allowbreak 8,\allowbreak 31,\allowbreak 7,\allowbreak 32,\allowbreak 6,\allowbreak 36,\allowbreak 4,\allowbreak 37,\allowbreak 3,\allowbreak 34,\allowbreak 5,\allowbreak 33)\par\smallskip (35,\allowbreak 17,\allowbreak 24,\allowbreak 15,\allowbreak 26,\allowbreak 14,\allowbreak 22,\allowbreak 16,\allowbreak 21,\allowbreak 18,\allowbreak 20,\allowbreak 19,\allowbreak 23,\allowbreak 13,\allowbreak 28,\allowbreak 12,\allowbreak 25,\allowbreak 11)}\tabularnewline[5pt]
{\raggedright\arraybackslash {\ttfamily\footnotesize (4,\allowbreak 17,\allowbreak 19)}\par $K=19$\par pinned W4} & {\ttfamily\footnotesize (40,\allowbreak 1,\allowbreak 39,\allowbreak 2)\par\smallskip (35,\allowbreak 3,\allowbreak 33,\allowbreak 14,\allowbreak 22,\allowbreak 19,\allowbreak 20,\allowbreak 18,\allowbreak 24,\allowbreak 9,\allowbreak 32,\allowbreak 5,\allowbreak 31,\allowbreak 6,\allowbreak 27,\allowbreak 11,\allowbreak 28)\par\smallskip (36,\allowbreak 8,\allowbreak 37,\allowbreak 4,\allowbreak 38,\allowbreak 7,\allowbreak 29,\allowbreak 15,\allowbreak 26,\allowbreak 17,\allowbreak 21,\allowbreak 16,\allowbreak 23,\allowbreak 13,\allowbreak 25,\allowbreak 12,\allowbreak 30,\allowbreak 10,\allowbreak 34)}\tabularnewline[5pt]
{\raggedright\arraybackslash {\ttfamily\footnotesize (4,\allowbreak 17,\allowbreak 20)}\par $K=20$\par pinned W4} & {\ttfamily\footnotesize (41,\allowbreak 1,\allowbreak 40,\allowbreak 2)\par\smallskip (33,\allowbreak 6,\allowbreak 29,\allowbreak 11,\allowbreak 31,\allowbreak 15,\allowbreak 27,\allowbreak 17,\allowbreak 22,\allowbreak 18,\allowbreak 24,\allowbreak 16,\allowbreak 30,\allowbreak 8,\allowbreak 34,\allowbreak 10,\allowbreak 35)\par\smallskip (37,\allowbreak 7,\allowbreak 28,\allowbreak 13,\allowbreak 32,\allowbreak 4,\allowbreak 36,\allowbreak 5,\allowbreak 39,\allowbreak 3,\allowbreak 38,\allowbreak 9,\allowbreak 26,\allowbreak 19,\allowbreak 21,\allowbreak 20,\allowbreak 23,\allowbreak 14,\allowbreak 25,\allowbreak 12)}\tabularnewline[5pt]
{\raggedright\arraybackslash {\ttfamily\footnotesize (4,\allowbreak 18,\allowbreak 18)}\par $K=20$\par pinned W4} & {\ttfamily\footnotesize (40,\allowbreak 1,\allowbreak 39,\allowbreak 2)\par\smallskip (33,\allowbreak 6,\allowbreak 25,\allowbreak 15,\allowbreak 26,\allowbreak 19,\allowbreak 22,\allowbreak 20,\allowbreak 21,\allowbreak 17,\allowbreak 29,\allowbreak 13,\allowbreak 31,\allowbreak 10,\allowbreak 32,\allowbreak 8,\allowbreak 28,\allowbreak 11)\par\smallskip (36,\allowbreak 4,\allowbreak 30,\allowbreak 16,\allowbreak 24,\allowbreak 18,\allowbreak 23,\allowbreak 14,\allowbreak 27,\allowbreak 12,\allowbreak 35,\allowbreak 5,\allowbreak 34,\allowbreak 9,\allowbreak 37,\allowbreak 3,\allowbreak 38,\allowbreak 7)}\tabularnewline[5pt]
{\raggedright\arraybackslash {\ttfamily\footnotesize (4,\allowbreak 18,\allowbreak 19)}\par $K=20$\par pinned W4\par $2e\le b+8$} & {\ttfamily\footnotesize (41,\allowbreak 1,\allowbreak 40,\allowbreak 2)\par\smallskip (35,\allowbreak 6,\allowbreak 36,\allowbreak 4,\allowbreak 38,\allowbreak 5,\allowbreak 33,\allowbreak 9,\allowbreak 29,\allowbreak 14,\allowbreak 22,\allowbreak 18,\allowbreak 23,\allowbreak 20,\allowbreak 21,\allowbreak 19,\allowbreak 25,\allowbreak 13,\allowbreak 30)\par\smallskip (37,\allowbreak 10,\allowbreak 28,\allowbreak 15,\allowbreak 24,\allowbreak 17,\allowbreak 39,\allowbreak 3,\allowbreak 34,\allowbreak 8,\allowbreak 31,\allowbreak 12,\allowbreak 26,\allowbreak 16,\allowbreak 27,\allowbreak 11,\allowbreak 32,\allowbreak 7)}\tabularnewline[5pt]
{\raggedright\arraybackslash {\ttfamily\footnotesize (4,\allowbreak 18,\allowbreak 20)}\par $K=21$\par pinned W4} & {\ttfamily\footnotesize (42,\allowbreak 1,\allowbreak 41,\allowbreak 2)\par\smallskip (37,\allowbreak 4,\allowbreak 35,\allowbreak 5,\allowbreak 30,\allowbreak 17,\allowbreak 24,\allowbreak 20,\allowbreak 28,\allowbreak 11,\allowbreak 31,\allowbreak 9,\allowbreak 32,\allowbreak 8,\allowbreak 40,\allowbreak 6,\allowbreak 34,\allowbreak 7)\par\smallskip (38,\allowbreak 3,\allowbreak 39,\allowbreak 10,\allowbreak 36,\allowbreak 15,\allowbreak 33,\allowbreak 14,\allowbreak 29,\allowbreak 13,\allowbreak 25,\allowbreak 19,\allowbreak 22,\allowbreak 21,\allowbreak 23,\allowbreak 18,\allowbreak 27,\allowbreak 16,\allowbreak 26,\allowbreak 12)}\tabularnewline[5pt]
{\raggedright\arraybackslash {\ttfamily\footnotesize (4,\allowbreak 18,\allowbreak 21)}\par $K=21$\par pinned W4} & {\ttfamily\footnotesize (43,\allowbreak 1,\allowbreak 42,\allowbreak 2)\par\smallskip (38,\allowbreak 8,\allowbreak 34,\allowbreak 13,\allowbreak 25,\allowbreak 18,\allowbreak 24,\allowbreak 19,\allowbreak 23,\allowbreak 20,\allowbreak 22,\allowbreak 21,\allowbreak 29,\allowbreak 10,\allowbreak 37,\allowbreak 3,\allowbreak 35,\allowbreak 7)\par\smallskip (39,\allowbreak 6,\allowbreak 30,\allowbreak 14,\allowbreak 36,\allowbreak 11,\allowbreak 28,\allowbreak 15,\allowbreak 33,\allowbreak 4,\allowbreak 41,\allowbreak 5,\allowbreak 40,\allowbreak 9,\allowbreak 32,\allowbreak 12,\allowbreak 27,\allowbreak 16,\allowbreak 26,\allowbreak 17,\allowbreak 31)}\tabularnewline[5pt]
{\raggedright\arraybackslash {\ttfamily\footnotesize (4,\allowbreak 18,\allowbreak 22)}\par $K=22$\par pinned W4} & {\ttfamily\footnotesize (44,\allowbreak 1,\allowbreak 43,\allowbreak 2)\par\smallskip (39,\allowbreak 4,\allowbreak 36,\allowbreak 3,\allowbreak 41,\allowbreak 5,\allowbreak 42,\allowbreak 8,\allowbreak 38,\allowbreak 15,\allowbreak 29,\allowbreak 20,\allowbreak 25,\allowbreak 19,\allowbreak 30,\allowbreak 9,\allowbreak 33,\allowbreak 7)\par\smallskip (40,\allowbreak 11,\allowbreak 27,\allowbreak 17,\allowbreak 24,\allowbreak 21,\allowbreak 23,\allowbreak 22,\allowbreak 26,\allowbreak 18,\allowbreak 31,\allowbreak 14,\allowbreak 32,\allowbreak 12,\allowbreak 34,\allowbreak 6,\allowbreak 37,\allowbreak 10,\allowbreak 35,\allowbreak 16,\allowbreak 28,\allowbreak 13)}\tabularnewline[5pt]
{\raggedright\arraybackslash {\ttfamily\footnotesize (4,\allowbreak 19,\allowbreak 19)}\par $K=20$\par pinned W4\par $2e\le b+7$} & {\ttfamily\footnotesize (42,\allowbreak 1,\allowbreak 41,\allowbreak 2)\par\smallskip (37,\allowbreak 3,\allowbreak 39,\allowbreak 10,\allowbreak 33,\allowbreak 9,\allowbreak 29,\allowbreak 12,\allowbreak 25,\allowbreak 19,\allowbreak 21,\allowbreak 20,\allowbreak 23,\allowbreak 15,\allowbreak 31,\allowbreak 6,\allowbreak 32,\allowbreak 11,\allowbreak 30)\par\smallskip (38,\allowbreak 7,\allowbreak 40,\allowbreak 5,\allowbreak 27,\allowbreak 18,\allowbreak 22,\allowbreak 17,\allowbreak 24,\allowbreak 13,\allowbreak 28,\allowbreak 14,\allowbreak 26,\allowbreak 16,\allowbreak 34,\allowbreak 4,\allowbreak 36,\allowbreak 8,\allowbreak 35)}\tabularnewline[5pt]
{\raggedright\arraybackslash {\ttfamily\footnotesize (4,\allowbreak 19,\allowbreak 20)}\par $K=21$\par pinned W4} & {\ttfamily\footnotesize (43,\allowbreak 1,\allowbreak 42,\allowbreak 2)\par\smallskip (35,\allowbreak 6,\allowbreak 37,\allowbreak 4,\allowbreak 40,\allowbreak 3,\allowbreak 41,\allowbreak 9,\allowbreak 34,\allowbreak 10,\allowbreak 30,\allowbreak 18,\allowbreak 25,\allowbreak 16,\allowbreak 26,\allowbreak 15,\allowbreak 38,\allowbreak 8,\allowbreak 36)\par\smallskip (39,\allowbreak 5,\allowbreak 32,\allowbreak 11,\allowbreak 28,\allowbreak 20,\allowbreak 22,\allowbreak 21,\allowbreak 24,\allowbreak 19,\allowbreak 23,\allowbreak 17,\allowbreak 33,\allowbreak 7,\allowbreak 29,\allowbreak 14,\allowbreak 27,\allowbreak 13,\allowbreak 31,\allowbreak 12)}\tabularnewline[5pt]
{\raggedright\arraybackslash {\ttfamily\footnotesize (4,\allowbreak 19,\allowbreak 21)}\par $K=21$\par pinned W4} & {\ttfamily\footnotesize (44,\allowbreak 1,\allowbreak 43,\allowbreak 2)\par\smallskip (30,\allowbreak 17,\allowbreak 25,\allowbreak 18,\allowbreak 23,\allowbreak 21,\allowbreak 22,\allowbreak 19,\allowbreak 28,\allowbreak 16,\allowbreak 27,\allowbreak 12,\allowbreak 33,\allowbreak 13,\allowbreak 29,\allowbreak 15,\allowbreak 34,\allowbreak 11,\allowbreak 36,\allowbreak 7,\allowbreak 31)\par\smallskip (40,\allowbreak 5,\allowbreak 39,\allowbreak 8,\allowbreak 35,\allowbreak 9,\allowbreak 26,\allowbreak 20,\allowbreak 24,\allowbreak 14,\allowbreak 32,\allowbreak 10,\allowbreak 38,\allowbreak 6,\allowbreak 42,\allowbreak 3,\allowbreak 41,\allowbreak 4,\allowbreak 37)}\tabularnewline[5pt]
{\raggedright\arraybackslash {\ttfamily\footnotesize (4,\allowbreak 19,\allowbreak 22)}\par $K=22$\par pinned W4} & {\ttfamily\footnotesize (45,\allowbreak 1,\allowbreak 44,\allowbreak 2)\par\smallskip (35,\allowbreak 14,\allowbreak 37,\allowbreak 5,\allowbreak 42,\allowbreak 4,\allowbreak 43,\allowbreak 3,\allowbreak 39,\allowbreak 8,\allowbreak 34,\allowbreak 16,\allowbreak 28,\allowbreak 19,\allowbreak 25,\allowbreak 22,\allowbreak 23,\allowbreak 21,\allowbreak 26,\allowbreak 18,\allowbreak 29,\allowbreak 13)\par\smallskip (41,\allowbreak 12,\allowbreak 31,\allowbreak 11,\allowbreak 33,\allowbreak 9,\allowbreak 36,\allowbreak 6,\allowbreak 40,\allowbreak 7,\allowbreak 32,\allowbreak 15,\allowbreak 30,\allowbreak 17,\allowbreak 27,\allowbreak 20,\allowbreak 24,\allowbreak 10,\allowbreak 38)}\tabularnewline[5pt]
{\raggedright\arraybackslash {\ttfamily\footnotesize (4,\allowbreak 20,\allowbreak 20)}\par $K=22$\par pinned W4} & {\ttfamily\footnotesize (44,\allowbreak 1,\allowbreak 43,\allowbreak 2)\par\smallskip (39,\allowbreak 3,\allowbreak 41,\allowbreak 4,\allowbreak 33,\allowbreak 15,\allowbreak 32,\allowbreak 7,\allowbreak 34,\allowbreak 11,\allowbreak 37,\allowbreak 6,\allowbreak 38,\allowbreak 10,\allowbreak 31,\allowbreak 17,\allowbreak 28,\allowbreak 9,\allowbreak 42,\allowbreak 8)\par\smallskip (40,\allowbreak 5,\allowbreak 35,\allowbreak 13,\allowbreak 23,\allowbreak 22,\allowbreak 24,\allowbreak 21,\allowbreak 25,\allowbreak 20,\allowbreak 26,\allowbreak 19,\allowbreak 27,\allowbreak 18,\allowbreak 30,\allowbreak 14,\allowbreak 29,\allowbreak 16,\allowbreak 36,\allowbreak 12)}\tabularnewline[5pt]
{\raggedright\arraybackslash {\ttfamily\footnotesize (4,\allowbreak 20,\allowbreak 21)}\par $K=22$\par pinned W4\par $2e\le b+8$} & {\ttfamily\footnotesize (45,\allowbreak 1,\allowbreak 44,\allowbreak 2)\par\smallskip (36,\allowbreak 16,\allowbreak 28,\allowbreak 15,\allowbreak 30,\allowbreak 19,\allowbreak 25,\allowbreak 20,\allowbreak 27,\allowbreak 18,\allowbreak 26,\allowbreak 22,\allowbreak 23,\allowbreak 21,\allowbreak 24,\allowbreak 14,\allowbreak 37,\allowbreak 9,\allowbreak 38,\allowbreak 8)\par\smallskip (41,\allowbreak 7,\allowbreak 39,\allowbreak 6,\allowbreak 32,\allowbreak 13,\allowbreak 40,\allowbreak 5,\allowbreak 42,\allowbreak 4,\allowbreak 43,\allowbreak 3,\allowbreak 34,\allowbreak 12,\allowbreak 29,\allowbreak 11,\allowbreak 35,\allowbreak 10,\allowbreak 31,\allowbreak 17,\allowbreak 33)}\tabularnewline[5pt]
{\raggedright\arraybackslash {\ttfamily\footnotesize (4,\allowbreak 20,\allowbreak 22)}\par $K=23$\par pinned W4} & {\ttfamily\footnotesize (46,\allowbreak 1,\allowbreak 45,\allowbreak 2)\par\smallskip (41,\allowbreak 6,\allowbreak 36,\allowbreak 10,\allowbreak 33,\allowbreak 19,\allowbreak 26,\allowbreak 21,\allowbreak 27,\allowbreak 18,\allowbreak 29,\allowbreak 17,\allowbreak 30,\allowbreak 14,\allowbreak 31,\allowbreak 16,\allowbreak 34,\allowbreak 15,\allowbreak 35,\allowbreak 13)\par\smallskip (42,\allowbreak 9,\allowbreak 37,\allowbreak 12,\allowbreak 39,\allowbreak 7,\allowbreak 38,\allowbreak 4,\allowbreak 28,\allowbreak 20,\allowbreak 24,\allowbreak 23,\allowbreak 25,\allowbreak 22,\allowbreak 32,\allowbreak 11,\allowbreak 40,\allowbreak 3,\allowbreak 43,\allowbreak 5,\allowbreak 44,\allowbreak 8)}\tabularnewline[5pt]
{\raggedright\arraybackslash {\ttfamily\footnotesize (4,\allowbreak 20,\allowbreak 23)}\par $K=23$\par pinned W4} & {\ttfamily\footnotesize (47,\allowbreak 1,\allowbreak 46,\allowbreak 2)\par\smallskip (40,\allowbreak 12,\allowbreak 35,\allowbreak 14,\allowbreak 32,\allowbreak 13,\allowbreak 37,\allowbreak 15,\allowbreak 30,\allowbreak 17,\allowbreak 31,\allowbreak 5,\allowbreak 44,\allowbreak 3,\allowbreak 45,\allowbreak 10,\allowbreak 41,\allowbreak 4,\allowbreak 42,\allowbreak 8)\par\smallskip (43,\allowbreak 7,\allowbreak 39,\allowbreak 6,\allowbreak 36,\allowbreak 11,\allowbreak 38,\allowbreak 9,\allowbreak 29,\allowbreak 19,\allowbreak 28,\allowbreak 20,\allowbreak 26,\allowbreak 22,\allowbreak 27,\allowbreak 16,\allowbreak 33,\allowbreak 21,\allowbreak 24,\allowbreak 23,\allowbreak 25,\allowbreak 18,\allowbreak 34)}\tabularnewline[5pt]
{\raggedright\arraybackslash {\ttfamily\footnotesize (4,\allowbreak 20,\allowbreak 24)}\par $K=24$\par pinned W4} & {\ttfamily\footnotesize (48,\allowbreak 1,\allowbreak 47,\allowbreak 2)\par\smallskip (37,\allowbreak 12,\allowbreak 34,\allowbreak 13,\allowbreak 32,\allowbreak 17,\allowbreak 28,\allowbreak 24,\allowbreak 25,\allowbreak 23,\allowbreak 26,\allowbreak 21,\allowbreak 27,\allowbreak 20,\allowbreak 29,\allowbreak 19,\allowbreak 33,\allowbreak 16,\allowbreak 39,\allowbreak 8)\par\smallskip (44,\allowbreak 9,\allowbreak 38,\allowbreak 11,\allowbreak 35,\allowbreak 15,\allowbreak 31,\allowbreak 18,\allowbreak 30,\allowbreak 22,\allowbreak 40,\allowbreak 7,\allowbreak 41,\allowbreak 5,\allowbreak 43,\allowbreak 4,\allowbreak 45,\allowbreak 3,\allowbreak 46,\allowbreak 6,\allowbreak 36,\allowbreak 10,\allowbreak 42,\allowbreak 14)}\tabularnewline[5pt]
{\raggedright\arraybackslash {\ttfamily\footnotesize (4,\allowbreak 21,\allowbreak 21)}\par $K=22$\par pinned W4\par $2e\le b+7$} & {\ttfamily\footnotesize (46,\allowbreak 1,\allowbreak 45,\allowbreak 2)\par\smallskip (40,\allowbreak 5,\allowbreak 37,\allowbreak 9,\allowbreak 29,\allowbreak 16,\allowbreak 26,\allowbreak 21,\allowbreak 23,\allowbreak 22,\allowbreak 25,\allowbreak 19,\allowbreak 30,\allowbreak 15,\allowbreak 27,\allowbreak 18,\allowbreak 34,\allowbreak 12,\allowbreak 39,\allowbreak 8,\allowbreak 38)\par\smallskip (42,\allowbreak 6,\allowbreak 44,\allowbreak 3,\allowbreak 36,\allowbreak 7,\allowbreak 41,\allowbreak 4,\allowbreak 43,\allowbreak 17,\allowbreak 24,\allowbreak 20,\allowbreak 28,\allowbreak 14,\allowbreak 31,\allowbreak 13,\allowbreak 32,\allowbreak 11,\allowbreak 35,\allowbreak 10,\allowbreak 33)}\tabularnewline[5pt]
{\raggedright\arraybackslash {\ttfamily\footnotesize (4,\allowbreak 21,\allowbreak 22)}\par $K=23$\par pinned W4} & {\ttfamily\footnotesize (47,\allowbreak 1,\allowbreak 46,\allowbreak 2)\par\smallskip (42,\allowbreak 7,\allowbreak 31,\allowbreak 18,\allowbreak 29,\allowbreak 22,\allowbreak 24,\allowbreak 23,\allowbreak 26,\allowbreak 12,\allowbreak 37,\allowbreak 4,\allowbreak 40,\allowbreak 8,\allowbreak 45,\allowbreak 5,\allowbreak 44,\allowbreak 3,\allowbreak 41,\allowbreak 10,\allowbreak 33,\allowbreak 13)\par\smallskip (43,\allowbreak 9,\allowbreak 35,\allowbreak 6,\allowbreak 36,\allowbreak 14,\allowbreak 32,\allowbreak 16,\allowbreak 28,\allowbreak 19,\allowbreak 27,\allowbreak 21,\allowbreak 25,\allowbreak 20,\allowbreak 30,\allowbreak 15,\allowbreak 34,\allowbreak 17,\allowbreak 38,\allowbreak 11,\allowbreak 39)}\tabularnewline[5pt]
{\raggedright\arraybackslash {\ttfamily\footnotesize (4,\allowbreak 21,\allowbreak 23)}\par $K=23$\par pinned W4} & {\ttfamily\footnotesize (48,\allowbreak 1,\allowbreak 47,\allowbreak 2)\par\smallskip (43,\allowbreak 5,\allowbreak 38,\allowbreak 14,\allowbreak 36,\allowbreak 4,\allowbreak 45,\allowbreak 3,\allowbreak 42,\allowbreak 8,\allowbreak 35,\allowbreak 12,\allowbreak 33,\allowbreak 13,\allowbreak 31,\allowbreak 17,\allowbreak 28,\allowbreak 18,\allowbreak 30,\allowbreak 15,\allowbreak 34)\par\smallskip (44,\allowbreak 7,\allowbreak 32,\allowbreak 16,\allowbreak 29,\allowbreak 20,\allowbreak 25,\allowbreak 23,\allowbreak 24,\allowbreak 21,\allowbreak 27,\allowbreak 19,\allowbreak 26,\allowbreak 22,\allowbreak 39,\allowbreak 10,\allowbreak 46,\allowbreak 6,\allowbreak 41,\allowbreak 11,\allowbreak 37,\allowbreak 9,\allowbreak 40)}\tabularnewline[5pt]
{\raggedright\arraybackslash {\ttfamily\footnotesize (4,\allowbreak 21,\allowbreak 24)}\par $K=24$\par pinned W4} & {\ttfamily\footnotesize (49,\allowbreak 1,\allowbreak 48,\allowbreak 2)\par\smallskip (43,\allowbreak 3,\allowbreak 47,\allowbreak 6,\allowbreak 42,\allowbreak 21,\allowbreak 27,\allowbreak 22,\allowbreak 26,\allowbreak 23,\allowbreak 25,\allowbreak 24,\allowbreak 31,\allowbreak 12,\allowbreak 37,\allowbreak 10,\allowbreak 36,\allowbreak 8,\allowbreak 39,\allowbreak 9,\allowbreak 41)\par\smallskip (45,\allowbreak 7,\allowbreak 44,\allowbreak 5,\allowbreak 40,\allowbreak 11,\allowbreak 35,\allowbreak 13,\allowbreak 46,\allowbreak 4,\allowbreak 38,\allowbreak 15,\allowbreak 30,\allowbreak 18,\allowbreak 32,\allowbreak 19,\allowbreak 29,\allowbreak 20,\allowbreak 28,\allowbreak 17,\allowbreak 33,\allowbreak 16,\allowbreak 34,\allowbreak 14)}\tabularnewline[5pt]
{\raggedright\arraybackslash {\ttfamily\footnotesize (4,\allowbreak 22,\allowbreak 22)}\par $K=24$\par pinned W4} & {\ttfamily\footnotesize (48,\allowbreak 1,\allowbreak 47,\allowbreak 2)\par\smallskip (40,\allowbreak 5,\allowbreak 36,\allowbreak 15,\allowbreak 35,\allowbreak 11,\allowbreak 37,\allowbreak 20,\allowbreak 30,\allowbreak 14,\allowbreak 39,\allowbreak 17,\allowbreak 31,\allowbreak 22,\allowbreak 27,\allowbreak 23,\allowbreak 26,\allowbreak 24,\allowbreak 25,\allowbreak 19,\allowbreak 32,\allowbreak 13)\par\smallskip (44,\allowbreak 6,\allowbreak 43,\allowbreak 7,\allowbreak 34,\allowbreak 16,\allowbreak 28,\allowbreak 21,\allowbreak 29,\allowbreak 18,\allowbreak 33,\allowbreak 10,\allowbreak 38,\allowbreak 8,\allowbreak 42,\allowbreak 3,\allowbreak 46,\allowbreak 4,\allowbreak 45,\allowbreak 12,\allowbreak 41,\allowbreak 9)}\tabularnewline[5pt]
{\raggedright\arraybackslash {\ttfamily\footnotesize (4,\allowbreak 22,\allowbreak 23)}\par $K=24$\par pinned W4\par $2e\le b+8$} & {\ttfamily\footnotesize (49,\allowbreak 1,\allowbreak 48,\allowbreak 2)\par\smallskip (44,\allowbreak 12,\allowbreak 30,\allowbreak 18,\allowbreak 34,\allowbreak 11,\allowbreak 41,\allowbreak 10,\allowbreak 35,\allowbreak 14,\allowbreak 33,\allowbreak 23,\allowbreak 26,\allowbreak 24,\allowbreak 25,\allowbreak 21,\allowbreak 27,\allowbreak 22,\allowbreak 29,\allowbreak 16,\allowbreak 38,\allowbreak 9)\par\smallskip (45,\allowbreak 6,\allowbreak 47,\allowbreak 7,\allowbreak 40,\allowbreak 4,\allowbreak 46,\allowbreak 3,\allowbreak 37,\allowbreak 13,\allowbreak 39,\allowbreak 19,\allowbreak 28,\allowbreak 20,\allowbreak 31,\allowbreak 17,\allowbreak 32,\allowbreak 15,\allowbreak 42,\allowbreak 5,\allowbreak 43,\allowbreak 8,\allowbreak 36)}\tabularnewline[5pt]
{\raggedright\arraybackslash {\ttfamily\footnotesize (4,\allowbreak 22,\allowbreak 24)}\par $K=25$\par pinned W4} & {\ttfamily\footnotesize (50,\allowbreak 1,\allowbreak 49,\allowbreak 2)\par\smallskip (40,\allowbreak 18,\allowbreak 35,\allowbreak 15,\allowbreak 34,\allowbreak 20,\allowbreak 30,\allowbreak 19,\allowbreak 32,\allowbreak 23,\allowbreak 31,\allowbreak 13,\allowbreak 37,\allowbreak 16,\allowbreak 39,\allowbreak 7,\allowbreak 41,\allowbreak 12,\allowbreak 42,\allowbreak 6,\allowbreak 44,\allowbreak 9)\par\smallskip (46,\allowbreak 3,\allowbreak 48,\allowbreak 4,\allowbreak 45,\allowbreak 8,\allowbreak 47,\allowbreak 5,\allowbreak 36,\allowbreak 11,\allowbreak 38,\allowbreak 10,\allowbreak 43,\allowbreak 17,\allowbreak 33,\allowbreak 21,\allowbreak 26,\allowbreak 25,\allowbreak 27,\allowbreak 24,\allowbreak 28,\allowbreak 22,\allowbreak 29,\allowbreak 14)}\tabularnewline[5pt]
{\raggedright\arraybackslash {\ttfamily\footnotesize (4,\allowbreak 22,\allowbreak 25)}\par $K=25$\par pinned W4} & {\ttfamily\footnotesize (51,\allowbreak 1,\allowbreak 50,\allowbreak 2)\par\smallskip (46,\allowbreak 3,\allowbreak 43,\allowbreak 10,\allowbreak 41,\allowbreak 15,\allowbreak 29,\allowbreak 20,\allowbreak 30,\allowbreak 24,\allowbreak 28,\allowbreak 21,\allowbreak 34,\allowbreak 18,\allowbreak 38,\allowbreak 16,\allowbreak 33,\allowbreak 22,\allowbreak 27,\allowbreak 25,\allowbreak 26,\allowbreak 23,\allowbreak 31,\allowbreak 19,\allowbreak 37)\par\smallskip (47,\allowbreak 17,\allowbreak 32,\allowbreak 13,\allowbreak 36,\allowbreak 12,\allowbreak 40,\allowbreak 11,\allowbreak 48,\allowbreak 6,\allowbreak 42,\allowbreak 8,\allowbreak 35,\allowbreak 14,\allowbreak 39,\allowbreak 7,\allowbreak 45,\allowbreak 4,\allowbreak 49,\allowbreak 5,\allowbreak 44,\allowbreak 9)}\tabularnewline[5pt]
{\raggedright\arraybackslash {\ttfamily\footnotesize (4,\allowbreak 22,\allowbreak 26)}\par $K=26$\par pinned W4} & {\ttfamily\footnotesize (52,\allowbreak 1,\allowbreak 51,\allowbreak 2)\par\smallskip (45,\allowbreak 4,\allowbreak 47,\allowbreak 7,\allowbreak 44,\allowbreak 17,\allowbreak 32,\allowbreak 21,\allowbreak 31,\allowbreak 22,\allowbreak 41,\allowbreak 5,\allowbreak 43,\allowbreak 12,\allowbreak 46,\allowbreak 11,\allowbreak 34,\allowbreak 18,\allowbreak 35,\allowbreak 14,\allowbreak 42,\allowbreak 9)\par\smallskip (48,\allowbreak 6,\allowbreak 50,\allowbreak 3,\allowbreak 49,\allowbreak 10,\allowbreak 40,\allowbreak 8,\allowbreak 37,\allowbreak 19,\allowbreak 33,\allowbreak 20,\allowbreak 28,\allowbreak 25,\allowbreak 27,\allowbreak 26,\allowbreak 30,\allowbreak 23,\allowbreak 29,\allowbreak 24,\allowbreak 36,\allowbreak 16,\allowbreak 38,\allowbreak 13,\allowbreak 39,\allowbreak 15)}\tabularnewline[5pt]
{\raggedright\arraybackslash {\ttfamily\footnotesize (4,\allowbreak 23,\allowbreak 23)}\par $K=24$\par pinned W4\par $2e\le b+7$} & {\ttfamily\footnotesize (50,\allowbreak 1,\allowbreak 49,\allowbreak 2)\par\smallskip (45,\allowbreak 11,\allowbreak 35,\allowbreak 16,\allowbreak 32,\allowbreak 19,\allowbreak 29,\allowbreak 21,\allowbreak 30,\allowbreak 18,\allowbreak 33,\allowbreak 15,\allowbreak 38,\allowbreak 7,\allowbreak 39,\allowbreak 12,\allowbreak 37,\allowbreak 8,\allowbreak 43,\allowbreak 10,\allowbreak 40,\allowbreak 14,\allowbreak 36)\par\smallskip (46,\allowbreak 3,\allowbreak 44,\allowbreak 4,\allowbreak 48,\allowbreak 9,\allowbreak 47,\allowbreak 5,\allowbreak 42,\allowbreak 6,\allowbreak 26,\allowbreak 24,\allowbreak 25,\allowbreak 22,\allowbreak 28,\allowbreak 23,\allowbreak 27,\allowbreak 20,\allowbreak 31,\allowbreak 17,\allowbreak 34,\allowbreak 13,\allowbreak 41)}\tabularnewline[5pt]
{\raggedright\arraybackslash {\ttfamily\footnotesize (4,\allowbreak 23,\allowbreak 24)}\par $K=25$\par pinned W4} & {\ttfamily\footnotesize (51,\allowbreak 1,\allowbreak 50,\allowbreak 2)\par\smallskip (43,\allowbreak 5,\allowbreak 49,\allowbreak 3,\allowbreak 48,\allowbreak 8,\allowbreak 36,\allowbreak 12,\allowbreak 35,\allowbreak 15,\allowbreak 34,\allowbreak 17,\allowbreak 33,\allowbreak 21,\allowbreak 32,\allowbreak 22,\allowbreak 30,\allowbreak 16,\allowbreak 31,\allowbreak 18,\allowbreak 40,\allowbreak 10,\allowbreak 42)\par\smallskip (47,\allowbreak 13,\allowbreak 44,\allowbreak 7,\allowbreak 46,\allowbreak 4,\allowbreak 45,\allowbreak 9,\allowbreak 38,\allowbreak 11,\allowbreak 37,\allowbreak 19,\allowbreak 28,\allowbreak 24,\allowbreak 29,\allowbreak 23,\allowbreak 26,\allowbreak 25,\allowbreak 27,\allowbreak 20,\allowbreak 41,\allowbreak 6,\allowbreak 39,\allowbreak 14)}\tabularnewline[5pt]
{\raggedright\arraybackslash {\ttfamily\footnotesize (4,\allowbreak 23,\allowbreak 25)}\par $K=25$\par pinned W4} & {\ttfamily\footnotesize (48,\allowbreak 5,\allowbreak 47,\allowbreak 6,\allowbreak 50,\allowbreak 3,\allowbreak 49,\allowbreak 4,\allowbreak 34,\allowbreak 19,\allowbreak 33,\allowbreak 21,\allowbreak 32,\allowbreak 23,\allowbreak 28,\allowbreak 24,\allowbreak 27,\allowbreak 25,\allowbreak 26,\allowbreak 20,\allowbreak 30,\allowbreak 22,\allowbreak 29,\allowbreak 13,\allowbreak 37)\par\smallskip (40,\allowbreak 15,\allowbreak 38,\allowbreak 17,\allowbreak 35,\allowbreak 16,\allowbreak 36,\allowbreak 14,\allowbreak 31,\allowbreak 18,\allowbreak 44,\allowbreak 8,\allowbreak 45,\allowbreak 7,\allowbreak 46,\allowbreak 11,\allowbreak 39,\allowbreak 12,\allowbreak 41,\allowbreak 10,\allowbreak 42,\allowbreak 9,\allowbreak 43)\par\smallskip (52,\allowbreak 1,\allowbreak 51,\allowbreak 2)}\tabularnewline[5pt]
{\raggedright\arraybackslash {\ttfamily\footnotesize (4,\allowbreak 23,\allowbreak 26)}\par $K=26$\par pinned W4} & {\ttfamily\footnotesize (53,\allowbreak 1,\allowbreak 52,\allowbreak 2)\par\smallskip (34,\allowbreak 19,\allowbreak 31,\allowbreak 20,\allowbreak 38,\allowbreak 13,\allowbreak 43,\allowbreak 12,\allowbreak 40,\allowbreak 11,\allowbreak 37,\allowbreak 17,\allowbreak 33,\allowbreak 23,\allowbreak 36,\allowbreak 14,\allowbreak 41,\allowbreak 8,\allowbreak 46,\allowbreak 5,\allowbreak 47,\allowbreak 7,\allowbreak 44)\par\smallskip (49,\allowbreak 6,\allowbreak 50,\allowbreak 4,\allowbreak 51,\allowbreak 3,\allowbreak 48,\allowbreak 9,\allowbreak 45,\allowbreak 10,\allowbreak 42,\allowbreak 21,\allowbreak 30,\allowbreak 25,\allowbreak 27,\allowbreak 26,\allowbreak 29,\allowbreak 22,\allowbreak 28,\allowbreak 24,\allowbreak 32,\allowbreak 18,\allowbreak 35,\allowbreak 16,\allowbreak 39,\allowbreak 15)}\tabularnewline[5pt]
{\raggedright\arraybackslash {\ttfamily\footnotesize (4,\allowbreak 24,\allowbreak 24)}\par $K=26$\par pinned W4} & {\ttfamily\footnotesize (48,\allowbreak 5,\allowbreak 47,\allowbreak 6,\allowbreak 50,\allowbreak 3,\allowbreak 49,\allowbreak 4,\allowbreak 32,\allowbreak 21,\allowbreak 33,\allowbreak 23,\allowbreak 29,\allowbreak 24,\allowbreak 31,\allowbreak 22,\allowbreak 30,\allowbreak 26,\allowbreak 27,\allowbreak 25,\allowbreak 28,\allowbreak 15,\allowbreak 38,\allowbreak 14)\par\smallskip (40,\allowbreak 13,\allowbreak 39,\allowbreak 17,\allowbreak 35,\allowbreak 18,\allowbreak 37,\allowbreak 16,\allowbreak 36,\allowbreak 20,\allowbreak 34,\allowbreak 19,\allowbreak 44,\allowbreak 8,\allowbreak 45,\allowbreak 7,\allowbreak 46,\allowbreak 11,\allowbreak 41,\allowbreak 12,\allowbreak 43,\allowbreak 9,\allowbreak 42,\allowbreak 10)\par\smallskip (52,\allowbreak 1,\allowbreak 51,\allowbreak 2)}\tabularnewline[5pt]
{\raggedright\arraybackslash {\ttfamily\footnotesize (4,\allowbreak 24,\allowbreak 25)}\par $K=26$\par pinned W4\par $2e\le b+8$} & {\ttfamily\footnotesize (41,\allowbreak 13,\allowbreak 40,\allowbreak 14,\allowbreak 35,\allowbreak 17,\allowbreak 37,\allowbreak 18,\allowbreak 34,\allowbreak 23,\allowbreak 32,\allowbreak 20,\allowbreak 45,\allowbreak 8,\allowbreak 42,\allowbreak 12,\allowbreak 43,\allowbreak 11,\allowbreak 44,\allowbreak 9,\allowbreak 47,\allowbreak 7,\allowbreak 46,\allowbreak 10)\par\smallskip (53,\allowbreak 1,\allowbreak 52,\allowbreak 2)\par\smallskip (49,\allowbreak 5,\allowbreak 48,\allowbreak 6,\allowbreak 51,\allowbreak 3,\allowbreak 50,\allowbreak 4,\allowbreak 33,\allowbreak 19,\allowbreak 36,\allowbreak 21,\allowbreak 31,\allowbreak 24,\allowbreak 30,\allowbreak 22,\allowbreak 27,\allowbreak 26,\allowbreak 28,\allowbreak 25,\allowbreak 29,\allowbreak 16,\allowbreak 38,\allowbreak 15,\allowbreak 39)}\tabularnewline[5pt]
{\raggedright\arraybackslash {\ttfamily\footnotesize (4,\allowbreak 24,\allowbreak 26)}\par $K=27$\par pinned W4} & {\ttfamily\footnotesize (42,\allowbreak 12,\allowbreak 47,\allowbreak 8,\allowbreak 48,\allowbreak 7,\allowbreak 45,\allowbreak 9,\allowbreak 46,\allowbreak 23,\allowbreak 32,\allowbreak 24,\allowbreak 31,\allowbreak 21,\allowbreak 34,\allowbreak 22,\allowbreak 33,\allowbreak 27,\allowbreak 28,\allowbreak 26,\allowbreak 29,\allowbreak 25,\allowbreak 30,\allowbreak 16,\allowbreak 40,\allowbreak 15)\par\smallskip (54,\allowbreak 1,\allowbreak 53,\allowbreak 2)\par\smallskip (50,\allowbreak 5,\allowbreak 49,\allowbreak 6,\allowbreak 52,\allowbreak 3,\allowbreak 51,\allowbreak 4,\allowbreak 37,\allowbreak 18,\allowbreak 38,\allowbreak 17,\allowbreak 35,\allowbreak 20,\allowbreak 36,\allowbreak 19,\allowbreak 41,\allowbreak 14,\allowbreak 43,\allowbreak 11,\allowbreak 39,\allowbreak 13,\allowbreak 44,\allowbreak 10)}\tabularnewline[5pt]
{\raggedright\arraybackslash {\ttfamily\footnotesize (4,\allowbreak 24,\allowbreak 27)}\par $K=27$\par pinned W4} & {\ttfamily\footnotesize (51,\allowbreak 5,\allowbreak 50,\allowbreak 6,\allowbreak 53,\allowbreak 3,\allowbreak 52,\allowbreak 4,\allowbreak 37,\allowbreak 22,\allowbreak 35,\allowbreak 21,\allowbreak 28,\allowbreak 27,\allowbreak 29,\allowbreak 26,\allowbreak 30,\allowbreak 25,\allowbreak 31,\allowbreak 13,\allowbreak 39,\allowbreak 17,\allowbreak 38,\allowbreak 18,\allowbreak 41,\allowbreak 16,\allowbreak 40)\par\smallskip (43,\allowbreak 14,\allowbreak 45,\allowbreak 9,\allowbreak 46,\allowbreak 11,\allowbreak 49,\allowbreak 7,\allowbreak 48,\allowbreak 8,\allowbreak 47,\allowbreak 19,\allowbreak 36,\allowbreak 20,\allowbreak 32,\allowbreak 24,\allowbreak 33,\allowbreak 23,\allowbreak 34,\allowbreak 15,\allowbreak 42,\allowbreak 12,\allowbreak 44,\allowbreak 10)\par\smallskip (55,\allowbreak 1,\allowbreak 54,\allowbreak 2)}\tabularnewline[5pt]
{\raggedright\arraybackslash {\ttfamily\footnotesize (4,\allowbreak 24,\allowbreak 28)}\par $K=28$\par pinned W4} & {\ttfamily\footnotesize (44,\allowbreak 15,\allowbreak 41,\allowbreak 14,\allowbreak 42,\allowbreak 22,\allowbreak 33,\allowbreak 21,\allowbreak 36,\allowbreak 23,\allowbreak 48,\allowbreak 12,\allowbreak 43,\allowbreak 13,\allowbreak 45,\allowbreak 11,\allowbreak 46,\allowbreak 9,\allowbreak 47,\allowbreak 8,\allowbreak 49,\allowbreak 7,\allowbreak 50,\allowbreak 10)\par\smallskip (56,\allowbreak 1,\allowbreak 55,\allowbreak 2)\par\smallskip (52,\allowbreak 5,\allowbreak 51,\allowbreak 6,\allowbreak 54,\allowbreak 3,\allowbreak 53,\allowbreak 4,\allowbreak 37,\allowbreak 20,\allowbreak 38,\allowbreak 19,\allowbreak 35,\allowbreak 28,\allowbreak 29,\allowbreak 27,\allowbreak 30,\allowbreak 26,\allowbreak 31,\allowbreak 25,\allowbreak 34,\allowbreak 24,\allowbreak 32,\allowbreak 18,\allowbreak 39,\allowbreak 17,\allowbreak 40,\allowbreak 16)}\tabularnewline[5pt]
{\raggedright\arraybackslash {\ttfamily\footnotesize (4,\allowbreak 25,\allowbreak 25)}\par $K=26$\par pinned W4\par $2e\le b+7$} & {\ttfamily\footnotesize (42,\allowbreak 12,\allowbreak 40,\allowbreak 13,\allowbreak 29,\allowbreak 25,\allowbreak 28,\allowbreak 26,\allowbreak 27,\allowbreak 22,\allowbreak 30,\allowbreak 24,\allowbreak 31,\allowbreak 21,\allowbreak 33,\allowbreak 20,\allowbreak 46,\allowbreak 7,\allowbreak 48,\allowbreak 8,\allowbreak 45,\allowbreak 11,\allowbreak 47,\allowbreak 9,\allowbreak 44)\par\smallskip (54,\allowbreak 1,\allowbreak 53,\allowbreak 2)\par\smallskip (50,\allowbreak 5,\allowbreak 49,\allowbreak 6,\allowbreak 52,\allowbreak 3,\allowbreak 51,\allowbreak 4,\allowbreak 36,\allowbreak 17,\allowbreak 37,\allowbreak 19,\allowbreak 34,\allowbreak 23,\allowbreak 32,\allowbreak 18,\allowbreak 35,\allowbreak 14,\allowbreak 43,\allowbreak 10,\allowbreak 41,\allowbreak 16,\allowbreak 38,\allowbreak 15,\allowbreak 39)}\tabularnewline[5pt]
{\raggedright\arraybackslash {\ttfamily\footnotesize (4,\allowbreak 25,\allowbreak 26)}\par $K=27$\par pinned W4} & {\ttfamily\footnotesize (51,\allowbreak 5,\allowbreak 50,\allowbreak 6,\allowbreak 53,\allowbreak 3,\allowbreak 52,\allowbreak 4,\allowbreak 37,\allowbreak 18,\allowbreak 38,\allowbreak 17,\allowbreak 35,\allowbreak 23,\allowbreak 32,\allowbreak 24,\allowbreak 31,\allowbreak 21,\allowbreak 34,\allowbreak 20,\allowbreak 36,\allowbreak 19,\allowbreak 41,\allowbreak 13,\allowbreak 45)\par\smallskip (43,\allowbreak 14,\allowbreak 40,\allowbreak 16,\allowbreak 39,\allowbreak 12,\allowbreak 42,\allowbreak 11,\allowbreak 48,\allowbreak 7,\allowbreak 49,\allowbreak 9,\allowbreak 44,\allowbreak 10,\allowbreak 46,\allowbreak 8,\allowbreak 47,\allowbreak 22,\allowbreak 33,\allowbreak 27,\allowbreak 28,\allowbreak 26,\allowbreak 29,\allowbreak 25,\allowbreak 30,\allowbreak 15)\par\smallskip (55,\allowbreak 1,\allowbreak 54,\allowbreak 2)}\tabularnewline[5pt]
{\raggedright\arraybackslash {\ttfamily\footnotesize (4,\allowbreak 25,\allowbreak 27)}\par $K=27$\par pinned W4} & {\ttfamily\footnotesize (56,\allowbreak 1,\allowbreak 55,\allowbreak 2)\par\smallskip (43,\allowbreak 14,\allowbreak 41,\allowbreak 13,\allowbreak 31,\allowbreak 25,\allowbreak 30,\allowbreak 26,\allowbreak 29,\allowbreak 27,\allowbreak 28,\allowbreak 21,\allowbreak 32,\allowbreak 24,\allowbreak 33,\allowbreak 23,\allowbreak 36,\allowbreak 22,\allowbreak 34,\allowbreak 15,\allowbreak 38,\allowbreak 18,\allowbreak 39,\allowbreak 17,\allowbreak 42,\allowbreak 16,\allowbreak 40)\par\smallskip (52,\allowbreak 5,\allowbreak 51,\allowbreak 6,\allowbreak 54,\allowbreak 3,\allowbreak 53,\allowbreak 4,\allowbreak 48,\allowbreak 12,\allowbreak 44,\allowbreak 11,\allowbreak 45,\allowbreak 10,\allowbreak 49,\allowbreak 8,\allowbreak 50,\allowbreak 19,\allowbreak 35,\allowbreak 20,\allowbreak 37,\allowbreak 7,\allowbreak 47,\allowbreak 9,\allowbreak 46)}\tabularnewline[5pt]
{\raggedright\arraybackslash {\ttfamily\footnotesize (4,\allowbreak 25,\allowbreak 28)}\par $K=28$\par pinned W4} & {\ttfamily\footnotesize (57,\allowbreak 1,\allowbreak 56,\allowbreak 2)\par\smallskip (51,\allowbreak 9,\allowbreak 46,\allowbreak 15,\allowbreak 43,\allowbreak 18,\allowbreak 34,\allowbreak 26,\allowbreak 31,\allowbreak 24,\allowbreak 33,\allowbreak 27,\allowbreak 30,\allowbreak 28,\allowbreak 29,\allowbreak 25,\allowbreak 35,\allowbreak 22,\allowbreak 36,\allowbreak 4,\allowbreak 54,\allowbreak 6,\allowbreak 45,\allowbreak 12,\allowbreak 47)\par\smallskip (53,\allowbreak 7,\allowbreak 50,\allowbreak 5,\allowbreak 49,\allowbreak 13,\allowbreak 42,\allowbreak 19,\allowbreak 39,\allowbreak 17,\allowbreak 38,\allowbreak 23,\allowbreak 41,\allowbreak 14,\allowbreak 44,\allowbreak 10,\allowbreak 48,\allowbreak 8,\allowbreak 55,\allowbreak 3,\allowbreak 52,\allowbreak 11,\allowbreak 37,\allowbreak 20,\allowbreak 32,\allowbreak 21,\allowbreak 40,\allowbreak 16)}\tabularnewline[5pt]
{\raggedright\arraybackslash {\ttfamily\footnotesize (4,\allowbreak 26,\allowbreak 26)}\par $K=28$\par pinned W4} & {\ttfamily\footnotesize (56,\allowbreak 1,\allowbreak 55,\allowbreak 2)\par\smallskip (52,\allowbreak 3,\allowbreak 54,\allowbreak 4,\allowbreak 51,\allowbreak 5,\allowbreak 53,\allowbreak 8,\allowbreak 50,\allowbreak 6,\allowbreak 49,\allowbreak 9,\allowbreak 48,\allowbreak 7,\allowbreak 45,\allowbreak 10,\allowbreak 47,\allowbreak 11,\allowbreak 44,\allowbreak 13,\allowbreak 43,\allowbreak 15,\allowbreak 42,\allowbreak 16,\allowbreak 41,\allowbreak 12)\par\smallskip (34,\allowbreak 17,\allowbreak 40,\allowbreak 18,\allowbreak 39,\allowbreak 19,\allowbreak 38,\allowbreak 20,\allowbreak 36,\allowbreak 21,\allowbreak 35,\allowbreak 23,\allowbreak 33,\allowbreak 24,\allowbreak 37,\allowbreak 26,\allowbreak 31,\allowbreak 25,\allowbreak 32,\allowbreak 28,\allowbreak 29,\allowbreak 27,\allowbreak 30,\allowbreak 22,\allowbreak 46,\allowbreak 14)}\tabularnewline[5pt]
{\raggedright\arraybackslash {\ttfamily\footnotesize (4,\allowbreak 26,\allowbreak 27)}\par $K=28$\par pinned W4\par $2e\le b+8$} & {\ttfamily\footnotesize (57,\allowbreak 1,\allowbreak 56,\allowbreak 2)\par\smallskip (52,\allowbreak 7,\allowbreak 51,\allowbreak 3,\allowbreak 54,\allowbreak 5,\allowbreak 55,\allowbreak 8,\allowbreak 47,\allowbreak 10,\allowbreak 48,\allowbreak 6,\allowbreak 49,\allowbreak 14,\allowbreak 40,\allowbreak 17,\allowbreak 36,\allowbreak 4,\allowbreak 50,\allowbreak 9,\allowbreak 45,\allowbreak 18,\allowbreak 32,\allowbreak 25,\allowbreak 34,\allowbreak 21,\allowbreak 42)\par\smallskip (53,\allowbreak 13,\allowbreak 41,\allowbreak 16,\allowbreak 46,\allowbreak 12,\allowbreak 43,\allowbreak 19,\allowbreak 39,\allowbreak 24,\allowbreak 30,\allowbreak 28,\allowbreak 29,\allowbreak 26,\allowbreak 31,\allowbreak 27,\allowbreak 35,\allowbreak 23,\allowbreak 33,\allowbreak 22,\allowbreak 38,\allowbreak 20,\allowbreak 37,\allowbreak 15,\allowbreak 44,\allowbreak 11)}\tabularnewline[5pt]
{\raggedright\arraybackslash {\ttfamily\footnotesize (4,\allowbreak 26,\allowbreak 28)}\par $K=29$\par pinned W4} & {\ttfamily\footnotesize (58,\allowbreak 1,\allowbreak 57,\allowbreak 2)\par\smallskip (53,\allowbreak 13,\allowbreak 44,\allowbreak 18,\allowbreak 40,\allowbreak 21,\allowbreak 37,\allowbreak 24,\allowbreak 33,\allowbreak 28,\allowbreak 32,\allowbreak 26,\allowbreak 34,\allowbreak 27,\allowbreak 30,\allowbreak 29,\allowbreak 31,\allowbreak 7,\allowbreak 56,\allowbreak 9,\allowbreak 52,\allowbreak 4,\allowbreak 48,\allowbreak 14,\allowbreak 43,\allowbreak 16)\par\smallskip (54,\allowbreak 3,\allowbreak 55,\allowbreak 5,\allowbreak 51,\allowbreak 6,\allowbreak 47,\allowbreak 15,\allowbreak 45,\allowbreak 17,\allowbreak 42,\allowbreak 19,\allowbreak 39,\allowbreak 22,\allowbreak 36,\allowbreak 25,\allowbreak 35,\allowbreak 23,\allowbreak 38,\allowbreak 20,\allowbreak 41,\allowbreak 8,\allowbreak 50,\allowbreak 12,\allowbreak 49,\allowbreak 10,\allowbreak 46,\allowbreak 11)}\tabularnewline[5pt]
{\raggedright\arraybackslash {\ttfamily\footnotesize (4,\allowbreak 26,\allowbreak 29)}\par $K=29$\par pinned W4} & {\ttfamily\footnotesize (59,\allowbreak 1,\allowbreak 58,\allowbreak 2)\par\smallskip (53,\allowbreak 3,\allowbreak 57,\allowbreak 12,\allowbreak 47,\allowbreak 6,\allowbreak 54,\allowbreak 10,\allowbreak 48,\allowbreak 8,\allowbreak 51,\allowbreak 5,\allowbreak 52,\allowbreak 16,\allowbreak 40,\allowbreak 24,\allowbreak 33,\allowbreak 25,\allowbreak 32,\allowbreak 28,\allowbreak 34,\allowbreak 21,\allowbreak 44,\allowbreak 13,\allowbreak 46,\allowbreak 9)\par\smallskip (55,\allowbreak 4,\allowbreak 56,\allowbreak 7,\allowbreak 49,\allowbreak 15,\allowbreak 45,\allowbreak 18,\allowbreak 50,\allowbreak 11,\allowbreak 36,\allowbreak 26,\allowbreak 31,\allowbreak 29,\allowbreak 30,\allowbreak 27,\allowbreak 38,\allowbreak 19,\allowbreak 39,\allowbreak 22,\allowbreak 37,\allowbreak 23,\allowbreak 35,\allowbreak 17,\allowbreak 43,\allowbreak 14,\allowbreak 42,\allowbreak 20,\allowbreak 41)}\tabularnewline[5pt]
{\raggedright\arraybackslash {\ttfamily\footnotesize (4,\allowbreak 26,\allowbreak 30)}\par $K=30$\par pinned W4} & {\ttfamily\footnotesize (60,\allowbreak 1,\allowbreak 59,\allowbreak 2)\par\smallskip (53,\allowbreak 7,\allowbreak 57,\allowbreak 5,\allowbreak 52,\allowbreak 8,\allowbreak 49,\allowbreak 10,\allowbreak 50,\allowbreak 12,\allowbreak 47,\allowbreak 22,\allowbreak 36,\allowbreak 27,\allowbreak 31,\allowbreak 30,\allowbreak 32,\allowbreak 29,\allowbreak 35,\allowbreak 28,\allowbreak 33,\allowbreak 25,\allowbreak 44,\allowbreak 16,\allowbreak 48,\allowbreak 11)\par\smallskip (56,\allowbreak 13,\allowbreak 46,\allowbreak 20,\allowbreak 42,\allowbreak 15,\allowbreak 45,\allowbreak 14,\allowbreak 43,\allowbreak 23,\allowbreak 39,\allowbreak 21,\allowbreak 38,\allowbreak 26,\allowbreak 37,\allowbreak 24,\allowbreak 34,\allowbreak 19,\allowbreak 40,\allowbreak 6,\allowbreak 55,\allowbreak 4,\allowbreak 58,\allowbreak 3,\allowbreak 51,\allowbreak 9,\allowbreak 54,\allowbreak 18,\allowbreak 41,\allowbreak 17)}\tabularnewline[5pt]
{\raggedright\arraybackslash {\ttfamily\footnotesize (4,\allowbreak 27,\allowbreak 27)}\par $K=28$\par pinned W4\par $2e\le b+7$} & {\ttfamily\footnotesize (58,\allowbreak 1,\allowbreak 57,\allowbreak 2)\par\smallskip (40,\allowbreak 25,\allowbreak 32,\allowbreak 27,\allowbreak 31,\allowbreak 23,\allowbreak 35,\allowbreak 22,\allowbreak 33,\allowbreak 24,\allowbreak 30,\allowbreak 28,\allowbreak 29,\allowbreak 26,\allowbreak 36,\allowbreak 20,\allowbreak 34,\allowbreak 9,\allowbreak 53,\allowbreak 7,\allowbreak 50,\allowbreak 8,\allowbreak 49,\allowbreak 11,\allowbreak 48,\allowbreak 15,\allowbreak 47)\par\smallskip (54,\allowbreak 4,\allowbreak 52,\allowbreak 5,\allowbreak 56,\allowbreak 3,\allowbreak 55,\allowbreak 6,\allowbreak 51,\allowbreak 12,\allowbreak 46,\allowbreak 10,\allowbreak 45,\allowbreak 16,\allowbreak 43,\allowbreak 21,\allowbreak 38,\allowbreak 17,\allowbreak 41,\allowbreak 18,\allowbreak 37,\allowbreak 19,\allowbreak 39,\allowbreak 13,\allowbreak 44,\allowbreak 14,\allowbreak 42)}\tabularnewline[5pt]
{\raggedright\arraybackslash {\ttfamily\footnotesize (4,\allowbreak 27,\allowbreak 28)}\par $K=29$\par pinned W4} & {\ttfamily\footnotesize (59,\allowbreak 1,\allowbreak 58,\allowbreak 2)\par\smallskip (54,\allowbreak 6,\allowbreak 56,\allowbreak 10,\allowbreak 53,\allowbreak 8,\allowbreak 50,\allowbreak 9,\allowbreak 49,\allowbreak 12,\allowbreak 47,\allowbreak 13,\allowbreak 45,\allowbreak 18,\allowbreak 40,\allowbreak 14,\allowbreak 43,\allowbreak 23,\allowbreak 37,\allowbreak 25,\allowbreak 34,\allowbreak 26,\allowbreak 36,\allowbreak 15,\allowbreak 51,\allowbreak 7,\allowbreak 46,\allowbreak 16)\par\smallskip (55,\allowbreak 4,\allowbreak 57,\allowbreak 5,\allowbreak 52,\allowbreak 3,\allowbreak 41,\allowbreak 24,\allowbreak 31,\allowbreak 27,\allowbreak 33,\allowbreak 28,\allowbreak 30,\allowbreak 29,\allowbreak 32,\allowbreak 21,\allowbreak 39,\allowbreak 11,\allowbreak 44,\allowbreak 20,\allowbreak 35,\allowbreak 22,\allowbreak 38,\allowbreak 19,\allowbreak 42,\allowbreak 17,\allowbreak 48)}\tabularnewline[5pt]
{\raggedright\arraybackslash {\ttfamily\footnotesize (4,\allowbreak 27,\allowbreak 29)}\par $K=29$\par pinned W4} & {\ttfamily\footnotesize (60,\allowbreak 1,\allowbreak 59,\allowbreak 2)\par\smallskip (55,\allowbreak 7,\allowbreak 45,\allowbreak 17,\allowbreak 40,\allowbreak 21,\allowbreak 35,\allowbreak 23,\allowbreak 36,\allowbreak 25,\allowbreak 33,\allowbreak 16,\allowbreak 57,\allowbreak 3,\allowbreak 54,\allowbreak 10,\allowbreak 50,\allowbreak 4,\allowbreak 53,\allowbreak 8,\allowbreak 42,\allowbreak 18,\allowbreak 39,\allowbreak 19,\allowbreak 41,\allowbreak 15,\allowbreak 44,\allowbreak 13,\allowbreak 43)\par\smallskip (56,\allowbreak 6,\allowbreak 58,\allowbreak 5,\allowbreak 52,\allowbreak 9,\allowbreak 51,\allowbreak 12,\allowbreak 48,\allowbreak 11,\allowbreak 46,\allowbreak 14,\allowbreak 47,\allowbreak 20,\allowbreak 38,\allowbreak 22,\allowbreak 37,\allowbreak 28,\allowbreak 30,\allowbreak 29,\allowbreak 32,\allowbreak 26,\allowbreak 31,\allowbreak 27,\allowbreak 34,\allowbreak 24,\allowbreak 49)}\tabularnewline[5pt]
{\raggedright\arraybackslash {\ttfamily\footnotesize (4,\allowbreak 27,\allowbreak 30)}\par $K=30$\par pinned W4} & {\ttfamily\footnotesize (61,\allowbreak 1,\allowbreak 60,\allowbreak 2)\par\smallskip (56,\allowbreak 8,\allowbreak 48,\allowbreak 13,\allowbreak 49,\allowbreak 15,\allowbreak 47,\allowbreak 14,\allowbreak 44,\allowbreak 21,\allowbreak 40,\allowbreak 20,\allowbreak 38,\allowbreak 26,\allowbreak 36,\allowbreak 23,\allowbreak 37,\allowbreak 28,\allowbreak 33,\allowbreak 27,\allowbreak 35,\allowbreak 24,\allowbreak 39,\allowbreak 22,\allowbreak 43,\allowbreak 18,\allowbreak 42,\allowbreak 16,\allowbreak 45,\allowbreak 17)\par\smallskip (57,\allowbreak 6,\allowbreak 59,\allowbreak 4,\allowbreak 58,\allowbreak 11,\allowbreak 52,\allowbreak 7,\allowbreak 53,\allowbreak 10,\allowbreak 54,\allowbreak 5,\allowbreak 55,\allowbreak 3,\allowbreak 34,\allowbreak 30,\allowbreak 31,\allowbreak 29,\allowbreak 32,\allowbreak 25,\allowbreak 41,\allowbreak 19,\allowbreak 46,\allowbreak 9,\allowbreak 51,\allowbreak 12,\allowbreak 50)}\tabularnewline[5pt]
{\raggedright\arraybackslash {\ttfamily\footnotesize (4,\allowbreak 28,\allowbreak 28)}\par $K=30$\par pinned W4} & {\ttfamily\footnotesize (48,\allowbreak 10,\allowbreak 49,\allowbreak 14,\allowbreak 45,\allowbreak 19,\allowbreak 43,\allowbreak 18,\allowbreak 41,\allowbreak 20,\allowbreak 42,\allowbreak 22,\allowbreak 38,\allowbreak 25,\allowbreak 37,\allowbreak 23,\allowbreak 52,\allowbreak 7,\allowbreak 54,\allowbreak 8,\allowbreak 51,\allowbreak 11,\allowbreak 53,\allowbreak 9,\allowbreak 50,\allowbreak 13,\allowbreak 46,\allowbreak 12)\par\smallskip (60,\allowbreak 1,\allowbreak 59,\allowbreak 2)\par\smallskip (56,\allowbreak 5,\allowbreak 55,\allowbreak 6,\allowbreak 58,\allowbreak 3,\allowbreak 57,\allowbreak 4,\allowbreak 40,\allowbreak 21,\allowbreak 39,\allowbreak 24,\allowbreak 35,\allowbreak 26,\allowbreak 36,\allowbreak 28,\allowbreak 33,\allowbreak 27,\allowbreak 34,\allowbreak 30,\allowbreak 31,\allowbreak 29,\allowbreak 32,\allowbreak 15,\allowbreak 47,\allowbreak 17,\allowbreak 44,\allowbreak 16)}\tabularnewline[5pt]
{\raggedright\arraybackslash {\ttfamily\footnotesize (4,\allowbreak 28,\allowbreak 29)}\par $K=30$\par pinned W4\par $2e\le b+8$} & {\ttfamily\footnotesize (61,\allowbreak 1,\allowbreak 60,\allowbreak 2)\par\smallskip (54,\allowbreak 8,\allowbreak 55,\allowbreak 5,\allowbreak 48,\allowbreak 16,\allowbreak 42,\allowbreak 24,\allowbreak 39,\allowbreak 20,\allowbreak 43,\allowbreak 22,\allowbreak 36,\allowbreak 28,\allowbreak 32,\allowbreak 27,\allowbreak 34,\allowbreak 25,\allowbreak 38,\allowbreak 26,\allowbreak 37,\allowbreak 17,\allowbreak 41,\allowbreak 14,\allowbreak 52,\allowbreak 10,\allowbreak 51,\allowbreak 12)\par\smallskip (57,\allowbreak 9,\allowbreak 46,\allowbreak 21,\allowbreak 49,\allowbreak 13,\allowbreak 47,\allowbreak 7,\allowbreak 58,\allowbreak 6,\allowbreak 50,\allowbreak 19,\allowbreak 35,\allowbreak 29,\allowbreak 31,\allowbreak 30,\allowbreak 33,\allowbreak 23,\allowbreak 40,\allowbreak 18,\allowbreak 53,\allowbreak 4,\allowbreak 59,\allowbreak 3,\allowbreak 56,\allowbreak 11,\allowbreak 44,\allowbreak 15,\allowbreak 45)}\tabularnewline[5pt]
{\raggedright\arraybackslash {\ttfamily\footnotesize (4,\allowbreak 28,\allowbreak 30)}\par $K=31$\par pinned W4} & {\ttfamily\footnotesize (62,\allowbreak 1,\allowbreak 61,\allowbreak 2)\par\smallskip (57,\allowbreak 13,\allowbreak 47,\allowbreak 22,\allowbreak 38,\allowbreak 28,\allowbreak 37,\allowbreak 25,\allowbreak 39,\allowbreak 15,\allowbreak 53,\allowbreak 10,\allowbreak 55,\allowbreak 4,\allowbreak 60,\allowbreak 5,\allowbreak 59,\allowbreak 9,\allowbreak 46,\allowbreak 11,\allowbreak 44,\allowbreak 24,\allowbreak 35,\allowbreak 30,\allowbreak 34,\allowbreak 26,\allowbreak 43,\allowbreak 17)\par\smallskip (58,\allowbreak 6,\allowbreak 52,\allowbreak 3,\allowbreak 56,\allowbreak 8,\allowbreak 49,\allowbreak 7,\allowbreak 54,\allowbreak 14,\allowbreak 50,\allowbreak 19,\allowbreak 42,\allowbreak 23,\allowbreak 41,\allowbreak 20,\allowbreak 48,\allowbreak 16,\allowbreak 45,\allowbreak 18,\allowbreak 40,\allowbreak 27,\allowbreak 33,\allowbreak 31,\allowbreak 32,\allowbreak 29,\allowbreak 36,\allowbreak 21,\allowbreak 51,\allowbreak 12)}\tabularnewline[5pt]
{\raggedright\arraybackslash {\ttfamily\footnotesize (4,\allowbreak 28,\allowbreak 31)}\par $K=31$\par pinned W4} & {\ttfamily\footnotesize (48,\allowbreak 16,\allowbreak 45,\allowbreak 20,\allowbreak 44,\allowbreak 18,\allowbreak 49,\allowbreak 15,\allowbreak 35,\allowbreak 29,\allowbreak 33,\allowbreak 30,\allowbreak 32,\allowbreak 31,\allowbreak 36,\allowbreak 24,\allowbreak 42,\allowbreak 26,\allowbreak 34,\allowbreak 27,\allowbreak 37,\allowbreak 28,\allowbreak 39,\allowbreak 25,\allowbreak 38,\allowbreak 23,\allowbreak 40,\allowbreak 17,\allowbreak 47,\allowbreak 19,\allowbreak 46)\par\smallskip (63,\allowbreak 1,\allowbreak 62,\allowbreak 2)\par\smallskip (59,\allowbreak 5,\allowbreak 58,\allowbreak 6,\allowbreak 61,\allowbreak 3,\allowbreak 60,\allowbreak 4,\allowbreak 55,\allowbreak 11,\allowbreak 53,\allowbreak 10,\allowbreak 57,\allowbreak 7,\allowbreak 56,\allowbreak 21,\allowbreak 43,\allowbreak 22,\allowbreak 41,\allowbreak 8,\allowbreak 54,\allowbreak 9,\allowbreak 50,\allowbreak 14,\allowbreak 51,\allowbreak 13,\allowbreak 52,\allowbreak 12)}\tabularnewline[5pt]
{\raggedright\arraybackslash {\ttfamily\footnotesize (4,\allowbreak 28,\allowbreak 32)}\par $K=32$\par pinned W4} & {\ttfamily\footnotesize (60,\allowbreak 5,\allowbreak 59,\allowbreak 6,\allowbreak 62,\allowbreak 3,\allowbreak 61,\allowbreak 4,\allowbreak 56,\allowbreak 10,\allowbreak 57,\allowbreak 8,\allowbreak 40,\allowbreak 23,\allowbreak 42,\allowbreak 26,\allowbreak 41,\allowbreak 28,\allowbreak 37,\allowbreak 30,\allowbreak 38,\allowbreak 27,\allowbreak 39,\allowbreak 25,\allowbreak 35,\allowbreak 29,\allowbreak 33,\allowbreak 32,\allowbreak 34,\allowbreak 31,\allowbreak 36,\allowbreak 18)\par\smallskip (50,\allowbreak 15,\allowbreak 49,\allowbreak 16,\allowbreak 45,\allowbreak 20,\allowbreak 46,\allowbreak 19,\allowbreak 47,\allowbreak 17,\allowbreak 48,\allowbreak 24,\allowbreak 44,\allowbreak 21,\allowbreak 43,\allowbreak 22,\allowbreak 58,\allowbreak 7,\allowbreak 55,\allowbreak 11,\allowbreak 54,\allowbreak 9,\allowbreak 51,\allowbreak 14,\allowbreak 52,\allowbreak 13,\allowbreak 53,\allowbreak 12)\par\smallskip (64,\allowbreak 1,\allowbreak 63,\allowbreak 2)}\tabularnewline[5pt]
{\raggedright\arraybackslash {\ttfamily\footnotesize (4,\allowbreak 29,\allowbreak 29)}\par $K=30$\par pinned W4\par $2e\le b+7$} & {\ttfamily\footnotesize (62,\allowbreak 1,\allowbreak 61,\allowbreak 2)\par\smallskip (57,\allowbreak 8,\allowbreak 54,\allowbreak 4,\allowbreak 56,\allowbreak 5,\allowbreak 52,\allowbreak 12,\allowbreak 51,\allowbreak 9,\allowbreak 50,\allowbreak 16,\allowbreak 41,\allowbreak 25,\allowbreak 33,\allowbreak 29,\allowbreak 35,\allowbreak 24,\allowbreak 36,\allowbreak 26,\allowbreak 39,\allowbreak 22,\allowbreak 42,\allowbreak 19,\allowbreak 40,\allowbreak 21,\allowbreak 43,\allowbreak 17,\allowbreak 47)\par\smallskip (58,\allowbreak 3,\allowbreak 59,\allowbreak 15,\allowbreak 44,\allowbreak 20,\allowbreak 38,\allowbreak 23,\allowbreak 37,\allowbreak 28,\allowbreak 31,\allowbreak 30,\allowbreak 32,\allowbreak 27,\allowbreak 34,\allowbreak 6,\allowbreak 60,\allowbreak 7,\allowbreak 55,\allowbreak 10,\allowbreak 53,\allowbreak 18,\allowbreak 45,\allowbreak 14,\allowbreak 46,\allowbreak 13,\allowbreak 49,\allowbreak 11,\allowbreak 48)}\tabularnewline[5pt]
{\raggedright\arraybackslash {\ttfamily\footnotesize (4,\allowbreak 29,\allowbreak 30)}\par $K=31$\par pinned W4} & {\ttfamily\footnotesize (63,\allowbreak 1,\allowbreak 62,\allowbreak 2)\par\smallskip (58,\allowbreak 6,\allowbreak 61,\allowbreak 4,\allowbreak 60,\allowbreak 7,\allowbreak 56,\allowbreak 12,\allowbreak 50,\allowbreak 10,\allowbreak 57,\allowbreak 3,\allowbreak 40,\allowbreak 26,\allowbreak 44,\allowbreak 18,\allowbreak 42,\allowbreak 21,\allowbreak 49,\allowbreak 19,\allowbreak 39,\allowbreak 23,\allowbreak 45,\allowbreak 14,\allowbreak 55,\allowbreak 5,\allowbreak 53,\allowbreak 11,\allowbreak 46,\allowbreak 17)\par\smallskip (59,\allowbreak 8,\allowbreak 54,\allowbreak 9,\allowbreak 52,\allowbreak 13,\allowbreak 47,\allowbreak 22,\allowbreak 41,\allowbreak 24,\allowbreak 37,\allowbreak 25,\allowbreak 36,\allowbreak 30,\allowbreak 33,\allowbreak 31,\allowbreak 32,\allowbreak 28,\allowbreak 38,\allowbreak 29,\allowbreak 34,\allowbreak 27,\allowbreak 35,\allowbreak 20,\allowbreak 43,\allowbreak 16,\allowbreak 48,\allowbreak 15,\allowbreak 51)}\tabularnewline[5pt]
{\raggedright\arraybackslash {\ttfamily\footnotesize (4,\allowbreak 29,\allowbreak 31)}\par $K=31$\par pinned W4} & {\ttfamily\footnotesize (64,\allowbreak 1,\allowbreak 63,\allowbreak 2)\par\smallskip (55,\allowbreak 14,\allowbreak 48,\allowbreak 19,\allowbreak 43,\allowbreak 22,\allowbreak 38,\allowbreak 25,\allowbreak 39,\allowbreak 24,\allowbreak 41,\allowbreak 23,\allowbreak 42,\allowbreak 16,\allowbreak 49,\allowbreak 10,\allowbreak 57,\allowbreak 6,\allowbreak 58,\allowbreak 13,\allowbreak 56,\allowbreak 7,\allowbreak 53,\allowbreak 9,\allowbreak 59,\allowbreak 5,\allowbreak 61,\allowbreak 8,\allowbreak 46,\allowbreak 21,\allowbreak 44)\par\smallskip (60,\allowbreak 3,\allowbreak 62,\allowbreak 4,\allowbreak 52,\allowbreak 15,\allowbreak 50,\allowbreak 18,\allowbreak 45,\allowbreak 17,\allowbreak 47,\allowbreak 11,\allowbreak 51,\allowbreak 20,\allowbreak 40,\allowbreak 28,\allowbreak 32,\allowbreak 31,\allowbreak 33,\allowbreak 30,\allowbreak 35,\allowbreak 29,\allowbreak 36,\allowbreak 27,\allowbreak 37,\allowbreak 26,\allowbreak 34,\allowbreak 12,\allowbreak 54)}\tabularnewline[5pt]
{\raggedright\arraybackslash {\ttfamily\footnotesize (4,\allowbreak 29,\allowbreak 32)}\par $K=32$\par pinned W4} & {\ttfamily\footnotesize (65,\allowbreak 1,\allowbreak 64,\allowbreak 2)\par\smallskip (58,\allowbreak 8,\allowbreak 56,\allowbreak 16,\allowbreak 47,\allowbreak 11,\allowbreak 49,\allowbreak 15,\allowbreak 42,\allowbreak 23,\allowbreak 46,\allowbreak 13,\allowbreak 59,\allowbreak 5,\allowbreak 57,\allowbreak 12,\allowbreak 53,\allowbreak 14,\allowbreak 51,\allowbreak 19,\allowbreak 54,\allowbreak 10,\allowbreak 52,\allowbreak 9,\allowbreak 60,\allowbreak 7,\allowbreak 62,\allowbreak 6,\allowbreak 55)\par\smallskip (61,\allowbreak 4,\allowbreak 63,\allowbreak 3,\allowbreak 50,\allowbreak 21,\allowbreak 39,\allowbreak 22,\allowbreak 44,\allowbreak 24,\allowbreak 45,\allowbreak 32,\allowbreak 33,\allowbreak 31,\allowbreak 34,\allowbreak 30,\allowbreak 35,\allowbreak 29,\allowbreak 36,\allowbreak 28,\allowbreak 37,\allowbreak 27,\allowbreak 38,\allowbreak 26,\allowbreak 40,\allowbreak 25,\allowbreak 41,\allowbreak 17,\allowbreak 43,\allowbreak 18,\allowbreak 48,\allowbreak 20)}\tabularnewline[5pt]
{\raggedright\arraybackslash {\ttfamily\footnotesize (4,\allowbreak 30,\allowbreak 30)}\par $K=32$\par pinned W4} & {\ttfamily\footnotesize (64,\allowbreak 1,\allowbreak 63,\allowbreak 2)\par\smallskip (58,\allowbreak 8,\allowbreak 53,\allowbreak 11,\allowbreak 45,\allowbreak 24,\allowbreak 40,\allowbreak 25,\allowbreak 39,\allowbreak 26,\allowbreak 38,\allowbreak 27,\allowbreak 37,\allowbreak 28,\allowbreak 36,\allowbreak 29,\allowbreak 35,\allowbreak 30,\allowbreak 34,\allowbreak 31,\allowbreak 33,\allowbreak 32,\allowbreak 49,\allowbreak 12,\allowbreak 44,\allowbreak 21,\allowbreak 47,\allowbreak 20,\allowbreak 42,\allowbreak 14)\par\smallskip (60,\allowbreak 6,\allowbreak 62,\allowbreak 3,\allowbreak 56,\allowbreak 5,\allowbreak 57,\allowbreak 19,\allowbreak 50,\allowbreak 17,\allowbreak 46,\allowbreak 22,\allowbreak 41,\allowbreak 23,\allowbreak 43,\allowbreak 18,\allowbreak 48,\allowbreak 13,\allowbreak 61,\allowbreak 4,\allowbreak 59,\allowbreak 10,\allowbreak 51,\allowbreak 7,\allowbreak 54,\allowbreak 15,\allowbreak 55,\allowbreak 9,\allowbreak 52,\allowbreak 16)}\tabularnewline[5pt]
{\raggedright\arraybackslash {\ttfamily\footnotesize (4,\allowbreak 30,\allowbreak 31)}\par $K=32$\par pinned W4\par $2e\le b+8$} & {\ttfamily\footnotesize (65,\allowbreak 1,\allowbreak 64,\allowbreak 2)\par\smallskip (57,\allowbreak 12,\allowbreak 52,\allowbreak 16,\allowbreak 44,\allowbreak 23,\allowbreak 43,\allowbreak 24,\allowbreak 42,\allowbreak 25,\allowbreak 41,\allowbreak 19,\allowbreak 46,\allowbreak 22,\allowbreak 45,\allowbreak 20,\allowbreak 35,\allowbreak 31,\allowbreak 37,\allowbreak 28,\allowbreak 39,\allowbreak 27,\allowbreak 40,\allowbreak 26,\allowbreak 36,\allowbreak 29,\allowbreak 34,\allowbreak 32,\allowbreak 33,\allowbreak 30,\allowbreak 38)\par\smallskip (61,\allowbreak 5,\allowbreak 60,\allowbreak 6,\allowbreak 59,\allowbreak 8,\allowbreak 58,\allowbreak 9,\allowbreak 56,\allowbreak 14,\allowbreak 53,\allowbreak 7,\allowbreak 55,\allowbreak 11,\allowbreak 54,\allowbreak 21,\allowbreak 47,\allowbreak 18,\allowbreak 48,\allowbreak 17,\allowbreak 49,\allowbreak 15,\allowbreak 50,\allowbreak 13,\allowbreak 51,\allowbreak 10,\allowbreak 62,\allowbreak 4,\allowbreak 63,\allowbreak 3)}\tabularnewline[5pt]
{\raggedright\arraybackslash {\ttfamily\footnotesize (4,\allowbreak 30,\allowbreak 32)}\par $K=33$\par pinned W4} & {\ttfamily\footnotesize (66,\allowbreak 1,\allowbreak 65,\allowbreak 2)\par\smallskip (55,\allowbreak 17,\allowbreak 51,\allowbreak 5,\allowbreak 61,\allowbreak 4,\allowbreak 63,\allowbreak 3,\allowbreak 64,\allowbreak 6,\allowbreak 60,\allowbreak 10,\allowbreak 58,\allowbreak 11,\allowbreak 52,\allowbreak 7,\allowbreak 59,\allowbreak 8,\allowbreak 57,\allowbreak 15,\allowbreak 54,\allowbreak 14,\allowbreak 50,\allowbreak 19,\allowbreak 45,\allowbreak 22,\allowbreak 46,\allowbreak 25,\allowbreak 40,\allowbreak 12,\allowbreak 56,\allowbreak 13)\par\smallskip (62,\allowbreak 9,\allowbreak 44,\allowbreak 24,\allowbreak 43,\allowbreak 21,\allowbreak 53,\allowbreak 16,\allowbreak 49,\allowbreak 20,\allowbreak 47,\allowbreak 30,\allowbreak 38,\allowbreak 26,\allowbreak 42,\allowbreak 29,\allowbreak 39,\allowbreak 28,\allowbreak 35,\allowbreak 33,\allowbreak 34,\allowbreak 31,\allowbreak 37,\allowbreak 32,\allowbreak 36,\allowbreak 27,\allowbreak 41,\allowbreak 23,\allowbreak 48,\allowbreak 18)}\tabularnewline[5pt]
{\raggedright\arraybackslash {\ttfamily\footnotesize (4,\allowbreak 30,\allowbreak 33)}\par $K=33$\par pinned W4} & {\ttfamily\footnotesize (67,\allowbreak 1,\allowbreak 66,\allowbreak 2)\par\smallskip (62,\allowbreak 8,\allowbreak 56,\allowbreak 13,\allowbreak 60,\allowbreak 3,\allowbreak 54,\allowbreak 14,\allowbreak 59,\allowbreak 24,\allowbreak 42,\allowbreak 28,\allowbreak 41,\allowbreak 32,\allowbreak 34,\allowbreak 33,\allowbreak 36,\allowbreak 30,\allowbreak 35,\allowbreak 31,\allowbreak 38,\allowbreak 22,\allowbreak 43,\allowbreak 18,\allowbreak 47,\allowbreak 21,\allowbreak 45,\allowbreak 23,\allowbreak 51,\allowbreak 12)\par\smallskip (63,\allowbreak 10,\allowbreak 52,\allowbreak 11,\allowbreak 55,\allowbreak 9,\allowbreak 58,\allowbreak 6,\allowbreak 61,\allowbreak 5,\allowbreak 64,\allowbreak 4,\allowbreak 65,\allowbreak 7,\allowbreak 57,\allowbreak 20,\allowbreak 40,\allowbreak 25,\allowbreak 44,\allowbreak 27,\allowbreak 39,\allowbreak 29,\allowbreak 37,\allowbreak 26,\allowbreak 49,\allowbreak 19,\allowbreak 46,\allowbreak 15,\allowbreak 53,\allowbreak 17,\allowbreak 50,\allowbreak 16,\allowbreak 48)}\tabularnewline[5pt]
{\raggedright\arraybackslash {\ttfamily\footnotesize (4,\allowbreak 30,\allowbreak 34)}\par $K=34$\par pinned W4} & {\ttfamily\footnotesize (68,\allowbreak 1,\allowbreak 67,\allowbreak 2)\par\smallskip (63,\allowbreak 3,\allowbreak 65,\allowbreak 4,\allowbreak 56,\allowbreak 15,\allowbreak 41,\allowbreak 29,\allowbreak 42,\allowbreak 31,\allowbreak 37,\allowbreak 34,\allowbreak 35,\allowbreak 33,\allowbreak 38,\allowbreak 28,\allowbreak 36,\allowbreak 32,\allowbreak 39,\allowbreak 30,\allowbreak 48,\allowbreak 20,\allowbreak 49,\allowbreak 18,\allowbreak 55,\allowbreak 9,\allowbreak 57,\allowbreak 6,\allowbreak 59,\allowbreak 19)\par\smallskip (64,\allowbreak 10,\allowbreak 60,\allowbreak 17,\allowbreak 51,\allowbreak 16,\allowbreak 54,\allowbreak 12,\allowbreak 61,\allowbreak 5,\allowbreak 62,\allowbreak 7,\allowbreak 66,\allowbreak 8,\allowbreak 53,\allowbreak 21,\allowbreak 45,\allowbreak 24,\allowbreak 44,\allowbreak 11,\allowbreak 58,\allowbreak 14,\allowbreak 50,\allowbreak 23,\allowbreak 46,\allowbreak 27,\allowbreak 43,\allowbreak 26,\allowbreak 40,\allowbreak 25,\allowbreak 47,\allowbreak 22,\allowbreak 52,\allowbreak 13)}\tabularnewline[5pt]
{\raggedright\arraybackslash {\ttfamily\footnotesize (4,\allowbreak 31,\allowbreak 31)}\par $K=32$\par pinned W4\par $2e\le b+7$} & {\ttfamily\footnotesize (66,\allowbreak 1,\allowbreak 65,\allowbreak 2)\par\smallskip (60,\allowbreak 5,\allowbreak 58,\allowbreak 9,\allowbreak 59,\allowbreak 7,\allowbreak 63,\allowbreak 4,\allowbreak 61,\allowbreak 3,\allowbreak 64,\allowbreak 19,\allowbreak 39,\allowbreak 30,\allowbreak 38,\allowbreak 25,\allowbreak 40,\allowbreak 24,\allowbreak 36,\allowbreak 32,\allowbreak 33,\allowbreak 31,\allowbreak 34,\allowbreak 29,\allowbreak 35,\allowbreak 28,\allowbreak 45,\allowbreak 17,\allowbreak 47,\allowbreak 22,\allowbreak 48)\par\smallskip (62,\allowbreak 8,\allowbreak 56,\allowbreak 10,\allowbreak 57,\allowbreak 6,\allowbreak 44,\allowbreak 26,\allowbreak 37,\allowbreak 27,\allowbreak 41,\allowbreak 20,\allowbreak 52,\allowbreak 18,\allowbreak 42,\allowbreak 23,\allowbreak 46,\allowbreak 15,\allowbreak 51,\allowbreak 14,\allowbreak 49,\allowbreak 16,\allowbreak 43,\allowbreak 21,\allowbreak 50,\allowbreak 11,\allowbreak 55,\allowbreak 12,\allowbreak 54,\allowbreak 13,\allowbreak 53)}\tabularnewline[5pt]
{\raggedright\arraybackslash {\ttfamily\footnotesize (4,\allowbreak 31,\allowbreak 32)}\par $K=33$\par pinned W4} & {\ttfamily\footnotesize (67,\allowbreak 1,\allowbreak 66,\allowbreak 2)\par\smallskip (61,\allowbreak 6,\allowbreak 56,\allowbreak 11,\allowbreak 53,\allowbreak 20,\allowbreak 45,\allowbreak 18,\allowbreak 46,\allowbreak 23,\allowbreak 44,\allowbreak 27,\allowbreak 38,\allowbreak 30,\allowbreak 39,\allowbreak 29,\allowbreak 42,\allowbreak 22,\allowbreak 51,\allowbreak 13,\allowbreak 54,\allowbreak 19,\allowbreak 43,\allowbreak 25,\allowbreak 41,\allowbreak 26,\allowbreak 40,\allowbreak 21,\allowbreak 47,\allowbreak 17,\allowbreak 48,\allowbreak 16)\par\smallskip (63,\allowbreak 10,\allowbreak 49,\allowbreak 5,\allowbreak 64,\allowbreak 4,\allowbreak 62,\allowbreak 8,\allowbreak 57,\allowbreak 9,\allowbreak 65,\allowbreak 3,\allowbreak 60,\allowbreak 14,\allowbreak 50,\allowbreak 28,\allowbreak 35,\allowbreak 33,\allowbreak 34,\allowbreak 31,\allowbreak 37,\allowbreak 32,\allowbreak 36,\allowbreak 24,\allowbreak 58,\allowbreak 7,\allowbreak 59,\allowbreak 12,\allowbreak 55,\allowbreak 15,\allowbreak 52)}\tabularnewline[5pt]
{\raggedright\arraybackslash {\ttfamily\footnotesize (4,\allowbreak 31,\allowbreak 33)}\par $K=33$\par pinned W4} & {\ttfamily\footnotesize (68,\allowbreak 1,\allowbreak 67,\allowbreak 2)\par\smallskip (63,\allowbreak 13,\allowbreak 59,\allowbreak 10,\allowbreak 57,\allowbreak 5,\allowbreak 61,\allowbreak 4,\allowbreak 65,\allowbreak 3,\allowbreak 62,\allowbreak 11,\allowbreak 56,\allowbreak 15,\allowbreak 54,\allowbreak 14,\allowbreak 51,\allowbreak 21,\allowbreak 48,\allowbreak 24,\allowbreak 42,\allowbreak 26,\allowbreak 41,\allowbreak 12,\allowbreak 55,\allowbreak 7,\allowbreak 60,\allowbreak 6,\allowbreak 66,\allowbreak 8,\allowbreak 46)\par\smallskip (64,\allowbreak 9,\allowbreak 53,\allowbreak 17,\allowbreak 52,\allowbreak 18,\allowbreak 50,\allowbreak 19,\allowbreak 47,\allowbreak 25,\allowbreak 44,\allowbreak 27,\allowbreak 40,\allowbreak 28,\allowbreak 39,\allowbreak 29,\allowbreak 38,\allowbreak 30,\allowbreak 37,\allowbreak 31,\allowbreak 36,\allowbreak 32,\allowbreak 35,\allowbreak 33,\allowbreak 34,\allowbreak 20,\allowbreak 45,\allowbreak 22,\allowbreak 43,\allowbreak 23,\allowbreak 49,\allowbreak 16,\allowbreak 58)}\tabularnewline[5pt]
{\raggedright\arraybackslash {\ttfamily\footnotesize (4,\allowbreak 31,\allowbreak 34)}\par $K=34$\par pinned W4} & {\ttfamily\footnotesize (69,\allowbreak 1,\allowbreak 68,\allowbreak 2)\par\smallskip (63,\allowbreak 10,\allowbreak 64,\allowbreak 3,\allowbreak 67,\allowbreak 11,\allowbreak 54,\allowbreak 16,\allowbreak 52,\allowbreak 24,\allowbreak 51,\allowbreak 20,\allowbreak 41,\allowbreak 28,\allowbreak 38,\allowbreak 26,\allowbreak 40,\allowbreak 33,\allowbreak 35,\allowbreak 34,\allowbreak 37,\allowbreak 32,\allowbreak 36,\allowbreak 30,\allowbreak 39,\allowbreak 31,\allowbreak 42,\allowbreak 22,\allowbreak 47,\allowbreak 25,\allowbreak 43,\allowbreak 9,\allowbreak 59,\allowbreak 13)\par\smallskip (65,\allowbreak 5,\allowbreak 62,\allowbreak 4,\allowbreak 66,\allowbreak 7,\allowbreak 58,\allowbreak 14,\allowbreak 61,\allowbreak 6,\allowbreak 55,\allowbreak 18,\allowbreak 53,\allowbreak 21,\allowbreak 45,\allowbreak 29,\allowbreak 44,\allowbreak 27,\allowbreak 46,\allowbreak 23,\allowbreak 49,\allowbreak 19,\allowbreak 48,\allowbreak 8,\allowbreak 60,\allowbreak 12,\allowbreak 57,\allowbreak 15,\allowbreak 56,\allowbreak 17,\allowbreak 50)}\tabularnewline[5pt]
{\raggedright\arraybackslash {\ttfamily\footnotesize (4,\allowbreak 32,\allowbreak 32)}\par $K=34$\par pinned W4} & {\ttfamily\footnotesize (68,\allowbreak 1,\allowbreak 67,\allowbreak 2)\par\smallskip (62,\allowbreak 13,\allowbreak 54,\allowbreak 17,\allowbreak 47,\allowbreak 24,\allowbreak 44,\allowbreak 29,\allowbreak 35,\allowbreak 34,\allowbreak 36,\allowbreak 33,\allowbreak 37,\allowbreak 32,\allowbreak 40,\allowbreak 27,\allowbreak 41,\allowbreak 25,\allowbreak 43,\allowbreak 26,\allowbreak 45,\allowbreak 23,\allowbreak 58,\allowbreak 12,\allowbreak 55,\allowbreak 16,\allowbreak 48,\allowbreak 22,\allowbreak 46,\allowbreak 19,\allowbreak 52,\allowbreak 18)\par\smallskip (64,\allowbreak 6,\allowbreak 51,\allowbreak 20,\allowbreak 49,\allowbreak 21,\allowbreak 42,\allowbreak 30,\allowbreak 39,\allowbreak 28,\allowbreak 38,\allowbreak 31,\allowbreak 56,\allowbreak 9,\allowbreak 57,\allowbreak 5,\allowbreak 60,\allowbreak 4,\allowbreak 65,\allowbreak 8,\allowbreak 59,\allowbreak 15,\allowbreak 53,\allowbreak 11,\allowbreak 61,\allowbreak 7,\allowbreak 66,\allowbreak 3,\allowbreak 63,\allowbreak 10,\allowbreak 50,\allowbreak 14)}\tabularnewline[5pt]
{\raggedright\arraybackslash {\ttfamily\footnotesize (4,\allowbreak 32,\allowbreak 33)}\par $K=34$\par pinned W4\par $2e\le b+8$} & {\ttfamily\footnotesize (69,\allowbreak 1,\allowbreak 68,\allowbreak 2)\par\smallskip (63,\allowbreak 5,\allowbreak 64,\allowbreak 3,\allowbreak 67,\allowbreak 7,\allowbreak 56,\allowbreak 15,\allowbreak 53,\allowbreak 18,\allowbreak 51,\allowbreak 19,\allowbreak 49,\allowbreak 22,\allowbreak 44,\allowbreak 24,\allowbreak 43,\allowbreak 26,\allowbreak 39,\allowbreak 31,\allowbreak 35,\allowbreak 34,\allowbreak 36,\allowbreak 33,\allowbreak 38,\allowbreak 32,\allowbreak 66,\allowbreak 4,\allowbreak 61,\allowbreak 6,\allowbreak 60,\allowbreak 13)\par\smallskip (65,\allowbreak 9,\allowbreak 62,\allowbreak 10,\allowbreak 58,\allowbreak 8,\allowbreak 59,\allowbreak 14,\allowbreak 57,\allowbreak 11,\allowbreak 55,\allowbreak 16,\allowbreak 52,\allowbreak 12,\allowbreak 54,\allowbreak 17,\allowbreak 48,\allowbreak 20,\allowbreak 45,\allowbreak 27,\allowbreak 42,\allowbreak 28,\allowbreak 37,\allowbreak 30,\allowbreak 40,\allowbreak 29,\allowbreak 41,\allowbreak 25,\allowbreak 46,\allowbreak 23,\allowbreak 47,\allowbreak 21,\allowbreak 50)}\tabularnewline[5pt]
{\raggedright\arraybackslash {\ttfamily\footnotesize (4,\allowbreak 32,\allowbreak 34)}\par $K=35$\par pinned W4} & {\ttfamily\footnotesize (70,\allowbreak 1,\allowbreak 69,\allowbreak 2)\par\smallskip (65,\allowbreak 3,\allowbreak 67,\allowbreak 4,\allowbreak 62,\allowbreak 12,\allowbreak 58,\allowbreak 10,\allowbreak 64,\allowbreak 5,\allowbreak 61,\allowbreak 6,\allowbreak 59,\allowbreak 15,\allowbreak 48,\allowbreak 22,\allowbreak 46,\allowbreak 26,\allowbreak 43,\allowbreak 29,\allowbreak 39,\allowbreak 31,\allowbreak 40,\allowbreak 33,\allowbreak 37,\allowbreak 32,\allowbreak 38,\allowbreak 35,\allowbreak 36,\allowbreak 34,\allowbreak 57,\allowbreak 14,\allowbreak 63,\allowbreak 16)\par\smallskip (66,\allowbreak 9,\allowbreak 60,\allowbreak 8,\allowbreak 68,\allowbreak 7,\allowbreak 52,\allowbreak 20,\allowbreak 51,\allowbreak 11,\allowbreak 53,\allowbreak 23,\allowbreak 50,\allowbreak 21,\allowbreak 49,\allowbreak 24,\allowbreak 45,\allowbreak 27,\allowbreak 42,\allowbreak 30,\allowbreak 41,\allowbreak 28,\allowbreak 44,\allowbreak 25,\allowbreak 47,\allowbreak 13,\allowbreak 54,\allowbreak 19,\allowbreak 55,\allowbreak 18,\allowbreak 56,\allowbreak 17)}\tabularnewline[5pt]
{\raggedright\arraybackslash {\ttfamily\footnotesize (4,\allowbreak 32,\allowbreak 35)}\par $K=35$\par pinned W4} & {\ttfamily\footnotesize (71,\allowbreak 1,\allowbreak 70,\allowbreak 2)\par\smallskip (65,\allowbreak 7,\allowbreak 69,\allowbreak 3,\allowbreak 64,\allowbreak 8,\allowbreak 59,\allowbreak 15,\allowbreak 56,\allowbreak 11,\allowbreak 60,\allowbreak 10,\allowbreak 62,\allowbreak 5,\allowbreak 68,\allowbreak 4,\allowbreak 63,\allowbreak 16,\allowbreak 48,\allowbreak 29,\allowbreak 47,\allowbreak 20,\allowbreak 49,\allowbreak 23,\allowbreak 46,\allowbreak 25,\allowbreak 55,\allowbreak 9,\allowbreak 57,\allowbreak 17,\allowbreak 50,\allowbreak 13)\par\smallskip (67,\allowbreak 14,\allowbreak 52,\allowbreak 21,\allowbreak 45,\allowbreak 28,\allowbreak 42,\allowbreak 32,\allowbreak 36,\allowbreak 35,\allowbreak 37,\allowbreak 34,\allowbreak 40,\allowbreak 27,\allowbreak 43,\allowbreak 31,\allowbreak 38,\allowbreak 33,\allowbreak 41,\allowbreak 30,\allowbreak 39,\allowbreak 24,\allowbreak 44,\allowbreak 22,\allowbreak 58,\allowbreak 19,\allowbreak 53,\allowbreak 18,\allowbreak 61,\allowbreak 6,\allowbreak 66,\allowbreak 12,\allowbreak 54,\allowbreak 26,\allowbreak 51)}\tabularnewline[5pt]
{\raggedright\arraybackslash {\ttfamily\footnotesize (4,\allowbreak 32,\allowbreak 36)}\par $K=36$\par pinned W4} & {\ttfamily\footnotesize (72,\allowbreak 1,\allowbreak 71,\allowbreak 2)\par\smallskip (66,\allowbreak 8,\allowbreak 65,\allowbreak 14,\allowbreak 67,\allowbreak 5,\allowbreak 69,\allowbreak 4,\allowbreak 64,\allowbreak 3,\allowbreak 70,\allowbreak 7,\allowbreak 63,\allowbreak 11,\allowbreak 53,\allowbreak 18,\allowbreak 59,\allowbreak 20,\allowbreak 51,\allowbreak 17,\allowbreak 60,\allowbreak 6,\allowbreak 61,\allowbreak 13,\allowbreak 57,\allowbreak 21,\allowbreak 46,\allowbreak 29,\allowbreak 43,\allowbreak 25,\allowbreak 48,\allowbreak 19,\allowbreak 52,\allowbreak 15,\allowbreak 62,\allowbreak 12)\par\smallskip (68,\allowbreak 9,\allowbreak 58,\allowbreak 26,\allowbreak 50,\allowbreak 23,\allowbreak 49,\allowbreak 30,\allowbreak 40,\allowbreak 31,\allowbreak 42,\allowbreak 35,\allowbreak 37,\allowbreak 36,\allowbreak 39,\allowbreak 34,\allowbreak 38,\allowbreak 32,\allowbreak 47,\allowbreak 27,\allowbreak 55,\allowbreak 10,\allowbreak 56,\allowbreak 16,\allowbreak 54,\allowbreak 24,\allowbreak 45,\allowbreak 33,\allowbreak 41,\allowbreak 28,\allowbreak 44,\allowbreak 22)}\tabularnewline[5pt]
{\raggedright\arraybackslash {\ttfamily\footnotesize (4,\allowbreak 33,\allowbreak 33)}\par $K=34$\par pinned W4\par $2e\le b+7$} & {\ttfamily\footnotesize (70,\allowbreak 1,\allowbreak 69,\allowbreak 2)\par\smallskip (64,\allowbreak 3,\allowbreak 68,\allowbreak 8,\allowbreak 57,\allowbreak 16,\allowbreak 50,\allowbreak 19,\allowbreak 56,\allowbreak 10,\allowbreak 58,\allowbreak 15,\allowbreak 41,\allowbreak 29,\allowbreak 40,\allowbreak 31,\allowbreak 39,\allowbreak 25,\allowbreak 49,\allowbreak 20,\allowbreak 52,\allowbreak 13,\allowbreak 60,\allowbreak 9,\allowbreak 63,\allowbreak 6,\allowbreak 65,\allowbreak 12,\allowbreak 48,\allowbreak 21,\allowbreak 51,\allowbreak 18,\allowbreak 53)\par\smallskip (66,\allowbreak 24,\allowbreak 43,\allowbreak 30,\allowbreak 37,\allowbreak 34,\allowbreak 35,\allowbreak 33,\allowbreak 38,\allowbreak 32,\allowbreak 36,\allowbreak 26,\allowbreak 42,\allowbreak 27,\allowbreak 47,\allowbreak 22,\allowbreak 44,\allowbreak 23,\allowbreak 46,\allowbreak 28,\allowbreak 45,\allowbreak 17,\allowbreak 55,\allowbreak 11,\allowbreak 61,\allowbreak 5,\allowbreak 67,\allowbreak 4,\allowbreak 62,\allowbreak 7,\allowbreak 59,\allowbreak 14,\allowbreak 54)}\tabularnewline[5pt]
{\raggedright\arraybackslash {\ttfamily\footnotesize (4,\allowbreak 33,\allowbreak 34)}\par $K=35$\par pinned W4} & {\ttfamily\footnotesize (71,\allowbreak 1,\allowbreak 70,\allowbreak 2)\par\smallskip (65,\allowbreak 7,\allowbreak 55,\allowbreak 14,\allowbreak 38,\allowbreak 33,\allowbreak 40,\allowbreak 32,\allowbreak 41,\allowbreak 30,\allowbreak 42,\allowbreak 26,\allowbreak 46,\allowbreak 25,\allowbreak 44,\allowbreak 31,\allowbreak 37,\allowbreak 34,\allowbreak 36,\allowbreak 35,\allowbreak 39,\allowbreak 29,\allowbreak 43,\allowbreak 28,\allowbreak 45,\allowbreak 27,\allowbreak 49,\allowbreak 10,\allowbreak 61,\allowbreak 15,\allowbreak 58,\allowbreak 13,\allowbreak 57)\par\smallskip (67,\allowbreak 17,\allowbreak 52,\allowbreak 20,\allowbreak 54,\allowbreak 18,\allowbreak 56,\allowbreak 23,\allowbreak 48,\allowbreak 22,\allowbreak 50,\allowbreak 21,\allowbreak 51,\allowbreak 24,\allowbreak 47,\allowbreak 16,\allowbreak 53,\allowbreak 11,\allowbreak 60,\allowbreak 5,\allowbreak 62,\allowbreak 9,\allowbreak 63,\allowbreak 4,\allowbreak 68,\allowbreak 6,\allowbreak 66,\allowbreak 3,\allowbreak 69,\allowbreak 8,\allowbreak 64,\allowbreak 12,\allowbreak 59,\allowbreak 19)}\tabularnewline[5pt]
{\raggedright\arraybackslash {\ttfamily\footnotesize (4,\allowbreak 33,\allowbreak 35)}\par $K=35$\par pinned W4} & {\ttfamily\footnotesize (72,\allowbreak 1,\allowbreak 71,\allowbreak 2)\par\smallskip (66,\allowbreak 8,\allowbreak 65,\allowbreak 6,\allowbreak 70,\allowbreak 3,\allowbreak 64,\allowbreak 4,\allowbreak 69,\allowbreak 25,\allowbreak 45,\allowbreak 26,\allowbreak 43,\allowbreak 30,\allowbreak 40,\allowbreak 31,\allowbreak 42,\allowbreak 28,\allowbreak 44,\allowbreak 29,\allowbreak 41,\allowbreak 33,\allowbreak 38,\allowbreak 32,\allowbreak 39,\allowbreak 35,\allowbreak 36,\allowbreak 34,\allowbreak 37,\allowbreak 19,\allowbreak 52,\allowbreak 22,\allowbreak 47,\allowbreak 20,\allowbreak 57)\par\smallskip (68,\allowbreak 5,\allowbreak 67,\allowbreak 14,\allowbreak 54,\allowbreak 18,\allowbreak 50,\allowbreak 16,\allowbreak 55,\allowbreak 13,\allowbreak 51,\allowbreak 27,\allowbreak 48,\allowbreak 17,\allowbreak 62,\allowbreak 7,\allowbreak 63,\allowbreak 9,\allowbreak 61,\allowbreak 10,\allowbreak 60,\allowbreak 11,\allowbreak 59,\allowbreak 12,\allowbreak 58,\allowbreak 15,\allowbreak 56,\allowbreak 21,\allowbreak 49,\allowbreak 23,\allowbreak 46,\allowbreak 24,\allowbreak 53)}\tabularnewline[5pt]
{\raggedright\arraybackslash {\ttfamily\footnotesize (4,\allowbreak 33,\allowbreak 36)}\par $K=36$\par pinned W4} & {\ttfamily\footnotesize (73,\allowbreak 1,\allowbreak 72,\allowbreak 2)\par\smallskip (65,\allowbreak 4,\allowbreak 70,\allowbreak 3,\allowbreak 71,\allowbreak 12,\allowbreak 60,\allowbreak 15,\allowbreak 61,\allowbreak 11,\allowbreak 66,\allowbreak 8,\allowbreak 48,\allowbreak 16,\allowbreak 57,\allowbreak 21,\allowbreak 46,\allowbreak 30,\allowbreak 43,\allowbreak 29,\allowbreak 41,\allowbreak 33,\allowbreak 40,\allowbreak 31,\allowbreak 42,\allowbreak 27,\allowbreak 44,\allowbreak 26,\allowbreak 50,\allowbreak 23,\allowbreak 49,\allowbreak 28,\allowbreak 47,\allowbreak 24,\allowbreak 58,\allowbreak 14)\par\smallskip (69,\allowbreak 5,\allowbreak 68,\allowbreak 6,\allowbreak 62,\allowbreak 10,\allowbreak 63,\allowbreak 9,\allowbreak 56,\allowbreak 13,\allowbreak 64,\allowbreak 7,\allowbreak 67,\allowbreak 18,\allowbreak 55,\allowbreak 17,\allowbreak 59,\allowbreak 20,\allowbreak 51,\allowbreak 22,\allowbreak 52,\allowbreak 19,\allowbreak 54,\allowbreak 32,\allowbreak 38,\allowbreak 34,\allowbreak 39,\allowbreak 36,\allowbreak 37,\allowbreak 35,\allowbreak 45,\allowbreak 25,\allowbreak 53)}\tabularnewline[5pt]
{\raggedright\arraybackslash {\ttfamily\footnotesize (4,\allowbreak 34,\allowbreak 34)}\par $K=36$\par pinned W4} & {\ttfamily\footnotesize (72,\allowbreak 1,\allowbreak 71,\allowbreak 2)\par\smallskip (44,\allowbreak 25,\allowbreak 42,\allowbreak 31,\allowbreak 37,\allowbreak 36,\allowbreak 38,\allowbreak 35,\allowbreak 39,\allowbreak 34,\allowbreak 46,\allowbreak 20,\allowbreak 49,\allowbreak 22,\allowbreak 53,\allowbreak 23,\allowbreak 57,\allowbreak 14,\allowbreak 56,\allowbreak 16,\allowbreak 55,\allowbreak 18,\allowbreak 54,\allowbreak 21,\allowbreak 59,\allowbreak 10,\allowbreak 66,\allowbreak 7,\allowbreak 65,\allowbreak 8,\allowbreak 61,\allowbreak 13,\allowbreak 60,\allowbreak 15)\par\smallskip (68,\allowbreak 5,\allowbreak 67,\allowbreak 6,\allowbreak 70,\allowbreak 3,\allowbreak 69,\allowbreak 4,\allowbreak 64,\allowbreak 9,\allowbreak 63,\allowbreak 11,\allowbreak 62,\allowbreak 12,\allowbreak 58,\allowbreak 17,\allowbreak 52,\allowbreak 24,\allowbreak 47,\allowbreak 27,\allowbreak 48,\allowbreak 30,\allowbreak 45,\allowbreak 29,\allowbreak 43,\allowbreak 33,\allowbreak 40,\allowbreak 32,\allowbreak 41,\allowbreak 28,\allowbreak 50,\allowbreak 26,\allowbreak 51,\allowbreak 19)}\tabularnewline[5pt]
{\raggedright\arraybackslash {\ttfamily\footnotesize (4,\allowbreak 34,\allowbreak 35)}\par $K=36$\par pinned W4\par $2e\le b+8$} & {\ttfamily\footnotesize (73,\allowbreak 1,\allowbreak 72,\allowbreak 2)\par\smallskip (68,\allowbreak 5,\allowbreak 64,\allowbreak 11,\allowbreak 57,\allowbreak 20,\allowbreak 53,\allowbreak 24,\allowbreak 51,\allowbreak 27,\allowbreak 47,\allowbreak 28,\allowbreak 46,\allowbreak 25,\allowbreak 50,\allowbreak 22,\allowbreak 52,\allowbreak 21,\allowbreak 55,\allowbreak 19,\allowbreak 54,\allowbreak 15,\allowbreak 58,\allowbreak 17,\allowbreak 59,\allowbreak 14,\allowbreak 61,\allowbreak 10,\allowbreak 67,\allowbreak 9,\allowbreak 65,\allowbreak 4,\allowbreak 70,\allowbreak 3)\par\smallskip (69,\allowbreak 7,\allowbreak 71,\allowbreak 6,\allowbreak 66,\allowbreak 12,\allowbreak 60,\allowbreak 8,\allowbreak 63,\allowbreak 13,\allowbreak 62,\allowbreak 18,\allowbreak 56,\allowbreak 16,\allowbreak 48,\allowbreak 26,\allowbreak 49,\allowbreak 23,\allowbreak 40,\allowbreak 36,\allowbreak 37,\allowbreak 35,\allowbreak 38,\allowbreak 31,\allowbreak 43,\allowbreak 30,\allowbreak 44,\allowbreak 29,\allowbreak 45,\allowbreak 34,\allowbreak 39,\allowbreak 33,\allowbreak 41,\allowbreak 32,\allowbreak 42)}\tabularnewline[5pt]
{\raggedright\arraybackslash {\ttfamily\footnotesize (4,\allowbreak 34,\allowbreak 36)}\par $K=37$\par pinned W4} & {\ttfamily\footnotesize (56,\allowbreak 24,\allowbreak 51,\allowbreak 25,\allowbreak 50,\allowbreak 22,\allowbreak 53,\allowbreak 23,\allowbreak 52,\allowbreak 28,\allowbreak 46,\allowbreak 29,\allowbreak 43,\allowbreak 32,\allowbreak 42,\allowbreak 33,\allowbreak 45,\allowbreak 30,\allowbreak 63,\allowbreak 14,\allowbreak 60,\allowbreak 16,\allowbreak 61,\allowbreak 13,\allowbreak 64,\allowbreak 11,\allowbreak 65,\allowbreak 10,\allowbreak 62,\allowbreak 12,\allowbreak 59,\allowbreak 19,\allowbreak 54,\allowbreak 18,\allowbreak 57,\allowbreak 20)\par\smallskip (74,\allowbreak 1,\allowbreak 73,\allowbreak 2)\par\smallskip (70,\allowbreak 5,\allowbreak 69,\allowbreak 6,\allowbreak 72,\allowbreak 3,\allowbreak 71,\allowbreak 4,\allowbreak 66,\allowbreak 9,\allowbreak 67,\allowbreak 8,\allowbreak 68,\allowbreak 7,\allowbreak 49,\allowbreak 26,\allowbreak 48,\allowbreak 27,\allowbreak 47,\allowbreak 31,\allowbreak 44,\allowbreak 36,\allowbreak 39,\allowbreak 37,\allowbreak 38,\allowbreak 34,\allowbreak 41,\allowbreak 35,\allowbreak 40,\allowbreak 21,\allowbreak 55,\allowbreak 17,\allowbreak 58,\allowbreak 15)}\tabularnewline[5pt]
{\raggedright\arraybackslash {\ttfamily\footnotesize (4,\allowbreak 34,\allowbreak 37)}\par $K=37$\par pinned W4} & {\ttfamily\footnotesize (57,\allowbreak 18,\allowbreak 54,\allowbreak 23,\allowbreak 53,\allowbreak 21,\allowbreak 56,\allowbreak 19,\allowbreak 59,\allowbreak 16,\allowbreak 60,\allowbreak 14,\allowbreak 65,\allowbreak 12,\allowbreak 61,\allowbreak 13,\allowbreak 63,\allowbreak 11,\allowbreak 45,\allowbreak 30,\allowbreak 43,\allowbreak 31,\allowbreak 47,\allowbreak 33,\allowbreak 44,\allowbreak 34,\allowbreak 42,\allowbreak 35,\allowbreak 40,\allowbreak 36,\allowbreak 39,\allowbreak 37,\allowbreak 38,\allowbreak 32,\allowbreak 41,\allowbreak 22,\allowbreak 55)\par\smallskip (75,\allowbreak 1,\allowbreak 74,\allowbreak 2)\par\smallskip (71,\allowbreak 5,\allowbreak 70,\allowbreak 6,\allowbreak 73,\allowbreak 3,\allowbreak 72,\allowbreak 4,\allowbreak 67,\allowbreak 9,\allowbreak 64,\allowbreak 10,\allowbreak 66,\allowbreak 7,\allowbreak 69,\allowbreak 8,\allowbreak 68,\allowbreak 26,\allowbreak 50,\allowbreak 27,\allowbreak 48,\allowbreak 28,\allowbreak 46,\allowbreak 29,\allowbreak 51,\allowbreak 25,\allowbreak 52,\allowbreak 24,\allowbreak 49,\allowbreak 20,\allowbreak 58,\allowbreak 17,\allowbreak 62,\allowbreak 15)}\tabularnewline[5pt]
{\raggedright\arraybackslash {\ttfamily\footnotesize (4,\allowbreak 34,\allowbreak 38)}\par $K=38$\par pinned W4} & {\ttfamily\footnotesize (76,\allowbreak 1,\allowbreak 75,\allowbreak 2)\par\smallskip (70,\allowbreak 14,\allowbreak 60,\allowbreak 16,\allowbreak 54,\allowbreak 19,\allowbreak 56,\allowbreak 9,\allowbreak 68,\allowbreak 4,\allowbreak 73,\allowbreak 10,\allowbreak 62,\allowbreak 20,\allowbreak 51,\allowbreak 31,\allowbreak 46,\allowbreak 29,\allowbreak 43,\allowbreak 34,\allowbreak 41,\allowbreak 37,\allowbreak 45,\allowbreak 32,\allowbreak 42,\allowbreak 36,\allowbreak 39,\allowbreak 38,\allowbreak 40,\allowbreak 35,\allowbreak 47,\allowbreak 28,\allowbreak 58,\allowbreak 15)\par\smallskip (72,\allowbreak 6,\allowbreak 71,\allowbreak 3,\allowbreak 74,\allowbreak 7,\allowbreak 65,\allowbreak 5,\allowbreak 67,\allowbreak 17,\allowbreak 57,\allowbreak 8,\allowbreak 69,\allowbreak 12,\allowbreak 66,\allowbreak 11,\allowbreak 64,\allowbreak 13,\allowbreak 61,\allowbreak 22,\allowbreak 44,\allowbreak 33,\allowbreak 49,\allowbreak 26,\allowbreak 50,\allowbreak 25,\allowbreak 52,\allowbreak 24,\allowbreak 53,\allowbreak 27,\allowbreak 48,\allowbreak 30,\allowbreak 63,\allowbreak 18,\allowbreak 59,\allowbreak 23,\allowbreak 55,\allowbreak 21)}\tabularnewline[5pt]
{\raggedright\arraybackslash {\ttfamily\footnotesize (4,\allowbreak 35,\allowbreak 36)}\par $K=37$\par pinned W4} & {\ttfamily\footnotesize (75,\allowbreak 1,\allowbreak 74,\allowbreak 2)\par\smallskip (69,\allowbreak 10,\allowbreak 64,\allowbreak 14,\allowbreak 58,\allowbreak 24,\allowbreak 51,\allowbreak 19,\allowbreak 59,\allowbreak 3,\allowbreak 73,\allowbreak 7,\allowbreak 68,\allowbreak 16,\allowbreak 63,\allowbreak 15,\allowbreak 57,\allowbreak 18,\allowbreak 55,\allowbreak 22,\allowbreak 53,\allowbreak 23,\allowbreak 61,\allowbreak 12,\allowbreak 65,\allowbreak 5,\allowbreak 72,\allowbreak 4,\allowbreak 67,\allowbreak 9,\allowbreak 66,\allowbreak 11,\allowbreak 62,\allowbreak 17,\allowbreak 60)\par\smallskip (71,\allowbreak 6,\allowbreak 70,\allowbreak 8,\allowbreak 54,\allowbreak 13,\allowbreak 41,\allowbreak 29,\allowbreak 47,\allowbreak 31,\allowbreak 44,\allowbreak 36,\allowbreak 39,\allowbreak 37,\allowbreak 38,\allowbreak 34,\allowbreak 40,\allowbreak 35,\allowbreak 42,\allowbreak 33,\allowbreak 43,\allowbreak 32,\allowbreak 46,\allowbreak 25,\allowbreak 49,\allowbreak 30,\allowbreak 45,\allowbreak 28,\allowbreak 48,\allowbreak 26,\allowbreak 52,\allowbreak 27,\allowbreak 50,\allowbreak 21,\allowbreak 56,\allowbreak 20)}\tabularnewline[5pt]
{\raggedright\arraybackslash {\ttfamily\footnotesize (4,\allowbreak 35,\allowbreak 37)}\par $K=37$\par pinned W4} & {\ttfamily\footnotesize (72,\allowbreak 5,\allowbreak 71,\allowbreak 6,\allowbreak 74,\allowbreak 3,\allowbreak 73,\allowbreak 4,\allowbreak 68,\allowbreak 9,\allowbreak 65,\allowbreak 13,\allowbreak 66,\allowbreak 11,\allowbreak 45,\allowbreak 31,\allowbreak 47,\allowbreak 26,\allowbreak 69,\allowbreak 8,\allowbreak 70,\allowbreak 7,\allowbreak 67,\allowbreak 10,\allowbreak 64,\allowbreak 14,\allowbreak 63,\allowbreak 12,\allowbreak 60,\allowbreak 18,\allowbreak 59,\allowbreak 15,\allowbreak 62,\allowbreak 16,\allowbreak 61)\par\smallskip (58,\allowbreak 19,\allowbreak 57,\allowbreak 17,\allowbreak 54,\allowbreak 21,\allowbreak 56,\allowbreak 20,\allowbreak 49,\allowbreak 24,\allowbreak 52,\allowbreak 25,\allowbreak 51,\allowbreak 27,\allowbreak 50,\allowbreak 28,\allowbreak 48,\allowbreak 30,\allowbreak 43,\allowbreak 33,\allowbreak 44,\allowbreak 29,\allowbreak 46,\allowbreak 34,\allowbreak 42,\allowbreak 35,\allowbreak 40,\allowbreak 36,\allowbreak 39,\allowbreak 37,\allowbreak 38,\allowbreak 32,\allowbreak 41,\allowbreak 22,\allowbreak 53,\allowbreak 23,\allowbreak 55)\par\smallskip (76,\allowbreak 1,\allowbreak 75,\allowbreak 2)}\tabularnewline[5pt]
{\raggedright\arraybackslash {\ttfamily\footnotesize (4,\allowbreak 35,\allowbreak 38)}\par $K=38$\par pinned W4} & {\ttfamily\footnotesize (77,\allowbreak 1,\allowbreak 76,\allowbreak 2)\par\smallskip (71,\allowbreak 8,\allowbreak 70,\allowbreak 6,\allowbreak 74,\allowbreak 4,\allowbreak 69,\allowbreak 9,\allowbreak 75,\allowbreak 3,\allowbreak 72,\allowbreak 5,\allowbreak 64,\allowbreak 14,\allowbreak 59,\allowbreak 17,\allowbreak 56,\allowbreak 22,\allowbreak 54,\allowbreak 23,\allowbreak 47,\allowbreak 35,\allowbreak 40,\allowbreak 38,\allowbreak 39,\allowbreak 36,\allowbreak 42,\allowbreak 33,\allowbreak 66,\allowbreak 11,\allowbreak 65,\allowbreak 7,\allowbreak 63,\allowbreak 16,\allowbreak 60,\allowbreak 19,\allowbreak 57,\allowbreak 21)\par\smallskip (73,\allowbreak 12,\allowbreak 58,\allowbreak 18,\allowbreak 55,\allowbreak 20,\allowbreak 50,\allowbreak 28,\allowbreak 48,\allowbreak 29,\allowbreak 43,\allowbreak 32,\allowbreak 49,\allowbreak 26,\allowbreak 51,\allowbreak 30,\allowbreak 46,\allowbreak 31,\allowbreak 44,\allowbreak 34,\allowbreak 41,\allowbreak 37,\allowbreak 45,\allowbreak 27,\allowbreak 53,\allowbreak 24,\allowbreak 52,\allowbreak 25,\allowbreak 68,\allowbreak 15,\allowbreak 67,\allowbreak 10,\allowbreak 61,\allowbreak 13,\allowbreak 62)}\tabularnewline[5pt]
{\raggedright\arraybackslash {\ttfamily\footnotesize (4,\allowbreak 36,\allowbreak 37)}\par $K=38$\par pinned W4\par $2e\le b+8$} & {\ttfamily\footnotesize (77,\allowbreak 1,\allowbreak 76,\allowbreak 2)\par\smallskip (71,\allowbreak 10,\allowbreak 64,\allowbreak 16,\allowbreak 61,\allowbreak 12,\allowbreak 65,\allowbreak 14,\allowbreak 70,\allowbreak 7,\allowbreak 75,\allowbreak 3,\allowbreak 46,\allowbreak 36,\allowbreak 43,\allowbreak 34,\allowbreak 42,\allowbreak 37,\allowbreak 40,\allowbreak 38,\allowbreak 39,\allowbreak 35,\allowbreak 41,\allowbreak 30,\allowbreak 45,\allowbreak 32,\allowbreak 44,\allowbreak 27,\allowbreak 51,\allowbreak 26,\allowbreak 52,\allowbreak 25,\allowbreak 47,\allowbreak 33,\allowbreak 49,\allowbreak 31,\allowbreak 54)\par\smallskip (73,\allowbreak 6,\allowbreak 72,\allowbreak 8,\allowbreak 66,\allowbreak 11,\allowbreak 68,\allowbreak 9,\allowbreak 69,\allowbreak 4,\allowbreak 74,\allowbreak 5,\allowbreak 67,\allowbreak 15,\allowbreak 62,\allowbreak 18,\allowbreak 60,\allowbreak 19,\allowbreak 59,\allowbreak 21,\allowbreak 58,\allowbreak 24,\allowbreak 57,\allowbreak 28,\allowbreak 48,\allowbreak 29,\allowbreak 50,\allowbreak 22,\allowbreak 53,\allowbreak 23,\allowbreak 55,\allowbreak 20,\allowbreak 56,\allowbreak 17,\allowbreak 63,\allowbreak 13)}\tabularnewline[5pt]
{\raggedright\arraybackslash {\ttfamily\footnotesize (4,\allowbreak 37,\allowbreak 37)}\par $K=38$\par pinned W4\par $2e\le b+7$} & {\ttfamily\footnotesize (78,\allowbreak 1,\allowbreak 77,\allowbreak 2)\par\smallskip (73,\allowbreak 3,\allowbreak 75,\allowbreak 4,\allowbreak 70,\allowbreak 9,\allowbreak 69,\allowbreak 14,\allowbreak 61,\allowbreak 19,\allowbreak 57,\allowbreak 22,\allowbreak 52,\allowbreak 30,\allowbreak 51,\allowbreak 24,\allowbreak 49,\allowbreak 29,\allowbreak 47,\allowbreak 28,\allowbreak 44,\allowbreak 37,\allowbreak 39,\allowbreak 38,\allowbreak 41,\allowbreak 35,\allowbreak 40,\allowbreak 36,\allowbreak 46,\allowbreak 32,\allowbreak 45,\allowbreak 33,\allowbreak 48,\allowbreak 31,\allowbreak 42,\allowbreak 34,\allowbreak 43)\par\smallskip (74,\allowbreak 6,\allowbreak 71,\allowbreak 8,\allowbreak 72,\allowbreak 5,\allowbreak 67,\allowbreak 11,\allowbreak 68,\allowbreak 10,\allowbreak 64,\allowbreak 18,\allowbreak 62,\allowbreak 12,\allowbreak 63,\allowbreak 15,\allowbreak 60,\allowbreak 20,\allowbreak 56,\allowbreak 17,\allowbreak 58,\allowbreak 21,\allowbreak 54,\allowbreak 23,\allowbreak 55,\allowbreak 26,\allowbreak 50,\allowbreak 27,\allowbreak 53,\allowbreak 25,\allowbreak 59,\allowbreak 16,\allowbreak 65,\allowbreak 13,\allowbreak 66,\allowbreak 7,\allowbreak 76)}\tabularnewline[5pt]
{\raggedright\arraybackslash {\ttfamily\footnotesize (4,\allowbreak 37,\allowbreak 38)}\par $K=39$\par pinned W4} & {\ttfamily\footnotesize (79,\allowbreak 1,\allowbreak 78,\allowbreak 2)\par\smallskip (73,\allowbreak 5,\allowbreak 69,\allowbreak 17,\allowbreak 66,\allowbreak 11,\allowbreak 76,\allowbreak 4,\allowbreak 71,\allowbreak 9,\allowbreak 70,\allowbreak 7,\allowbreak 77,\allowbreak 3,\allowbreak 74,\allowbreak 14,\allowbreak 51,\allowbreak 30,\allowbreak 46,\allowbreak 34,\allowbreak 49,\allowbreak 29,\allowbreak 52,\allowbreak 27,\allowbreak 62,\allowbreak 8,\allowbreak 54,\allowbreak 28,\allowbreak 55,\allowbreak 24,\allowbreak 56,\allowbreak 23,\allowbreak 53,\allowbreak 25,\allowbreak 59,\allowbreak 18,\allowbreak 63)\par\smallskip (75,\allowbreak 6,\allowbreak 72,\allowbreak 13,\allowbreak 60,\allowbreak 20,\allowbreak 64,\allowbreak 16,\allowbreak 67,\allowbreak 10,\allowbreak 68,\allowbreak 12,\allowbreak 65,\allowbreak 15,\allowbreak 58,\allowbreak 22,\allowbreak 61,\allowbreak 19,\allowbreak 57,\allowbreak 33,\allowbreak 47,\allowbreak 37,\allowbreak 44,\allowbreak 35,\allowbreak 43,\allowbreak 38,\allowbreak 41,\allowbreak 39,\allowbreak 40,\allowbreak 36,\allowbreak 42,\allowbreak 31,\allowbreak 48,\allowbreak 26,\allowbreak 45,\allowbreak 32,\allowbreak 50,\allowbreak 21)}\tabularnewline[5pt]
{\raggedright\arraybackslash {\ttfamily\footnotesize (4,\allowbreak 37,\allowbreak 40)}\par $K=40$\par pinned W4} & {\ttfamily\footnotesize (77,\allowbreak 5,\allowbreak 76,\allowbreak 6,\allowbreak 79,\allowbreak 3,\allowbreak 78,\allowbreak 4,\allowbreak 73,\allowbreak 9,\allowbreak 72,\allowbreak 10,\allowbreak 75,\allowbreak 7,\allowbreak 74,\allowbreak 8,\allowbreak 54,\allowbreak 29,\allowbreak 52,\allowbreak 31,\allowbreak 49,\allowbreak 34,\allowbreak 50,\allowbreak 30,\allowbreak 47,\allowbreak 33,\allowbreak 46,\allowbreak 35,\allowbreak 45,\allowbreak 36,\allowbreak 48,\allowbreak 40,\allowbreak 41,\allowbreak 39,\allowbreak 42,\allowbreak 38,\allowbreak 43,\allowbreak 37,\allowbreak 44,\allowbreak 22)\par\smallskip (81,\allowbreak 1,\allowbreak 80,\allowbreak 2)\par\smallskip (61,\allowbreak 20,\allowbreak 62,\allowbreak 23,\allowbreak 58,\allowbreak 24,\allowbreak 57,\allowbreak 21,\allowbreak 59,\allowbreak 19,\allowbreak 63,\allowbreak 18,\allowbreak 67,\allowbreak 16,\allowbreak 64,\allowbreak 17,\allowbreak 60,\allowbreak 28,\allowbreak 56,\allowbreak 25,\allowbreak 55,\allowbreak 26,\allowbreak 53,\allowbreak 27,\allowbreak 51,\allowbreak 32,\allowbreak 69,\allowbreak 11,\allowbreak 71,\allowbreak 12,\allowbreak 68,\allowbreak 14,\allowbreak 66,\allowbreak 13,\allowbreak 70,\allowbreak 15,\allowbreak 65)}\tabularnewline[5pt]
{\raggedright\arraybackslash {\ttfamily\footnotesize (4,\allowbreak 42,\allowbreak 42)}\par $K=44$\par pinned W4} & {\ttfamily\footnotesize (84,\allowbreak 5,\allowbreak 83,\allowbreak 6,\allowbreak 86,\allowbreak 3,\allowbreak 85,\allowbreak 4,\allowbreak 80,\allowbreak 9,\allowbreak 79,\allowbreak 10,\allowbreak 82,\allowbreak 7,\allowbreak 81,\allowbreak 8,\allowbreak 76,\allowbreak 12,\allowbreak 74,\allowbreak 16,\allowbreak 73,\allowbreak 17,\allowbreak 72,\allowbreak 13,\allowbreak 78,\allowbreak 33,\allowbreak 56,\allowbreak 32,\allowbreak 59,\allowbreak 30,\allowbreak 58,\allowbreak 11,\allowbreak 77,\allowbreak 14,\allowbreak 75,\allowbreak 15,\allowbreak 69,\allowbreak 20,\allowbreak 68,\allowbreak 22,\allowbreak 66,\allowbreak 23)\par\smallskip (88,\allowbreak 1,\allowbreak 87,\allowbreak 2)\par\smallskip (67,\allowbreak 25,\allowbreak 64,\allowbreak 24,\allowbreak 65,\allowbreak 27,\allowbreak 62,\allowbreak 26,\allowbreak 63,\allowbreak 29,\allowbreak 60,\allowbreak 28,\allowbreak 61,\allowbreak 31,\allowbreak 57,\allowbreak 35,\allowbreak 54,\allowbreak 34,\allowbreak 55,\allowbreak 37,\allowbreak 52,\allowbreak 36,\allowbreak 53,\allowbreak 39,\allowbreak 50,\allowbreak 38,\allowbreak 51,\allowbreak 41,\allowbreak 48,\allowbreak 40,\allowbreak 49,\allowbreak 43,\allowbreak 45,\allowbreak 44,\allowbreak 47,\allowbreak 42,\allowbreak 46,\allowbreak 21,\allowbreak 71,\allowbreak 18,\allowbreak 70,\allowbreak 19)}\tabularnewline[5pt]
{\raggedright\arraybackslash {\ttfamily\footnotesize (4,\allowbreak 42,\allowbreak 43)}\par $K=44$\par pinned W4\par $2e\le b+8$} & {\ttfamily\footnotesize (89,\allowbreak 1,\allowbreak 88,\allowbreak 2)\par\smallskip (84,\allowbreak 6,\allowbreak 87,\allowbreak 8,\allowbreak 82,\allowbreak 7,\allowbreak 83,\allowbreak 10,\allowbreak 80,\allowbreak 9,\allowbreak 81,\allowbreak 4,\allowbreak 86,\allowbreak 3,\allowbreak 72,\allowbreak 12,\allowbreak 79,\allowbreak 15,\allowbreak 69,\allowbreak 18,\allowbreak 70,\allowbreak 24,\allowbreak 65,\allowbreak 30,\allowbreak 55,\allowbreak 28,\allowbreak 61,\allowbreak 21,\allowbreak 64,\allowbreak 26,\allowbreak 63,\allowbreak 27,\allowbreak 58,\allowbreak 32,\allowbreak 50,\allowbreak 35,\allowbreak 59,\allowbreak 29,\allowbreak 57,\allowbreak 36,\allowbreak 68,\allowbreak 19)\par\smallskip (85,\allowbreak 5,\allowbreak 73,\allowbreak 17,\allowbreak 76,\allowbreak 11,\allowbreak 77,\allowbreak 14,\allowbreak 71,\allowbreak 23,\allowbreak 62,\allowbreak 20,\allowbreak 75,\allowbreak 25,\allowbreak 78,\allowbreak 16,\allowbreak 74,\allowbreak 13,\allowbreak 60,\allowbreak 31,\allowbreak 53,\allowbreak 34,\allowbreak 54,\allowbreak 37,\allowbreak 47,\allowbreak 44,\allowbreak 45,\allowbreak 43,\allowbreak 48,\allowbreak 42,\allowbreak 46,\allowbreak 39,\allowbreak 51,\allowbreak 38,\allowbreak 52,\allowbreak 41,\allowbreak 49,\allowbreak 40,\allowbreak 56,\allowbreak 33,\allowbreak 67,\allowbreak 22,\allowbreak 66)}\tabularnewline[5pt]
{\raggedright\arraybackslash {\ttfamily\footnotesize (4,\allowbreak 42,\allowbreak 44)}\par $K=45$\par pinned W4} & {\ttfamily\footnotesize (90,\allowbreak 1,\allowbreak 89,\allowbreak 2)\par\smallskip (80,\allowbreak 13,\allowbreak 79,\allowbreak 11,\allowbreak 75,\allowbreak 12,\allowbreak 77,\allowbreak 22,\allowbreak 72,\allowbreak 20,\allowbreak 67,\allowbreak 19,\allowbreak 78,\allowbreak 8,\allowbreak 83,\allowbreak 7,\allowbreak 84,\allowbreak 10,\allowbreak 81,\allowbreak 9,\allowbreak 82,\allowbreak 4,\allowbreak 87,\allowbreak 3,\allowbreak 88,\allowbreak 6,\allowbreak 85,\allowbreak 16,\allowbreak 69,\allowbreak 28,\allowbreak 56,\allowbreak 34,\allowbreak 54,\allowbreak 36,\allowbreak 55,\allowbreak 31,\allowbreak 60,\allowbreak 29,\allowbreak 71,\allowbreak 17,\allowbreak 68,\allowbreak 25)\par\smallskip (86,\allowbreak 5,\allowbreak 62,\allowbreak 37,\allowbreak 48,\allowbreak 43,\allowbreak 49,\allowbreak 42,\allowbreak 46,\allowbreak 45,\allowbreak 47,\allowbreak 44,\allowbreak 52,\allowbreak 39,\allowbreak 53,\allowbreak 38,\allowbreak 50,\allowbreak 41,\allowbreak 51,\allowbreak 35,\allowbreak 58,\allowbreak 23,\allowbreak 57,\allowbreak 40,\allowbreak 61,\allowbreak 21,\allowbreak 66,\allowbreak 30,\allowbreak 63,\allowbreak 33,\allowbreak 59,\allowbreak 32,\allowbreak 64,\allowbreak 27,\allowbreak 65,\allowbreak 26,\allowbreak 70,\allowbreak 24,\allowbreak 73,\allowbreak 15,\allowbreak 76,\allowbreak 14,\allowbreak 74,\allowbreak 18)}\tabularnewline[5pt]
{\raggedright\arraybackslash {\ttfamily\footnotesize (4,\allowbreak 42,\allowbreak 45)}\par $K=45$\par pinned W4} & {\ttfamily\footnotesize (91,\allowbreak 1,\allowbreak 90,\allowbreak 2)\par\smallskip (76,\allowbreak 18,\allowbreak 66,\allowbreak 25,\allowbreak 65,\allowbreak 26,\allowbreak 68,\allowbreak 22,\allowbreak 74,\allowbreak 24,\allowbreak 69,\allowbreak 12,\allowbreak 78,\allowbreak 34,\allowbreak 52,\allowbreak 39,\allowbreak 54,\allowbreak 38,\allowbreak 57,\allowbreak 35,\allowbreak 58,\allowbreak 37,\allowbreak 49,\allowbreak 42,\allowbreak 48,\allowbreak 43,\allowbreak 47,\allowbreak 44,\allowbreak 46,\allowbreak 45,\allowbreak 53,\allowbreak 36,\allowbreak 50,\allowbreak 41,\allowbreak 51,\allowbreak 40,\allowbreak 60,\allowbreak 33,\allowbreak 59,\allowbreak 16,\allowbreak 80,\allowbreak 19)\par\smallskip (87,\allowbreak 5,\allowbreak 86,\allowbreak 6,\allowbreak 89,\allowbreak 3,\allowbreak 88,\allowbreak 4,\allowbreak 83,\allowbreak 9,\allowbreak 82,\allowbreak 10,\allowbreak 85,\allowbreak 7,\allowbreak 84,\allowbreak 15,\allowbreak 71,\allowbreak 17,\allowbreak 77,\allowbreak 14,\allowbreak 81,\allowbreak 11,\allowbreak 79,\allowbreak 8,\allowbreak 73,\allowbreak 20,\allowbreak 75,\allowbreak 13,\allowbreak 72,\allowbreak 21,\allowbreak 70,\allowbreak 23,\allowbreak 61,\allowbreak 28,\allowbreak 64,\allowbreak 29,\allowbreak 63,\allowbreak 31,\allowbreak 62,\allowbreak 32,\allowbreak 56,\allowbreak 27,\allowbreak 55,\allowbreak 30,\allowbreak 67)}\tabularnewline[5pt]
{\raggedright\arraybackslash {\ttfamily\footnotesize (4,\allowbreak 42,\allowbreak 46)}\par $K=46$\par pinned W4} & {\ttfamily\footnotesize (92,\allowbreak 1,\allowbreak 91,\allowbreak 2)\par\smallskip (87,\allowbreak 6,\allowbreak 90,\allowbreak 8,\allowbreak 77,\allowbreak 23,\allowbreak 66,\allowbreak 26,\allowbreak 67,\allowbreak 39,\allowbreak 54,\allowbreak 32,\allowbreak 61,\allowbreak 30,\allowbreak 62,\allowbreak 29,\allowbreak 64,\allowbreak 28,\allowbreak 65,\allowbreak 31,\allowbreak 58,\allowbreak 44,\allowbreak 50,\allowbreak 45,\allowbreak 48,\allowbreak 46,\allowbreak 47,\allowbreak 43,\allowbreak 53,\allowbreak 40,\allowbreak 52,\allowbreak 41,\allowbreak 49,\allowbreak 42,\allowbreak 51,\allowbreak 35,\allowbreak 60,\allowbreak 34,\allowbreak 57,\allowbreak 33,\allowbreak 63,\allowbreak 24,\allowbreak 68,\allowbreak 20,\allowbreak 76,\allowbreak 25)\par\smallskip (88,\allowbreak 5,\allowbreak 73,\allowbreak 27,\allowbreak 69,\allowbreak 22,\allowbreak 72,\allowbreak 12,\allowbreak 75,\allowbreak 17,\allowbreak 78,\allowbreak 16,\allowbreak 80,\allowbreak 3,\allowbreak 89,\allowbreak 4,\allowbreak 84,\allowbreak 9,\allowbreak 83,\allowbreak 10,\allowbreak 86,\allowbreak 7,\allowbreak 85,\allowbreak 13,\allowbreak 79,\allowbreak 14,\allowbreak 71,\allowbreak 18,\allowbreak 70,\allowbreak 21,\allowbreak 59,\allowbreak 38,\allowbreak 55,\allowbreak 37,\allowbreak 56,\allowbreak 36,\allowbreak 81,\allowbreak 11,\allowbreak 82,\allowbreak 15,\allowbreak 74,\allowbreak 19)}\tabularnewline[5pt]
{\raggedright\arraybackslash {\ttfamily\footnotesize (4,\allowbreak 43,\allowbreak 45)}\par $K=45$\par pinned W4} & {\ttfamily\footnotesize (71,\allowbreak 24,\allowbreak 61,\allowbreak 28,\allowbreak 64,\allowbreak 29,\allowbreak 63,\allowbreak 31,\allowbreak 62,\allowbreak 32,\allowbreak 60,\allowbreak 33,\allowbreak 55,\allowbreak 37,\allowbreak 57,\allowbreak 38,\allowbreak 54,\allowbreak 39,\allowbreak 52,\allowbreak 40,\allowbreak 51,\allowbreak 41,\allowbreak 50,\allowbreak 36,\allowbreak 53,\allowbreak 45,\allowbreak 46,\allowbreak 44,\allowbreak 47,\allowbreak 43,\allowbreak 48,\allowbreak 42,\allowbreak 49,\allowbreak 23,\allowbreak 69,\allowbreak 20,\allowbreak 70,\allowbreak 22,\allowbreak 66,\allowbreak 25,\allowbreak 68,\allowbreak 26,\allowbreak 65,\allowbreak 27,\allowbreak 67)\par\smallskip (88,\allowbreak 5,\allowbreak 87,\allowbreak 6,\allowbreak 90,\allowbreak 3,\allowbreak 89,\allowbreak 4,\allowbreak 84,\allowbreak 9,\allowbreak 83,\allowbreak 10,\allowbreak 86,\allowbreak 7,\allowbreak 85,\allowbreak 8,\allowbreak 80,\allowbreak 13,\allowbreak 76,\allowbreak 17,\allowbreak 75,\allowbreak 19,\allowbreak 74,\allowbreak 21,\allowbreak 72,\allowbreak 18,\allowbreak 78,\allowbreak 14,\allowbreak 79,\allowbreak 11,\allowbreak 56,\allowbreak 35,\allowbreak 58,\allowbreak 34,\allowbreak 59,\allowbreak 30,\allowbreak 82,\allowbreak 12,\allowbreak 81,\allowbreak 15,\allowbreak 77,\allowbreak 16,\allowbreak 73)\par\smallskip (92,\allowbreak 1,\allowbreak 91,\allowbreak 2)}\tabularnewline[5pt]
{\raggedright\arraybackslash {\ttfamily\footnotesize (4,\allowbreak 43,\allowbreak 46)}\par $K=46$\par pinned W4} & {\ttfamily\footnotesize (93,\allowbreak 1,\allowbreak 92,\allowbreak 2)\par\smallskip (73,\allowbreak 19,\allowbreak 72,\allowbreak 21,\allowbreak 71,\allowbreak 22,\allowbreak 70,\allowbreak 23,\allowbreak 68,\allowbreak 24,\allowbreak 67,\allowbreak 26,\allowbreak 66,\allowbreak 27,\allowbreak 69,\allowbreak 31,\allowbreak 65,\allowbreak 28,\allowbreak 64,\allowbreak 29,\allowbreak 62,\allowbreak 30,\allowbreak 61,\allowbreak 32,\allowbreak 59,\allowbreak 37,\allowbreak 54,\allowbreak 39,\allowbreak 55,\allowbreak 36,\allowbreak 57,\allowbreak 43,\allowbreak 53,\allowbreak 40,\allowbreak 52,\allowbreak 41,\allowbreak 49,\allowbreak 42,\allowbreak 51,\allowbreak 46,\allowbreak 47,\allowbreak 45,\allowbreak 48,\allowbreak 44,\allowbreak 50,\allowbreak 25)\par\smallskip (89,\allowbreak 5,\allowbreak 88,\allowbreak 6,\allowbreak 91,\allowbreak 3,\allowbreak 90,\allowbreak 4,\allowbreak 85,\allowbreak 9,\allowbreak 84,\allowbreak 10,\allowbreak 87,\allowbreak 7,\allowbreak 86,\allowbreak 8,\allowbreak 60,\allowbreak 34,\allowbreak 58,\allowbreak 38,\allowbreak 56,\allowbreak 33,\allowbreak 63,\allowbreak 35,\allowbreak 81,\allowbreak 11,\allowbreak 83,\allowbreak 12,\allowbreak 80,\allowbreak 14,\allowbreak 79,\allowbreak 18,\allowbreak 75,\allowbreak 20,\allowbreak 76,\allowbreak 17,\allowbreak 77,\allowbreak 13,\allowbreak 82,\allowbreak 15,\allowbreak 78,\allowbreak 16,\allowbreak 74)}\tabularnewline[5pt]
{\raggedright\arraybackslash {\ttfamily\footnotesize (4,\allowbreak 45,\allowbreak 45)}\par $K=46$\par pinned W4\par $2e\le b+7$} & {\ttfamily\footnotesize (94,\allowbreak 1,\allowbreak 93,\allowbreak 2)\par\smallskip (83,\allowbreak 14,\allowbreak 80,\allowbreak 8,\allowbreak 87,\allowbreak 7,\allowbreak 88,\allowbreak 10,\allowbreak 85,\allowbreak 9,\allowbreak 86,\allowbreak 4,\allowbreak 91,\allowbreak 3,\allowbreak 92,\allowbreak 6,\allowbreak 63,\allowbreak 31,\allowbreak 76,\allowbreak 17,\allowbreak 78,\allowbreak 18,\allowbreak 74,\allowbreak 21,\allowbreak 69,\allowbreak 22,\allowbreak 64,\allowbreak 30,\allowbreak 67,\allowbreak 27,\allowbreak 66,\allowbreak 28,\allowbreak 61,\allowbreak 33,\allowbreak 60,\allowbreak 29,\allowbreak 58,\allowbreak 36,\allowbreak 56,\allowbreak 39,\allowbreak 65,\allowbreak 24,\allowbreak 68,\allowbreak 25,\allowbreak 71)\par\smallskip (90,\allowbreak 5,\allowbreak 89,\allowbreak 15,\allowbreak 77,\allowbreak 13,\allowbreak 81,\allowbreak 16,\allowbreak 79,\allowbreak 12,\allowbreak 82,\allowbreak 11,\allowbreak 84,\allowbreak 26,\allowbreak 62,\allowbreak 32,\allowbreak 57,\allowbreak 34,\allowbreak 53,\allowbreak 37,\allowbreak 55,\allowbreak 40,\allowbreak 52,\allowbreak 41,\allowbreak 54,\allowbreak 44,\allowbreak 51,\allowbreak 42,\allowbreak 50,\allowbreak 46,\allowbreak 47,\allowbreak 45,\allowbreak 48,\allowbreak 43,\allowbreak 49,\allowbreak 35,\allowbreak 59,\allowbreak 38,\allowbreak 73,\allowbreak 19,\allowbreak 70,\allowbreak 20,\allowbreak 75,\allowbreak 23,\allowbreak 72)}\tabularnewline[5pt]
{\raggedright\arraybackslash {\ttfamily\footnotesize (4,\allowbreak 45,\allowbreak 46)}\par $K=47$\par pinned W4} & {\ttfamily\footnotesize (95,\allowbreak 1,\allowbreak 94,\allowbreak 2)\par\smallskip (90,\allowbreak 6,\allowbreak 93,\allowbreak 8,\allowbreak 81,\allowbreak 20,\allowbreak 70,\allowbreak 26,\allowbreak 72,\allowbreak 18,\allowbreak 78,\allowbreak 27,\allowbreak 68,\allowbreak 28,\allowbreak 67,\allowbreak 33,\allowbreak 62,\allowbreak 34,\allowbreak 61,\allowbreak 30,\allowbreak 63,\allowbreak 41,\allowbreak 60,\allowbreak 36,\allowbreak 54,\allowbreak 42,\allowbreak 53,\allowbreak 39,\allowbreak 56,\allowbreak 40,\allowbreak 55,\allowbreak 45,\allowbreak 52,\allowbreak 43,\allowbreak 51,\allowbreak 46,\allowbreak 49,\allowbreak 47,\allowbreak 48,\allowbreak 44,\allowbreak 50,\allowbreak 37,\allowbreak 84,\allowbreak 17,\allowbreak 75)\par\smallskip (91,\allowbreak 5,\allowbreak 80,\allowbreak 12,\allowbreak 74,\allowbreak 15,\allowbreak 85,\allowbreak 13,\allowbreak 82,\allowbreak 11,\allowbreak 77,\allowbreak 21,\allowbreak 76,\allowbreak 23,\allowbreak 66,\allowbreak 29,\allowbreak 65,\allowbreak 35,\allowbreak 58,\allowbreak 38,\allowbreak 59,\allowbreak 24,\allowbreak 69,\allowbreak 31,\allowbreak 57,\allowbreak 32,\allowbreak 64,\allowbreak 22,\allowbreak 71,\allowbreak 19,\allowbreak 83,\allowbreak 3,\allowbreak 92,\allowbreak 4,\allowbreak 87,\allowbreak 9,\allowbreak 86,\allowbreak 10,\allowbreak 89,\allowbreak 7,\allowbreak 88,\allowbreak 14,\allowbreak 79,\allowbreak 16,\allowbreak 73,\allowbreak 25)}\tabularnewline[5pt]
{\raggedright\arraybackslash {\ttfamily\footnotesize (4,\allowbreak 45,\allowbreak 47)}\par $K=47$\par pinned W4} & {\ttfamily\footnotesize (96,\allowbreak 1,\allowbreak 95,\allowbreak 2)\par\smallskip (86,\allowbreak 18,\allowbreak 79,\allowbreak 19,\allowbreak 75,\allowbreak 12,\allowbreak 83,\allowbreak 17,\allowbreak 81,\allowbreak 22,\allowbreak 77,\allowbreak 26,\allowbreak 71,\allowbreak 24,\allowbreak 74,\allowbreak 20,\allowbreak 68,\allowbreak 25,\allowbreak 69,\allowbreak 27,\allowbreak 73,\allowbreak 21,\allowbreak 62,\allowbreak 34,\allowbreak 63,\allowbreak 33,\allowbreak 65,\allowbreak 32,\allowbreak 66,\allowbreak 30,\allowbreak 67,\allowbreak 29,\allowbreak 64,\allowbreak 40,\allowbreak 57,\allowbreak 38,\allowbreak 56,\allowbreak 16,\allowbreak 78,\allowbreak 11,\allowbreak 84,\allowbreak 14,\allowbreak 72,\allowbreak 23,\allowbreak 76)\par\smallskip (92,\allowbreak 5,\allowbreak 80,\allowbreak 15,\allowbreak 91,\allowbreak 6,\allowbreak 94,\allowbreak 3,\allowbreak 93,\allowbreak 4,\allowbreak 88,\allowbreak 9,\allowbreak 87,\allowbreak 10,\allowbreak 90,\allowbreak 7,\allowbreak 89,\allowbreak 8,\allowbreak 82,\allowbreak 13,\allowbreak 85,\allowbreak 28,\allowbreak 59,\allowbreak 39,\allowbreak 60,\allowbreak 37,\allowbreak 52,\allowbreak 41,\allowbreak 55,\allowbreak 42,\allowbreak 54,\allowbreak 44,\allowbreak 53,\allowbreak 45,\allowbreak 50,\allowbreak 43,\allowbreak 49,\allowbreak 46,\allowbreak 48,\allowbreak 47,\allowbreak 51,\allowbreak 35,\allowbreak 61,\allowbreak 36,\allowbreak 58,\allowbreak 31,\allowbreak 70)}\tabularnewline[5pt]
{\raggedright\arraybackslash {\ttfamily\footnotesize (4,\allowbreak 45,\allowbreak 48)}\par $K=48$\par pinned W4} & {\ttfamily\footnotesize (97,\allowbreak 1,\allowbreak 96,\allowbreak 2)\par\smallskip (87,\allowbreak 12,\allowbreak 86,\allowbreak 15,\allowbreak 92,\allowbreak 6,\allowbreak 95,\allowbreak 3,\allowbreak 94,\allowbreak 4,\allowbreak 89,\allowbreak 9,\allowbreak 88,\allowbreak 10,\allowbreak 91,\allowbreak 7,\allowbreak 90,\allowbreak 8,\allowbreak 84,\allowbreak 11,\allowbreak 81,\allowbreak 13,\allowbreak 85,\allowbreak 16,\allowbreak 83,\allowbreak 17,\allowbreak 82,\allowbreak 18,\allowbreak 74,\allowbreak 24,\allowbreak 71,\allowbreak 27,\allowbreak 70,\allowbreak 25,\allowbreak 73,\allowbreak 32,\allowbreak 68,\allowbreak 29,\allowbreak 67,\allowbreak 30,\allowbreak 65,\allowbreak 31,\allowbreak 64,\allowbreak 33,\allowbreak 63,\allowbreak 23,\allowbreak 72,\allowbreak 26)\par\smallskip (93,\allowbreak 5,\allowbreak 66,\allowbreak 37,\allowbreak 69,\allowbreak 44,\allowbreak 60,\allowbreak 34,\allowbreak 62,\allowbreak 35,\allowbreak 59,\allowbreak 39,\allowbreak 58,\allowbreak 36,\allowbreak 57,\allowbreak 40,\allowbreak 52,\allowbreak 45,\allowbreak 54,\allowbreak 48,\allowbreak 49,\allowbreak 47,\allowbreak 50,\allowbreak 46,\allowbreak 51,\allowbreak 43,\allowbreak 53,\allowbreak 42,\allowbreak 55,\allowbreak 41,\allowbreak 56,\allowbreak 38,\allowbreak 61,\allowbreak 19,\allowbreak 79,\allowbreak 28,\allowbreak 80,\allowbreak 21,\allowbreak 78,\allowbreak 20,\allowbreak 75,\allowbreak 22,\allowbreak 76,\allowbreak 14,\allowbreak 77)}\tabularnewline[5pt]
{\raggedright\arraybackslash {\ttfamily\footnotesize (4,\allowbreak 46,\allowbreak 50)}\par $K=50$\par pinned W4} & {\ttfamily\footnotesize (100,\allowbreak 1,\allowbreak 99,\allowbreak 2)\par\smallskip (96,\allowbreak 5,\allowbreak 95,\allowbreak 6,\allowbreak 98,\allowbreak 3,\allowbreak 97,\allowbreak 4,\allowbreak 92,\allowbreak 9,\allowbreak 91,\allowbreak 10,\allowbreak 94,\allowbreak 7,\allowbreak 93,\allowbreak 8,\allowbreak 88,\allowbreak 15,\allowbreak 86,\allowbreak 17,\allowbreak 83,\allowbreak 18,\allowbreak 85,\allowbreak 13,\allowbreak 87,\allowbreak 36,\allowbreak 63,\allowbreak 37,\allowbreak 65,\allowbreak 11,\allowbreak 90,\allowbreak 12,\allowbreak 89,\allowbreak 14,\allowbreak 84,\allowbreak 16,\allowbreak 80,\allowbreak 22,\allowbreak 81,\allowbreak 20,\allowbreak 82,\allowbreak 19,\allowbreak 79,\allowbreak 23,\allowbreak 78,\allowbreak 21)\par\smallskip (76,\allowbreak 28,\allowbreak 72,\allowbreak 29,\allowbreak 74,\allowbreak 24,\allowbreak 77,\allowbreak 25,\allowbreak 54,\allowbreak 48,\allowbreak 51,\allowbreak 50,\allowbreak 52,\allowbreak 47,\allowbreak 57,\allowbreak 49,\allowbreak 53,\allowbreak 46,\allowbreak 55,\allowbreak 44,\allowbreak 56,\allowbreak 43,\allowbreak 58,\allowbreak 42,\allowbreak 60,\allowbreak 38,\allowbreak 62,\allowbreak 45,\allowbreak 59,\allowbreak 40,\allowbreak 61,\allowbreak 41,\allowbreak 64,\allowbreak 39,\allowbreak 69,\allowbreak 35,\allowbreak 66,\allowbreak 34,\allowbreak 67,\allowbreak 32,\allowbreak 68,\allowbreak 31,\allowbreak 70,\allowbreak 30,\allowbreak 71,\allowbreak 33,\allowbreak 75,\allowbreak 26,\allowbreak 73,\allowbreak 27)}\tabularnewline[5pt]
{\raggedright\arraybackslash {\ttfamily\footnotesize (4,\allowbreak 47,\allowbreak 50)}\par $K=50$\par pinned W4} & {\ttfamily\footnotesize (101,\allowbreak 1,\allowbreak 100,\allowbreak 2)\par\smallskip (97,\allowbreak 5,\allowbreak 96,\allowbreak 6,\allowbreak 99,\allowbreak 3,\allowbreak 98,\allowbreak 4,\allowbreak 93,\allowbreak 9,\allowbreak 92,\allowbreak 10,\allowbreak 95,\allowbreak 7,\allowbreak 94,\allowbreak 8,\allowbreak 89,\allowbreak 11,\allowbreak 91,\allowbreak 12,\allowbreak 87,\allowbreak 17,\allowbreak 84,\allowbreak 18,\allowbreak 86,\allowbreak 15,\allowbreak 88,\allowbreak 14,\allowbreak 90,\allowbreak 13,\allowbreak 69,\allowbreak 35,\allowbreak 66,\allowbreak 34,\allowbreak 67,\allowbreak 32,\allowbreak 68,\allowbreak 31,\allowbreak 70,\allowbreak 30,\allowbreak 71,\allowbreak 33,\allowbreak 75,\allowbreak 25,\allowbreak 76,\allowbreak 23,\allowbreak 80)\par\smallskip (72,\allowbreak 29,\allowbreak 73,\allowbreak 28,\allowbreak 77,\allowbreak 22,\allowbreak 82,\allowbreak 19,\allowbreak 83,\allowbreak 21,\allowbreak 79,\allowbreak 20,\allowbreak 81,\allowbreak 16,\allowbreak 85,\allowbreak 39,\allowbreak 61,\allowbreak 38,\allowbreak 63,\allowbreak 37,\allowbreak 64,\allowbreak 36,\allowbreak 65,\allowbreak 41,\allowbreak 62,\allowbreak 45,\allowbreak 59,\allowbreak 40,\allowbreak 60,\allowbreak 42,\allowbreak 58,\allowbreak 43,\allowbreak 56,\allowbreak 44,\allowbreak 55,\allowbreak 46,\allowbreak 53,\allowbreak 49,\allowbreak 57,\allowbreak 47,\allowbreak 52,\allowbreak 50,\allowbreak 51,\allowbreak 48,\allowbreak 54,\allowbreak 24,\allowbreak 78,\allowbreak 26,\allowbreak 74,\allowbreak 27)}\tabularnewline[5pt]
{\raggedright\arraybackslash {\ttfamily\footnotesize (4,\allowbreak 54,\allowbreak 57)}\par $K=57$\par pinned W4} & {\ttfamily\footnotesize (115,\allowbreak 1,\allowbreak 114,\allowbreak 2)\par\smallskip (87,\allowbreak 28,\allowbreak 88,\allowbreak 27,\allowbreak 92,\allowbreak 24,\allowbreak 93,\allowbreak 23,\allowbreak 90,\allowbreak 26,\allowbreak 89,\allowbreak 31,\allowbreak 84,\allowbreak 29,\allowbreak 86,\allowbreak 30,\allowbreak 79,\allowbreak 34,\allowbreak 82,\allowbreak 35,\allowbreak 81,\allowbreak 37,\allowbreak 80,\allowbreak 38,\allowbreak 78,\allowbreak 39,\allowbreak 77,\allowbreak 36,\allowbreak 73,\allowbreak 40,\allowbreak 76,\allowbreak 41,\allowbreak 75,\allowbreak 44,\allowbreak 69,\allowbreak 45,\allowbreak 99,\allowbreak 17,\allowbreak 98,\allowbreak 18,\allowbreak 101,\allowbreak 15,\allowbreak 100,\allowbreak 16,\allowbreak 95,\allowbreak 21,\allowbreak 94,\allowbreak 22,\allowbreak 97,\allowbreak 19,\allowbreak 96,\allowbreak 20,\allowbreak 91,\allowbreak 25)\par\smallskip (111,\allowbreak 5,\allowbreak 110,\allowbreak 6,\allowbreak 113,\allowbreak 3,\allowbreak 112,\allowbreak 4,\allowbreak 107,\allowbreak 9,\allowbreak 106,\allowbreak 10,\allowbreak 109,\allowbreak 7,\allowbreak 108,\allowbreak 8,\allowbreak 103,\allowbreak 13,\allowbreak 102,\allowbreak 14,\allowbreak 105,\allowbreak 11,\allowbreak 104,\allowbreak 12,\allowbreak 74,\allowbreak 42,\allowbreak 72,\allowbreak 46,\allowbreak 67,\allowbreak 47,\allowbreak 70,\allowbreak 43,\allowbreak 71,\allowbreak 49,\allowbreak 68,\allowbreak 50,\allowbreak 66,\allowbreak 51,\allowbreak 64,\allowbreak 52,\allowbreak 63,\allowbreak 53,\allowbreak 62,\allowbreak 48,\allowbreak 65,\allowbreak 57,\allowbreak 58,\allowbreak 56,\allowbreak 59,\allowbreak 55,\allowbreak 60,\allowbreak 54,\allowbreak 61,\allowbreak 32,\allowbreak 83,\allowbreak 33,\allowbreak 85)}\tabularnewline[5pt]
{\raggedright\arraybackslash {\ttfamily\footnotesize (4,\allowbreak 55,\allowbreak 57)}\par $K=57$\par pinned W4} & {\ttfamily\footnotesize (88,\allowbreak 28,\allowbreak 89,\allowbreak 23,\allowbreak 94,\allowbreak 24,\allowbreak 93,\allowbreak 25,\allowbreak 92,\allowbreak 20,\allowbreak 97,\allowbreak 19,\allowbreak 98,\allowbreak 22,\allowbreak 95,\allowbreak 21,\allowbreak 96,\allowbreak 16,\allowbreak 101,\allowbreak 15,\allowbreak 102,\allowbreak 18,\allowbreak 99,\allowbreak 17,\allowbreak 100,\allowbreak 42,\allowbreak 71,\allowbreak 44,\allowbreak 72,\allowbreak 47,\allowbreak 67,\allowbreak 45,\allowbreak 75,\allowbreak 41,\allowbreak 76,\allowbreak 40,\allowbreak 73,\allowbreak 36,\allowbreak 77,\allowbreak 39,\allowbreak 78,\allowbreak 38,\allowbreak 80,\allowbreak 37,\allowbreak 81,\allowbreak 35,\allowbreak 82,\allowbreak 34,\allowbreak 79,\allowbreak 30,\allowbreak 83,\allowbreak 33,\allowbreak 84,\allowbreak 32,\allowbreak 87,\allowbreak 31,\allowbreak 85)\par\smallskip (112,\allowbreak 5,\allowbreak 111,\allowbreak 6,\allowbreak 114,\allowbreak 3,\allowbreak 113,\allowbreak 4,\allowbreak 108,\allowbreak 9,\allowbreak 107,\allowbreak 10,\allowbreak 110,\allowbreak 7,\allowbreak 109,\allowbreak 8,\allowbreak 104,\allowbreak 13,\allowbreak 103,\allowbreak 14,\allowbreak 106,\allowbreak 11,\allowbreak 105,\allowbreak 12,\allowbreak 74,\allowbreak 43,\allowbreak 69,\allowbreak 46,\allowbreak 70,\allowbreak 49,\allowbreak 68,\allowbreak 50,\allowbreak 66,\allowbreak 51,\allowbreak 64,\allowbreak 52,\allowbreak 63,\allowbreak 53,\allowbreak 62,\allowbreak 48,\allowbreak 65,\allowbreak 57,\allowbreak 58,\allowbreak 56,\allowbreak 59,\allowbreak 55,\allowbreak 60,\allowbreak 54,\allowbreak 61,\allowbreak 29,\allowbreak 86,\allowbreak 27,\allowbreak 90,\allowbreak 26,\allowbreak 91)\par\smallskip (116,\allowbreak 1,\allowbreak 115,\allowbreak 2)}\tabularnewline[5pt]
\end{longtable}
\endgroup

\subsection{The 20 small minimum-four/five helper cores}
\label{aux:min45}
\begingroup\footnotesize
\begin{longtable}{@{}p{25mm}p{115mm}@{}}
Profile / interface & Complete outward words\\\hline\endhead
{\raggedright\arraybackslash {\ttfamily\footnotesize (4,\allowbreak 4,\allowbreak 7)}\par $K=7$} & {\ttfamily\footnotesize (15,\allowbreak 1,\allowbreak 14,\allowbreak 2)\par\smallskip (11,\allowbreak 4,\allowbreak 9,\allowbreak 7,\allowbreak 8,\allowbreak 5,\allowbreak 13)\par\smallskip (10,\allowbreak 6,\allowbreak 12,\allowbreak 3)}\tabularnewline[5pt]
{\raggedright\arraybackslash {\ttfamily\footnotesize (4,\allowbreak 4,\allowbreak 8)}\par $K=8$} & {\ttfamily\footnotesize (16,\allowbreak 2,\allowbreak 14,\allowbreak 1)\par\smallskip (15,\allowbreak 5,\allowbreak 13,\allowbreak 6,\allowbreak 9,\allowbreak 8,\allowbreak 10,\allowbreak 4)\par\smallskip (11,\allowbreak 7,\allowbreak 12,\allowbreak 3)}\tabularnewline[5pt]
{\raggedright\arraybackslash {\ttfamily\footnotesize (4,\allowbreak 4,\allowbreak 17)}\par $K=12$} & {\ttfamily\footnotesize (25,\allowbreak 1,\allowbreak 24,\allowbreak 2)\par\smallskip (21,\allowbreak 3,\allowbreak 23,\allowbreak 4)\par\smallskip (17,\allowbreak 5,\allowbreak 20,\allowbreak 6,\allowbreak 22,\allowbreak 9,\allowbreak 16,\allowbreak 7,\allowbreak 18,\allowbreak 8,\allowbreak 14,\allowbreak 10,\allowbreak 15,\allowbreak 12,\allowbreak 13,\allowbreak 11,\allowbreak 19)}\tabularnewline[5pt]
{\raggedright\arraybackslash {\ttfamily\footnotesize (4,\allowbreak 5,\allowbreak 5)}\par $K=6$\par pinned W4\par $2e\le b+7$} & {\ttfamily\footnotesize (14,\allowbreak 1,\allowbreak 13,\allowbreak 2)\par\smallskip (10,\allowbreak 3,\allowbreak 8,\allowbreak 4,\allowbreak 12)\par\smallskip (9,\allowbreak 6,\allowbreak 7,\allowbreak 5,\allowbreak 11)}\tabularnewline[5pt]
{\raggedright\arraybackslash {\ttfamily\footnotesize (4,\allowbreak 5,\allowbreak 7)}\par $K=7$} & {\ttfamily\footnotesize (16,\allowbreak 1,\allowbreak 13,\allowbreak 3)\par\smallskip (14,\allowbreak 5,\allowbreak 9,\allowbreak 6,\allowbreak 8,\allowbreak 7,\allowbreak 12)\par\smallskip (11,\allowbreak 4,\allowbreak 10,\allowbreak 2,\allowbreak 15)}\tabularnewline[5pt]
{\raggedright\arraybackslash {\ttfamily\footnotesize (4,\allowbreak 5,\allowbreak 8)}\par $K=8$\par pinned W4} & {\ttfamily\footnotesize (17,\allowbreak 1,\allowbreak 16,\allowbreak 2)\par\smallskip (13,\allowbreak 4,\allowbreak 12,\allowbreak 6,\allowbreak 10,\allowbreak 5,\allowbreak 15,\allowbreak 3)\par\smallskip (11,\allowbreak 8,\allowbreak 9,\allowbreak 7,\allowbreak 14)}\tabularnewline[5pt]
{\raggedright\arraybackslash {\ttfamily\footnotesize (4,\allowbreak 6,\allowbreak 7)}\par $K=8$\par pinned W4\par $2e\le b+8$} & {\ttfamily\footnotesize (17,\allowbreak 1,\allowbreak 16,\allowbreak 2)\par\smallskip (13,\allowbreak 4,\allowbreak 15,\allowbreak 5,\allowbreak 11,\allowbreak 6,\allowbreak 14)\par\smallskip (12,\allowbreak 8,\allowbreak 9,\allowbreak 7,\allowbreak 10,\allowbreak 3)}\tabularnewline[5pt]
{\raggedright\arraybackslash {\ttfamily\footnotesize (4,\allowbreak 6,\allowbreak 8)}\par $K=9$\par pinned W4} & {\ttfamily\footnotesize (18,\allowbreak 1,\allowbreak 17,\allowbreak 2)\par\smallskip (14,\allowbreak 5,\allowbreak 16,\allowbreak 6,\allowbreak 12,\allowbreak 7,\allowbreak 15,\allowbreak 3)\par\smallskip (13,\allowbreak 9,\allowbreak 10,\allowbreak 8,\allowbreak 11,\allowbreak 4)}\tabularnewline[5pt]
{\raggedright\arraybackslash {\ttfamily\footnotesize (4,\allowbreak 6,\allowbreak 9)}\par $K=9$\par pinned W4} & {\ttfamily\footnotesize (19,\allowbreak 1,\allowbreak 18,\allowbreak 2)\par\smallskip (15,\allowbreak 3,\allowbreak 17,\allowbreak 6,\allowbreak 11,\allowbreak 9,\allowbreak 10,\allowbreak 7,\allowbreak 16)\par\smallskip (13,\allowbreak 5,\allowbreak 12,\allowbreak 8,\allowbreak 14,\allowbreak 4)}\tabularnewline[5pt]
{\raggedright\arraybackslash {\ttfamily\footnotesize (4,\allowbreak 6,\allowbreak 10)}\par $K=10$\par pinned W4} & {\ttfamily\footnotesize (20,\allowbreak 1,\allowbreak 19,\allowbreak 2)\par\smallskip (16,\allowbreak 4,\allowbreak 17,\allowbreak 6,\allowbreak 15,\allowbreak 7,\allowbreak 13,\allowbreak 8,\allowbreak 18,\allowbreak 3)\par\smallskip (14,\allowbreak 10,\allowbreak 11,\allowbreak 9,\allowbreak 12,\allowbreak 5)}\tabularnewline[5pt]
{\raggedright\arraybackslash {\ttfamily\footnotesize (4,\allowbreak 7,\allowbreak 7)}\par $K=8$\par pinned W4\par $2e\le b+7$} & {\ttfamily\footnotesize (18,\allowbreak 1,\allowbreak 17,\allowbreak 2)\par\smallskip (14,\allowbreak 4,\allowbreak 15,\allowbreak 6,\allowbreak 11,\allowbreak 3,\allowbreak 16)\par\smallskip (12,\allowbreak 5,\allowbreak 9,\allowbreak 8,\allowbreak 10,\allowbreak 7,\allowbreak 13)}\tabularnewline[5pt]
{\raggedright\arraybackslash {\ttfamily\footnotesize (4,\allowbreak 7,\allowbreak 8)}\par $K=9$\par pinned W4} & {\ttfamily\footnotesize (19,\allowbreak 1,\allowbreak 18,\allowbreak 2)\par\smallskip (15,\allowbreak 3,\allowbreak 17,\allowbreak 7,\allowbreak 10,\allowbreak 9,\allowbreak 11,\allowbreak 4)\par\smallskip (13,\allowbreak 8,\allowbreak 12,\allowbreak 6,\allowbreak 14,\allowbreak 5,\allowbreak 16)}\tabularnewline[5pt]
{\raggedright\arraybackslash {\ttfamily\footnotesize (4,\allowbreak 7,\allowbreak 9)}\par $K=9$\par pinned W4} & {\ttfamily\footnotesize (20,\allowbreak 1,\allowbreak 19,\allowbreak 2)\par\smallskip (16,\allowbreak 3,\allowbreak 18,\allowbreak 4,\allowbreak 14,\allowbreak 5,\allowbreak 13,\allowbreak 6,\allowbreak 17)\par\smallskip (12,\allowbreak 7,\allowbreak 11,\allowbreak 8,\allowbreak 10,\allowbreak 9,\allowbreak 15)}\tabularnewline[5pt]
{\raggedright\arraybackslash {\ttfamily\footnotesize (4,\allowbreak 7,\allowbreak 10)}\par $K=10$\par pinned W4} & {\ttfamily\footnotesize (21,\allowbreak 1,\allowbreak 20,\allowbreak 2)\par\smallskip (17,\allowbreak 3,\allowbreak 19,\allowbreak 4,\allowbreak 16,\allowbreak 8,\allowbreak 12,\allowbreak 6,\allowbreak 15,\allowbreak 5)\par\smallskip (13,\allowbreak 10,\allowbreak 11,\allowbreak 9,\allowbreak 14,\allowbreak 7,\allowbreak 18)}\tabularnewline[5pt]
{\raggedright\arraybackslash {\ttfamily\footnotesize (4,\allowbreak 8,\allowbreak 8)}\par $K=10$\par pinned W4} & {\ttfamily\footnotesize (20,\allowbreak 1,\allowbreak 19,\allowbreak 2)\par\smallskip (16,\allowbreak 4,\allowbreak 17,\allowbreak 6,\allowbreak 15,\allowbreak 8,\allowbreak 18,\allowbreak 3)\par\smallskip (14,\allowbreak 9,\allowbreak 12,\allowbreak 10,\allowbreak 11,\allowbreak 7,\allowbreak 13,\allowbreak 5)}\tabularnewline[5pt]
{\raggedright\arraybackslash {\ttfamily\footnotesize (4,\allowbreak 8,\allowbreak 9)}\par $K=10$\par pinned W4\par $2e\le b+8$} & {\ttfamily\footnotesize (21,\allowbreak 1,\allowbreak 20,\allowbreak 2)\par\smallskip (17,\allowbreak 3,\allowbreak 19,\allowbreak 4,\allowbreak 16,\allowbreak 6,\allowbreak 15,\allowbreak 7,\allowbreak 18)\par\smallskip (13,\allowbreak 8,\allowbreak 14,\allowbreak 10,\allowbreak 11,\allowbreak 9,\allowbreak 12,\allowbreak 5)}\tabularnewline[5pt]
{\raggedright\arraybackslash {\ttfamily\footnotesize (4,\allowbreak 8,\allowbreak 10)}\par $K=11$\par pinned W4} & {\ttfamily\footnotesize (22,\allowbreak 1,\allowbreak 21,\allowbreak 2)\par\smallskip (18,\allowbreak 3,\allowbreak 20,\allowbreak 6,\allowbreak 19,\allowbreak 7,\allowbreak 17,\allowbreak 8,\allowbreak 15,\allowbreak 4)\par\smallskip (16,\allowbreak 10,\allowbreak 14,\allowbreak 9,\allowbreak 12,\allowbreak 11,\allowbreak 13,\allowbreak 5)}\tabularnewline[5pt]
{\raggedright\arraybackslash {\ttfamily\footnotesize (4,\allowbreak 8,\allowbreak 11)}\par $K=11$\par pinned W4} & {\ttfamily\footnotesize (23,\allowbreak 1,\allowbreak 22,\allowbreak 2)\par\smallskip (19,\allowbreak 3,\allowbreak 21,\allowbreak 4,\allowbreak 18,\allowbreak 5,\allowbreak 17,\allowbreak 7,\allowbreak 16,\allowbreak 9,\allowbreak 20)\par\smallskip (15,\allowbreak 10,\allowbreak 13,\allowbreak 11,\allowbreak 12,\allowbreak 8,\allowbreak 14,\allowbreak 6)}\tabularnewline[5pt]
{\raggedright\arraybackslash {\ttfamily\footnotesize (4,\allowbreak 8,\allowbreak 12)}\par $K=12$\par pinned W4} & {\ttfamily\footnotesize (24,\allowbreak 1,\allowbreak 23,\allowbreak 2)\par\smallskip (20,\allowbreak 3,\allowbreak 22,\allowbreak 7,\allowbreak 21,\allowbreak 5,\allowbreak 17,\allowbreak 8,\allowbreak 16,\allowbreak 9,\allowbreak 19,\allowbreak 6)\par\smallskip (18,\allowbreak 12,\allowbreak 13,\allowbreak 11,\allowbreak 14,\allowbreak 10,\allowbreak 15,\allowbreak 4)}\tabularnewline[5pt]
{\raggedright\arraybackslash {\ttfamily\footnotesize (5,\allowbreak 5,\allowbreak 7)}\par $K=7$} & {\ttfamily\footnotesize (17,\allowbreak 1,\allowbreak 16,\allowbreak 2,\allowbreak 14)\par\smallskip (13,\allowbreak 3,\allowbreak 11,\allowbreak 6,\allowbreak 10,\allowbreak 4,\allowbreak 15)\par\smallskip (9,\allowbreak 7,\allowbreak 8,\allowbreak 5,\allowbreak 12)}\tabularnewline[5pt]
\end{longtable}
\endgroup

\subsection{Fixed-endpoint alpha path prefixes and bases}
\label{aux:endpoints}
\begingroup\footnotesize
\begin{longtable}{@{}p{25mm}p{115mm}@{}}
Profile / interface & Complete path permutation (no extra hub)\\\hline\endhead
{\raggedright\arraybackslash $e=1$, prefix\par $n=4$} & {\ttfamily\footnotesize (1,\allowbreak 2,\allowbreak 0,\allowbreak 3)}\tabularnewline[5pt]
{\raggedright\arraybackslash $e=1$, $n=3$\par alpha path} & {\ttfamily\footnotesize (1,\allowbreak 2,\allowbreak 0)}\tabularnewline[5pt]
{\raggedright\arraybackslash $e=1$, $n=4$\par alpha path} & {\ttfamily\footnotesize (1,\allowbreak 2,\allowbreak 0,\allowbreak 3)}\tabularnewline[5pt]
{\raggedright\arraybackslash $e=1$, $n=6$\par alpha path} & {\ttfamily\footnotesize (1,\allowbreak 5,\allowbreak 0,\allowbreak 3,\allowbreak 2,\allowbreak 4)}\tabularnewline[5pt]
{\raggedright\arraybackslash $e=1$, $n=9$\par alpha path} & {\ttfamily\footnotesize (1,\allowbreak 8,\allowbreak 0,\allowbreak 6,\allowbreak 2,\allowbreak 7,\allowbreak 4,\allowbreak 5,\allowbreak 3)}\tabularnewline[5pt]
{\raggedright\arraybackslash $e=2$, prefix\par $n=8$} & {\ttfamily\footnotesize (2,\allowbreak 5,\allowbreak 3,\allowbreak 4,\allowbreak 0,\allowbreak 7,\allowbreak 1,\allowbreak 6)}\tabularnewline[5pt]
{\raggedright\arraybackslash $e=2$, $n=5$\par alpha path} & {\ttfamily\footnotesize (2,\allowbreak 3,\allowbreak 1,\allowbreak 4,\allowbreak 0)}\tabularnewline[5pt]
{\raggedright\arraybackslash $e=2$, $n=6$\par alpha path} & {\ttfamily\footnotesize (2,\allowbreak 3,\allowbreak 1,\allowbreak 4,\allowbreak 0,\allowbreak 5)}\tabularnewline[5pt]
{\raggedright\arraybackslash $e=2$, $n=7$\par alpha path} & {\ttfamily\footnotesize (2,\allowbreak 6,\allowbreak 0,\allowbreak 5,\allowbreak 3,\allowbreak 4,\allowbreak 1)}\tabularnewline[5pt]
{\raggedright\arraybackslash $e=2$, $n=8$\par alpha path} & {\ttfamily\footnotesize (2,\allowbreak 5,\allowbreak 3,\allowbreak 4,\allowbreak 0,\allowbreak 7,\allowbreak 1,\allowbreak 6)}\tabularnewline[5pt]
{\raggedright\arraybackslash $e=2$, $n=10$\par alpha path} & {\ttfamily\footnotesize (2,\allowbreak 9,\allowbreak 0,\allowbreak 8,\allowbreak 3,\allowbreak 6,\allowbreak 4,\allowbreak 5,\allowbreak 1,\allowbreak 7)}\tabularnewline[5pt]
{\raggedright\arraybackslash $e=2$, $n=11$\par alpha path} & {\ttfamily\footnotesize (2,\allowbreak 10,\allowbreak 0,\allowbreak 9,\allowbreak 4,\allowbreak 8,\allowbreak 1,\allowbreak 7,\allowbreak 5,\allowbreak 6,\allowbreak 3)}\tabularnewline[5pt]
{\raggedright\arraybackslash $e=2$, $n=12$\par alpha path} & {\ttfamily\footnotesize (2,\allowbreak 11,\allowbreak 0,\allowbreak 10,\allowbreak 3,\allowbreak 9,\allowbreak 1,\allowbreak 6,\allowbreak 5,\allowbreak 7,\allowbreak 4,\allowbreak 8)}\tabularnewline[5pt]
{\raggedright\arraybackslash $e=2$, $n=17$\par alpha path} & {\ttfamily\footnotesize (2,\allowbreak 16,\allowbreak 0,\allowbreak 15,\allowbreak 3,\allowbreak 14,\allowbreak 1,\allowbreak 11,\allowbreak 4,\allowbreak 13,\allowbreak 5,\allowbreak 10,\allowbreak 7,\allowbreak 9,\allowbreak 8,\allowbreak 12,\allowbreak 6)}\tabularnewline[5pt]
{\raggedright\arraybackslash $e=3$, prefix\par $n=12$} & {\ttfamily\footnotesize (3,\allowbreak 11,\allowbreak 0,\allowbreak 10,\allowbreak 1,\allowbreak 8,\allowbreak 2,\allowbreak 7,\allowbreak 4,\allowbreak 6,\allowbreak 5,\allowbreak 9)}\tabularnewline[5pt]
{\raggedright\arraybackslash $e=3$, $n=7$\par alpha path} & {\ttfamily\footnotesize (3,\allowbreak 4,\allowbreak 2,\allowbreak 5,\allowbreak 1,\allowbreak 6,\allowbreak 0)}\tabularnewline[5pt]
{\raggedright\arraybackslash $e=3$, $n=8$\par alpha path} & {\ttfamily\footnotesize (3,\allowbreak 4,\allowbreak 2,\allowbreak 5,\allowbreak 1,\allowbreak 6,\allowbreak 0,\allowbreak 7)}\tabularnewline[5pt]
{\raggedright\arraybackslash $e=3$, $n=9$\par alpha path} & {\ttfamily\footnotesize (3,\allowbreak 7,\allowbreak 0,\allowbreak 8,\allowbreak 2,\allowbreak 5,\allowbreak 4,\allowbreak 6,\allowbreak 1)}\tabularnewline[5pt]
{\raggedright\arraybackslash $e=3$, $n=10$\par alpha path} & {\ttfamily\footnotesize (3,\allowbreak 7,\allowbreak 0,\allowbreak 9,\allowbreak 1,\allowbreak 6,\allowbreak 4,\allowbreak 5,\allowbreak 2,\allowbreak 8)}\tabularnewline[5pt]
{\raggedright\arraybackslash $e=3$, $n=11$\par alpha path} & {\ttfamily\footnotesize (3,\allowbreak 10,\allowbreak 0,\allowbreak 9,\allowbreak 1,\allowbreak 6,\allowbreak 5,\allowbreak 7,\allowbreak 4,\allowbreak 8,\allowbreak 2)}\tabularnewline[5pt]
{\raggedright\arraybackslash $e=3$, $n=12$\par alpha path} & {\ttfamily\footnotesize (3,\allowbreak 11,\allowbreak 0,\allowbreak 10,\allowbreak 1,\allowbreak 8,\allowbreak 2,\allowbreak 7,\allowbreak 4,\allowbreak 6,\allowbreak 5,\allowbreak 9)}\tabularnewline[5pt]
{\raggedright\arraybackslash $e=3$, $n=14$\par alpha path} & {\ttfamily\footnotesize (3,\allowbreak 13,\allowbreak 0,\allowbreak 12,\allowbreak 1,\allowbreak 9,\allowbreak 2,\allowbreak 11,\allowbreak 6,\allowbreak 7,\allowbreak 5,\allowbreak 8,\allowbreak 4,\allowbreak 10)}\tabularnewline[5pt]
{\raggedright\arraybackslash $e=3$, $n=15$\par alpha path} & {\ttfamily\footnotesize (3,\allowbreak 14,\allowbreak 0,\allowbreak 13,\allowbreak 1,\allowbreak 11,\allowbreak 2,\allowbreak 10,\allowbreak 5,\allowbreak 12,\allowbreak 6,\allowbreak 9,\allowbreak 7,\allowbreak 8,\allowbreak 4)}\tabularnewline[5pt]
{\raggedright\arraybackslash $e=3$, $n=16$\par alpha path} & {\ttfamily\footnotesize (3,\allowbreak 15,\allowbreak 0,\allowbreak 14,\allowbreak 1,\allowbreak 12,\allowbreak 2,\allowbreak 10,\allowbreak 5,\allowbreak 9,\allowbreak 6,\allowbreak 8,\allowbreak 7,\allowbreak 13,\allowbreak 4,\allowbreak 11)}\tabularnewline[5pt]
{\raggedright\arraybackslash $e=3$, $n=17$\par alpha path} & {\ttfamily\footnotesize (3,\allowbreak 16,\allowbreak 0,\allowbreak 15,\allowbreak 1,\allowbreak 13,\allowbreak 2,\allowbreak 12,\allowbreak 4,\allowbreak 11,\allowbreak 6,\allowbreak 10,\allowbreak 7,\allowbreak 9,\allowbreak 8,\allowbreak 14,\allowbreak 5)}\tabularnewline[5pt]
{\raggedright\arraybackslash $e=3$, $n=18$\par alpha path} & {\ttfamily\footnotesize (3,\allowbreak 17,\allowbreak 0,\allowbreak 16,\allowbreak 1,\allowbreak 14,\allowbreak 2,\allowbreak 13,\allowbreak 4,\allowbreak 11,\allowbreak 7,\allowbreak 15,\allowbreak 5,\allowbreak 10,\allowbreak 8,\allowbreak 9,\allowbreak 6,\allowbreak 12)}\tabularnewline[5pt]
{\raggedright\arraybackslash $e=3$, $n=25$\par alpha path} & {\ttfamily\footnotesize (3,\allowbreak 24,\allowbreak 0,\allowbreak 23,\allowbreak 1,\allowbreak 21,\allowbreak 2,\allowbreak 20,\allowbreak 4,\allowbreak 19,\allowbreak 6,\allowbreak 18,\allowbreak 7,\allowbreak 17,\allowbreak 8,\allowbreak 22,\allowbreak 5,\allowbreak 13,\allowbreak 12,\allowbreak 14,\allowbreak 11,\allowbreak 15,\allowbreak 10,\allowbreak 16,\allowbreak 9)}\tabularnewline[5pt]
\end{longtable}
\endgroup

\subsection{Pairs with minimum four: offset two}
\label{aux:pair4-2}
\begingroup\footnotesize
\begin{longtable}{@{}p{25mm}p{115mm}@{}}
Profile / interface & Complete outward words\\\hline\endhead
{\raggedright\arraybackslash {\ttfamily\footnotesize (4,\allowbreak 4)}\par offset $k=2$} & {\ttfamily\footnotesize (8,\allowbreak 3,\allowbreak 10,\allowbreak 7)\par\smallskip (6,\allowbreak 4,\allowbreak 5,\allowbreak 9)}\tabularnewline[5pt]
{\raggedright\arraybackslash {\ttfamily\footnotesize (4,\allowbreak 5)}\par offset $k=2$} & {\ttfamily\footnotesize (9,\allowbreak 3,\allowbreak 11,\allowbreak 4)\par\smallskip (5,\allowbreak 8,\allowbreak 10,\allowbreak 6,\allowbreak 7)}\tabularnewline[5pt]
{\raggedright\arraybackslash {\ttfamily\footnotesize (4,\allowbreak 6)}\par offset $k=2$} & {\ttfamily\footnotesize (10,\allowbreak 3,\allowbreak 12,\allowbreak 4)\par\smallskip (6,\allowbreak 11,\allowbreak 7,\allowbreak 9,\allowbreak 8,\allowbreak 5)}\tabularnewline[5pt]
{\raggedright\arraybackslash {\ttfamily\footnotesize (4,\allowbreak 7)}\par offset $k=2$} & {\ttfamily\footnotesize (11,\allowbreak 3,\allowbreak 13,\allowbreak 4)\par\smallskip (7,\allowbreak 12,\allowbreak 6,\allowbreak 8,\allowbreak 5,\allowbreak 9,\allowbreak 10)}\tabularnewline[5pt]
{\raggedright\arraybackslash {\ttfamily\footnotesize (4,\allowbreak 8)}\par offset $k=2$} & {\ttfamily\footnotesize (12,\allowbreak 3,\allowbreak 14,\allowbreak 4)\par\smallskip (8,\allowbreak 11,\allowbreak 6,\allowbreak 13,\allowbreak 7,\allowbreak 5,\allowbreak 9,\allowbreak 10)}\tabularnewline[5pt]
{\raggedright\arraybackslash {\ttfamily\footnotesize (4,\allowbreak 9)}\par offset $k=2$} & {\ttfamily\footnotesize (13,\allowbreak 3,\allowbreak 15,\allowbreak 4)\par\smallskip (9,\allowbreak 14,\allowbreak 6,\allowbreak 12,\allowbreak 5,\allowbreak 8,\allowbreak 10,\allowbreak 11,\allowbreak 7)}\tabularnewline[5pt]
{\raggedright\arraybackslash {\ttfamily\footnotesize (4,\allowbreak 10)}\par offset $k=2$} & {\ttfamily\footnotesize (14,\allowbreak 3,\allowbreak 16,\allowbreak 4)\par\smallskip (10,\allowbreak 15,\allowbreak 6,\allowbreak 13,\allowbreak 5,\allowbreak 11,\allowbreak 9,\allowbreak 12,\allowbreak 8,\allowbreak 7)}\tabularnewline[5pt]
{\raggedright\arraybackslash {\ttfamily\footnotesize (4,\allowbreak 11)}\par offset $k=2$} & {\ttfamily\footnotesize (15,\allowbreak 3,\allowbreak 17,\allowbreak 4)\par\smallskip (11,\allowbreak 5,\allowbreak 14,\allowbreak 6,\allowbreak 16,\allowbreak 9,\allowbreak 13,\allowbreak 8,\allowbreak 7,\allowbreak 10,\allowbreak 12)}\tabularnewline[5pt]
{\raggedright\arraybackslash {\ttfamily\footnotesize (4,\allowbreak 12)}\par offset $k=2$} & {\ttfamily\footnotesize (16,\allowbreak 3,\allowbreak 18,\allowbreak 4)\par\smallskip (12,\allowbreak 5,\allowbreak 15,\allowbreak 6,\allowbreak 17,\allowbreak 9,\allowbreak 14,\allowbreak 8,\allowbreak 7,\allowbreak 11,\allowbreak 13,\allowbreak 10)}\tabularnewline[5pt]
{\raggedright\arraybackslash {\ttfamily\footnotesize (4,\allowbreak 13)}\par offset $k=2$} & {\ttfamily\footnotesize (17,\allowbreak 3,\allowbreak 19,\allowbreak 4)\par\smallskip (13,\allowbreak 5,\allowbreak 16,\allowbreak 6,\allowbreak 18,\allowbreak 9,\allowbreak 15,\allowbreak 8,\allowbreak 11,\allowbreak 10,\allowbreak 14,\allowbreak 12,\allowbreak 7)}\tabularnewline[5pt]
{\raggedright\arraybackslash {\ttfamily\footnotesize (4,\allowbreak 14)}\par offset $k=2$} & {\ttfamily\footnotesize (18,\allowbreak 3,\allowbreak 20,\allowbreak 4)\par\smallskip (14,\allowbreak 5,\allowbreak 17,\allowbreak 6,\allowbreak 19,\allowbreak 9,\allowbreak 16,\allowbreak 8,\allowbreak 13,\allowbreak 7,\allowbreak 11,\allowbreak 10,\allowbreak 12,\allowbreak 15)}\tabularnewline[5pt]
{\raggedright\arraybackslash {\ttfamily\footnotesize (4,\allowbreak 15)}\par offset $k=2$} & {\ttfamily\footnotesize (19,\allowbreak 3,\allowbreak 21,\allowbreak 4)\par\smallskip (15,\allowbreak 5,\allowbreak 18,\allowbreak 6,\allowbreak 20,\allowbreak 9,\allowbreak 17,\allowbreak 8,\allowbreak 14,\allowbreak 7,\allowbreak 12,\allowbreak 16,\allowbreak 13,\allowbreak 11,\allowbreak 10)}\tabularnewline[5pt]
{\raggedright\arraybackslash {\ttfamily\footnotesize (4,\allowbreak 16)}\par offset $k=2$} & {\ttfamily\footnotesize (20,\allowbreak 3,\allowbreak 22,\allowbreak 4)\par\smallskip (16,\allowbreak 5,\allowbreak 19,\allowbreak 6,\allowbreak 21,\allowbreak 9,\allowbreak 18,\allowbreak 8,\allowbreak 15,\allowbreak 7,\allowbreak 13,\allowbreak 17,\allowbreak 12,\allowbreak 14,\allowbreak 11,\allowbreak 10)}\tabularnewline[5pt]
{\raggedright\arraybackslash {\ttfamily\footnotesize (4,\allowbreak 17)}\par offset $k=2$} & {\ttfamily\footnotesize (21,\allowbreak 3,\allowbreak 23,\allowbreak 4)\par\smallskip (17,\allowbreak 5,\allowbreak 20,\allowbreak 6,\allowbreak 22,\allowbreak 9,\allowbreak 19,\allowbreak 8,\allowbreak 16,\allowbreak 7,\allowbreak 14,\allowbreak 18,\allowbreak 12,\allowbreak 15,\allowbreak 10,\allowbreak 11,\allowbreak 13)}\tabularnewline[5pt]
\end{longtable}
\endgroup

\subsection{Pairs with minimum four: offset three}
\label{aux:pair4-3}
\begingroup\footnotesize
\begin{longtable}{@{}p{25mm}p{115mm}@{}}
Profile / interface & Complete outward words\\\hline\endhead
{\raggedright\arraybackslash {\ttfamily\footnotesize (4,\allowbreak 4)}\par offset $k=3$} & {\ttfamily\footnotesize (8,\allowbreak 4,\allowbreak 11,\allowbreak 9)\par\smallskip (6,\allowbreak 7,\allowbreak 10,\allowbreak 5)}\tabularnewline[5pt]
{\raggedright\arraybackslash {\ttfamily\footnotesize (4,\allowbreak 5)}\par offset $k=3$} & {\ttfamily\footnotesize (9,\allowbreak 4,\allowbreak 12,\allowbreak 5)\par\smallskip (6,\allowbreak 7,\allowbreak 11,\allowbreak 8,\allowbreak 10)}\tabularnewline[5pt]
{\raggedright\arraybackslash {\ttfamily\footnotesize (4,\allowbreak 6)}\par offset $k=3$} & {\ttfamily\footnotesize (10,\allowbreak 4,\allowbreak 13,\allowbreak 5)\par\smallskip (7,\allowbreak 11,\allowbreak 6,\allowbreak 8,\allowbreak 9,\allowbreak 12)}\tabularnewline[5pt]
{\raggedright\arraybackslash {\ttfamily\footnotesize (4,\allowbreak 7)}\par offset $k=3$} & {\ttfamily\footnotesize (11,\allowbreak 4,\allowbreak 14,\allowbreak 5)\par\smallskip (8,\allowbreak 13,\allowbreak 7,\allowbreak 6,\allowbreak 10,\allowbreak 12,\allowbreak 9)}\tabularnewline[5pt]
{\raggedright\arraybackslash {\ttfamily\footnotesize (4,\allowbreak 8)}\par offset $k=3$} & {\ttfamily\footnotesize (12,\allowbreak 4,\allowbreak 15,\allowbreak 5)\par\smallskip (9,\allowbreak 14,\allowbreak 8,\allowbreak 10,\allowbreak 13,\allowbreak 6,\allowbreak 7,\allowbreak 11)}\tabularnewline[5pt]
{\raggedright\arraybackslash {\ttfamily\footnotesize (4,\allowbreak 9)}\par offset $k=3$} & {\ttfamily\footnotesize (13,\allowbreak 4,\allowbreak 16,\allowbreak 5)\par\smallskip (10,\allowbreak 15,\allowbreak 7,\allowbreak 14,\allowbreak 8,\allowbreak 12,\allowbreak 11,\allowbreak 9,\allowbreak 6)}\tabularnewline[5pt]
{\raggedright\arraybackslash {\ttfamily\footnotesize (4,\allowbreak 10)}\par offset $k=3$} & {\ttfamily\footnotesize (14,\allowbreak 4,\allowbreak 17,\allowbreak 5)\par\smallskip (11,\allowbreak 16,\allowbreak 7,\allowbreak 15,\allowbreak 8,\allowbreak 9,\allowbreak 13,\allowbreak 10,\allowbreak 12,\allowbreak 6)}\tabularnewline[5pt]
{\raggedright\arraybackslash {\ttfamily\footnotesize (4,\allowbreak 11)}\par offset $k=3$} & {\ttfamily\footnotesize (15,\allowbreak 4,\allowbreak 18,\allowbreak 5)\par\smallskip (12,\allowbreak 6,\allowbreak 16,\allowbreak 7,\allowbreak 14,\allowbreak 17,\allowbreak 9,\allowbreak 13,\allowbreak 8,\allowbreak 10,\allowbreak 11)}\tabularnewline[5pt]
{\raggedright\arraybackslash {\ttfamily\footnotesize (4,\allowbreak 12)}\par offset $k=3$} & {\ttfamily\footnotesize (16,\allowbreak 4,\allowbreak 19,\allowbreak 5)\par\smallskip (13,\allowbreak 6,\allowbreak 17,\allowbreak 7,\allowbreak 15,\allowbreak 10,\allowbreak 14,\allowbreak 12,\allowbreak 18,\allowbreak 9,\allowbreak 8,\allowbreak 11)}\tabularnewline[5pt]
{\raggedright\arraybackslash {\ttfamily\footnotesize (4,\allowbreak 13)}\par offset $k=3$} & {\ttfamily\footnotesize (17,\allowbreak 4,\allowbreak 20,\allowbreak 5)\par\smallskip (14,\allowbreak 6,\allowbreak 18,\allowbreak 7,\allowbreak 16,\allowbreak 9,\allowbreak 19,\allowbreak 13,\allowbreak 8,\allowbreak 10,\allowbreak 11,\allowbreak 15,\allowbreak 12)}\tabularnewline[5pt]
{\raggedright\arraybackslash {\ttfamily\footnotesize (4,\allowbreak 14)}\par offset $k=3$} & {\ttfamily\footnotesize (18,\allowbreak 4,\allowbreak 21,\allowbreak 5)\par\smallskip (15,\allowbreak 6,\allowbreak 19,\allowbreak 7,\allowbreak 17,\allowbreak 9,\allowbreak 20,\allowbreak 13,\allowbreak 8,\allowbreak 14,\allowbreak 11,\allowbreak 10,\allowbreak 12,\allowbreak 16)}\tabularnewline[5pt]
{\raggedright\arraybackslash {\ttfamily\footnotesize (4,\allowbreak 15)}\par offset $k=3$} & {\ttfamily\footnotesize (19,\allowbreak 4,\allowbreak 22,\allowbreak 5)\par\smallskip (16,\allowbreak 6,\allowbreak 20,\allowbreak 7,\allowbreak 18,\allowbreak 9,\allowbreak 21,\allowbreak 13,\allowbreak 8,\allowbreak 15,\allowbreak 12,\allowbreak 14,\allowbreak 10,\allowbreak 11,\allowbreak 17)}\tabularnewline[5pt]
{\raggedright\arraybackslash {\ttfamily\footnotesize (4,\allowbreak 16)}\par offset $k=3$} & {\ttfamily\footnotesize (20,\allowbreak 4,\allowbreak 23,\allowbreak 5)\par\smallskip (17,\allowbreak 6,\allowbreak 21,\allowbreak 7,\allowbreak 19,\allowbreak 9,\allowbreak 22,\allowbreak 13,\allowbreak 18,\allowbreak 10,\allowbreak 16,\allowbreak 12,\allowbreak 14,\allowbreak 15,\allowbreak 8,\allowbreak 11)}\tabularnewline[5pt]
{\raggedright\arraybackslash {\ttfamily\footnotesize (4,\allowbreak 17)}\par offset $k=3$} & {\ttfamily\footnotesize (21,\allowbreak 4,\allowbreak 24,\allowbreak 5)\par\smallskip (18,\allowbreak 6,\allowbreak 22,\allowbreak 7,\allowbreak 20,\allowbreak 9,\allowbreak 23,\allowbreak 13,\allowbreak 19,\allowbreak 10,\allowbreak 17,\allowbreak 12,\allowbreak 14,\allowbreak 11,\allowbreak 15,\allowbreak 16,\allowbreak 8)}\tabularnewline[5pt]
{\raggedright\arraybackslash {\ttfamily\footnotesize (4,\allowbreak 18)}\par offset $k=3$} & {\ttfamily\footnotesize (22,\allowbreak 4,\allowbreak 25,\allowbreak 5)\par\smallskip (19,\allowbreak 6,\allowbreak 23,\allowbreak 7,\allowbreak 21,\allowbreak 9,\allowbreak 24,\allowbreak 13,\allowbreak 20,\allowbreak 10,\allowbreak 18,\allowbreak 12,\allowbreak 17,\allowbreak 8,\allowbreak 11,\allowbreak 15,\allowbreak 16,\allowbreak 14)}\tabularnewline[5pt]
{\raggedright\arraybackslash {\ttfamily\footnotesize (4,\allowbreak 19)}\par offset $k=3$} & {\ttfamily\footnotesize (23,\allowbreak 4,\allowbreak 26,\allowbreak 5)\par\smallskip (20,\allowbreak 6,\allowbreak 24,\allowbreak 7,\allowbreak 22,\allowbreak 9,\allowbreak 25,\allowbreak 13,\allowbreak 21,\allowbreak 10,\allowbreak 19,\allowbreak 12,\allowbreak 17,\allowbreak 11,\allowbreak 8,\allowbreak 18,\allowbreak 14,\allowbreak 16,\allowbreak 15)}\tabularnewline[5pt]
{\raggedright\arraybackslash {\ttfamily\footnotesize (4,\allowbreak 20)}\par offset $k=3$} & {\ttfamily\footnotesize (24,\allowbreak 4,\allowbreak 27,\allowbreak 5)\par\smallskip (21,\allowbreak 6,\allowbreak 25,\allowbreak 7,\allowbreak 23,\allowbreak 9,\allowbreak 26,\allowbreak 13,\allowbreak 22,\allowbreak 10,\allowbreak 20,\allowbreak 12,\allowbreak 19,\allowbreak 8,\allowbreak 14,\allowbreak 18,\allowbreak 15,\allowbreak 17,\allowbreak 16,\allowbreak 11)}\tabularnewline[5pt]
{\raggedright\arraybackslash {\ttfamily\footnotesize (4,\allowbreak 21)}\par offset $k=3$} & {\ttfamily\footnotesize (25,\allowbreak 4,\allowbreak 28,\allowbreak 5)\par\smallskip (22,\allowbreak 6,\allowbreak 26,\allowbreak 7,\allowbreak 24,\allowbreak 9,\allowbreak 27,\allowbreak 13,\allowbreak 23,\allowbreak 10,\allowbreak 21,\allowbreak 12,\allowbreak 20,\allowbreak 8,\allowbreak 15,\allowbreak 19,\allowbreak 14,\allowbreak 11,\allowbreak 17,\allowbreak 18,\allowbreak 16)}\tabularnewline[5pt]
{\raggedright\arraybackslash {\ttfamily\footnotesize (4,\allowbreak 22)}\par offset $k=3$} & {\ttfamily\footnotesize (26,\allowbreak 4,\allowbreak 29,\allowbreak 5)\par\smallskip (23,\allowbreak 6,\allowbreak 27,\allowbreak 7,\allowbreak 25,\allowbreak 9,\allowbreak 28,\allowbreak 13,\allowbreak 24,\allowbreak 10,\allowbreak 22,\allowbreak 12,\allowbreak 21,\allowbreak 8,\allowbreak 16,\allowbreak 11,\allowbreak 18,\allowbreak 19,\allowbreak 15,\allowbreak 17,\allowbreak 20,\allowbreak 14)}\tabularnewline[5pt]
{\raggedright\arraybackslash {\ttfamily\footnotesize (4,\allowbreak 23)}\par offset $k=3$} & {\ttfamily\footnotesize (27,\allowbreak 4,\allowbreak 30,\allowbreak 5)\par\smallskip (24,\allowbreak 6,\allowbreak 28,\allowbreak 7,\allowbreak 26,\allowbreak 9,\allowbreak 29,\allowbreak 13,\allowbreak 25,\allowbreak 10,\allowbreak 23,\allowbreak 12,\allowbreak 22,\allowbreak 8,\allowbreak 17,\allowbreak 11,\allowbreak 19,\allowbreak 14,\allowbreak 21,\allowbreak 18,\allowbreak 20,\allowbreak 16,\allowbreak 15)}\tabularnewline[5pt]
\end{longtable}
\endgroup

\subsection{Pairs with minimum five: finite offsets}
\label{aux:pair5}
\begingroup\footnotesize
\begin{longtable}{@{}p{25mm}p{115mm}@{}}
Profile / interface & Complete outward words\\\hline\endhead
{\raggedright\arraybackslash {\ttfamily\footnotesize (5,\allowbreak 5)}\par offset $k=1$} & {\ttfamily\footnotesize (10,\allowbreak 2,\allowbreak 11,\allowbreak 4,\allowbreak 9)\par\smallskip (6,\allowbreak 7,\allowbreak 3,\allowbreak 5,\allowbreak 8)}\tabularnewline[5pt]
{\raggedright\arraybackslash {\ttfamily\footnotesize (5,\allowbreak 6)}\par offset $k=1$} & {\ttfamily\footnotesize (11,\allowbreak 2,\allowbreak 12,\allowbreak 4,\allowbreak 10)\par\smallskip (7,\allowbreak 9,\allowbreak 5,\allowbreak 6,\allowbreak 3,\allowbreak 8)}\tabularnewline[5pt]
{\raggedright\arraybackslash {\ttfamily\footnotesize (5,\allowbreak 7)}\par offset $k=1$} & {\ttfamily\footnotesize (12,\allowbreak 2,\allowbreak 13,\allowbreak 4,\allowbreak 11)\par\smallskip (8,\allowbreak 3,\allowbreak 9,\allowbreak 7,\allowbreak 10,\allowbreak 6,\allowbreak 5)}\tabularnewline[5pt]
{\raggedright\arraybackslash {\ttfamily\footnotesize (5,\allowbreak 8)}\par offset $k=1$} & {\ttfamily\footnotesize (13,\allowbreak 2,\allowbreak 14,\allowbreak 4,\allowbreak 12)\par\smallskip (9,\allowbreak 3,\allowbreak 10,\allowbreak 5,\allowbreak 8,\allowbreak 6,\allowbreak 7,\allowbreak 11)}\tabularnewline[5pt]
{\raggedright\arraybackslash {\ttfamily\footnotesize (5,\allowbreak 9)}\par offset $k=1$} & {\ttfamily\footnotesize (14,\allowbreak 2,\allowbreak 15,\allowbreak 4,\allowbreak 13)\par\smallskip (10,\allowbreak 3,\allowbreak 11,\allowbreak 5,\allowbreak 9,\allowbreak 12,\allowbreak 7,\allowbreak 8,\allowbreak 6)}\tabularnewline[5pt]
{\raggedright\arraybackslash {\ttfamily\footnotesize (5,\allowbreak 10)}\par offset $k=1$} & {\ttfamily\footnotesize (15,\allowbreak 2,\allowbreak 16,\allowbreak 4,\allowbreak 14)\par\smallskip (11,\allowbreak 3,\allowbreak 12,\allowbreak 5,\allowbreak 10,\allowbreak 7,\allowbreak 13,\allowbreak 9,\allowbreak 8,\allowbreak 6)}\tabularnewline[5pt]
{\raggedright\arraybackslash {\ttfamily\footnotesize (5,\allowbreak 11)}\par offset $k=1$} & {\ttfamily\footnotesize (16,\allowbreak 2,\allowbreak 17,\allowbreak 4,\allowbreak 15)\par\smallskip (12,\allowbreak 3,\allowbreak 13,\allowbreak 5,\allowbreak 11,\allowbreak 6,\allowbreak 9,\allowbreak 8,\allowbreak 10,\allowbreak 14,\allowbreak 7)}\tabularnewline[5pt]
{\raggedright\arraybackslash {\ttfamily\footnotesize (5,\allowbreak 12)}\par offset $k=1$} & {\ttfamily\footnotesize (17,\allowbreak 2,\allowbreak 18,\allowbreak 4,\allowbreak 16)\par\smallskip (13,\allowbreak 3,\allowbreak 14,\allowbreak 5,\allowbreak 12,\allowbreak 6,\allowbreak 11,\allowbreak 15,\allowbreak 7,\allowbreak 10,\allowbreak 8,\allowbreak 9)}\tabularnewline[5pt]
{\raggedright\arraybackslash {\ttfamily\footnotesize (5,\allowbreak 5)}\par offset $k=3$} & {\ttfamily\footnotesize (10,\allowbreak 4,\allowbreak 13,\allowbreak 5,\allowbreak 9)\par\smallskip (7,\allowbreak 12,\allowbreak 11,\allowbreak 8,\allowbreak 6)}\tabularnewline[5pt]
{\raggedright\arraybackslash {\ttfamily\footnotesize (5,\allowbreak 6)}\par offset $k=3$} & {\ttfamily\footnotesize (11,\allowbreak 4,\allowbreak 14,\allowbreak 5,\allowbreak 13)\par\smallskip (6,\allowbreak 10,\allowbreak 8,\allowbreak 9,\allowbreak 12,\allowbreak 7)}\tabularnewline[5pt]
{\raggedright\arraybackslash {\ttfamily\footnotesize (5,\allowbreak 7)}\par offset $k=3$} & {\ttfamily\footnotesize (12,\allowbreak 4,\allowbreak 15,\allowbreak 5,\allowbreak 14)\par\smallskip (7,\allowbreak 13,\allowbreak 8,\allowbreak 11,\allowbreak 9,\allowbreak 10,\allowbreak 6)}\tabularnewline[5pt]
{\raggedright\arraybackslash {\ttfamily\footnotesize (5,\allowbreak 8)}\par offset $k=3$} & {\ttfamily\footnotesize (13,\allowbreak 4,\allowbreak 16,\allowbreak 5,\allowbreak 15)\par\smallskip (8,\allowbreak 14,\allowbreak 7,\allowbreak 12,\allowbreak 9,\allowbreak 11,\allowbreak 10,\allowbreak 6)}\tabularnewline[5pt]
{\raggedright\arraybackslash {\ttfamily\footnotesize (5,\allowbreak 9)}\par offset $k=3$} & {\ttfamily\footnotesize (14,\allowbreak 4,\allowbreak 17,\allowbreak 5,\allowbreak 16)\par\smallskip (9,\allowbreak 15,\allowbreak 7,\allowbreak 12,\allowbreak 8,\allowbreak 6,\allowbreak 13,\allowbreak 10,\allowbreak 11)}\tabularnewline[5pt]
{\raggedright\arraybackslash {\ttfamily\footnotesize (5,\allowbreak 10)}\par offset $k=3$} & {\ttfamily\footnotesize (15,\allowbreak 4,\allowbreak 18,\allowbreak 5,\allowbreak 17)\par\smallskip (10,\allowbreak 16,\allowbreak 7,\allowbreak 14,\allowbreak 6,\allowbreak 9,\allowbreak 11,\allowbreak 12,\allowbreak 8,\allowbreak 13)}\tabularnewline[5pt]
{\raggedright\arraybackslash {\ttfamily\footnotesize (5,\allowbreak 11)}\par offset $k=3$} & {\ttfamily\footnotesize (16,\allowbreak 4,\allowbreak 19,\allowbreak 5,\allowbreak 18)\par\smallskip (11,\allowbreak 17,\allowbreak 7,\allowbreak 15,\allowbreak 6,\allowbreak 13,\allowbreak 8,\allowbreak 9,\allowbreak 12,\allowbreak 14,\allowbreak 10)}\tabularnewline[5pt]
{\raggedright\arraybackslash {\ttfamily\footnotesize (5,\allowbreak 12)}\par offset $k=3$} & {\ttfamily\footnotesize (17,\allowbreak 4,\allowbreak 20,\allowbreak 5,\allowbreak 19)\par\smallskip (12,\allowbreak 18,\allowbreak 7,\allowbreak 16,\allowbreak 6,\allowbreak 14,\allowbreak 9,\allowbreak 13,\allowbreak 15,\allowbreak 8,\allowbreak 11,\allowbreak 10)}\tabularnewline[5pt]
{\raggedright\arraybackslash {\ttfamily\footnotesize (5,\allowbreak 13)}\par offset $k=3$} & {\ttfamily\footnotesize (18,\allowbreak 4,\allowbreak 21,\allowbreak 5,\allowbreak 20)\par\smallskip (13,\allowbreak 6,\allowbreak 17,\allowbreak 7,\allowbreak 19,\allowbreak 10,\allowbreak 16,\allowbreak 8,\allowbreak 12,\allowbreak 15,\allowbreak 14,\allowbreak 9,\allowbreak 11)}\tabularnewline[5pt]
{\raggedright\arraybackslash {\ttfamily\footnotesize (5,\allowbreak 14)}\par offset $k=3$} & {\ttfamily\footnotesize (19,\allowbreak 4,\allowbreak 22,\allowbreak 5,\allowbreak 21)\par\smallskip (14,\allowbreak 6,\allowbreak 18,\allowbreak 7,\allowbreak 20,\allowbreak 10,\allowbreak 17,\allowbreak 8,\allowbreak 13,\allowbreak 16,\allowbreak 12,\allowbreak 11,\allowbreak 9,\allowbreak 15)}\tabularnewline[5pt]
{\raggedright\arraybackslash {\ttfamily\footnotesize (5,\allowbreak 15)}\par offset $k=3$} & {\ttfamily\footnotesize (20,\allowbreak 4,\allowbreak 23,\allowbreak 5,\allowbreak 22)\par\smallskip (15,\allowbreak 6,\allowbreak 19,\allowbreak 7,\allowbreak 21,\allowbreak 10,\allowbreak 18,\allowbreak 8,\allowbreak 14,\allowbreak 17,\allowbreak 12,\allowbreak 11,\allowbreak 13,\allowbreak 9,\allowbreak 16)}\tabularnewline[5pt]
{\raggedright\arraybackslash {\ttfamily\footnotesize (5,\allowbreak 16)}\par offset $k=3$} & {\ttfamily\footnotesize (21,\allowbreak 4,\allowbreak 24,\allowbreak 5,\allowbreak 23)\par\smallskip (16,\allowbreak 6,\allowbreak 20,\allowbreak 7,\allowbreak 22,\allowbreak 10,\allowbreak 19,\allowbreak 8,\allowbreak 15,\allowbreak 9,\allowbreak 17,\allowbreak 12,\allowbreak 13,\allowbreak 11,\allowbreak 14,\allowbreak 18)}\tabularnewline[5pt]
{\raggedright\arraybackslash {\ttfamily\footnotesize (5,\allowbreak 17)}\par offset $k=3$} & {\ttfamily\footnotesize (22,\allowbreak 4,\allowbreak 25,\allowbreak 5,\allowbreak 24)\par\smallskip (17,\allowbreak 6,\allowbreak 21,\allowbreak 7,\allowbreak 23,\allowbreak 10,\allowbreak 20,\allowbreak 8,\allowbreak 16,\allowbreak 9,\allowbreak 18,\allowbreak 12,\allowbreak 15,\allowbreak 19,\allowbreak 14,\allowbreak 13,\allowbreak 11)}\tabularnewline[5pt]
{\raggedright\arraybackslash {\ttfamily\footnotesize (5,\allowbreak 18)}\par offset $k=3$} & {\ttfamily\footnotesize (23,\allowbreak 4,\allowbreak 26,\allowbreak 5,\allowbreak 25)\par\smallskip (18,\allowbreak 6,\allowbreak 22,\allowbreak 7,\allowbreak 24,\allowbreak 10,\allowbreak 21,\allowbreak 8,\allowbreak 17,\allowbreak 9,\allowbreak 19,\allowbreak 12,\allowbreak 16,\allowbreak 15,\allowbreak 20,\allowbreak 14,\allowbreak 11,\allowbreak 13)}\tabularnewline[5pt]
{\raggedright\arraybackslash {\ttfamily\footnotesize (5,\allowbreak 19)}\par offset $k=3$} & {\ttfamily\footnotesize (24,\allowbreak 4,\allowbreak 27,\allowbreak 5,\allowbreak 26)\par\smallskip (19,\allowbreak 6,\allowbreak 23,\allowbreak 7,\allowbreak 25,\allowbreak 10,\allowbreak 22,\allowbreak 8,\allowbreak 18,\allowbreak 9,\allowbreak 20,\allowbreak 12,\allowbreak 17,\allowbreak 13,\allowbreak 11,\allowbreak 14,\allowbreak 21,\allowbreak 15,\allowbreak 16)}\tabularnewline[5pt]
{\raggedright\arraybackslash {\ttfamily\footnotesize (5,\allowbreak 20)}\par offset $k=3$} & {\ttfamily\footnotesize (25,\allowbreak 4,\allowbreak 28,\allowbreak 5,\allowbreak 27)\par\smallskip (20,\allowbreak 6,\allowbreak 24,\allowbreak 7,\allowbreak 26,\allowbreak 10,\allowbreak 23,\allowbreak 8,\allowbreak 19,\allowbreak 9,\allowbreak 21,\allowbreak 12,\allowbreak 18,\allowbreak 11,\allowbreak 16,\allowbreak 13,\allowbreak 17,\allowbreak 15,\allowbreak 14,\allowbreak 22)}\tabularnewline[5pt]
{\raggedright\arraybackslash {\ttfamily\footnotesize (5,\allowbreak 5)}\par offset $k=4$} & {\ttfamily\footnotesize (10,\allowbreak 5,\allowbreak 14,\allowbreak 6,\allowbreak 12)\par\smallskip (7,\allowbreak 8,\allowbreak 11,\allowbreak 13,\allowbreak 9)}\tabularnewline[5pt]
{\raggedright\arraybackslash {\ttfamily\footnotesize (5,\allowbreak 6)}\par offset $k=4$} & {\ttfamily\footnotesize (11,\allowbreak 5,\allowbreak 15,\allowbreak 6,\allowbreak 14)\par\smallskip (7,\allowbreak 12,\allowbreak 10,\allowbreak 13,\allowbreak 9,\allowbreak 8)}\tabularnewline[5pt]
{\raggedright\arraybackslash {\ttfamily\footnotesize (5,\allowbreak 7)}\par offset $k=4$} & {\ttfamily\footnotesize (12,\allowbreak 5,\allowbreak 16,\allowbreak 6,\allowbreak 15)\par\smallskip (8,\allowbreak 14,\allowbreak 9,\allowbreak 13,\allowbreak 11,\allowbreak 10,\allowbreak 7)}\tabularnewline[5pt]
{\raggedright\arraybackslash {\ttfamily\footnotesize (5,\allowbreak 8)}\par offset $k=4$} & {\ttfamily\footnotesize (13,\allowbreak 5,\allowbreak 17,\allowbreak 6,\allowbreak 16)\par\smallskip (9,\allowbreak 15,\allowbreak 8,\allowbreak 11,\allowbreak 10,\allowbreak 14,\allowbreak 12,\allowbreak 7)}\tabularnewline[5pt]
{\raggedright\arraybackslash {\ttfamily\footnotesize (5,\allowbreak 9)}\par offset $k=4$} & {\ttfamily\footnotesize (14,\allowbreak 5,\allowbreak 18,\allowbreak 6,\allowbreak 17)\par\smallskip (10,\allowbreak 16,\allowbreak 9,\allowbreak 13,\allowbreak 15,\allowbreak 7,\allowbreak 12,\allowbreak 11,\allowbreak 8)}\tabularnewline[5pt]
{\raggedright\arraybackslash {\ttfamily\footnotesize (5,\allowbreak 10)}\par offset $k=4$} & {\ttfamily\footnotesize (15,\allowbreak 5,\allowbreak 19,\allowbreak 6,\allowbreak 18)\par\smallskip (11,\allowbreak 17,\allowbreak 8,\allowbreak 16,\allowbreak 9,\allowbreak 13,\allowbreak 14,\allowbreak 12,\allowbreak 7,\allowbreak 10)}\tabularnewline[5pt]
{\raggedright\arraybackslash {\ttfamily\footnotesize (5,\allowbreak 11)}\par offset $k=4$} & {\ttfamily\footnotesize (16,\allowbreak 5,\allowbreak 20,\allowbreak 6,\allowbreak 19)\par\smallskip (12,\allowbreak 18,\allowbreak 8,\allowbreak 17,\allowbreak 9,\allowbreak 14,\allowbreak 7,\allowbreak 11,\allowbreak 10,\allowbreak 13,\allowbreak 15)}\tabularnewline[5pt]
{\raggedright\arraybackslash {\ttfamily\footnotesize (5,\allowbreak 12)}\par offset $k=4$} & {\ttfamily\footnotesize (17,\allowbreak 5,\allowbreak 21,\allowbreak 6,\allowbreak 20)\par\smallskip (13,\allowbreak 19,\allowbreak 8,\allowbreak 18,\allowbreak 9,\allowbreak 16,\allowbreak 11,\allowbreak 12,\allowbreak 14,\allowbreak 10,\allowbreak 7,\allowbreak 15)}\tabularnewline[5pt]
{\raggedright\arraybackslash {\ttfamily\footnotesize (5,\allowbreak 13)}\par offset $k=4$} & {\ttfamily\footnotesize (18,\allowbreak 5,\allowbreak 22,\allowbreak 6,\allowbreak 21)\par\smallskip (14,\allowbreak 7,\allowbreak 19,\allowbreak 8,\allowbreak 17,\allowbreak 9,\allowbreak 15,\allowbreak 20,\allowbreak 10,\allowbreak 13,\allowbreak 11,\allowbreak 12,\allowbreak 16)}\tabularnewline[5pt]
{\raggedright\arraybackslash {\ttfamily\footnotesize (5,\allowbreak 14)}\par offset $k=4$} & {\ttfamily\footnotesize (19,\allowbreak 5,\allowbreak 23,\allowbreak 6,\allowbreak 22)\par\smallskip (15,\allowbreak 7,\allowbreak 20,\allowbreak 8,\allowbreak 18,\allowbreak 9,\allowbreak 16,\allowbreak 10,\allowbreak 21,\allowbreak 17,\allowbreak 12,\allowbreak 13,\allowbreak 11,\allowbreak 14)}\tabularnewline[5pt]
{\raggedright\arraybackslash {\ttfamily\footnotesize (5,\allowbreak 15)}\par offset $k=4$} & {\ttfamily\footnotesize (20,\allowbreak 5,\allowbreak 24,\allowbreak 6,\allowbreak 23)\par\smallskip (16,\allowbreak 7,\allowbreak 21,\allowbreak 8,\allowbreak 19,\allowbreak 9,\allowbreak 17,\allowbreak 11,\allowbreak 18,\allowbreak 22,\allowbreak 10,\allowbreak 15,\allowbreak 12,\allowbreak 14,\allowbreak 13)}\tabularnewline[5pt]
{\raggedright\arraybackslash {\ttfamily\footnotesize (5,\allowbreak 16)}\par offset $k=4$} & {\ttfamily\footnotesize (21,\allowbreak 5,\allowbreak 25,\allowbreak 6,\allowbreak 24)\par\smallskip (17,\allowbreak 7,\allowbreak 22,\allowbreak 8,\allowbreak 20,\allowbreak 9,\allowbreak 18,\allowbreak 10,\allowbreak 23,\allowbreak 16,\allowbreak 11,\allowbreak 15,\allowbreak 12,\allowbreak 14,\allowbreak 13,\allowbreak 19)}\tabularnewline[5pt]
{\raggedright\arraybackslash {\ttfamily\footnotesize (5,\allowbreak 17)}\par offset $k=4$} & {\ttfamily\footnotesize (22,\allowbreak 5,\allowbreak 26,\allowbreak 6,\allowbreak 25)\par\smallskip (18,\allowbreak 7,\allowbreak 23,\allowbreak 8,\allowbreak 21,\allowbreak 9,\allowbreak 19,\allowbreak 10,\allowbreak 24,\allowbreak 16,\allowbreak 11,\allowbreak 15,\allowbreak 17,\allowbreak 14,\allowbreak 20,\allowbreak 13,\allowbreak 12)}\tabularnewline[5pt]
{\raggedright\arraybackslash {\ttfamily\footnotesize (5,\allowbreak 18)}\par offset $k=4$} & {\ttfamily\footnotesize (23,\allowbreak 5,\allowbreak 27,\allowbreak 6,\allowbreak 26)\par\smallskip (19,\allowbreak 7,\allowbreak 24,\allowbreak 8,\allowbreak 22,\allowbreak 9,\allowbreak 20,\allowbreak 10,\allowbreak 25,\allowbreak 16,\allowbreak 21,\allowbreak 13,\allowbreak 17,\allowbreak 11,\allowbreak 18,\allowbreak 15,\allowbreak 14,\allowbreak 12)}\tabularnewline[5pt]
{\raggedright\arraybackslash {\ttfamily\footnotesize (5,\allowbreak 19)}\par offset $k=4$} & {\ttfamily\footnotesize (24,\allowbreak 5,\allowbreak 28,\allowbreak 6,\allowbreak 27)\par\smallskip (20,\allowbreak 7,\allowbreak 25,\allowbreak 8,\allowbreak 23,\allowbreak 9,\allowbreak 21,\allowbreak 10,\allowbreak 26,\allowbreak 16,\allowbreak 22,\allowbreak 13,\allowbreak 18,\allowbreak 17,\allowbreak 14,\allowbreak 12,\allowbreak 19,\allowbreak 11,\allowbreak 15)}\tabularnewline[5pt]
{\raggedright\arraybackslash {\ttfamily\footnotesize (5,\allowbreak 20)}\par offset $k=4$} & {\ttfamily\footnotesize (25,\allowbreak 5,\allowbreak 29,\allowbreak 6,\allowbreak 28)\par\smallskip (21,\allowbreak 7,\allowbreak 26,\allowbreak 8,\allowbreak 24,\allowbreak 9,\allowbreak 22,\allowbreak 10,\allowbreak 27,\allowbreak 16,\allowbreak 23,\allowbreak 13,\allowbreak 19,\allowbreak 11,\allowbreak 20,\allowbreak 15,\allowbreak 12,\allowbreak 14,\allowbreak 18,\allowbreak 17)}\tabularnewline[5pt]
{\raggedright\arraybackslash {\ttfamily\footnotesize (5,\allowbreak 21)}\par offset $k=4$} & {\ttfamily\footnotesize (26,\allowbreak 5,\allowbreak 30,\allowbreak 6,\allowbreak 29)\par\smallskip (22,\allowbreak 7,\allowbreak 27,\allowbreak 8,\allowbreak 25,\allowbreak 9,\allowbreak 23,\allowbreak 10,\allowbreak 28,\allowbreak 16,\allowbreak 24,\allowbreak 13,\allowbreak 20,\allowbreak 11,\allowbreak 21,\allowbreak 15,\allowbreak 12,\allowbreak 17,\allowbreak 19,\allowbreak 18,\allowbreak 14)}\tabularnewline[5pt]
{\raggedright\arraybackslash {\ttfamily\footnotesize (5,\allowbreak 22)}\par offset $k=4$} & {\ttfamily\footnotesize (27,\allowbreak 5,\allowbreak 31,\allowbreak 6,\allowbreak 30)\par\smallskip (23,\allowbreak 7,\allowbreak 28,\allowbreak 8,\allowbreak 26,\allowbreak 9,\allowbreak 24,\allowbreak 10,\allowbreak 29,\allowbreak 16,\allowbreak 25,\allowbreak 13,\allowbreak 21,\allowbreak 11,\allowbreak 22,\allowbreak 15,\allowbreak 20,\allowbreak 19,\allowbreak 17,\allowbreak 14,\allowbreak 18,\allowbreak 12)}\tabularnewline[5pt]
{\raggedright\arraybackslash {\ttfamily\footnotesize (5,\allowbreak 23)}\par offset $k=4$} & {\ttfamily\footnotesize (28,\allowbreak 5,\allowbreak 32,\allowbreak 6,\allowbreak 31)\par\smallskip (24,\allowbreak 7,\allowbreak 29,\allowbreak 8,\allowbreak 27,\allowbreak 9,\allowbreak 25,\allowbreak 10,\allowbreak 30,\allowbreak 16,\allowbreak 26,\allowbreak 13,\allowbreak 22,\allowbreak 11,\allowbreak 23,\allowbreak 15,\allowbreak 21,\allowbreak 20,\allowbreak 18,\allowbreak 14,\allowbreak 17,\allowbreak 12,\allowbreak 19)}\tabularnewline[5pt]
{\raggedright\arraybackslash {\ttfamily\footnotesize (5,\allowbreak 24)}\par offset $k=4$} & {\ttfamily\footnotesize (29,\allowbreak 5,\allowbreak 33,\allowbreak 6,\allowbreak 32)\par\smallskip (25,\allowbreak 7,\allowbreak 30,\allowbreak 8,\allowbreak 28,\allowbreak 9,\allowbreak 26,\allowbreak 10,\allowbreak 31,\allowbreak 16,\allowbreak 27,\allowbreak 13,\allowbreak 23,\allowbreak 11,\allowbreak 24,\allowbreak 15,\allowbreak 22,\allowbreak 14,\allowbreak 18,\allowbreak 12,\allowbreak 17,\allowbreak 20,\allowbreak 21,\allowbreak 19)}\tabularnewline[5pt]
{\raggedright\arraybackslash {\ttfamily\footnotesize (5,\allowbreak 25)}\par offset $k=4$} & {\ttfamily\footnotesize (30,\allowbreak 5,\allowbreak 34,\allowbreak 6,\allowbreak 33)\par\smallskip (26,\allowbreak 7,\allowbreak 31,\allowbreak 8,\allowbreak 29,\allowbreak 9,\allowbreak 27,\allowbreak 10,\allowbreak 32,\allowbreak 16,\allowbreak 28,\allowbreak 13,\allowbreak 24,\allowbreak 11,\allowbreak 25,\allowbreak 15,\allowbreak 23,\allowbreak 14,\allowbreak 21,\allowbreak 17,\allowbreak 22,\allowbreak 19,\allowbreak 20,\allowbreak 18,\allowbreak 12)}\tabularnewline[5pt]
{\raggedright\arraybackslash {\ttfamily\footnotesize (5,\allowbreak 26)}\par offset $k=4$} & {\ttfamily\footnotesize (31,\allowbreak 5,\allowbreak 35,\allowbreak 6,\allowbreak 34)\par\smallskip (27,\allowbreak 7,\allowbreak 32,\allowbreak 8,\allowbreak 30,\allowbreak 9,\allowbreak 28,\allowbreak 10,\allowbreak 33,\allowbreak 16,\allowbreak 29,\allowbreak 13,\allowbreak 25,\allowbreak 11,\allowbreak 26,\allowbreak 15,\allowbreak 24,\allowbreak 14,\allowbreak 22,\allowbreak 17,\allowbreak 21,\allowbreak 23,\allowbreak 20,\allowbreak 19,\allowbreak 12,\allowbreak 18)}\tabularnewline[5pt]
\end{longtable}
\endgroup

\subsection{The twelve offset-two pairs with a two-arm}
\label{aux:offset2-twos}
\begingroup\footnotesize
\begin{longtable}{@{}p{25mm}p{115mm}@{}}
Profile / interface & Complete outward words\\\hline\endhead
{\raggedright\arraybackslash {\ttfamily\footnotesize (2,\allowbreak 3)}\par offset $k=2$} & {\ttfamily\footnotesize (5,\allowbreak 7)\par\smallskip (4,\allowbreak 3,\allowbreak 6)}\tabularnewline[5pt]
{\raggedright\arraybackslash {\ttfamily\footnotesize (2,\allowbreak 4)}\par offset $k=2$} & {\ttfamily\footnotesize (6,\allowbreak 3)\par\smallskip (5,\allowbreak 7,\allowbreak 8,\allowbreak 4)}\tabularnewline[5pt]
{\raggedright\arraybackslash {\ttfamily\footnotesize (2,\allowbreak 5)}\par offset $k=2$} & {\ttfamily\footnotesize (7,\allowbreak 3)\par\smallskip (6,\allowbreak 8,\allowbreak 5,\allowbreak 4,\allowbreak 9)}\tabularnewline[5pt]
{\raggedright\arraybackslash {\ttfamily\footnotesize (2,\allowbreak 6)}\par offset $k=2$} & {\ttfamily\footnotesize (8,\allowbreak 3)\par\smallskip (7,\allowbreak 10,\allowbreak 4,\allowbreak 6,\allowbreak 5,\allowbreak 9)}\tabularnewline[5pt]
{\raggedright\arraybackslash {\ttfamily\footnotesize (2,\allowbreak 7)}\par offset $k=2$} & {\ttfamily\footnotesize (9,\allowbreak 3)\par\smallskip (8,\allowbreak 4,\allowbreak 11,\allowbreak 6,\allowbreak 5,\allowbreak 7,\allowbreak 10)}\tabularnewline[5pt]
{\raggedright\arraybackslash {\ttfamily\footnotesize (2,\allowbreak 8)}\par offset $k=2$} & {\ttfamily\footnotesize (10,\allowbreak 3)\par\smallskip (9,\allowbreak 4,\allowbreak 12,\allowbreak 6,\allowbreak 5,\allowbreak 7,\allowbreak 11,\allowbreak 8)}\tabularnewline[5pt]
{\raggedright\arraybackslash {\ttfamily\footnotesize (2,\allowbreak 9)}\par offset $k=2$} & {\ttfamily\footnotesize (11,\allowbreak 3)\par\smallskip (10,\allowbreak 4,\allowbreak 13,\allowbreak 6,\allowbreak 9,\allowbreak 8,\allowbreak 12,\allowbreak 7,\allowbreak 5)}\tabularnewline[5pt]
{\raggedright\arraybackslash {\ttfamily\footnotesize (2,\allowbreak 10)}\par offset $k=2$} & {\ttfamily\footnotesize (12,\allowbreak 3)\par\smallskip (11,\allowbreak 4,\allowbreak 14,\allowbreak 6,\allowbreak 10,\allowbreak 5,\allowbreak 8,\allowbreak 9,\allowbreak 7,\allowbreak 13)}\tabularnewline[5pt]
{\raggedright\arraybackslash {\ttfamily\footnotesize (2,\allowbreak 11)}\par offset $k=2$} & {\ttfamily\footnotesize (13,\allowbreak 3)\par\smallskip (12,\allowbreak 4,\allowbreak 15,\allowbreak 6,\allowbreak 11,\allowbreak 5,\allowbreak 9,\allowbreak 8,\allowbreak 10,\allowbreak 7,\allowbreak 14)}\tabularnewline[5pt]
{\raggedright\arraybackslash {\ttfamily\footnotesize (2,\allowbreak 12)}\par offset $k=2$} & {\ttfamily\footnotesize (14,\allowbreak 3)\par\smallskip (13,\allowbreak 4,\allowbreak 16,\allowbreak 6,\allowbreak 12,\allowbreak 5,\allowbreak 10,\allowbreak 7,\allowbreak 15,\allowbreak 11,\allowbreak 9,\allowbreak 8)}\tabularnewline[5pt]
{\raggedright\arraybackslash {\ttfamily\footnotesize (2,\allowbreak 13)}\par offset $k=2$} & {\ttfamily\footnotesize (15,\allowbreak 3)\par\smallskip (14,\allowbreak 4,\allowbreak 17,\allowbreak 6,\allowbreak 13,\allowbreak 5,\allowbreak 11,\allowbreak 16,\allowbreak 7,\allowbreak 10,\allowbreak 12,\allowbreak 8,\allowbreak 9)}\tabularnewline[5pt]
{\raggedright\arraybackslash {\ttfamily\footnotesize (2,\allowbreak 14)}\par offset $k=2$} & {\ttfamily\footnotesize (16,\allowbreak 3)\par\smallskip (15,\allowbreak 4,\allowbreak 18,\allowbreak 6,\allowbreak 14,\allowbreak 5,\allowbreak 12,\allowbreak 17,\allowbreak 7,\allowbreak 13,\allowbreak 9,\allowbreak 10,\allowbreak 8,\allowbreak 11)}\tabularnewline[5pt]
\end{longtable}
\endgroup

\subsection{Small equal-even alpha sources}
\label{aux:equal-even}
\begingroup\footnotesize
\begin{longtable}{@{}p{25mm}p{115mm}@{}}
Profile / interface & Complete outward words\\\hline\endhead
{\raggedright\arraybackslash $L=6$; $L^3$\par $K=9$\par complementary arm ends} & {\ttfamily\footnotesize (18,\allowbreak 1,\allowbreak 17,\allowbreak 5,\allowbreak 16,\allowbreak 2)\par\smallskip (15,\allowbreak 6,\allowbreak 12,\allowbreak 7,\allowbreak 14,\allowbreak 4)\par\smallskip (13,\allowbreak 9,\allowbreak 10,\allowbreak 8,\allowbreak 11,\allowbreak 3)}\tabularnewline[5pt]
{\raggedright\arraybackslash $L=8$; $L^3$\par $K=12$} & {\ttfamily\footnotesize (24,\allowbreak 1,\allowbreak 23,\allowbreak 2,\allowbreak 22,\allowbreak 6,\allowbreak 21,\allowbreak 3)\par\smallskip (19,\allowbreak 7,\allowbreak 20,\allowbreak 9,\allowbreak 16,\allowbreak 8,\allowbreak 18,\allowbreak 4)\par\smallskip (17,\allowbreak 11,\allowbreak 15,\allowbreak 10,\allowbreak 13,\allowbreak 12,\allowbreak 14,\allowbreak 5)}\tabularnewline[5pt]
{\raggedright\arraybackslash $L=10$; $L^3$\par $K=15$} & {\ttfamily\footnotesize (30,\allowbreak 1,\allowbreak 29,\allowbreak 2,\allowbreak 28,\allowbreak 3,\allowbreak 27,\allowbreak 7,\allowbreak 26,\allowbreak 4)\par\smallskip (23,\allowbreak 8,\allowbreak 25,\allowbreak 9,\allowbreak 22,\allowbreak 10,\allowbreak 20,\allowbreak 11,\allowbreak 19,\allowbreak 5)\par\smallskip (21,\allowbreak 14,\allowbreak 18,\allowbreak 12,\allowbreak 17,\allowbreak 15,\allowbreak 16,\allowbreak 13,\allowbreak 24,\allowbreak 6)}\tabularnewline[5pt]
{\raggedright\arraybackslash $L=14$; $L^3$\par $K=21$\par complementary arm ends} & {\ttfamily\footnotesize (42,\allowbreak 1,\allowbreak 41,\allowbreak 2,\allowbreak 40,\allowbreak 3,\allowbreak 33,\allowbreak 11,\allowbreak 39,\allowbreak 4,\allowbreak 38,\allowbreak 5,\allowbreak 37,\allowbreak 6)\par\smallskip (36,\allowbreak 15,\allowbreak 25,\allowbreak 16,\allowbreak 28,\allowbreak 17,\allowbreak 24,\allowbreak 18,\allowbreak 32,\allowbreak 12,\allowbreak 31,\allowbreak 13,\allowbreak 30,\allowbreak 7)\par\smallskip (29,\allowbreak 14,\allowbreak 27,\allowbreak 19,\allowbreak 23,\allowbreak 20,\allowbreak 22,\allowbreak 21,\allowbreak 26,\allowbreak 10,\allowbreak 34,\allowbreak 9,\allowbreak 35,\allowbreak 8)}\tabularnewline[5pt]
{\raggedright\arraybackslash $L=18$; $L^3$\par $K=27$} & {\ttfamily\footnotesize (54,\allowbreak 1,\allowbreak 53,\allowbreak 2,\allowbreak 52,\allowbreak 3,\allowbreak 51,\allowbreak 4,\allowbreak 42,\allowbreak 22,\allowbreak 31,\allowbreak 21,\allowbreak 36,\allowbreak 20,\allowbreak 33,\allowbreak 19,\allowbreak 38,\allowbreak 9)\par\smallskip (46,\allowbreak 18,\allowbreak 35,\allowbreak 17,\allowbreak 39,\allowbreak 16,\allowbreak 40,\allowbreak 15,\allowbreak 41,\allowbreak 14,\allowbreak 50,\allowbreak 5,\allowbreak 49,\allowbreak 6,\allowbreak 48,\allowbreak 7,\allowbreak 47,\allowbreak 8)\par\smallskip (37,\allowbreak 25,\allowbreak 29,\allowbreak 26,\allowbreak 28,\allowbreak 27,\allowbreak 32,\allowbreak 24,\allowbreak 30,\allowbreak 23,\allowbreak 34,\allowbreak 13,\allowbreak 43,\allowbreak 12,\allowbreak 44,\allowbreak 11,\allowbreak 45,\allowbreak 10)}\tabularnewline[5pt]
{\raggedright\arraybackslash $L=22$; $L^3$\par $K=33$\par complementary arm ends} & {\ttfamily\footnotesize (66,\allowbreak 1,\allowbreak 65,\allowbreak 2,\allowbreak 64,\allowbreak 3,\allowbreak 63,\allowbreak 4,\allowbreak 62,\allowbreak 5,\allowbreak 51,\allowbreak 17,\allowbreak 61,\allowbreak 6,\allowbreak 60,\allowbreak 7,\allowbreak 59,\allowbreak 8,\allowbreak 58,\allowbreak 9,\allowbreak 57,\allowbreak 10)\par\smallskip (56,\allowbreak 23,\allowbreak 41,\allowbreak 24,\allowbreak 44,\allowbreak 25,\allowbreak 39,\allowbreak 26,\allowbreak 42,\allowbreak 27,\allowbreak 37,\allowbreak 28,\allowbreak 50,\allowbreak 18,\allowbreak 49,\allowbreak 19,\allowbreak 48,\allowbreak 20,\allowbreak 47,\allowbreak 21,\allowbreak 46,\allowbreak 11)\par\smallskip (45,\allowbreak 22,\allowbreak 43,\allowbreak 31,\allowbreak 35,\allowbreak 32,\allowbreak 34,\allowbreak 33,\allowbreak 38,\allowbreak 30,\allowbreak 36,\allowbreak 29,\allowbreak 40,\allowbreak 16,\allowbreak 52,\allowbreak 15,\allowbreak 53,\allowbreak 14,\allowbreak 54,\allowbreak 13,\allowbreak 55,\allowbreak 12)}\tabularnewline[5pt]
\end{longtable}
\endgroup

\subsection{All finite four-equal-odd alpha bases}
\label{aux:equal-four-odd}
\begingroup\footnotesize
\begin{longtable}{@{}p{25mm}p{115mm}@{}}
Profile / interface & Complete outward words\\\hline\endhead
{\raggedright\arraybackslash $L=7$; $L^4$\par $K=12$} & {\ttfamily\footnotesize (20,\allowbreak 4,\allowbreak 17,\allowbreak 8,\allowbreak 19,\allowbreak 5,\allowbreak 24)\par\smallskip (28,\allowbreak 1,\allowbreak 27,\allowbreak 2,\allowbreak 26,\allowbreak 3,\allowbreak 25)\par\smallskip (21,\allowbreak 9,\allowbreak 14,\allowbreak 10,\allowbreak 16,\allowbreak 6,\allowbreak 23)\par\smallskip (18,\allowbreak 11,\allowbreak 13,\allowbreak 12,\allowbreak 15,\allowbreak 7,\allowbreak 22)}\tabularnewline[5pt]
{\raggedright\arraybackslash $L=9$; $L^4$\par $K=16$} & {\ttfamily\footnotesize (36,\allowbreak 1,\allowbreak 35,\allowbreak 2,\allowbreak 34,\allowbreak 3,\allowbreak 33,\allowbreak 4,\allowbreak 32)\par\smallskip (22,\allowbreak 14,\allowbreak 17,\allowbreak 16,\allowbreak 18,\allowbreak 13,\allowbreak 24,\allowbreak 9,\allowbreak 29)\par\smallskip (25,\allowbreak 12,\allowbreak 19,\allowbreak 15,\allowbreak 21,\allowbreak 7,\allowbreak 28,\allowbreak 5,\allowbreak 31)\par\smallskip (27,\allowbreak 8,\allowbreak 26,\allowbreak 10,\allowbreak 20,\allowbreak 11,\allowbreak 23,\allowbreak 6,\allowbreak 30)}\tabularnewline[5pt]
{\raggedright\arraybackslash $L=11$; $L^4$\par $K=20$} & {\ttfamily\footnotesize (30,\allowbreak 13,\allowbreak 23,\allowbreak 19,\allowbreak 25,\allowbreak 11,\allowbreak 35,\allowbreak 10,\allowbreak 32,\allowbreak 9,\allowbreak 37)\par\smallskip (44,\allowbreak 1,\allowbreak 43,\allowbreak 2,\allowbreak 42,\allowbreak 3,\allowbreak 41,\allowbreak 4,\allowbreak 40,\allowbreak 5,\allowbreak 39)\par\smallskip (31,\allowbreak 12,\allowbreak 28,\allowbreak 16,\allowbreak 27,\allowbreak 14,\allowbreak 29,\allowbreak 8,\allowbreak 34,\allowbreak 7,\allowbreak 36)\par\smallskip (33,\allowbreak 15,\allowbreak 24,\allowbreak 17,\allowbreak 22,\allowbreak 20,\allowbreak 21,\allowbreak 18,\allowbreak 26,\allowbreak 6,\allowbreak 38)}\tabularnewline[5pt]
{\raggedright\arraybackslash $L=13$; $L^4$\par $K=24$} & {\ttfamily\footnotesize (39,\allowbreak 10,\allowbreak 28,\allowbreak 21,\allowbreak 29,\allowbreak 19,\allowbreak 33,\allowbreak 14,\allowbreak 36,\allowbreak 8,\allowbreak 41,\allowbreak 11,\allowbreak 43)\par\smallskip (38,\allowbreak 12,\allowbreak 35,\allowbreak 15,\allowbreak 32,\allowbreak 20,\allowbreak 31,\allowbreak 16,\allowbreak 37,\allowbreak 13,\allowbreak 40,\allowbreak 9,\allowbreak 45)\par\smallskip (34,\allowbreak 18,\allowbreak 27,\allowbreak 22,\allowbreak 26,\allowbreak 23,\allowbreak 25,\allowbreak 24,\allowbreak 30,\allowbreak 17,\allowbreak 42,\allowbreak 7,\allowbreak 44)\par\smallskip (52,\allowbreak 1,\allowbreak 51,\allowbreak 2,\allowbreak 50,\allowbreak 3,\allowbreak 49,\allowbreak 4,\allowbreak 48,\allowbreak 5,\allowbreak 47,\allowbreak 6,\allowbreak 46)}\tabularnewline[5pt]
{\raggedright\arraybackslash $L=15$; $L^4$\par $K=28$} & {\ttfamily\footnotesize (42,\allowbreak 8,\allowbreak 49,\allowbreak 11,\allowbreak 46,\allowbreak 13,\allowbreak 37,\allowbreak 22,\allowbreak 35,\allowbreak 23,\allowbreak 40,\allowbreak 12,\allowbreak 48,\allowbreak 9,\allowbreak 52)\par\smallskip (45,\allowbreak 15,\allowbreak 47,\allowbreak 16,\allowbreak 39,\allowbreak 21,\allowbreak 31,\allowbreak 28,\allowbreak 29,\allowbreak 27,\allowbreak 32,\allowbreak 25,\allowbreak 36,\allowbreak 10,\allowbreak 50)\par\smallskip (60,\allowbreak 1,\allowbreak 59,\allowbreak 2,\allowbreak 58,\allowbreak 3,\allowbreak 57,\allowbreak 4,\allowbreak 56,\allowbreak 5,\allowbreak 55,\allowbreak 6,\allowbreak 54,\allowbreak 7,\allowbreak 53)\par\smallskip (44,\allowbreak 17,\allowbreak 33,\allowbreak 24,\allowbreak 30,\allowbreak 26,\allowbreak 34,\allowbreak 20,\allowbreak 41,\allowbreak 19,\allowbreak 38,\allowbreak 18,\allowbreak 43,\allowbreak 14,\allowbreak 51)}\tabularnewline[5pt]
{\raggedright\arraybackslash $L=17$; $L^4$\par $K=32$} & {\ttfamily\footnotesize (48,\allowbreak 20,\allowbreak 44,\allowbreak 21,\allowbreak 41,\allowbreak 28,\allowbreak 36,\allowbreak 29,\allowbreak 40,\allowbreak 26,\allowbreak 38,\allowbreak 23,\allowbreak 39,\allowbreak 17,\allowbreak 46,\allowbreak 11,\allowbreak 58)\par\smallskip (51,\allowbreak 15,\allowbreak 52,\allowbreak 13,\allowbreak 53,\allowbreak 19,\allowbreak 45,\allowbreak 18,\allowbreak 50,\allowbreak 12,\allowbreak 55,\allowbreak 14,\allowbreak 56,\allowbreak 10,\allowbreak 54,\allowbreak 9,\allowbreak 59)\par\smallskip (68,\allowbreak 1,\allowbreak 67,\allowbreak 2,\allowbreak 66,\allowbreak 3,\allowbreak 65,\allowbreak 4,\allowbreak 64,\allowbreak 5,\allowbreak 63,\allowbreak 6,\allowbreak 62,\allowbreak 7,\allowbreak 61,\allowbreak 8,\allowbreak 60)\par\smallskip (49,\allowbreak 16,\allowbreak 47,\allowbreak 22,\allowbreak 43,\allowbreak 24,\allowbreak 42,\allowbreak 25,\allowbreak 34,\allowbreak 30,\allowbreak 35,\allowbreak 32,\allowbreak 33,\allowbreak 31,\allowbreak 37,\allowbreak 27,\allowbreak 57)}\tabularnewline[5pt]
{\raggedright\arraybackslash $L=19$; $L^4$\par $K=36$} & {\ttfamily\footnotesize (54,\allowbreak 23,\allowbreak 59,\allowbreak 16,\allowbreak 62,\allowbreak 15,\allowbreak 56,\allowbreak 26,\allowbreak 50,\allowbreak 17,\allowbreak 52,\allowbreak 18,\allowbreak 58,\allowbreak 14,\allowbreak 63,\allowbreak 11,\allowbreak 61,\allowbreak 10,\allowbreak 66)\par\smallskip (57,\allowbreak 20,\allowbreak 43,\allowbreak 33,\allowbreak 41,\allowbreak 32,\allowbreak 39,\allowbreak 34,\allowbreak 40,\allowbreak 36,\allowbreak 37,\allowbreak 35,\allowbreak 38,\allowbreak 25,\allowbreak 53,\allowbreak 21,\allowbreak 60,\allowbreak 12,\allowbreak 65)\par\smallskip (55,\allowbreak 13,\allowbreak 51,\allowbreak 24,\allowbreak 45,\allowbreak 28,\allowbreak 44,\allowbreak 30,\allowbreak 42,\allowbreak 31,\allowbreak 46,\allowbreak 27,\allowbreak 49,\allowbreak 29,\allowbreak 47,\allowbreak 22,\allowbreak 48,\allowbreak 19,\allowbreak 64)\par\smallskip (76,\allowbreak 1,\allowbreak 75,\allowbreak 2,\allowbreak 74,\allowbreak 3,\allowbreak 73,\allowbreak 4,\allowbreak 72,\allowbreak 5,\allowbreak 71,\allowbreak 6,\allowbreak 70,\allowbreak 7,\allowbreak 69,\allowbreak 8,\allowbreak 68,\allowbreak 9,\allowbreak 67)}\tabularnewline[5pt]
{\raggedright\arraybackslash $L=21$; $L^4$\par $K=40$} & {\ttfamily\footnotesize (60,\allowbreak 21,\allowbreak 55,\allowbreak 27,\allowbreak 47,\allowbreak 33,\allowbreak 46,\allowbreak 37,\allowbreak 45,\allowbreak 35,\allowbreak 52,\allowbreak 26,\allowbreak 67,\allowbreak 16,\allowbreak 70,\allowbreak 12,\allowbreak 68,\allowbreak 15,\allowbreak 65,\allowbreak 13,\allowbreak 72)\par\smallskip (84,\allowbreak 1,\allowbreak 83,\allowbreak 2,\allowbreak 82,\allowbreak 3,\allowbreak 81,\allowbreak 4,\allowbreak 80,\allowbreak 5,\allowbreak 79,\allowbreak 6,\allowbreak 78,\allowbreak 7,\allowbreak 77,\allowbreak 8,\allowbreak 76,\allowbreak 9,\allowbreak 75,\allowbreak 10,\allowbreak 74)\par\smallskip (63,\allowbreak 25,\allowbreak 58,\allowbreak 23,\allowbreak 59,\allowbreak 29,\allowbreak 53,\allowbreak 30,\allowbreak 48,\allowbreak 36,\allowbreak 43,\allowbreak 39,\allowbreak 44,\allowbreak 38,\allowbreak 41,\allowbreak 40,\allowbreak 42,\allowbreak 31,\allowbreak 56,\allowbreak 11,\allowbreak 73)\par\smallskip (61,\allowbreak 17,\allowbreak 64,\allowbreak 18,\allowbreak 66,\allowbreak 24,\allowbreak 51,\allowbreak 32,\allowbreak 54,\allowbreak 22,\allowbreak 62,\allowbreak 19,\allowbreak 50,\allowbreak 34,\allowbreak 49,\allowbreak 28,\allowbreak 57,\allowbreak 20,\allowbreak 69,\allowbreak 14,\allowbreak 71)}\tabularnewline[5pt]
{\raggedright\arraybackslash $L=23$; $L^4$\par $K=44$} & {\ttfamily\footnotesize (69,\allowbreak 21,\allowbreak 66,\allowbreak 27,\allowbreak 54,\allowbreak 38,\allowbreak 51,\allowbreak 42,\allowbreak 49,\allowbreak 37,\allowbreak 52,\allowbreak 34,\allowbreak 58,\allowbreak 32,\allowbreak 65,\allowbreak 19,\allowbreak 70,\allowbreak 23,\allowbreak 63,\allowbreak 22,\allowbreak 71,\allowbreak 14,\allowbreak 79)\par\smallskip (68,\allowbreak 24,\allowbreak 74,\allowbreak 15,\allowbreak 57,\allowbreak 35,\allowbreak 55,\allowbreak 36,\allowbreak 53,\allowbreak 39,\allowbreak 50,\allowbreak 40,\allowbreak 48,\allowbreak 43,\allowbreak 46,\allowbreak 44,\allowbreak 45,\allowbreak 41,\allowbreak 47,\allowbreak 26,\allowbreak 64,\allowbreak 12,\allowbreak 78)\par\smallskip (92,\allowbreak 1,\allowbreak 91,\allowbreak 2,\allowbreak 90,\allowbreak 3,\allowbreak 89,\allowbreak 4,\allowbreak 88,\allowbreak 5,\allowbreak 87,\allowbreak 6,\allowbreak 86,\allowbreak 7,\allowbreak 85,\allowbreak 8,\allowbreak 84,\allowbreak 9,\allowbreak 83,\allowbreak 10,\allowbreak 82,\allowbreak 11,\allowbreak 81)\par\smallskip (67,\allowbreak 30,\allowbreak 61,\allowbreak 25,\allowbreak 60,\allowbreak 28,\allowbreak 62,\allowbreak 33,\allowbreak 56,\allowbreak 31,\allowbreak 59,\allowbreak 29,\allowbreak 72,\allowbreak 18,\allowbreak 73,\allowbreak 20,\allowbreak 76,\allowbreak 16,\allowbreak 77,\allowbreak 13,\allowbreak 75,\allowbreak 17,\allowbreak 80)}\tabularnewline[5pt]
{\raggedright\arraybackslash $L=25$; $L^4$\par $K=48$} & {\ttfamily\footnotesize (100,\allowbreak 1,\allowbreak 99,\allowbreak 2,\allowbreak 98,\allowbreak 3,\allowbreak 97,\allowbreak 4,\allowbreak 96,\allowbreak 5,\allowbreak 95,\allowbreak 6,\allowbreak 94,\allowbreak 7,\allowbreak 93,\allowbreak 8,\allowbreak 92,\allowbreak 9,\allowbreak 91,\allowbreak 10,\allowbreak 90,\allowbreak 11,\allowbreak 89,\allowbreak 12,\allowbreak 88)\par\smallskip (75,\allowbreak 23,\allowbreak 77,\allowbreak 22,\allowbreak 65,\allowbreak 31,\allowbreak 67,\allowbreak 32,\allowbreak 64,\allowbreak 33,\allowbreak 66,\allowbreak 21,\allowbreak 78,\allowbreak 20,\allowbreak 79,\allowbreak 19,\allowbreak 80,\allowbreak 18,\allowbreak 82,\allowbreak 17,\allowbreak 83,\allowbreak 16,\allowbreak 84,\allowbreak 15,\allowbreak 85)\par\smallskip (71,\allowbreak 27,\allowbreak 57,\allowbreak 46,\allowbreak 49,\allowbreak 48,\allowbreak 50,\allowbreak 45,\allowbreak 51,\allowbreak 47,\allowbreak 54,\allowbreak 44,\allowbreak 52,\allowbreak 43,\allowbreak 56,\allowbreak 42,\allowbreak 69,\allowbreak 29,\allowbreak 68,\allowbreak 30,\allowbreak 81,\allowbreak 28,\allowbreak 70,\allowbreak 14,\allowbreak 86)\par\smallskip (73,\allowbreak 24,\allowbreak 74,\allowbreak 26,\allowbreak 72,\allowbreak 25,\allowbreak 62,\allowbreak 36,\allowbreak 61,\allowbreak 37,\allowbreak 60,\allowbreak 38,\allowbreak 59,\allowbreak 39,\allowbreak 55,\allowbreak 40,\allowbreak 58,\allowbreak 41,\allowbreak 53,\allowbreak 34,\allowbreak 63,\allowbreak 35,\allowbreak 76,\allowbreak 13,\allowbreak 87)}\tabularnewline[5pt]
{\raggedright\arraybackslash $L=27$; $L^4$\par $K=52$} & {\ttfamily\footnotesize (108,\allowbreak 1,\allowbreak 107,\allowbreak 2,\allowbreak 106,\allowbreak 3,\allowbreak 105,\allowbreak 4,\allowbreak 104,\allowbreak 5,\allowbreak 103,\allowbreak 6,\allowbreak 102,\allowbreak 7,\allowbreak 101,\allowbreak 8,\allowbreak 100,\allowbreak 9,\allowbreak 99,\allowbreak 10,\allowbreak 98,\allowbreak 11,\allowbreak 97,\allowbreak 12,\allowbreak 96,\allowbreak 13,\allowbreak 95)\par\smallskip (77,\allowbreak 29,\allowbreak 76,\allowbreak 30,\allowbreak 75,\allowbreak 33,\allowbreak 61,\allowbreak 46,\allowbreak 56,\allowbreak 47,\allowbreak 59,\allowbreak 48,\allowbreak 55,\allowbreak 49,\allowbreak 57,\allowbreak 52,\allowbreak 53,\allowbreak 51,\allowbreak 54,\allowbreak 50,\allowbreak 66,\allowbreak 40,\allowbreak 67,\allowbreak 27,\allowbreak 78,\allowbreak 15,\allowbreak 93)\par\smallskip (79,\allowbreak 26,\allowbreak 80,\allowbreak 28,\allowbreak 65,\allowbreak 41,\allowbreak 62,\allowbreak 39,\allowbreak 64,\allowbreak 42,\allowbreak 60,\allowbreak 43,\allowbreak 63,\allowbreak 44,\allowbreak 58,\allowbreak 45,\allowbreak 74,\allowbreak 31,\allowbreak 88,\allowbreak 19,\allowbreak 89,\allowbreak 18,\allowbreak 90,\allowbreak 17,\allowbreak 91,\allowbreak 16,\allowbreak 92)\par\smallskip (81,\allowbreak 25,\allowbreak 83,\allowbreak 24,\allowbreak 84,\allowbreak 23,\allowbreak 85,\allowbreak 21,\allowbreak 86,\allowbreak 20,\allowbreak 87,\allowbreak 32,\allowbreak 73,\allowbreak 34,\allowbreak 72,\allowbreak 22,\allowbreak 71,\allowbreak 35,\allowbreak 70,\allowbreak 36,\allowbreak 69,\allowbreak 37,\allowbreak 68,\allowbreak 38,\allowbreak 82,\allowbreak 14,\allowbreak 94)}\tabularnewline[5pt]
{\raggedright\arraybackslash $L=29$; $L^4$\par $K=56$} & {\ttfamily\footnotesize (87,\allowbreak 42,\allowbreak 71,\allowbreak 43,\allowbreak 70,\allowbreak 54,\allowbreak 60,\allowbreak 55,\allowbreak 58,\allowbreak 56,\allowbreak 57,\allowbreak 53,\allowbreak 61,\allowbreak 50,\allowbreak 59,\allowbreak 52,\allowbreak 62,\allowbreak 47,\allowbreak 68,\allowbreak 46,\allowbreak 65,\allowbreak 45,\allowbreak 69,\allowbreak 44,\allowbreak 67,\allowbreak 29,\allowbreak 84,\allowbreak 16,\allowbreak 100)\par\smallskip (83,\allowbreak 31,\allowbreak 82,\allowbreak 32,\allowbreak 81,\allowbreak 33,\allowbreak 80,\allowbreak 34,\allowbreak 94,\allowbreak 35,\allowbreak 79,\allowbreak 36,\allowbreak 77,\allowbreak 37,\allowbreak 76,\allowbreak 23,\allowbreak 92,\allowbreak 22,\allowbreak 93,\allowbreak 21,\allowbreak 95,\allowbreak 20,\allowbreak 96,\allowbreak 19,\allowbreak 97,\allowbreak 18,\allowbreak 98,\allowbreak 17,\allowbreak 99)\par\smallskip (116,\allowbreak 1,\allowbreak 115,\allowbreak 2,\allowbreak 114,\allowbreak 3,\allowbreak 113,\allowbreak 4,\allowbreak 112,\allowbreak 5,\allowbreak 111,\allowbreak 6,\allowbreak 110,\allowbreak 7,\allowbreak 109,\allowbreak 8,\allowbreak 108,\allowbreak 9,\allowbreak 107,\allowbreak 10,\allowbreak 106,\allowbreak 11,\allowbreak 105,\allowbreak 12,\allowbreak 104,\allowbreak 13,\allowbreak 103,\allowbreak 14,\allowbreak 102)\par\smallskip (85,\allowbreak 28,\allowbreak 86,\allowbreak 30,\allowbreak 72,\allowbreak 41,\allowbreak 73,\allowbreak 40,\allowbreak 74,\allowbreak 39,\allowbreak 75,\allowbreak 38,\allowbreak 64,\allowbreak 51,\allowbreak 63,\allowbreak 49,\allowbreak 66,\allowbreak 48,\allowbreak 78,\allowbreak 24,\allowbreak 91,\allowbreak 25,\allowbreak 90,\allowbreak 26,\allowbreak 89,\allowbreak 27,\allowbreak 88,\allowbreak 15,\allowbreak 101)}\tabularnewline[5pt]
{\raggedright\arraybackslash $L=31$; $L^4$\par $K=60$} & {\ttfamily\footnotesize (124,\allowbreak 1,\allowbreak 123,\allowbreak 2,\allowbreak 122,\allowbreak 3,\allowbreak 121,\allowbreak 4,\allowbreak 120,\allowbreak 5,\allowbreak 119,\allowbreak 6,\allowbreak 118,\allowbreak 7,\allowbreak 117,\allowbreak 8,\allowbreak 116,\allowbreak 9,\allowbreak 115,\allowbreak 10,\allowbreak 114,\allowbreak 11,\allowbreak 113,\allowbreak 12,\allowbreak 112,\allowbreak 13,\allowbreak 111,\allowbreak 14,\allowbreak 110,\allowbreak 15,\allowbreak 109)\par\smallskip (93,\allowbreak 43,\allowbreak 79,\allowbreak 42,\allowbreak 80,\allowbreak 40,\allowbreak 68,\allowbreak 56,\allowbreak 63,\allowbreak 57,\allowbreak 65,\allowbreak 60,\allowbreak 61,\allowbreak 59,\allowbreak 62,\allowbreak 58,\allowbreak 77,\allowbreak 46,\allowbreak 75,\allowbreak 45,\allowbreak 78,\allowbreak 44,\allowbreak 76,\allowbreak 32,\allowbreak 90,\allowbreak 31,\allowbreak 91,\allowbreak 30,\allowbreak 94,\allowbreak 16,\allowbreak 106)\par\smallskip (88,\allowbreak 34,\allowbreak 87,\allowbreak 35,\allowbreak 86,\allowbreak 38,\allowbreak 83,\allowbreak 37,\allowbreak 84,\allowbreak 49,\allowbreak 73,\allowbreak 50,\allowbreak 71,\allowbreak 51,\allowbreak 66,\allowbreak 55,\allowbreak 64,\allowbreak 54,\allowbreak 67,\allowbreak 53,\allowbreak 69,\allowbreak 52,\allowbreak 70,\allowbreak 48,\allowbreak 74,\allowbreak 47,\allowbreak 72,\allowbreak 33,\allowbreak 89,\allowbreak 17,\allowbreak 108)\par\smallskip (92,\allowbreak 29,\allowbreak 95,\allowbreak 28,\allowbreak 96,\allowbreak 27,\allowbreak 97,\allowbreak 26,\allowbreak 81,\allowbreak 39,\allowbreak 82,\allowbreak 41,\allowbreak 98,\allowbreak 25,\allowbreak 99,\allowbreak 24,\allowbreak 100,\allowbreak 23,\allowbreak 85,\allowbreak 36,\allowbreak 101,\allowbreak 22,\allowbreak 102,\allowbreak 21,\allowbreak 103,\allowbreak 20,\allowbreak 104,\allowbreak 19,\allowbreak 105,\allowbreak 18,\allowbreak 107)}\tabularnewline[5pt]
{\raggedright\arraybackslash $L=33$; $L^4$\par $K=64$} & {\ttfamily\footnotesize (95,\allowbreak 35,\allowbreak 77,\allowbreak 62,\allowbreak 65,\allowbreak 64,\allowbreak 66,\allowbreak 61,\allowbreak 67,\allowbreak 63,\allowbreak 70,\allowbreak 60,\allowbreak 68,\allowbreak 59,\allowbreak 72,\allowbreak 58,\allowbreak 69,\allowbreak 57,\allowbreak 74,\allowbreak 56,\allowbreak 91,\allowbreak 39,\allowbreak 90,\allowbreak 40,\allowbreak 107,\allowbreak 38,\allowbreak 92,\allowbreak 37,\allowbreak 93,\allowbreak 36,\allowbreak 94,\allowbreak 18,\allowbreak 114)\par\smallskip (97,\allowbreak 32,\allowbreak 98,\allowbreak 34,\allowbreak 96,\allowbreak 33,\allowbreak 82,\allowbreak 48,\allowbreak 81,\allowbreak 49,\allowbreak 80,\allowbreak 50,\allowbreak 79,\allowbreak 51,\allowbreak 75,\allowbreak 52,\allowbreak 78,\allowbreak 53,\allowbreak 73,\allowbreak 54,\allowbreak 76,\allowbreak 55,\allowbreak 71,\allowbreak 44,\allowbreak 85,\allowbreak 45,\allowbreak 84,\allowbreak 46,\allowbreak 83,\allowbreak 47,\allowbreak 100,\allowbreak 17,\allowbreak 115)\par\smallskip (99,\allowbreak 31,\allowbreak 101,\allowbreak 30,\allowbreak 102,\allowbreak 29,\allowbreak 103,\allowbreak 28,\allowbreak 87,\allowbreak 41,\allowbreak 89,\allowbreak 42,\allowbreak 86,\allowbreak 43,\allowbreak 88,\allowbreak 27,\allowbreak 104,\allowbreak 26,\allowbreak 105,\allowbreak 25,\allowbreak 106,\allowbreak 24,\allowbreak 108,\allowbreak 23,\allowbreak 109,\allowbreak 22,\allowbreak 110,\allowbreak 21,\allowbreak 111,\allowbreak 20,\allowbreak 112,\allowbreak 19,\allowbreak 113)\par\smallskip (132,\allowbreak 1,\allowbreak 131,\allowbreak 2,\allowbreak 130,\allowbreak 3,\allowbreak 129,\allowbreak 4,\allowbreak 128,\allowbreak 5,\allowbreak 127,\allowbreak 6,\allowbreak 126,\allowbreak 7,\allowbreak 125,\allowbreak 8,\allowbreak 124,\allowbreak 9,\allowbreak 123,\allowbreak 10,\allowbreak 122,\allowbreak 11,\allowbreak 121,\allowbreak 12,\allowbreak 120,\allowbreak 13,\allowbreak 119,\allowbreak 14,\allowbreak 118,\allowbreak 15,\allowbreak 117,\allowbreak 16,\allowbreak 116)}\tabularnewline[5pt]
{\raggedright\arraybackslash $L=35$; $L^4$\par $K=68$} & {\ttfamily\footnotesize (101,\allowbreak 37,\allowbreak 100,\allowbreak 38,\allowbreak 99,\allowbreak 39,\allowbreak 98,\allowbreak 40,\allowbreak 97,\allowbreak 43,\allowbreak 79,\allowbreak 60,\allowbreak 74,\allowbreak 61,\allowbreak 77,\allowbreak 62,\allowbreak 72,\allowbreak 63,\allowbreak 75,\allowbreak 64,\allowbreak 71,\allowbreak 65,\allowbreak 73,\allowbreak 68,\allowbreak 69,\allowbreak 67,\allowbreak 70,\allowbreak 66,\allowbreak 86,\allowbreak 52,\allowbreak 87,\allowbreak 35,\allowbreak 102,\allowbreak 19,\allowbreak 121)\par\smallskip (103,\allowbreak 34,\allowbreak 104,\allowbreak 36,\allowbreak 85,\allowbreak 53,\allowbreak 82,\allowbreak 51,\allowbreak 84,\allowbreak 54,\allowbreak 80,\allowbreak 55,\allowbreak 83,\allowbreak 56,\allowbreak 78,\allowbreak 57,\allowbreak 81,\allowbreak 58,\allowbreak 76,\allowbreak 59,\allowbreak 96,\allowbreak 41,\allowbreak 114,\allowbreak 25,\allowbreak 115,\allowbreak 24,\allowbreak 116,\allowbreak 23,\allowbreak 117,\allowbreak 22,\allowbreak 118,\allowbreak 21,\allowbreak 119,\allowbreak 20,\allowbreak 120)\par\smallskip (105,\allowbreak 33,\allowbreak 107,\allowbreak 32,\allowbreak 108,\allowbreak 31,\allowbreak 109,\allowbreak 30,\allowbreak 110,\allowbreak 29,\allowbreak 111,\allowbreak 27,\allowbreak 112,\allowbreak 26,\allowbreak 113,\allowbreak 42,\allowbreak 95,\allowbreak 44,\allowbreak 94,\allowbreak 28,\allowbreak 93,\allowbreak 45,\allowbreak 92,\allowbreak 46,\allowbreak 91,\allowbreak 47,\allowbreak 90,\allowbreak 48,\allowbreak 89,\allowbreak 49,\allowbreak 88,\allowbreak 50,\allowbreak 106,\allowbreak 18,\allowbreak 122)\par\smallskip (140,\allowbreak 1,\allowbreak 139,\allowbreak 2,\allowbreak 138,\allowbreak 3,\allowbreak 137,\allowbreak 4,\allowbreak 136,\allowbreak 5,\allowbreak 135,\allowbreak 6,\allowbreak 134,\allowbreak 7,\allowbreak 133,\allowbreak 8,\allowbreak 132,\allowbreak 9,\allowbreak 131,\allowbreak 10,\allowbreak 130,\allowbreak 11,\allowbreak 129,\allowbreak 12,\allowbreak 128,\allowbreak 13,\allowbreak 127,\allowbreak 14,\allowbreak 126,\allowbreak 15,\allowbreak 125,\allowbreak 16,\allowbreak 124,\allowbreak 17,\allowbreak 123)}\tabularnewline[5pt]
{\raggedright\arraybackslash $L=37$; $L^4$\par $K=72$} & {\ttfamily\footnotesize (148,\allowbreak 1,\allowbreak 147,\allowbreak 2,\allowbreak 146,\allowbreak 3,\allowbreak 145,\allowbreak 4,\allowbreak 144,\allowbreak 5,\allowbreak 143,\allowbreak 6,\allowbreak 142,\allowbreak 7,\allowbreak 141,\allowbreak 8,\allowbreak 140,\allowbreak 9,\allowbreak 139,\allowbreak 10,\allowbreak 138,\allowbreak 11,\allowbreak 137,\allowbreak 12,\allowbreak 136,\allowbreak 13,\allowbreak 135,\allowbreak 14,\allowbreak 134,\allowbreak 15,\allowbreak 133,\allowbreak 16,\allowbreak 132,\allowbreak 17,\allowbreak 131,\allowbreak 18,\allowbreak 130)\par\smallskip (111,\allowbreak 36,\allowbreak 109,\allowbreak 35,\allowbreak 113,\allowbreak 34,\allowbreak 114,\allowbreak 33,\allowbreak 115,\allowbreak 32,\allowbreak 116,\allowbreak 31,\allowbreak 117,\allowbreak 29,\allowbreak 98,\allowbreak 48,\allowbreak 97,\allowbreak 30,\allowbreak 119,\allowbreak 43,\allowbreak 120,\allowbreak 28,\allowbreak 118,\allowbreak 27,\allowbreak 121,\allowbreak 26,\allowbreak 122,\allowbreak 25,\allowbreak 123,\allowbreak 24,\allowbreak 124,\allowbreak 23,\allowbreak 125,\allowbreak 22,\allowbreak 126,\allowbreak 21,\allowbreak 129)\par\smallskip (110,\allowbreak 38,\allowbreak 108,\allowbreak 37,\allowbreak 92,\allowbreak 53,\allowbreak 93,\allowbreak 55,\allowbreak 90,\allowbreak 54,\allowbreak 91,\allowbreak 70,\allowbreak 74,\allowbreak 69,\allowbreak 75,\allowbreak 72,\allowbreak 73,\allowbreak 71,\allowbreak 78,\allowbreak 68,\allowbreak 76,\allowbreak 67,\allowbreak 81,\allowbreak 47,\allowbreak 99,\allowbreak 46,\allowbreak 100,\allowbreak 49,\allowbreak 96,\allowbreak 50,\allowbreak 95,\allowbreak 51,\allowbreak 94,\allowbreak 52,\allowbreak 112,\allowbreak 19,\allowbreak 128)\par\smallskip (106,\allowbreak 40,\allowbreak 105,\allowbreak 41,\allowbreak 104,\allowbreak 42,\allowbreak 103,\allowbreak 44,\allowbreak 102,\allowbreak 45,\allowbreak 101,\allowbreak 60,\allowbreak 86,\allowbreak 61,\allowbreak 84,\allowbreak 62,\allowbreak 79,\allowbreak 66,\allowbreak 77,\allowbreak 65,\allowbreak 80,\allowbreak 64,\allowbreak 82,\allowbreak 63,\allowbreak 83,\allowbreak 59,\allowbreak 88,\allowbreak 58,\allowbreak 85,\allowbreak 57,\allowbreak 89,\allowbreak 56,\allowbreak 87,\allowbreak 39,\allowbreak 107,\allowbreak 20,\allowbreak 127)}\tabularnewline[5pt]
{\raggedright\arraybackslash $L=39$; $L^4$\par $K=76$} & {\ttfamily\footnotesize (117,\allowbreak 55,\allowbreak 99,\allowbreak 54,\allowbreak 100,\allowbreak 53,\allowbreak 101,\allowbreak 52,\allowbreak 102,\allowbreak 50,\allowbreak 86,\allowbreak 70,\allowbreak 80,\allowbreak 71,\allowbreak 83,\allowbreak 72,\allowbreak 79,\allowbreak 73,\allowbreak 81,\allowbreak 76,\allowbreak 77,\allowbreak 75,\allowbreak 78,\allowbreak 74,\allowbreak 97,\allowbreak 58,\allowbreak 95,\allowbreak 57,\allowbreak 98,\allowbreak 56,\allowbreak 96,\allowbreak 40,\allowbreak 114,\allowbreak 39,\allowbreak 115,\allowbreak 38,\allowbreak 118,\allowbreak 20,\allowbreak 134)\par\smallskip (112,\allowbreak 42,\allowbreak 111,\allowbreak 43,\allowbreak 110,\allowbreak 44,\allowbreak 109,\allowbreak 45,\allowbreak 108,\allowbreak 48,\allowbreak 105,\allowbreak 47,\allowbreak 106,\allowbreak 63,\allowbreak 91,\allowbreak 64,\allowbreak 89,\allowbreak 65,\allowbreak 84,\allowbreak 69,\allowbreak 82,\allowbreak 68,\allowbreak 85,\allowbreak 67,\allowbreak 87,\allowbreak 66,\allowbreak 88,\allowbreak 62,\allowbreak 93,\allowbreak 61,\allowbreak 90,\allowbreak 60,\allowbreak 94,\allowbreak 59,\allowbreak 92,\allowbreak 41,\allowbreak 113,\allowbreak 21,\allowbreak 136)\par\smallskip (116,\allowbreak 37,\allowbreak 119,\allowbreak 36,\allowbreak 120,\allowbreak 35,\allowbreak 121,\allowbreak 34,\allowbreak 122,\allowbreak 33,\allowbreak 123,\allowbreak 32,\allowbreak 103,\allowbreak 49,\allowbreak 104,\allowbreak 51,\allowbreak 124,\allowbreak 31,\allowbreak 125,\allowbreak 30,\allowbreak 126,\allowbreak 29,\allowbreak 107,\allowbreak 46,\allowbreak 127,\allowbreak 28,\allowbreak 128,\allowbreak 27,\allowbreak 129,\allowbreak 26,\allowbreak 130,\allowbreak 25,\allowbreak 131,\allowbreak 24,\allowbreak 132,\allowbreak 23,\allowbreak 133,\allowbreak 22,\allowbreak 135)\par\smallskip (156,\allowbreak 1,\allowbreak 155,\allowbreak 2,\allowbreak 154,\allowbreak 3,\allowbreak 153,\allowbreak 4,\allowbreak 152,\allowbreak 5,\allowbreak 151,\allowbreak 6,\allowbreak 150,\allowbreak 7,\allowbreak 149,\allowbreak 8,\allowbreak 148,\allowbreak 9,\allowbreak 147,\allowbreak 10,\allowbreak 146,\allowbreak 11,\allowbreak 145,\allowbreak 12,\allowbreak 144,\allowbreak 13,\allowbreak 143,\allowbreak 14,\allowbreak 142,\allowbreak 15,\allowbreak 141,\allowbreak 16,\allowbreak 140,\allowbreak 17,\allowbreak 139,\allowbreak 18,\allowbreak 138,\allowbreak 19,\allowbreak 137)}\tabularnewline[5pt]
{\raggedright\arraybackslash $L=41$; $L^4$\par $K=80$} & {\ttfamily\footnotesize (119,\allowbreak 43,\allowbreak 97,\allowbreak 78,\allowbreak 81,\allowbreak 80,\allowbreak 82,\allowbreak 77,\allowbreak 83,\allowbreak 79,\allowbreak 86,\allowbreak 76,\allowbreak 84,\allowbreak 75,\allowbreak 88,\allowbreak 74,\allowbreak 85,\allowbreak 73,\allowbreak 90,\allowbreak 72,\allowbreak 87,\allowbreak 71,\allowbreak 92,\allowbreak 70,\allowbreak 113,\allowbreak 49,\allowbreak 112,\allowbreak 50,\allowbreak 133,\allowbreak 48,\allowbreak 114,\allowbreak 47,\allowbreak 115,\allowbreak 46,\allowbreak 116,\allowbreak 45,\allowbreak 117,\allowbreak 44,\allowbreak 118,\allowbreak 22,\allowbreak 142)\par\smallskip (121,\allowbreak 40,\allowbreak 122,\allowbreak 42,\allowbreak 120,\allowbreak 41,\allowbreak 102,\allowbreak 60,\allowbreak 101,\allowbreak 61,\allowbreak 100,\allowbreak 62,\allowbreak 99,\allowbreak 63,\allowbreak 95,\allowbreak 64,\allowbreak 98,\allowbreak 65,\allowbreak 93,\allowbreak 66,\allowbreak 96,\allowbreak 67,\allowbreak 91,\allowbreak 68,\allowbreak 94,\allowbreak 69,\allowbreak 89,\allowbreak 54,\allowbreak 107,\allowbreak 55,\allowbreak 106,\allowbreak 56,\allowbreak 105,\allowbreak 57,\allowbreak 104,\allowbreak 58,\allowbreak 103,\allowbreak 59,\allowbreak 124,\allowbreak 21,\allowbreak 143)\par\smallskip (123,\allowbreak 39,\allowbreak 125,\allowbreak 38,\allowbreak 126,\allowbreak 37,\allowbreak 127,\allowbreak 36,\allowbreak 128,\allowbreak 35,\allowbreak 129,\allowbreak 34,\allowbreak 109,\allowbreak 51,\allowbreak 111,\allowbreak 52,\allowbreak 108,\allowbreak 53,\allowbreak 110,\allowbreak 33,\allowbreak 130,\allowbreak 32,\allowbreak 131,\allowbreak 31,\allowbreak 132,\allowbreak 30,\allowbreak 134,\allowbreak 29,\allowbreak 135,\allowbreak 28,\allowbreak 136,\allowbreak 27,\allowbreak 137,\allowbreak 26,\allowbreak 138,\allowbreak 25,\allowbreak 139,\allowbreak 24,\allowbreak 140,\allowbreak 23,\allowbreak 141)\par\smallskip (164,\allowbreak 1,\allowbreak 163,\allowbreak 2,\allowbreak 162,\allowbreak 3,\allowbreak 161,\allowbreak 4,\allowbreak 160,\allowbreak 5,\allowbreak 159,\allowbreak 6,\allowbreak 158,\allowbreak 7,\allowbreak 157,\allowbreak 8,\allowbreak 156,\allowbreak 9,\allowbreak 155,\allowbreak 10,\allowbreak 154,\allowbreak 11,\allowbreak 153,\allowbreak 12,\allowbreak 152,\allowbreak 13,\allowbreak 151,\allowbreak 14,\allowbreak 150,\allowbreak 15,\allowbreak 149,\allowbreak 16,\allowbreak 148,\allowbreak 17,\allowbreak 147,\allowbreak 18,\allowbreak 146,\allowbreak 19,\allowbreak 145,\allowbreak 20,\allowbreak 144)}\tabularnewline[5pt]
{\raggedright\arraybackslash $L=43$; $L^4$\par $K=84$} & {\ttfamily\footnotesize (172,\allowbreak 1,\allowbreak 171,\allowbreak 2,\allowbreak 170,\allowbreak 3,\allowbreak 169,\allowbreak 4,\allowbreak 168,\allowbreak 5,\allowbreak 167,\allowbreak 6,\allowbreak 166,\allowbreak 7,\allowbreak 165,\allowbreak 8,\allowbreak 164,\allowbreak 9,\allowbreak 163,\allowbreak 10,\allowbreak 162,\allowbreak 11,\allowbreak 161,\allowbreak 12,\allowbreak 160,\allowbreak 13,\allowbreak 159,\allowbreak 14,\allowbreak 158,\allowbreak 15,\allowbreak 157,\allowbreak 16,\allowbreak 156,\allowbreak 17,\allowbreak 155,\allowbreak 18,\allowbreak 154,\allowbreak 19,\allowbreak 153,\allowbreak 20,\allowbreak 152,\allowbreak 21,\allowbreak 151)\par\smallskip (125,\allowbreak 45,\allowbreak 124,\allowbreak 46,\allowbreak 123,\allowbreak 47,\allowbreak 122,\allowbreak 48,\allowbreak 121,\allowbreak 49,\allowbreak 120,\allowbreak 50,\allowbreak 119,\allowbreak 53,\allowbreak 97,\allowbreak 74,\allowbreak 92,\allowbreak 75,\allowbreak 95,\allowbreak 76,\allowbreak 90,\allowbreak 77,\allowbreak 93,\allowbreak 78,\allowbreak 88,\allowbreak 79,\allowbreak 91,\allowbreak 80,\allowbreak 87,\allowbreak 81,\allowbreak 89,\allowbreak 84,\allowbreak 85,\allowbreak 83,\allowbreak 86,\allowbreak 82,\allowbreak 106,\allowbreak 64,\allowbreak 107,\allowbreak 43,\allowbreak 126,\allowbreak 23,\allowbreak 149)\par\smallskip (127,\allowbreak 42,\allowbreak 128,\allowbreak 44,\allowbreak 105,\allowbreak 65,\allowbreak 102,\allowbreak 63,\allowbreak 104,\allowbreak 66,\allowbreak 100,\allowbreak 67,\allowbreak 103,\allowbreak 68,\allowbreak 98,\allowbreak 69,\allowbreak 101,\allowbreak 70,\allowbreak 96,\allowbreak 71,\allowbreak 99,\allowbreak 72,\allowbreak 94,\allowbreak 73,\allowbreak 118,\allowbreak 51,\allowbreak 140,\allowbreak 31,\allowbreak 141,\allowbreak 30,\allowbreak 142,\allowbreak 29,\allowbreak 143,\allowbreak 28,\allowbreak 144,\allowbreak 27,\allowbreak 145,\allowbreak 26,\allowbreak 146,\allowbreak 25,\allowbreak 147,\allowbreak 24,\allowbreak 148)\par\smallskip (129,\allowbreak 41,\allowbreak 131,\allowbreak 40,\allowbreak 132,\allowbreak 39,\allowbreak 133,\allowbreak 38,\allowbreak 134,\allowbreak 37,\allowbreak 135,\allowbreak 36,\allowbreak 136,\allowbreak 35,\allowbreak 137,\allowbreak 33,\allowbreak 138,\allowbreak 32,\allowbreak 139,\allowbreak 52,\allowbreak 117,\allowbreak 54,\allowbreak 116,\allowbreak 34,\allowbreak 115,\allowbreak 55,\allowbreak 114,\allowbreak 56,\allowbreak 113,\allowbreak 57,\allowbreak 112,\allowbreak 58,\allowbreak 111,\allowbreak 59,\allowbreak 110,\allowbreak 60,\allowbreak 109,\allowbreak 61,\allowbreak 108,\allowbreak 62,\allowbreak 130,\allowbreak 22,\allowbreak 150)}\tabularnewline[5pt]
{\raggedright\arraybackslash $L=45$; $L^4$\par $K=88$} & {\ttfamily\footnotesize (135,\allowbreak 44,\allowbreak 133,\allowbreak 43,\allowbreak 137,\allowbreak 42,\allowbreak 138,\allowbreak 41,\allowbreak 139,\allowbreak 40,\allowbreak 140,\allowbreak 39,\allowbreak 141,\allowbreak 38,\allowbreak 142,\allowbreak 37,\allowbreak 143,\allowbreak 35,\allowbreak 120,\allowbreak 58,\allowbreak 119,\allowbreak 36,\allowbreak 145,\allowbreak 53,\allowbreak 146,\allowbreak 34,\allowbreak 144,\allowbreak 33,\allowbreak 147,\allowbreak 32,\allowbreak 148,\allowbreak 31,\allowbreak 149,\allowbreak 30,\allowbreak 150,\allowbreak 29,\allowbreak 151,\allowbreak 28,\allowbreak 152,\allowbreak 27,\allowbreak 153,\allowbreak 26,\allowbreak 154,\allowbreak 25,\allowbreak 157)\par\smallskip (130,\allowbreak 48,\allowbreak 129,\allowbreak 49,\allowbreak 128,\allowbreak 50,\allowbreak 127,\allowbreak 51,\allowbreak 126,\allowbreak 52,\allowbreak 125,\allowbreak 54,\allowbreak 124,\allowbreak 55,\allowbreak 123,\allowbreak 74,\allowbreak 104,\allowbreak 75,\allowbreak 102,\allowbreak 76,\allowbreak 97,\allowbreak 80,\allowbreak 95,\allowbreak 79,\allowbreak 98,\allowbreak 78,\allowbreak 100,\allowbreak 77,\allowbreak 101,\allowbreak 73,\allowbreak 106,\allowbreak 72,\allowbreak 103,\allowbreak 71,\allowbreak 108,\allowbreak 70,\allowbreak 105,\allowbreak 69,\allowbreak 109,\allowbreak 68,\allowbreak 107,\allowbreak 47,\allowbreak 131,\allowbreak 24,\allowbreak 155)\par\smallskip (180,\allowbreak 1,\allowbreak 179,\allowbreak 2,\allowbreak 178,\allowbreak 3,\allowbreak 177,\allowbreak 4,\allowbreak 176,\allowbreak 5,\allowbreak 175,\allowbreak 6,\allowbreak 174,\allowbreak 7,\allowbreak 173,\allowbreak 8,\allowbreak 172,\allowbreak 9,\allowbreak 171,\allowbreak 10,\allowbreak 170,\allowbreak 11,\allowbreak 169,\allowbreak 12,\allowbreak 168,\allowbreak 13,\allowbreak 167,\allowbreak 14,\allowbreak 166,\allowbreak 15,\allowbreak 165,\allowbreak 16,\allowbreak 164,\allowbreak 17,\allowbreak 163,\allowbreak 18,\allowbreak 162,\allowbreak 19,\allowbreak 161,\allowbreak 20,\allowbreak 160,\allowbreak 21,\allowbreak 159,\allowbreak 22,\allowbreak 158)\par\smallskip (134,\allowbreak 46,\allowbreak 132,\allowbreak 45,\allowbreak 112,\allowbreak 65,\allowbreak 113,\allowbreak 67,\allowbreak 110,\allowbreak 66,\allowbreak 111,\allowbreak 86,\allowbreak 90,\allowbreak 85,\allowbreak 91,\allowbreak 88,\allowbreak 89,\allowbreak 87,\allowbreak 94,\allowbreak 84,\allowbreak 92,\allowbreak 83,\allowbreak 96,\allowbreak 82,\allowbreak 93,\allowbreak 81,\allowbreak 99,\allowbreak 57,\allowbreak 121,\allowbreak 56,\allowbreak 122,\allowbreak 59,\allowbreak 118,\allowbreak 60,\allowbreak 117,\allowbreak 61,\allowbreak 116,\allowbreak 62,\allowbreak 115,\allowbreak 63,\allowbreak 114,\allowbreak 64,\allowbreak 136,\allowbreak 23,\allowbreak 156)}\tabularnewline[5pt]
{\raggedright\arraybackslash $L=47$; $L^4$\par $K=92$} & {\ttfamily\footnotesize (188,\allowbreak 1,\allowbreak 187,\allowbreak 2,\allowbreak 186,\allowbreak 3,\allowbreak 185,\allowbreak 4,\allowbreak 184,\allowbreak 5,\allowbreak 183,\allowbreak 6,\allowbreak 182,\allowbreak 7,\allowbreak 181,\allowbreak 8,\allowbreak 180,\allowbreak 9,\allowbreak 179,\allowbreak 10,\allowbreak 178,\allowbreak 11,\allowbreak 177,\allowbreak 12,\allowbreak 176,\allowbreak 13,\allowbreak 175,\allowbreak 14,\allowbreak 174,\allowbreak 15,\allowbreak 173,\allowbreak 16,\allowbreak 172,\allowbreak 17,\allowbreak 171,\allowbreak 18,\allowbreak 170,\allowbreak 19,\allowbreak 169,\allowbreak 20,\allowbreak 168,\allowbreak 21,\allowbreak 167,\allowbreak 22,\allowbreak 166,\allowbreak 23,\allowbreak 165)\par\smallskip (141,\allowbreak 67,\allowbreak 119,\allowbreak 66,\allowbreak 120,\allowbreak 65,\allowbreak 121,\allowbreak 64,\allowbreak 122,\allowbreak 63,\allowbreak 123,\allowbreak 62,\allowbreak 124,\allowbreak 60,\allowbreak 104,\allowbreak 84,\allowbreak 98,\allowbreak 85,\allowbreak 101,\allowbreak 86,\allowbreak 96,\allowbreak 87,\allowbreak 99,\allowbreak 88,\allowbreak 95,\allowbreak 89,\allowbreak 97,\allowbreak 92,\allowbreak 93,\allowbreak 91,\allowbreak 94,\allowbreak 90,\allowbreak 117,\allowbreak 70,\allowbreak 115,\allowbreak 69,\allowbreak 118,\allowbreak 68,\allowbreak 116,\allowbreak 48,\allowbreak 138,\allowbreak 47,\allowbreak 139,\allowbreak 46,\allowbreak 142,\allowbreak 24,\allowbreak 162)\par\smallskip (136,\allowbreak 50,\allowbreak 135,\allowbreak 51,\allowbreak 134,\allowbreak 52,\allowbreak 133,\allowbreak 53,\allowbreak 132,\allowbreak 54,\allowbreak 131,\allowbreak 55,\allowbreak 130,\allowbreak 58,\allowbreak 127,\allowbreak 57,\allowbreak 128,\allowbreak 77,\allowbreak 109,\allowbreak 78,\allowbreak 107,\allowbreak 79,\allowbreak 102,\allowbreak 83,\allowbreak 100,\allowbreak 82,\allowbreak 103,\allowbreak 81,\allowbreak 105,\allowbreak 80,\allowbreak 106,\allowbreak 76,\allowbreak 111,\allowbreak 75,\allowbreak 108,\allowbreak 74,\allowbreak 113,\allowbreak 73,\allowbreak 110,\allowbreak 72,\allowbreak 114,\allowbreak 71,\allowbreak 112,\allowbreak 49,\allowbreak 137,\allowbreak 25,\allowbreak 164)\par\smallskip (140,\allowbreak 45,\allowbreak 143,\allowbreak 44,\allowbreak 144,\allowbreak 43,\allowbreak 145,\allowbreak 42,\allowbreak 146,\allowbreak 41,\allowbreak 147,\allowbreak 40,\allowbreak 148,\allowbreak 39,\allowbreak 149,\allowbreak 38,\allowbreak 125,\allowbreak 59,\allowbreak 126,\allowbreak 61,\allowbreak 150,\allowbreak 37,\allowbreak 151,\allowbreak 36,\allowbreak 152,\allowbreak 35,\allowbreak 129,\allowbreak 56,\allowbreak 153,\allowbreak 34,\allowbreak 154,\allowbreak 33,\allowbreak 155,\allowbreak 32,\allowbreak 156,\allowbreak 31,\allowbreak 157,\allowbreak 30,\allowbreak 158,\allowbreak 29,\allowbreak 159,\allowbreak 28,\allowbreak 160,\allowbreak 27,\allowbreak 161,\allowbreak 26,\allowbreak 163)}\tabularnewline[5pt]
\end{longtable}
\endgroup


\section{Exact eight-arm construction predicates}
\label{eight:specification}
This section specifies the predicates $\mathsf A_8^\epsilon$, $\mathsf O_8^\epsilon$ and $\mathsf H_8^\epsilon$ used in the article, with $\epsilon=o$ or $e$. Within each of the two domains below we suppress the subscript and superscript: $\A$, $\mathsf O$ and $\Hh$ always mean that domain's eight-arm predicates. They do not mean the ten-family predicate $\mathsf A_7$. Set $K(I)=\sum_{i\in I}\floor{L_i/2}$; Outer-route complements are taken in the eight-element index set; within a recursive rule, a complementary pair or triple is taken within the selected subset. All offsets are nonnegative integers.

\begin{lemma}\label{eight:single-equal-offsets}
An arm of length $L$ admits every offset $0\le k<L$.  More generally,
for $r,L\ge1$, the equal-arm spider $\Sp{L^r}$ admits every such offset.
\end{lemma}
\begin{proof}
For the single arm, take an ordinary graceful permutation of
$1,\ldots,L$ beginning at $L-k$, and add $k$ to each entry.  The root
edge has color $L$, and the internal colors are $1,\ldots,L-1$.
The prescribed-first-label path theorem is
Lemma~\ref{common:prescribed-first-path}; it includes the ordinary
non-alpha permutations needed at the exceptional odd-order requests.
For the equal-arm assertion, subdivide the root-zero $r$-leaf unit star
by factor $L$.  Theorem~\ref{common:offset-subdivision}, applied at source
offset zero, gives every residue $0,\ldots,L-1$.
\end{proof}

\subsection{The twelve odd-pair clauses}\label{eight:odd-pairs}

For $a=2t+1\le b=2u+1$, where $t,u\ge1$, define $P_o(a,b;k)$ to mean
$k\ge0$ and at least one of the following twelve conditions.  Parameters
introduced in one condition are local to that condition.

\begin{enumerate}
\item[(P1)] $u\le k\le a-1$, and either
  $4k\ne4a-b-3$ or $3\mid a$ and $3\mid b$.
\item[(P2)] $k\le t$, and either $4k\ne a-1$ or $2b\le3a-1$.
\item[(P3)] $k\le u$, $k+u\le a-1$, and $4k\ne b-1$.
\item[(P4)] $a=b$ and $k<a$.
\item[(P5)] $a=b\ge5$ and $2k\le3a-5$.
\item[(P6)] $a\ge7$, $b=a+2$, and $k=2a$.
\item[(P7)] $a\ge7$, $b=a+2$, $a\le k$, and $2k\le3a-1$.
\item[(P8)] With $S=k-2u-2$, $\delta=u-t$, and $V=k-t-u+1$,
  \[
   S\ge2,\quad t\ge2S+2,\quad 0\le\delta\le S-1,
   \quad\delta\ne S-3,\quad1\le V\le t.
  \]
\item[(P9)] With $S=k-2t$ and $V=k-t-u+1$,
  \[
   1\le S\le t\le u\le2S-2,\quad u>S,
   \quad u\ne2t-S,\quad1\le V\le u.
  \]
\item[(P10)] With $S=k-2t$, $g=u-2S$, $h=k-t-u$, and $V=k-t-u+1$,
  \[
   g\ge1,\quad h\ge2,\quad S\ge g+h+1,\quad1\le V\le u.
  \]
\item[(P11)] With $S=u-t+1$ and $g=2t-u-2$,
  \[
   k=t+u+2,\quad g\ge1,\quad S\ge g+3,\quad u\ge3.
  \]
\item[(P12)] With $S=3u-k+3$, $h=t+2u-k+1$,
  $g=2k-4u-t-4$, and $V=k-t-u+1$,
  \[
   g\ge1,\quad h\ge2,\quad S\ge g+h+1,\quad1\le V\le t.
  \]
\end{enumerate}

\begin{lemma}\label{eight:odd-pair-sufficiency}
If $P_o(a,b;k)$ holds, then $k\in\Om(\Sp{a,b})$.
\end{lemma}
\begin{proof}
Clauses (P1)--(P3) are the HIGH, LOW, and reversed-LOW constructions,
with their exact exceptions and repairs, proved in
Section~\ref{app:pair-offsets}.  Clause (P4) is
Lemma~\ref{eight:single-equal-offsets}.  The extended equal-odd interval
(P5), the consecutive interval (P7), the two general kernel bands
(P8)--(P9), and the three asymmetric bands (P10)--(P12) are constructed
in Section~\ref{construction:odd-pair-kernels}.  The consecutive maximum
(P6) also follows from the complete maximal-pair theorem in
Section~\ref{app:pair-offsets}.

In (P12) the kernel construction interchanges the order of the two
arms.  This does not change the rooted spider, but the restriction
$V\le t$ is part of its even-path premise and must be retained.
The consecutive construction includes its fifteen literal small kernels;
the kernel-band proofs include all source endpoint and hole conditions.

The divisibility repair in (P1) is used only within its source scope.
Indeed, if $a=3$, HIGH gives $u\le k\le2$.  If also $3\mid b$ and
$b\ge a$, this forces $b=3$.  Its exceptional equality would be
$4k=6$, impossible for integer $k$.  Every active repaired exceptional
request therefore has $a,b\ge9$, so its quotient lengths are at least
three.  No small forbidden source request is used.
\end{proof}

\subsection{Finite recursive construction rules}\label{eight:recursions}

This section first treats odd lengths $L_i\ge3$.  Every index subset
inherits the increasing order of the full list.  The same recursions
will later be used for one even minimum, with the pair predicate and
the insertion capacities changed as explicitly stated there.

\begin{definition}\label{eight:alpha-rule}
The predicate $\A(I)$ is defined by the following finite recursion.
\begin{enumerate}
\item For $|I|=1$ or $2$, it is true.
\item For $|I|=3$, it means $\max_{i\in I}L_i\ge5$.
\item For $|I|\ge4$, it holds if either $|I|=4$ and all four lengths
  are equal odd integers at least seven, or some $j\in I$ satisfies
  \[
   \A(I\setminus\{j\}),\qquad
   L_j>2K(I\setminus\{j\}),\qquad
   L_j\ne4K(I\setminus\{j\})+1.
  \]
\end{enumerate}
\end{definition}

These are root-zero alpha constructions of index $K(I)$.  The triple
rule uses the complete odd triple theorem: among odd lengths at least
three, the only omitted profile is $(3,3,3)$
(Section~\ref{construction:all-odd-alpha3}).  The equal-four rule uses
Section~\ref{construction:equal-odd-alpha4}; its lower bound seven is
retained.  The last rule is Lemma~\ref{common:alpha-extension}, with
both its strict threshold and its exceptional length intact.

\begin{definition}\label{eight:offset-rule}
Require $k\ge0$.  For one arm, $\mathsf O(\{i\},k)$ means $k<L_i$;
for two arms, it means $P_o(L_i,L_j;k)$.  For three ordered lengths
$x\le y\le z$, let $T(x,y,z;k)$ be the disjunction of:
\begin{enumerate}
\item $k=0$;
\item $y=z$, $\floor{y/2}\le k\le x-1$, and either $x=y$ or both
  \[
   4k\ne3(y-1),\qquad 4k\ne4x-y-3;
  \]
\item $x=z$ and $k<x$;
\item some arm $L_j$ satisfies
  \[
   k\le\floor{L_j/2},\qquad4k\ne L_j-1,
  \]
  and its complementary pair has predicate $P_o$ at offset
  $k+\floor{L_j/2}$.
\end{enumerate}
This is $\mathsf O(I,k)$ when $|I|=3$.  For $|I|=4$, let
$\mathsf O(I,k)$ hold if all four lengths agree and $k$ is below that
common length, or if some $j\in I$ satisfies the same two displayed
conditions and its complementary triple satisfies $T$ at offset
$k+\floor{L_j/2}$.  Put $\mathsf O(I,k)=\mathrm{false}$ for $|I|\ge5$.
\end{definition}

The triple rule at $k=0$ uses ordinary hub-zero gracefulness, not an
alpha assertion about $(3,3,3)$.  The repeated-long rule and the one-leg
insertion rule, including their excluded equalities, are proved in
Section~\ref{app:triple-offsets}.  The equal-arm rules follow from
Lemma~\ref{eight:single-equal-offsets}.  Thus each true offset predicate
has the stated offset realization.

\begin{definition}\label{eight:hub-rule}
Put $\Hh(I)=\mathrm{true}$ for $|I|\le3$.  For four ordered lengths
$a\le b\le c\le d$, let $\Hh(I)$ be the disjunction of:
\begin{enumerate}
\item $a\le5$;
\item $a=b$ and $c=d$;
\item some triple $J\subset I$ has $\A(J)$, and its complementary
  arm $j$ satisfies $L_j>K(J)$;
\item some arm $j$ has
  $\mathsf O(I\setminus\{j\},\floor{L_j/2})$;
\item some pair $J\subset I$ has $\mathsf O(I\setminus J,K(J))$.
\end{enumerate}
For $|I|\ge5$, put
\[
 \Hh(I)=\A(I)\ \lor\!
 \bigvee_{\varnothing\ne J\subsetneq I}
 \bigl(\A(J)\land\mathsf O(I\setminus J,K(J))\bigr).
\]
\end{definition}

\begin{lemma}\label{eight:predicate-soundness}
Every true $\A(I)$ supplies a root-zero alpha labeling of index $K(I)$;
every true $\mathsf O(I,k)$ supplies offset $k$; and every true $\Hh(I)$
supplies an ordinary hub-zero graceful labeling.
\end{lemma}
\begin{proof}
The alpha and offset assertions follow by induction on arity from the
construction lemmas just specified.  The first two four-arm hub rules
follow from Lemma~\ref{seven:four-short} and
Section~\ref{po:paired-four}, respectively. The other three are
alpha-offset amalgamation, using respectively an offset single arm,
an alpha single arm, or an alpha pair.  The pair-root exception $(2,2)$
cannot occur on odd lengths.  The higher-arity hub rule is exactly
alpha-offset amalgamation, or the alpha labeling itself.  Its
singleton-offset term is the ordinary arm-attachment condition $L>K$;
it must not be confused with the alpha-preserving condition
$L>2K$, $L\ne4K+1$.
\end{proof}

\subsection{One even minimum}\label{eight:even-minimum}

For this section put
\[
 L_0=2h_0,\qquad L_i=2h_i+1\quad(1\le i\le7),
 \qquad2\le L_0<L_1\le\cdots\le L_7,
 \qquad L_1\ge3.
\]
The even length is the unique even arm and the shortest arm of this
eight-arm core.  The equality $L_0=L_1$ is impossible by parity.

Every one- or two-arm selection is root-zero alpha, since $(2,2)$
cannot occur.  For three arms the exact sufficient predicate remains
\begin{equation}\label{eight:small-triple-guard}
 \A(I)\quad\Longleftrightarrow\quad\max_{i\in I}L_i\ge5.
\end{equation}
Here this is a statement about the specified sufficient predicate.
A triple avoiding index zero is odd and its only omitted profile is
$(3,3,3)$.  A triple containing zero has EOO parity with an even
minimum, and its only omitted profile is $(2,3,3)$.
The complete positive constructions, including minima two and four,
are in Sections~\ref{construction:all-odd-alpha3},
\ref{construction:all-eoo-alpha3}, and~\ref{app:small-alpha}.
Thus no unconditional all-triples alpha shortcut is used.
For larger alpha pieces keep Definition~\ref{eight:alpha-rule}; its
equal-four clause is available only on four odd arms of common length
at least seven.

\subsubsection{The mixed-pair predicate}\label{eight:mixed-pairs}

For any selected ordered pair $a\le b$, put
$t=\floor{a/2}$, $u=\floor{b/2}$ and
$c(a)=\ceil{a/2}-1$, $c(b)=\ceil{b/2}-1$.  Define $P_e(a,b;k)$ to
require $k\ge0$ and one of the following six base clauses:
\begin{enumerate}
\item[(M1)] $u\le k\le a-1$, and either $4k\ne4a-b-3$ or
  $a\ge6$, $3\mid a$, and $3\mid b$.
\item[(M2)] $k\le c(a)$, and either $4k\ne a-1$ or $2b\le3a-1$.
\item[(M3)] $k\le c(b)$, $k+u\le a-1$, and $4k\ne b-1$.
\item[(M4)] $a=b$ and $k<a$.
\item[(M5)] $k=a+b-2$, with $(a,b)\notin\{(2,2),(3,3)\}$.
\item[(M6)] $3\mid a$, $3\mid b$, $(a,b)\notin\{(6,6),(9,9)\}$,
  and $a+b-6\le k\le a+b-4$.
\end{enumerate}
If both selected arms are odd, also allow exactly clauses
(P5)--(P12), with their displayed guards.  None of those eight clauses
is applied to a mixed pair.

\begin{lemma}\label{eight:mixed-pair-sufficiency}
Each clause defining $P_e(a,b;k)$ supplies an offset-$k$ labeling.
\end{lemma}
\begin{proof}
Clauses (M1)--(M3) are the mixed HIGH, LOW and reversed-LOW
constructions in Section~\ref{app:pair-offsets}.  The extra guard
$a\ge6$ in the divisible-by-three repair gives shortest quotient at
least two; it is not removed at the present minimum.  Clause (M4)
is Lemma~\ref{eight:single-equal-offsets}.  Clause (M5) is the complete
maximal-pair theorem, including both small exceptions, in
Section~\ref{app:pair-offsets}.  For (M6), apply that theorem to
$(a/3,b/3)$ and then apply factor-three offset subdivision.  The
three residues give the interval
\[
 3\bigl((a+b)/3-2\bigr)+\{0,1,2\}
   =[a+b-6,a+b-4].
\]
The two excluded target pairs are exactly the forbidden quotient pairs
$(2,2)$ and $(3,3)$.  Quotient $(1,1)$ is allowed; thus this clause
also validly supplies offsets $0,1,2$ on $(3,3)$.  It does not supply
the forbidden maximal offset four there.  The remaining odd-only
clauses follow from Lemma~\ref{eight:odd-pair-sufficiency}.
\end{proof}

Use the same offset and hub recursions as before, replacing $P_o$ by
$P_e$ and replacing each one-leg insertion capacity
$\floor{L_j/2}$ by
\[
 c(L_j)=\ceil{L_j/2}-1.
\]
The shifted source offset is still $k+\floor{L_j/2}$, and the guard
$4k\ne L_j-1$ is retained.  Equal-arm and repeated-long direct clauses
are unchanged.  The ordinary $k=0$ triple rule refers to hub-zero
gracefulness, so it makes no alpha claim about either exceptional
triple.  In the four-arm hub rule the alpha-triple test
\eqref{eight:small-triple-guard} remains explicit.


\section{Further attachment and subdivision consequences}\label{eight:extensions}

\begin{corollary}\label{eight:common-short}
Let $q\ge1$, $u\ge0$, and $0\le r\le8$.  If $b_1,\ldots,b_r\ge3$
are odd, then
\[
 \Sp{qb_1,\ldots,qb_r,q^u}
\]
is graceful.  Likewise, under the hypotheses of
Article Theorem~\ref*{main-art:even-eight},
$\Sp{qa,qb_1,\ldots,qb_7,q^u}$ is graceful.
In particular every eight-arm spider whose lengths have the same
$2$-adic valuation is graceful.
\end{corollary}
\begin{proof}
Apply the respective theorem before subdivision, with $u$ unit arms,
and then use ordinary full $q$-subdivision
(Theorem~\ref{common:offset-subdivision}, at source offset zero).
Root the ordinary source labeling at whichever vertex has label zero;
that vertex need not be the hub.  The underlying subdivided spider has
the displayed lengths.  For the final assertion divide all lengths by
their common power of two to obtain eight positive odd quotients.
\end{proof}

\begin{proposition}[Four-path attachment]\label{eight:four-paths}
Let $G$ be any root-zero alpha-labeled tree of index $K$.  Adjoin four
positive-length paths $a,b,x,y$ at its root.  If $a+b=K+2$, the
resulting tree is graceful, also with any number of additional unit
paths at that root.  When $(a,b)\notin\{(2,2),(3,3)\}$, the two
paths $x,y$ may instead be replaced by any root-zero alpha-labeled
tree $T$.
\end{proposition}
\begin{proof}
Suppose first that $(a,b)\notin\{(2,2),(3,3)\}$.  The maximal-pair
theorem in Section~\ref{app:pair-offsets} gives offset $a+b-2=K$ on
$\Sp{a,b}$.  Alpha-offset amalgamation with $G$ gives a hub-zero
component.  If $(x,y)\ne(2,2)$, the two further arms form a root-zero
alpha pair, which can be ordinarily amalgamated with that component.
Append the units to the hub-zero component first.

If $(x,y)=(2,2)$, the four attached paths have a short arm, so their
whole four-arm spider is hub-zero.  Ordinarily amalgamate it with $G$,
after adjoining the units to the hub-zero piece.  The same direct
short-arm argument applies if $(a,b)$ is either maximal-offset
exception: in those cases the four-arm piece has an arm at most three.
Thus neither a nonexistent maximal offset nor the alpha-pair exception
is used.  The assertion with a general $T$ follows from the first
hub-zero component and ordinary alpha amalgamation; its stated
restriction on $(a,b)$ is retained.
\end{proof}

\begin{corollary}\label{eight:equal-four-six-core}
For every odd $L\ge7$ and all positive $a,b$ with $a+b=2L$,
$\Sp{L,L,L,L,a,b}$ is hub-zero graceful.  Adjoining arbitrary positive
arms $x,y$ and any number of unit arms gives a graceful spider.
\end{corollary}
\begin{proof}
The equal-odd-four alpha component has index $2L-2$, and the pair has
maximal offset $a+b-2=2L-2$.  The two forbidden maximal pairs have
sums four and six, below $2L$, so neither occurs.  This proves the
hub-zero claim.  The ordinary extension also follows directly from
Proposition~\ref{eight:four-paths}.
\end{proof}

\begin{proposition}[Five long paths]\label{eight:five-long-paths}
Let $G$ be any root-zero alpha-labeled tree of index $K$.  Attach five
paths with lengths $L_1,\ldots,L_5$ satisfying
\[
 L_i>2K,\qquad L_i\ne4K+1\quad(1\le i\le5).
\]
The resulting tree is graceful, also after adjoining arbitrary unit
paths at the same root.
\end{proposition}
\begin{proof}
Select a hub-zero four-arm piece among these five paths by
Section~\ref{app:five-to-four}.  The unused path attaches to $G$ as a
root-zero alpha extension by Lemma~\ref{common:alpha-extension}.
Ordinarily amalgamate it with the selected hub-zero piece, after
appending the units to that piece.  The hypotheses refer only to the
original index $K$; no successive growth or hierarchy among the five
lengths is needed.
\end{proof}

\begin{proposition}[A mixed five-arm core]\label{eight:double-five-core}
For every $r\ge2$, the spider
\[
 F_r=\Sp{2r,4r,2r+1,2r+1,2r+1}
\]
is hub-zero graceful.  Ordinarily amalgamating $F_r$ with any
root-zero alpha tree $T$, and adjoining any number of unit arms, gives
a graceful tree.  In particular, for arbitrary odd $a,b,c\ge5$,
\[
 \Sp{2r,4r,(2r+1)^3,a,b,c,1^u}
\]
is graceful for every $u\ge0$.
\end{proposition}
\begin{proof}
The literal pair $(2)\ ;\ (3,4)$ has offset one on $\Sp{1,2}$.
Subdivide by $2r$ and choose residue $r$.  The new pair has lengths
$2r,4r$ and offset $2r+r=3r$.  The equal odd alpha triple has common
length $2r+1\ge5$ and index $3r$.  Alpha-offset amalgamation yields
$F_r$.  Append the units there before ordinary amalgamation with $T$.
For the last statement use the universal odd alpha triple $\Sp{a,b,c}$.
The restriction $r\ge2$ avoids the unavailable source alpha triple
$\Sp{3,3,3}$.

For completeness the offset pair needs no path search.  Form $A_r$
by concatenating the pairs $(4r-j,8r+j+1)$ in the order
$j=0,\ldots,r-1$.  Form $B_r$ by first concatenating all pairs
$(6r-j,6r+j+1)$ in that order, and then concatenating all pairs
$(5r-j,7r+j+1)$ in that order.
These use $3r+1,\ldots,9r$.  The internal colors of $A_r$ are
$4r+1,\ldots,6r-1$ and its root color is $4r$.  The internal colors
of $B_r$ are $1,\ldots,4r-1$ and its root color is $6r$.
\end{proof}

\begin{proposition}[Three equal paths]\label{eight:three-equal-offsets}
For every $q\ge1$,
\[
 [0,7q-1]\subseteq\Om(\Sp{3q,3q,3q}).
\]
Consequently, if a root-zero alpha tree $G$ has index $K\le7q-1$,
attaching three paths of length $3q$ gives a hub-zero graceful tree.
Ordinarily amalgamating any second root-zero alpha tree $T$ gives a
graceful tree, also with arbitrary unit arms.
\end{proposition}
\begin{proof}
Here are offset-$s$ certificates on $\Sp{3,3,3}$ for each $s=0,\ldots,6$:
\[
\begin{array}{c|ccc}
s&\multicolumn{3}{c}{\text{words outward from root zero}}\\ \hline
0&(4,7,5)&(8,1,6)&(9,3,2)\\
1&(7,3,6)&(8,2,4)&(9,10,5)\\
2&(6,5,10)&(7,4,8)&(9,11,3)\\
3&(6,4,12)&(7,11,8)&(9,10,5)\\
4&(5,12,6)&(8,7,10)&(9,13,11)\\
5&(7,13,14)&(8,12,10)&(9,6,11)\\
6&(7,13,15)&(8,11,12)&(9,14,10)
\end{array}
\]
Each row has labels $s+1,\ldots,s+9$ and colors $1,\ldots,9$.
Factor-$q$ offset subdivision supplies every offset in
$[qs,q(s+1)-1]$.  The seven intervals partition $[0,7q-1]$.
The attachment and ordinary-gluing conclusions follow from the common
amalgamation lemmas.  Append units to a hub-zero component first.
\end{proof}

\begin{corollary}\label{eight:nine-family}
Suppose a nine-arm spider has three equal arms of length $3q$, with
$q\ge1$, and its other six lengths are odd integers at least five.
If three of those six have half-length sum at most $7q-1$, the spider
is graceful, also with arbitrary unit arms.  The remaining three odd
lengths have no upper bound.  For odd $q$ all nine arms are odd; for
even $q$ the three designated arms are even.
\end{corollary}
\begin{proof}
Use the bounded-index odd triple as $G$ in
Proposition~\ref{eight:three-equal-offsets} and the other odd triple as
$T$.  Both are root-zero alpha since their minima are at least five.
\end{proof}

\begin{proposition}[Upper pair bands]\label{eight:dyadic-pair-band}
Let $a,b$ be positive and let $d=2^h$ divide both, where $h\ge0$.
If $a\ne b$, then
\[
 [a+b-2d,a+b-2]\subseteq\Om(\Sp{a,b}).
\]
The same conclusion holds for $a=b$ if $a\ge4d$.
\end{proposition}
\begin{proof}
Write $m=a+b$.  For $q=1,2,4,\ldots,d$, the quotient pair is unequal
and hence avoids $(2,2),(3,3)$.  Its maximal offset is $m/q-2$.
Factor-$q$ subdivision gives
\[
 [m-2q,m-q-1].
\]
These intervals abut and their union is $[m-2d,m-2]$.
For equal arms the smallest quotient is $a/d\ge4$, so every quotient
still avoids both forbidden maximal pairs and the same proof applies.
\end{proof}

Uniform full subdivision may be applied to each ordinary graceful
family proved in this section.  This multiplies every arm length by
the subdivision factor and turns each unit arm into an arm of that
common length.  It uses ordinary subdivision, with no assertion that
alpha separation is preserved.

\section{Arithmetic certificates and reproducibility}\label{app:seven-certificates}
This section specifies the computer-assisted part of the proof. A certificate refutes a formula over arbitrary integers. It proves that one of the displayed construction routes is applicable; the mathematical construction appendices separately prove that each route produces a labeling. Neither a bounded search nor the success of a solver is substituted for these two obligations.

\subsection{Exact logical statements}


The complete mathematical statement is Article equation~\textup{(\ref*{main-art:seven-failure})}. For each of the $128$ binary vectors $p\in\{0,1\}^7$, introduce integer constants $q_0,\ldots,q_6$, substitute $L_i=2q_i+p_i$, and expand the ten sufficient alpha predicates in Article Table~\ref*{main-art:ten-table}, the four sufficient hub-zero predicates in Article Section~\ref*{main-art:seven-predicates}, and the fifteen triples in Article equation~\textup{(\ref*{main-art:fifteen-list})}. All half-length floors become $q_i$. The assertion is the conjunction of the ordered-domain conditions and failure of all fifteen routes. This finite prescription uniquely defines every assertion and contains no bound above on any variable.

The formulas are defined in \file{proof/seven/formulas.py}. The saved assertions are \file{proof/seven/inputs/P.smt2}, with \code{P} the seven-bit word. The refutations are checked for these saved assertions themselves.


The exact odd assertion is Article equation~\textup{(\ref*{main-art:odd-failure})}, using all twelve pair clauses and all stated finite recursive rules. The mixed assertion is the domain for one even shortest arm, together with failure of all stated construction routes, excluding the separately constructed profile $(2,3^7)$. Section~\ref{eight:specification} gives the complete predicate definitions and the guarded triple condition. The exact source-to-normalized-input chain is specified below.

The mixed primary input has fourteen top-level assertions. Replace them, in their original order, by the single assertion of their conjunction; declarations, terms and tokens are preserved. Satisfiability of the original list is by definition satisfiability of that conjunction. The verifier checks exact identity of this normalization. In particular no arithmetic simplification or change of quantifier domain occurs. The CPC proof is bound to the normalized single assertion. The odd input already has a single assertion.



\subsection{Checking the actual premise}
An external refutation is useful only if its premise is the intended assertion. The restricted input language here consists of integer constant declarations, ordinary arithmetic and Boolean terms, lexical SMT \code{let} expressions, and acyclic zero-argument CPC definitions. The binder \file{proof/binder.py} performs the following exact procedure.
\begin{enumerate}
\item Parse every top-level S-expression with an explicit stack. Strip comments; reject quoted symbols, strings, malformed parentheses and unsupported forms. Require distinct integer declarations and exactly one source assertion.
\item Expand SMT lexical lets using parallel binding: evaluate every right-hand side in the old environment, install all distinct names together for the body, then restore the previous environment. This excludes both sequential-let substitution and accidental capture.
\item Expand each CPC alias at its definition, using only earlier aliases. Require fresh names with no collision with declared variables. Require the CPC declaration map to equal the source declaration map and require exactly one global \code{assume} in the entire file.
\item Intern every atom and ordered application with the exact tuple consisting of its operator and complete child identifiers. The source assertion and the CPC global assumption must have the same resulting identifier. This is structural equality, not a probabilistic hash comparison. No Boolean rewriting or arithmetic simplification is permitted.
\end{enumerate}

The soundness of this comparison follows by induction on the expression DAG. An atom identifier denotes exactly that atom. An application key determines its operator and ordered list of previously fixed subterms; equal keys therefore have identical expanded expressions. Lexical expansion preserves the language semantics because every let right-hand side is evaluated in the outer scope and aliases are frozen before later local bindings can capture their variables. Thus equal final identifiers establish literal identity after only the permitted abbreviation expansion.

The binder checks premise bodies, rather than the proof calculus. The complete Logos checker must also parse the whole CPC proof, check declarations and types, verify every proof step, close every local assumption scope and accept a closed contradiction. Its global \code{assume} is distinct from a local \code{assume-push}. A malformed assumption name, an integer used as a proposition, an undeclared symbol, or an unclosed local contradiction cannot be accepted merely because the binder finds the same body. These obligations belong to the full checker, which is always required together with the binder.

\subsection{External checker and saved proof objects}
The checker is Logos, pinned to commit
\begin{center}\small\texttt{a4ab2f8deada090d96178ee9da22a47b5c834e33}.\end{center}
The bundled executable has SHA-256
\begin{center}\small\texttt{30dd5bf36cc5e4d522c5ff3f8fd2a22c8}\par
\texttt{ca3b1387393cb058f01d6897156ae0a}.\end{center}
The premise binder has SHA-256
\begin{center}\small\texttt{ba1d2722931f99954e5f83a7606747b5a}\par
\texttt{4346f03ca064a37e33c979a1248bbc9}.\end{center}
The supporting repository contains the executable and the checker source at that commit. Every canonical proof is required to return the literal verdict \code{correct} with exit status zero. None contains an explicit CPC \code{TRUST} step. The trusted base consequently includes the actual checker implementation and its compilation/runtime, the exact premise binder, and the interpretation of the displayed finite formulas. The graph constructions are mathematical proofs and exact arithmetic audits; this work does not assert that they have been encoded in a proof assistant.

The $128$ seven-arm proofs are \file{proof/seven/proofs/P.cpc}. Their total size is $121\,028\,536$ bytes; all $130$ canonical proofs together occupy $162\,974\,929$ bytes.

The odd proof object is \file{proof/eight_odd/proof.cpc}, of $18\,306\,737$ bytes, with SHA-256
\begin{center}\small\texttt{b1428c3833d075a4dc3792aaca145f4b8}\par
\texttt{02499c030407e607321f64c91105ad9}.\end{center}
The odd input is \file{proof/eight_odd/input.smt2}, with SHA-256
\begin{center}\small\texttt{808107ab5ce53e70e36a0b89e500b497}\par
\texttt{02496ae4b779832b3543b72b11690a36}.\end{center}
The mixed proof object is \file{proof/eight_even_min/proof.cpc}, of $23\,639\,656$ bytes, with SHA-256
\begin{center}\small\texttt{9bac0ecc93b014fc93821760fb85829c}\par
\texttt{b6ac31f5c814b874d2995f97502d6baa}.\end{center}
The normalized mixed input is \file{proof/eight_even_min/input.smt2}, with SHA-256
\begin{center}\small\texttt{9175765aab3210bc0f22b15c3381d8ae}\par
\texttt{f363eb15acae36ba7068a847712ba4a9}.\end{center}
The fourteen-assertion mixed source is \file{proof/eight_even_min/input_fourteen_assertions.smt2}, with SHA-256
\begin{center}\small\texttt{1f2973517759e8be4e200103798a66d1}\par
\texttt{c4cf8b694f58977fd35a1d2ea56b2d6d}.\end{center}

\subsection{Reproduction}
The GitHub repository of the article provides this supplement,
the affine construction data, and the directory \file{proof/}. That
directory contains the canonical arithmetic inputs, all $130$ CPC
proof objects, the premise binder, the checker, and the builders of
the finite tables. From \file{proof/} run
\begin{Code}
python3 verify.py --finite
python3 verify.py --arithmetic
\end{Code}
The first command regenerates both complete finite-word appendices, verifying exact required domains, labels, colors, alpha classes, controlled leaves, and the $311$ printed seven-core decompositions, and confirms that the regenerated tables are identical to those printed here. The second binds each actual CPC premise and replays all $128+1+1$ refutations. Its optional \code{--family seven}, \code{--family eight-odd}, and \code{--family eight-even-min} switches select one group. It also verifies exact identity of the fourteen mixed assertions with their normalized conjunction. Python must be run without \code{-O}.

Replay needs only Python's standard library and the bundled checker; it requires no search solver or network access. The checker executable targets x86-64 Linux with glibc. Its source is supplied for rebuilding on other systems.

\subsection{The alpha-triple guard in the even-minimum domain}
\label{eight:verification}
In one redundant four-arm hub alternative, the saved mixed input omits the alpha-triple guard of Section~\ref{eight:specification}. The only triples it adds are $(3,3,3)$ and $(2,3,3)$, and every four-arm piece containing one of them already satisfies the short-arm hub alternative. Consequently the omission does not enlarge the hub predicate, and the saved input expresses exactly the displayed predicate.

\bibliographystyle{plain}
\bibliography{references}
\authoraddresses
\end{document}
